% =============================================================================
% Books/JKSSB/notes-book.tex — JKSSB Computer Awareness · Detailed Notes
% Master file. Content is generated by scripts/build-jkssb-books.py into
% _generated/notes-content.tex from Notes/JKSSB/*-notes.tex (single source of
% truth). Compile with XeLaTeX from this directory.
% =============================================================================
\documentclass[11pt,a4paper,oneside]{report}
% =============================================================================
% Books/JKSSB/book-style.tex — branded book styling for the JKSSB Computer
% Awareness books (Detailed Notes + MCQ Bank).
%
% Derived-artifact styling only: this file never defines subject content.
% Palette: deep teal #0F5C5C + gold #C9A227. Fonts: Georgia + Segoe UI.
% Compile with XeLaTeX from inside Books/JKSSB/ (see scripts/build-jkssb-books.py).
% =============================================================================

\usepackage[a4paper,margin=2.1cm,top=2.3cm,bottom=2.4cm,headsep=10pt]{geometry}
\usepackage{microtype}

% Shared palette + convenience macros (\key, \muted, \good, \bad) + tikz.
\input{../../Preambles/header}

% -----------------------------------------------------------------------------
% Fonts (XeLaTeX system fonts)
% -----------------------------------------------------------------------------
\usepackage{fontspec}
\setmainfont{Georgia}
\setsansfont{Segoe UI}
\setmonofont{Consolas}[Scale=0.92]

% -----------------------------------------------------------------------------
% Book palette
% -----------------------------------------------------------------------------
\definecolor{bookteal}{HTML}{0F5C5C}
\definecolor{booktealdark}{HTML}{0A4444}
\definecolor{booktealmid}{HTML}{3E8E8E}
\definecolor{bookteallight}{HTML}{E4EFEF}
\definecolor{bookgold}{HTML}{C9A227}
\definecolor{bookgolddark}{HTML}{9C7C12}
\definecolor{bookgoldlight}{HTML}{F7F0D6}
\definecolor{ink}{HTML}{20262E}

\hypersetup{colorlinks=true, linkcolor=bookteal, urlcolor=bookgolddark,
            citecolor=bookteal, pdfborder={0 0 0}}

% -----------------------------------------------------------------------------
% Extra packages
% -----------------------------------------------------------------------------
\usepackage{colortbl}
\usepackage{array}
\usepackage{tcolorbox}
\tcbuselibrary{skins,breakable}
\usepackage{tocloft}
\usepackage{fancyhdr}
\usepackage{enumitem}
\usepackage{needspace}
\usepackage{caption}
\usepackage{eso-pic}

\arrayrulecolor{bookteal}
\renewcommand{\arraystretch}{1.15}
\captionsetup{font=small, labelfont={bf,color=bookteal}}

% -----------------------------------------------------------------------------
% Book metadata macros (masters set these)
% -----------------------------------------------------------------------------
\newcommand{\BTitle}{}
\newcommand{\BSubtitle}{}
\newcommand{\BShort}{}
\newcommand{\BAuthor}{}
\newcommand{\BBrand}{Concepts That Click}
\newcommand{\BBrandURL}{https://www.youtube.com/@concepts.thatclick}
\newcommand{\BWeb}{}
\newcommand{\BWebURL}{}
\newcommand{\BEdition}{}

% -----------------------------------------------------------------------------
% Section heading styling. Defined via \@startsection (NOT titlesec): titlesec
% advances the heading counter before the before-skip, so a heading that lands
% at the top of a page leaves its hyperref anchor on the PREVIOUS page, which
% breaks the table-of-contents links. \@startsection keeps the anchor on the
% heading's own page.
% -----------------------------------------------------------------------------
\makeatletter
\renewcommand{\section}{\@startsection{section}{1}{\z@}%
  {1.6ex plus .2ex}{0.7ex}{\normalfont\sffamily\large\bfseries\color{bookteal}}}
\renewcommand{\subsection}{\@startsection{subsection}{2}{\z@}%
  {1.3ex plus .2ex}{0.5ex}{\normalfont\sffamily\normalsize\bfseries\color{booktealdark}}}
\renewcommand{\subsubsection}{\@startsection{subsubsection}{3}{\z@}%
  {1.2ex plus .2ex}{0.4ex}{\normalfont\sffamily\normalsize\bfseries\color{bookgolddark}}}
\makeatother

% -----------------------------------------------------------------------------
% Example theorem (teal head). Bodies rely on this counter resetting per lesson.
% -----------------------------------------------------------------------------
\newtheoremstyle{tealexample}{6pt}{6pt}{\itshape}{}{\sffamily\bfseries\color{bookteal}}{.}{.5em}{}
\theoremstyle{tealexample}
\newtheorem{example}{Example}[section]

% Unique hyperref anchors per lesson (see \lessonheading). The section/example
% counters are reset every lesson, so their hyperref \theH... forms must be
% lesson-scoped too -- otherwise every lesson's "section.0.1" collides and the
% table-of-contents links all jump back to Lesson 01.
\providecommand{\theHsection}{}
\providecommand{\theHsubsection}{}
\providecommand{\theHsubsubsection}{}
\providecommand{\theHexample}{}

% -----------------------------------------------------------------------------
% Definition box (bodies call \begin{definitionbox}{term} ... \end{definitionbox})
% -----------------------------------------------------------------------------
\newtcolorbox{jkssbdefbox}[1]{%
  enhanced, breakable, colback=bookteallight, colframe=bookteal,
  colbacktitle=bookteal, coltitle=white, boxrule=0.6pt, arc=2pt,
  left=8pt, right=8pt, top=6pt, bottom=6pt,
  fonttitle=\sffamily\bfseries\small,
  title={Definition \textemdash\ #1}
}
\newenvironment{definitionbox}[1]
  {\begin{jkssbdefbox}{#1}\itshape}
  {\end{jkssbdefbox}}

% -----------------------------------------------------------------------------
% Table of contents styling
% -----------------------------------------------------------------------------
\renewcommand{\cftpartfont}{\sffamily\bfseries\color{bookteal}}
\renewcommand{\cftpartpagefont}{\sffamily\bfseries\color{bookteal}}
\renewcommand{\cftchapfont}{\sffamily\bfseries\color{booktealdark}}
\renewcommand{\cftchappagefont}{\sffamily\color{booktealdark}}
\renewcommand{\cftsecfont}{\color{ink}}
\renewcommand{\cftsecpagefont}{\color{bookteal}}
\setlength{\cftbeforechapskip}{4pt}
\setlength{\cftbeforepartskip}{10pt}
\setlength{\cftchapindent}{0pt}
\cftsetpnumwidth{2.8em}
\cftsetrmarg{3.4em}
\renewcommand{\cftchapfont}{\sffamily\small\bfseries\color{booktealdark}}
\renewcommand{\cftchappagefont}{\sffamily\small\color{booktealdark}}
\renewcommand{\cftpartleader}{\cftdotfill{\cftdotsep}}
\renewcommand{\cftchapleader}{\cftdotfill{\cftdotsep}}

% -----------------------------------------------------------------------------
% Running heads / feet
% -----------------------------------------------------------------------------
\pagestyle{fancy}
\fancyhf{}
\fancyhead[L]{\sffamily\footnotesize\color{bookteal}\BShort}
\fancyhead[R]{\sffamily\footnotesize\color{booktealmid}\rightmark}
\fancyfoot[C]{\sffamily\footnotesize\color{bookgolddark}\BBrand\ \textperiodcentered\ \href{\BBrandURL}{@concepts.thatclick}}
\fancyfoot[R]{\sffamily\footnotesize\color{bookteal}\thepage}
\renewcommand{\headrulewidth}{0.4pt}
\renewcommand{\footrulewidth}{0pt}
\renewcommand{\headrule}{\hbox to\headwidth{\color{bookteal!40}\leaders\hrule height \headrulewidth\hfill}}
\fancypagestyle{plain}{%
  \fancyhf{}%
  \fancyfoot[C]{\sffamily\footnotesize\color{bookgolddark}\BBrand\ \textperiodcentered\ \href{\BBrandURL}{@concepts.thatclick}}%
  \fancyfoot[R]{\sffamily\footnotesize\color{bookteal}\thepage}%
  \renewcommand{\headrulewidth}{0pt}%
}

% -----------------------------------------------------------------------------
% Cover toggle — masters call \ifshowcover\makecover\fi. The build script
% compiles a no-cover variant by injecting \def\BOOKNOCOVER{} before book-style.
% -----------------------------------------------------------------------------
\newif\ifshowcover
\showcovertrue
\ifdefined\BOOKNOCOVER\showcoverfalse\fi

% -----------------------------------------------------------------------------
% Watermark — subtle diagonal "Concepts That Click" (YouTube) on every page.
% Drawn on the BACKGROUND layer (behind all text/boxes) in a solid pale tint,
% so it never obscures reading and never relies on transparency (ExtGState).
% -----------------------------------------------------------------------------
\AddToShipoutPictureBG{%
  \begin{tikzpicture}[remember picture, overlay]
    \node[rotate=45, align=center, anchor=center]
      at (current page.center)
      {\sffamily\bfseries\color{bookteal!13}\fontsize{46}{52}\selectfont
       Concepts That Click\\[3mm]
       \fontsize{17}{21}\selectfont\color{bookgold!22}%
       YouTube\ \textperiodcentered\ @concepts.thatclick};
  \end{tikzpicture}%
}

% -----------------------------------------------------------------------------
% Cover — uses covers/<key>-cover.png if present, else a teal fallback.
% -----------------------------------------------------------------------------
\newcommand{\CoverFile}{}
\newcommand{\makecover}{%
  \clearpage
  \thispagestyle{empty}%
  \IfFileExists{\CoverFile}{%
    \begin{tikzpicture}[remember picture, overlay]
      \node[anchor=center, inner sep=0pt] at (current page.center)
        {\includegraphics[width=\paperwidth,height=\paperheight]{\CoverFile}};
    \end{tikzpicture}%
  }{%
    \begin{tikzpicture}[remember picture, overlay]
      \fill[bookteal] (current page.north west) rectangle (current page.south east);
      \fill[booktealdark] (current page.south west) rectangle ([yshift=1.4cm]current page.south east);
      \node[rotate=32, opacity=0.08, anchor=center] at (current page.center)
        {\sffamily\bfseries\fontsize{58}{58}\selectfont\color{white}Concepts That Click};
      \node[anchor=north west] at ([xshift=2.0cm,yshift=-2.2cm]current page.north west)
        {\sffamily\bfseries\color{bookgold}\large COMPUTER AWARENESS};
      \node[anchor=north west, text width=16cm] at ([xshift=2.0cm,yshift=-4.0cm]current page.north west)
        {\sffamily\bfseries\color{white}\fontsize{34}{38}\selectfont \BTitle};
      \draw[bookgold, line width=2pt] ([xshift=2.0cm,yshift=-8.4cm]current page.north west)
        -- ([xshift=9cm,yshift=-8.4cm]current page.north west);
      \node[anchor=north west, text width=15cm] at ([xshift=2.0cm,yshift=-9.0cm]current page.north west)
        {\sffamily\color{bookteallight}\Large \BSubtitle};
      \node[anchor=south west] at ([xshift=2.0cm,yshift=4.2cm]current page.south west)
        {\sffamily\color{white}\large \BAuthor};
      \node[anchor=south west] at ([xshift=2.0cm,yshift=3.0cm]current page.south west)
        {\sffamily\color{bookgold}\large \BBrand};
      \node[anchor=south west] at ([xshift=2.0cm,yshift=2.0cm]current page.south west)
        {\sffamily\color{bookteallight}\small \BWeb};
      \node[anchor=south east] at ([xshift=-2.0cm,yshift=2.0cm]current page.south east)
        {\sffamily\color{white}\small \BEdition};
    \end{tikzpicture}%
  }%
  \null
  \clearpage
}

% -----------------------------------------------------------------------------
% Unit opener — full teal page + TOC part entry; sets the running head.
% -----------------------------------------------------------------------------
\newcommand{\unitopener}[2]{%
  \clearpage
  \thispagestyle{empty}%
  \begin{tikzpicture}[remember picture, overlay]
    \fill[bookteal] (current page.north west) rectangle (current page.south east);
    \node[rotate=32, opacity=0.09, anchor=center] at (current page.center)
      {\sffamily\bfseries\fontsize{62}{62}\selectfont\color{white}Concepts That Click};
    \node[anchor=south west] at ([xshift=2.2cm,yshift=6.5cm]current page.south west)
      {\sffamily\bfseries\color{bookgold}\Large UNIT #1};
    \node[anchor=south west, text width=15cm] at ([xshift=2.2cm,yshift=4.6cm]current page.south west)
      {\sffamily\bfseries\color{white}\fontsize{30}{34}\selectfont #2};
    \draw[bookgold, line width=2pt] ([xshift=2.2cm,yshift=4.0cm]current page.south west)
      -- ([xshift=11cm,yshift=4.0cm]current page.south west);
  \end{tikzpicture}%
  \markright{Unit #1\ \textperiodcentered\ #2}%
  \phantomsection
  \addcontentsline{toc}{part}{Unit #1\ \textemdash\ #2}%
  \null
  \clearpage
}

% -----------------------------------------------------------------------------
% Lesson heading — teal banner; resets section/example counters to lesson number.
% -----------------------------------------------------------------------------
\newcommand{\lessonheading}[2]{%
  \clearpage
  \setcounter{section}{0}%
  \setcounter{example}{0}%
  \renewcommand{\thesection}{#1.\arabic{section}}%
  \renewcommand{\theHsection}{#1.\arabic{section}}%
  \renewcommand{\theexample}{#1.\arabic{example}}%
  \renewcommand{\theHexample}{#1.\arabic{example}}%
  \phantomsection
  \begin{tcolorbox}[enhanced, colback=bookteal, colframe=bookteal, coltitle=white,
      sharp corners, boxrule=0pt, left=12pt, right=12pt, top=8pt, bottom=8pt,
      borderline west={3.5pt}{0pt}{bookgold}]
    {\sffamily\bfseries\color{bookgold}\small LESSON #1}\\[2pt]
    {\sffamily\bfseries\color{white}\Large #2}%
  \end{tcolorbox}%
  \addcontentsline{toc}{chapter}{Lesson #1\ \textemdash\ #2}%
  \vspace{0.5em}%
}

% Optional video-lesson line, emitted by the build script when a URL exists.
\newcommand{\videoline}[1]{%
  {\sffamily\footnotesize\color{booktealdark}Video lesson:\ \href{#1}{\nolinkurl{#1}}\par}%
  \vspace{0.4em}%
}

% Answer-key banner (MCQ book)
\newcommand{\answersection}[1]{%
  \clearpage
  \begin{tcolorbox}[enhanced, colback=bookgoldlight, colframe=bookgold,
      sharp corners, boxrule=0pt, left=12pt, right=12pt, top=8pt, bottom=8pt,
      borderline west={3.5pt}{0pt}{booktealdark}]
    {\sffamily\bfseries\color{booktealdark}\Large Answer Key \& Rationales}\\[2pt]
    {\sffamily\color{bookgolddark}Unit #1}
  \end{tcolorbox}%
  \vspace{0.5em}%
}
\newcommand{\answersubheading}[2]{%
  \needspace{2\baselineskip}%
  \vspace{0.7em}%
  {\sffamily\bfseries\color{bookteal}Lesson #1\ \textemdash\ #2\par}%
  \vspace{0.2em}%
}

% -----------------------------------------------------------------------------
% Front matter — coloured welcome + about-the-author components.
% Sourced from the author's site (carries the "Concepts That Click" brand).
% -----------------------------------------------------------------------------
\newcommand{\fmHeading}[1]{%
  {\sffamily\bfseries\color{bookteal}\LARGE #1\par}%
  \vspace{0.15em}%
  {\color{bookgold}\rule{9.5cm}{2.2pt}}\par\vspace{0.7em}%
}

% Titled coloured panel for prose blocks.
\newtcolorbox{fmbox}[1]{%
  enhanced, breakable, colback=bookteallight, colframe=bookteal,
  colbacktitle=bookteal, coltitle=white, boxrule=0.6pt, arc=2pt,
  left=9pt, right=9pt, top=7pt, bottom=7pt,
  fonttitle=\sffamily\bfseries, title={#1}
}

% Gold call-out panel (quotes, highlights).
\newtcolorbox{goldbox}[1]{%
  enhanced, breakable, colback=bookgoldlight, colframe=bookgold,
  colbacktitle=bookgolddark, coltitle=white, boxrule=0.6pt, arc=2pt,
  left=9pt, right=9pt, top=7pt, bottom=7pt,
  fonttitle=\sffamily\bfseries, title={#1}
}

% Inline bullet for the welcome panels.
\newcommand{\fmbullet}[2]{%
  \noindent{\color{bookgold}\rule[0.15em]{0.55em}{0.55em}}\ \textbf{\color{booktealdark}#1}\ #2\par\vspace{0.35em}%
}

% Small coloured "chip" used on the about page highlight strip.
\newcommand{\brandchip}[1]{%
  \tcbox[on line, arc=2pt, colback=bookteal, colframe=bookteal, coltext=white,
         left=4pt, right=4pt, top=2pt, bottom=2pt,
         fontupper=\sffamily\bfseries\footnotesize]{#1}\hspace{0.5em}%
}


\renewcommand{\BTitle}{JKSSB Computer Awareness}
\renewcommand{\BSubtitle}{Concise Exam-Focused Notes}
\renewcommand{\BShort}{JKSSB Computer Awareness \textperiodcentered\ Concise Notes}
\renewcommand{\BAuthor}{Dr.\ Auqib Hamid Lone}
\renewcommand{\BWeb}{auqib92.github.io/claude-code-my-workflow/courses/jkssb}
\renewcommand{\BWebURL}{https://auqib92.github.io/claude-code-my-workflow/courses/jkssb/}
\renewcommand{\BEdition}{First Edition \textperiodcentered\ 2026}
\renewcommand{\CoverFile}{covers/notes-cover.png}

\begin{document}

\ifshowcover\makecover\fi

% ------------------------------- title page -------------------------------
\pagenumbering{roman}
\thispagestyle{plain}
\begin{center}
  \vspace*{2.2cm}
  {\sffamily\bfseries\color{bookteal}\Huge JKSSB Computer Awareness\par}
  \vspace{0.4cm}
  {\sffamily\color{bookgolddark}\Large Concise Notes\par}
  \vspace{0.2cm}
  {\color{bookteal}\rule{6cm}{1.5pt}}
  \vspace{0.6cm}

  {\large Key definitions, comparisons, and examples for the JKSSB Computer
  Awareness syllabus.\par}
  \vspace{1.4cm}

  {\sffamily\bfseries\large \BAuthor\par}
  \vspace{0.2cm}
  {\color{bookgolddark}\large Concepts That Click\par}
  \vspace{0.1cm}
  {\small \href{\BWebURL}{\BWeb}\par}
  \vspace{1.0cm}

  {\small First Edition \textperiodcentered\ 2026\par}
\end{center}
\clearpage

% ------------------------------ copyright ---------------------------------
\thispagestyle{plain}
\vspace*{\fill}
{\small
\noindent\textbf{JKSSB Computer Awareness \textemdash\ Concise Notes}\\
First Edition, 2026. \textcopyright\ \BAuthor.

\medskip
\noindent Prepared for the Computer Awareness component of JKSSB recruitment
examinations and the SSC computer-knowledge syllabus.

\medskip
\noindent\textbf{Source of truth.} These concise revision notes are derived from
the course lessons. The lesson materials remain authoritative.

\medskip
\noindent\textbf{Provenance of practice items.} Any question presented as a past
year question is labelled \texttt{[PYQ~\ldots]} with its paper and key; adapted
items are labelled \texttt{[PYQ-adapted~\ldots]}; generated items are labelled
\texttt{[Original]}. Generated practice is never represented as an authentic exam
question.

\medskip
\noindent\textbf{Companion volumes.} \emph{MCQ Bank with Answer Keys}
(practice sets, unit-wise) and the video course on \emph{Concepts That Click}.

\medskip
\noindent\href{\BWebURL}{\BWeb}
}
\vspace*{\fill}
\clearpage

% ----------------------------- introduction -------------------------------
\markright{Welcome}
\phantomsection
\addcontentsline{toc}{chapter}{Welcome}

\fmHeading{Welcome}

\noindent\emph{Unit-wise revision notes for JKSSB Computer Awareness --- from
computer fundamentals and hardware to office applications, networking, and
cybersecurity.}

\bigskip
This volume gathers the complete JKSSB Computer Awareness course into a single,
distributable revision companion. Each lesson distils its key definitions,
comparisons, tables, and selected worked examples into one scannable pass, kept
1:1 with the course slides.

\vspace{0.6em}
\begin{fmbox}{What this volume contains}
\fmbullet{Eleven units:}{Fundamentals \& data representation, CPU \& hardware,
memory \& storage, I/O \& peripherals, software \& operating systems, programming,
office applications, the internet, networking, cybersecurity, and e-governance.}
\fmbullet{One lesson per topic:}{Definitions boxed for quick recall, alongside
comparison tables and worked examples.}
\fmbullet{A matching MCQ set:}{Practise recall, comparison, and one-step
application with the companion \emph{MCQ Bank with Answer Keys}.}
\end{fmbox}

\vspace{0.7em}
\begin{goldbox}{How to revise}
Read a lesson, cover the boxed definitions, then recall them from memory. Use the
companion MCQ set to test the same lesson under exam conditions.
\end{goldbox}

\clearpage
% =============================================================================
% Books/JKSSB/about-author.tex — shared "About the Author" front matter.
% Sourced from the author's site homepage:
%   https://auqib92.github.io/claude-code-my-workflow/index.html
% Included by both master files (notes-book.tex, mcq-book.tex).
% =============================================================================

\fmHeading{About the Author}

\begin{fmbox}{Dr.\ Auqib Hamid Lone \textemdash\ Assistant Professor \& Blockchain Researcher}
\noindent Dr.\ Auqib Hamid Lone is an \textbf{Assistant Professor of Computer
Science and Engineering} at GCET Kashmir, teaching core undergraduate courses to
B.Tech CSE students. His PhD, at NIT Srinagar, explored blockchain technology
beyond cryptocurrencies --- harnessing it to remove problems of trust and security
in real applications.

\medskip
Before academia he was a \textbf{Research Scientist in blockchain and
cryptography at Parity Financial (London)}, where he worked on custody
infrastructure, transaction construction, and zero-knowledge privacy layers. He
publishes his teaching materials openly so students and fellow instructors can
reuse, adapt, and improve them.
\end{fmbox}

\vspace{0.7em}
\begin{goldbox}{At a glance}
\noindent\brandchip{GCET Kashmir}\brandchip{PhD, NIT Srinagar}\brandchip{Blockchain \& Crypto}\\[0.5em]
\brandchip{GATE 98.5\%ile}\brandchip{AIR 31 UGC NET-JRF}\brandchip{6+ Publications}\brandchip{3 yrs Industry R\&D, London}
\end{goldbox}

\vspace{0.7em}
\begin{fmbox}{Field notes}
\fmbullet{Role:}{Assistant Professor, Computer Science \& Engineering}
\fmbullet{Institution:}{Government College of Engineering \& Technology (GCET), Kashmir}
\fmbullet{Focus:}{Blockchain, applied cryptography, and teaching computing from first principles}
\fmbullet{Education:}{PhD, NIT Srinagar}
\fmbullet{Office hours:}{Mon--Sat, 2--3 PM}
\end{fmbox}

\vspace{0.7em}
\begin{tcolorbox}[enhanced, breakable, colback=bookteal, colframe=booktealdark,
    boxrule=0.6pt, arc=2pt, left=10pt, right=10pt, top=9pt, bottom=9pt]
{\itshape\color{white}\large ``Security is the balanced mix of feeling and reality.''\par}
\vspace{0.35em}
{\sffamily\bfseries\color{bookgold}\small --- Dr.\ Auqib Hamid Lone\par}
\end{tcolorbox}

\vspace{1.0em}
\begin{center}
{\sffamily\bfseries\color{booktealdark}Concepts That Click}\quad
{\color{bookgold}\textbullet}\quad
\href{mailto:auqib.cse@gcetkashmir.ac.in}{auqib.cse@gcetkashmir.ac.in}\\[0.3em]
{\small\color{bookteal}%
\href{https://auqib92.github.io/claude-code-my-workflow/}{Course materials}\quad\textbullet\quad
\href{https://github.com/AUQIB92}{GitHub}\quad\textbullet\quad
\href{https://www.linkedin.com/in/dr-auqib-hamid-lone-45996067}{LinkedIn}\quad\textbullet\quad
\href{https://scholar.google.com/citations?user=Dr.-Auqib-Hamid-Lone}{Google Scholar}}
\end{center}

\clearpage

\begingroup
\setlength{\parskip}{0pt}
\tableofcontents
\endgroup
\clearpage

\pagenumbering{arabic}

\unitopener{1}{Fundamentals \& Data Representation}
\lessonheading{01}{What Is a Computer?}
\videoline{https://youtu.be/EdRZK57oqpo}
\section{A computer is a data-processing system}
Start with an office task. A clerk enters marks into a spreadsheet, a formula
adds the values and calculates an average, the result appears on screen, and
the workbook can be saved. A computer is not merely the screen or the box on
the desk. The task depends on data, instructions, processing, an output path,
and possibly storage.

\begin{definitionbox}{Computer}
An electronic, programmable system that accepts data, processes it under
instructions, produces information, and can store data or results.
\end{definitionbox}

Each part of the definition matters. \emph{Electronic} describes the
technology used by the system. \emph{Programmable} means that instructions
specify the operations to perform. The system accepts inputs, applies those
instructions, and makes results available. It can also retain data or results
for later use.

\begin{definitionbox}{Data}
Raw facts or symbols supplied to a system before they have been organized for
the particular question being asked.
\end{definitionbox}

\begin{definitionbox}{Information}
Data that have been organized or processed so that they are useful for a
particular purpose.
\end{definitionbox}

\begin{definitionbox}{Instruction}
A rule that tells the computer what operation to perform or what action to
take.
\end{definitionbox}

\begin{definitionbox}{Program}
An organized set of instructions that directs a computer through a task.
\end{definitionbox}

Data and information are defined by their role in a task, not by whether the
values are numbers or text. Individual marks are data when the question is
``What are the three scores?'' Their average is information when the question
is ``What was the student's average?'' The average can itself become input
data for a later task, so the same value may play different roles in
different contexts.

\begin{example}
A clerk enters $68$, $74$, and $81$. These three marks are the input data.
Adding them gives $223$; dividing by three gives
$223/3=74.333\ldots$, or approximately $74.33$ to two decimal places. The
total and average are processed results. They are useful information if the
task is to summarize performance. Saving the workbook retains the records.
\end{example}

\begin{example}
Suppose the same three marks are each out of $100$. The total possible score
is therefore $300$. The percentage is
\[
  \frac{68+74+81}{300}\times 100
  =\frac{223}{300}\times 100
  =74.333\ldots\%.
\]
Rounded to two decimal places, the percentage is $74.33\%$. The total,
maximum possible score, formula, and rounding rule all matter; the computer
cannot infer a missing maximum score from the marks alone.
\end{example}

\section{Input, processing, output, and storage}
\begin{definitionbox}{Input}
Data and instructions that enter a computer system through a device, file,
network, or software interface.
\end{definitionbox}

\begin{definitionbox}{Processing}
Operations carried out on input data according to instructions, usually by
the processor and program involved in the task.
\end{definitionbox}

\begin{definitionbox}{Output}
The result made available by a computer system, for example on a display, in
a printed document, as sound, or through another system.
\end{definitionbox}

\begin{definitionbox}{Storage}
The retention of data, instructions, or results for current or later use.
\end{definitionbox}

The basic three-stage model is IPO: Input, Processing, Output. Many
introductory diagrams add Storage, giving IPOS. Storage is a distinct
function, not another name for processing. The stages describe roles in a
task; they do not require every computer job to follow a single rigid,
one-direction path. A saved result can be read back as input to a later
operation.

\begin{center}
\begin{tabular}{p{0.17\textwidth}p{0.34\textwidth}p{0.37\textwidth}}
\toprule
\textbf{Stage} & \textbf{Question to ask} & \textbf{Marks-spreadsheet example}\\
\midrule
Input & What data or instructions enter? & The three marks and the average formula\\
Processing & What operation changes or organizes the data? & Addition, division, and percentage calculation\\
Output & What result is presented, and how? & The total or percentage shown on screen or printed\\
Storage & What is retained for later use? & The saved workbook and its records\\
\bottomrule
\end{tabular}
\end{center}

The physical item does not determine the stage by itself. A monitor is
normally an output device, but the classification question concerns the
monitor's role in the described task. A saved file is stored data; opening
that same file later makes its contents input to the new task. Likewise,
information describes the usefulness of processed data, whereas output
describes how the result is made available. A screen can present useful
information, but the terms answer different questions.

\begin{example}
In the marks spreadsheet, typing or importing the scores is input. The
formula performs processing. The displayed total and average are output.
Saving the workbook is storage. If the saved workbook is opened tomorrow,
reading it is input to that later session.
\end{example}

\section{Characteristics and limitations}
\begin{definitionbox}{Speed}
The ability to carry out many operations in a short time.
\end{definitionbox}

\begin{definitionbox}{Accuracy}
The ability to produce consistent, precise results when the data, instructions,
and operation are correct.
\end{definitionbox}

\begin{definitionbox}{Diligence}
The ability to repeat work without tiring in the human sense, provided the
system continues to receive power and function correctly.
\end{definitionbox}

\begin{definitionbox}{Versatility}
The ability to perform different kinds of tasks by running different
programs.
\end{definitionbox}

\begin{definitionbox}{Automation}
The ability to continue a programmed operation after it has been started and
the required inputs or conditions have been supplied.
\end{definitionbox}

\begin{definitionbox}{Storage as a characteristic}
The ability to retain data, instructions, and results for later access.
\end{definitionbox}

These characteristics describe capabilities, not a guarantee that every
result is correct or useful. Speed concerns how quickly operations are
performed; accuracy concerns whether the operation is correctly carried out
on the supplied data. Storage concerns retaining material; it does not
validate that material. Versatility comes from using different programs,
while automation means a started procedure can continue according to its
instructions and conditions.

Speed and accuracy must not be confused with judgement. A computer can
calculate quickly and repeat the same operation consistently, but it cannot
infer that a deliberately or accidentally incorrect mark is false merely
because the value is stored in a spreadsheet. Accuracy is conditional on the
quality of data and instructions.

\begin{definitionbox}{GIGO}
``Garbage In, Garbage Out'': unsuitable or incorrect input can lead to
unsuitable or incorrect output, even when the computer follows its
instructions.
\end{definitionbox}

\begin{example}
Suppose the correct marks are $68$, $74$, and $81$, but a user enters $18$
instead of $81$. The formula may correctly compute
$68+74+18=160$ and $160/3\approx 53.33$. The arithmetic is consistent with
the entered values, but the output is not a valid summary of the actual
marks. The error entered at the input stage was not corrected by processing.
\end{example}

The example illustrates the limitation often summarized by GIGO. If the
input is wrong, a correct formula can still produce an incorrect result. If
the input is correct but the formula is wrong, the result can also be wrong.
The machine may execute either set of instructions consistently; it does not
automatically decide that the source data are truthful, complete, or fair.
Computers do not possess human judgement or common-sense understanding by
default, so a person must interpret output in context and decide what action
it justifies.

\section{Applications and role classification}
The same IPO/IPOS model applies across fields. What changes is the type of
data and the useful result; the roles of input, processing, output, and
storage remain identifiable.

\begin{center}
\begin{tabular}{p{0.15\textwidth}p{0.24\textwidth}p{0.27\textwidth}p{0.22\textwidth}}
\toprule
\textbf{Field} & \textbf{Input} & \textbf{Processing / output} & \textbf{Storage}\\
\midrule
Office & Text, figures, and instructions & Create a document or calculate a total & Document or workbook file\\
Banking & Account and transaction details & Validate and process a transaction; show a status & Account and transaction records\\
Education & Learner responses and assessment rules & Score work and present results or learning resources & Assessment and learner records\\
Healthcare & Patient details and appointment requests & Organize records, schedules, or diagnostic support & Patient and appointment records\\
Governance & Online form fields and supporting details & Process a request and provide a citizen service & Submitted forms and service records\\
\bottomrule
\end{tabular}
\end{center}

For exam questions, classify the role within the described task. A keyboard
is an input device; the formula is processing; a monitor presents output; a
saved file is stored data. A physical object may support more than one
stage, but the question usually asks for the role it is playing in that
particular scenario.

\begin{example}
An online expense form accepts entered amounts, checks and sums them using
the programmed rules, displays the total, and stores the submitted record.
If one amount is mistyped, the calculation may still be executed correctly
on the wrong data. The user should correct the input before relying on the
total.
\end{example}

When a question gives a complete workflow, mark the stage of each action
rather than assigning one label to the entire system. Entering a value is
input; checking or summing it is processing; showing a result is output;
keeping the transaction record is storage. If the stored record is opened
and processed again, its role changes to input for that later operation.

\subsection*{Exercises}
\begin{enumerate}[label=\arabic*.]
  \item A user enters four prices, applies a tax formula, prints the invoice,
    and saves the file. Identify the input, processing, output, and storage.
  \item Explain why a fast, consistent computer can still produce an
    incorrect result.
  \item Three temperature readings are $18$, $20$, and $22$ degrees. For a
    task asking for their average, identify the data and the resulting
    information.
  \item Give one example from banking or governance and trace it through
    IPOS.
\end{enumerate}

\subsection*{Solutions}
\begingroup\small
\begin{enumerate}[label=\arabic*.]
  \item The four prices are input; applying the tax formula is processing;
    the printed invoice is output; saving the invoice file is storage.
  \item The computer follows its instructions and operates on the supplied
    data. Incorrect input or a wrong formula can therefore produce an
    incorrect result consistently; speed does not establish truth.
  \item The three readings are input data. Their average is
    $(18+20+22)/3=20$ degrees and is information for the average-reading
    task.
  \item Example: an online fee form accepts the user's details and amount
    (input), calculates the fee (processing), shows a receipt (output), and
    saves the transaction record (storage).
\end{enumerate}
\endgroup

\lessonheading{02}{Types of Computers}
\videoline{https://youtu.be/d6zL3YS69kM}
\section{Choose the classification axis first}
Computer labels are useful only when the basis of classification is clear.
``Digital'' describes how data are represented; ``microcomputer'' describes
a traditional size/use class; ``laptop'' describes a form factor. These
labels do not compete. One laptop can correctly be described using several
axes at once.

\begin{center}
\begin{tabular}{p{0.22\textwidth}p{0.34\textwidth}p{0.34\textwidth}}
\toprule
\textbf{Axis} & \textbf{Question} & \textbf{Clue / example}\\
\midrule
Signal representation & How are values represented? & Continuous signal: analog; distinct states: digital\\
Scale and workload & What kind of work or volume is handled? & Many transactions: mainframe; very large calculations: supercomputer\\
Purpose & How broad is the intended task range? & Many different programs: general-purpose; narrow task: special-purpose\\
Integration & Where is the computer installed? & Built into a larger product: embedded\\
Form factor & What physical form does it take? & Desktop, laptop, tablet, or smartphone\\
\bottomrule
\end{tabular}
\end{center}

\begin{definitionbox}{Classification axis}
The feature used to group computers, such as signal representation,
capability/workload, purpose, or physical form.
\end{definitionbox}

Before choosing an answer, identify the question being asked:
\begin{enumerate}
  \item Is the stem about continuous or discrete signals?
  \item Is it about personal use, transaction volume, or very large
    calculations?
  \item Is it about a computer's general purpose or its built-in role in a
    larger product?
  \item Is it naming a shape, such as desktop, laptop, or tablet?
\end{enumerate}
Selecting a familiar label without checking the axis is a common source of
wrong answers.

\section{Classification by signal representation}
\begin{definitionbox}{Analog computer}
A computer that represents or processes quantities as continuously varying
signals over a range.
\end{definitionbox}

\begin{definitionbox}{Digital computer}
A computer that represents or processes data as discrete values, commonly
using binary digits 0 and 1.
\end{definitionbox}

\begin{definitionbox}{Hybrid computer}
A system combining analog measurement or input with digital processing,
control, or display.
\end{definitionbox}

An analog quantity can take intermediate values within a range. For
example, a moving pointer can indicate many positions along a speed range,
not just a fixed list of allowed readings. A digital system represents
information using distinct states; modern computers commonly use binary
states, written as 0 and 1. The important contrast is continuous versus
discrete representation, not old versus new or mechanical versus electronic.

\begin{center}
\begin{tabular}{p{0.18\textwidth}p{0.35\textwidth}p{0.35\textwidth}}
\toprule
\textbf{Class} & \textbf{Representation} & \textbf{Recognition clue}\\
\midrule
Analog & Continuously varying quantity & A value can lie anywhere within a range\\
Digital & Separate symbols or states & A value is encoded using distinct states, commonly 0 and 1\\
Hybrid & Uses both forms at different stages & Measurement or input is analog; processing or display is digital\\
\bottomrule
\end{tabular}
\end{center}

A hybrid arrangement can sense a physical signal and then use digital
circuitry to process and present the measurement. Its classification
describes the complete path, not just the display or one internal component.

\begin{example}
A digital patient monitor may measure a changing physical signal through a
sensor. The sensor/front end handles an analog quantity; conversion allows
digital processing; the screen shows a numeric result. When a question asks
about the complete measurement-and-processing arrangement, ``hybrid'' is
the useful classification. The monitor's display alone does not determine
the signal class of its entire system.
\end{example}

\section{Traditional size and capability classes}
\begin{definitionbox}{Microcomputer}
A personal-use computer built around a microprocessor; common examples
include desktop computers, laptops, tablets, and smartphones.
\end{definitionbox}

\begin{definitionbox}{Minicomputer}
A historical term for a mid-range computer class, often described as
supporting multiple users; it is less common in current product marketing.
\end{definitionbox}

\begin{definitionbox}{Mainframe}
A computer class associated with reliable, high-volume input/output and
large numbers of concurrent users or transactions.
\end{definitionbox}

\begin{definitionbox}{Supercomputer}
A high-performance computer designed for very large computational
workloads, often using parallel processing.
\end{definitionbox}

\begin{center}
\begin{tabular}{p{0.2\textwidth}p{0.34\textwidth}p{0.34\textwidth}}
\toprule
\textbf{Class} & \textbf{Primary clue} & \textbf{Example / caution}\\
\midrule
Microcomputer & Personal-use computer built around a microprocessor & Desktop, laptop, tablet, or smartphone; not defined only by its outer size\\
Minicomputer & Historical mid-range, multi-user class & A traditional category; less common in current product marketing\\
Mainframe & High transaction and input/output throughput for many users & Large banking or account-processing workload\\
Supercomputer & Very high computational performance, often parallel & Large climate or engineering simulation\\
\bottomrule
\end{tabular}
\end{center}

Mainframes and supercomputers may both be powerful. The distinction in an
exam question is usually the workload. A mainframe is associated with
throughput, reliability, and many transactions or users. A supercomputer is
associated with computationally intensive scientific or engineering
problems. ``Throughput'' here means the amount of work completed over time;
the mainframe clue is handling many transactions reliably, rather than
solving one unusually large numerical problem. Parallel processing divides
computational work so multiple processing units can perform parts of it at
the same time. Physical size alone is not a dependable classification rule.

\begin{example}
A bank processes a large number of account transactions for many customers.
The emphasis is concurrent transaction throughput and service reliability,
so the mainframe class is the best match. A research centre running a
large-scale climate simulation emphasizes high-performance computation, so
the supercomputer class is the better match.
\end{example}

\begin{example}
A question describes a ``mid-range, multi-user computer'' and asks for its
traditional category. The intended answer is minicomputer. Do not infer that
the label is common in modern consumer advertisements; it is a historical
classification term.
\end{example}

\section{Classification by purpose and form}
\begin{definitionbox}{General-purpose computer}
A computer capable of performing many different tasks by running different
programs.
\end{definitionbox}

\begin{definitionbox}{Special-purpose computer}
A computer designed primarily for a narrower task or a limited set of
related tasks.
\end{definitionbox}

\begin{definitionbox}{Embedded computer}
A computer built into a larger device to control or support that device's
function.
\end{definitionbox}

\begin{center}
\begin{tabular}{p{0.2\textwidth}p{0.37\textwidth}p{0.3\textwidth}}
\toprule
\textbf{Term} & \textbf{What it describes} & \textbf{Example}\\
\midrule
General-purpose & Broad range of tasks through different programs & Personal computer used for documents, browsing, and media\\
Special-purpose & Narrow task or small set of related tasks & Controller dedicated to an appliance function\\
Embedded & Integration into a larger product & Washing-machine controller\\
Form factor & Physical arrangement or shape & Desktop, laptop, tablet, smartphone\\
\bottomrule
\end{tabular}
\end{center}

An embedded computer is commonly special-purpose, but these terms name
different aspects: ``embedded'' describes where the computer is integrated;
``special-purpose'' describes the range of tasks it is intended to perform.
A desktop computer that runs office, browsing, and media programs is a
general-purpose system. A washing-machine controller is built into a larger
appliance and handles a limited control task. A computer can therefore be
both embedded and special-purpose; the terms are compatible, not competing
answers. In contrast, ``laptop'' says how a computer is physically arranged
and does not by itself tell whether the data representation is analog or
digital.

\begin{example}
A smartphone can be classified as digital by signal representation, a
microcomputer by personal-use class, and a mobile device by form factor.
These labels can all be true at the same time because they answer different
questions.
\end{example}

\begin{example}
A washing machine's controller is embedded in the appliance and is
special-purpose because it manages a limited set of appliance functions.
Calling it a mainframe because it controls a whole machine confuses physical
integration with workload class.
\end{example}

\section{A reliable matching procedure}
When an exam stem describes a device, underline the clue that indicates the
axis. ``Continuous'' points to analog; ``discrete'' or binary points to
digital; ``sensor plus digital processing'' may point to hybrid. ``Many
concurrent transactions'' suggests mainframe. ``Large scientific
calculations'' suggests supercomputer. ``Built into an appliance'' points to
embedded. ``Desktop'' or ``laptop'' identifies form factor, not signal type.

\begin{example}
Classify a laptop used by an accounts clerk. By representation it is digital;
by traditional size/use class it is a microcomputer; by form it is a laptop;
and by purpose it is general-purpose if it runs many kinds of programs.
None of these labels invalidates the others.
\end{example}

\subsection*{Computer type and generation}
A computer's type may describe its signal class, scale, purpose, or form.
Generation instead names a historical technology grouping. Keep these axes
separate: integrated circuits are associated with the third generation,
while microprocessors are associated with the fourth.

This distinction prevents a common category error. A laptop can be called a
digital microcomputer and can belong to a particular historical generation
based on the technology used in its development. ``Laptop'' identifies its
form, ``microcomputer'' its traditional user/capability class, ``digital''
its representation, and ``generation'' its place in a historical grouping.
Each answer must match the axis named in the question.

\begin{example}
For a washing-machine controller, ``embedded'' answers where the computer is
installed. ``Special-purpose'' answers the range of work it is designed to
perform. The controller can correctly have both labels at once. If a
question asks how it represents data, that is a separate signal
classification and must be answered from that clue.
\end{example}

\subsection*{Exercises}
\begin{enumerate}[label=\arabic*.]
  \item Classify a smartphone on the signal, traditional class, and form
    axes.
  \item A system is described as handling many simultaneous banking
    transactions reliably. Which class is the strongest match, and why?
  \item Distinguish an embedded computer from a special-purpose computer.
  \item Why is ``laptop'' not an answer to a question asking whether a
    computer is analog or digital?
  \item A digital instrument measures a continuously changing physical
    quantity and processes the result digitally. Explain when ``hybrid'' is
    an appropriate label.
\end{enumerate}

\subsection*{Solutions}
\begingroup\small
\begin{enumerate}[label=\arabic*.]
  \item It is digital by representation, a microcomputer by personal-use
    class, and a smartphone by form factor.
  \item Mainframe. The clues are many simultaneous transactions and
    reliability/throughput, rather than a single extremely large calculation.
  \item Embedded describes integration into a larger product. Special-purpose
    describes a narrow intended task. An embedded controller is often
    special-purpose, but the terms describe different properties.
  \item Laptop names physical form. Analog/digital names the representation
    of data, so the question asks about a different classification axis.
  \item It is appropriate when the complete arrangement combines analog
    sensing or measurement with digital conversion and processing.
\end{enumerate}
\endgroup

\lessonheading{03}{Bits, Bytes \& Units}
\videoline{https://youtu.be/j_62CGCBElo}
\section{The basic units of digital data}
Digital information is represented using bits. Groups of bits form larger
units, but the terms do not all describe the same kind of quantity. A bit is
a single binary digit; a word is an architecture-dependent processor unit;
a character is an encoded symbol.

\begin{definitionbox}{Bit}
A binary digit that has one of two values, 0 or 1; the smallest unit of data
in the usual computer-fundamentals description.
\end{definitionbox}

\begin{definitionbox}{Nibble}
A group of four bits.
\end{definitionbox}

\begin{definitionbox}{Byte}
A group of eight bits, equal to two nibbles.
\end{definitionbox}

\begin{definitionbox}{Word}
The natural data size a processor handles as a unit; its size depends on the
processor architecture.
\end{definitionbox}

\begin{definitionbox}{Character}
A text symbol represented in digital storage by a character encoding; its
storage size is not necessarily one byte.
\end{definitionbox}

\begin{center}
\begin{tabular}{p{0.19\textwidth}p{0.2\textwidth}p{0.48\textwidth}}
\toprule
\textbf{Unit} & \textbf{Size / basis} & \textbf{Meaning}\\
\midrule
Bit & 1 binary digit & One of the two values 0 or 1\\
Nibble & 4 bits & Half of an 8-bit byte\\
Byte & 8 bits & Two nibbles; common unit for file/storage amounts\\
Word & Architecture-dependent & Natural processor data size, such as 32 or 64 bits\\
Character & Encoding-dependent & A text symbol; its encoded storage may take more than one byte\\
\bottomrule
\end{tabular}
\end{center}

The bit-to-byte relationships are fixed in this lesson:
\[
  1\text{ nibble}=4\text{ bits},\qquad
  1\text{ byte}=8\text{ bits}=2\text{ nibbles}.
\]
The word size is not defined by this ladder. A processor may work with a
natural word size such as 32 or 64 bits, but the term depends on the
architecture and should not be treated as one universal number.

\begin{example}
Write the byte \texttt{10110110} as two nibbles: \texttt{1011} and
\texttt{0110}. Each group contains four bits; together the two groups still
contain eight bits, so the original value is one byte.
\end{example}

\section{Bits and bytes in names and data rates}
Capitalization changes the unit. Lowercase \texttt{b} denotes a bit, while
uppercase \texttt{B} denotes a byte. A byte contains eight bits. Therefore a
rate reported in megabits per second is not numerically equal to a rate in
megabytes per second.

\begin{definitionbox}{Transfer rate}
The amount of data transferred in a stated time interval, commonly written
in bits per second or bytes per second.
\end{definitionbox}

\begin{definitionbox}{Capacity}
The amount of data that a storage medium can hold.
\end{definitionbox}

The slash in a rate such as Mbps means ``per second.'' It is not part of the
data amount itself. Keep the numerical amount, whether it is bits or bytes,
and the time unit visible while calculating. Under matching decimal prefixes,
\[
  1\text{ byte}=8\text{ bits},\qquad
  1\text{ MB}=8\text{ Mb}.
\]
Thus dividing a rate in Mb/s by eight gives MB/s, while multiplying a file
size in MB by eight gives Mb. These conversions change the unit, not the
underlying amount of data.

To convert a bit rate to a byte rate, divide by eight, provided the prefixes
match and the question is asking for the ideal conversion. Protocol overhead,
congestion, and other factors can reduce actual network throughput, so a
nominal rate is not a promise of application-level speed.

\begin{example}
A link is rated at $8$ Mbps. Using matching decimal prefixes and ignoring
overhead, convert bits to bytes:
\[
8\text{ megabits/s}\div 8=1\text{ megabyte/s}.
\]
This is a transfer rate, not a storage capacity.
\end{example}

\begin{example}
Convert an ideal rate of $24$ Mbps to megabytes per second. The prefixes
already match, so only the bit-to-byte conversion is needed:
\[
  24\text{ Mb/s}\div 8=3\text{ MB/s}.
\]
The answer is a rate because the time denominator remains. If the question
asked how long a file takes to transfer, the file amount would first need to
be expressed in the same unit as the rate, then divided by that rate.
\end{example}

\begin{example}
A $4$ MB file contains $32$ megabits under matching decimal prefixes because
$4\text{ MB}\times 8=32\text{ Mb}$. If transferred over an ideal $20$ Mbps
link, the simple lower-bound calculation is $32/20=1.6$ seconds. Actual
transfer can take longer because the ideal calculation ignores protocol
overhead and other delays.
\end{example}

\section{Capacity units and the 1024 convention}
For the conventional JKSSB-style ladder used in this lesson, each step is a
factor of 1024:
\[
1\text{ KB}=1024\text{ bytes},\quad
1\text{ MB}=1024\text{ KB},\quad
1\text{ GB}=1024\text{ MB},\quad
1\text{ TB}=1024\text{ GB}.
\]
Thus $1$ MB under this convention is $1024\times1024=1{,}048{,}576$ bytes.
The exam-style abbreviations KB, MB, GB, and TB are often used with this
traditional binary ladder. Modern technical standards distinguish decimal
SI units (powers of 1000) from binary IEC units such as KiB (powers of 1024).
If a question specifies decimal units or the IEC prefixes, use those
definitions instead. Do not silently mix conventions.

\begin{center}
\begin{tabular}{p{0.18\textwidth}p{0.25\textwidth}p{0.43\textwidth}}
\toprule
\textbf{Unit} & \textbf{Previous unit} & \textbf{Bytes under repeated 1024 convention}\\
\midrule
1 KB & 1024 bytes & $1{,}024$ bytes\\
1 MB & 1024 KB & $1{,}048{,}576$ bytes\\
1 GB & 1024 MB & $1{,}073{,}741{,}824$ bytes\\
1 TB & 1024 GB & $1{,}099{,}511{,}627{,}776$ bytes\\
\bottomrule
\end{tabular}
\end{center}

To move from a larger unit to the next smaller unit, multiply by 1024.
Repeat once for every step crossed. To move from a smaller unit to the next
larger unit, divide by 1024. Writing the units beside the numbers helps
prevent using the wrong direction.

\begin{example}
Convert $3$ KB to bytes under the 1024 convention:
\[
3\text{ KB}\times 1024\text{ bytes/KB}=3072\text{ bytes}.
\]
The KB units cancel, leaving bytes.
\end{example}

\begin{example}
Convert $2$ MB to KB under the same convention:
\[
2\text{ MB}\times 1024\text{ KB/MB}=2048\text{ KB}.
\]
This is one step down the capacity ladder, so multiply by 1024.
\end{example}

\begin{example}
Convert $5$ GB to MB using the same convention. GB is one step larger than
MB, so multiply by 1024:
\[
  5\text{ GB}\times 1024\text{ MB/GB}=5120\text{ MB}.
\]
The units GB cancel, leaving MB. To convert $5120$ MB back to GB, divide by
1024 and recover $5$ GB.
\end{example}

\section{Keep capacity separate from rate}
Capacity answers ``how much can be stored?'' Transfer rate answers ``how
much is moved per unit of time?'' MB is a byte-based amount; Mbps is a
bit-based rate. The lowercase/uppercase distinction and the \texttt{/s}
time component carry information. A drive with a larger capacity does not
necessarily transfer data faster than a smaller drive.

\begin{center}
\begin{tabular}{p{0.2\textwidth}p{0.31\textwidth}p{0.31\textwidth}}
\toprule
\textbf{Question} & \textbf{Capacity} & \textbf{Transfer rate}\\
\midrule
What does it measure? & Amount that can be held & Amount moved per unit of time\\
Common example & A 4 MB file or 1 TB drive & A 20 Mbps connection\\
Unit clue & No time denominator & Includes ``per second'' or \texttt{/s}\\
\bottomrule
\end{tabular}
\end{center}

\begin{example}
In the lesson's opening comparison, $4$ MB describes the file's storage
amount, while $20$ Mbps describes the network's transfer rate. The first
has no time denominator; the second is measured per second, so they are
different kinds of quantities.
\end{example}

Before calculating, perform a quick unit check: if the final question asks
for bytes, convert bits to bytes; if it asks for seconds, divide an amount
of data by a compatible rate; if it asks for a storage amount, do not attach
a time unit. A numerically correct calculation with mismatched units is not
a correct answer.

\subsection*{Exercises}
\begin{enumerate}[label=\arabic*.]
  \item How many bits are in $7$ bytes? Show the conversion.
  \item How many bytes are in $5$ KB under the 1024 convention?
  \item Convert $24$ Mbps to MB/s using matching decimal prefixes and
    ignoring overhead.
  \item A device stores $4$ GB. Is this a capacity or a transfer rate? What
    unit clue supports your answer?
  \item A processor has a 64-bit word size. How many bytes are in one word?
\end{enumerate}

\subsection*{Solutions}
\begingroup\small
\begin{enumerate}[label=\arabic*.]
  \item $7\times 8=56$ bits.
  \item $5\times 1024=5120$ bytes.
  \item $24/8=3$ MB/s before protocol overhead.
  \item It is capacity: GB is a data amount, and there is no ``per second''
    time component.
  \item $64/8=8$ bytes per word.
\end{enumerate}
\endgroup

\lessonheading{49}{Computer Generations \& History}
\section{What a computer generation represents}
Computer generations are conventional historical groupings based on the
dominant hardware technology or design development associated with a period.
They help organize the major changes in computer hardware, programming, and
use. They are not exact universal year boundaries: a date range can differ
between accounts, so the technology association is the safer cue for an
objective question.

\begin{definitionbox}{Computer generation}
A conventional historical grouping of computers associated with a dominant
technology or design development.
\end{definitionbox}

\begin{center}
\begin{tabular}{p{0.22\textwidth}p{0.68\textwidth}}
\toprule
\textbf{Term} & \textbf{Meaning in this lesson}\\
\midrule
Vacuum tube & An early electronic component used in first-generation computers; these systems were large and produced substantial heat.\\
Transistor & A smaller electronic component that replaced vacuum tubes in the conventional second-generation account.\\
Integrated circuit (IC) & A compact circuit incorporating multiple electronic components; the key third-generation association.\\
Microprocessor & A processor implemented as an integrated chip; the key fourth-generation association.\\
VLSI & Very Large Scale Integration: integrating a very large number of circuit components on a chip.\\
Machine language & Instructions expressed in the low-level codes directly used by the machine.\\
Assembly language & A symbolic, low-level notation for machine operations; it must be translated before execution.\\
High-level language & A programming language using more abstract statements than machine code or assembly.\\
Operating system & System software that manages the computer's operation and provides services to programs and users.\\
Multiprogramming & Keeping multiple programs available so the system can organize processor use while one program is waiting.\\
\bottomrule
\end{tabular}
\end{center}

The language and system-software terms help explain what changed beyond the
hardware labels. ``Machine language'' identifies the first generation's
programming cue. Assembly and early high-level languages appear with the
second-generation association. Operating systems and multiprogramming
develop with the third-generation systems described here. These terms are
not alternate names for the hardware generation; they are developments
associated with it.

\section{The first four technology associations}
\begin{center}
\begin{tabular}{p{0.20\textwidth}p{0.27\textwidth}p{0.43\textwidth}}
\textbf{Generation} & \textbf{Main association} & \textbf{Exam cue}\\
\hline
First & Vacuum tubes & Large, heat-producing systems; machine-language programming.\\
Second & Transistors & Smaller and more reliable; assembly and early high-level languages.\\
Third & Integrated circuits (ICs) & Multiple electronic components integrated
in compact circuits; smaller, more reliable systems; operating systems and
multiprogramming develop.\\
Fourth & Microprocessors & VLSI integration; personal computers become
widespread.\\
\end{tabular}
\end{center}

The progression is a change in how electronic components are built and
combined. Vacuum tubes made early systems large and heat-producing.
Transistors replaced those tubes and supported smaller, more reliable
systems. Integrated circuits placed multiple components together in a
compact circuit, continuing the trend toward reduced size and improved
reliability. In the fourth-generation association, a microprocessor brings
processor functions together on a chip; VLSI integration and the spread of
personal computers are the accompanying cues.

The sequence is the important memory structure:
\[
\text{vacuum tubes}\ \longrightarrow\ \text{transistors}\
\longrightarrow\ \text{integrated circuits}\ \longrightarrow\
\text{microprocessors}.
\]
The most common near-neighbour error is to switch the third and fourth
generations. An integrated circuit is the third-generation cue; the
microprocessor is the fourth-generation cue.

\begin{example}
An item asks which generation is associated with integrated circuits. Trace
the sequence: vacuum tubes are first, transistors second, ICs third, and
microprocessors fourth. Therefore the answer is the third generation. If
the stem instead names a microprocessor, the conventional answer is fourth.
\end{example}

\section{The fifth-generation association}
The conventional fifth-generation association commonly includes Artificial
Intelligence (AI), parallel processing, and natural-language systems. This
is a broad category, not a claim that every modern computer has one
universally agreed start date or a single defining design. AI applications
can run on systems that are not otherwise classified by one precise
``fifth-generation'' boundary. For an exam stem that explicitly uses this
generation scheme, match the conventional association and avoid treating
the label as a sharply dated technical standard.

\begin{example}
If the stem asks which broad generation is commonly linked to AI and
parallel processing in the conventional classification, select the fifth.
If it asks for the microprocessor generation, select the fourth instead;
``modern'' alone is not a sufficient clue to erase the distinction.
\end{example}

\section{India's PARAM 8000 milestone}
For the named national milestone, C-DAC's official high-performance
computing history identifies:
\href{https://cdac.in/index.aspx?id=hpc_ss_supercomputing_systems}
{C-DAC, HPC Supercomputing Systems}.

\begin{definitionbox}{C-DAC}
Centre for Development of Advanced Computing.
\end{definitionbox}

\begin{definitionbox}{PARAM 8000}
An Indian supercomputer developed by C-DAC and unveiled in 1991; C-DAC
identifies it as India's first supercomputer.
\end{definitionbox}

Keep the three parts of this fact separate. C-DAC is the organization;
PARAM 8000 is the system; 1991 is the unveiling year. This milestone
belongs to India's high-performance-computing history. It is not a
generation label, and it should not be used to redefine the conventional
technology sequence.

\begin{example}
For a matching item with ``Centre for Development of Advanced Computing,''
select C-DAC. For a prompt asking for India's first supercomputer, select
PARAM 8000. If the prompt asks when the system was unveiled, use 1991.
These answers refer to related but different parts of the same fact.
\end{example}

\section{Exam retrieval and common confusions}
Use two retrieval directions. Given the technology, recall the generation;
given the generation, recall the technology. Then isolate named milestones
that sit outside the generation sequence.
\begin{itemize}
  \item Vacuum tubes $\leftrightarrow$ first generation.
  \item Transistors $\leftrightarrow$ second generation.
  \item Integrated circuits $\leftrightarrow$ third generation.
  \item Microprocessors $\leftrightarrow$ fourth generation.
  \item AI and parallel processing $\leftrightarrow$ common fifth-generation
    description.
  \item C-DAC / PARAM 8000 / 1991 $\leftrightarrow$ separate Indian
    supercomputing milestone.
\end{itemize}

Three distinctions prevent most errors in this topic:
\begin{enumerate}
  \item Do not swap the adjacent third- and fourth-generation cues: ICs are
    third; microprocessors are fourth.
  \item Do not treat ``fifth generation'' as a single precisely dated
    hardware boundary. AI, parallel processing, and natural-language
    systems are the conventional association used in this lesson.
  \item Do not mix a generation label with the PARAM 8000 fact. The latter
    identifies a named Indian supercomputer, its developer, and its
    unveiling year.
\end{enumerate}

\subsection*{Exercises}
\begin{enumerate}[label=\arabic*.]
  \item Match vacuum tubes, transistors, integrated circuits, and
    microprocessors to the first four generations.
  \item Explain why the presence of an IC inside a modern microprocessor
    does not change the conventional microprocessor/fourth-generation
    association.
  \item Expand C-DAC and identify PARAM 8000's significance.
  \item Why should fifth-generation date boundaries be treated cautiously?
  \item Distinguish a generation milestone from the PARAM 8000 milestone.
\end{enumerate}

\subsection*{Solutions}
\begingroup\small
\begin{enumerate}[label=\arabic*.]
  \item First: vacuum tubes; second: transistors; third: integrated
    circuits; fourth: microprocessors.
  \item The generation scheme uses the dominant historical milestone as its
    label. A later microprocessor can contain ICs without changing the
    conventional fourth-generation association.
  \item C-DAC is Centre for Development of Advanced Computing. PARAM 8000
    was developed by C-DAC and is identified as India's first
    supercomputer.
  \item Generation boundaries are broad historical groupings and do not
    have one universally fixed transition date; technology associations
    are the more stable recall cue.
  \item A generation milestone classifies a historical technology period.
    PARAM 8000 is a named Indian supercomputer and national HPC milestone,
    not one of the generation labels.
\end{enumerate}
\endgroup

\lessonheading{50}{Data Representation, Number Systems \& Logic}
\section{Bases and positional notation}
Digital circuits commonly use two reliably distinguishable states, written
as 0 and 1. A number system specifies a base (also called a radix) and the
digits allowed in each position. The base determines the place values.

\begin{definitionbox}{Number-system base}
The number of distinct digit symbols in a positional number system, with
place values given by powers of that base.
\end{definitionbox}

\begin{center}
\begin{tabular}{llll}
\textbf{Name} & \textbf{Base} & \textbf{Allowed digits} &
\textbf{Example}\\
\hline
Binary & 2 & 0--1 & $1011_2$\\
Octal & 8 & 0--7 & $27_8$\\
Decimal & 10 & 0--9 & $45_{10}$\\
Hexadecimal & 16 & 0--9, A--F & $2F_{16}$\\
\end{tabular}
\end{center}

The subscript identifies the base. In hexadecimal, A represents 10, B 11,
C 12, D 13, E 14, and F 15. A digit that is not allowed in the stated base
makes the numeral invalid: for example, digit 8 cannot appear in an octal
number.

To test whether a numeral is valid, compare every digit with the allowed
range for its base. Binary permits only 0 and 1; octal permits 0 through 7;
decimal permits 0 through 9; hexadecimal permits 0 through 9 and A through
F. A numeral can be converted only after it passes this check. For example,
$128_8$ is not a valid octal numeral because it contains the digit 8, even
though the same characters could be interpreted under a different base.

\begin{definitionbox}{Binary}
A base-2 positional number system using digits 0 and 1.
\end{definitionbox}

\begin{definitionbox}{Octal}
A base-8 positional number system using digits 0 through 7.
\end{definitionbox}

\begin{definitionbox}{Decimal}
A base-10 positional number system using digits 0 through 9.
\end{definitionbox}

\begin{definitionbox}{Hexadecimal}
A base-16 positional number system using digits 0 through 9 and letters A
through F, where A represents 10 and F represents 15.
\end{definitionbox}

Each position contributes its digit multiplied by a power of the base. For
$1011_2$, the rightmost position is $2^0$, the next is $2^1$, then $2^2$,
and then $2^3$:
\[
1011_2=1(2^3)+0(2^2)+1(2^1)+1(2^0)=8+0+2+1=11_{10}.
\]
The rightmost position is the least-significant position because it has the
smallest power.

\begin{center}
\begin{tabular}{ccccc}
\toprule
\textbf{Digit} & 1 & 0 & 1 & 1\\
\textbf{Position from right} & 3 & 2 & 1 & 0\\
\textbf{Place value} & $2^3=8$ & $2^2=4$ & $2^1=2$ & $2^0=1$\\
\textbf{Contribution} & 8 & 0 & 2 & 1\\
\bottomrule
\end{tabular}
\end{center}

The place-value method works in every base. Multiply each digit by the
power assigned to its position, then add the contributions. The rightmost
digit always has power zero; moving one place left increases the exponent
by one. This is why leading zeros do not change a value: each added zero
contributes zero at a new, higher place.

\section{Converting decimal and binary}
To convert a positive decimal integer to binary, divide repeatedly by two.
At each step, keep the remainder, which must be either zero or one. Stop
when the quotient becomes zero, then read the remainders from the last
division upward. Reading them in the order produced gives the wrong
direction.

Use this sequence every time:
\begin{enumerate}
  \item Divide the current integer by 2.
  \item Record the remainder, either 0 or 1.
  \item Replace the current integer with the quotient and repeat until the
    quotient is 0.
  \item Read the recorded remainders from last to first.
  \item Verify the result by adding its binary place-value contributions.
\end{enumerate}

\begin{example}
Convert $45_{10}$ to binary:
\[
\begin{array}{rcl}
45&=&2(22)+1\\
22&=&2(11)+0\\
11&=&2(5)+1\\
5&=&2(2)+1\\
2&=&2(1)+0\\
1&=&2(0)+1
\end{array}
\]
The remainders from top to bottom are $1,0,1,1,0,1$. Reading bottom to top
gives $101101_2$. Check by place value:
$1(32)+0(16)+1(8)+1(4)+0(2)+1(1)=45$, so
$45_{10}=101101_2$.
\end{example}

\begin{example}
Convert $13_{10}$ to binary. Repeated division gives
$13=2(6)+1$, $6=2(3)+0$, $3=2(1)+1$, and $1=2(0)+1$. Reading the
remainders upward gives $1101_2$. Check:
$1(8)+1(4)+0(2)+1(1)=13$.
\end{example}

\section{Binary grouping for octal and hexadecimal}
Three binary bits encode one octal digit because $2^3=8$. Four binary bits
encode one hexadecimal digit because $2^4=16$. To convert a binary numeral,
start at the right and group bits into threes for octal or fours for
hexadecimal. If the leftmost group is short, add leading zeros; leading
zeros do not change the value.

\begin{example}
Convert $1111111_2$ to octal. Group from the right:
$001\ 111\ 111$. The groups are $001_2=1_8$, $111_2=7_8$, and
$111_2=7_8$, so the answer is $177_8$. Check:
$1(8^2)+7(8)+7=64+56+7=127$, the same value as $1111111_2$.
\end{example}

\begin{example}
Convert $1111111_2$ to hexadecimal. Group from the right:
$0111\ 1111$. These groups are $7$ and $15$, and hexadecimal digit 15 is
F. Therefore the answer is $7F_{16}$. Check:
$7(16)+15=112+15=127$.
\end{example}

\begin{example}
Convert $2F_{16}$ to decimal. The hexadecimal digit F is 15:
\[
2F_{16}=2(16^1)+15(16^0)=32+15=47_{10}.
\]
\end{example}

\section{Character codes, Unicode, and decimal digits}
\begin{definitionbox}{ASCII}
A character-encoding standard whose original character set uses 7-bit
codes.
\end{definitionbox}

\begin{definitionbox}{Unicode}
A standard that assigns code points to characters across many writing
systems.
\end{definitionbox}

\begin{definitionbox}{UTF-8}
A variable-length encoding form for Unicode code points.
\end{definitionbox}

\begin{definitionbox}{Binary-Coded Decimal (BCD)}
A representation in which each decimal digit is encoded separately using a
4-bit code.
\end{definitionbox}

\begin{center}
\begin{tabular}{p{0.2\textwidth}p{0.31\textwidth}p{0.37\textwidth}}
\toprule
\textbf{Representation} & \textbf{What it identifies or encodes} & \textbf{Key property}\\
\midrule
ASCII & Characters & Original ASCII uses 7-bit codes\\
Unicode & Characters from many writing systems & Assigns a code point to a character\\
UTF-8 & Unicode code points as encoded bytes & Variable-length encoding\\
BCD & Decimal digits & Uses a separate 4-bit code for each digit\\
\bottomrule
\end{tabular}
\end{center}

ASCII and Unicode concern character representation. Unicode code points
identify characters; UTF-8 is one way to encode those code points as bytes.
The code point and its byte encoding are different things: a character's
Unicode identity does not by itself state how many bytes a particular
encoding uses. UTF-8 is variable-length, so a character is not always stored
using one byte. BCD has a different purpose: it represents each decimal
digit separately using four bits. It is not the same as converting the
entire multi-digit decimal value to binary.

\begin{example}
Represent decimal $59$ in BCD. Encode the digit 5 as $0101$ and the digit 9
as $1001$, giving $0101\ 1001$. Ordinary binary for the numeric value 59 is
$111011_2$. These bit strings differ because BCD preserves the two decimal
digits separately, while ordinary binary represents the value 59 as one
base-2 numeral.
\end{example}

The grouping explains why BCD and ordinary binary use different numbers of
bits. BCD uses one 4-bit group for each decimal digit, so a two-digit
decimal value uses two groups. Ordinary binary instead places the entire
quantity in powers of 2. When an item asks for BCD, split the decimal
number into its digits before encoding; do not convert the whole number
directly to binary.

\section{RGB colour representation}
\begin{definitionbox}{RGB}
A colour model representing colour through red, green, and blue channel
values.
\end{definitionbox}

In a common 24-bit RGB representation, each of the three channels uses 8
bits. Therefore the three channels require $8+8+8=24$ bits per pixel.
Eight bits can encode $2^8=256$ possible values per channel, from 0 through
255. The total number of possible channel combinations in that representation
is $256^3=16{,}777{,}216$. This does not mean every image format or display
uses exactly 24-bit RGB; other bit depths, alpha channels, palettes, or
colour spaces are possible.

The channel count and the bit count are separate quantities. RGB has three
channels: red, green, and blue. In the common representation described
above, each channel has 8 bits, so one pixel uses 24 bits total. A value
within one channel ranges from 0 to 255 because an 8-bit field has
$2^8=256$ possible patterns. To count all channel combinations, multiply
the number of choices for each independent channel:
\[
  256\times256\times256=256^3=16{,}777{,}216.
\]
This count follows only under the stated 8-bits-per-channel assumption.

\begin{example}
In common 24-bit RGB, the red, green, and blue channels each use 8 bits.
Adding the channel widths gives $8+8+8=24$ bits per pixel.
\end{example}

\section{Logic gates and truth tables}
\begin{definitionbox}{AND gate}
A logic operation that outputs 1 only when every input is 1.
\end{definitionbox}

\begin{definitionbox}{OR gate}
A logic operation that outputs 1 when at least one input is 1.
\end{definitionbox}

\begin{definitionbox}{XOR gate}
A logic operation that outputs 1 when its two inputs differ.
\end{definitionbox}

\begin{definitionbox}{NAND gate}
The negation of AND: it outputs 0 only when every input is 1.
\end{definitionbox}

\begin{definitionbox}{NOT gate}
A one-input logic operation that reverses its input: 0 becomes 1 and 1
becomes 0.
\end{definitionbox}

\begin{center}
\begin{tabular}{cc|cccc|c}
\toprule
$A$ & $B$ & AND & OR & XOR & NAND & NOT $A$\\
\midrule
0 & 0 & 0 & 0 & 0 & 1 & 1\\
0 & 1 & 0 & 1 & 1 & 1 & 1\\
1 & 0 & 0 & 1 & 1 & 1 & 0\\
1 & 1 & 1 & 1 & 0 & 0 & 0\\
\bottomrule
\end{tabular}
\end{center}

Read the truth table from the rule, not by memorizing an unlabelled row.
AND requires both inputs to be 1. OR needs at least one 1. XOR checks
whether the inputs differ. NAND first evaluates AND and then reverses that
output. NOT has only one input and reverses it.

For two-input gates there are exactly four possible input pairs:
$00$, $01$, $10$, and $11$. Evaluate every pair using the gate's rule:
\begin{itemize}
  \item AND is 1 only on the $11$ row.
  \item OR is 0 only on the $00$ row.
  \item XOR is 1 on $01$ and $10$, where the inputs differ.
  \item NAND is the opposite of AND on every row.
  \item NOT does not use $B$; it reverses $A$ in each row.
\end{itemize}
OR and XOR are easy to confuse. For input $A=1$, $B=1$, OR is 1 because
at least one input is 1, but XOR is 0 because the two inputs are equal.
NAND is checked only after evaluating AND and then taking the opposite
result.

\begin{example}
For $A=1$ and $B=0$, AND is 0 because both inputs are not 1; OR is 1
because one input is 1; XOR is 1 because the values differ; and NAND is 1
because it reverses the AND result. NOT($A$) is 0.
\end{example}

\begin{example}
For $A=1$ and $B=1$, AND is 1 and OR is 1. XOR is 0 because the inputs do
not differ. NAND is the opposite of AND, so it is 0. NOT($A$) is 0.
\end{example}

\subsection*{Exercises}
\begin{enumerate}[label=\arabic*.]
  \item Convert $29_{10}$ to binary and verify the result by place value.
  \item Convert $101101_2$ to octal and hexadecimal.
  \item Explain the difference between BCD for decimal 25 and ordinary
    binary for the value 25.
  \item In common 24-bit RGB, how many bits are assigned to each channel?
    Is that rule universal for all images?
  \item For $A=0$, $B=1$ and for $A=1$, $B=1$, find AND, OR, XOR, NAND,
    and NOT($A$).
\end{enumerate}

\subsection*{Solutions}
\begingroup\small
\begin{enumerate}[label=\arabic*.]
  \item Repeated division gives $29=2(14)+1$, $14=2(7)+0$,
    $7=2(3)+1$, $3=2(1)+1$, and $1=2(0)+1$. Reading upward gives
    $11101_2$. Check: $16+8+4+0+1=29$.
  \item Group from the right: $101\ 101$, so the octal result is $55_8$.
    Pad and group by four: $0010\ 1101$, so the hexadecimal result is
    $2D_{16}$.
  \item BCD for 25 is $0010\ 0101$, encoding 2 and 5 separately.
    Ordinary binary for the value 25 is $11001_2$.
  \item Eight bits per channel, since $8+8+8=24$. No; this describes a
    common 24-bit representation, not every image or display.
  \item For $A=0$, $B=1$: AND=0, OR=1, XOR=1, NAND=1, and NOT($A$)=1.
    For $A=1$, $B=1$: AND=1, OR=1, XOR=0, NAND=0, and NOT($A$)=0.
\end{enumerate}
\endgroup

\unitopener{2}{CPU, Hardware \& Architecture}
\lessonheading{04}{Organization of a Computer}
\videoline{https://youtu.be/92XQ_n1t_cc}
\section{The functional units of a computer}
Start with the office desktop. A clerk types marks on the keyboard, the CPU
processes them through a spreadsheet formula, the monitor and printer present
the result, and the file is kept on the SSD for later use. This single machine
contains every functional unit that the exam asks you to name.

\begin{definitionbox}{Functional unit}
A major part of a computer that performs one broad job --- accepting input,
producing output, holding data, performing operations, or controlling the
machine.
\end{definitionbox}

\begin{center}
\begin{tabular}{p{0.25\textwidth}p{0.63\textwidth}}
\toprule
\textbf{Unit} & \textbf{Job in this lesson}\\
\midrule
Input unit & Accepts data and instructions from the outside world.\\
Output unit & Presents the result to the user.\\
Memory (primary storage) & Holds data and instructions currently in use.\\
Arithmetic Logic Unit (ALU) & Performs arithmetic and logical operations.\\
Control unit & Directs the sequence of operations and the flow of data.\\
\bottomrule
\end{tabular}
\end{center}

The memory unit described here is \emph{primary storage}, that is, main
memory. Secondary storage such as a hard disk or an SSD is a separate idea and
is treated in Section 4.

\section{The CPU: ALU, control unit and registers}
The \key{Central Processing Unit (CPU)} is the processing centre of the
computer. It is built from three parts: the arithmetic logic unit, the control
unit, and registers.

\begin{definitionbox}{Central Processing Unit (CPU)}
The unit that fetches, decodes and executes instructions; it is formed by the
ALU, the control unit and the registers.
\end{definitionbox}

\begin{definitionbox}{Arithmetic Logic Unit (ALU)}
The part of the CPU that performs arithmetic operations (such as addition and
subtraction) and logical operations (such as comparison) on binary data.
\end{definitionbox}

\begin{definitionbox}{Control unit}
The part of the CPU that fetches and decodes instructions, issues control
signals, and coordinates the operation of all other units.
\end{definitionbox}

\begin{definitionbox}{Register}
A small, very fast storage location inside the CPU used to hold data,
addresses or control information during processing.
\end{definitionbox}

The ALU does the calculating, but it does not decide what to do. The control
unit decides the sequence of operations and the flow of data, and then tells
the ALU and the other units when to act. Registers are the CPU's own working
storage; they are not the same as main memory or secondary storage.

\begin{example}
In the marks spreadsheet, the formula that adds three marks is carried out by
the ALU. The control unit has already fetched the formula's instruction from
memory, decoded it, and signalled the ALU to act. The intermediate values are
kept in registers while the operation runs. The final total is then written
back to memory and later shown on the monitor.
\end{example}

\section{Input and output units}
The input unit brings data and instructions into the system. The output unit
presents the result.

\begin{center}
\begin{tabular}{p{0.30\textwidth}p{0.58\textwidth}}
\toprule
\textbf{Input devices} & \textbf{Output devices}\\
\midrule
Keyboard, mouse, scanner, microphone, light pen, digitizer & Monitor,
printer, plotter, speaker, sound-card output\\
\bottomrule
\end{tabular}
\end{center}

A device is classified by the role it plays in the described task, not by its
shape. A monitor is normally an output device because it presents results. A
scanner is an input device because it brings data in. When a question lists
devices, mark each one as input or output before choosing the option.

\section{Memory (primary storage) versus secondary storage}
This is a high-value contrast.

\begin{definitionbox}{Primary storage (main memory)}
The memory the CPU accesses directly while processing; it holds the programs
and data currently in use. Random Access Memory (RAM) is volatile; Read Only
Memory (ROM) is non-volatile.
\end{definitionbox}

\begin{definitionbox}{Secondary storage}
Storage that keeps data and files permanently but is not accessed directly by
the CPU during processing; examples are the hard disk, the Solid State Drive
(SSD) and the pen drive.
\end{definitionbox}

The key rule for the exam is that the \key{CPU works with main memory, not
with secondary storage}. When a program on the hard disk is run, it is first
loaded into main memory; the CPU then works from main memory. Secondary
storage is reached through main memory and the operating system.

\begin{example}
The spreadsheet file is saved on the SSD. When the clerk opens it, the file is
first loaded into RAM. All the additions and comparisons are then performed
on data held in RAM and in the CPU registers. Writing the result back to the
SSD happens later, through the operating system. At no point does the ALU
process the file straight from the SSD.
\end{example}

\section{The system bus: address, data and control}
The \key{system bus} is the set of wires that carries information between the
CPU, memory and the other units. It has three parts, and their directions
differ --- this is a common trap.

\begin{center}
\begin{tabular}{p{0.20\textwidth}p{0.26\textwidth}p{0.46\textwidth}}
\toprule
\textbf{Bus} & \textbf{Direction} & \textbf{Role}\\
\midrule
Address bus & CPU $\rightarrow$ memory & Carries the address of the memory
location the CPU wants to use.\\
Data bus & Bidirectional & Carries the actual data and instructions between
the CPU and memory.\\
Control bus & CPU $\rightarrow$ units & Carries control and timing signals,
such as read and write.\\
\bottomrule
\end{tabular}
\end{center}

\begin{definitionbox}{System bus}
The address bus, data bus and control bus taken together; they connect the
CPU with memory and the input/output units.
\end{definitionbox}

Remember the three questions: the address bus answers \emph{where}, the data
bus answers \emph{what}, and the control bus answers \emph{when and how}. The
address bus is one-way (from the CPU outwards) and carries the address, not
the contents of the location. The data bus is two-way. The control bus
carries signals, not data or addresses.

\begin{example}
The CPU wants the value stored at memory location 200. It places the address
200 on the address bus and sends a read signal on the control bus. The memory
puts the contents of location 200 onto the data bus, and the CPU receives it.
The address travelled out on the address bus; the data travelled back on the
data bus.
\end{example}

\section{The motherboard, expansion slots and interfaces}
The \key{motherboard} is the main circuit board that connects and links the
CPU, main memory and the buses. Expansion slots on the motherboard allow
extra cards to be added, such as a graphics card or a network card. Interfaces
and ports connect external devices --- keyboard, printer, monitor --- to the
system. The system bus runs along the motherboard between the CPU and the
other units, so the motherboard is the physical platform on which the units
communicate.

\section{The stored-program (von Neumann) idea}
In early machines, changing the task often meant rewiring the machine. The
\key{stored-program concept} changed this.

\begin{definitionbox}{Stored-program concept (von Neumann architecture)}
The design in which a program's instructions are stored in main memory
together with the data, so the CPU can fetch and execute them without
rewiring the machine.
\end{definitionbox}

Because instructions and data share the same memory, the CPU fetches an
instruction from memory, executes it, and moves on to the next. This repeating
pattern is the instruction cycle, described next.

\section{The instruction cycle}
The \key{instruction cycle} is the sequence the CPU follows for each
instruction.

\begin{enumerate}
  \item \key{Fetch}: the CPU obtains the next instruction from main memory.
  \item \key{Decode}: the control unit works out what the instruction means.
  \item \key{Execute}: the ALU or another unit performs the operation.
  \item \key{Store} (when required): the result is written back to memory or
    to a register.
\end{enumerate}

The order is fixed: fetch, then decode, then execute. The optional store step
writes the result back. Data and instructions move between the units over the
system bus, and the whole cycle repeats for the next instruction.

\begin{example}
Consider an instruction ``add the value in memory to a register''. During
fetch, the CPU gets the instruction from memory. During decode, the control
unit identifies it as an addition and signals the ALU. During execute, the ALU
performs the addition. The updated value is then stored back in the register
or in memory, and the cycle begins again.
\end{example}

Three distinctions prevent most errors in this topic:
\begin{enumerate}
  \item The CPU accesses main memory directly, not secondary storage.
  \item Address, data and control buses are not interchangeable: the address
    bus is one-way and carries addresses; the data bus is two-way and carries
    data and instructions; the control bus carries control signals.
  \item Memory (primary) holds current work; storage (secondary) keeps files
    permanently.
\end{enumerate}

\subsection*{Exercises}
\begin{enumerate}[label=\arabic*.]
  \item Name the five functional units of a computer and state one job of
    each.
  \item What are the three parts of the CPU, and what does each do?
  \item State the direction and role of the address bus, the data bus and the
    control bus.
  \item A file is opened from a hard disk and edited. Explain why the CPU does
    not process the file directly from the hard disk.
  \item List the steps of the instruction cycle in order.
\end{enumerate}

\subsection*{Solutions}
\begingroup\small
\begin{enumerate}[label=\arabic*.]
  \item Input unit (accepts data and instructions); output unit (presents
    results); memory/primary storage (holds data and instructions in use);
    ALU (performs arithmetic and logical operations); control unit (directs
    operations and data flow).
  \item The ALU, the control unit and the registers. The ALU computes; the
    control unit fetches, decodes and signals; the registers hold data and
    addresses during processing.
  \item Address bus: one-way from CPU to memory, carrying the address of the
    location. Data bus: bidirectional, carrying data and instructions.
    Control bus: carries control and timing signals from the CPU to other
    units.
  \item The CPU accesses only main memory directly. The file is first loaded
    into main memory, the CPU processes the copy in main memory and registers,
    and the operating system writes the changes back to the hard disk.
  \item Fetch, decode, execute (and store when a result must be written back).
\end{enumerate}
\endgroup

\lessonheading{05}{CPU, ALU, Control Unit \& Registers}
\videoline{https://youtu.be/b6w8uT4GBYQ}
\section{The CPU and its functional units}
The Central Processing Unit (CPU) is the part of the computer that executes
instructions. It repeatedly fetches an instruction from main memory, decodes
it, and then executes it. To do this work the CPU is organised into three
functional units: the Arithmetic Logic Unit (ALU), the Control Unit (CU),
and the registers. Data and control signals move between these units over
internal buses. The CPU works directly with main memory (RAM); it does not
reach into secondary storage such as a hard disk by itself.

\begin{definitionbox}{CPU (Central Processing Unit)}
The processor that fetches, decodes, and executes instructions, made up of
the ALU, the Control Unit, and registers.
\end{definitionbox}

A simple way to keep the roles apart is this: the ALU computes, the Control
Unit coordinates, and the registers hold the working values. Almost every
objective question about this topic turns on not swapping those three roles.

\section{The Arithmetic Logic Unit (ALU)}
The ALU is the computing part of the CPU. It performs two broad families of
operations. \textbf{Arithmetic} operations include addition, subtraction,
increment, decrement, multiplication, and division. \textbf{Logical}
operations include AND, OR, NOT, XOR, shifts, and comparisons. The ALU
operates on data held in binary form, and its operations are carried out with
logic-gate operations on those binary values. When the ALU finishes, it
returns its result to a register or the accumulator; it does not keep a
permanent copy.

\begin{definitionbox}{ALU (Arithmetic Logic Unit)}
The CPU unit that performs arithmetic and logical operations on binary data
using logic gates.
\end{definitionbox}

\begin{example}
For inputs $1$ and $1$, an addition produces the binary result for two, and
a logical XOR produces $0$ because the two bits are equal. Both results are
produced by the same unit --- the ALU --- and handed back for further use.
\end{example}

\section{The Control Unit (CU)}
The Control Unit directs and sequences the operations of the computer. It
decodes each instruction and then issues the control signals that time and
trigger the other units. It also controls the flow of data between the ALU,
the registers, and memory. What the Control Unit does \textbf{not} do is
compute: it does not add, subtract, or carry out logical operations. Those
are the ALU's job. Students frequently write ``the Control Unit performs the
calculations,'' which is exactly the error to avoid.

\begin{definitionbox}{Control Unit (CU)}
The CPU unit that directs and sequences operations, decodes instructions,
and issues control signals; it does not perform arithmetic or logic.
\end{definitionbox}

\begin{example}
When an addition instruction is executed, the Control Unit fetches and
decodes the instruction and signals the ALU to add the two values. The ALU
performs the addition; the Control Unit only directs the sequence.
\end{example}

\section{Registers}
Registers are the smallest and fastest storage in a computer. They are
located inside the CPU, so they are faster than cache and far faster than
main memory. Because they are very small, they hold only the values needed
for the work currently in progress: instructions, addresses, data, and
intermediate results. The following table lists the key registers and their
purposes.

\begin{center}
\begin{tabular}{p{0.30\textwidth}p{0.60\textwidth}}
\toprule
\textbf{Register} & \textbf{Purpose} \\
\midrule
Program Counter (PC / IP) & Holds the address of the \textbf{next} instruction to be fetched. \\
Instruction Register (IR) & Holds the instruction currently being decoded or executed. \\
Memory Address Register (MAR) & Holds the address of the memory location to be accessed. \\
Memory Data Register (MDR) & Holds the data being transferred to or from memory. \\
Accumulator (ACC) & Holds intermediate arithmetic and logic results. \\
General-purpose registers & Hold temporary operands and results. \\
Status / flags register & Holds condition flags such as zero, carry, sign, and overflow. \\
\bottomrule
\end{tabular}
\end{center}

The Program Counter is the register most often confused. It holds the
address of the \textbf{next} instruction to fetch, not the address of the
instruction currently being executed. The instruction currently being
decoded or executed is held in the Instruction Register. Memory access is
split between two registers: the MAR carries the address, and the MDR carries
the data itself. The accumulator is the natural home for an ALU result that
will be used again immediately.

\begin{example}
Suppose the CPU is about to fetch the instruction at address 200. The
Program Counter holds 200. Once the instruction is fetched, it is placed in
the Instruction Register for decoding. If the instruction needs to read a
value stored at address 900, the MAR holds 900 and the MDR holds the value
brought back from memory.
\end{example}

\section{Clock speed, cycles, cores, and threads}
The CPU's operations are paced by a clock. \textbf{Clock speed} is measured
in hertz (Hz); one gigahertz (GHz) is one billion clock cycles per second. A
\textbf{clock cycle} is a single tick of that clock. Do not assume every
instruction takes exactly one cycle: different instructions can require
different numbers of clock cycles. A \textbf{core} is an independent
processing unit within the CPU; a quad-core chip contains four cores. A
\textbf{thread} is a stream of instructions, and a system can run multiple
threads across the available cores.

\begin{definitionbox}{Clock cycle}
One tick of the system clock; the basic unit of timing for CPU operations.
\end{definitionbox}

\begin{example}
The office desktop in our running example is a quad-core machine running at
2.4 GHz. It has four independent cores, and each core can complete clock
cycles at about 2.4 billion per second. Running a spreadsheet calculation and
a browser at the same time lets different threads keep the cores busy.
\end{example}

\section{Exam retrieval and common confusions}
Three distinctions prevent most errors on this topic.
\begin{enumerate}
  \item Do not swap the ALU and the Control Unit. The ALU performs
    arithmetic and logic; the Control Unit only directs operations and
    issues control signals.
  \item Do not say the Program Counter holds the current instruction. It
    holds the address of the \textbf{next} instruction; the current
    instruction is in the Instruction Register.
  \item Do not place registers below cache or RAM in speed. Registers are
    the fastest storage, and the accumulator is where intermediate
    arithmetic and logic results are kept.
\end{enumerate}

\subsection*{Exercises}
\begin{enumerate}[label=\arabic*.]
  \item State the three main functional units of the CPU and one role of
    each.
  \item Expand ALU, CU, PC, IR, MAR, and MDR.
  \item Does the Control Unit perform arithmetic? Explain in one line.
  \item Which register holds the address of the next instruction, and which
    register holds the instruction currently being executed?
  \item Why are registers described as the fastest storage, and what does
    the accumulator hold?
\end{enumerate}

\subsection*{Solutions}
\begingroup\small
\begin{enumerate}[label=\arabic*.]
  \item ALU --- performs arithmetic and logical operations on binary data;
    Control Unit --- directs and sequences operations and issues control
    signals; registers --- hold instructions, addresses, data, and
    intermediate results.
  \item ALU = Arithmetic Logic Unit; CU = Control Unit; PC = Program
    Counter; IR = Instruction Register; MAR = Memory Address Register;
    MDR = Memory Data Register.
  \item No. The Control Unit directs and coordinates operations; the ALU
    performs the arithmetic and logical calculations.
  \item The Program Counter (PC/IP) holds the address of the next
    instruction; the Instruction Register (IR) holds the instruction
    currently being decoded or executed.
  \item Registers are located inside the CPU and are smaller and faster than
    cache and RAM. The accumulator (ACC) holds intermediate arithmetic and
    logic results.
\end{enumerate}
\endgroup

\lessonheading{06}{CPU vs GPU \& Microprocessor}
\videoline{https://youtu.be/Vaod5Ox7lzM}
\section{Serial and parallel processing}
\textbf{Serial processing} means that one task is carried out at a time, in
sequence: the first operation finishes before the next begins. \textbf{Parallel
processing} means that several tasks or operations proceed at the same time.
A serial workload rewards a processor with a few very fast units; a parallel
workload rewards a processor with many units that can work simultaneously.

\begin{definitionbox}{Serial processing}
Carrying out one task at a time, in sequence, so that each step completes
before the next starts.
\end{definitionbox}

\begin{definitionbox}{Parallel processing}
Carrying out several tasks or operations at the same time.
\end{definitionbox}

The difference matters because it explains why two processors can both be
described as ``fast'' while being designed for opposite jobs. A processor that
excels at finishing a single chain of instructions quickly is a
latency-optimised design. A processor that excels at completing a very large
number of small independent operations per second is a
throughput-optimised design.

\section{CPU versus GPU}
The CPU (Central Processing Unit) is the general-purpose processor of a
computer. It has \textbf{few, powerful cores}, and each core is capable of
running a sequential instruction stream quickly. Its design goal is
\textbf{low latency}: to complete one task with the least possible delay. In
the office-desktop running example (EX-01), the CPU executes the operating
system, \mbox{MS Word}, and spreadsheet logic step by step.

The GPU (Graphics Processing Unit) is a processor built around
\textbf{many simpler cores}. Each core does a small amount of work, and the
design goal is \textbf{high throughput}: a very large number of independent
operations per second. This makes the GPU a natural fit for pixels and
graphics, and also for scientific and machine-learning workloads that can be
expressed as many parallel calculations.

\begin{center}
\begin{tabular}{p{0.20\textwidth}p{0.34\textwidth}p{0.34\textwidth}}
\toprule
\textbf{Feature} & \textbf{CPU} & \textbf{GPU}\\
\midrule
Cores & Few, powerful & Many, simple\\
Design goal & Low latency & High throughput\\
Best suited to & Sequential, general-purpose work & Parallel work (pixels, ML)\\
Core type & General-purpose & Specialised, simpler\\
Example use & EX-01 desktop logic & Image and graphics operations\\
\bottomrule
\end{tabular}
\end{center}

\begin{example}
An item asks which processor is better suited to an image-processing task
that repeats the same operation on millions of pixels. Because those pixel
operations are independent and can run in parallel, the GPU --- with its many
simple cores --- is the better fit. The CPU's few powerful cores would each
handle the chain of pixels one after another.
\end{example}

The most common near-neighbour error is to treat GPU cores as if they were
interchangeable general-purpose CPU cores. They are not. GPU cores are
simpler and optimised for parallel throughput, so they are not a full
replacement for the CPU's general-purpose cores.

\section{Cores, threads, and hyper-threading}
A \textbf{core} is a physical processing unit. A quad-core CPU contains four
such units. A \textbf{thread} is a stream of instructions that a core
executes. \textbf{Hyper-threading} lets a single physical core present two
logical threads, so the core can stay busy while one thread is waiting.

Two logical threads on one core are not two physical cores. A quad-core
processor with hyper-threading may report eight logical threads, but it still
has four physical cores. Clock speed and core count together, not thread
count alone, determine raw capacity.

\begin{definitionbox}{Hyper-threading}
A technique in which one physical core presents two logical threads to the
operating system, improving utilisation while a thread waits.
\end{definitionbox}

\section{Cache depth: L1, L2, and L3}
Cache is small, fast memory held close to the core. It is organised in levels.
\textbf{L1} is the smallest and fastest level and sits nearest the core.
\textbf{L2} is larger and a little slower than L1. \textbf{L3} is the largest
and slowest of the three and is typically shared across cores. The order of
speed and closeness is L1, then L2, then L3, then main memory.

\begin{center}
\begin{tabular}{p{0.12\textwidth}p{0.30\textwidth}p{0.42\textwidth}}
\toprule
\textbf{Level} & \textbf{Relative size} & \textbf{Relative speed and location}\\
\midrule
L1 & Smallest & Fastest; closest to the core\\
L2 & Larger & Slower than L1\\
L3 & Largest & Slowest of the three; often shared across cores\\
\bottomrule
\end{tabular}
\end{center}

\section{Clock speed and multi-core}
\textbf{Clock speed} is measured in \textbf{GHz} (gigahertz), which counts
billions of cycles per second. A higher clock speed generally means more work
completed per second, all else being equal. A \textbf{multi-core} processor
places several cores on one chip, which enables parallel work. However, clock
speed alone never decides real performance: the number of cores, the cache
size, the amount of memory, and the input/output speed all contribute.

\section{Microprocessor versus microcontroller}
A \textbf{microprocessor} is a central processing unit implemented on a single
chip. It is a general-purpose processor, and software decides what it does.
On its own it still needs external memory and input/output devices to form a
complete working system. The EX-01 office desktop is built around a
microprocessor.

A \textbf{microcontroller} integrates more on the same chip: a microprocessor
\textbf{plus} memory \textbf{plus} input/output peripherals. It therefore
forms a near-complete computer on one chip and is used for
\textbf{embedded/special-purpose} control, for example in washing machines,
remote controls, and appliance or vehicle controllers.

\begin{definitionbox}{Microprocessor}
A central processing unit implemented on a single chip; a general-purpose
device that needs external memory and I/O.
\end{definitionbox}

\begin{definitionbox}{Microcontroller}
A microprocessor together with integrated memory and input/output peripherals
on one chip, used in embedded and special-purpose systems.
\end{definitionbox}

\begin{center}
\begin{tabular}{p{0.28\textwidth}p{0.30\textwidth}p{0.30\textwidth}}
\toprule
\textbf{Feature} & \textbf{Microprocessor} & \textbf{Microcontroller}\\
\midrule
Contents & CPU on a chip & CPU + memory + I/O on a chip\\
Purpose & General-purpose & Embedded / special-purpose\\
System & Needs external parts & Self-contained on one chip\\
Example & EX-01 office desktop & Appliance controller\\
\bottomrule
\end{tabular}
\end{center}

The key difference is the \textbf{level of integration}, not clock speed. A
microcontroller is not ``just a microprocessor'': it is a more integrated
device that already includes memory and I/O on the same chip.

\section{Diagnosing a ``slow'' computer (EX-01)}
The office desktop in this running example is a quad-core 2.4\,GHz machine
with 16\,GB of RAM and an SSD. Suppose one particular program still runs
slowly. A common but incorrect conclusion is that the CPU must be slow,
because ``more cores should mean faster''. More cores do not help a
single-threaded program, which uses only one core at a time. The real cause is
usually a \textbf{bottleneck}: memory pressure, disk input/output, or software
that is not multithreaded.

\begin{example}
A question describes a quad-core 2.4\,GHz desktop with 16\,GB RAM and an SSD,
and asks why one application is slow. The best answer is to diagnose a stated
bottleneck --- a single-threaded program, insufficient memory, or slow disk
input/output --- rather than to declare the CPU slow. A slow program is not
automatically a ``slow CPU.''
\end{example}

\section{Exam retrieval and common confusions}
Use two retrieval directions. Given the processor, recall its core profile;
given the workload, recall the better processor.
\begin{itemize}
  \item CPU $\leftrightarrow$ few powerful general-purpose cores, low latency.
  \item GPU $\leftrightarrow$ many simple cores, high throughput, parallel.
  \item Serial $\leftrightarrow$ one task at a time in sequence.
  \item Parallel $\leftrightarrow$ several tasks or operations at once.
  \item Microprocessor $\leftrightarrow$ CPU on a single chip, general-purpose.
  \item Microcontroller $\leftrightarrow$ microprocessor + memory + I/O on one
    chip, embedded.
  \item Cache order $\leftrightarrow$ L1 (smallest, fastest) to L3 (largest,
    slowest).
\end{itemize}

Three distinctions prevent most errors in this topic:
\begin{enumerate}
  \item Do not treat GPU cores as general-purpose CPU cores. They are simpler
    and optimised for parallel throughput.
  \item Do not confuse serial with parallel. Serial is sequential, one task at
    a time; parallel is simultaneous.
  \item Do not attribute a slow program on a capable machine automatically to
    a slow CPU. Check the multithreading, memory, and disk-I/O bottleneck; and
    do not confuse a microprocessor (CPU on a chip) with a microcontroller
    (CPU + memory + I/O on one chip).
\end{enumerate}

\subsection*{Exercises}
\begin{enumerate}[label=\arabic*.]
  \item Distinguish serial processing from parallel processing.
  \item State one design goal of the CPU and one of the GPU, and give the
    core profile of each.
  \item Explain the difference between a microprocessor and a microcontroller.
  \item What does hyper-threading allow, and why is it not the same as adding
    a physical core?
  \item A quad-core 2.4\,GHz desktop with 16\,GB RAM and an SSD runs one
    program slowly. Give the most appropriate diagnosis and say why.
\end{enumerate}

\subsection*{Solutions}
\begingroup\small
\begin{enumerate}[label=\arabic*.]
  \item Serial processing performs one task at a time in sequence; parallel
    processing performs several tasks or operations at the same time.
  \item The CPU is optimised for low latency and has few powerful
    general-purpose cores. The GPU is optimised for high throughput and has
    many simpler cores for parallel work.
  \item A microprocessor is a CPU on a single chip and is general-purpose,
    needing external memory and I/O. A microcontroller adds memory and I/O
    peripherals on the same chip, forming an embedded, special-purpose system;
    the difference is the level of integration.
  \item Hyper-threading lets one physical core present two logical threads, so
    it can stay busy while one thread waits. It does not add a physical core;
    a quad-core chip with hyper-threading still has four physical cores.
  \item The program is likely limited by a bottleneck --- a single-threaded
    workload, insufficient memory, or slow disk input/output --- rather than
    by a slow CPU, since extra cores do not speed up a single-threaded
    program. A slow program is not automatically a ``slow CPU.''
\end{enumerate}
\endgroup

\lessonheading{59}{Current CPU Architecture, Data Representation \& RISC/CISC}
\section{von Neumann architecture and the bottleneck}
A \key{von Neumann} (stored-program) computer keeps both instructions and
data in the same memory. The processor fetches an instruction, decodes it,
and executes it; when the instruction needs a data operand, that operand is
read from the same memory. Because a single shared path carries instructions
and data, the processor frequently has to wait for memory. This limitation is
called the \key{von Neumann bottleneck}: the rate at which work can proceed is
capped by the rate at which instructions and data can be transferred between
the CPU and memory over that shared path. Faster memory, a cache, and wider
paths all reduce --- but do not remove --- the bottleneck.

\begin{definitionbox}{von Neumann architecture}
A stored-program design in which instructions and data are held together in
the same memory and fetched over a shared path between the CPU and memory.
\end{definitionbox}

The important contrast is \emph{same memory for instructions and data}. A
system that keeps them physically separate is a different (Harvard-style)
design and is not the von Neumann model.

\section{Addressing modes: immediate addressing}
An \key{addressing mode} describes how an instruction states where its
operand is. \key{Immediate addressing} embeds the operand value directly
inside the instruction. For example, an instruction ``add 5 to the register''
carries the value 5 in the instruction encoding itself. Because the value is
already present, no extra memory access is required to obtain the operand.

The near-neighbour contrasts matter for exams:
\begin{itemize}
  \item \emph{Immediate}: the operand is in the instruction.
  \item \emph{Register}: the operand is in a named register.
  \item \emph{Direct}: the instruction names a memory address holding the
    operand.
\end{itemize}
The single line to remember is: immediate addressing means the operand is a
constant carried with the instruction.

\section{Cache and the cache-hit ratio}
A \key{cache} is a small, fast memory placed between the CPU and main memory.
It holds copies of recently used instructions and data. When the CPU asks for
an item and the cache has it, that access is a \key{hit}; if the cache does not
have it, the access is a \key{miss} and the item must come from main memory
(and is usually placed in the cache for next time).

\begin{definitionbox}{Cache-hit ratio}
\[
\text{hit ratio} = \frac{\text{number of hits}}{\text{total accesses}}
= \frac{\text{hits}}{\text{hits} + \text{misses}}.
\]
\end{definitionbox}

A \key{hit} is faster because cache access time is much lower than the access
time of main memory. On a hit the CPU avoids the slow main-memory access; on
a miss it pays that cost. This is why a higher hit ratio improves average
performance.

\begin{example}
If a program makes 900 cache hits and 100 cache misses, the total accesses
are $900+100=1000$ and the hit ratio is $900/1000 = 0.90$ (90\%).
\end{example}

\section{Data hazards even with register addressing}
A \key{data hazard} is a dependency that prevents an instruction from
proceeding correctly in the cycle originally planned. A common case is a
read-after-write dependency: an instruction needs the value of a register
that an earlier instruction has not yet written. Pipelining is what exposes
these dependencies, because several instructions are in progress at once.

Register addressing avoids an extra memory access and removes some hazards
caused by memory operands, but it does \bad{not} eliminate all hazards. Two
instructions can still depend on the same register. This is the trap: a
candidate may reason that ``no memory operand means no hazard,'' which is
false. Hazards arise from dependencies between instructions, whatever kind of
operand is used.

\section{Virtual memory versus cache}
Cache and virtual memory are frequently confused, but they solve different
problems.

\begin{center}
\begin{tabular}{p{0.20\textwidth}p{0.36\textwidth}p{0.36\textwidth}}
\toprule
\textbf{Point} & \textbf{Cache} & \textbf{Virtual memory} \\
\midrule
Location & Small fast memory next to the CPU. & Uses main memory plus secondary storage as backing store. \\
Purpose & Reduce average memory access time by keeping recently used blocks close. & Make the memory available to programs appear larger than physical RAM. \\
Management & Largely hardware. & Operating system with hardware support. \\
\bottomrule
\end{tabular}
\end{center}

Both involve small ``fast'' storage for larger ``slow'' storage, but the cache
shortens access time while virtual memory expands the apparent size of memory.
They are not the same mechanism and are not interchangeable.

\section{Instruction cycle: Program Counter and Instruction Register}
During the fetch stage the CPU reads an instruction from memory. The
\key{Program Counter (PC)} holds the address of the \key{next} instruction to
be fetched; once the current instruction has been fetched, the PC is advanced
to point at the following one. The \key{Instruction Register (IR)} holds the
instruction that is currently being fetched, decoded, and executed.

The exam trap is the PC. It is easy to say the PC holds ``the address of the
current instruction,'' but at the moment the instruction is executing, the PC
already points at the \emph{next} instruction. The correct statement is: the
PC holds the address of the next instruction to be fetched.

\section{Interrupts are processed at instruction boundaries}
An \key{interrupt} is a signal that causes the CPU to suspend the current
program temporarily and transfer control to a handler routine, after which
execution usually resumes. Interrupts are recognised at \key{instruction
boundaries} --- between one instruction and the next --- during the program's
execution, not only when the whole program finishes.

The trap is the claim that ``an interrupt is processed only after the entire
program completes.'' That is false. Mid-program interrupts are the normal
case: the processor finishes the instruction in progress, saves its state, and
services the interrupt before continuing.

\section{Instruction timing: clock cycles vary}
Different machine instructions perform different amounts of work. A simple
register-to-register operation may complete in a single clock cycle, while a
division, a complex operation, or an instruction that must access slow memory
can take many cycles. Instruction timing therefore depends on the
instruction, its operands, and the memory system.

The trap is the claim that ``all machine instructions take the same number of
clock cycles.'' The count varies by instruction; averaging over a program
gives an average cycles-per-instruction, not a constant per instruction.

\section{Two's complement and overflow detection}
In an $n$-bit \key{two's-complement} representation the most significant bit
is the \key{sign} bit: 0 for non-negative and 1 for negative values. A
negative number is formed by inverting the bits of its magnitude and adding
one. With $n$ bits the representable range is
$-2^{\,n-1}$ through $+2^{\,n-1}-1$.

\key{Overflow} occurs when the arithmetic result cannot be represented in the
available bits. The reliable test is the \key{sign} rule:

\begin{definitionbox}{Overflow test}
Overflow occurs when two operands of the \key{same sign} are added and the
result has the \key{opposite sign}. If the two operands have different signs,
no overflow occurs.
\end{definitionbox}

Overflow is \bad{not} the same as a carry out of the most significant bit. A
carry can be produced without overflow, and overflow can occur without the
final carry telling the full story. Use the sign test, not the carry flag.

\begin{example}[Overflow]
Work in 8-bit two's complement. Add $+7$ and $+1$:
\[
0111\ (+\!7) \;+\; 0001\ (+\!1) \;=\; 1000.
\]
The operands are both positive, but the sum $1000$ has sign bit 1 and
represents $-8$. Two same-sign operands produced an opposite-sign result, so
this is \good{overflow}.
\end{example}

\begin{example}[Carry without overflow]
Again in 8 bits, add $-1$ and $-1$:
\[
1111\,1111\ (-\!1) \;+\; 1111\,1111\ (-\!1) \;=\; 1\,1111\,1110.
\]
There is a carry out of the most significant bit, but the discarded 9th bit
leaves $1111\,1110$, which is exactly $-2$. The result is correct, so there is
\good{no overflow}. This shows that a carry out of the MSB is not the same as
overflow.
\end{example}

\section{Zero-extension versus sign-extension}
\key{Extension} widens a value from fewer bits to more bits while preserving
its meaning. The rule depends on whether the value is unsigned or signed:
zero-extension fills the new leftmost bits with 0 and is used for
\key{unsigned} values; sign-extension copies the sign bit into the new
leftmost bits and is used for \key{signed} values. The common trap is to swap
the two.

\begin{center}
\begin{tabular}{llll}
\toprule
\textbf{8-bit pattern} & \textbf{Unsigned} & \textbf{Signed} & \textbf{16-bit sign-extended} \\
\midrule
$0000\,0101$ & $5$ & $+5$ & $0000\,0000\,0000\,0101$ \\
$1111\,1011$ & $251$ & $-5$ & $1111\,1111\,1111\,1011$ \\
\bottomrule
\end{tabular}
\end{center}

Zero-extending $1111\,1011$ would give $0000\,0000\,1111\,1011$ (value $251$),
which preserves the \emph{unsigned} value. Sign-extending it gives
$1111\,1111\,1111\,1011$ (value $-5$), which preserves the \emph{signed}
value. Using the wrong rule changes the value.

\section{BCD: binary-coded decimal}
\key{BCD} (Binary-Coded Decimal) represents each \key{decimal digit} with its
own group of \key{4 bits}. The ten digit codes are:

\begin{center}
\begin{tabular}{cccccccccc}
\toprule
$0$ & $1$ & $2$ & $3$ & $4$ & $5$ & $6$ & $7$ & $8$ & $9$ \\
\midrule
$0000$ & $0001$ & $0010$ & $0011$ & $0100$ & $0101$ & $0110$ & $0111$ & $1000$ & $1001$ \\
\bottomrule
\end{tabular}
\end{center}

\begin{example}[BCD encoding]
Encode the decimal number $59$ in BCD. Take each digit separately: $5$ is
$0101$ and $9$ is $1001$. So $59$ becomes
\[
0101\ 1001 \quad\text{(BCD)}.
\]
For contrast, the ordinary binary value of $59$ is $111011_2$. BCD keeps one
4-bit code per decimal digit; it does not pack the whole number as plain
binary.
\end{example}

The trap is the claim that BCD is ``just the binary of the whole number.''
It is not: each decimal digit is coded in 4 bits, and the codes $1010$ through
$1111$ are not valid BCD digits.

\section{Byte-addressability and effective address vs instruction length}
Memory is \key{byte-addressable} when every individual byte has its own unique
address. A word is a group of bytes at consecutive addresses. Because bytes
are individually addressable, a word access begins at the address of its
first byte.

The \key{effective address} is the actual location that an instruction uses
for its operand. Depending on the addressing mode, it may be given directly
or computed at run time from registers and constants. Instruction length is
the number of bytes the instruction itself occupies. These are different
quantities: a long instruction does not imply a distant operand. The address
of the operand is set by the addressing mode, not by the instruction's size.
Confusing the two is the trap; keep ``how big is the instruction'' separate
from ``where is the operand.''

\section{RISC versus CISC}
The two broad instruction-set philosophies are \key{RISC} (Reduced Instruction
Set Computer) and \key{CISC} (Complex Instruction Set Computer).

\begin{center}
\begin{tabular}{p{0.22\textwidth}p{0.34\textwidth}p{0.34\textwidth}}
\toprule
\textbf{Aspect} & \textbf{RISC} & \textbf{CISC} \\
\midrule
Instruction set & Small, simple, regular. & Large and complex. \\
Instruction length & Fixed. & Variable. \\
Registers & Many general-purpose registers. & Fewer registers. \\
Memory access & Load/store architecture: only load and store touch memory. & Instructions may operate directly on memory operands. \\
Cycles per instruction & Typically one, suited to pipelining. & Often multiple cycles. \\
Control emphasis & Hardwired control; work shifted to the compiler. & Microcoded control; work done in hardware. \\
Typical examples & ARM, MIPS, RISC-V. & x86 and older mainframe designs. \\
\bottomrule
\end{tabular}
\end{center}

RISC keeps instructions simple and fixed-length so they pipeline well, relying
on many registers and a load/store discipline. CISC offers powerful,
variable-length instructions that can do more in one step, at the cost of
more complex decoding. The exam cue is the pair: \emph{RISC = reduced, fixed,
many registers, load/store}; \emph{CISC = complex, variable-length, memory
operands allowed}.

\subsection*{Exercises}
\begin{enumerate}[label=\arabic*.]
  \item State what is shared in von Neumann architecture and define the von
    Neumann bottleneck.
  \item What is immediate addressing, and why does it need no extra memory
    access for the operand?
  \item A program makes 850 cache hits and 150 misses. Compute the
    cache-hit ratio.
  \item Explain why register addressing does not eliminate all data hazards.
  \item In 8-bit two's complement, add $+7$ and $+1$ and state whether
    overflow occurs. Justify using the sign test.
  \item Encode the decimal number $59$ in BCD, and contrast it with the plain
    binary value.
\end{enumerate}

\subsection*{Solutions}
\begingroup\small
\begin{enumerate}[label=\arabic*.]
  \item Instructions and data share the same memory. The bottleneck is the
    limited rate at which instructions and data can be transferred between
    the CPU and memory over the shared path.
  \item Immediate addressing carries the operand inside the instruction
    encoding, so the value is already available and no memory fetch is
    required for it.
  \item Total accesses $=850+150=1000$, so the hit ratio is
    $850/1000 = 0.85$ (85\%).
  \item Register operands avoid a memory access, but two instructions can
    still depend on the same register (for example a read-after-write
    dependency), so hazards remain.
  \item $0111\ (+\!7) + 0001\ (+\!1) = 1000$, which has sign bit 1 and
    represents $-8$. Two same-sign operands produced an opposite-sign
    result, so overflow occurs.
  \item BCD: $5 \rightarrow 0101$, $9 \rightarrow 1001$, giving
    $0101\ 1001$. Plain binary of $59$ is $111011_2$; BCD keeps one 4-bit
    code per decimal digit instead of packing the number as binary.
\end{enumerate}
\endgroup

\unitopener{3}{Memory, Storage \& Backup}
\lessonheading{07}{Computer Memory Hierarchy}
\section{Why a memory hierarchy exists}
A running computer never keeps all its data and instructions in one place.
Instead it uses several stores arranged in a hierarchy. The reason is a
three-way trade-off: the memory that is fastest is also the smallest and the
most expensive per bit, while the memory that is largest and cheapest is the
slowest. No single memory can be fast, large, and cheap at the same time.
The hierarchy places the small, fast, costly memories close to the CPU and
the large, slow, cheap memories further away.

\begin{definitionbox}{Memory hierarchy}
The ordered arrangement of storage levels in a computer, from the fastest
and smallest (registers) through cache and main memory down to the slowest
and largest (secondary and offline storage), chosen to balance speed, cost,
and capacity.
\end{definitionbox}

\begin{center}
\begin{tabular}{p{0.24\textwidth}p{0.66\textwidth}}
\toprule
\textbf{Term} & \textbf{Meaning in this lesson} \\
\midrule
Register & The fastest and smallest storage, located inside the CPU. \\
Cache & A small, very fast memory that holds recently and frequently used data. \\
Main memory (RAM) & The primary memory that holds active programs and data; it is volatile. \\
ROM & Non-volatile read-only primary memory that holds firmware such as the BIOS. \\
Secondary storage & Non-volatile storage such as an SSD or HDD, used for files. \\
Tertiary / offline & Backup and archival storage such as magnetic tape or cloud backup. \\
\bottomrule
\end{tabular}
\end{center}

The rest of this lesson uses these terms in one consistent sense. A
\textbf{register} is the CPU's own storage. \textbf{Cache} is the small fast
memory between the CPU and main memory. \textbf{Main memory} means RAM: the
primary memory that holds what is running. \textbf{Secondary storage} means
SSD and HDD, which keep files without power. \textbf{Tertiary or offline}
storage is the slowest level, used for backup and archive.

\section{The ordered hierarchy}
The levels, read from fastest to slowest, are registers, cache, main memory
(RAM), secondary storage, and then tertiary or offline storage.

\begin{center}
\begin{tabular}{p{0.20\textwidth}p{0.18\textwidth}p{0.16\textwidth}p{0.36\textwidth}}
\toprule
\textbf{Level} & \textbf{Speed} & \textbf{Size} & \textbf{Typical role} \\
\midrule
Registers & Fastest & Smallest & Values in use right now \\
Cache (L1/L2/L3) & Very fast & Small & Recently and frequently used data \\
Main memory (RAM) & Moderately fast & Larger & Active programs and data \\
Secondary (SSD/HDD) & Slower & Large & Permanent file storage \\
Tertiary/offline (tape, cloud) & Slowest & Largest & Backup and archive \\
\bottomrule
\end{tabular}
\end{center}

Reading downward, capacity tends to grow and speed tends to fall. The
anchor to remember is at the top: \textbf{registers are the fastest} storage
in the computer. Below them cache is still very fast but slower than the
registers. Main memory (RAM) is larger than cache but slower. Secondary
storage (SSD or HDD) is larger again, cheaper per bit, and slower, but it
keeps its contents without power. Tertiary or offline storage (magnetic
tape, cloud backup) is the slowest and cheapest level and is used for backup
and long-term archive.

\section{Registers and cache}
Registers are located inside the CPU and are the fastest and smallest
storage in the machine. Because they are so small, they hold only the values
needed at this instant: instructions, addresses, data, and intermediate
results.

Cache is a small, very fast memory that sits between the CPU and main
memory. It holds recently and frequently used data and instructions, and it
is faster than main memory. Cache does not store new information of its own;
it keeps \emph{copies} of data that really live in RAM. It is therefore not
a backup, and it loses its contents when power is removed, just as RAM does.

\begin{definitionbox}{Cache}
A small, very fast memory that holds recently and frequently used data and
instructions so the CPU can avoid the slower trip to main memory.
\end{definitionbox}

Cache is usually described in levels. \textbf{L1} is the smallest and
fastest and sits closest to the CPU core. \textbf{L2} is larger and a little
slower than L1. \textbf{L3} is the largest and slowest of the three, is
often shared across cores, and is still faster than main memory. For exam
purposes, the key point is the ordering: L1 is fastest and smallest, and L3
is largest and slowest among the three, but all cache is faster than RAM.

\begin{example}
While a student edits a document on the office desktop, the same lines and
the nearby rows are read again and again. The cache keeps those recently
used and nearby values ready, so the CPU does not have to wait for main
memory on every access. The cache holds copies of those values; the master
copy remains in RAM.
\end{example}

\section{Main memory (RAM) and ROM}
Main memory is the computer's primary memory. The programs and data that
are currently in use are held here while the machine runs, and the CPU works
directly with main memory. \textbf{RAM} (Random Access Memory) is
\emph{volatile}: its contents are lost when the power is switched off. RAM
is larger than cache but slower.

\textbf{ROM} (Read-Only Memory) is the other part of primary memory. It is
\emph{non-volatile}: it keeps its contents without power. In normal
operation it can be read but not written; changing it needs a special
process. ROM holds firmware such as the startup (BIOS) instructions.

\begin{definitionbox}{Volatile memory}
Memory that loses its contents when the power is removed. RAM is the common
example.
\end{definitionbox}

\begin{definitionbox}{Non-volatile memory}
Memory that retains its contents without power. ROM, SSD, and HDD are common
examples.
\end{definitionbox}

\begin{example}
An item states, ``RAM is non-volatile and will keep the file after
shutdown.'' This is wrong. RAM is volatile; its contents are lost when the
power goes off. It is \textbf{ROM} that is non-volatile. The wrong option has
swapped the two primary-memory types.
\end{example}

\section{Primary memory vs secondary storage}
\textbf{Primary memory} is directly accessible by the CPU and holds the
active programs and data. It includes RAM, which is volatile, and ROM, which
is non-volatile. \textbf{Secondary storage}, such as an SSD or HDD, is
non-volatile and keeps files and programs between sessions. It is larger and
cheaper per bit but slower than primary memory. The CPU does not access
secondary storage directly; data moves through main memory under the control
of the operating system.

\section{Secondary and tertiary storage}
Secondary storage comes in two common forms. A \textbf{solid-state drive
(SSD)} uses flash memory and has no moving parts; a \textbf{hard disk drive
(HDD)} uses rotating magnetic disks. Both are non-volatile and keep data
without power. \textbf{Tertiary or offline} storage, such as magnetic tape
and cloud backup, is the slowest and cheapest level and is used for backup
and long-term archive.

\section{Locality of reference}
Programs do not access memory at random; they tend to reuse the same data
and instructions. This behaviour is called \textbf{locality of reference},
and it comes in two forms. \textbf{Temporal locality} means the same item is
accessed again soon after it was last used, such as a loop counter used on
every iteration. \textbf{Spatial locality} means items stored at nearby
addresses are accessed next, such as the consecutive elements of an array.
This reuse is exactly why cache works: by keeping recently used and nearby
data ready, cache reduces the average time the CPU waits for data.

\begin{definitionbox}{Locality of reference}
The tendency of a running program to reuse recently accessed data
(temporal locality) and data at nearby addresses (spatial locality).
\end{definitionbox}

\section{Exam retrieval and common confusions}
Five distinctions prevent most errors on this topic.
\begin{enumerate}
  \item RAM is \textbf{volatile}; ROM is \textbf{non-volatile}. Do not swap
    them (TRAP-01).
  \item Cache is \textbf{not} a backup. It holds copies of data that live in
    RAM.
  \item RAM is main memory, \textbf{not} storage for files. Files belong in
    secondary storage.
  \item Cache is \textbf{faster} than main memory; reversing this is a common
    wrong answer (PYQ-03).
  \item The order of speed is registers, then cache, then RAM, then
    secondary storage. Registers are fastest.
\end{enumerate}

\begin{example}
An item asks which memory is fastest. Trace the hierarchy from the top:
registers are fastest and smallest, then cache, then main memory (RAM), then
secondary storage. The answer is registers. If the stem instead asks which
level holds recently used data, the answer is cache, and if it asks which is
volatile primary memory, the answer is RAM.
\end{example}

\subsection*{Exercises}
\begin{enumerate}[label=\arabic*.]
  \item List the levels of the memory hierarchy from fastest to slowest.
  \item Explain why the faster memories of the hierarchy are placed close to
    the CPU.
  \item Distinguish volatile memory from non-volatile memory, giving one
    example of each.
  \item Why is cache described as faster than main memory, and why is it not
    a backup?
  \item Define temporal locality and spatial locality, and explain how they
    relate to cache.
\end{enumerate}

\subsection*{Solutions}
\begingroup\small
\begin{enumerate}[label=\arabic*.]
  \item Registers, cache (L1, L2, L3), main memory (RAM), secondary storage
    (SSD/HDD), then tertiary/offline storage.
  \item Faster memory is smaller and more expensive per bit. Placing it close
    to the CPU keeps the values the CPU needs soonest at the fastest level,
    so the average access time is low while cost and capacity are balanced.
  \item Volatile memory loses its contents when power is removed --- RAM is
    the example. Non-volatile memory retains its contents without power ---
    ROM, or an SSD/HDD, is the example.
  \item Cache sits between the CPU and main memory and is a smaller, faster
    technology, so it is faster than RAM. It only keeps copies of data that
    already exist in RAM, so if the cache is lost the data is still in RAM;
    it is therefore not a backup.
  \item Temporal locality is the reuse of the same item soon after it was
    last accessed; spatial locality is the use of data at nearby addresses.
    Because programs show both patterns, keeping recently used and nearby
    data in the small fast cache pays off.
\end{enumerate}
\endgroup

\lessonheading{08}{RAM, ROM \& Cache}
\section{Memory, volatility, and why the distinction matters}
\textbf{Main memory} (also called primary memory) holds the data and
instructions that the CPU uses directly. Secondary storage such as a hard
disk or an SSD is not accessed by the CPU directly; its contents are brought
into main memory first. This lesson concentrates on the primary-memory
components a computer-awareness paper actually asks about: RAM, ROM, and
cache.

The single most tested distinction in this lesson is \textbf{volatility}.

\begin{definitionbox}{Volatile memory}
Memory whose stored contents are lost when the power supply is removed.
\end{definitionbox}

\begin{definitionbox}{Non-volatile memory}
Memory whose stored contents remain available after the power supply is
removed.
\end{definitionbox}

RAM is volatile; ROM is non-volatile. A practice item may phrase this as a
NOT/incorrect-statement question, so learn both directions: ``RAM loses its
contents on power-off,'' and ``ROM retains its contents without power.''

\section{RAM: volatile main memory}
\textbf{RAM} stands for \textbf{Random Access Memory}. It is the main memory
in which running programs and their active data are held. The term ``random
access'' means that any storage location can be reached in roughly the same
amount of time; it is not a statement about volatility. RAM is
\textbf{volatile}: when the office desktop (EX-01) is switched off, whatever
was in RAM is lost unless it was saved to secondary storage.

\begin{example}
A stem states that RAM keeps its data intact after the computer is switched
off. This is false. RAM is volatile, so its contents are lost on power-off;
it is ROM (and flash) that retains data without power. If the stem instead
asks which memory holds the running programs, select RAM.
\end{example}

\section{DRAM versus SRAM}
RAM comes in two main types, and the exam routinely contrasts them. Both are
RAM and both are volatile; the difference lies in speed, refresh, cost, and
use.

\begin{center}
\begin{tabular}{p{0.22\textwidth}p{0.33\textwidth}p{0.33\textwidth}}
\toprule
\textbf{Feature} & \textbf{DRAM} & \textbf{SRAM}\\
\midrule
Full form & Dynamic Random Access Memory & Static Random Access Memory\\
Refresh & Needs periodic refresh & No refresh needed\\
Speed & Slower & Faster\\
Density and cost & Denser and cheaper per bit & Less dense and costlier\\
Typical use & Main memory (the computer's RAM) & Cache memory\\
\bottomrule
\end{tabular}
\end{center}

\textbf{DRAM} stores each bit in a tiny capacitor that leaks charge, so it
must be \textbf{refreshed} periodically. That refresh overhead is the price
of its high density and low cost, which is why DRAM is the natural choice
for large main memory. \textbf{SRAM} uses a different cell design that holds
its state without refresh, making it faster but more expensive and less
dense; it is therefore used for the small, fast \textbf{cache}.

\begin{example}
An item asks which RAM type is used to build cache memory. SRAM is faster
and needs no refresh, so it is used for cache; DRAM is denser and cheaper
and serves as main memory. If the stem asks which type needs periodic
refreshing, the answer is DRAM.
\end{example}

\section{ROM and the ROM family}
\textbf{ROM} stands for \textbf{Read-Only Memory}. It is \textbf{non-volatile}
and is \textbf{read-only in normal operation}: the processor reads its
contents, and changing them requires a special programming process rather
than ordinary memory writes. ROM commonly stores \textbf{firmware}, such as
the start-up instructions a computer runs before an operating system loads.

\begin{definitionbox}{ROM}
Read-Only Memory: non-volatile primary memory that is normally read but not
written during ordinary operation.
\end{definitionbox}

Several ROM variants appear in exams. They differ mainly in \emph{when} they
are programmed and \emph{how} they are erased.

\begin{center}
\begin{tabular}{p{0.16\textwidth}p{0.74\textwidth}}
\toprule
\textbf{Type} & \textbf{Key property}\\
\midrule
PROM & Programmable Read-Only Memory: can be programmed \textbf{once} by the user; it cannot be erased and reprogrammed.\\
EPROM & Erasable Programmable Read-Only Memory: erased by exposing the chip to \textbf{ultraviolet (UV)} light.\\
EEPROM & Electrically Erasable Programmable Read-Only Memory: erased and rewritten \textbf{electrically}, without removing the chip.\\
Flash & An \textbf{EEPROM-derived} non-volatile memory erased in \textbf{blocks}; widely used in USB drives and SSDs.\\
\bottomrule
\end{tabular}
\end{center}

\begin{example}
An item asks which memory can be erased by ultraviolet light. PROM is
one-time programmable and cannot be erased; EEPROM is erased electrically;
flash is erased in blocks. The UV-erasable type is EPROM. If the stem
instead asks which type is programmable only once, select PROM.
\end{example}

\section{Cache memory and its levels}
\textbf{Cache} is a small, very fast memory placed between the CPU and main
memory. It holds frequently or recently used data and instructions so the
CPU does not have to wait as often for slower main memory. Cache is
\textbf{faster than main memory} --- a fact the paper has tested directly
(PYQ-03). It is temporary working memory, \textbf{not} a backup and not a
substitute for storage.

Cache is organised in levels.

\begin{definitionbox}{L1 cache}
The level-1 cache: the smallest, fastest cache, located closest to the CPU
core.
\end{definitionbox}

\begin{definitionbox}{L2 cache}
The level-2 cache: larger and slightly slower than L1.
\end{definitionbox}

\begin{definitionbox}{L3 cache}
The level-3 cache: the largest and slowest of the three levels, typically
shared across the cores of the processor.
\end{definitionbox}

As the level number rises, size grows and speed falls, and distance from the
core increases: L1 $<$ L2 $<$ L3 in size, with the speed order reversed.
Even L3, however, is faster than main memory.

\section{Cache hits, misses, and hit ratio}
Performance depends on how often the CPU finds what it needs already in the
cache.

\begin{definitionbox}{Cache hit}
A memory access in which the required data is found in the cache.
\end{definitionbox}

\begin{definitionbox}{Cache miss}
A memory access in which the required data is not found in the cache, so it
must be fetched from main memory.
\end{definitionbox}

The \textbf{hit ratio} measures how often accesses succeed in the cache:
\[
\text{hit ratio} = \frac{\text{hits}}{\text{total accesses}}.
\]

\begin{example}
If the cache receives 100 memory accesses and 90 of them are found in the
cache, then the hit ratio is $90/100 = 0.90$. The remaining 10 accesses are
misses and must go to main memory. A higher hit ratio means the CPU waits
less, so the system performs better.
\end{example}

\section{Exam retrieval and common confusions}
Use the following distinctions to answer both recall and NOT/incorrect
items reliably:
\begin{itemize}
  \item RAM is \textbf{volatile}; ROM is \textbf{non-volatile} (TRAP-01).
  \item ROM is \textbf{read-only in normal operation}; writing needs a
    special process (TRAP-02).
  \item \textbf{Cache is faster} than main memory, not slower (PYQ-03).
  \item \textbf{SRAM} is used for cache; \textbf{DRAM} is used for main
    memory --- and DRAM is the type that needs refresh.
  \item Cache is \textbf{not} a backup; it is temporary working memory.
  \item ROM family: PROM (once), EPROM (UV), EEPROM (electrical), flash
    (block erase).
  \item Hit ratio $=$ hits $\div$ total accesses.
\end{itemize}

\subsection*{Exercises}
\begin{enumerate}[label=\arabic*.]
  \item Explain why RAM is described as volatile and ROM as non-volatile.
  \item Distinguish DRAM from SRAM in terms of speed, refresh, and typical
    use.
  \item Expand PROM, EPROM, EEPROM, and flash, and state how each is erased
    (or whether it can be).
  \item Order L1, L2, and L3 cache by size and by speed, and state which
    level is typically shared across cores.
  \item A cache handles 200 memory accesses and 150 are hits. Compute the
    hit ratio.
\end{enumerate}

\subsection*{Solutions}
\begingroup\small
\begin{enumerate}[label=\arabic*.]
  \item RAM is volatile because it loses its contents when power is
    removed; ROM is non-volatile because it retains its contents without
    power.
  \item DRAM needs periodic refresh, is slower but denser and cheaper, and
    is used for main memory. SRAM needs no refresh, is faster but costlier
    and less dense, and is used for cache.
  \item PROM: Programmable Read-Only Memory, programmable once and not
    erasable. EPROM: Erasable Programmable Read-Only Memory, erased by
    ultraviolet light. EEPROM: Electrically Erasable Programmable Read-Only
    Memory, erased electrically. Flash: EEPROM-derived, erased in blocks.
  \item By size: L1 $<$ L2 $<$ L3, so L3 is the largest. By speed: L1 is
    fastest and L3 the slowest of the three. L3 is typically shared across
    cores.
  \item Hit ratio $= 150/200 = 0.75$.
\end{enumerate}
\endgroup

\lessonheading{09}{Virtual Memory}
\section{What virtual memory is and why it exists}
Physical RAM is limited. A modern application, its data, and the other
programs open at the same time can together demand more memory than the
machine physically has. \textbf{Virtual memory} is the technique that lets the
system keep running anyway: it uses secondary storage (the disk) to hold parts
of a program that are not currently needed in RAM. In this way the \emph{addressable}
memory seen by a program can be much larger than the physical memory actually
installed.

\begin{definitionbox}{Virtual memory}
A memory-management technique that uses secondary storage to extend the
addressable memory beyond physical RAM, giving each process the impression of a
large, contiguous memory space.
\end{definitionbox}

The office-desktop running example (EX-01) makes this concrete. The desktop has
16\,GB of RAM and runs MS Word, a spreadsheet, and a web browser at once. If the
combined working set of those programs exceeds 16\,GB, the operating system
moves some of the least-used pages to disk and brings them back when they are
referenced again. No single program is told that RAM ran out; each continues to
use its own large address space.

Two ideas are worth separating at the start:
\begin{itemize}
  \item Virtual memory extends \textbf{capacity}. It lets a program use more
    memory than the physical RAM provides.
  \item Virtual memory does \textbf{not} add speed. Because it relies on the
    disk, it is \textbf{much slower} than RAM.
\end{itemize}

\section{Virtual addresses, physical addresses, and the MMU}
A program does not work with the real location of its data in RAM. It works
with \textbf{virtual addresses}, which form a large, uniform address space.
The hardware translates each virtual address into the actual \textbf{physical
address} in RAM when the instruction executes.

\begin{definitionbox}{Virtual address}
The address generated by a running program; it belongs to the program's large
virtual address space and is not the real RAM location.
\end{definitionbox}

\begin{definitionbox}{Physical address}
The actual location of data in main memory (RAM), which the hardware uses after
translation.
\end{definitionbox}

The translator is the \textbf{MMU}, the Memory Management Unit.

\begin{definitionbox}{MMU (Memory Management Unit)}
The hardware that translates virtual addresses into physical addresses at run
time, using the page table maintained by the operating system.
\end{definitionbox}

The MMU is central to every memory access. It does \emph{not} store files and it
does not itself create storage; it maps one address form to another. Keeping
this role straight avoids two common exam errors: attributing file storage to
the MMU, or describing translation as a software-only operation with no
hardware component.

\begin{example}
A question asks what the MMU does. The correct answer is that it translates
virtual addresses into physical addresses. An option saying that the MMU stores
files permanently is wrong: file storage is the job of secondary storage, not
the address translator. An option saying the MMU holds the BIOS is also wrong;
that describes firmware, not the unit that maps addresses.
\end{example}

\section{Paging: fixed-size pages and frames}
\textbf{Paging} divides memory into equal, fixed-size blocks. The virtual
address space of a process is split into \textbf{pages}, and physical memory is
split into \textbf{frames} of exactly the same size. Any page can be placed in
any free frame, and the page table records which page currently sits in which
frame.

\begin{definitionbox}{Page}
A fixed-size block of the virtual address space in paging.
\end{definitionbox}

\begin{definitionbox}{Frame}
A fixed-size block of physical memory (RAM) into which a page is loaded; frames
are the same size as pages.
\end{definitionbox}

Because every block is the same size, paging makes memory allocation simple and
predictable. There is no awkward leftover space caused by blocks of many
different sizes. The cost is a small amount of unused space in the last page of
a program, but the mechanism itself is uniform.

Two related terms must not be swapped. \textbf{Pages} are the \emph{virtual}
blocks; \textbf{frames} are the \emph{physical} blocks. They are the same size,
but a page lives in the program's virtual view while a frame lives in RAM.

\section{Segmentation: variable-size logical segments}
\textbf{Segmentation} divides a program into \textbf{variable-size segments},
each corresponding to a logical unit of the program: the code, the data, the
stack, a procedure, or a module. Because these logical units differ in size,
segment sizes differ as well.

\begin{definitionbox}{Segment}
A variable-size block of a program that corresponds to a logical unit such as
code, data, stack, or a module.
\end{definitionbox}

Segmentation reflects how a programmer thinks about a program: one segment per
logical part. Paging, by contrast, reflects no logical unit at all --- its pages
are simply equal slices chosen for uniform allocation. This is the key
near-neighbour distinction between the two techniques.

\begin{example}
A stem asks which technique divides a program according to logical units such as
code and data. The answer is segmentation, because those units are logical and
of different sizes. If the stem instead describes equal fixed-size blocks and
the terms ``page'' and ``frame'', the answer is paging.
\end{example}

\section{Paging versus segmentation at a glance}
\begin{center}
\begin{tabular}{p{0.30\textwidth}p{0.30\textwidth}p{0.30\textwidth}}
\toprule
\textbf{Feature} & \textbf{Paging} & \textbf{Segmentation}\\
\midrule
Block size & Fixed & Variable\\
Blocks & Pages (virtual) and frames (physical) & Logical segments\\
Sizes equal? & Yes, all equal & No, they differ\\
Logical view & Not visible to the programmer & Matches program structure\\
Unit & Equal slices & Code, data, stack, module\\
\bottomrule
\end{tabular}
\end{center}

The single most common error is to call paging's blocks ``segments'' or to say
that segmentation uses equal-sized blocks. Keep the pairing firm:
\textit{paging is fixed-size; segmentation is variable-size and logical}.

\section{Page faults and demand paging}
When a program references a page that is \emph{not} currently in RAM, the
hardware raises a \textbf{page fault}. The operating system then loads the
required page from the disk into a free frame and lets the program continue.

\begin{definitionbox}{Page fault}
An event that occurs when a referenced page is not currently in RAM; the page
must be loaded from secondary storage (disk) into a free frame.
\end{definitionbox}

\begin{definitionbox}{Demand paging}
A policy in which a page is loaded into memory only when it is actually
referenced, rather than loading the whole program in advance.
\end{definitionbox}

A page fault is a normal, expected event --- \textbf{not a hardware fault}. It
simply means the needed page is on the disk for the moment. Demand paging is
what makes a program able to start before every one of its pages is in RAM: the
remaining pages are brought in on demand, as execution reaches them.

\begin{example}
A question says a running program references a page that is not in RAM and asks
what happens. The correct answer is a page fault followed by loading that page
from the disk. An option describing a physical hardware failure is wrong,
because a page fault is a routine memory-management event, not a broken
component.
\end{example}

\section{Swapping and the disk-backed hierarchy}
\textbf{Swapping} is the movement of pages or whole processes between RAM and
secondary storage (the disk). Because virtual memory depends on this movement,
it is \textbf{disk-backed}: the disk supplies the extra capacity that RAM does
not have.

\begin{definitionbox}{Swapping}
Moving pages or processes between main memory (RAM) and secondary storage
(disk) so that the memory actually needed at a moment can reside in RAM.
\end{definitionbox}

The memory hierarchy, from fastest and smallest to slowest and largest, runs
from CPU registers, to cache, to RAM, to disk-backed virtual memory. Two facts
follow directly:
\begin{enumerate}
  \item Virtual memory is \textbf{much slower than RAM}, because the disk is far
    slower than main memory.
  \item Virtual memory is \textbf{not additional physical RAM}. It borrows disk
    space; it does not add RAM chips.
\end{enumerate}

This is also where virtual memory is most often confused with \textbf{cache}.
They are opposite in purpose and location. Cache is small, fast memory sitting
close to the CPU and it \emph{speeds up} repeated access to hot data. Virtual
memory is large, disk-backed, and it \emph{extends capacity} using slower
storage. Cache helps speed; virtual memory helps size.

\begin{example}
A stem contrasts cache and virtual memory. The correct statement is that cache
is small, fast memory close to the CPU, while virtual memory is slower and
disk-backed. Options that call both ``secondary storage'', or that say virtual
memory speeds up the CPU, reverse the roles and are wrong.
\end{example}

\section{Thrashing}
\textbf{Thrashing} is \emph{excessive paging}. When too many processes compete
for too little RAM, the system spends most of its time moving pages between RAM
and disk, and very little time doing useful work. The CPU stays busy, yet
throughput collapses.

\begin{definitionbox}{Thrashing}
A condition of excessive paging in which the system spends most of its time
swapping pages between RAM and disk, leaving little time for useful work.
\end{definitionbox}

Thrashing is a symptom of memory pressure, not a hardware failure. Adding
physical RAM, or reducing the number of active processes, usually relieves it.

\section{Exam retrieval and common confusions}
Use two retrieval directions. Given the technique, recall its block rule; given
a description, name the technique.
\begin{itemize}
  \item Virtual memory $\leftrightarrow$ extends memory using disk; slower than
    RAM; not extra physical RAM.
  \item MMU $\leftrightarrow$ Memory Management Unit; translates virtual to
    physical addresses.
  \item Paging $\leftrightarrow$ fixed-size pages (virtual) and frames
    (physical).
  \item Segmentation $\leftrightarrow$ variable-size logical segments.
  \item Page fault $\leftrightarrow$ referenced page not in RAM; load it from
    disk.
  \item Demand paging $\leftrightarrow$ load a page only when it is referenced.
  \item Swapping $\leftrightarrow$ move pages/processes between RAM and disk.
  \item Thrashing $\leftrightarrow$ excessive paging.
\end{itemize}

Three distinctions prevent most errors in this topic:
\begin{enumerate}
  \item Do not say virtual memory is \textbf{extra physical RAM}. It is
    disk-backed and \textbf{slower} than RAM; it adds capacity, not speed.
  \item Do not confuse \textbf{paging} (fixed-size pages and frames) with
    \textbf{segmentation} (variable-size logical segments).
  \item Do not call a \textbf{page fault} a hardware fault. It means the
    referenced page must be loaded from the disk; and do not confuse
    \textbf{cache} (small, fast, near the CPU) with virtual memory (large,
    disk-backed).
\end{enumerate}

\subsection*{Exercises}
\begin{enumerate}[label=\arabic*.]
  \item Explain why virtual memory is said to extend memory but not to add
    physical RAM.
  \item Expand MMU and state its function in one sentence.
  \item Distinguish paging from segmentation in terms of block size and the
    unit being divided.
  \item What is a page fault, and what happens after one occurs?
  \item Define thrashing and give one practical way to reduce it.
\end{enumerate}

\subsection*{Solutions}
\begingroup\small
\begin{enumerate}[label=\arabic*.]
  \item Virtual memory extends the addressable memory by using secondary
    storage (disk) to hold parts of a program not currently in RAM. It does not
    add RAM chips, so it is not additional physical RAM; because it is
    disk-backed, it is much slower than RAM.
  \item MMU is the Memory Management Unit. It translates virtual addresses
    generated by a program into physical addresses in RAM at run time.
  \item Paging uses fixed-size, equal blocks: pages in the virtual space and
    frames in physical memory. Segmentation uses variable-size blocks that
    correspond to logical program units such as code, data, or a module.
  \item A page fault occurs when a program references a page that is not
    currently in RAM. The operating system then loads that page from the disk
    into a free frame and resumes execution.
  \item Thrashing is excessive paging, in which the system spends most of its
    time swapping pages between RAM and disk and little time on useful work.
    Adding physical RAM (or reducing the number of active processes) usually
    reduces it.
\end{enumerate}
\endgroup

\lessonheading{10}{Storage Devices}
\section{The storage hierarchy}
Storage is organized into levels according to how close it sits to the
processor and what job it does. The CPU works directly with primary storage;
secondary storage holds data permanently until it is needed and feeds it to
main memory; tertiary storage is an archival level used for large, rarely
accessed collections.

\begin{definitionbox}{Primary storage}
The directly CPU-accessible working memory of a computer, such as RAM, ROM,
and cache; it holds the data and instructions currently in use.
\end{definitionbox}

\begin{definitionbox}{Secondary storage}
Non-volatile storage that retains data when the power is off and is not
accessed directly by the CPU --- for example a hard disk, SSD, optical disc,
or USB flash drive.
\end{definitionbox}

\begin{definitionbox}{Tertiary storage}
An archival storage level for large volumes of infrequently used data,
typically robotic tape libraries whose media are loaded on demand.
\end{definitionbox}

\begin{center}
\begin{tabular}{p{0.2\textwidth}p{0.4\textwidth}p{0.3\textwidth}}
\toprule
\textbf{Level} & \textbf{Examples} & \textbf{Key property}\\
\midrule
Primary & RAM, ROM, cache & Directly accessible by the CPU\\
Secondary & HDD, SSD, optical disc, USB drive & Non-volatile; holds files\\
Tertiary & Robotic tape library & Archival; loaded on demand\\
\bottomrule
\end{tabular}
\end{center}

The CPU does not read a file straight from secondary storage. The operating
system brings the required data into main memory first, and the CPU then works
with that copy. This is why the hierarchy matters for an exam scenario: a
large program on disk is limited by the memory and storage that actually feed
it, not only by processor speed.

\section{Magnetic storage}
Magnetic storage records data by magnetizing a surface, and it is read back by
detecting those magnetic patterns. Hard disks and magnetic tape are the two
common magnetic media, but they differ sharply in how they are reached.

\subsection{Hard disk drive}
A hard disk drive (HDD) stores data on \emph{rotating platters} coated with a
magnetic material. A \emph{read/write head} mounted on an actuator arm floats
just above the platter surface and reads or writes as the platter spins. The
head moves inward and outward to reach different tracks. Because the platters
spin and the arm moves, an HDD is a \emph{mechanical} device with moving parts.

\begin{definitionbox}{Hard disk drive (HDD)}
Magnetic secondary storage that stores data on rotating platters read and
written by a moving read/write head; it is mechanical and contains moving
parts.
\end{definitionbox}

\subsection{Magnetic tape}
Magnetic tape stores data on a long, thin magnetic ribbon wound in a reel or
cartridge. The tape is read or written in order, position by position, so a
record far from the current position takes longer to reach. This makes tape a
\emph{sequential-access} medium. It is inexpensive and dense, which is why it
is used for backup and long-term archival rather than for routine access.

\begin{definitionbox}{Magnetic tape}
Magnetic storage on a long tape that is read and written in order; it is a
sequential-access medium used mainly for backup and archival.
\end{definitionbox}

\section{Optical storage}
Optical storage records data as \emph{pits and lands} --- tiny physical
markings --- on a disc and reads them with a \emph{laser}. The reflected light
pattern encodes the data. Optical discs are removable and non-volatile, so
they are convenient for distributing and archiving files.

The three common families are the compact disc (CD), the digital versatile
disc (DVD), and the Blu-ray disc. Blu-ray uses a shorter-wavelength laser and
tighter tracks, allowing it to store more than a DVD, which in turn stores
more than a CD.

\begin{definitionbox}{Optical storage}
Storage in which a laser reads data encoded as pits and lands on a disc, such
as a CD, DVD, or Blu-ray disc.
\end{definitionbox}

\begin{center}
\begin{tabular}{lll}
\toprule
\textbf{Medium} & \textbf{Approximate capacity (single layer)} &
\textbf{Reading method}\\
\midrule
CD & 700 MB & Laser\\
DVD & 4.7 GB & Laser (shorter wavelength)\\
Blu-ray & 25 GB & Laser (shortest wavelength)\\
\bottomrule
\end{tabular}
\end{center}

Order the capacities from smallest to largest as CD $<$ DVD $<$ Blu-ray.
A common test error is to swap the middle entries or to reverse the order
entirely; fix the sequence by remembering that CD is the oldest and smallest
and Blu-ray is the newest and largest.

\begin{example}
An item asks which optical disc holds the most data. Compare the standard
single-layer figures: CD $\approx$ 700 MB, DVD $\approx$ 4.7 GB, and Blu-ray
$\approx$ 25 GB. Since $25\ \text{GB} > 4.7\ \text{GB} > 700\ \text{MB}$, the
Blu-ray disc holds the most.
\end{example}

\section{Flash and solid-state storage}
Flash memory stores data in semiconductor cells that hold their state without
power. There are no rotating platters and no moving heads, so flash devices
are \emph{solid-state}. Common examples are the solid-state drive (SSD), the
USB flash drive, and the memory card.

\begin{definitionbox}{Flash memory}
Non-volatile semiconductor storage that keeps data in cells with no moving
parts; the technology behind SSDs, USB flash drives, and memory cards.
\end{definitionbox}

\begin{definitionbox}{Solid-state drive (SSD)}
A flash-based storage device with no moving parts. It is faster and more
shock-resistant than a mechanical hard disk and has a high transfer rate.
\end{definitionbox}

Because nothing spins and no head has to move, flash can start a transfer
almost immediately and is far less sensitive to being bumped or dropped.
Among a hard disk drive, an optical drive, magnetic tape, and an SSD, the
\key{SSD} has the highest transfer rate.

\section{Cloud storage}
Cloud storage keeps data on remote servers in a data centre and makes it
available over the Internet. The user sees a folder or an application, while
the provider manages the underlying hardware. Even though the storage feels
abstract, it ultimately rests on physical secondary storage --- arrays of
disks and solid-state drives --- on the provider's machines.

\begin{definitionbox}{Cloud storage}
Storage in which data is held on remote servers and accessed over the
Internet rather than on a locally attached device.
\end{definitionbox}

Cloud storage is not main memory and it is not the same as primary storage.
It is remote secondary storage reached through a network connection, useful
for backup, sharing, and access from several devices.

\section{Hard disk versus solid-state drive}
The comparison below is the most frequently tested near-neighbour pair in this
topic. The hard disk is magnetic and mechanical; the SSD is flash and
solid-state. Keep the technology and the consequences together.

\begin{center}
\begin{tabular}{p{0.2\textwidth}p{0.35\textwidth}p{0.35\textwidth}}
\toprule
\textbf{Aspect} & \textbf{Hard disk drive (HDD)} & \textbf{Solid-state drive (SSD)}\\
\midrule
Technology & Magnetic & Flash memory\\
Moving parts & Yes --- spinning platters, moving head & None\\
Speed & Usually slower transfer rate & Higher transfer rate (fastest here)\\
Shock resistance & More vulnerable to shock & More shock-resistant\\
Typical use & Bulk data at low cost per GB & Fast system and boot drive\\
\bottomrule
\end{tabular}
\end{center}

\begin{example}
A candidate reads ``the fastest storage device is the hard disk.'' That is
incorrect. The hard disk is magnetic and mechanical, with a spinning platter
and a moving head. The SSD is flash-based with no moving parts and has the
higher transfer rate. The correct choice is the SSD (KB PYQ-27).
\end{example}

\section{Access methods: random and sequential}
An access method describes how quickly a device can reach a stored location.

\begin{definitionbox}{Random (direct) access}
An access method in which any location can be reached in approximately equal
time; RAM and hard disks behave this way.
\end{definitionbox}

\begin{definitionbox}{Sequential access}
An access method in which locations are reached in order, so the time depends
on how far the target is from the current position; magnetic tape works this
way.
\end{definitionbox}

\begin{center}
\begin{tabular}{p{0.22\textwidth}p{0.3\textwidth}p{0.38\textwidth}}
\toprule
\textbf{Access method} & \textbf{Behavior} & \textbf{Examples}\\
\midrule
Random (direct) & Any location in about equal time & RAM, hard disk\\
Sequential & Reached in order & Magnetic tape\\
\bottomrule
\end{tabular}
\end{center}

Memory and disks are approximately random access because the electronics or
the head can be directed to a location without stepping through everything
before it. Tape is sequential because the ribbon must physically pass over the
read head in order.

\section{Non-volatility of storage}
Non-volatility is the property that stored contents survive when the power is
switched off. Storage devices are designed to be non-volatile so that files
remain available after shutdown. Primary storage such as RAM is
\emph{volatile}: its contents disappear when power is removed.

\begin{definitionbox}{Non-volatile storage}
Storage that retains its contents without power; hard disks, SSDs, optical
discs, and USB drives are non-volatile.
\end{definitionbox}

The one-line rule to carry into the exam is: RAM is volatile; storage is
non-volatile. This is exactly why unsaved work in RAM can be lost in a power
failure while a file already saved to disk remains.

\section{Exam retrieval and common confusions}
Use a two-direction recall drill. Given a device, state its technology and its
access method; given a technology or capacity, state the matching device.
\begin{itemize}
  \item HDD $\leftrightarrow$ magnetic, mechanical, random access.
  \item Magnetic tape $\leftrightarrow$ magnetic, sequential access, backup.
  \item CD / DVD / Blu-ray $\leftrightarrow$ optical, capacities 700 MB /
    4.7 GB / 25 GB (single layer).
  \item SSD / USB drive / memory card $\leftrightarrow$ flash, no moving
    parts; the SSD has the highest transfer rate.
  \item Cloud storage $\leftrightarrow$ remote servers over the Internet.
  \item RAM volatile versus storage non-volatile.
\end{itemize}

Three distinctions prevent most errors in this topic:
\begin{enumerate}
  \item Do not call the SSD magnetic or mechanical. The magnetic and
    mechanical device is the hard disk; the SSD is flash with no moving parts.
  \item Do not claim that the hard disk has the highest transfer rate; the
    SSD does (KB PYQ-27).
  \item Do not reverse the optical ordering. It is CD $<$ DVD $<$ Blu-ray
    for standard single-layer capacities, and do not confuse volatile RAM
    with non-volatile storage.
\end{enumerate}

\begin{example}
An item asks which device stores data without any moving parts. Eliminate the
hard disk (spinning platters and a moving head) and magnetic tape (a moving
ribbon). An optical drive also spins its disc. The flash-based SSD stores data
in semiconductor cells with no moving parts, so it is the answer.
\end{example}

\subsection*{Exercises}
\begin{enumerate}[label=\arabic*.]
  \item Distinguish primary, secondary, and tertiary storage and give one
    example of each.
  \item Explain why an SSD is described as having no moving parts, and why it
    is more shock-resistant than a hard disk.
  \item Order CD, DVD, and Blu-ray by single-layer capacity from smallest to
    largest, with the approximate figures.
  \item Classify RAM, a hard disk, and magnetic tape as random access or
    sequential access.
  \item A student writes: ``RAM is non-volatile, and the hard disk is the
    fastest storage device.'' Correct both errors.
\end{enumerate}

\subsection*{Solutions}
\begingroup\small
\begin{enumerate}[label=\arabic*.]
  \item Primary storage is directly accessible by the CPU (for example RAM).
    Secondary storage is non-volatile and holds data when power is off (for
    example a hard disk). Tertiary storage is archival storage loaded on
    demand (for example a robotic tape library).
  \item An SSD stores data in flash semiconductor cells, with no spinning
    platter and no moving read/write head. With no moving parts to be
    disturbed, it tolerates shock better than a hard disk and transfers data
    faster.
  \item CD $\approx$ 700 MB, then DVD $\approx$ 4.7 GB, then Blu-ray
    $\approx$ 25 GB, for single-layer media.
  \item RAM is random access; a hard disk is approximately random access;
    magnetic tape is sequential access.
  \item RAM is \emph{volatile}, not non-volatile: its contents are lost
    without power. The device with the highest transfer rate is the
    \emph{SSD}, not the hard disk; the hard disk is magnetic and mechanical
    with moving parts.
\end{enumerate}
\endgroup

\lessonheading{11}{Backup Types}
\section{Backup, storage and archive}
These three words describe different jobs, and exam questions often test the
distinction. \textbf{Storage} is where your working data lives day to day ---
the primary copy you actually open and edit. A \textbf{backup} is a separate
copy kept so that data can be recovered after loss, damage, corruption or an
attack. An \textbf{archive} is older data retained for the long term, usually
moved out of routine use. A backup exists for recovery; an archive exists for
retention. Neither is simply ``another folder'' on the same disk.

\begin{definitionbox}{Backup}
A separate copy of data kept for the purpose of recovery after loss, damage,
corruption or attack.
\end{definitionbox}

\begin{definitionbox}{Archive}
Older data kept for long-term retention, normally not used in everyday work.
\end{definitionbox}

The office desktop of the running example (KB EX-01) has one primary copy of
its documents. If the only copy is the one on the desktop's own drive, then
the user has storage but no backup. The moment a second, independent copy
exists, recovery becomes possible.

\section{Full, incremental and differential backups}
A backup type decides \emph{how much} is copied on each run. The three basic
types differ only in the answer to one question: since which earlier point are
changes included?

\begin{definitionbox}{Full backup}
Copies every selected file and folder (or the entire volume) on every run.
It is the recovery baseline for the other types.
\end{definitionbox}

\begin{definitionbox}{Incremental backup}
Copies only the files changed or added since the \emph{last backup of any
type} --- whether that was a full, an incremental or a differential backup.
\end{definitionbox}

\begin{definitionbox}{Differential backup}
Copies everything changed or added since the \emph{last full backup}, not
since the previous backup.
\end{definitionbox}

The single phrase that separates incremental from differential is therefore
the word \emph{full}. An incremental looks only as far back as the
immediately preceding backup; a differential always looks back to the last
full backup, so it keeps growing until the next full backup resets it.

\begin{center}
\begin{tabular}{p{0.20\textwidth}p{0.34\textwidth}p{0.36\textwidth}}
\toprule
\textbf{Type} & \textbf{What it copies} & \textbf{Restore requirement}\\
\midrule
Full & Everything selected, every run. & One complete set. Fastest single restore, but largest and slowest to create.\\
Incremental & Changes since the last backup of any type. & The last full backup plus every incremental backup in sequence.\\
Differential & Changes since the last full backup. & The last full backup plus only the most recent differential.\\
\bottomrule
\end{tabular}
\end{center}

There is a consistent trade-off. Smaller and faster backups produce longer or
more complex restores. A full backup is heavy to create but trivial to
restore; an incremental is light to create but may need many files chained
together at restore time; a differential sits in between.

\begin{example}
On Monday the office desktop is given a full backup. On Tuesday, a single
document changes. With an incremental strategy, Tuesday's backup contains
only that one change. With a differential strategy, Tuesday's backup contains
that change measured against Monday's full backup. If only that one file
changed, the two results look identical on Tuesday --- the difference appears
on Wednesday, when the incremental again holds only Wednesday's change, while
the differential holds Tuesday's change plus Wednesday's change.
\end{example}

\section{Restore scenarios}
Restore behaviour is the part candidates most often miss, because it is not
visible on the day of the backup --- only when recovery is needed.

\begin{example}[Restore from an incremental chain]
A full backup is taken Sunday. Incrementals are taken Monday, Tuesday and
Wednesday. On Thursday the disk fails. To recover Wednesday's data you need
the Sunday full backup and then \emph{all three} incrementals --- Monday,
Tuesday and Wednesday --- applied in order. Missing any one link loses the
changes it carried.
\end{example}

\begin{example}[Restore from a differential chain]
A full backup is taken Sunday. Differentials are taken Monday, Tuesday and
Wednesday. On Thursday the disk fails. To recover Wednesday's data you need
only the Sunday full backup and the \emph{most recent} differential ---
Wednesday's. The earlier differentials are not required, because Wednesday's
already includes all changes since Sunday.
\end{example}

Both scenarios produce the same recovered data. The difference is operational:
the incremental restore is longer and depends on every tape or file in the
chain; the differential restore is shorter and depends on only two sets.

\section{Image and file-level backups}
Beyond the three time-based types, two coverage-based terms appear in exams.

\begin{definitionbox}{Image backup (system image)}
A complete clone of a system --- operating system, settings, applications and
data --- that can restore a whole machine.
\end{definitionbox}

\begin{definitionbox}{File-level backup}
A backup of selected files and folders rather than a whole system.
\end{definitionbox}

An image restores a working machine in one step, which makes it ideal after a
drive failure or a corrupted operating system. A file-level backup is smaller
and more flexible for everyday documents, but recovering a machine from it
means reinstalling the operating system first and then restoring the files.
The office desktop could use an image for disaster recovery and file-level
backups for its working documents.

\section{The 3-2-1 rule, off-site and offline copies}
The \textbf{3-2-1 rule} is the standard exam reminder for resilient backup:
\begin{enumerate}
  \item \textbf{3} copies of the data --- the original plus two backups.
  \item On \textbf{2} different types of media (for example, an external disk
    and cloud storage).
  \item With \textbf{1} copy kept \textbf{off-site}, at a different location.
\end{enumerate}
The numbers map to copies, media and location, in that order. A common trap
is to swap them, or to read the ``2'' as two off-site copies.

An \textbf{off-site} copy is kept at a different physical place, so a single
fire, flood or theft cannot destroy every copy at once. An \textbf{offline}
copy is disconnected from the network; it is especially useful against
ransomware, which spreads through connected drives. A \textbf{cloud backup}
is an off-site copy held on remote servers. These are not mutually exclusive:
the office desktop's cloud backup is both off-site and, when not mounted, can
serve as an offline copy as well.

\begin{example}
The office desktop keeps its working files on the local SSD (copy 1), an
external USB disk (copy 2, different medium, kept offline), and a cloud
account (copy 3, off-site). This satisfies the 3-2-1 rule: three copies, two
media types, one off-site.
\end{example}

\section{Why RAID is not a backup}
RAID (\textbf{Redundant Array of Independent Disks}) combines several disks
for performance or redundancy. If one disk fails, a mirrored or parity RAID
keeps the system running. That protects \emph{uptime}, not \emph{recovery}.
RAID is not a backup because it does not protect against the things a backup
exists for: accidental deletion, corruption or ransomware. When a file is
deleted, or an encrypted copy is written by ransomware, the change is written
to every member of the RAID array at once. There is no earlier clean copy to
go back to. Redundancy duplicates the current state --- including the damage.

\section{Backup rotation}
\textbf{Backup rotation} means keeping several dated backup sets and reusing
them on a fixed schedule, commonly the grandfather--father--son scheme (daily,
weekly, monthly). Rotation matters because a problem such as silent
corruption may not be noticed for days; if the single backup set has already
been overwritten with the corrupted data, there is no clean copy left.
Keeping multiple versions across time ensures that an older, clean copy
remains available.

\subsection*{Exercises}
\begin{enumerate}[label=\arabic*.]
  \item Define full, incremental and differential backup in one line each.
  \item A full backup runs Sunday; differentials run Monday to Wednesday. The
    disk fails on Thursday. Which sets are needed to restore Wednesday's
    data?
  \item State the 3-2-1 rule and explain what each number refers to.
  \item Why is RAID not a substitute for a backup? Give one concrete danger.
  \item Explain the purpose of backup rotation.
\end{enumerate}

\subsection*{Solutions}
\begingroup\small
\begin{enumerate}[label=\arabic*.]
  \item Full: copies everything selected each run. Incremental: copies changes
    since the last backup of any type. Differential: copies changes since the
    last full backup.
  \item The Sunday full backup and the most recent (Wednesday) differential
    only. The earlier differentials are unnecessary because Wednesday's
    includes all changes since the full backup.
  \item Three copies of the data, on two different types of media, with one
    copy kept off-site.
  \item RAID provides redundancy for uptime, not recovery. A deletion,
    corruption or ransomware write is applied to every RAID member at once,
    so no earlier clean copy survives.
  \item Rotation retains several dated backup versions so that an older, clean
    copy is still available if corruption is discovered only after the newest
    backup has already been overwritten.
\end{enumerate}
\endgroup

\unitopener{4}{I/O, Peripherals \& Multimedia}
\lessonheading{12}{Input Devices}
\section{What is an input device}
A computer works in a cycle: input, process, and output. In the input step,
data and instructions are fed into the machine. An \textbf{input device} is
the hardware that accepts data or control signals from the user or from the
environment and passes them to the computer, converting a human action or a
physical signal into machine-readable form. Without an input device the
computer would have nothing to process.

\begin{definitionbox}{Input device}
Hardware that accepts data or control signals from the user or the
environment and supplies them to the CPU, converting them into a form the
computer can process.
\end{definitionbox}

Common input devices include the keyboard, mouse, touchscreen, scanner,
microphone, digital camera, barcode reader, and biometric sensor. The
following table fixes the meaning of each term used in this lesson.

\begin{center}
\begin{tabular}{p{0.26\textwidth}p{0.66\textwidth}}
\toprule
\textbf{Term} & \textbf{Meaning in this lesson} \\
\midrule
Keyboard & Text input device that converts a keystroke into a character code. \\
Mouse & Pointing device moved on a surface to control the on-screen pointer. \\
Trackball & Pointing device whose ball is rotated in place; the device stays still. \\
Touchpad & Flat, touch-sensitive pad on a laptop; the pointer moves on the screen. \\
Joystick & Lever that moves on two or three axes, used for games and control. \\
Light pen & Pen-shaped input device with a light sensor, used on a CRT screen. \\
Stylus / digitizer & Pen plus tablet that converts strokes into digital coordinates. \\
Scanner & Converts a paper document or picture into a digital image. \\
Biometrics & Input from a physical or behavioural human trait (fingerprint, iris, face). \\
\bottomrule
\end{tabular}
\end{center}

\section{Classifying input devices}
Input devices are grouped by the type of data they capture. This is the
classification JKSSB question-setters use most often, so learn one example
for each row.

\begin{center}
\begin{tabular}{p{0.26\textwidth}p{0.60\textwidth}}
\toprule
\textbf{Category} & \textbf{Typical devices} \\
\midrule
Text input & Keyboard \\
Pointing input & Mouse, trackball, touchpad, joystick, light pen, stylus, touchscreen \\
Audio input & Microphone \\
Image / video input & Scanner, digital camera, webcam \\
Biometric input & Fingerprint, iris, and face sensors \\
Code / tag input & Barcode, QR code, RFID \\
\bottomrule
\end{tabular}
\end{center}

A single device may belong to more than one row if it captures more than one
kind of data, but for exam purposes use the primary role: a touchscreen is a
pointing input device, a webcam is an image/video input device, and a
microphone is an audio input device.

\section{Direct vs indirect input}
A second useful classification is whether the data is captured at its source
or encoded by the user.

\textbf{Direct input devices} capture the data themselves. Examples are the
touchscreen, light pen, scanner, barcode reader, biometric sensor, and
microphone: the device reads the object, code, trait, or sound and turns it
directly into computer data.

\textbf{Indirect input devices} require the user to act on the device to
encode the data. Examples are the keyboard, mouse, trackball, touchpad, and
joystick: the user translates an intention into keystrokes or pointer
movements, which the device then sends to the computer.

\begin{definitionbox}{Direct input}
Input in which the device captures the data at its source (for example, a
scanner reading a page or a biometric sensor reading a fingerprint) without
the user encoding it by hand.
\end{definitionbox}

The distinction matters because it explains why some devices are fast and
error-free while others depend on the user. A barcode reader or a biometric
sensor captures the data exactly as it is, so there is no typing mistake; a
keyboard, being indirect, is only as accurate as the person typing. In a
question, if the stem stresses that the device \emph{reads} or \emph{captures}
an object, code, trait, or sound, think of a direct input device; if it
stresses that the \emph{user types or moves} something to enter the data,
think of an indirect input device.

\section{The keyboard: text input}
The keyboard is the most common text input device. Its standard letter
layout is called \textbf{QWERTY}, after the first six letters of the top row.
A keyboard has alphanumeric keys, function keys (F1 to F12), modifier keys
such as Ctrl, Alt, and Shift, and usually a numeric keypad. Each time a key
is pressed, the keyboard sends the corresponding character code to the
computer. The keyboard is an indirect input device: the user types the data,
and the device encodes it.

Keyboards differ in construction even though their role is the same. A
\textbf{membrane keyboard} uses a flexible membrane under the keys and is
cheap and quiet. A \textbf{mechanical keyboard} uses a separate physical
switch under each key and gives a firmer, more definite feel. A
\textbf{wireless keyboard} connects to the computer by Bluetooth or by a USB
receiver instead of a cable. Whatever the construction, a keyboard is still a
text input device: it sends character codes for the keys that are pressed.

\section{Pointing devices}
A pointing device controls the position of an on-screen pointer (cursor) and
is used to point, click, and drag. The table below fixes each pointing device
and its nearest neighbour, because JKSSB questions often ask candidates to
tell two similar pointing devices apart.

\begin{center}
\begin{tabular}{p{0.16\textwidth}p{0.38\textwidth}p{0.34\textwidth}}
\toprule
\textbf{Device} & \textbf{How it works} & \textbf{Nearest neighbour} \\
\midrule
Mouse & The device slides on a surface and the pointer follows. & Trackball: the ball rotates while the device stays still. \\
Trackball & The ball is rotated in place; the device is stationary. & Mouse: the whole device moves. \\
Touchpad & A flat finger pad moves the pointer on the screen (indirect). & Touchscreen: the display itself is touched (direct). \\
Joystick & A lever moves on two or three axes for games and control. & Mouse: flat-surface pointing. \\
Light pen & A light sensor is pointed at a CRT screen. & Stylus: used on a digitizer surface. \\
Stylus + digitizer & A pen stroke is converted into digital coordinates. & Light pen: uses a light sensor, not a tablet. \\
\bottomrule
\end{tabular}
\end{center}

\subsection{The mouse}
The \textbf{mouse} is a hand-held pointing device moved on a flat surface. It
has one or more buttons and usually a scroll wheel. An \textbf{optical or
laser mouse} tracks movement with light, which is the common modern design;
an older \textbf{mechanical mouse} used a rolling ball. A mouse is an input
device (PYQ-08).

\subsection{The trackball}
The \textbf{trackball} is, in effect, an upside-down mouse: the ball is on
top and is rotated with the fingers while the device itself stays still. This
saves desk space because the device does not have to be moved. The
nearest-neighbour trap is that a trackball and a mouse both point, but with a
mouse the \emph{device} slides, while with a trackball the \emph{ball}
rotates in place.

\subsection{The trackpad or touchpad}
The \textbf{trackpad} (also called the touchpad) is a flat, touch-sensitive
pad found on laptops. The user slides a finger across the pad to move the
pointer on the screen and taps to click. It is an indirect device because the
finger does not touch the object on the screen; the pad only controls the
pointer.

\subsection{The joystick}
The \textbf{joystick} is a lever that can be moved on two or three axes. It is
used mainly for games and for some industrial control tasks.

\subsection{The light pen}
The \textbf{light pen} is a pen-shaped input device with a light sensor at its
tip. It is pointed at a CRT screen to select items or to draw. A light pen is
an input device, not an output device (TRAP-12).

\section{Stylus and digitizer tablet}
A \textbf{stylus} is a pen-like tool used to write or draw on a tablet or a
touch screen. A \textbf{digitizer} (graphics tablet) is a flat surface that
converts the movement of the stylus into digital coordinates. Together they
are used for freehand drawing, capturing signatures, computer-aided design
(CAD), and handwriting recognition. Because a digitizer captures a drawing as
data, it is an input device (TRAP-12). Do not confuse the stylus with the
light pen: the light pen works with a light sensor, while the stylus works
with a digitizer surface.

\section{Touchscreen}
A \textbf{touchscreen} is a display that also works as an input device. The
user touches the screen directly to select and enter data. Because the item
touched is the item selected, a touchscreen is a direct input device. The two
common technologies are \textbf{resistive} and \textbf{capacitive}
touchscreens. Touchscreens are used in smartphones, tablets, ATMs, and
information kiosks.

The nearest-neighbour trap is the touchscreen versus the touchpad. On a
touchscreen the user touches the display itself, so the input is direct. On a
touchpad the user touches a separate pad and the pointer moves on the screen,
so the input is indirect. A touchpad is not a touchscreen.

\section{Scanners, cameras, and webcams}
A \textbf{scanner} converts a paper document or a picture into a digital
image. Common types are the \textbf{flatbed} scanner (the page is placed on a
glass surface), the sheet-fed scanner, the handheld scanner, and the drum
scanner. Software such as OCR (Optical Character Recognition) or OMR (Optical
Mark Recognition) can then read text or marks from the scanned image. A
scanner is an input device (TRAP-12).

A \textbf{digital camera} captures photographs as digital images and stores
them on a memory card. A \textbf{webcam} is a small digital camera connected
to a computer and used mainly for live video, such as video calls and
streaming. Both are image/video input devices.

\section{Barcode, QR code, and RFID}
A \textbf{barcode} is a one-dimensional pattern of parallel lines of varying
width, read by a barcode scanner. It is used on product labels and on book
ISBNs.

A \textbf{QR code} (Quick Response code) is a two-dimensional matrix of black
and white squares. It can store more data than a barcode and can be read from
any direction. The full form of QR is \textbf{Quick Response}.

\textbf{RFID} stands for \textbf{Radio Frequency Identification}. An RFID
system has a tag and a reader; the two exchange radio waves, so no line of
sight is needed, unlike a barcode. RFID is used in toll collection, access
cards, and inventory tracking. A barcode, a QR code, and an RFID tag are all
read by an input device that feeds the data into the computer.

\section{Microphone and biometric input}
A \textbf{microphone} is an audio input device. It converts sound, such as
voice or music, into an electrical signal that the computer digitises. It is
used for recording, voice calls, and speech recognition. The matching output
device is the \textbf{speaker}: a microphone takes audio in, a speaker sends
audio out.

\textbf{Biometric} input identifies a person from a physical or behavioural
trait. Common traits are the \textbf{fingerprint}, the \textbf{iris}, the
\textbf{face}, the voice, and the palm. A sensor captures the trait and feeds
the measured data to the computer. Biometrics are used for attendance
systems, device unlocking, and authentication. Because the sensor reads a
human trait and passes the data to the system, biometric input is a direct
input form.

\begin{example}
A student wants to scan a printed page into the computer. The correct device
is a scanner, which converts the page into a digital image. If the student
instead wanted to type the same text by hand, the device would be a keyboard,
which is an indirect text input device. If the goal were to capture a live
video of the student, the device would be a webcam.
\end{example}

\section{Summary of input devices}
The table below gathers the devices in this lesson for quick revision. Use it
as a checklist: for each device, be ready to state its category and its
one-line function, and to name the device it is most often confused with.

\begin{center}
\begin{tabular}{p{0.24\textwidth}p{0.24\textwidth}p{0.42\textwidth}}
\toprule
\textbf{Device} & \textbf{Category} & \textbf{Function} \\
\midrule
Keyboard & Text & Converts a keystroke into a character code. \\
Mouse & Pointing & Moves the on-screen pointer; point, click, drag. \\
Trackball & Pointing & Ball rotated in place; device stays still. \\
Touchpad & Pointing & Flat finger pad; pointer moves on the screen. \\
Joystick & Pointing & Lever on two or three axes for games/control. \\
Light pen & Pointing & Light sensor pointed at a CRT screen. \\
Stylus + digitizer & Pointing & Converts pen strokes into digital coordinates. \\
Touchscreen & Pointing (direct) & Touch the display to select and enter. \\
Scanner & Image & Converts a document or picture into a digital image. \\
Digital camera & Image/video & Captures photographs as digital images. \\
Webcam & Image/video & Small camera for live video on a computer. \\
Microphone & Audio & Converts sound into an electrical signal. \\
Biometric sensor & Biometric & Reads a fingerprint, iris, or face trait. \\
Barcode / QR code & Code/tag & Read optical codes that identify an item. \\
RFID & Code/tag & Reads a tag by radio waves; no line of sight. \\
\bottomrule
\end{tabular}
\end{center}

\section{Points of confusion between similar devices}
Beyond classifying a device, JKSSB items often test a single distinction
between two devices that look or sound alike. The table below collects the
pairs used in this lesson.

\begin{center}
\begin{tabular}{p{0.30\textwidth}p{0.60\textwidth}}
\toprule
\textbf{Pair} & \textbf{The distinction to state} \\
\midrule
Mouse vs trackball & Mouse: the device slides. Trackball: the ball rotates, device still. \\
Touchscreen vs touchpad & Touchscreen: touch the display (direct). Touchpad: separate pad moves the pointer (indirect). \\
Light pen vs stylus & Light pen: light sensor on a CRT screen. Stylus: used with a digitizer tablet. \\
Scanner vs digital camera & Scanner: converts a flat document or picture into an image. Camera: captures a scene. \\
Barcode vs QR code & Barcode: 1D parallel lines. QR code: 2D matrix, holds more data. \\
Barcode/QR vs RFID & Barcode/QR are read optically and need line of sight; RFID uses radio waves. \\
Microphone vs speaker & Microphone: audio input. Speaker: audio output. \\
Data glove vs haptic glove & Data glove: captures hand movement (input). Haptic glove: gives force feedback (output). \\
\bottomrule
\end{tabular}
\end{center}

\section{Exam retrieval and common confusions}
Six distinctions prevent most errors on this topic.
\begin{enumerate}
  \item A light pen is an \textbf{input} device; a plotter is an \textbf{output}
    device. Do not swap them (TRAP-12).
  \item A punch-card reader, a digitizer tablet, and a scanner are all
    \textbf{input} devices. A drum printer and a sound card are
    \textbf{output} devices (TRAP-12).
  \item A \textbf{trackball} rotates its ball in place; a \textbf{mouse}
    slides over a surface.
  \item A \textbf{touchscreen} is touched directly; a \textbf{touchpad} is a
    separate pad that moves the pointer indirectly.
  \item \textbf{RFID} uses radio waves and needs no line of sight; a
    \textbf{barcode} and a \textbf{QR code} are read optically.
  \item \textbf{Haptic gloves}, a \textbf{tactile display}, and a
    \textbf{Voice Response Unit (VRU)} give feedback to the user and are
    \textbf{output}; a \textbf{data glove} captures hand movement and is
    \textbf{input} (PYQ-09).
\end{enumerate}

\begin{example}
An item lists four devices --- light pen, punch-card reader, plotter, and
scanner --- and asks which are input devices. Classify each one: the light
pen is input, the punch-card reader is input, the plotter is output, and the
scanner is input. The answer is the light pen, the punch-card reader, and the
scanner. The plotter is the only output device in the list.
\end{example}

\subsection*{Exercises}
\begin{enumerate}[label=\arabic*.]
  \item Define an input device and give three examples.
  \item Classify these devices by the type of data they capture: keyboard,
    microphone, scanner, fingerprint sensor.
  \item Distinguish direct input from indirect input, giving one example of
    each.
  \item Explain the difference between a trackball and a mouse, and between a
    touchscreen and a touchpad.
  \item Differentiate a barcode, a QR code, and RFID.
  \item State whether each is input or output: light pen, plotter, digitizer
    tablet, drum printer, sound card.
\end{enumerate}

\subsection*{Solutions}
\begingroup\small
\begin{enumerate}[label=\arabic*.]
  \item An input device accepts data or control signals from the user or the
    environment and passes them to the CPU in machine-readable form.
    Examples: keyboard, mouse, scanner (any three valid devices).
  \item Keyboard --- text input; microphone --- audio input; scanner ---
    image input; fingerprint sensor --- biometric input.
  \item Direct input captures the data at its source, such as a scanner or a
    biometric sensor. Indirect input needs the user to encode the data, such
    as a keyboard or a mouse.
  \item With a mouse, the device slides over a surface; with a trackball, the
    ball is rotated in place and the device stays still. A touchscreen is
    touched directly on the display; a touchpad is a separate pad that moves
    the pointer indirectly.
  \item A barcode is a 1D pattern of lines read optically; a QR code is a 2D
    matrix that holds more data; RFID uses radio waves between a tag and a
    reader and needs no line of sight.
  \item Light pen --- input; plotter --- output; digitizer tablet --- input;
    drum printer --- output; sound card --- output.
\end{enumerate}
\endgroup

\lessonheading{13}{OCR, OMR, MICR}
\section{Reading machines and why they matter}
A computer works only with data that has already been converted into
machine-readable form. When the data exists on paper --- as printed text, as
a filled-in bubble on a form, or as a line of characters on a cheque --- the
machine needs a way to read it. A \textbf{reading machine} is the device or
technology that captures this physical information and turns it into digital
data the system can process. The physical reader (often a scanner, sometimes
a dedicated reader) performs the capture; the \emph{recognition method}
decides what the captured pattern means. This lesson covers the three most
examined recognition methods --- OCR, OMR, and MICR --- together with the
code-and-tag devices (barcode, QR code, RFID) and biometric input, because
questions frequently ask candidates to tell these near-neighbours apart.

The single most important idea is the \emph{read-target} of each method.
\textbf{OCR} reads \emph{text}. \textbf{OMR} reads \emph{marks}.
\textbf{MICR} reads \emph{characters printed in magnetic ink}. Almost every
item on this topic can be solved by fixing this one distinction in memory.

All the devices in this lesson are \emph{input} devices or reading
\emph{technologies}: they bring data into the system. Most of them are also
\emph{direct} input, because the device reads the object, mark, code, or
trait itself, instead of asking the user to type it. This is why a reading
machine is faster and less error-prone than a keyboard for a repetitive task
such as reading a thousand answer sheets or a thousand cheques.

\begin{definitionbox}{Reading machine}
A device or technology that captures printed characters, marks, codes, or
tags from a physical medium and converts them into machine-readable data for
the computer to process.
\end{definitionbox}

\section{Terminology at a glance}
The table below fixes the meaning of every term used in this lesson. Read it
down once; then, for each term, be ready to state the read-target and one
place it is used.

\begin{center}
\begin{tabular}{p{0.24\textwidth}p{0.66\textwidth}}
\toprule
\textbf{Term} & \textbf{Meaning in this lesson} \\
\midrule
OCR & Optical Character Recognition --- converts an image of printed or written text into editable text. \\
OMR & Optical Mark Recognition --- detects the presence or absence of marks at fixed positions on a form. \\
MICR & Magnetic Ink Character Recognition --- reads characters printed in magnetic ink, used on bank cheques. \\
Barcode & A one-dimensional pattern of parallel lines read optically to identify an item. \\
QR code & A two-dimensional matrix of squares (Quick Response) that holds more data than a barcode. \\
RFID & Radio Frequency Identification --- a tag and reader exchange radio waves; no line of sight is needed. \\
Biometrics & Input that identifies a person from a physical or behavioural trait (fingerprint, iris, face). \\
\bottomrule
\end{tabular}
\end{center}

\section{OCR --- Optical Character Recognition}
\textbf{OCR} stands for \textbf{Optical Character Recognition}. It takes a
scanned image of \emph{printed or handwritten text} and converts it into
machine-encoded, \emph{editable text}. The key word is
\emph{character}: OCR recognises individual letters, digits, and symbols, so
the output is text that a computer can search, copy, and edit. This is what
makes OCR different from OMR, which recognises only whether a position has
been marked and never identifies a letter.

OCR is used to digitise the printed word. A book, a bundle of invoices, a
passport page, or an old form can be scanned into an image, and OCR software
then extracts the text from that image. A useful way to remember the pairing
is: the \textbf{scanner} is the hardware that creates the image, while
\textbf{OCR} is the technology that turns the image of text into actual,
selectable text.

\begin{definitionbox}{OCR}
A recognition technology that reads characters from a scanned image of
printed or handwritten text and converts them into machine-encoded,
editable text.
\end{definitionbox}

OCR works best on clean, clearly printed text in a known font; its accuracy
drops when the page is skewed, smudged, or handwritten. When the input is
handwritten rather than printed, the same idea is sometimes called
\emph{Intelligent Character Recognition} (ICR), but for exam purposes both
count as OCR-style reading of text. OCR is \emph{not} a storage device and
\emph{not} a printer: it only converts an existing image into text.

\begin{example}
A student scans a page of a textbook into the computer. The scanner produces
an \emph{image}, so the words on the page cannot yet be copied or searched.
Running OCR on that image produces a text file in which every word can be
selected and edited. The scanner supplied the image; OCR supplied the text.
\end{example}

\section{OMR --- Optical Mark Recognition}
\textbf{OMR} stands for \textbf{Optical Mark Recognition}. It reads a form
by detecting the \emph{presence or absence of a mark} at fixed positions.
Typically the candidate fills a small bubble or shades a box with a pencil or
pen, and the OMR reader checks which positions have been darkened. Crucially,
OMR reads \emph{marks}, not characters: it does not know what a letter
``says'', only that a particular position has been filled.

Because it detects marks at fixed coordinates, OMR is ideal for forms where
the answer is one of a small set of choices. Its classic use is the
\textbf{MCQ answer sheet}: the machine reads which bubble has been shaded for
each question. It is also used for questionnaires, surveys, census forms, and
ballots. A very common exam trap is to confuse OMR with OCR; remember that
OMR works on \emph{marks}, while OCR works on \emph{text}.

\begin{definitionbox}{OMR}
A recognition technology that detects the presence or absence of marks at
fixed positions on a form, used mainly to read objective (MCQ) answer
sheets and survey forms.
\end{definitionbox}

OMR is extremely fast because it does not have to interpret a shape as a
character; it only has to decide whether a small region is dark or light.
That speed is why OMR is the standard method for large-scale examinations
where millions of answer sheets must be processed quickly. Its natural limit
is that it can record only \emph{which} option was marked, not any written
answer: if the candidate writes a sentence instead of shading a bubble, OMR
cannot read it.

\section{MICR --- Magnetic Ink Character Recognition}
\textbf{MICR} stands for \textbf{Magnetic Ink Character Recognition}. It
reads characters printed in a special \emph{magnetic ink} that contains iron
oxide. The characters are printed in a standard font along the bottom of a
\emph{bank cheque} as the \textbf{MICR line}, which normally encodes the
cheque number, the bank and branch code, and the account number. A MICR
reader magnetises the ink and detects the shapes of the characters as they
pass the read head.

Banking is the defining application of MICR. Cheque processing must be very
fast and highly secure, and MICR satisfies both needs: the reading is quick,
and because the ink is magnetic it is difficult to forge or to alter without
detection. When a question mentions a \emph{bank cheque} or \emph{magnetic
ink}, the answer is almost always MICR.

\begin{definitionbox}{MICR}
A recognition technology that reads characters printed in magnetic ink,
used mainly in banking to process the MICR line printed at the bottom of a
cheque.
\end{definitionbox}

Unlike ordinary ink, magnetic ink can be read even if the cheque is slightly
crumpled or over-stamped, because the reader senses the magnetic field of
the ink rather than a purely visual image. This ruggedness and the difficulty
of counterfeiting magnetic ink are exactly the reasons banks adopted MICR
for cheque clearing.

\section{The crucial trio: OCR vs OMR vs MICR}
These three names are the heart of the lesson, and JKSSB questions routinely
swap their meanings in the options. The table below is the single most useful
thing to memorise.

\begin{center}
\begin{tabular}{p{0.12\textwidth}p{0.26\textwidth}p{0.50\textwidth}}
\toprule
\textbf{Method} & \textbf{What it reads} & \textbf{Where it is used} \\
\midrule
OCR & Printed or written characters & Digitising books, documents, and passports \\
OMR & Marks on a form & MCQ answer sheets, surveys, ballots \\
MICR & Characters in magnetic ink & Bank cheques and cheque clearing \\
\bottomrule
\end{tabular}
\end{center}

The distinction to hold is: \textbf{MICR} uses \emph{magnetic ink};
\textbf{OCR} reads \emph{printed text}; \textbf{OMR} reads \emph{marks}. This
is the JKSSB trap TRAP-25, and the real paper item behind it (PYQ-30) asks
simply that magnetic ink means MICR.

It also helps to picture \emph{how} each one captures its data. OMR and OCR
both read an optical (light) signal: the reader sees the paper. MICR is the
odd one out, because its reader senses a \emph{magnetic} signal from the ink.
So a quick way to separate them in a hurry is: two are \emph{optical}, one is
\emph{magnetic}.

\begin{center}
\begin{tabular}{p{0.11\textwidth}p{0.20\textwidth}p{0.25\textwidth}p{0.31\textwidth}}
\toprule
\textbf{Method} & \textbf{Signal read} & \textbf{Input pattern} & \textbf{Cannot do} \\
\midrule
OCR & Optical (light) & Text characters & Cannot read marks alone \\
OMR & Optical (light) & Marks / shaded bubbles & Cannot read written characters \\
MICR & Magnetic & Magnetic-ink characters & Not used for forms or plain text \\
\bottomrule
\end{tabular}
\end{center}

\section{Barcodes and QR codes}
A \textbf{barcode} is a one-dimensional (1D) pattern of parallel lines of
varying width and spacing. A barcode scanner shines a light beam across the
lines and reads the reflected pattern, which represents a number or code
that identifies a product. Because the scanner must see the lines, a barcode
requires \emph{line of sight}. Barcodes are used on product labels, on the
ISBN of books, and on postal items.

A \textbf{QR code} --- \textbf{Quick Response} code --- is a
\emph{two-dimensional} (2D) matrix of black and white squares. It stores
more data than a barcode and can be read from any direction. Because it can
hold a URL, scanning a QR code can open a website or an app directly. Both a
barcode and a QR code are read \emph{optically}, so the nearest-neighbour
trap is simply that a QR code is two-dimensional and carries more
information, while a barcode is one-dimensional.

\begin{definitionbox}{QR code}
A two-dimensional matrix of black and white squares (Quick Response code)
that stores more data than a barcode and can be read from any direction.
\end{definitionbox}

\begin{example}
A shop assistant passes a product over a laser at the billing counter: the
scanner reads the 1D \emph{barcode} on the packet and looks up its price. A
customer points a phone camera at the square \emph{QR code} on a poster: the
phone reads the 2D code and opens the shop's website. The barcode carried a
short number; the QR code carried a whole web address.
\end{example}

\section{RFID --- Radio Frequency Identification}
\textbf{RFID} stands for \textbf{Radio Frequency Identification}. An RFID
system has two parts: a \emph{tag} attached to the object and a \emph{reader}
that interrogates it. The two exchange \emph{radio waves}, so no line of
sight is required, unlike a barcode. A passive tag carries no battery and is
powered by the reader's signal; an active tag has its own battery and can
broadcast over a longer range.

RFID is used where the tag may be out of sight or where many tags must be
read quickly: toll collection (for example FASTag), access-control cards,
passports, library books, and inventory tracking. The exam contrast is
always barcode/QR versus RFID: barcode and QR are read \emph{optically} and
need line of sight, while RFID uses \emph{radio waves} and can read a tag
without it being visible.

\section{Biometric input}
\textbf{Biometrics} identifies a person from a physical or behavioural
trait. The common physical traits are the \textbf{fingerprint}, the
\textbf{iris}, the \textbf{face}, and the palm; the voice is a common
behavioural trait. A sensor captures the trait and feeds the measured data to
the computer as \emph{input}; the system then compares it with the stored
template. Biometric input is a \emph{direct} form of input because the sensor
reads the trait itself rather than asking the user to type or encode it.

Biometrics are used for attendance systems, device unlocking, and
authentication, and they appear in national identity and benefit systems. A
useful checklist item for the exam is that a fingerprint reader, an iris
scanner, and a face-recognition camera are all \emph{input} devices; a device
that merely prints or displays the result is \emph{output}. Do not confuse
biometric traits (fingerprint, iris, face) with the optical reading machines:
a fingerprint sensor reads a part of the human body, not a paper code.

\section{Full-form ladder and memory hooks}
Direct-recall items on this topic usually ask for a full form, so learn the
expansions exactly. The same expansions appear on the slides; repeat them
until they can be produced without hesitation.

\begin{center}
\begin{tabular}{p{0.16\textwidth}p{0.72\textwidth}}
\toprule
\textbf{Short form} & \textbf{Full form} \\
\midrule
OCR & Optical Character Recognition \\
OMR & Optical Mark Recognition \\
MICR & Magnetic Ink Character Recognition \\
QR & Quick Response \\
RFID & Radio Frequency Identification \\
\bottomrule
\end{tabular}
\end{center}

Three one-line hooks carry most of the content of this lesson. First,
\emph{OCR reads text, OMR reads marks, MICR reads magnetic ink}. Second,
\emph{a barcode is one-dimensional, a QR code is two-dimensional, and RFID
uses radio waves}. Third, \emph{a biometric sensor reads a human trait ---
fingerprint, iris, face, or voice}. If a question seems unfamiliar, reduce
it to one of these three sentences and the answer usually becomes clear.

\section{Applications at a glance}
The table gathers the reading machines and code devices with their defining
application, for quick revision.

\begin{center}
\begin{tabular}{p{0.20\textwidth}p{0.66\textwidth}}
\toprule
\textbf{Technology} & \textbf{Defining application} \\
\midrule
OCR & Converting scanned printed text into editable text \\
OMR & Reading marks on objective (MCQ) answer sheets \\
MICR & Reading the magnetic-ink line at the bottom of a bank cheque \\
Barcode & Reading 1D product labels and book ISBNs \\
QR code & Reading 2D codes that can open a link or store more data \\
RFID & Reading tags by radio waves (tolls, access cards, inventory) \\
Biometrics & Identifying a person from a fingerprint, iris, face, or voice \\
\bottomrule
\end{tabular}
\end{center}

\section{Exam retrieval and common confusions}
Six distinctions prevent most errors on this topic.
\begin{enumerate}
  \item \textbf{MICR} uses \textbf{magnetic ink}; \textbf{OCR} reads
    \textbf{printed text}; \textbf{OMR} reads \textbf{marks}. Do not
    interchange them (TRAP-25, PYQ-30).
  \item A real PYQ asks which technology uses magnetic ink on cheques: the
    answer is \textbf{MICR} (PYQ-30).
  \item \textbf{OMR}, not OCR, is used for \textbf{MCQ answer sheets},
    because it reads marks at fixed positions.
  \item A \textbf{QR code} is two-dimensional and holds more data than a
    one-dimensional \textbf{barcode}; both are read optically.
  \item \textbf{RFID} uses \textbf{radio waves} and needs no line of sight;
    a barcode or QR code must be \emph{seen} by the reader.
  \item A \textbf{fingerprint} or \textbf{iris} scanner is \textbf{input};
    biometrics identify a person from a physical or behavioural trait.
\end{enumerate}

\begin{example}
An item states that a bank reads the characters at the bottom of a cheque
using OCR. This is wrong. The cheque characters are printed in magnetic ink,
so the technology is \textbf{MICR}. OCR would be used to read ordinary
\emph{printed text} from a scanned document, and OMR would be used to read
\emph{marks} on a form --- neither of which matches a cheque's magnetic-ink
line.
\end{example}

\begin{example}
An item lists four reading methods and asks which two are \emph{optical}.
OCR and OMR both use a light signal to read the paper, so they are the two
optical methods. MICR is \emph{magnetic}, not optical, and RFID is neither
optical nor magnetic --- it uses radio waves. Separating the methods by their
signal (optical, magnetic, radio) is a fast way to answer such a question.
\end{example}

\subsection*{Exercises}
\begin{enumerate}[label=\arabic*.]
  \item Expand OCR, OMR, and MICR, and state in one line what each reads.
  \item Which technology would read an objective (MCQ) answer sheet, and
    why?
  \item Explain why a bank uses MICR rather than OCR for cheque processing.
  \item Distinguish a barcode from a QR code on two points.
  \item State one way in which RFID differs from a barcode.
  \item Give three examples of biometric traits and state the purpose for
    which biometric input is commonly used.
  \item Classify these as optical, magnetic, or radio reading: OMR, MICR,
    RFID, OCR.
\end{enumerate}

\subsection*{Solutions}
\begingroup\small
\begin{enumerate}[label=\arabic*.]
  \item OCR = Optical Character Recognition, which reads printed or written
    characters; OMR = Optical Mark Recognition, which reads marks on a form;
    MICR = Magnetic Ink Character Recognition, which reads characters in
    magnetic ink.
  \item OMR, because it detects the presence or absence of a mark at fixed
    positions, which is exactly how a candidate fills bubbles on an MCQ
    answer sheet.
  \item The cheque characters are printed in magnetic ink, which MICR reads
    directly; the reading is fast, and the magnetic ink is difficult to
    forge or alter. OCR reads ordinary printed text, not magnetic-ink
    characters.
  \item A barcode is a one-dimensional pattern of parallel lines and holds
    less data; a QR code is a two-dimensional matrix of squares, holds more
    data, and can be read from any direction.
  \item RFID exchanges radio waves between a tag and a reader and needs no
    line of sight, whereas a barcode is read optically and must be visible
    to the scanner. RFID can also read many tags at once.
  \item Fingerprint, iris, and face (voice and palm are also valid).
    Biometric input is commonly used for attendance, authentication, and
    device unlocking.
  \item OMR and OCR are \emph{optical}; MICR is \emph{magnetic}; RFID is
    \emph{radio}.
\end{enumerate}
\endgroup

\lessonheading{14}{Output Devices}
\section{What is an output device}
A computer works in a cycle: input, process, and output. In the output step,
the result of processing is handed back to the user or to another device. An
\textbf{output device} is the hardware that receives processed data from the
CPU and presents it in a form a human can see, hear, or read. Without an
output device the result of the computation would stay locked inside the
machine.

\begin{definitionbox}{Output device}
Hardware that receives processed data from the CPU and presents it to the
user or to another device in a readable form, such as on a screen, through a
speaker, or on paper.
\end{definitionbox}

Common output devices include the monitor, projector, speaker, headphones,
printer, and plotter. The following table fixes the meaning of each term used
in this lesson.

\begin{center}
\begin{tabular}{p{0.24\textwidth}p{0.66\textwidth}}
\toprule
\textbf{Term} & \textbf{Meaning in this lesson} \\
\midrule
Monitor & The main visual output device; shows text, images, and video as soft copy. \\
Projector & Throws a large image onto a screen or wall as soft copy. \\
Speaker & Converts an electrical signal into sound for a group of listeners. \\
Headphones & Produce sound for a single listener; worn on or over the ears. \\
Printer & Produces hard copy on paper; impact or non-impact. \\
Plotter & Draws large, precise line drawings such as engineering designs. \\
MFP & Multifunction Printer; combines printing, scanning, copying, and often faxing. \\
\bottomrule
\end{tabular}
\end{center}

\section{Classifying output devices}
Output devices are grouped by the form in which they present the result. This
is the classification JKSSB question-setters use most often, so learn one
example for each row.

\begin{center}
\begin{tabular}{p{0.26\textwidth}p{0.60\textwidth}}
\toprule
\textbf{Category} & \textbf{Typical devices} \\
\midrule
Visual / soft copy & Monitor, projector \\
Audio & Speaker, headphones \\
Printed / hard copy & Printer, plotter \\
\bottomrule
\end{tabular}
\end{center}

A single device may belong to more than one row if it produces more than one
kind of output, but for exam purposes use the primary role: a monitor is a
visual output device, a speaker is an audio output device, and a printer is a
printed (hard copy) output device.

\section{Soft copy vs hard copy}
Output comes in two forms. \textbf{Soft copy} is output that is displayed
electronically, such as an image on a monitor or a projected slide; it is not
printed and can usually still be edited. \textbf{Hard copy} is output that is
printed on a physical medium such as paper, produced by a printer or a
plotter; once printed it is fixed.

\begin{definitionbox}{Soft copy}
Output that is shown on a screen or projected electronically and is not
printed, such as the image on a monitor or projector.
\end{definitionbox}

\begin{definitionbox}{Hard copy}
Output that is printed on a physical medium such as paper, produced by a
printer or a plotter.
\end{definitionbox}

The distinction matters because a single item can test it in either
direction. If the stem mentions a screen, a display, or a projected image,
the answer is soft copy; if it mentions paper, printing, or a physical
printout, the answer is hard copy.

\begin{center}
\begin{tabular}{p{0.20\textwidth}p{0.34\textwidth}p{0.34\textwidth}}
\toprule
\textbf{Point} & \textbf{Soft copy} & \textbf{Hard copy} \\
\midrule
Form & Electronic, on a screen & Physical, on paper \\
Examples & Monitor, projector & Printer, plotter \\
Editable & Easily edited & Fixed once printed \\
\bottomrule
\end{tabular}
\end{center}

\subsection{Uses and limits of each form}
Soft copy and hard copy serve different needs, and an exam item may ask which
is better for a stated purpose. Soft copy is convenient because it can be
edited, searched, copied, and shared electronically, and it costs nothing to
display; its limit is that it needs power and a screen, and it is easy to lose
if it is not saved. Hard copy is convenient because it needs no power to read,
is permanent, and can be signed or filed as a legal record; its limits are
that it can be bulky, is fixed once printed, and consumes paper and ink. The
rule of thumb is: to read, edit, or send, use soft copy; to sign, file, or
archive, use hard copy.

\section{Monitors}
A \textbf{monitor} is the main visual output device of a computer. It displays
text, images, and video as soft copy. Screen size is measured diagonally in
inches, and picture sharpness is measured in pixels, which together give the
resolution of the display. A monitor shows output to the user; it does not
take input, with the single exception of a touchscreen, which senses touch in
addition to displaying.

\subsection{CRT displays}
A \textbf{CRT} display uses a \textbf{Cathode Ray Tube}. An electron beam
sweeps across a phosphor-coated screen inside a vacuum tube to form the image.
CRT monitors are bulky and deep, and they are now largely replaced by flat
panels. For exam purposes the key point is the full form: CRT = \textbf{Cathode
Ray Tube}, and it is the older, bulky technology.

\subsection{Flat-panel displays: LCD, LED, and OLED}
A \textbf{flat-panel display} is thin, light, and flat. The common
technologies are LCD, LED, and OLED.

\begin{center}
\begin{tabular}{p{0.14\textwidth}p{0.44\textwidth}p{0.24\textwidth}}
\toprule
\textbf{Type} & \textbf{What it means} & \textbf{Backlight} \\
\midrule
LCD & Liquid Crystal Display; cells of liquid crystal control light & Needs a backlight \\
LED & Light Emitting Diode display; an LCD panel lit by LEDs & LED backlight \\
OLED & Organic Light Emitting Diode; each pixel emits its own light & No backlight (self-emitting) \\
\bottomrule
\end{tabular}
\end{center}

The nearest-neighbour trap is LCD versus LED versus OLED. An \textbf{LCD}
needs a separate backlight to be visible. An \textbf{LED} display is really an
LCD panel whose backlight is made of LEDs. An \textbf{OLED} display needs no
backlight at all, because each pixel makes its own light.

\subsection{Display size, resolution, and refresh}
The size of a monitor is given as the diagonal length of the screen in inches.
The picture is made of tiny dots called \textbf{pixels}; the number of pixels
across and down the screen is the \textbf{resolution}, such as 1920 by 1080. A
higher resolution shows more detail in the same screen area. The
\textbf{refresh rate} is how many times per second the screen is redrawn,
measured in hertz (Hz); a higher refresh rate gives smoother motion. These are
the three numbers a display specification quotes, and they apply to monitors
and to laptop screens alike.

\section{Projector}
A \textbf{projector} throws a large image onto a screen or a wall. It is used
in classrooms, meetings, and cinemas to show the same computer output to many
people at once. The projected image is soft copy, because it is shown on a
surface and not printed. A projector is an output device; it does not capture
anything, so it must not be confused with a scanner or a camera.

\section{Audio output: speakers and headphones}
A \textbf{speaker} converts an electrical signal into sound that a group of
listeners can hear. \textbf{Headphones} produce sound for a single listener
and are worn on or over the ears. Both are audio output devices.

The nearest-neighbour pair is the speaker against the microphone. A
\textbf{microphone} takes sound in and is an audio input device; a
\textbf{speaker} or a \textbf{headphone} sends sound out and is an audio
output device.

\begin{example}
A student plays a recorded lecture on the office desktop. The sound leaves the
computer through the speakers, so the speakers are audio output devices. If
the student then records a reply into the same computer, the microphone is
used; that is an audio input device. The same machine performs both jobs
through two different devices.
\end{example}

\section{Printers: overview}
A \textbf{printer} produces hard copy on paper. Printers fall into two broad
classes according to how they mark the paper. An \textbf{impact printer}
strikes a ribbon against the paper; a \textbf{non-impact printer} marks the
paper without striking it, using ink, toner, or heat. The common types are the
dot-matrix printer (impact) and the inkjet and laser printers (non-impact).
Laser printers are common in offices, while inkjet printers are common in
homes. A printer is an output device (TRAP-12).

\begin{center}
\begin{tabular}{p{0.18\textwidth}p{0.34\textwidth}p{0.34\textwidth}}
\toprule
\textbf{Point} & \textbf{Impact} & \textbf{Non-impact} \\
\midrule
Method & Strikes a ribbon and the paper & No striking; uses ink, toner, or heat \\
Examples & Dot-matrix & Inkjet, laser \\
Noise & Noisier & Quieter \\
Graphics quality & Lower & Higher \\
\bottomrule
\end{tabular}
\end{center}

A \textbf{dot-matrix} printer is an impact printer and is the noisiest of the
common types. An \textbf{inkjet} printer sprays tiny drops of ink onto the
paper, and a \textbf{laser} printer forms the image with toner that is fused
onto the paper by heat. Both are non-impact printers.

The choice follows the job. A dot-matrix printer is cheap to run and can print
multi-part carbon forms, but its quality is low and it is noisy. An inkjet
printer is quiet and inexpensive to buy and is suited to colour photographs
and home use, though the ink can be costly. A laser printer is fast, quiet,
and gives sharp text, which is why it is the usual office choice. Whatever the
type, every printer is an output device producing hard copy.

\begin{example}
An office prints a hundred letters every morning. The best choice is a laser
printer, because it is fast, quiet, and gives sharp text. If the same office
had to print only the occasional colour photograph at home, an inkjet printer
would be enough. If it needed multi-part carbon forms for a carbon copy, a
dot-matrix printer would be the right impact printer.
\end{example}

\section{Plotter}
A \textbf{plotter} draws large, precise line drawings, such as engineering and
architectural designs, on paper larger than a normal printer's. It draws using
a pen or a vector-based method, which makes it suitable for the continuous
lines of technical drawings. A plotter is an output device: it produces hard
copy but it does not read anything, so it must not be confused with a scanner
(TRAP-12).

\begin{example}
An item asks which output device is used to produce large engineering
drawings. The answer is a \textbf{plotter}, because it is designed for large,
precise line drawings. A normal printer is meant for ordinary pages, a monitor
shows soft copy, and a scanner is an input device.
\end{example}

\section{Multifunction printer (MFP)}
An \textbf{MFP} is a \textbf{Multifunction Printer}, also called an
\textbf{all-in-one} device. It combines \textbf{printing, scanning, copying},
and often \textbf{faxing} in a single unit. Because it does several jobs, an
MFP is the one common device that is both an input and an output device: its
scanner part takes input, while its printer part gives output. MFPs are common
in offices and small businesses.

\section{Summary of output devices}
The table below gathers the devices in this lesson for quick revision. Use it
as a checklist: for each device, be ready to state its category (visual,
audio, or printed) and its one-line function, and to name the device it is
most often confused with.

\begin{center}
\begin{tabular}{p{0.20\textwidth}p{0.22\textwidth}p{0.46\textwidth}}
\toprule
\textbf{Device} & \textbf{Category} & \textbf{Function} \\
\midrule
Monitor & Visual / soft copy & Shows text, images, and video on a screen. \\
Projector & Visual / soft copy & Throws a large image onto a screen or wall. \\
Speaker & Audio & Converts an electrical signal into sound for a group. \\
Headphones & Audio & Produces sound for a single listener. \\
Printer & Printed / hard copy & Prints text and graphics on paper; impact or non-impact. \\
Plotter & Printed / hard copy & Draws large, precise line drawings. \\
MFP & Hybrid & Combines printing, scanning, copying, and often faxing. \\
\bottomrule
\end{tabular}
\end{center}

\section{Points of confusion between similar devices}
Beyond classifying a device, JKSSB items often test a single distinction
between two devices that look or sound alike. The table below collects the
pairs used in this lesson.

\begin{center}
\begin{tabular}{p{0.30\textwidth}p{0.60\textwidth}}
\toprule
\textbf{Pair} & \textbf{The distinction to state} \\
\midrule
Printer vs plotter & Printer: standard pages of text and graphics. Plotter: large, precise line drawings. \\
Monitor vs projector & Monitor: image on its own screen. Projector: image thrown onto a screen or wall. \\
Speaker vs microphone & Speaker: audio output. Microphone: audio input. \\
Soft copy vs hard copy & Soft copy: shown on a screen. Hard copy: printed on paper. \\
LCD vs OLED & LCD: needs a backlight. OLED: self-emitting, no backlight. \\
Impact vs non-impact & Impact: strikes the paper (dot-matrix). Non-impact: no striking (inkjet, laser). \\
\bottomrule
\end{tabular}
\end{center}

\begin{example}
An office buys one machine that must print letters, scan signed forms, and
photocopy pages. The right choice is a Multifunction Printer (MFP), because it
does all three jobs in one unit. If the office instead needed very large
engineering plans, the correct device would be a plotter, and if it only
needed to show a presentation to a room, the correct device would be a
projector. Matching the output device to the job is the practical form of this
topic's classification.
\end{example}

\section{Other output devices on the exam}
A few output devices appear in real JKSSB items without being part of the
main list. A \textbf{sound card} produces audio output: it turns the digital
sound data into a signal for the speakers. A \textbf{Voice Response Unit
(VRU)} speaks responses to a caller and so gives audio output. \textbf{Haptic
gloves} and a \textbf{tactile display} give touch or force feedback to the
user and are therefore output devices. All of these feed information
\emph{back} to the user and must be told apart from input devices of similar
names, such as a \textbf{data glove}, which captures hand movement as input
(PYQ-09).

\begin{definitionbox}{Sound card}
An expansion card that produces audio output, converting digital sound data
into a signal that speakers or headphones can play.
\end{definitionbox}

\section{Exam retrieval and common confusions}
Seven distinctions prevent most errors on this topic.
\begin{enumerate}
  \item \textbf{Soft copy} is on a screen (monitor, projector); \textbf{hard
    copy} is printed on paper (printer, plotter). Do not swap them.
  \item A \textbf{plotter} and a \textbf{printer} are \textbf{output} devices;
    a \textbf{scanner}, a \textbf{light pen}, and a \textbf{digitizer} are
    \textbf{input} devices (TRAP-12).
  \item \textbf{CRT} = Cathode Ray Tube and is bulky; modern monitors are flat
    panels (LCD, LED, OLED).
  \item An \textbf{LCD} needs a backlight, an \textbf{LED} display is an LCD
    with an LED backlight, and an \textbf{OLED} display is self-emitting with
    no backlight.
  \item A \textbf{dot-matrix} printer is an \textbf{impact} printer; inkjet
    and laser printers are \textbf{non-impact}.
  \item A \textbf{plotter} draws large engineering drawings (PYQ-29), while a
    normal printer is meant for ordinary pages.
  \item A \textbf{sound card}, a \textbf{Voice Response Unit (VRU)}, a
    \textbf{haptic glove}, and a \textbf{tactile display} give output to the
    user; a \textbf{data glove} captures hand movement as input (PYQ-09).
\end{enumerate}

\begin{example}
An item lists four devices --- scanner, plotter, drum printer, and light pen
--- and asks which are output devices. Classify each one: the scanner is
input, the plotter is output, the drum printer is output, and the light pen is
input. The answer is the plotter and the drum printer. This is exactly the
input/output classification tested by TRAP-12.
\end{example}

\subsection*{Exercises}
\begin{enumerate}[label=\arabic*.]
  \item Define an output device and give three examples.
  \item Distinguish soft copy from hard copy, giving one example of each.
  \item Expand CRT, LCD, and OLED, and state which of them needs a backlight.
  \item Explain the difference between an impact printer and a non-impact
    printer, giving one example of each.
  \item Which output device is used for large engineering drawings, and why?
  \item State whether each is input or output: plotter, scanner, projector,
    digitizer tablet, speaker.
\end{enumerate}

\subsection*{Solutions}
\begingroup\small
\begin{enumerate}[label=\arabic*.]
  \item An output device receives processed data from the CPU and presents it
    to the user in a readable form. Examples: monitor, printer, speaker (any
    three valid devices).
  \item Soft copy is shown on a screen and is not printed, such as an image on
    a monitor or projector. Hard copy is printed on paper, such as a page from
    a printer or plotter.
  \item CRT = Cathode Ray Tube; LCD = Liquid Crystal Display; OLED = Organic
    Light Emitting Diode. The LCD needs a backlight; the OLED does not,
    because each pixel emits its own light.
  \item An impact printer strikes a ribbon against the paper, such as a
    dot-matrix printer. A non-impact printer marks the paper without striking
    it, such as an inkjet or a laser printer.
  \item A plotter is used, because it is designed to draw large, precise line
    drawings such as engineering and architectural designs (PYQ-29).
  \item Plotter --- output; scanner --- input; projector --- output; digitizer
    tablet --- input; speaker --- output.
\end{enumerate}
\endgroup

\lessonheading{15}{Printers and Displays}
\section{Printers and hard-copy output}
A printer is an output device. It accepts data from the computer and records
it as \textbf{hard copy} --- output on a physical medium, normally paper.
Output that stays on a screen is called \textbf{soft copy}; a printer is the
device that turns soft copy into hard copy. Printers are described from three
angles: by the \textbf{mechanism} they use (dot matrix, inkjet, laser,
thermal, 3D), by whether they touch the paper (\textbf{impact} or
\textbf{non-impact}), and by the \textbf{unit printed} (a character, a line,
or a page).

\begin{definitionbox}{Printer}
An output device that records computer output as a permanent, human-readable
copy on paper or another medium.
\end{definitionbox}

\begin{center}
\begin{tabular}{p{0.26\textwidth}p{0.62\textwidth}}
\toprule
\textbf{Basis} & \textbf{Types} \\
\midrule
By mechanism & Dot matrix, inkjet, laser, thermal, 3D printer \\
By contact & Impact (strikes the paper) vs non-impact (no contact) \\
By unit printed & Character printer, line printer, page printer \\
\bottomrule
\end{tabular}
\end{center}

The same machine can carry several labels at once. A laser printer, for
example, is a non-impact printer, a page printer, and a laser-mechanism
printer. In an exam item, read which label is being asked for before choosing.

\section{Impact and non-impact printers}
An \textbf{impact printer} forms characters by physically striking an inked
ribbon against the paper. The contact makes these printers relatively noisy,
but because the type presses through the paper, they can produce
\textbf{multiple carbon copies} of multi-part forms. Common impact printers
are the dot matrix printer, the daisy wheel printer, and line printers such
as the drum printer and chain printer.

A \textbf{non-impact printer} forms characters without striking the paper.
These printers are quieter and generally give better print quality, but they
cannot make carbon copies. Common non-impact printers are the inkjet printer,
the laser printer, and the thermal printer.

\begin{definitionbox}{Impact printer}
A printer that strikes an inked ribbon against the paper to form a character.
\end{definitionbox}

\begin{definitionbox}{Non-impact printer}
A printer that forms characters without touching the paper with a striking
mechanism.
\end{definitionbox}

\begin{center}
\begin{tabular}{lll}
\toprule
\textbf{Point} & \textbf{Impact} & \textbf{Non-impact} \\
\midrule
Contact with paper & Yes (strikes) & No \\
Noise & Louder & Quieter \\
Carbon copies & Possible & Not possible \\
Examples & Dot matrix, daisy wheel, & Inkjet, laser, thermal \\
           & drum/chain line printers & \\
\bottomrule
\end{tabular}
\end{center}

\begin{example}
An item lists ``inkjet, laser, thermal, dot matrix'' and asks which is an
impact printer. Trace the test: which device strikes an inked ribbon against
the paper? Only the dot matrix printer does. The other three form the image
without contact, so they are non-impact.
\end{example}

\section{Impact printers in detail}
\subsection*{Dot matrix printer}
The dot matrix printer is an impact printer. Its print head holds rows of
tiny pins or hammers. The pins strike an inked ribbon to build each character
as a pattern of dots inside a matrix. It is noisy and its print quality is
low, but it is cheap to run and can print carbon copies and multi-part
invoices and receipts. For this reason it survives in billing and point-of-sale
work.

\subsection*{Daisy wheel printer}
A daisy wheel printer is an impact printer that uses a spinning wheel with
raised type characters, much like an electric typewriter. It produces good
text quality but cannot print graphics or images well.

\subsection*{Line printers}
A line printer prints an entire line of text in one operation instead of one
character at a time. A \textbf{drum printer} uses a rotating drum that carries
the full character set, with a hammer striking the paper wherever a character
is needed. A \textbf{chain (band) printer} moves a chain or band of type
characters past the hammers. Line printers are high-speed impact printers
associated with older mainframes and large batch printing.

\begin{definitionbox}{Line printer}
A printer that prints one full line of characters at a time; the drum printer
and chain printer are line printers.
\end{definitionbox}

\subsection*{Print units}
Printers are often grouped by the unit they print. A \textbf{character
printer} prints one character at a time (dot matrix, daisy wheel). A
\textbf{line printer} prints one line at a time (drum, chain). A \textbf{page
printer} prints one complete page at a time (laser and LED printers). The
unit grows from character to line to page.

\section{Inkjet printers}
An inkjet printer is a non-impact printer. Its print head sprays tiny
\textbf{droplets of liquid ink} through fine nozzles onto the page. Because
there is no striking, it is quieter than an impact printer, and it prints
good-quality full-colour images and photographs on many kinds of paper. Its
drawbacks are that it is slower per page than a laser printer, that wet ink
can \textbf{smudge}, and that its ink cartridges must be replaced when empty.
Inkjet printers use liquid ink, not toner.

\begin{definitionbox}{Inkjet printer}
A non-impact printer that forms characters and images by spraying fine
droplets of liquid ink onto the paper.
\end{definitionbox}

\section{Laser printers}
A laser printer is a non-impact printer that uses a \textbf{laser beam} and
\textbf{toner}, a fine dry powder, in a xerographic (electrostatic) process.
The laser draws the page image onto a charged drum; the toner sticks to the
charged areas; the toner is then transferred to the paper and fused in place
with heat and pressure. The result is fast, sharp, high-quality text, which
is why laser printers are common in offices. A laser printer is a
\textbf{page printer}: it prints a whole page at a time. Because it uses
toner rather than liquid ink, its output does not smudge when it gets wet.

\begin{definitionbox}{Laser printer}
A non-impact page printer that uses a laser beam to draw an image on a
charged drum and toner powder to form the printed page.
\end{definitionbox}

\begin{center}
\begin{tabular}{lll}
\toprule
\textbf{Point} & \textbf{Inkjet} & \textbf{Laser} \\
\midrule
Contact & Non-impact & Non-impact \\
Image formed by & Liquid ink droplets & Laser beam + toner \\
Unit printed & Character/line & Page printer \\
Speed & Slower & Faster \\
Strengths & Colour photos & Sharp text, offices \\
\bottomrule
\end{tabular}
\end{center}

\section{Thermal printers}
A thermal printer is a non-impact printer that uses \textbf{heat}. It heats
specially treated \textbf{thermal paper} (or a heat-sensitive ribbon) to form
characters. In its simplest form it needs no ink or toner cartridge, and it
operates silently. It is used for point-of-sale receipts, tickets, barcode
labels, and older fax machines. It is usually limited to a single colour.
Thermal, laser, and inkjet printers are all non-impact; dot matrix is impact.

\section{3D printers}
A 3D printer builds a solid three-dimensional object \textbf{layer by layer}
from a digital model. This process is called \textbf{additive
manufacturing}. It does not print ink onto paper; instead it adds material
such as plastic, resin, or metal until the object is complete. 3D printers
are used for prototypes, models, spare parts, and some medical implants. The
word ``printer'' is used, but the output is a physical object rather than a
document.

\section{Plotters}
A plotter is an output device that \textbf{draws} line graphics, usually with
one or more pens that move over the paper. It works with \textbf{vector}
drawings, so it is ideal for large engineering drawings, maps, and CAD plans,
and it handles sheet sizes that ordinary printers cannot. A plotter draws
lines; a printer lays down dots or characters. Both are output devices. A
plotter must not be confused with a scanner or a digitizer, which are input
devices.

\begin{example}
An item says, ``A plotter is an input device.'' This is wrong. A plotter
produces output by drawing lines on paper, so it is an \textbf{output}
device. The items that capture drawings as data --- the scanner and the
digitizer tablet --- are the input devices.
\end{example}

\section{Display terminology}
A monitor displays soft copy, and several exact terms describe it.

\begin{definitionbox}{Pixel}
The smallest addressable element of a display; short for ``picture element''.
Every image on the screen is built from pixels.
\end{definitionbox}

\begin{definitionbox}{Resolution}
The number of pixels a display can show, written as columns $\times$ rows,
for example $1920 \times 1080$. More pixels mean a sharper image.
\end{definitionbox}

The sharpness of a printed page is measured in \textbf{DPI}, which stands for
\textbf{Dots Per Inch}.

\begin{definitionbox}{Refresh rate}
How many times per second a display redraws the image. It is measured in
\textbf{hertz (Hz)}. A higher refresh rate (for example 120 Hz) gives
smoother motion and less flicker.
\end{definitionbox}

\begin{definitionbox}{Response time}
How long a pixel takes to change from one state to another. It is measured
in \textbf{milliseconds (ms)}. A lower response time means less motion blur
or ghosting.
\end{definitionbox}

\textbf{Brightness} is the amount of light a display emits, commonly measured
in \textbf{nits} (candela per square metre, cd/m$^2$). \textbf{Contrast} is
the difference between the brightest white and the darkest black the display
can produce. The \textbf{contrast ratio} compares the two, for example
$1000:1$; a higher ratio gives deeper blacks and more vivid images. In short:
brightness is how much light, and contrast is the range between light and
dark.

\begin{center}
\begin{tabular}{lll}
\toprule
\textbf{Display term} & \textbf{Meaning} & \textbf{Unit} \\
\midrule
Pixel & Smallest addressable element & --- \\
Resolution & Number of pixels shown & pixels (or DPI for print) \\
Refresh rate & Screen redraws per second & hertz (Hz) \\
Response time & Time for a pixel to change & milliseconds (ms) \\
Brightness & Light emitted by the screen & nits (cd/m$^2$) \\
Contrast ratio & Brightest white : darkest black & ratio \\
\bottomrule
\end{tabular}
\end{center}

\section{Exam retrieval and common confusions}
Six distinctions prevent most errors on this topic.
\begin{enumerate}
  \item \textbf{Impact vs non-impact}. Only a printer that \emph{strikes} an
    inked ribbon is impact. Dot matrix, daisy wheel, and drum/chain line
    printers are impact; inkjet, laser, and thermal are non-impact.
  \item \textbf{Dot matrix is impact}. It is easy to group it with inkjet and
    laser by ``modern-ness'', but it strikes the ribbon, so it is impact.
  \item \textbf{Inkjet vs laser}. Inkjet sprays \emph{liquid ink}; laser uses
    a \emph{laser beam and toner}. Do not swap the consumables.
  \item \textbf{Page vs line printer}. A laser printer is a page printer
    (one page at a time); a drum or chain printer is a line printer (one line
    at a time).
  \item \textbf{Printer vs plotter}. A plotter \emph{draws} vector line
    graphics for engineering drawings; it is an output device, not a printer
    of characters and not an input device.
  \item \textbf{Display units}. Refresh rate is in \emph{Hz} (higher is
    better); response time is in \emph{ms} (lower is better); brightness is
    in \emph{nits}; resolution is a pixel count.
\end{enumerate}

\begin{example}
An item asks for the difference between an inkjet and a laser printer. The
correct contrast is the image-forming method and consumable: the inkjet sprays
liquid ink droplets, while the laser uses a laser beam and toner powder.
Both are non-impact. A choice that calls the laser a character printer, or
says the inkjet uses toner, is wrong.
\end{example}

\subsection*{Exercises}
\begin{enumerate}[label=\arabic*.]
  \item Define an impact printer and a non-impact printer, and give two
    examples of each.
  \item Explain how a dot matrix printer forms a character, and state whether
    it is impact or non-impact.
  \item Distinguish an inkjet printer from a laser printer in terms of the
    image-forming method and the consumable used.
  \item What is a page printer? Name one example and contrast it with a line
    printer.
  \item Define pixel, resolution, refresh rate, response time, brightness,
    and contrast, giving the unit of measurement for each where one applies.
\end{enumerate}

\subsection*{Solutions}
\begingroup\small
\begin{enumerate}[label=\arabic*.]
  \item An impact printer forms characters by striking an inked ribbon
    against the paper (dot matrix, daisy wheel, drum/chain line printers). A
    non-impact printer forms characters without striking the paper (inkjet,
    laser, thermal).
  \item The dot matrix print head has rows of pins that strike an inked
    ribbon to build each character as a pattern of dots in a matrix. Because
    it strikes the ribbon, it is an impact printer.
  \item An inkjet printer sprays fine droplets of liquid ink through nozzles;
    a laser printer uses a laser beam to draw an image on a charged drum and
    toner powder to form the page. The inkjet consumable is liquid ink; the
    laser consumable is toner.
  \item A page printer prints a whole page at a time; a laser printer is an
    example. A line printer prints a whole line at a time (drum or chain
    printer), a smaller unit than a page.
  \item A pixel is the smallest addressable element of a display (no unit);
    resolution is the number of pixels shown, as columns $\times$ rows (pixel
    count, or DPI for print); refresh rate is the number of screen redraws per
    second (hertz, Hz); response time is the time for a pixel to change
    (milliseconds, ms); brightness is the light emitted by the screen (nits,
    cd/m$^2$); contrast is the difference between brightest white and darkest
    black, expressed as a contrast ratio.
\end{enumerate}
\endgroup

\lessonheading{16}{Computer Ports}
\section{What a port is}
Every device that connects to a computer does so through a \textbf{port}. A
port is a socket or interface on the machine into which a cable or a device
is plugged. It lets data and, in many cases, electric power pass between the
computer and the attached device. Ports are not all alike: they differ in
their shape, in the type of signal they carry (digital or analog), and in the
job they are meant to do.

The common physical ports found on a personal computer are USB (including
USB-C), HDMI, DisplayPort, VGA, Ethernet through an RJ-45 connector, the
3.5~mm audio jack, and SATA. Knowing which port suits which device is the
central skill of this lesson.

\begin{definitionbox}{Port}
A socket or interface on a device where a cable or another device is
connected, allowing data (and often power) to flow between the computer and
the attached device.
\end{definitionbox}

\begin{center}
\begin{tabular}{p{0.26\textwidth}p{0.64\textwidth}}
\toprule
\textbf{Term} & \textbf{Meaning in this lesson} \\
\midrule
Port & The socket on the device (a USB port, an HDMI port). \\
Connector & The plug end of a cable that fits the port. \\
Cable & The wire that joins two connectors. \\
Physical port & A hardware socket on the machine. \\
Logical port & A software number that names a network service. \\
\bottomrule
\end{tabular}
\end{center}

A first view of the ports covered in this lesson, with their signal type and
their normal use, is given below. The rest of the lesson explains each one in
turn.

\begin{center}
\begin{tabular}{p{0.20\textwidth}p{0.22\textwidth}p{0.48\textwidth}}
\toprule
\textbf{Port} & \textbf{Signal} & \textbf{Typical use} \\
\midrule
USB / USB-C & Digital (data + power) & Keyboard, mouse, pen drive, phone \\
HDMI & Digital (video + audio) & TV, monitor, projector \\
DisplayPort & Digital (video + audio) & PC monitors, graphics cards \\
VGA & Analog (video only) & Older monitors and projectors \\
RJ-45 (Ethernet) & Digital (network) & Wired LAN cable, router, switch \\
3.5 mm audio jack & Analog (audio) & Headphones, speakers, microphone \\
SATA / eSATA & Digital (storage) & Internal and external drives \\
\bottomrule
\end{tabular}
\end{center}

\section{The two meanings of ``port''}
The word \textbf{port} has two entirely different meanings in computing, and
the exam frequently tests the difference. A \textbf{physical port} is a piece
of hardware --- a socket such as a USB port, an HDMI port, or an RJ-45
network port. A \textbf{logical port}, also called a network port, is a
software number that identifies a particular service running on a computer
during network communication.

A web server, for example, listens on logical port 80 for ordinary web
requests and on port 443 for secure web requests. The physical network cable,
meanwhile, is plugged into an RJ-45 port on the network card. Thus one
machine uses both meanings at the same time without any contradiction: the
socket is hardware, and the port number is software.

\begin{definitionbox}{Logical (network) port}
A number, in the range 0 to 65535, that identifies a specific service or
process on a computer for the purpose of network communication. It is not a
physical socket.
\end{definitionbox}

\begin{example}
An item states, ``The monitor is connected to logical port 80.'' This is
nonsense. Port 80 is a logical port used by web (HTTP) traffic, not a socket
for a monitor. A monitor connects to a physical video port such as HDMI,
DisplayPort, or VGA. The wrong statement has mixed the two meanings of the
word \emph{port}.
\end{example}

A few well-known logical port numbers are worth remembering, because the exam
sometimes asks which service uses a given port. The list below is
illustrative; the concept being tested is that these are \emph{numbers}, not
sockets.

\begin{center}
\begin{tabular}{p{0.30\textwidth}p{0.16\textwidth}p{0.40\textwidth}}
\toprule
\textbf{Service} & \textbf{Logical port} & \textbf{Full form} \\
\midrule
HTTP & 80 & HyperText Transfer Protocol \\
HTTPS & 443 & HTTP Secure \\
FTP & 21 & File Transfer Protocol \\
SMTP & 25 & Simple Mail Transfer Protocol \\
DNS & 53 & Domain Name System \\
\bottomrule
\end{tabular}
\end{center}

\section{USB: the Universal Serial Bus}
\textbf{USB} stands for \textbf{Universal Serial Bus}. It is a standard
interface for connecting a wide range of peripherals to a computer. Two of
its best-known features are \textbf{plug and play} and \textbf{hot swapping}:
a device can be attached or removed while the computer is running, without
switching the machine off. A USB connection carries both \textbf{data} and
\textbf{power}, which is why many small devices can be used or charged
directly from a USB port.

Typical USB devices include the keyboard, the mouse, a printer, a pen drive
(flash drive), an external hard disk, a webcam, and a smartphone. The
connector at the end of a USB cable comes in several shapes:

\begin{itemize}
  \item \textbf{USB Type-A}: the familiar rectangular connector found on most
    PCs.
  \item \textbf{USB Type-B}: often used for printers and scanners.
  \item \textbf{Mini-USB} and \textbf{Micro-USB}: smaller connectors used on
    older phones and devices.
  \item \textbf{USB Type-C}: a modern, small, reversible connector.
\end{itemize}

\textbf{USB-C}, or USB Type-C, deserves special attention. It is a
\emph{connector type}, not a separate bus. It is \textbf{reversible}, so it
can be plugged in either way up, and a single USB-C connection can carry
data, power, and video together. It is now common on smartphones, laptops,
external SSDs, and docking stations. The key exam point is precision: USB is
the standard, USB-C is one connector shape for that standard, and a USB port
is the socket on the machine.

The USB standard itself has developed over time through versions such as USB
1.1, USB 2.0, USB 3.0 (and its later 3.1 and 3.2 refinements), and USB4. Each
newer version generally offers a higher data-transfer rate, and the version
is independent of the connector shape: a Type-A or Type-C socket can appear
on a machine that supports different USB versions. For exam purposes, the
safe points are the meaning of the abbreviation, the plug-and-play and
hot-swap behaviour, the combined data-and-power carrying, and the
reversibility of USB-C.

\section{Video ports: HDMI, DisplayPort, and VGA}
Displays are attached through video ports. Three names dominate the syllabus.

\textbf{HDMI} stands for \textbf{High-Definition Multimedia Interface}. It
carries \textbf{digital video and digital audio} over a single cable, which
is why one HDMI cable can replace a separate video cable and a separate audio
cable. It is very common on televisions, monitors, projectors, set-top boxes,
and game consoles.

\textbf{DisplayPort} is also a \textbf{digital} audio/video interface, but it
is designed mainly for computers and monitors. It appears on graphics cards,
laptops, and desktop monitors, and it supports high refresh rates and
multiple monitors through daisy-chaining. \textbf{Mini DisplayPort} is a
smaller variant of the same interface.

\begin{center}
\begin{tabular}{p{0.22\textwidth}p{0.30\textwidth}p{0.38\textwidth}}
\toprule
\textbf{Point} & \textbf{HDMI} & \textbf{DisplayPort} \\
\midrule
Signal & Digital video + audio & Digital video + audio \\
Main use & TV, home theatre, monitor & PC monitors, graphics cards \\
Note & Very common on TVs & Common on PCs; multi-monitor \\
\bottomrule
\end{tabular}
\end{center}

Both HDMI and DisplayPort are digital. The right choice in practice usually
depends on which sockets the devices offer, not on one being inherently
better than the other.

\textbf{VGA} stands for \textbf{Video Graphics Array}. It is an
\textbf{analog} video interface: it carries video only and no audio. It uses
a 15-pin D-sub connector, often coloured blue and held in place by two
screws. VGA is still seen on older monitors and projectors but has largely
been replaced by HDMI and DisplayPort.

\begin{definitionbox}{Digital vs analog video}
Digital video ports (HDMI, DisplayPort) send the picture as discrete numbers
and also carry audio; the analog port (VGA) sends a continuously varying
signal and carries video only.
\end{definitionbox}

\begin{example}
An item reads, ``VGA carries both video and audio.'' This is false. VGA is an
analog video-only interface with no audio channel. HDMI and DisplayPort, by
contrast, carry digital video together with digital audio. The wrong option
has transferred the audio capability of the digital ports onto VGA.
\end{example}

\section{Ethernet and the RJ-45 connector}
\textbf{Ethernet} is the standard for \textbf{wired} computer networking, and
it is commonly associated with the IEEE 802.3 family of standards. The
physical socket used for a wired network connection is the \textbf{RJ-45}
connector. RJ-45 stands for \textbf{Registered Jack 45}, and it has eight
pins, used with twisted-pair cable such as Cat5e or Cat6.

Typical devices attached through an RJ-45 port are a router, a switch, a
modem, and the network card of a desktop computer. A near-neighbour that
causes confusion is \textbf{RJ-11}, the smaller telephone jack. RJ-45 is the
network connector with eight pins; RJ-11 is the telephone connector and is
smaller. The two look similar but serve different purposes.

\section{Audio jacks}
The common audio socket on a personal computer is the \textbf{3.5~mm audio
jack}. It connects headphones, earphones, speakers, and microphones. The plug
wiring has standard forms: \textbf{TRS} (Tip-Ring-Sleeve) carries stereo
audio, while \textbf{TRRS} (Tip-Ring-Ring-Sleeve) adds a microphone channel,
which is why a phone headset with a microphone uses a TRRS plug. On desktop
PCs a colour convention is often followed: green for audio output, such as
headphones or speakers, and pink for a microphone. Some equipment uses
2.5~mm or 6.35~mm jacks instead of the 3.5~mm size.

\section{Storage interfaces: SATA and eSATA}
\textbf{SATA} stands for \textbf{Serial Advanced Technology Attachment}. It
is the interface that connects \textbf{internal storage} --- a hard disk
drive, a solid-state drive, or an optical drive --- to the motherboard. The
SATA data cable is thin and has seven pins, and a separate power connector
feeds the drive. \textbf{eSATA} stands for \textbf{external SATA}, and it is
used to connect an external drive. It is important to classify eSATA
correctly: it is a \emph{storage interface}, not a general-purpose peripheral
port like USB. SATA is about attaching drives, whereas USB, HDMI, and the
audio jack are about attaching other kinds of peripherals and displays.

\section{Classifying a port: four questions to ask}
When a question shows a port or names it, four questions settle almost every
case. First, \emph{is it a physical port or a logical one?} A socket is
physical; a number such as 80 or 443 is logical. Second, \emph{what kind of
signal does it carry?} HDMI and DisplayPort are digital audio/video, VGA is
analog video, the 3.5~mm jack is analog audio, and RJ-45 is a digital network
link. Third, \emph{what class of device does it attach?} Displays go to
video ports, drives go to SATA, network equipment goes to RJ-45, and general
peripherals go to USB. Fourth, \emph{which near-neighbour is it often
confused with?} USB-C is confused with USB generally; RJ-45 is confused with
RJ-11; VGA is confused with HDMI and DisplayPort.

\begin{definitionbox}{Port class}
A way of grouping ports by what they attach: peripheral ports (USB), display
ports (HDMI, DisplayPort, VGA), the network port (RJ-45), audio ports (3.5~mm
jack), and storage interfaces (SATA, eSATA).
\end{definitionbox}

Two ports are easily confused with a near-neighbour. The table below keeps the
distinctions in one place.

\begin{center}
\begin{tabular}{p{0.16\textwidth}p{0.26\textwidth}p{0.44\textwidth}}
\toprule
\textbf{Port} & \textbf{Nearest confusion} & \textbf{The distinction} \\
\midrule
USB-C & USB (the standard) & USB-C is a connector shape; USB is the bus. \\
HDMI & DisplayPort & Both digital A/V; HDMI is common on TVs, DisplayPort on PCs. \\
VGA & HDMI / DisplayPort & VGA is analog, video-only; the others are digital with audio. \\
RJ-45 & RJ-11 & RJ-45 is the 8-pin network jack; RJ-11 is the smaller phone jack. \\
SATA & USB & SATA attaches drives (storage); USB is a general peripheral port. \\
Logical port & Physical port & Logical port is a number; physical port is a socket. \\
\bottomrule
\end{tabular}
\end{center}

\section{Exam retrieval and common confusions}
The following distinctions prevent most errors on this topic.
\begin{enumerate}
  \item A \textbf{physical port} is a hardware socket; a \textbf{logical
    (network) port} is a software number, such as HTTP 80 or HTTPS 443. Never
    mix the two meanings of ``port''.
  \item \textbf{VGA is analog} and video-only; \textbf{HDMI and DisplayPort
    are digital} and carry both video and audio.
  \item \textbf{RJ-45} (network, eight pins) is not \textbf{RJ-11} (the
    smaller telephone jack).
  \item \textbf{USB-C is a connector type} and is reversible; ``USB'' is the
    standard, and a USB port is the socket.
  \item \textbf{SATA is a storage interface}; do not call it a display or
    network port.
  \item Say the term precisely: a \emph{port} is the socket, a
    \emph{connector} is the plug, and a \emph{cable} is the wire. This is the
    same discipline as file \emph{format} vs file \emph{extension} (TRM-09).
\end{enumerate}

\begin{example}
An item asks which port carries digital audio and digital video over one
cable. Name the digital audio/video ports: HDMI and DisplayPort both qualify,
whereas VGA is analog video-only, RJ-45 is networking, and SATA is storage.
If the stem prefers the port common on television sets, choose HDMI; if it
stresses PC monitors and multi-monitor setups, choose DisplayPort.
\end{example}

The full forms below are the ones most likely to be asked directly, so they
should be memorised exactly.

\begin{center}
\begin{tabular}{p{0.18\textwidth}p{0.56\textwidth}}
\toprule
\textbf{Abbreviation} & \textbf{Full form} \\
\midrule
USB & Universal Serial Bus \\
HDMI & High-Definition Multimedia Interface \\
VGA & Video Graphics Array \\
RJ-45 & Registered Jack 45 \\
SATA & Serial Advanced Technology Attachment \\
eSATA & external SATA \\
TRS & Tip-Ring-Sleeve \\
TRRS & Tip-Ring-Ring-Sleeve \\
HTTP & HyperText Transfer Protocol \\
HTTPS & HTTP Secure \\
\bottomrule
\end{tabular}
\end{center}

\subsection*{Exercises}
\begin{enumerate}[label=\arabic*.]
  \item Distinguish a physical port from a logical port, giving one example
    of each.
  \item Expand USB, HDMI, VGA, RJ-45, and SATA.
  \item Which video port is analog and video-only, and which two common video
    ports are digital and also carry audio?
  \item What is the difference between RJ-45 and RJ-11?
  \item What is the difference between SATA and eSATA, and why is SATA not a
    general peripheral port?
  \item Which port is commonly used to connect a monitor and supports high
    refresh rates and multiple monitors, and which connector is used for a
    wired Ethernet connection?
  \item Explain the difference between TRS and TRRS audio plugs.
  \item Name two features of USB that make attaching devices convenient.
\end{enumerate}

\subsection*{Solutions}
\begingroup\small
\begin{enumerate}[label=\arabic*.]
  \item A physical port is a hardware socket such as a USB port or an RJ-45
    port. A logical port is a software number that identifies a network
    service, such as HTTP port 80.
  \item USB = Universal Serial Bus; HDMI = High-Definition Multimedia
    Interface; VGA = Video Graphics Array; RJ-45 = Registered Jack 45;
    SATA = Serial Advanced Technology Attachment.
  \item VGA is the analog, video-only port. HDMI and DisplayPort are digital
    and carry both video and audio.
  \item RJ-45 is the eight-pin connector used for wired Ethernet networking;
    RJ-11 is the smaller telephone jack.
  \item SATA connects internal drives (HDD, SSD, optical drive) to the
    motherboard. eSATA is external SATA, used for an external drive. SATA is a
    storage interface --- it attaches drives --- so it is not a
    general-purpose peripheral port such as USB.
  \item DisplayPort is the monitor port for high refresh rates and multiple
    monitors; RJ-45 is the connector for a wired Ethernet connection.
  \item TRS (Tip-Ring-Sleeve) has three contact areas and carries stereo
    audio; TRRS (Tip-Ring-Ring-Sleeve) has four and adds a microphone channel.
  \item USB supports plug and play and hot swapping, and it carries both data
    and power over the same connection.
\end{enumerate}
\endgroup

\lessonheading{65}{Current Multimedia, GUI and Media Formats}
\section{What multimedia means}
\textbf{Multimedia} is the presentation of content that combines more than
one medium in a single document or program. The media may be \textbf{text,
graphics, audio, video}, and \textbf{animation}. Text and graphics are
\emph{static} media: they do not change with time. Audio, video, and animation
are \emph{time-based} media: their meaning depends on being played in the
correct order and at the correct speed. Because a computer stores everything
as binary data, every medium --- sound, picture, or moving image --- must
first be captured and converted into numbers before it can be stored,
compressed, and replayed.

\begin{definitionbox}{Multimedia}
Content that combines two or more media types --- such as text, graphics,
audio, video, and animation --- in a single digital presentation.
\end{definitionbox}

\begin{center}
\begin{tabular}{p{0.20\textwidth}p{0.68\textwidth}}
\toprule
\textbf{Element} & \textbf{Meaning} \\
\midrule
Text & Readable characters; the base medium of almost every document. \\
Graphics & Still images, drawings, and photographs. \\
Audio & Sound, music, and speech. \\
Video & Moving pictures, usually accompanied by sound. \\
Animation & A sequence of images played fast enough to create motion. \\
\bottomrule
\end{tabular}
\end{center}

\section{Digitising media: sampling and quantization}
An analogue signal, such as a sound wave or a scene's light, is
\emph{continuous}: it has a value at every instant. A computer, however,
stores only discrete numbers. Converting a continuous signal into digital
form is done by \textbf{analogue-to-digital conversion}, and it has two
distinct steps, \textbf{sampling} and \textbf{quantization}.

\begin{definitionbox}{Sampling}
The process of measuring a continuous signal at regular intervals and
recording each measurement. The number of measurements taken per second is
the \textbf{sampling rate}, measured in hertz (Hz).
\end{definitionbox}

A \textbf{higher sampling rate} captures the signal more often, so the
reproduction is more faithful, but the file becomes larger. A lower sampling
rate saves space but loses detail. The sampling rate is the horizontal
resolution of the digitised signal: how often we look at it.

\begin{definitionbox}{Quantization}
The process of mapping each measured sample to one of a fixed set of discrete
levels, because a computer can store only a limited number of distinct values
per sample.
\end{definitionbox}

The number of available levels is set by the \textbf{bit depth} of each
sample. With a higher bit depth --- more bits per sample --- there are more
levels, so each recorded value is closer to the true value; the rounding
error, called \emph{quantization noise}, is smaller, and quality is better,
but again the file is larger.

\begin{example}
Think of drawing a smooth curve on graph paper. \textbf{Sampling} decides how
many points across the page you mark --- more points follow the curve more
closely. \textbf{Quantization} decides how finely the vertical axis is ruled
--- more lines let you record each point's height more precisely. The two
choices together determine how closely the digitised version matches the
original, and how big the file becomes.
\end{example}

The common exam confusion is to swap the two. Remember the one-line
distinction: \textbf{sampling is how often we measure; quantization is how
precisely each measurement is stored.}

\section{Compression: lossy and lossless}
Raw multimedia files are large, so they are almost always \textbf{compressed}
--- stored in a form that takes less space. There are two kinds of
compression, and the difference is whether the original can be recovered
exactly.

\begin{definitionbox}{Lossless compression}
Compression that allows the original data to be reconstructed exactly. No
information is discarded. PNG and GIF images and ZIP archives are lossless.
\end{definitionbox}

\begin{definitionbox}{Lossy compression}
Compression that discards some data to reduce the file size further, so the
reconstruction is only an approximation of the original. JPEG images and MP3
audio are lossy.
\end{definitionbox}

Lossless compression works by removing redundancy --- repeated or predictable
data --- while keeping every value. Lossy compression goes further: it removes
detail that a human eye or ear is unlikely to notice. Lossy files are usually
much smaller, which is why lossy formats dominate the internet for
photographs, music, and video. When exactness matters --- for a diagram,
screenshot, or archive --- a lossless format is used instead.

\section{Synchronization}
A multimedia file often carries more than one \textbf{stream} at once: a video
stream and an audio stream, for instance. \textbf{Synchronization} is the
process of keeping those streams aligned in time so they play together
correctly. If the streams drift out of step, the sound no longer matches the
picture --- the narrator's voice is heard before or after the lips move, which
is called a loss of \emph{sync}. Synchronization is what makes a video play
with the sound in the right place at the right moment; it is a timing
requirement peculiar to time-based media.

\begin{definitionbox}{Synchronization}
The alignment in time of separate media streams --- such as audio and video
--- so that they are presented together correctly during playback.
\end{definitionbox}

\section{The graphical user interface and its elements}
A \textbf{GUI} is a \textbf{Graphical User Interface}: it lets the user work
with pictures, icons, and menus instead of typing commands. A \textbf{CLI} is
a \textbf{Command-Line Interface}: it accepts typed text commands. The GUI is
easier for beginners and is the standard on modern desktops; the CLI is often
faster for an expert and is still used for administration.

\begin{definitionbox}{GUI (Graphical User Interface)}
A user interface in which the user interacts with the computer through visual
elements such as windows, icons, menus, and a pointer, rather than by typing
text commands.
\end{definitionbox}

The four classic GUI elements are remembered by the short form \textbf{WIMP}:
\textbf{Windows, Icons, Menus, Pointers}.

\begin{center}
\begin{tabular}{p{0.16\textwidth}p{0.72\textwidth}}
\toprule
\textbf{Element} & \textbf{Meaning} \\
\midrule
Window & A rectangular on-screen area that holds a program or a document. \\
Icon & A small picture that represents a file, folder, or program. \\
Menu & A list of commands that the software offers. \\
Pointer & The on-screen arrow, moved by the mouse or touchpad, used to select. \\
\bottomrule
\end{tabular}
\end{center}

\begin{definitionbox}{Dialogue box}
A small window that appears to ask the user for input or to confirm an action,
such as Save As or Print. It usually contains text boxes, drop-down lists,
check boxes, radio buttons, and OK/Cancel buttons.
\end{definitionbox}

\section{WYSIWYG}
\textbf{WYSIWYG} stands for \textbf{What You See Is What You Get}. It means
that what is shown on the screen is exactly how the document will look when
printed --- the same layout, fonts, spacing, and page breaks. A word processor
such as MS Word works in WYSIWYG mode, so the page visible on the monitor
matches the printed page.

\begin{example}
While editing a document in MS Word, the text, margins, and page breaks on
screen match the printed output. This is WYSIWYG. When the user opens the
Print dialogue box and chooses options, those settings are entered in a
dialogue box, not typed on the command line. The two ideas are often asked
together: WYSIWYG describes the on-screen view; a dialogue box is one GUI
element that collects input.
\end{example}

\section{Image formats: PNG, JPEG, and GIF}
Three image formats appear most often in competitive-exam questions.

\begin{center}
\begin{tabular}{p{0.12\textwidth}p{0.26\textwidth}p{0.50\textwidth}}
\toprule
\textbf{Format} & \textbf{Full form} & \textbf{Nature and typical use} \\
\midrule
PNG & Portable Network Graphics & Lossless; supports transparency; used for logos and web graphics. \\
JPEG & Joint Photographic Experts Group & Lossy; best for photographs and complex pictures. \\
GIF & Graphics Interchange Format & Lossless but limited to 256 colours; used for simple graphics and short animation. \\
\bottomrule
\end{tabular}
\end{center}

\textbf{JPEG} is chosen for photographs because its lossy compression makes
them much smaller with little visible loss. \textbf{PNG} is chosen when exact
colours or transparency are needed, because it is lossless. \textbf{GIF} is
chosen for simple graphics and small animations, but it can show only 256
colours, so it is unsuitable for photographs. \textbf{BMP} (Bitmap) is an
older, largely uncompressed image format that produces large files and is not
used for the web.

\section{Audio and video formats: MP3 and MP4}
\textbf{MP3} stands for \textbf{MPEG-1 Audio Layer III} (also written MPEG
Audio Layer III). It is a \textbf{lossy} audio format: it removes sound detail
that the human ear is unlikely to notice, which is why an MP3 file is far
smaller than an uncompressed one. Its small size made MP3 the standard for
portable music and downloads. By contrast, \textbf{WAV} is an uncompressed
audio format.

\textbf{MP4} stands for \textbf{MPEG-4 Part 14}. It is a \textbf{container}
format: a single MP4 file can hold video, audio, and subtitles together. It is
widely used for streaming, downloads, and mobile video. A container such as
MP4 holds the streams; the \emph{codec} inside the container is what actually
compresses the video and audio. \textbf{AVI} (Audio Video Interleave) is an
older video container.

\begin{definitionbox}{MPEG}
MPEG stands for Moving Picture Experts Group --- the group that defines the
MP3 audio layer and the MP4 container, among other standards.
\end{definitionbox}

\section{Frame rate, resolution, and bit rate}
Video is a rapid sequence of still frames. The \textbf{frame rate} is the
number of frames shown per second (fps); a higher frame rate produces smoother
motion but a larger file. The \textbf{resolution} is the number of pixels in
each frame (for example, $1920 \times 1080$); more pixels mean a sharper image
and a larger file. The \textbf{bit rate} is the amount of data stored or sent
per second; a higher bit rate keeps more detail but needs more space or
bandwidth. These three settings, together with the choice of lossy or lossless
compression, decide the final quality and size of a video or audio file.

\begin{definitionbox}{Frame rate}
The number of still frames displayed per second in a video. A higher frame
rate gives smoother motion at the cost of a larger file.
\end{definitionbox}

\begin{example}
A video at 30 frames per second with a high bit rate looks smooth and sharp
but takes much more space than the same clip at a lower frame rate and bit
rate. Streaming services lower the bit rate when the network is slow, trading
quality for smooth playback --- a daily example of the quality-versus-size
trade-off that runs through the whole lesson.
\end{example}

\section{Plotter}
A \textbf{plotter} is an \textbf{output} device that draws on paper. Unlike a
printer, which lays down text and small images, a plotter produces
\textbf{vector graphics}: lines and curves defined by coordinates rather than
by a grid of dots. It is used for \textbf{large engineering drawings}, maps,
blueprints, and banners, where continuous, accurate lines are needed. Early
pen plotters moved a pen across the paper; modern inkjet plotters spray ink
while the sheet moves. A plotter is always an output device, never an input
device (PYQ-29, PYQ-08).

\section{Impact versus non-impact printers (revisit)}
The basic classification of printers reappears in later units, so keep it
firmly in mind. An \textbf{impact printer} forms an image by striking the
paper; a \textbf{non-impact printer} does so without any striking action.

\begin{center}
\begin{tabular}{p{0.16\textwidth}p{0.36\textwidth}p{0.36\textwidth}}
\toprule
\textbf{Point} & \textbf{Impact} & \textbf{Non-impact} \\
\midrule
Action & Strikes the paper to print & Forms the image without striking \\
Examples & Dot-matrix, daisy-wheel, line printer & Inkjet, laser, thermal printer \\
Carbon copies & Possible (same strike hits several sheets) & Not possible \\
Noise & Comparatively noisy & Comparatively quiet \\
\bottomrule
\end{tabular}
\end{center}

The quick test is the mechanism: if the device physically hits the paper, it
is impact; if it does not, it is non-impact. Both kinds are output devices. A
\emph{dot-matrix} printer is impact; an \emph{inkjet} and a \emph{laser}
printer are non-impact.

\section{Multimedia hardware and applications}
Playing and creating multimedia needs both hardware and software. On the
hardware side, a \textbf{sound card} processes audio, \textbf{speakers} and
\textbf{headphones} play it, and a \textbf{microphone} captures it. A
\textbf{graphics card (GPU)} renders images and video, a \textbf{monitor or
projector} displays them, and a \textbf{camera} or \textbf{webcam} captures
still and moving pictures. Because audio and video files are large, multimedia
also places heavy demands on \textbf{memory and storage}.

On the software side, editors for images, audio, and video are used to
capture, edit, compress, and synchronize the media streams. The main
application areas of multimedia are education (e-learning and tutorials),
entertainment (music, games, and streaming films), business (presentations
and advertising), and medicine (medical imaging and simulation).

\begin{center}
\begin{tabular}{p{0.22\textwidth}p{0.66\textwidth}}
\toprule
\textbf{Short form} & \textbf{Full form} \\
\midrule
GUI & Graphical User Interface \\
CLI & Command-Line Interface \\
WYSIWYG & What You See Is What You Get \\
WIMP & Windows, Icons, Menus, Pointers \\
PNG & Portable Network Graphics \\
JPEG & Joint Photographic Experts Group \\
GIF & Graphics Interchange Format \\
MP3 & MPEG-1 Audio Layer III \\
MP4 & MPEG-4 Part 14 \\
MPEG & Moving Picture Experts Group \\
AVI & Audio Video Interleave \\
BMP & Bitmap \\
\bottomrule
\end{tabular}
\end{center}

\section{Exam retrieval and common confusions}
Several distinctions prevent most errors on this topic.
\begin{enumerate}
  \item \textbf{Sampling} is how often a signal is measured; \textbf{quantization}
    is how precisely each sample is stored. Do not swap them.
  \item \textbf{Lossless} formats restore the original exactly (PNG, GIF);
    \textbf{lossy} formats discard data (JPEG, MP3, MP4).
  \item \textbf{JPEG} is lossy and best for photographs; \textbf{PNG} is
    lossless and supports transparency; \textbf{GIF} is lossless and supports
    simple animation.
  \item \textbf{MP3} is audio (MPEG-1 Audio Layer III); \textbf{MP4} is video
    (MPEG-4 Part 14). Do not swap the full forms.
  \item \textbf{Synchronization} keeps audio and video aligned in time; it is
    not compression and not sampling.
  \item \textbf{GUI} is Graphical User Interface; its elements are
    \textbf{Windows, Icons, Menus, Pointers} (WIMP).
  \item \textbf{WYSIWYG} means the screen matches the print; a \textbf{dialogue
    box} is a small window that collects input.
  \item A \textbf{plotter} is an output device for large drawings; an
    \textbf{impact} printer strikes the paper and a \textbf{non-impact} printer
    does not.
\end{enumerate}

\begin{example}
An item asks, ``Which format is lossless and supports transparency?'' Trace
the formats: JPEG is lossy, MP3 is lossy audio, MP4 is lossy video, and GIF is
lossless but limited to 256 colours without transparency support. The answer
is PNG, the lossless format that supports transparency.
\end{example}

\subsection*{Exercises}
\begin{enumerate}[label=\arabic*.]
  \item Define multimedia and name its five building blocks.
  \item Distinguish sampling from quantization, and explain how each affects
    file size.
  \item Differentiate lossy and lossless compression, giving one example of
    each.
  \item What is synchronization in multimedia, and why is it needed?
  \item Expand GUI, WYSIWYG, and WIMP.
  \item State one suitable use each for PNG, JPEG, and GIF.
  \item Expand MP3 and MP4, and state the medium each is used for.
  \item Explain why a plotter is preferred over a printer for large engineering
    drawings.
  \item Distinguish impact printers from non-impact printers with one example
    of each.
\end{enumerate}

\subsection*{Solutions}
\begingroup\small
\begin{enumerate}[label=\arabic*.]
  \item Multimedia is content combining more than one medium in a single
    presentation. Its five building blocks are text, graphics, audio, video,
    and animation.
  \item Sampling measures a continuous signal at regular intervals; a higher
    sampling rate captures more detail but enlarges the file. Quantization maps
    each sample to one of a fixed set of levels set by the bit depth; more
    levels improve quality but also enlarge the file.
  \item Lossless compression restores the original exactly (for example PNG);
    lossy compression discards some data and is only an approximation (for
    example JPEG).
  \item Synchronization is the alignment in time of separate media streams such
    as audio and video. It is needed so that sound and picture play together
    correctly rather than drifting apart.
  \item GUI = Graphical User Interface; WYSIWYG = What You See Is What You Get;
    WIMP = Windows, Icons, Menus, Pointers.
  \item PNG: logos and web graphics needing transparency or exact colours.
    JPEG: photographs and complex pictures. GIF: simple graphics and short
    animations.
  \item MP3 = MPEG-1 Audio Layer III, used for audio (music and speech). MP4 =
    MPEG-4 Part 14, used for video (a container holding video, audio, and
    subtitles).
  \item A plotter produces vector graphics --- continuous lines defined by
    coordinates --- on large sheets, which suits precise engineering drawings,
    maps, and blueprints better than a printer.
  \item Impact printers strike the paper (for example dot-matrix); non-impact
    printers form the image without striking (for example inkjet or laser).
\end{enumerate}
\endgroup

\unitopener{5}{Software, Operating Systems, Windows \& Utilities}
\lessonheading{17}{Hardware, Software and Firmware}
\section{Hardware}
Hardware is the set of physical, tangible parts of a computer system --- the
parts you can see and touch. A computer's hardware includes the central
processing unit (CPU), the motherboard, random access memory (RAM), the hard
disk drive or solid-state drive, the keyboard, the mouse, the monitor, the
printer, and every other physical component. Hardware carries out the actual
work of the machine: it executes instructions, stores data, and moves signals.

Hardware alone does nothing useful. A processor with no instructions to run
is only silicon and metal. The instructions come from software, and the two
together make a working computer. For exam purposes, the test for hardware is
simple: if it is a physical object that you can touch, it is hardware.

\begin{definitionbox}{Hardware}
The physical equipment of a computer system --- the tangible components such
as the CPU, memory, storage, input devices, output devices, and communication
hardware.
\end{definitionbox}

Hardware is usually grouped by function. Input devices, such as the keyboard,
mouse, and scanner, bring data into the computer. Output devices, such as the
monitor, printer, and speaker, present results. Processing hardware, mainly
the CPU and the graphics processing unit (GPU), performs computation. Memory
hardware, such as RAM, ROM, and cache, holds data and instructions. Storage
hardware, such as the hard disk drive (HDD) and solid-state drive (SSD),
keeps files without power. Communication hardware, such as the network
interface card and the motherboard, connects the parts together.

Hardware is also described as internal or external. \textbf{Internal
hardware} is fixed inside the computer case --- the motherboard, CPU, RAM, and
internal drives. \textbf{External hardware}, also called a peripheral, is
connected from outside, such as the keyboard, mouse, monitor, printer, and an
external drive. A question may ask whether a named part is internal or
external; the touch test still applies, because both kinds are physical and
both are hardware.

A single concrete machine runs through this course: the \textbf{office
desktop} --- a personal computer with a CPU, memory, a disk, a keyboard, a
mouse, a monitor, and a printer, running an operating system and MS Word.
Every term in this lesson can be pointed at on that one machine, which makes
the abstract layers easier to remember.

\section{Software}
Software is the set of programs and instructions that tell the hardware what
to do. Unlike hardware, software is intangible: you cannot touch it. It is
kept on secondary storage, such as a hard disk or solid-state drive, and it
is loaded into main memory when it runs. Software is what turns a collection
of physical parts into a usable computer.

\begin{definitionbox}{Software}
The collection of programs, procedures, and related data that direct the
operation of a computer. Software is logical and intangible, and it instructs
the hardware.
\end{definitionbox}

Software falls into two broad kinds. \textbf{System software} runs and manages
the computer itself --- the operating system, device drivers, and utility
programs. \textbf{Application software} helps the user perform a task --- a
word processor such as MS Word, a web browser, a spreadsheet, or a game. Every
program a candidate meets in JKSSB questions belongs to one of these two
kinds, but the boundary matters: the operating system is system software, not
application software.

Software is kept on secondary storage as files. When it is launched, the
operating system copies it into main memory, and the CPU executes it from
there. This is why a program must be \emph{loaded} before it can run, and why
data held only in RAM is lost when the power goes off. The permanent copy
lives on disk; the running copy lives in memory.

\section{Hardware vs software}
The two are easy to confuse only when the wording is careless, so hold the
contrasts in a table.

\begin{center}
\begin{tabular}{p{0.18\textwidth}p{0.36\textwidth}p{0.36\textwidth}}
\toprule
\textbf{Point} & \textbf{Hardware} & \textbf{Software} \\
\midrule
Nature & Physical and tangible & Logical and intangible \\
Can be touched? & Yes & No \\
Role & Executes the work & Gives the instructions \\
Stored or fixed & Fixed as physical parts & Stored as files, loaded to run \\
Example & CPU, printer, keyboard & MS Word, a driver, an OS \\
How it fails & Wears out, breaks, burns out & Bugs, corruption, wrong logic \\
\bottomrule
\end{tabular}
\end{center}

Neither can work alone. Hardware and software together form a computer system,
and a question that names one and describes the other is usually testing this
single distinction. A candidate should be able to sort a mixed list --- a CPU,
MS Word, a printer driver, the BIOS, and a user --- into its correct class at a
glance.

\begin{example}
A question lists ``CPU, MS Word, printer driver, BIOS, operator'' and asks
which are software. \textbf{MS Word} and the \textbf{printer driver} are
software. The CPU is hardware; the BIOS is firmware (software stored in
hardware); and the operator is a person, that is, humanware. Sorting is done
by asking, for each item, whether it is physical, a program, embedded code, or
a person.
\end{example}

\section{System software and application software}
\textbf{System software} manages the machine and provides services to other
programs. The operating system is the main example, together with device
drivers and system utilities. \textbf{Application software} is written to let
the user accomplish a task, and it depends on the system software beneath it.

For the exam, the stable rule is: the operating system is system software.
This is the point tested by the real Junior Assistant 2026 concept recorded as
PYQ-06 (device driver = software; OS = system software). A device driver, too,
is system software --- it is a program, not a physical part.

\section{Firmware}
Firmware is software that is embedded in hardware and stored in non-volatile
memory. It is held in read-only memory (ROM), erasable programmable ROM
(EPROM), or flash memory inside the device, so it is available the moment the
device is powered on. Firmware provides the low-level control instructions a
device needs and starts it up.

\begin{definitionbox}{Firmware}
Software permanently stored in a hardware device, usually in non-volatile
memory such as ROM or flash. It is software by nature but is treated as part
of the hardware because it is fixed to a specific device.
\end{definitionbox}

Common examples of firmware include the \textbf{BIOS} (Basic Input Output
System) or \textbf{UEFI} on a personal computer, which initialises the
hardware and begins loading the operating system; the firmware inside a
printer, a router, a keyboard, or a solid-state drive. Firmware is usually
described as sitting between hardware and software. It is a program, so it is
software; it is stored on the physical device, so it behaves like part of the
hardware.

Firmware differs from ordinary software in how tightly it is bound to the
device. Firmware is written for one specific piece of hardware, is stored in
ROM or flash on that device, and changes rarely (it is updated by a process
called flashing). Ordinary software is general-purpose, is stored on secondary
storage, and is installed and updated frequently.

Firmware is the \textbf{first software to run} when a device is switched on. On
a personal computer, the CPU executes the BIOS or UEFI firmware from ROM before
any operating system is available; that firmware checks the hardware and then
hands control to the operating system. This is exactly why firmware must be
non-volatile: it has to be present even before anything can be loaded. A
candidate should not confuse firmware with the operating system --- the
firmware runs first and \emph{starts} the computer, while the operating system
is what the machine is finally running.

\section{Device drivers}
A device driver is a small program that lets the operating system communicate
with a piece of hardware. It acts as a translator: the operating system sends
a general request such as ``print this page'', and the driver converts it into
the exact signals that a specific printer understands. A printer driver, for
example, lets the operating system print on one particular printer model.

\begin{definitionbox}{Device driver}
A program that enables the operating system to communicate with, and control,
a specific hardware device by translating general commands into
device-specific signals.
\end{definitionbox}

The single most important exam point about a driver is that it is
\textbf{software}, even though it is tied to hardware and is often named after
a device. A candidate may be tempted to call a driver hardware because it
belongs to a printer or a keyboard, but that is wrong. The driver is a program,
it is stored like software, and a device will not work properly without the
right driver installed. This is TRAP-10, and it is also the heart of the real
Junior Assistant 2026 concept PYQ-06.

\begin{example}
A question states, ``A device driver is a hardware component that connects the
printer to the computer.'' This is wrong on two counts: the driver is a
\textbf{program}, not a physical component, and it is software, not hardware.
The physical cable and port are hardware; the driver that makes the printer
work is software. The trap is that the driver carries the device's name.
\end{example}

\section{Humanware}
Humanware, also called \textbf{peopleware}, is the human component of a
computer system. It covers everyone involved: the end users who operate the
machine, the programmers who write software, the system administrators who
maintain it, and the operators who run it. A computer system is usually
described as hardware, software, and humanware together.

\begin{definitionbox}{Humanware}
The people who design, operate, maintain, and use a computer system; the human
component that hardware and software ultimately serve.
\end{definitionbox}

Humanware is not a device and not a program. In a classification question, if
the named item is a person or a role --- user, programmer, operator,
administrator --- it is humanware. It is the layer that makes decisions and
uses the results, and without it the hardware and software would have no
purpose.

\section{The layered computing stack}
The parts of a computer system can be arranged as a stack of layers, each one
resting on the layer below. The order, from the bottom up, is:

\begin{center}
\begin{tabular}{p{0.30\textwidth}p{0.58\textwidth}}
\toprule
\textbf{Layer} & \textbf{Role} \\
\midrule
Hardware & The physical components at the base of the system \\
Firmware & Stored on the hardware; starts it and controls it at a low level \\
Device driver & Lets the operating system communicate with the hardware \\
Operating system & Manages all resources and provides services to programs \\
Application software & Lets the user perform a task \\
\bottomrule
\end{tabular}
\end{center}

Reading the stack from the bottom, the physical parts come first. Firmware sits
directly on them and brings the hardware to life. The device driver sits above
the firmware and gives the operating system a clean way to talk to each device.
The operating system manages memory, files, processes, and devices, and offers
services to the programs above it. The application software at the top is what
the user actually opens and works in. A command travels \emph{down} the stack
--- from an application, through the operating system and the driver, to the
hardware --- and results travel back up.

This stack is the single mental picture for the whole lesson: hardware at the
bottom, software above it, and firmware and drivers as the thin bridges that
connect the two.

\section{Why the stack matters}
The layered stack is not only a way to sort terms; it explains the order in
which work is done. When a user clicks Print in a word processor, the command
begins at the application-software layer. The application asks the operating
system to print; the operating system passes the job to the device driver for
that printer; the driver translates the request into the signals the printer
hardware understands; and the printer's own firmware controls the physical
printing. The result then travels back up the stack to the user. Every layer
depends on the one below it, and a fault in any single layer can stop the
whole chain --- which is why ``reinstall the driver'' is such a common fix.

\begin{example}
A printer fails to print although the paper and ink are fine. Because the
system has layers, the fault could lie in the application (a wrong page range),
the operating system (a stuck print queue), the driver (missing or outdated),
or the hardware (a loose cable). Testing from the bottom up checks each layer
in turn. In exam terms, the layer that ``lets the operating system talk to the
printer'' is the \textbf{device driver}.
\end{example}

\section{Exam retrieval and common confusions}
Five distinctions prevent most errors on this topic.
\begin{enumerate}
  \item Hardware is \textbf{physical and tangible}; software is
    \textbf{logical and intangible}. The simple test is touch: if you can
    touch it, it is hardware.
  \item A \textbf{device driver} is \textbf{software}, not hardware. Its name
    follows the device, but it is a program (TRAP-10; PYQ-06).
  \item The \textbf{operating system} is \textbf{system software}, not
    application software (PYQ-06).
  \item \textbf{Firmware} is software stored in non-volatile memory inside
    hardware, such as the BIOS; it is software by nature, hardware by storage.
  \item \textbf{Humanware} (peopleware) is the \textbf{people} in the system,
    not a device or a program.
\end{enumerate}

\begin{example}
An item asks, ``Which of the following is software?'' with options CPU, a
printer, a device driver, and a keyboard. The answer is the \textbf{device
driver}: it is a program, whereas the CPU, the printer, and the keyboard are
physical. If the options instead are an operating system, MS Word, a BIOS
chip, and a game, the item is testing system software versus application
software and firmware, and the operating system is the system software.
\end{example}

\subsection*{Exercises}
\begin{enumerate}[label=\arabic*.]
  \item Define hardware and software, and give one example of each.
  \item Distinguish system software from application software. To which kind
    does the operating system belong?
  \item What is firmware, where is it stored, and why is it called ``between
    hardware and software''?
  \item Explain why a device driver is classified as software even though it
    is named after a hardware device.
  \item What is humanware? Give two roles it includes.
  \item List the layers of the computing stack from the bottom (hardware) to
    the top (application).
  \item Give one example each of internal hardware and external hardware.
  \item Using the print command as an example, explain how a request travels
    down the computing stack.
\end{enumerate}

\subsection*{Solutions}
\begingroup\small
\begin{enumerate}[label=\arabic*.]
  \item Hardware is the physical, tangible equipment of a computer, such as
    the CPU or a printer. Software is the programs and instructions that tell
    the hardware what to do, such as MS Word or a device driver.
  \item System software runs and manages the computer itself --- the operating
    system, device drivers, and utilities. Application software lets the user
    do a task, such as a word processor or a game. The operating system is
    \textbf{system software}.
  \item Firmware is software embedded in a hardware device and stored in
    non-volatile memory such as ROM or flash, for example the BIOS. It is
    software by nature but is fixed to a specific device, so it is described
    as sitting between hardware and software.
  \item A device driver is a program, so it is software. It is named after the
    device it serves, which tempts a candidate to call it hardware, but it is
    stored and works like other software; it merely translates the operating
    system's commands for one specific device.
  \item Humanware, or peopleware, is the human component of a computer system.
    It includes roles such as end users and system administrators (also
    programmers or operators).
  \item Hardware, then firmware, then device driver, then operating system,
    then application software.
  \item Internal hardware: the motherboard (also the CPU, RAM, or an internal
    drive). External hardware: the keyboard (also the mouse, monitor, printer,
    or an external drive).
  \item The user clicks Print in the application software. The application
    asks the operating system, which passes the job to the printer's device
    driver. The driver translates the command into the signals the printer
    hardware understands, and the printer's firmware controls the physical
    printing. The result returns up the stack to the user.
\end{enumerate}
\endgroup

\lessonheading{18}{Software Categories}
\section{What is software?}
A computer system has two parts: hardware and software. \textbf{Hardware} is
the physical part that you can see and touch, such as the CPU, memory, and
input/output devices. \textbf{Software} is the set of programs, that is, the
instructions that tell the computer what to do. Programs are stored in memory
and executed by the CPU, and without software the hardware cannot perform any
useful task. Where hardware is the body of the computer, software is the set
of instructions that gives it purpose.

\begin{definitionbox}{Software}
A set of programs or instructions that directs the computer to perform a
task. Hardware is the physical equipment; software is the logical
instructions that run on it.
\end{definitionbox}

Software does not exist as a physical object; it exists as information. A
programmer writes the instructions in a programming language, and these
instructions are then translated, stored in memory, and executed by the CPU.
The same piece of hardware --- a desktop computer, for example --- can behave
completely differently depending on which software is loaded onto it. This is
why software is described as the logical part of a computer system: it
supplies the behaviour that the physical parts only make possible.

\begin{center}
\begin{tabular}{p{0.20\textwidth}p{0.32\textwidth}p{0.36\textwidth}}
\toprule
\textbf{Point} & \textbf{Hardware} & \textbf{Software} \\
\midrule
Nature & Physical; can be touched & Logical; a set of instructions \\
Made of & Electronic and mechanical components & Programs and data \\
Example & CPU, RAM, keyboard, monitor & Operating system, MS Word, a browser \\
Failure & Fails due to wear, damage, or age & Fails due to errors (bugs) or corruption \\
\bottomrule
\end{tabular}
\end{center}

Software is commonly grouped by \emph{what it is for} into six categories:
system software, application software, utility software, programming
software, embedded software, and middleware. These categories describe purpose
and they can overlap in practice, but each answers a different question about
a program's role.

\begin{center}
\begin{tabular}{p{0.26\textwidth}p{0.60\textwidth}}
\toprule
\textbf{Category} & \textbf{Role in one line} \\
\midrule
System software & Controls the hardware and provides a platform for other software. \\
Application software & Helps the user perform a specific task. \\
Utility software & Performs maintenance and housekeeping tasks. \\
Programming software & Used to write, translate, and test programs. \\
Embedded software & Built into a device to control that device. \\
Middleware & Connects different applications and systems. \\
\bottomrule
\end{tabular}
\end{center}

\section{System software}
\textbf{System software} manages and controls the computer hardware and
provides the \emph{platform} on which application software runs. It works
close to the machine, usually starts when the computer is switched on, and is
not written for one user's task. Common examples are the operating system,
device drivers, firmware such as the BIOS, and utility programs.

The \textbf{operating system (OS)} is the most important system software. It
acts as the interface between the user and the hardware and manages memory,
files, processes, and input/output devices. Windows, Linux, macOS, and Android
are familiar operating systems. Its main functions are process management
(starting, scheduling, and stopping programs), memory management (allocating
and freeing main memory), file management (creating, reading, and organising
files), device management (controlling hardware through drivers), and
providing a user interface. \textbf{Firmware} such as the BIOS or UEFI also
belongs to the system-software layer: it initialises the hardware when the
machine is switched on and then loads the operating system.

A \textbf{device driver} is also system
software: it allows the operating system to communicate with a particular
hardware device. It is easy to mistake a device driver for hardware, but it is
a piece of software (TRAP-10).

\begin{definitionbox}{System software}
Software that controls and coordinates the hardware and provides the services
and platform that other programs need in order to run.
\end{definitionbox}

\section{Application software}
\textbf{Application software} helps the user perform a specific task. It runs
on top of the system software and is written for the user's work rather than
to manage the machine. Word processors, spreadsheets, web browsers, media
players, and accounting packages are all application software.

Application software is often divided into two kinds. \textbf{General-purpose}
software is a ready-made package sold to many users, such as MS Word or MS
Excel; it is also called an \emph{application package}. \textbf{Specific-purpose}
software, also called \emph{tailor-made} or custom software, is built for one
organisation or one job, such as a payroll system or a railway reservation
system.

Application software covers many kinds of user work. Word processors prepare
documents; spreadsheets handle calculations and tables; presentation tools
build slide shows; database packages store and query records; graphics and
multimedia packages edit images, audio, and video; browsers and email clients
handle communication; and accounting packages keep business records. Most
modern application software has a graphical user interface (GUI), so the user
works with windows, icons, and menus instead of typing commands.

\begin{definitionbox}{Application software}
Software that is written to help the user perform a specific task or set of
tasks, running on top of the system software rather than managing the machine.
\end{definitionbox}

\begin{example}
The office desktop runs MS Windows and MS Office. Windows is system software
because it manages the hardware; MS Word and MS Excel are application software
because they help the user create documents and spreadsheets. The same machine
shows all the main categories in use at once.
\end{example}

\section{Utility software}
\textbf{Utility software} is a small program that performs maintenance or
housekeeping tasks on the system. Antivirus programs, disk cleanup tools,
disk defragmentation tools, backup utilities, and file-compression tools are
common examples. Utility software keeps the system healthy, fast, and safe,
and it belongs to the system software family.

A simple way to keep utilities and applications apart is to ask what the
program is for. A utility \emph{maintains the system} (for example, scanning
for viruses or cleaning temporary files), whereas an application helps the
user do \emph{work} (for example, typing a letter or preparing a budget).

\begin{definitionbox}{Utility software}
A small program within the system software family that carries out
maintenance and housekeeping tasks such as virus scanning, disk cleanup,
defragmentation, backup, and compression.
\end{definitionbox}

\begin{example}
A user notices that the office desktop has become slow and that the free disk
space keeps shrinking. Running a disk-cleanup tool to remove temporary files
and a disk-defragmentation tool to reorganise the hard disk restores speed and
space. Both tools are utility software: they maintain the system rather than
help the user produce a document.
\end{example}

\section{Programming software and language translators}
\textbf{Programming software} is used to write, translate, and test computer
programs. The core members are the \textbf{language translators}: a
\textbf{compiler}, an \textbf{interpreter}, and an \textbf{assembler}. A
compiler translates a whole program into machine code at once, while an
interpreter translates and runs the program line by line. An assembler
translates assembly language into machine code.

\begin{center}
\begin{tabular}{p{0.22\textwidth}p{0.64\textwidth}}
\toprule
\textbf{Translator} & \textbf{How it works} \\
\midrule
Compiler & Translates the whole program into machine code before it runs. \\
Interpreter & Translates and executes the program one line at a time. \\
Assembler & Translates assembly language (mnemonic codes) into machine code. \\
\bottomrule
\end{tabular}
\end{center}

Because a compiler translates the whole program first, it reports errors
together and the translated program generally runs faster; because an
interpreter works line by line, it can start running immediately and is
handy for testing small pieces of code.

Around the translators sit other supporting tools: the \textbf{linker}, which
joins separate pieces of a program; the \textbf{loader}, which brings a program
into memory to run; the \textbf{debugger}, which helps find errors; and the
text editor used to write the code. An \textbf{IDE} (Integrated Development
Environment) bundles an editor, a compiler or interpreter, and a debugger into
one package, which makes programming more convenient.

\section{Embedded software}
\textbf{Embedded software} is built into a device in order to control that
device. It is usually stored in read-only memory as \textbf{firmware}, runs
automatically when the device is switched on, and is written for one dedicated
job. Washing machines, microwave ovens, network routers, and the
engine-control units in modern cars all contain embedded software. The user
normally never sees this software directly; it simply makes the device behave
correctly.

\begin{definitionbox}{Embedded software}
Software built into a device and stored as firmware to control the device's
fixed, dedicated functions.
\end{definitionbox}

\section{Middleware}
\textbf{Middleware} is software that connects different applications and
systems so that they can work together. It sits between the operating system
and the application software and provides common services such as messaging,
authentication, and data exchange. An application server that links a web
front end to a back-end database is a typical example. Middleware allows
programs that were built separately to communicate without being rewritten
for each other.

\begin{definitionbox}{Middleware}
Software that sits between the operating system and applications and connects
different programs or systems, providing common services such as messaging,
authentication, and data exchange.
\end{definitionbox}

Middleware is common in large or distributed systems. On a web application,
for example, the browser (client), the web server, and the database are three
different programs. Middleware passing requests and replies between them lets
these separate programs behave as one service, without each program needing to
know the internal details of the others.

\section{Software licensing models}
A software \textbf{licence} states how a program may be used, copied, modified,
and distributed. The main models are summarised below.

\begin{center}
\begin{tabular}{p{0.24\textwidth}p{0.62\textwidth}}
\toprule
\textbf{Model} & \textbf{Key point} \\
\midrule
Proprietary & Source code is kept closed; use is governed by the vendor's licence. \\
Open-source & Source code is available to view, modify, and redistribute under its licence. \\
Free software & Gives the freedom to run, study, modify, and share the software. \\
Freeware & Distributed free of cost, but the source code may not be available. \\
Shareware & Offered as a limited trial; payment is needed for full or continued use. \\
Public domain & Has no copyright and may be used freely for any purpose. \\
Commercial & Sold for profit; usually proprietary. \\
\bottomrule
\end{tabular}
\end{center}

\textbf{Proprietary software} keeps its source code secret and restricts
copying and modification. MS Windows and MS Office are examples.
\textbf{Open-source software} publishes its source code and its licence allows
users to view, change, and share it; Linux, Mozilla Firefox, and LibreOffice
are examples. Note that ``open-source'' refers to \emph{source-code
availability}, not to price. Common open-source licences include the
\textbf{GNU General Public License (GPL)}, the \textbf{Apache License}, and the
\textbf{MIT License}; each sets out the conditions under which the source code
may be used and changed.

Proprietary software is usually supplied under an \textbf{EULA} (End User
License Agreement), which the user accepts before installing the product and
which forbids copying, modifying, or redistributing the program. Because the
source code is closed, only the vendor can improve or fix the software.

\textbf{Free software} is a related but distinct idea: it refers to freedom,
not to price. Free software gives the user the freedom to run the program,
study and modify it, and redistribute copies. It is important to see that both
free software and open-source software make the source code available; they
differ mainly in the philosophy they express, with free software emphasising
the user's freedom. \textbf{Freeware}, by contrast,
simply means free of cost; its source code may remain closed.
\textbf{Shareware} is a trial version distributed free for a limited period,
after which the user must pay for continued or full use. \textbf{Public-domain
software} carries no copyright at all and may be used freely.
\textbf{Commercial software} is sold for profit; it is usually also
proprietary, although a commercial product can occasionally be open-source.

\begin{definitionbox}{Open-source vs freeware vs shareware}
Open-source means the \emph{source code} is available; freeware means the
software is free of \emph{cost} (source may be closed); shareware is a
\emph{trial} that must later be paid for. The three are not the same.
\end{definitionbox}

\begin{example}
An item states, ``Freeware, open-source, and shareware all mean that the
source code is freely available.'' This is wrong. Only open-source guarantees
source-code availability. Freeware is free of cost but may be closed-source,
and shareware is a trial that later requires payment. The wrong option has
merged three different licensing models (TRAP-23).
\end{example}

\section{System software vs application software}
The two broadest groups are worth comparing side by side, because many exam
items turn on exactly this distinction.

\begin{center}
\begin{tabular}{p{0.28\textwidth}p{0.30\textwidth}p{0.30\textwidth}}
\toprule
\textbf{Point} & \textbf{System software} & \textbf{Application software} \\
\midrule
Purpose & Runs and manages the machine & Helps the user do a task \\
Runs & Close to the hardware; starts at boot & On top of the system software \\
Needed for & The computer to work at all & Performing the user's work \\
Examples & Operating system, device driver, BIOS & MS Word, MS Excel, a web browser \\
\bottomrule
\end{tabular}
\end{center}

The same physical computer can run one operating system (system software) and
many applications at the same time. The operating system is what makes the
hardware usable; the applications are what the user actually works with.

\section{Exam retrieval and common confusions}
Five distinctions prevent most errors on this topic.
\begin{enumerate}
  \item \textbf{System vs application}: system software manages the hardware;
    application software helps the user perform a task. The operating system
    is system software, while MS Word is application software.
  \item \textbf{Utility vs application}: a utility \emph{maintains} the system
    (antivirus, disk cleanup); an application helps the user do \emph{work}.
  \item \textbf{A device driver is software}, not hardware; the operating
    system is system software (TRAP-10).
  \item \textbf{Open-source $\neq$ freeware $\neq$ shareware} (TRAP-23):
    open-source gives the source code, freeware is free of cost, and shareware
    is a paid trial.
  \item \textbf{Free software is about freedom, not price}; it should not be
    confused with freeware, which is about price alone.
\end{enumerate}

\begin{example}
An item asks which of a compiler, an antivirus program, MS Excel, and the
operating system is \emph{not} system software. Trace the categories: a
compiler is programming software (system family), an antivirus program is
utility software (system family), and the operating system is system software.
MS Excel is application software, so it is the odd one out.
\end{example}

\subsection*{Exercises}
\begin{enumerate}[label=\arabic*.]
  \item Define software and distinguish it from hardware.
  \item List the six main categories of software and give one example of each.
  \item Distinguish system software from application software, with one
    example of each.
  \item Explain the difference between utility software and application
    software.
  \item Distinguish open-source software, freeware, and shareware, and state
    which one guarantees access to the source code.
  \item Classify each of the following as system, application, utility,
    programming, embedded, or middleware software: MS Word; an antivirus
    program; a compiler; the firmware in a router; an application server.
\end{enumerate}

\subsection*{Solutions}
\begingroup\small
\begin{enumerate}[label=\arabic*.]
  \item Software is a set of programs (instructions) that tells the computer
    what to do; hardware is the physical equipment that runs those
    instructions.
  \item System software (operating system), application software (MS Word),
    utility software (antivirus), programming software (compiler), embedded
    software (washing-machine firmware), and middleware (application server).
  \item System software manages and controls the hardware and provides a
    platform; application software helps the user perform a specific task.
    The operating system is system software; MS Excel is application
    software.
  \item Utility software performs maintenance and housekeeping tasks such as
    virus scanning or disk cleanup; application software helps the user carry
    out work such as writing a document or preparing a budget.
  \item Open-source software publishes its source code and allows it to be
    viewed, modified, and shared; freeware is free of cost but may keep its
    source code closed; shareware is a trial that later requires payment.
    Only open-source software guarantees access to the source code.
  \item MS Word is application software; an antivirus program is utility
    software; a compiler is programming software; the firmware in a router is
    embedded software; and an application server is middleware.
\end{enumerate}
\endgroup

\lessonheading{19}{Operating System}
\section{What an operating system is}
A computer's hardware cannot be used directly by a person or by ordinary
application programs. Something must sit between the user and the machine,
accept the user's instructions, and control the hardware on the user's behalf.
That something is the \textbf{operating system}. An operating system (OS) is
the main piece of \textbf{system software} that controls the whole computer. It
starts the machine, runs other programs, and manages every resource. It makes
the computer usable, because without an OS the hardware alone cannot run
application software.

\begin{definitionbox}{Operating system}
The main system software that sits between the user and the computer hardware,
provides a user-friendly interface, and manages the computer's resources such
as the processor, memory, storage, and input/output devices.
\end{definitionbox}

\begin{center}
\begin{tabular}{p{0.24\textwidth}p{0.66\textwidth}}
\toprule
\textbf{Term} & \textbf{Meaning in this lesson} \\
\midrule
Kernel & The core part of the OS that interacts directly with the hardware. \\
Booting & Starting the computer and loading the OS into main memory. \\
GUI & Graphical User Interface --- interaction through icons, windows, and clicks. \\
CLI & Command Line Interface --- interaction by typing text commands. \\
Process & A program that is currently executing. \\
Multitasking & Running more than one program at a time. \\
Time-sharing & Sharing one computer among several users at the same time. \\
\bottomrule
\end{tabular}
\end{center}

The rest of this lesson uses these terms in one consistent sense. The
\textbf{kernel} is the privileged core of the OS. \textbf{Booting} is the
startup sequence. A \textbf{GUI} is a graphical interface and a \textbf{CLI}
is a text-command interface. A \textbf{process} is a running program.
\textbf{Multitasking} and \textbf{time-sharing} are two ways of sharing the
CPU.

\section{The OS as interface and resource manager}
The operating system has \textbf{two main jobs}. The first is to provide a
\textbf{user-friendly interface}: the user gives instructions and receives
results through the OS, without needing to understand the hardware. This is
the definition of the OS as an interface layer, and it is a regularly tested
idea (PYQ-21). The second job is to act as a \textbf{resource manager}: the OS
shares the CPU, main memory, storage, and devices among all the programs that
are running, so that each one gets a fair turn and none interferes with the
others.

It is important to know \textbf{where} the operating system is. The OS is
stored permanently on secondary storage, such as an SSD or HDD. At startup it
is loaded into \textbf{primary storage}, that is, main memory. While the
computer is running, the OS is \emph{resident in primary storage} (PYQ-22): it
stays in main memory for the whole session so that it can respond instantly
whenever the hardware or an application needs it.

\begin{example}
An item states that the OS is kept in secondary storage while the computer is
running. This is wrong. The OS is stored on secondary storage when the machine
is off, but once the computer starts it is copied into primary storage (main
memory) and remains there while running. The wrong option confuses the place
where the OS is stored with the place where it runs.
\end{example}

\section{Functions of an operating system}
The work of an OS is usually grouped into a few large functions, each of which
is a favourite exam topic:

\begin{enumerate}
  \item \textbf{Process management}: creating processes, scheduling the CPU
    among them, and removing them when they finish.
  \item \textbf{Memory management}: deciding how much main memory each
    program gets and keeping programs apart.
  \item \textbf{File management}: organising data into files and folders and
    handling their creation, reading, writing, and deletion.
  \item \textbf{Device management}: controlling input and output devices
    through device drivers.
  \item \textbf{Security}: managing user accounts, logins, and permissions, and
    applying updates.
  \item \textbf{User interface}: presenting either a GUI or a CLI to the user.
\end{enumerate}

\section{The kernel}
The \textbf{kernel} is the \textbf{core} of the operating system. It is the
part of the OS that interacts \textbf{directly with the hardware}, and it
always stays in main memory while the computer runs. The kernel provides the
most important services: process scheduling, memory handling, and device
control. Other software, including the user interface, relies on the kernel to
reach the hardware.

\begin{definitionbox}{Kernel}
The central, privileged part of the operating system that manages the hardware
and provides the fundamental services on which all other software depends.
\end{definitionbox}

\section{User interfaces: GUI and CLI}
Two broad kinds of interface are used to work with a computer. A
\textbf{Graphical User Interface (GUI)} lets the user interact through
pictures, icons, windows, menus, and clicks; it is easier for beginners. A
\textbf{Command Line Interface (CLI)} requires the user to type text commands;
it gives precise control but demands that the commands be known. Windows,
macOS, and Linux desktops are examples of GUIs, while the command prompt and
the terminal are examples of CLIs.

\begin{center}
\begin{tabular}{p{0.24\textwidth}p{0.30\textwidth}p{0.30\textwidth}}
\toprule
\textbf{Point} & \textbf{GUI} & \textbf{CLI} \\
\midrule
Full form & Graphical User Interface & Command Line Interface \\
Input method & Clicks, icons, windows & Typed text commands \\
Ease of use & Easier for beginners & Needs the commands to be known \\
Example & Windows desktop, icons & Command prompt, terminal \\
\bottomrule
\end{tabular}
\end{center}

\section{Booting the computer}
\textbf{Booting} is the process of starting a computer and loading the
operating system into main memory. A \textbf{cold boot} starts from the
power-off state, while a \textbf{warm boot} restarts a computer that is already
on. The startup firmware, \textbf{BIOS} (Basic Input/Output System) or its
modern replacement \textbf{UEFI}, first checks the hardware and then finds and
loads the operating system from secondary storage. BIOS/UEFI initialises the
hardware and hands control to the OS; it is not the kernel and it is not a
device driver.

\section{Process management}
A \textbf{process} is a program that is currently executing. A program sitting
on disk is passive; once it is loaded and running it becomes an active
process. The operating system creates processes, gives each a turn on the CPU
through \textbf{CPU scheduling}, lets them communicate where necessary, and
removes them when they finish. Scheduling decides which process runs next and
for how long, which is what allows several programs to make progress.

\section{Memory management}
The operating system decides \textbf{how much main memory} each program
receives. It keeps one program's memory \textbf{separate} from another's so
that programs cannot overwrite each other's data. It frees memory when a
program ends so that the space can be reused. Where needed, the OS uses
\textbf{virtual memory} to give programs more space than the physical RAM
provides.

\section{File management}
The OS organises data on storage devices into \textbf{files} and
\textbf{folders}. It handles creating, reading, writing, renaming, and deleting
files, and it keeps a \textbf{directory} that records where each file is
stored. It also controls \textbf{permissions}: which accounts may open, change,
or delete a given file. Together these tasks are called the \textbf{file
system}, the way the OS names, stores, and finds files.

\section{Device management}
The operating system controls all \textbf{input/output devices}: the keyboard,
mouse, printer, disk, and so on. Each device communicates with the OS through a
\textbf{device driver}. A device driver is \textbf{software}, not hardware ---
a point that is regularly tested (TRAP-10). The OS also uses \textbf{spooling}
to queue print jobs and feed the printer at the printer's own speed, so that
the computer does not have to wait for a slow device to finish.

\section{Multitasking and time-sharing}
\textbf{Multitasking} means that the operating system runs \textbf{more than
one program at a time} (PYQ-23). The CPU switches between programs so quickly
that they appear to run together, even though a single CPU core does only one
thing at any instant. \textbf{Time-sharing} lets \textbf{many users} share one
computer at the same time, each receiving a slice of CPU time in turn. Both
multitasking and time-sharing depend on the OS \textbf{scheduler}.

\begin{example}
An item defines multitasking as using more than one monitor, or as storing more
than one file on a disk. Both are wrong. Multitasking is about the
\textbf{processor} running more than one program at a time. The number of
monitors and the number of files stored are unrelated to multitasking.
\end{example}

\section{Security, user accounts and updates}
The operating system provides \textbf{user accounts}, each with a username and
password. \textbf{Authentication} (the login step) verifies who is using the
computer. \textbf{Permissions} decide which files and settings each account may
use, so that one user cannot disturb another's work. Finally, system and
software \textbf{updates} should be applied regularly, because they fix
security holes and add improvements. Keeping the OS up to date is one of the
simplest ways to protect the machine.

\section{Examples of operating systems}
Common operating systems include \textbf{Microsoft Windows}, the open-source
\textbf{Linux} family, and Apple's \textbf{macOS} on desktops and laptops, plus
the mobile operating systems \textbf{Android} and \textbf{iOS}. Android is
built on the Linux kernel. An application such as MS Word or a web browser is
\emph{not} an operating system; it runs on top of one.

\section{Exam retrieval and common confusions}
The following distinctions prevent most errors on this topic.
\begin{enumerate}
  \item The OS is \textbf{system software}, not application software. A device
    driver is \textbf{software}, not hardware (TRAP-10).
  \item While running, the OS is resident in \textbf{primary storage} (main
    memory); it is only \emph{stored} on secondary storage (PYQ-22).
  \item The OS provides a \textbf{user-friendly interface} between the user and
    the hardware (PYQ-21).
  \item \textbf{Multitasking} means running more than one program at a time;
    \textbf{time-sharing} shares one machine among several users (PYQ-23).
  \item \textbf{BIOS/UEFI} initialises the hardware and loads the OS; it is not
    the kernel and not a device driver (TRAP-24).
  \item A \textbf{program} on disk is passive; once running it becomes a
    \textbf{process}.
\end{enumerate}

\begin{example}
An item asks where the operating system is loaded when the computer starts.
Trace the startup sequence: the machine is switched on, BIOS/UEFI checks the
hardware, and the OS is loaded from secondary storage into \textbf{primary
storage} (main memory), where it remains while the computer runs. If the stem
instead asks which component is the core of the OS, the answer is the
\textbf{kernel}.
\end{example}

\subsection*{Exercises}
\begin{enumerate}[label=\arabic*.]
  \item Define an operating system and state its two main roles.
  \item Where is the operating system stored, and where is it resident while
    the computer is running?
  \item List four core functions of an operating system.
  \item What is the kernel, and why is it important?
  \item Distinguish a GUI from a CLI.
  \item What does multitasking mean, and how does it differ from
    time-sharing?
  \item State three ways an operating system helps to keep a computer secure.
\end{enumerate}

\subsection*{Solutions}
\begingroup\small
\begin{enumerate}[label=\arabic*.]
  \item An operating system is the main system software that sits between the
    user and the hardware. Its two main roles are to provide a user-friendly
    interface and to manage the computer's resources.
  \item It is stored permanently on secondary storage (such as an SSD or HDD)
    and is resident in primary storage (main memory) while the computer is
    running.
  \item Any four of: process management, memory management, file management,
    device management, security, and the user interface.
  \item The kernel is the core part of the operating system that interacts
    directly with the hardware. It is important because it provides the
    fundamental services --- scheduling, memory handling, and device control
    --- on which all other software depends.
  \item A GUI (Graphical User Interface) uses icons, windows, and clicks and is
    easier for beginners; a CLI (Command Line Interface) requires the user to
    type text commands and gives precise control.
  \item Multitasking means the OS runs more than one program at a time by
    switching the CPU among them. Time-sharing means many users share one
    computer at the same time, each getting a slice of CPU time.
  \item Three of: user accounts with usernames and passwords; authentication at
    login; permissions that limit access to files and settings; and regular
    security updates.
\end{enumerate}
\endgroup

\lessonheading{20}{Windows Explorer}
\section{What Windows Explorer is}
\textbf{Windows Explorer} is the file-management program built into the
Windows operating system. Its job is to let the user see and manage
everything stored on the computer: drives, folders, and files. In Windows 10
and Windows 11 the same program is presented to the user as \textbf{File
Explorer}. Internally it is the same program, whose executable file is
\textbf{explorer.exe}. From one Explorer window a user can browse the entire
tree of storage and perform the everyday operations of copying, moving,
renaming, and deleting files.

\begin{definitionbox}{Windows Explorer (File Explorer)}
The file-management program of Windows that displays drives, folders, and
files in a graphical window and lets the user copy, move, rename, and delete
them. It is launched from the executable \texttt{explorer.exe}.
\end{definitionbox}

The main parts of the Explorer window are:
\begin{center}
\begin{tabular}{p{0.30\textwidth}p{0.62\textwidth}}
\toprule
\textbf{Part} & \textbf{What it shows} \\
\midrule
Navigation pane & The tree of drives and folders on the left. \\
Address bar & The full path of the folder that is currently open. \\
File list & The files and subfolders inside the current folder. \\
Ribbon / command bar & The tabs (New, Cut, Copy, Paste, Delete, View) for actions. \\
Search box & Files found inside the current folder matching the typed text. \\
\bottomrule
\end{tabular}
\end{center}

\section{Drives, folders, subfolders, and files}
Storage is organised as a tree. At the top is a \textbf{drive}, identified by
a drive letter. The hard disk is usually \textbf{C:}, and a USB flash drive
might appear as \textbf{E:} or \textbf{F:}. A drive contains
\textbf{folders}. A folder (also called a \textbf{directory}) is a container
that holds files and other folders. A folder placed inside another folder is
a \textbf{subfolder}. A \textbf{file} is a single item of stored data --- a
document, an image, a song, or a program.

\begin{definitionbox}{Drive}
A storage device or partition on which folders and files are kept, named by
a drive letter such as \texttt{C:} or \texttt{E:}.
\end{definitionbox}

\begin{definitionbox}{Folder (directory)}
A named container that holds files and other folders. A folder inside a
folder is called a subfolder.
\end{definitionbox}

Because folders can hold subfolders, and those can hold more subfolders, the
whole arrangement looks like an inverted tree --- the drive is the root, the
folders are branches, and the files are the leaves. The top-most folder of a
drive, directly under the drive letter, is called the \textbf{root folder}.

\section{The file path}
Every file and folder has one exact \textbf{path}: the address that tells
Windows where it lives. The parts of a path are separated by a
\textbf{backslash} (\verb|\|). For example:
\[
\texttt{C:\textbackslash Users\textbackslash Student\textbackslash
Documents\textbackslash report.docx}
\]
Reading left to right, \texttt{C:} is the drive, \texttt{Users},
\texttt{Student}, and \texttt{Documents} are the successive folders, and
\texttt{report.docx} is the file name including its extension. The path is
shown in the Explorer \textbf{address bar}, and the same information is given
in a file's Properties as its \textbf{location}.

\begin{definitionbox}{Path}
The complete address of a file or folder, from the drive letter through the
chain of folders to the item's own name, with parts separated by a backslash.
\end{definitionbox}

\section{File names and extensions}
A file name normally has two parts joined by a dot: the \textbf{name} and the
\textbf{extension}. In \texttt{report.docx}, \texttt{report} is the name and
\texttt{.docx} is the extension. The extension tells Windows which program
should open the file and what kind of data it holds. Common extensions
include \texttt{.txt} for a text file, \texttt{.docx} for a Word document,
\texttt{.xlsx} for an Excel workbook, \texttt{.pptx} for a PowerPoint
presentation, \texttt{.pdf} for a portable document, \texttt{.jpg} for a
photo, and \texttt{.exe} for a program.

\begin{center}
\begin{tabular}{p{0.22\textwidth}p{0.62\textwidth}}
\toprule
\textbf{Extension} & \textbf{Type of file / format} \\
\midrule
\texttt{.txt} & Plain text file \\
\texttt{.docx} & Microsoft Word document \\
\texttt{.xlsx} & Microsoft Excel workbook \\
\texttt{.pptx} & Microsoft PowerPoint presentation \\
\texttt{.pdf} & Portable Document Format \\
\texttt{.jpg} & JPEG image \\
\texttt{.exe} & Executable program \\
\bottomrule
\end{tabular}
\end{center}

\begin{definitionbox}{Extension}
The letters after the last dot in a file name (for example \texttt{.docx})
that indicate the file's type and the program that opens it.
\end{definitionbox}

One common confusion is between an \emph{extension} and a \emph{format}. An
extension is the label at the end of the name, such as \texttt{.png}; the
format is the way the data is actually encoded, such as PNG. The two usually
match, but the words name different things, so they should not be used as
interchangeable (TRM-09, TRAP-34).

\section{Hidden files and attributes}
Every file and folder carries \textbf{attributes}: settings that tell Windows
how to treat it. The common attributes are \textbf{Read-only}, \textbf{Hidden},
\textbf{System}, and \textbf{Archive}. A \textbf{hidden file} or folder has
the \textbf{Hidden} attribute set and is therefore not displayed by default.
Windows hides important system files so that a user does not move or delete
them by mistake. To display hidden items, open the \textbf{View} tab, choose
\textbf{Show}, and then \textbf{Hidden items} (in Windows 10 and 11). In older
versions the setting is reached through \textbf{Folder Options}.

\begin{definitionbox}{File attribute}
A property that controls how Windows treats a file or folder, such as
Read-only, Hidden, System, or Archive.
\end{definitionbox}

\section{Properties}
\textbf{Properties} is the dialog box that displays the details of a selected
file, folder, or drive. It is opened by pressing \textbf{Alt}\,+\,
\textbf{Enter} or by choosing \textbf{right-click $\rightarrow$ Properties}.
For a file, Properties shows the type (format), the size, the date created,
the date modified, the location (path), and the attributes. For a drive, it
shows the total size, the used space, and the free space. Properties is a
read-and-sometimes-edit view: changes to attributes may be made there, but
the file's contents are not shown.

\section{Selecting files}
Before a file can be copied, moved, renamed, or deleted it must be selected.
A single file is selected by clicking it once. A \textbf{contiguous range} of
files is selected by clicking the first file, holding \textbf{Shift}, and
clicking the last file. \textbf{Non-contiguous} files are selected by holding
\textbf{Ctrl} and clicking each file in turn. All files in a folder are
selected with \textbf{Ctrl}\,+\,\textbf{A}.

\begin{center}
\begin{tabular}{p{0.42\textwidth}p{0.50\textwidth}}
\toprule
\textbf{Selection} & \textbf{Method} \\
\midrule
One file & Click once. \\
A block of files in a row & Click the first, hold Shift, click the last. \\
Separate files scattered in the list & Hold Ctrl and click each. \\
Every file in the folder & Ctrl\,+\,A. \\
\bottomrule
\end{tabular}
\end{center}

\section{Copy}
\textbf{Copy} makes a duplicate of the selected item while leaving the
original where it is. The usual method is to press \textbf{Ctrl}\,+\,\textbf{C}
to copy, open the destination folder, and press \textbf{Ctrl}\,+\,\textbf{V}
to paste. The same steps can be done with \textbf{right-click $\rightarrow$
Copy} followed by \textbf{Paste}. Holding \textbf{Ctrl} while dragging an
item also copies it. The key point about copy is that after the operation
there are \emph{two} copies: the original and the new one.

\begin{example}
A student copies \texttt{report.docx} from \texttt{Documents} to a USB drive.
After pasting, the file exists in both places. The copy in
\texttt{Documents} is unchanged, because Copy never removes the original.
\end{example}

\section{Move}
\textbf{Move} shifts the selected item to a new location and removes it from
the old one. The standard way is to \textbf{cut} with
\textbf{Ctrl}\,+\,\textbf{X}, open the destination, and \textbf{paste} with
\textbf{Ctrl}\,+\,\textbf{V}. Cut followed by Paste is therefore the standard
move operation. Alternatively, dragging a file to a folder on the
\textbf{same drive} moves it. Dragging to a folder on a \textbf{different
drive} copies it instead, because Windows assumes you intend to transfer a
copy between separate devices.

\begin{definitionbox}{Move}
Relocating a file or folder from one place to another so that it no longer
exists at the old location. Cut\,+\,Paste performs a move.
\end{definitionbox}

\section{Rename}
\textbf{Rename} changes only the \textbf{name} of a file or folder; its
contents are untouched. The quickest method is to select the item and press
\textbf{F2}, type the new name, and press \textbf{Enter}. The same command is
available by right-clicking the item and choosing \textbf{Rename}. A warning
worth remembering: changing the text after the dot can change the file's
recognised type, so the extension should be left alone unless the type is
genuinely being changed.

\section{Delete and the Recycle Bin}
Pressing the \textbf{Delete} key on a selected item moves it to the
\textbf{Recycle Bin}. The Recycle Bin is a special folder that temporarily
holds deleted items, giving the user a chance to recover them. An item in the
Recycle Bin may be \textbf{restored} to its original location; the
\textbf{Restore} command puts it back. Choosing \textbf{Empty Recycle Bin}
deletes everything held in the Bin permanently.

\begin{definitionbox}{Recycle Bin}
A special Windows folder that temporarily stores files and folders deleted
with the Delete key, so that they can be restored or removed permanently.
\end{definitionbox}

\section{Recycle Bin vs permanent deletion}
An item deleted with \textbf{Delete} is recoverable, because it sits in the
Recycle Bin until the Bin is emptied. Permanent deletion happens when the
user presses \textbf{Shift}\,+\,\textbf{Delete}, when the Recycle Bin is
emptied, or when an item is deleted from a \textbf{USB drive}, a
\textbf{network drive}, or another location that does not use the Recycle
Bin. Once deleted permanently, a file cannot be restored through the normal
Windows interface.

\begin{center}
\begin{tabular}{p{0.30\textwidth}p{0.30\textwidth}p{0.30\textwidth}}
\toprule
\textbf{Point} & \textbf{Recycle Bin (Delete)} & \textbf{Permanent deletion} \\
\midrule
Shortcut & Delete key & Shift\,+\,Delete \\
Recoverable & Yes, can be restored & No, cannot be restored \\
Stored where & Recycle Bin folder & Removed from the disk \\
\bottomrule
\end{tabular}
\end{center}

Permanent deletion does not mean the storage space is instantly overwritten;
the freed space is simply marked as available for reuse. For exam purposes
the distinction is the one that matters: Delete is recoverable, Shift+Delete
is not.

\section{Keyboard shortcuts and drag-and-drop}
The following shortcuts are high-yield for the exam and should be memorised
exactly.

\begin{center}
\begin{tabular}{p{0.40\textwidth}p{0.52\textwidth}}
\toprule
\textbf{Action} & \textbf{Shortcut} \\
\midrule
Copy & Ctrl\,+\,C \\
Cut (move) & Ctrl\,+\,X \\
Paste & Ctrl\,+\,V \\
Undo & Ctrl\,+\,Z \\
Select all & Ctrl\,+\,A \\
Rename & F2 \\
Refresh & F5 \\
Properties & Alt\,+\,Enter \\
New folder & Ctrl\,+\,Shift\,+\,N \\
Delete to Recycle Bin & Delete \\
Permanent delete & Shift\,+\,Delete \\
\bottomrule
\end{tabular}
\end{center}

Drag-and-drop follows a fixed rule. Dragging an item to a folder on the
\textbf{same drive moves} it; dragging to a folder on a \textbf{different
drive copies} it. Holding \textbf{Ctrl} while dragging always \textbf{copies},
and holding \textbf{Shift} while dragging always \textbf{moves}, regardless
of the drives involved.

\section{Exam retrieval and common confusions}
Five distinctions prevent most errors on this topic.
\begin{enumerate}
  \item \textbf{Copy} leaves the original in place; \textbf{move}
    (Cut\,+\,Paste) removes it. Do not confuse the two.
  \item The \textbf{Delete} key sends a file to the \textbf{Recycle Bin}; it
    is not a permanent erase. \textbf{Shift\,+\,Delete} is the permanent
    erase and bypasses the Bin.
  \item Deleting from a \textbf{USB drive} or \textbf{network drive} bypasses
    the Recycle Bin even when the plain Delete key is used.
  \item An \textbf{extension} is the label after the dot (for example
    \texttt{.png}); it is not the same as the \textbf{format} (PNG)
    (TRM-09).
  \item Same-drive drag \textbf{moves}; different-drive drag \textbf{copies}.
    \textbf{Ctrl+drag} forces a copy; \textbf{Shift+drag} forces a move.
\end{enumerate}

\begin{example}
An item asks where a file goes when the Delete key is pressed. Trace the rule:
the Delete key moves the item to the Recycle Bin, from which it can still be
restored. If the stem instead says the file must be removed so that it cannot
be recovered, the required operation is Shift+Delete (or emptying the Recycle
Bin), not the plain Delete key.
\end{example}

\subsection*{Exercises}
\begin{enumerate}[label=\arabic*.]
  \item What is Windows Explorer called in Windows 11, and what is the name
    of its executable program?
  \item Distinguish between a folder and a subfolder, and define a file path.
  \item What is a file extension, and why is an extension not the same as a
    file format?
  \item Explain the difference between Copy and Move, and give the keyboard
    shortcut for each.
  \item A file is deleted with the Delete key. Can it be recovered? Explain,
    and state one way to delete a file permanently.
\end{enumerate}

\subsection*{Solutions}
\begingroup\small
\begin{enumerate}[label=\arabic*.]
  \item In Windows 11 it is called \textbf{File Explorer}. Its executable
    program is \textbf{explorer.exe}.
  \item A folder is a container that holds files and other folders; a folder
    placed inside another folder is a subfolder. A path is the complete
    address of a file or folder, from the drive letter through the chain of
    folders to the item's own name, with parts separated by a backslash.
  \item An extension is the letters after the last dot in a file name (for
    example \texttt{.docx}) that indicate the file's type. It is not the same
    as the format, because the extension is only the label at the end of the
    name, while the format is the way the data is actually encoded.
  \item \textbf{Copy} duplicates the item and leaves the original in place
    (Ctrl\,+\,C, then Ctrl\,+\,V). \textbf{Move} relocates the item and
    removes it from the old location (Ctrl\,+\,X, then Ctrl\,+\,V).
  \item Yes, a file deleted with the Delete key can normally be recovered,
    because it goes to the Recycle Bin and can be restored. It can be deleted
    permanently with \textbf{Shift\,+\,$\boxed{\text{Delete}}$} (or by
    emptying the Recycle Bin).
\end{enumerate}
\endgroup

\lessonheading{21}{Keyboard Shortcuts}
\section{How to read a keyboard shortcut}
A keyboard shortcut is a key, or a combination of keys, that performs a
command directly without opening a menu. Shortcuts save time and reduce the
need for the mouse, and they are a favourite topic in the JKSSB Computer
Awareness paper because the answer is a single exact letter or key. In a
combination written with a plus sign, such as \textbf{Ctrl+C}, the plus means:
hold the first key down, then press the second key. You do not type the plus
sign. Where three keys are written, such as \texttt{Ctrl+Shift+N}, hold the
first two and tap the third.

\begin{definitionbox}{Keyboard shortcut}
A key or key combination that runs a command directly --- the same command a
menu offers --- so the user reaches it from the keyboard instead of with the
mouse.
\end{definitionbox}

The keys used to build combinations are called modifier keys. They are
\textbf{Ctrl} (Control), \textbf{Alt} (Alternate), \textbf{Shift}, and the
\textbf{Windows (Win)} key. A modifier key does very little on its own; it
changes what the next key does. The \textbf{function keys} \textbf{F1} through
\textbf{F12} run along the top row of the keyboard and are usually pressed on
their own, without a modifier.

\begin{center}
\begin{tabular}{p{0.16\textwidth}p{0.28\textwidth}p{0.48\textwidth}}
\toprule
\textbf{Key} & \textbf{Full form} & \textbf{Main role} \\
\midrule
Ctrl & Control & Acts on content: copy, cut, paste, save, find. \\
Alt & Alternate & Acts on windows and menus: switch, close. \\
Shift & Shift & Changes what is typed or selected; builds ranges. \\
Win & Windows key & Opens or controls system features: desktop, Explorer, lock. \\
Fn & Function (laptop) & Gives a second job to the F-keys on many laptops. \\
\bottomrule
\end{tabular}
\end{center}

A practical split helps memory. The \textbf{Ctrl} combinations mostly act on
\emph{content} --- copying, cutting, saving, finding. The \textbf{Alt}
combinations mostly act on \emph{windows} --- switching and closing. The
\textbf{Windows} key opens or controls system-level features such as the
desktop, File Explorer, and the lock screen. Function keys are single keys
whose job depends on the application that is active.

Most shortcut letters are chosen to be easy to remember. \textbf{C} stands for
Copy, and \textbf{X} looks like a cross over the text, so it stands for Cut.
\textbf{V} sits next to X on the keyboard and stands for Paste. \textbf{Z} is
Undo and \textbf{Y} is Redo, the two editing-history keys. \textbf{A} is All
(Select All), \textbf{S} is Save, \textbf{P} is Print, \textbf{F} is Find, and
\textbf{H} opens the find-and-replace box. Learning the letter is more reliable
than memorising a list, because the letter usually hints at the command.

\section{Copy, cut, and paste}
The three most common shortcuts form the clipboard group. The
\textbf{clipboard} is the temporary holding area in memory where copied or cut
items wait before being pasted.

\begin{center}
\begin{tabular}{p{0.20\textwidth}p{0.70\textwidth}}
\toprule
\textbf{Shortcut} & \textbf{Action} \\
\midrule
Ctrl+C & Copy the selected item; the original stays where it is. \\
Ctrl+X & Cut the selected item; the original is removed. \\
Ctrl+V & Paste the copied or cut item at the current position. \\
\bottomrule
\end{tabular}
\end{center}

\begin{definitionbox}{Clipboard}
A temporary area in the computer's memory that holds text, images, or files
after they are copied or cut, until they are pasted.
\end{definitionbox}

The distinction to hold is that \textbf{Copy} and \textbf{Cut} both place the
selection on the clipboard, but only \textbf{Cut} removes the original. A
question that offers \texttt{Ctrl+C} and \texttt{Ctrl+X} as two options is
testing exactly this difference: \texttt{Ctrl+C} leaves the original in place,
and \texttt{Ctrl+X} does not. \texttt{Ctrl+V} only reads the clipboard; it
never decides whether the item was copied or cut.

\begin{example}
A candidate reads, ``Which shortcut copies selected text without removing it
from its original place?'' The answer is \texttt{Ctrl+C}. \texttt{Ctrl+X} also
puts the text on the clipboard, but it deletes the selection from the
document, so it does not satisfy the phrase ``without removing it''.
\end{example}

\section{Undo, redo, and select all}
Three Ctrl shortcuts control the editing history and the current selection.

\begin{itemize}
  \item \textbf{Ctrl+Z} --- Undo the last action. It steps backward through
    recent changes.
  \item \textbf{Ctrl+Y} --- Redo the action that was just undone. It steps
    forward again.
  \item \textbf{Ctrl+A} --- Select All the content of the current window or
    document.
\end{itemize}

Undo and redo are a matched pair and must not be swapped. \texttt{Ctrl+Z}
reverses a change; \texttt{Ctrl+Y} reapplies a change that was reversed.
\texttt{Ctrl+A} does not select one word; it selects everything that can be
selected in the current context. In MS Office the redo key is \texttt{Ctrl+Y};
a few non-Office programs use \texttt{Ctrl+Shift+Z}, but for this exam the
standard answer is \texttt{Ctrl+Y}.

\section{Save, print, find, and replace}
These four shortcuts cover saving a file and working with text inside it.

\begin{center}
\begin{tabular}{p{0.20\textwidth}p{0.70\textwidth}}
\toprule
\textbf{Shortcut} & \textbf{Action} \\
\midrule
Ctrl+S & Save the current file. \\
Ctrl+P & Print. \\
Ctrl+F & Find (open the search box for text). \\
Ctrl+H & Replace (find a word and substitute another). \\
\bottomrule
\end{tabular}
\end{center}

\textbf{Find} and \textbf{Replace} are near neighbours. \texttt{Ctrl+F} only
locates the search text and highlights it; \texttt{Ctrl+H} opens the
find-and-replace function, which can also change every occurrence to a new
word. Similarly, \texttt{Ctrl+S} saves and \texttt{Ctrl+P} prints; the two are
often offered together as distractors. \texttt{Ctrl+S} does not print, and
\texttt{Ctrl+P} does not save.

\begin{example}
A user types a report and wants to change every occurrence of the word
``received'' to ``recd''. \texttt{Ctrl+F} would only find the word and stop
there. The correct shortcut is \texttt{Ctrl+H}, which opens Replace and lets
the user substitute the new text, one occurrence at a time or all at once.
\end{example}

\section{Alt-key combinations}
The Alt key builds the window-management shortcuts.

\begin{itemize}
  \item \textbf{Alt+Tab} --- Switch between open windows or running
    applications. Repeatedly pressing Tab while holding Alt cycles through the
    open windows.
  \item \textbf{Alt+F4} --- Close the active window or application.
\end{itemize}

On the desktop with no application window active, \texttt{Alt+F4} opens the
shutdown or sign-out dialog. The concept to keep clear is that \texttt{Alt+Tab}
\emph{switches} between windows and does not close any of them, while
\texttt{Alt+F4} \emph{closes} the active window. A common wrong answer reverses
the two. Note that \texttt{Alt+F4} closes the whole window; saving the file is
still the job of \texttt{Ctrl+S}.

\section{Function keys}
A function key is a single key, \texttt{F1} to \texttt{F12}, on the top row of
the keyboard. Its action depends on the application that is running. For the
exam, four function keys are high-yield.

\begin{center}
\begin{tabular}{p{0.14\textwidth}p{0.78\textwidth}}
\toprule
\textbf{Key} & \textbf{Standard action} \\
\midrule
F1 & Help. \\
F2 & Rename the selected file or folder (for example in File Explorer). \\
F5 & Refresh the window or folder; in PowerPoint, start the slide show from
     the beginning. \\
F7 & Spelling and grammar check (in MS Word). \\
\bottomrule
\end{tabular}
\end{center}

\begin{definitionbox}{Function key}
One of the keys F1 to F12 on the top row of the keyboard, whose effect is
decided by the active application rather than by a fixed system command.
\end{definitionbox}

In \textbf{PowerPoint}, \texttt{F5} starts a slide show from the beginning and
\texttt{Shift+F5} starts it from the current slide. In \textbf{Windows}, on a
folder or desktop, \texttt{F5} refreshes the view. \texttt{F7} in MS Word
launches the spelling and grammar check. A documented JKSSB item asked the key
that starts a slide show from the beginning, and the answer was \texttt{F5}
(PYQ WO 2026 Q80, provisional key).

\begin{example}
A candidate sees four keys: F1, F2, F5, F7. If the stem asks for Help, the
answer is F1; for Rename, F2; for Refresh or slide show, F5; for the spell
check, F7. Read the stem's verb first, then pick the key that matches that verb.
\end{example}

The table below fixes the near neighbours that JKSSB likes to mix. Each row is
one function key; the middle column is its job; the last column is the trap
(the job a careless candidate might wrongly attach to it).

\begin{center}
\begin{tabular}{p{0.10\textwidth}p{0.32\textwidth}p{0.50\textwidth}}
\toprule
\textbf{Key} & \textbf{Correct job} & \textbf{Common wrong answer} \\
\midrule
F1 & Help & Rename or Refresh \\
F2 & Rename a selected file or folder & Help or New \\
F5 & Refresh; start slide show & Spell check or Save \\
F7 & Spelling and grammar check & Slide show or Help \\
\bottomrule
\end{tabular}
\end{center}

\section{Windows-key shortcuts}
The Windows key (the key with the Windows logo) builds system-level shortcuts.

\begin{center}
\begin{tabular}{p{0.20\textwidth}p{0.70\textwidth}}
\toprule
\textbf{Shortcut} & \textbf{Action} \\
\midrule
Win+D & Show the desktop (minimise all open windows). \\
Win+E & Open File Explorer. \\
Win+L & Lock the computer. \\
\bottomrule
\end{tabular}
\end{center}

Mnemonic: \textbf{D} for Desktop, \textbf{E} for Explorer, \textbf{L} for
Lock. These three are system shortcuts, not content shortcuts, so they do not
depend on which file is open. \texttt{Win+D} does not delete files, and
\texttt{Win+L} does not shut the computer down; it only locks the session, and
the user must sign in again.

The \textbf{Shift} key is a modifier too. It produces capital letters and the
upper symbol of a key, and \texttt{Shift+Click} selects a range of items
between two clicks. \texttt{Shift+F5} starts a PowerPoint slide show from the
current slide. On its own with a letter, Shift does not save, close, or print;
it changes what is typed or selected.

\section{Other high-yield keys}
The shortcuts above are the core of the lesson. A few adjacent keys and
combinations appear often enough in Office documents that they are worth
learning alongside them.

\begin{center}
\begin{tabular}{p{0.22\textwidth}p{0.64\textwidth}}
\toprule
\textbf{Shortcut / key} & \textbf{Action} \\
\midrule
Ctrl+O & Open an existing file. \\
Ctrl+N & Create a new file (new document or window). \\
Ctrl+W & Close the current document or window. \\
Ctrl+B / Ctrl+I / Ctrl+U & Bold / Italic / Underline the selection. \\
Esc & Cancel the current action or close a dialog. \\
Delete & Delete the selected item; in text, erase forward. \\
Backspace & Erase the character to the left of the cursor. \\
Home / End & Move to the start / end of the line. \\
Ctrl+Home / Ctrl+End & Move to the top / bottom of the document. \\
Page Up / Page Down & Move up / down one screen. \\
\bottomrule
\end{tabular}
\end{center}

Two cautions. First, \texttt{Ctrl+O} opens and \texttt{Ctrl+N} creates; do not
confuse the two. Second, \texttt{Ctrl+W} closes a \emph{document}, while
\texttt{Alt+F4} closes the \emph{window or application}; in a browser,
\texttt{Ctrl+W} closes the current tab.

\section{Special keys that stand alone}
A few keys send a command on their own, so they are never written with a plus
sign. They are not combinations, but they are learned together with shortcuts.

\begin{center}
\begin{tabular}{p{0.20\textwidth}p{0.70\textwidth}}
\toprule
\textbf{Key} & \textbf{Action} \\
\midrule
Enter & Confirm a command or start a new line. \\
Esc & Cancel the current action or close a dialog. \\
Tab & Move to the next field, cell, or control. \\
Caps Lock & Toggle capital letters on or off. \\
Backspace & Erase the character to the left of the cursor. \\
Delete & Erase the selected item, or the character to the right. \\
Insert & Toggle Insert and Overtype modes in text editors. \\
Home / End & Move to the start or the end of the current line. \\
\bottomrule
\end{tabular}
\end{center}

\textbf{Caps Lock} and \textbf{Shift} both produce capitals, and this is a
common confusion. Caps Lock changes every letter until it is switched off
again, while Shift changes only the key pressed at that moment. \textbf{Delete}
removes the selection or the character ahead of the cursor; \textbf{Backspace}
removes what is behind the cursor.

\section{Shortcuts in the office desktop workflow}
A short walk-through on the office desktop ties the keys together. A clerk
opens an existing letter with \texttt{Ctrl+O}, types the body, and underlines a
heading with \texttt{Ctrl+U}. The heading is copied with \texttt{Ctrl+C} and
pasted elsewhere with \texttt{Ctrl+V}; a wrong sentence is removed with
\texttt{Ctrl+X}. A mistake is reversed with \texttt{Ctrl+Z}. The clerk searches
the letter for a name with \texttt{Ctrl+F} and changes an old address to a new
one with \texttt{Ctrl+H}. The finished letter is saved with \texttt{Ctrl+S} and
printed with \texttt{Ctrl+P}. While the letter is open, the clerk switches to
the spreadsheet with \texttt{Alt+Tab}, and when the letter is done, the window
is closed with \texttt{Alt+F4}. Later, \texttt{Win+D} shows the desktop,
\texttt{Win+E} opens File Explorer to find last month's folder, and \texttt{Win+L}
locks the machine before the clerk steps away. This single sequence uses almost
every shortcut in the lesson, which is why a scenario item is easy to build
from it.

\section{Exam retrieval and common confusions}
Seven distinctions prevent most errors on this topic.
\begin{enumerate}
  \item \textbf{Ctrl+C} copies and keeps the original; \textbf{Ctrl+X} cuts and
    removes the original. Do not treat them as the same command.
  \item \textbf{Ctrl+Z} is Undo; \textbf{Ctrl+Y} is Redo. Reversing this pair
    is a common wrong answer.
  \item \textbf{Ctrl+F} is Find; \textbf{Ctrl+H} is Replace. Find only locates;
    Replace also substitutes.
  \item \textbf{Ctrl+S} is Save; \textbf{Ctrl+P} is Print. Save stores the
    file; Print sends it to the printer.
  \item \textbf{Alt+Tab} switches between windows; \textbf{Alt+F4} closes the
    active window.
  \item \textbf{F1} is Help, \textbf{F2} is Rename, \textbf{F5} is Refresh or
    slide show, and \textbf{F7} is the spell check (TRAP-32). Do not move a
    job from one key to another.
  \item \textbf{Win+D} shows the desktop, \textbf{Win+E} opens File Explorer,
    and \textbf{Win+L} locks the computer.
\end{enumerate}

\begin{example}
An item asks which key is used to check spelling and grammar in MS Word. Trace
the function keys: F1 is Help, F2 is Rename, F5 is Refresh or slide show, F7 is
the spelling and grammar check. The answer is F7. If the stem instead asked for
Rename, the answer would be F2, and if it asked for Help, the answer would be
F1.
\end{example}

\section{Quick-reference sheet}
Keep this single table for last-minute revision. It gathers every shortcut in
this lesson.

\begin{center}
\begin{tabular}{p{0.15\textwidth}p{0.28\textwidth}p{0.15\textwidth}p{0.28\textwidth}}
\toprule
\textbf{Shortcut} & \textbf{Action} & \textbf{Shortcut} & \textbf{Action} \\
\midrule
Ctrl+C & Copy & Ctrl+S & Save \\
Ctrl+X & Cut & Ctrl+P & Print \\
Ctrl+V & Paste & Ctrl+F & Find \\
Ctrl+Z & Undo & Ctrl+H & Replace \\
Ctrl+Y & Redo & F1 & Help \\
Ctrl+A & Select All & F2 & Rename \\
Alt+Tab & Switch window & F5 & Refresh / slide show \\
Alt+F4 & Close window & F7 & Spell check \\
Win+D & Show desktop & Win+E & File Explorer \\
Win+L & Lock computer & Shift+F5 & Slide show from current \\
\bottomrule
\end{tabular}
\end{center}

\subsection*{Exercises}
\begin{enumerate}[label=\arabic*.]
  \item State the action of Ctrl+C, Ctrl+X, and Ctrl+V.
  \item Distinguish Ctrl+Z from Ctrl+Y, and Ctrl+F from Ctrl+H.
  \item What do Alt+Tab and Alt+F4 do, and how are they different?
  \item Give the standard action of F1, F2, F5, and F7.
  \item What do Win+D, Win+E, and Win+L do? What does Shift+F5 do in
    PowerPoint?
  \item Why is the plus sign in \texttt{Ctrl+C} not typed by the user?
\end{enumerate}

\subsection*{Solutions}
\begingroup\small
\begin{enumerate}[label=\arabic*.]
  \item Ctrl+C copies the selection (the original stays); Ctrl+X cuts it (the
    original is removed); Ctrl+V pastes the copied or cut item at the current
    position.
  \item Ctrl+Z undoes the last action, while Ctrl+Y redoes the action that was
    undone. Ctrl+F opens Find, which only locates text; Ctrl+H opens Replace,
    which can also substitute new text for the found text.
  \item Alt+Tab switches between open windows or applications without closing
    any of them; Alt+F4 closes the active window or application. One switches,
    the other closes.
  \item F1 is Help; F2 renames the selected file or folder; F5 refreshes the
    window (or, in PowerPoint, starts the slide show from the beginning); F7
    runs the spelling and grammar check in MS Word.
  \item Win+D shows the desktop, Win+E opens File Explorer, and Win+L locks
    the computer. Shift+F5 starts a PowerPoint slide show from the current
    slide.
  \item The plus sign is only a way of writing the combination. It means
    ``hold the first key, then press the second''; it is not a character to be
    typed.
\end{enumerate}
\endgroup

\lessonheading{45}{Open Source}
\section{Source code and what it means}
Every program begins as \textbf{source code}: the set of instructions a
programmer writes in a programming language such as C, Java, or Python. Source
code is \emph{human-readable} --- a person can open the file, read it, and
change it. Before the computer can run the program, a \textbf{compiler} or an
\textbf{interpreter} converts the source code into a machine-readable form.
That machine-readable form is called \textbf{object code}, or simply the
executable, and it is what the processor actually executes. Ordinary users of
a finished program normally receive only this compiled form and never see the
source code.

\begin{definitionbox}{Source code}
The human-readable set of instructions, written in a programming language,
that defines what a program does; it is converted into machine-readable
object code by a compiler or an interpreter.
\end{definitionbox}

The important consequence is that whoever holds the source code controls how
the software can be studied, changed, and shared. The whole topic of open
source versus proprietary software grows out of this single fact.

\section{Open source versus proprietary software}
Software falls into two broad camps according to what is done with its source
code.

\begin{center}
\begin{tabular}{p{0.22\textwidth}p{0.34\textwidth}p{0.34\textwidth}}
\toprule
\textbf{Point} & \textbf{Open source} & \textbf{Proprietary (closed source)} \\
\midrule
Source code & Publicly available & Kept secret \\
View or modify & Allowed, subject to the licence & Not allowed without permission \\
Redistribute & Allowed, subject to the licence & Not allowed \\
Typical user gets & Source code plus executable & Only the compiled executable \\
\bottomrule
\end{tabular}
\end{center}

\textbf{Open-source software} makes its source code public. Anyone is allowed
to inspect it, modify it, and redistribute it, subject to the terms of its
licence. \textbf{Proprietary software} keeps its source code confidential; the
user receives only the compiled executable and may not view, change, or share
the source without the owner's permission. Examples of proprietary software
include Microsoft Windows, Apple macOS, Microsoft Office, Adobe Photoshop, and
Oracle Database. The deciding factor between the two is the source code, not
the price: some open-source software is sold, and some proprietary software is
given away free of cost.

\begin{definitionbox}{Open source}
A software development model in which the source code is made publicly
available so that anyone may view, modify, and redistribute it under the
terms of its licence.
\end{definitionbox}

\section{The open-source idea: view, modify, share}
The meaning that examiners most often test is the three-part permission of an
open-source licence. First, a user may \textbf{view} the source code and see
exactly how the software works. Second, a user may \textbf{modify} it --- fix a
bug, improve it, or add a new feature. Third, a user may \textbf{share} it,
redistributing copies and modified versions to others. These permissions are
granted by the software's licence, so the exact conditions differ from one
project to another; but the underlying spirit is the same three freedoms in
the code.

\begin{example}
An item states, ``An open-source licence allows a user to view, change, and
share the source code.'' This is the standard, correct description and is the
answer expected in the exam. An option that says only ``software available
free of cost'' describes \emph{freeware}, not open source.
\end{example}

\section{Freeware, free software, public domain, commercial, shareware}
Several labels are often mixed up with open source. Each has a distinct
meaning.

\begin{center}
\begin{tabular}{p{0.20\textwidth}p{0.34\textwidth}p{0.36\textwidth}}
\toprule
\textbf{Term} & \textbf{What is free} & \textbf{Source code} \\
\midrule
Open source & Use, study, change, share & Available \\
Freeware & Cost only & Usually not provided \\
Free software & The user's freedom & Available \\
Shareware & A limited trial period & Not provided \\
Public domain & No copyright at all & No restrictions \\
Commercial & Nothing (it is sold) & Usually not provided \\
\bottomrule
\end{tabular}
\end{center}

\textbf{Freeware} is software given away free of cost, but its source code is
usually not provided. \textbf{Free software} is about freedom, not price: the
user is free to study, change, and share it. \textbf{Public domain} software
has no copyright owner, so anyone may use, modify, and redistribute it without
any restriction. \textbf{Shareware} is distributed as a free trial, working
for a limited period or with limited features; to keep using it, the user must
pay. \textbf{Commercial software} is developed and sold for profit and is
often, though not always, proprietary.

\begin{definitionbox}{Freeware}
Software distributed free of cost, whose source code is usually not made
available; it is free in price, not in freedom.
\end{definitionbox}

\begin{definitionbox}{Shareware}
Software distributed as a trial version, free for a limited time or with
limited features, after which payment is required to continue using it.
\end{definitionbox}

\section{Free software and its four freedoms}
The term \textbf{free software} comes from the \textbf{Free Software
Foundation (FSF)}, founded by \textbf{Richard Stallman}. In this context,
``free'' refers to liberty, not price --- the well-known phrase is ``free as in
freedom, not free as in free beer''. The FSF defines four freedoms: the freedom
to run the program for any purpose; the freedom to study how it works and
change it; the freedom to redistribute copies to help others; and the freedom
to distribute modified versions to the public. Free software and open source
are closely related and frequently overlap, but they are not identical labels.
Free software stresses the ethical freedom of the user; open source stresses
the practical benefits of a publicly available source code.

\begin{definitionbox}{Free software}
Software that grants the user the freedom to run, study, change, and
redistribute it and its modified versions; ``free'' here means freedom, not
zero cost.
\end{definitionbox}

\section{The Open Source Initiative}
Alongside the Free Software Foundation stands the \textbf{Open Source
Initiative (OSI)}, founded in 1998. The OSI maintains the \textbf{Open Source
Definition}, the set of criteria a licence must satisfy to be called an
open-source licence, and it reviews and approves licences against those
criteria. The term \textbf{open source} itself was popularised in 1998 as a
more business-friendly label than ``free software''. The two movements
overlap heavily: most licences approved by the OSI are also free-software
licences. An umbrella term often seen in questions is \textbf{FOSS}, meaning
Free and Open-Source Software, or \textbf{FLOSS} (Free/Libre and Open-Source
Software), which covers both traditions together.

The Open Source Definition lays down the conditions a licence must meet. The
software must be freely redistributable; the source code must be included or
easily obtainable; derived works and modifications must be allowed; the
licence must not discriminate against any person, group, or field of use; and
the rights must apply to everyone who receives the software. A licence that
satisfies these conditions is called an \textbf{OSI-approved} licence. This is
why ``open source'' is a matter of the licence, not simply of whether the
software is free of cost.

\begin{example}
Two programs are both free to download. Program A ships with its source code
under an OSI-approved licence that permits modification and redistribution;
Program B ships only compiled files and forbids copying. A is open source; B
is \emph{freeware}. The download price is identical, but the licences differ,
and the licence is what decides the category.
\end{example}

\section{Benefits and limitations of open source}
Open source is chosen for several practical reasons. The licence cost is
usually zero, so organisations avoid per-seat fees. The source code can be
inspected, which lets many people review it for bugs and security weaknesses.
Users are not locked in to a single vendor, and the software can be
customised to local needs. A large community often provides documentation and
support.

Open source also has limitations that a fair answer should acknowledge. A
popular project may have no formal commercial support contract, so support
depends on the community or on paid third-party vendors. Compatibility with
proprietary formats and integration with existing systems can be difficult.
The total cost of ownership is not always zero, because training, migration,
and maintenance still cost money. These points are occasionally tested as
``which of the following is a limitation of open source''.

\section{Full forms and key names to remember}
Direct-recall items often ask for an expansion or for the name behind a term.
The following are the ones most worth memorising.

\begin{center}
\begin{tabular}{p{0.22\textwidth}p{0.68\textwidth}}
\toprule
\textbf{Short form} & \textbf{Full form / meaning} \\
\midrule
GPL & GNU General Public License \\
FSF & Free Software Foundation \\
OSI & Open Source Initiative \\
FOSS & Free and Open-Source Software \\
FLOSS & Free/Libre and Open-Source Software \\
AOSP & Android Open Source Project \\
MIT & MIT License (a permissive software licence) \\
BSD & Berkeley Software Distribution (its licence family) \\
\bottomrule
\end{tabular}
\end{center}

Remember too the names behind the ideas: the Free Software Foundation was
founded by \textbf{Richard Stallman}, and \textbf{Git} was created by Linus
Torvalds to manage the Linux kernel's source code. \textbf{GitHub},
\textbf{GitLab}, \textbf{Bitbucket}, and \textbf{SourceForge} are the
best-known hosting sites for open-source repositories.

\section{Repositories and community development}
Open-source projects live in a \textbf{repository} (often shortened to
``repo''). A repository is a storage location that holds a project's source
code together with every change ever made to it. \textbf{Version control}
records this history, so an earlier version can be restored if needed.
\textbf{Git} is the most widely used distributed version-control system, and
popular hosting sites for repositories include \textbf{GitHub}, \textbf{GitLab},
\textbf{Bitbucket}, and \textbf{SourceForge}.

Open-source software is usually built by a \textbf{community} of contributors
spread around the world. A contributor may \textbf{fork} a project, that is,
make a personal copy for experimenting or improving. A \textbf{pull request}
is then a proposal to merge those changes back into the main project, where
\textbf{maintainers} review them. An \textbf{issue tracker} records bugs and
feature requests. Because many people can inspect and improve the code,
community review is treated as a key strength of open-source development.

\section{Basic licensing concepts}
By default, \textbf{copyright} gives the creator of a work exclusive rights
over it. A \textbf{licence} is the legal permission that states what others
may do with the software. Open-source and free-software licences fall into two
broad families.

\textbf{Copyleft} licences, such as the GNU General Public License (GPL),
require that modified versions of the software be released under the
\emph{same} licence. This keeps the software and its derivatives open.
\textbf{Permissive} licences, such as the MIT License, the Apache License 2.0,
and the BSD License, allow the code to be reused with very few conditions,
including inside proprietary projects, provided the licence notice is kept.
The licence, not the price, decides whether software is truly open source.

\begin{definitionbox}{Copyleft}
A licence condition that requires any modified or extended version of the
software to be distributed under the same licence, so the work cannot be
closed.
\end{definitionbox}

\section{Copyright, licences, and piracy}
It helps to separate three ideas. \textbf{Copyright} is the creator's legal
ownership of the work. A \textbf{licence} is the permission the owner grants
for specific uses. \textbf{Piracy} is the use, copying, or distribution of
software in a way the licence forbids --- for example, installing a
proprietary program without a valid licence, or copying it to machines the
licence does not cover. Open-source software is not ``unlicensed'': it is
licensed, and its licence grants broad permissions. Public-domain software,
by contrast, has no owner and therefore needs no licence at all. The exam
sometimes tests this by asking whether open-source software is ``free of
copyright''; the answer is no, because the copyright remains with the authors
and the licence grants the freedoms.

\begin{example}
A student installs the same licensed copy of a proprietary office suite on
five computers when the licence permits only one. This is \textbf{piracy}.
Installing an open-source office suite such as LibreOffice on all five
computers, in line with its licence, is not piracy. The difference is not the
price paid but whether the use matches what the licence allows.
\end{example}

\section{Common licences at a glance}
\begin{center}
\begin{tabular}{p{0.20\textwidth}p{0.30\textwidth}p{0.40\textwidth}}
\toprule
\textbf{Licence} & \textbf{Type} & \textbf{Main condition} \\
\midrule
GNU GPL & Copyleft & Derivative works must use the same licence \\
MIT License & Permissive & Reuse allowed; keep the licence notice \\
Apache License 2.0 & Permissive & Few limits, plus an explicit patent grant \\
BSD License & Permissive & Very few conditions \\
Creative Commons & Content licence & For creative works, not usually software \\
\bottomrule
\end{tabular}
\end{center}

Note that \textbf{Creative Commons} licences are designed for creative works
such as text, images, and music; they are not typically used to license
software. For software, the GPL, MIT, Apache, and BSD licences are the ones to
recognise.

\section{Well-known open-source and proprietary software}
Familiar examples help fix the concepts.

\begin{center}
\begin{tabular}{p{0.46\textwidth}p{0.46\textwidth}}
\toprule
\textbf{Open source} & \textbf{Proprietary} \\
\midrule
Linux (operating system) & Microsoft Windows \\
Android (mobile operating system) & Apple macOS, iOS \\
Mozilla Firefox (browser) & Microsoft Office \\
Chromium (Chrome's open-source base) & Adobe Photoshop \\
LibreOffice (office suite) & Oracle Database \\
VLC media player, GIMP, Blender & \\
Apache HTTP Server, MySQL, Python, Git, WordPress & \\
\bottomrule
\end{tabular}
\end{center}

\textbf{Android} is the open-source \textbf{mobile operating system}, built on
the open-source Linux kernel. This is a directly tested fact: when a question
asks for the open-source mobile OS, the answer is Android. In contrast, iOS is
a proprietary mobile operating system, and Windows and macOS are proprietary
desktop operating systems.

\section{Exam retrieval and common confusions}
Five distinctions prevent most errors on this topic.
\begin{enumerate}
  \item The test of open source is the \textbf{source code}, not the price. An
    open-source licence lets users \textbf{view, change, and share} the source
    (PYQ-26).
  \item \textbf{Freeware} is free of \textbf{cost}; \textbf{free software} is
    about \textbf{freedom}. Do not treat the two as the same (TRAP-23).
  \item \textbf{Open source}, \textbf{freeware}, and \textbf{shareware} are
    three different things: one shares the source, one is free of cost, one is
    a paid trial (TRAP-23).
  \item The open-source \textbf{mobile operating system} is \textbf{Android}
    (PYQ-24); iOS is proprietary.
  \item \textbf{GPL} is copyleft (derivatives stay open); \textbf{MIT},
    \textbf{Apache 2.0}, and \textbf{BSD} are permissive.
\end{enumerate}

\begin{example}
An item asks, ``Which of the following is an open-source mobile operating
system?'' Trace the options: Android is open source (built on the Linux
kernel), while iOS, Windows, and macOS are proprietary. The answer is Android.
If the stem instead asks what an open-source licence permits, the answer is the
right to view, change, and share the source code.
\end{example}

\subsection*{Revision at a glance}
\begin{center}
\begin{tabular}{p{0.20\textwidth}p{0.32\textwidth}p{0.38\textwidth}}
\toprule
\textbf{Term} & \textbf{What it is} & \textbf{What it is not} \\
\midrule
Open source & Source code is public and licensed for view/change/share & Not merely free of cost \\
Proprietary & Source kept secret; only the executable is distributed & Not necessarily paid \\
Freeware & Free of cost & Not necessarily open source \\
Free software & Freedom to run, study, change, share & Not about zero price \\
Public domain & No copyright owner & Not the same as open source \\
Shareware & Free trial, then pay & Not free forever \\
Commercial & Sold for profit & Not necessarily proprietary or open \\
Copyleft (GPL) & Derivatives must stay open & Not permissive \\
Permissive (MIT) & Reuse with few conditions & Not copyleft \\
\bottomrule
\end{tabular}
\end{center}

\subsection*{Exercises}
\begin{enumerate}[label=\arabic*.]
  \item Define source code and explain how it differs from object code.
  \item State the three things an open-source licence typically permits a user
    to do with the source code.
  \item Distinguish freeware from free software.
  \item Explain what public domain software is and how it differs from
    open-source software.
  \item Name one copyleft licence and one permissive licence, and state the
    main condition of each.
\end{enumerate}

\subsection*{Solutions}
\begingroup\small
\begin{enumerate}[label=\arabic*.]
  \item Source code is the human-readable set of instructions written by a
    programmer in a programming language. It is converted by a compiler or
    interpreter into object code, the machine-readable form the computer
    actually runs.
  \item View the source code, modify it, and share (redistribute) it, subject
    to the licence terms.
  \item Freeware is free of cost but its source code is usually not provided.
    Free software is about freedom: the user may study, change, and share the
    software; ``free'' means liberty, not price.
  \item Public domain software has no copyright owner and carries no
    restrictions at all. Open-source software is still copyrighted, but its
    licence grants the freedom to view, change, and share the source.
  \item Copyleft: the GNU GPL, which requires modified versions to be released
    under the same licence. Permissive: the MIT License (or Apache 2.0 / BSD),
    which allows reuse with few conditions, such as keeping the licence
    notice.
\end{enumerate}
\endgroup

\lessonheading{46}{Firmware, BIOS, UEFI and Drivers}
\section{What firmware is}
Every computer has two broad parts: the hardware, which is the physical
circuitry, and the software, which is the set of instructions. Between them
sits a special kind of software called \textbf{firmware}. Firmware is
software that is stored in non-volatile memory on the hardware itself and
that gives the device its basic, low-level instructions. It is not opened by
the user like a word processor or a web browser; it is fixed in the device
and runs by itself.

Firmware runs \emph{before} the operating system. When a computer is switched
on, its memory is empty and no program is running, so something must tell the
hardware how to begin. That ``something'' is the firmware. It starts the
hardware, checks it, and then hands control to the operating system.

\begin{definitionbox}{Firmware}
Software stored in non-volatile memory that is programmed into a hardware
device and provides its basic, low-level startup and control instructions,
running before the operating system.
\end{definitionbox}

Firmware is different from ordinary application software in three ways. It is
stored on the device rather than installed like a program; it is tied closely
to the exact hardware; and it starts the machine rather than being started by
the user.

\section{Where firmware is stored}
Firmware must be available the instant the machine is switched on, so it is
kept in \textbf{non-volatile} memory --- memory that keeps its contents
without power. If firmware were stored in RAM, it would be lost the moment
the power was switched off, and the machine could never start.

Common homes for firmware are \textbf{ROM} (Read-Only Memory), \textbf{PROM}
(Programmable ROM), \textbf{EPROM} (Erasable PROM), \textbf{EEPROM}
(Electrically Erasable PROM), and \textbf{flash memory}. Flash memory is the
modern choice because it can be erased and rewritten while still installed,
which is what makes firmware updates possible.

\begin{definitionbox}{Non-volatile memory}
Memory that retains its contents when the power is removed. ROM, flash
memory, an SSD, and an HDD are non-volatile. RAM is the common volatile
memory.
\end{definitionbox}

The important contrast is with RAM. RAM is volatile and is empty at
power-on; firmware is non-volatile and is present at power-on. A common wrong
statement is that firmware is loaded into RAM before the machine starts --- in
fact the firmware resides in ROM/flash and runs from there.

It is worth naming the members of the ROM family, because exam papers often
ask which memory type can be rewritten. \textbf{ROM} is programmed at the
factory and cannot be changed. \textbf{PROM} (Programmable ROM) can be
written once by the user. \textbf{EPROM} (Erasable PROM) can be erased with
ultraviolet light and reprogrammed. \textbf{EEPROM} (Electrically Erasable
PROM) can be erased and rewritten electrically. \textbf{Flash memory} is a
fast, block-erasable form of EEPROM and is the type used for modern firmware,
which is why firmware can be updated in place.

\begin{example}
An item states, ``The firmware of a modern PC is stored in RAM so that it can
be updated.'': This is wrong on the first point. Firmware is stored in
non-volatile ROM or flash memory, not in volatile RAM. Flash memory is what
makes the firmware updatable, not RAM.
\end{example}

\section{Hardware, software, and firmware}
It helps to place firmware among the three basic layers of a computer.

\begin{center}
\begin{tabular}{p{0.16\textwidth}p{0.36\textwidth}p{0.38\textwidth}}
\toprule
\textbf{Layer} & \textbf{What it is} & \textbf{Example} \\
\midrule
Hardware & The physical parts of the computer & Processor, motherboard, keyboard \\
Firmware & Software fixed in non-volatile memory that controls the hardware & BIOS, UEFI \\
Software & Programs that the user runs or that run on the OS & MS Word, a web browser \\
\bottomrule
\end{tabular}
\end{center}

Hardware is the body, software is the instruction set, and firmware is the
low-level bridge that tells the hardware how to begin. Firmware is software
by nature, but it is treated separately because it is stored on the device
and is not installed or removed like an ordinary program.

\section{Firmware on everyday devices}
System firmware is not unique to computers. Many devices carry firmware that
gives them their basic behaviour. A printer has firmware that controls
printing; a Wi-Fi router has firmware that runs its networking; a smartphone
has firmware plus an operating system; a modern television has firmware for
its tuner and menu. Updating such a device usually means updating its
firmware, again by flashing. This is why a firmware fix can repair a device
without any change to the user's applications.

\section{System firmware: the BIOS}
On a personal computer, the most familiar firmware is the \textbf{BIOS},
which stands for \textbf{Basic Input/Output System}. It is stored on a chip
on the motherboard. When the computer is switched on, the BIOS runs first. It
performs the power-on self-test, initialises the basic hardware, reads its
settings, and then loads the operating system.

The single most important exam fact on this topic is that the firmware which
initialises the hardware and loads the operating system is the BIOS. This is
directly tested (PYQ-28). The BIOS is sometimes confused with a device driver
or with the kernel; neither of these loads the operating system. A device
driver is loaded \emph{by} the operating system, and the kernel is part
\emph{of} the operating system.

\begin{definitionbox}{BIOS}
The Basic Input/Output System: the system firmware on a PC motherboard that
runs first at power-on, initialises the hardware, and loads the operating
system.
\end{definitionbox}

\section{The power-on boot sequence}
The machine follows a fixed chain each time it starts. Learning this chain
makes the roles of firmware, the OS, and drivers easy to separate.

\begin{enumerate}
  \item \textbf{Power on.} Electrical power reaches the motherboard.
  \item \textbf{Firmware runs.} The BIOS or UEFI begins executing from
    ROM/flash.
  \item \textbf{Hardware initialisation and POST.} The firmware tests and
    configures the basic hardware.
  \item \textbf{Operating system loads.} The firmware finds the boot device,
    reads the bootloader, and hands control to the OS.
  \item \textbf{Drivers load.} Once the OS is running, it loads the device
    drivers it needs.
  \item \textbf{Applications run.} User programs start on top of the OS.
\end{enumerate}

The chain can be written compactly as: power-on $\rightarrow$ firmware
$\rightarrow$ initialisation $\rightarrow$ OS $\rightarrow$ driver
$\rightarrow$ application. Any question that asks ``which comes first?'' is
answered by walking this chain from left to right.

The step in which the firmware hands control to the operating system is worth
a little more detail. The firmware locates the boot device according to its
configured boot order, reads a small program called the
\textbf{bootloader}, and runs it. The bootloader then loads the main part of
the operating system into memory. Under the legacy BIOS, the bootloader is
found through the MBR (Master Boot Record); under UEFI it is found through
the GPT (GUID Partition Table) scheme and a dedicated EFI system partition.

\begin{definitionbox}{Bootloader}
A small program that the firmware runs to load the operating system into
memory and start it.
\end{definitionbox}

\section{Boot devices and boot order}
The firmware does not decide by itself which device to start from. It follows
the \textbf{boot order} stored in its settings: a list such as hard disk
first, then optical drive, then USB. The firmware tries each device in turn
until it finds one that can start the operating system. The boot order is the
reason a machine can be made to start from a USB drive by moving that drive
to the top of the list, and it is one of the first things a technician checks
when a computer reports that no boot device is found.

\begin{example}
A computer is switched on and shows, ``No boot device found.'' Walking the
chain helps: the firmware runs, the POST passes, and the firmware then looks
through the boot order for a startable device. If the disk has failed or its
entry has moved down the list, no operating system is found and the firmware
stops with this message. The failing layer is the firmware's boot
configuration or the boot disk, not the application software.
\end{example}

\section{POST: the power-on self-test}
\textbf{POST} stands for \textbf{Power-On Self-Test}. It is one of the first
things the BIOS does after power-on. During POST the firmware checks that the
essential hardware --- CPU, RAM, keyboard, disk controller, and video --- is
present and responding. If a component fails, the machine may beep in a code,
show an error message, or refuse to continue.

POST is a step \emph{inside} the firmware's work. It runs \emph{before} the
operating system, not after it. A frequent trap is to place the POST after
the OS has started; the correct order is firmware (including POST) first,
then the OS.

\begin{definitionbox}{POST}
The Power-On Self-Test: the hardware check the BIOS runs at startup before
loading the operating system.
\end{definitionbox}

\section{BIOS setup and CMOS}
The BIOS keeps a small set of settings: the system date and time, the order
in which the machine should look for a boot device, hardware options, and
sometimes passwords. These settings are stored in a tiny battery-backed
memory chip called \textbf{CMOS} (Complementary
Metal-Oxide-Semiconductor). Because the chip is battery-backed, the clock
and boot order are remembered between sessions, even when the main power is
off.

The BIOS Setup utility (sometimes called the CMOS setup) is opened by
pressing a key during startup. That key is commonly \textbf{Del} or
\textbf{F2}, but the exact key varies by manufacturer. Inside setup, the
user can change the boot order, enable or disable built-in devices, and set
passwords.

\section{UEFI: the modern firmware}
\textbf{UEFI} stands for \textbf{Unified Extensible Firmware Interface}. It
is the modern replacement for the legacy BIOS and performs the same basic
job: it starts the hardware and loads the operating system. UEFI adds
several improvements.

\begin{itemize}
  \item It boots the machine \textbf{faster}.
  \item It supports very large disks through the \textbf{GPT} (GUID
    Partition Table) scheme, instead of the older MBR scheme with its
    practical limit of about 2 TB.
  \item It provides a \textbf{graphical} setup utility, often with mouse
    support.
  \item It offers \textbf{Secure Boot}, which helps prevent unauthorised
    software from loading at startup.
\end{itemize}

For exam questions, remember that ``the firmware that boots the OS'' may be
either BIOS or UEFI depending on the age of the system. Both initialise the
hardware and load the operating system; they are firmware, never drivers.

\section{BIOS versus UEFI}
The two system-firmware generations can be compared on a few clear points.

\begin{center}
\begin{tabular}{p{0.18\textwidth}p{0.34\textwidth}p{0.40\textwidth}}
\toprule
\textbf{Point} & \textbf{BIOS} & \textbf{UEFI} \\
\midrule
Full form & Basic Input/Output System & Unified Extensible Firmware Interface \\
Age & Legacy (older) & Modern \\
Boot speed & Slower & Faster \\
Disk support & MBR (about 2 TB limit) & GPT (very large disks) \\
Interface & Mostly text-based & Graphical \\
Security & No Secure Boot & Secure Boot available \\
\bottomrule
\end{tabular}
\end{center}

\section{Device drivers}
A \textbf{device driver} is software that lets the operating system
communicate with a specific hardware device. Each device --- a printer, a
graphics card, a network card, a sound card --- needs its own driver, because
a driver translates the general commands of the operating system into the
particular instructions that one device understands.

Drivers are loaded \textbf{by the operating system}, after the firmware has
started the machine. They are matched to the exact device and the exact OS
version. Because a driver is software, it is wrong to call it hardware, and
because it is loaded by the OS, it is wrong to say that a driver initialises
the hardware and loads the OS.

\begin{definitionbox}{Device driver}
Software that enables the operating system to communicate with a particular
hardware device by translating general OS commands into device-specific
instructions.
\end{definitionbox}

The three layers are easy to mix up, so keep their roles separate:

\begin{center}
\begin{tabular}{p{0.20\textwidth}p{0.30\textwidth}p{0.34\textwidth}}
\toprule
\textbf{Layer} & \textbf{Stored where} & \textbf{Role} \\
\midrule
Firmware (BIOS/UEFI) & On the hardware (ROM/flash) & Starts the machine; initialises hardware; loads the OS \\
Device driver & With the OS, on disk & Talks to one specific device \\
Operating system & On disk & Manages all resources and provides the user interface \\
\bottomrule
\end{tabular}
\end{center}

\begin{example}
A question asks, ``Which software initialises the hardware and loads the
operating system?'' A device driver is offered as an option. This is the
classic near-neighbour trap. The correct answer is the \textbf{BIOS/UEFI}
firmware. A device driver cannot do this job: it is itself loaded by the OS,
and it controls one device rather than starting the whole machine.
\end{example}

\section{Firmware, driver, OS, and kernel}
Four terms are often mixed up because they all describe ``software close to
the hardware''. Their roles are different, and the exam tests the
differences.

\begin{center}
\begin{tabular}{p{0.15\textwidth}p{0.40\textwidth}p{0.35\textwidth}}
\toprule
\textbf{Term} & \textbf{What it does} & \textbf{When it runs} \\
\midrule
Firmware (BIOS/UEFI) & Initialises hardware and loads the OS & First, at power-on \\
Kernel & The core of the OS; manages memory, processes, and devices & After the OS starts loading \\
Device driver & Communicates with one specific device & After the OS is running \\
Operating system & Manages all resources and provides the interface & After firmware, before applications \\
\bottomrule
\end{tabular}
\end{center}

The kernel is part of the operating system, not separate firmware. A driver
is loaded by the OS once the kernel is running. Only firmware runs before the
OS. If an item asks which layer ``starts the machine and loads the OS'', the
answer must be firmware --- BIOS or UEFI --- and never a driver or the
kernel.

\section{Firmware updates and flashing}
A \textbf{firmware update} replaces the firmware stored on a device with a
newer version. The process is commonly called \textbf{flashing}, because the
new firmware is written into flash memory. Updates are issued to fix bugs,
add features, and patch security flaws.

Firmware updates are much rarer than ordinary software updates, and they
carry more risk. Firmware is low-level, so a failed update can leave a device
unusable --- often described as ``bricking'' it. Interrupting the power
during flashing, or loading the wrong firmware for the model, can damage the
device permanently. For this reason the standard advice is to use only the
manufacturer's update for the exact model and to keep a stable power supply
throughout the update.

\begin{example}
A candidate reads, ``A firmware update is called formatting.'' This is wrong.
Formatting prepares a storage volume for files; it does not rewrite firmware.
The correct term for a firmware update is \textbf{flashing}.
\end{example}

\section{Exam retrieval and common confusions}
The following distinctions prevent most errors on this topic.

\begin{enumerate}
  \item Firmware that initialises hardware and loads the OS is
    \textbf{BIOS/UEFI}, not a device driver or the kernel (TRAP-24, PYQ-28).
  \item A \textbf{device driver} is \textbf{software}; the OS is system
    software (TRAP-10).
  \item Firmware is stored in \textbf{ROM/flash} and is non-volatile; RAM is
    volatile and empty at power-on.
  \item \textbf{POST} runs \textbf{before} the operating system, as part of
    the firmware's startup work.
  \item A driver is loaded \textbf{by the operating system}, after the
    firmware has started the machine.
  \item The startup order is power-on, firmware, initialisation, OS, driver,
    application.
  \item A firmware update is called \textbf{flashing}; interrupting it can
    damage the device.
\end{enumerate}

\begin{example}
An item asks which layer runs first when a computer is switched on. Trace the
chain: firmware runs before the OS, and the OS runs before any driver, which
runs before any application. The answer is the firmware (BIOS/UEFI). If the
stem instead asks what tests the hardware at startup, the answer is the POST;
if it asks what lets the OS talk to one device, the answer is the device
driver.
\end{example}

\subsection*{Exercises}
\begin{enumerate}[label=\arabic*.]
  \item Define firmware and state where it is normally stored.
  \item List the steps of the power-on boot sequence from power-on to
    application.
  \item Expand BIOS, UEFI, and POST, and give one role of each.
  \item Distinguish firmware, a device driver, and an operating system.
  \item What is a firmware update called, and what is the main risk during
    one?
\end{enumerate}

\subsection*{Solutions}
\begingroup\small
\begin{enumerate}[label=\arabic*.]
  \item Firmware is software stored in non-volatile memory that gives a
    device its basic, low-level startup and control instructions. It is
    normally stored in ROM or flash memory on the hardware itself.
  \item Power on $\rightarrow$ firmware (BIOS/UEFI) runs $\rightarrow$
    hardware initialisation and POST $\rightarrow$ the operating system
    loads $\rightarrow$ device drivers load $\rightarrow$ application
    programs run.
  \item BIOS = Basic Input/Output System: system firmware that initialises
    hardware and loads the OS. UEFI = Unified Extensible Firmware Interface:
    the modern firmware with fast boot, GPT disk support, and Secure Boot.
    POST = Power-On Self-Test: the hardware check the BIOS runs at startup.
  \item Firmware is stored on the hardware in ROM/flash and starts the
    machine and loads the OS. A device driver is software loaded by the OS
    that talks to one specific device. The operating system is the large
    system-software layer that manages all resources and provides the user
    interface.
  \item A firmware update is called flashing. The main risk is that
    interrupting the power during flashing, or using the wrong firmware, can
    leave the device unusable.
\end{enumerate}
\endgroup

\lessonheading{52}{OS PYQ Traps and System Utilities}
\section{Why this lesson is a trap drill}
The computer section of the JKSSB papers does not only test whether a
candidate knows a term. It tests whether the candidate can tell a term apart
from its \emph{nearest neighbour}. A typical wrong option is not a random
word; it is the close confusion that ``sounds right'' until the exact wording
is checked. Device driver versus hardware, DMA versus ordinary CPU copying,
system software versus application software, and the Windows Run command
\texttt{winword} are all examples of this pattern.

Because of this, the present lesson is a \textbf{trap drill}. Every question
in the practice set is built around a real near-neighbour or defective-option
pattern catalogued in the course knowledge base. The aim is not to learn a
large amount of new material, but to build the habit of reading each option
as a \emph{claim to be tested}. The rest of these notes explain the small set
of facts the traps are built on, one distinction at a time.

\begin{definitionbox}{Trap drill}
A study method in which each practice option is treated as a factual claim,
and the candidate decides whether the claim is exactly correct, closely
correct but wrong, or plainly wrong.
\end{definitionbox}

\section{The trap catalogue for this lesson}
The traps this lesson drills on are collected below. Each row states the
tempting wrong belief, the correct distinction that defeats it, and the
knowledge-base reference where the pattern is recorded. Read the table as a
checklist: if an option matches a row on the left, it is the trap, and the
right-hand column tells you what to say instead.

\begin{center}
\begin{tabular}{p{0.32\textwidth}p{0.42\textwidth}p{0.10\textwidth}}
\toprule
\textbf{Tempting wrong belief} & \textbf{Correct distinction} & \textbf{Ref.} \\
\midrule
A device driver is hardware. & A device driver is software (system software); the device is hardware. & TRAP-10 \\
DMA makes the CPU copy data ``faster''. & DMA moves data between an I/O device and main memory, bypassing the CPU. & \S 52.5 \\
The OS runs from secondary storage. & While running, the OS resides in primary memory (RAM). & PYQ-22 \\
Utility software is application software. & Utility software maintains the system; application software does the user's task. & \S 52.3 \\
MS Word is system software. & MS Word is application software. & \S 52.3 \\
MS Access is an operating system. & MS Access is a database application, not an OS. & \S 52.7 \\
Windows or macOS is a Linux distribution. & Ubuntu, Fedora, and Debian are Linux distributions. & \S 52.7 \\
\texttt{winword} opens WordPad or Notepad. & \texttt{winword} opens MS Word. & \S 52.9 \\
A component is internal or external by importance. & The classification follows position relative to the system unit. & \S 52.8 \\
An RTOS is chosen for office or graphics work. & An RTOS is chosen for fixed time deadlines. & \S 52.6 \\
\bottomrule
\end{tabular}
\end{center}

\section{The operating system: role and classification}
The \textbf{operating system (OS)} is the master \textbf{system software} of a
computer. It sits between the user (and the application programs) on one side
and the hardware on the other. Its two central jobs are to provide a
\emph{user-friendly interface} and to manage the machine's resources.

The interface may be a graphical one, such as the Windows desktop, or a
command line, such as the MS-DOS prompt or the Linux shell. Either way, the
user does not need to drive the hardware directly. In its resource-management
role, the OS schedules the processor, allocates main memory, organises
storage, and controls input and output devices from a single set of
instructions. It also provides services --- opening a file, printing a
document, sending data over a network --- that application programs call
instead of talking to the hardware themselves.

\begin{definitionbox}{Operating system}
The system software that controls the computer's hardware and provides a
user-friendly interface and a set of services to application programs.
\end{definitionbox}

Common operating systems are Windows, Linux, UNIX, MS-DOS, Android, and
macOS. An application program is not an operating system. MS Word, MS Excel,
MS Access, a web browser, and a media player are applications, not operating
systems. This is one of the exact-recognition distinctions the exam uses.

A frequently tested additional fact is \emph{where} the OS is when the
computer is running. The OS is stored permanently on secondary storage (the
disk), but while the machine is running it is loaded into \textbf{primary
memory (RAM)}. The CPU executes instructions only from main memory, so the
running OS must be resident there. When power is removed, the RAM copy is
lost and the OS is loaded again from disk at the next boot.

\section{System, application, and utility software}
Software is grouped by its purpose. \textbf{System software} runs and manages
the machine: the operating system, device drivers, utility programs, and
language translators (compiler, interpreter, assembler) all belong here.
\textbf{Application software} performs a task the user wants done: typing a
document, preparing a spreadsheet, browsing the web. \textbf{Utility
software} is a small maintenance category: tools such as Disk Cleanup, Disk
Defragmenter, and antivirus programs that keep the system healthy and
efficient. Utility software is closely tied to system software, but it is
useful to name it separately because exam options mix the three categories.

\begin{center}
\begin{tabular}{p{0.24\textwidth}p{0.34\textwidth}p{0.32\textwidth}}
\toprule
\textbf{Category} & \textbf{Role} & \textbf{Examples} \\
\midrule
System software & Runs and manages the machine & OS, device drivers, translators \\
Application software & Does a user's task & MS Word, Excel, a browser \\
Utility software & Maintains and optimises the system & Disk Cleanup, Defragmenter, antivirus \\
\bottomrule
\end{tabular}
\end{center}

\begin{example}
An item asks, ``Which of the following is application software?'' with the
options Windows, MS Word, Disk Defragmenter, and a device driver. Windows is
an operating system (system software); Disk Defragmenter is utility software;
a device driver is system software too. Only \textbf{MS Word} is application
software. The trap is that all four options are well-known pieces of
software, and only the \emph{category} separates them.
\end{example}

\section{Device drivers are software}
A \textbf{device driver} is a small software program that allows the
operating system to communicate with a specific hardware device. Each device
--- a printer, a keyboard, a graphics card, a network adapter --- needs a
matching driver. The driver receives general commands from the OS and
translates them into the device-specific instructions that the hardware
understands. Without the correct driver, the OS cannot use the device
properly.

The trap here is a classification trap. A device driver \emph{controls}
hardware, but the driver itself is \textbf{software}, and because it extends
the operating system's control of the machine it is counted as \textbf{system
software}. The device it drives is hardware. The OS itself is system
software. An option that says ``device driver is hardware'' has confused the
software with the device.

\begin{definitionbox}{Device driver}
A software program that lets the operating system control a particular
hardware device by translating OS commands into device-specific
instructions.
\end{definitionbox}

\section{DMA: transferring data without the CPU}
Normally, moving data between an input/output device and main memory keeps
the CPU busy: the CPU reads each unit from the device and writes it to memory
(or the reverse), one step at a time. \textbf{DMA} --- \textbf{Direct Memory
Access} --- replaces this with a dedicated transfer. A DMA controller moves
data \emph{directly} between the I/O device and main memory while the CPU is
bypassed for the duration of the transfer. The CPU starts the transfer and is
notified by an interrupt when it is complete.

\begin{definitionbox}{Direct Memory Access (DMA)}
A method of data transfer in which a controller moves data between an
input/output device and main memory directly, without the CPU handling each
individual byte.
\end{definitionbox}

The advantages are clear. Because the CPU is not occupied with copying each
byte, it is free to run other instructions, and large block transfers finish
faster. Several traps grow from this description. DMA moves data between a
\emph{device and main memory}, not between the CPU and the device. The CPU is
\emph{bypassed} --- it is not made ``faster''. DMA does not increase the size
of RAM and does not encrypt data. The CPU still initiates the transfer and
handles the completion interrupt; it is simply not in the per-byte path.

\section{Real-time operating systems}
A \textbf{real-time operating system (RTOS)} is an operating system designed
so that tasks complete within \textbf{fixed time deadlines}. In an ordinary
desktop OS, a task may take as long as the scheduler permits as long as it
finishes eventually. In an RTOS, meeting the deadline is the whole point: a
late response is a failure, even if the answer is correct.

An RTOS is therefore chosen for embedded and control systems, such as
industrial machinery, medical devices, and vehicle electronics, where a
delayed action can be dangerous. It is not chosen for typing a document,
designing a website, or editing a photograph; those are ordinary
application tasks that do not carry hard deadlines.

\begin{definitionbox}{Real-Time Operating System (RTOS)}
An operating system that guarantees that critical tasks complete within a
bounded, fixed time so that deadlines are always met.
\end{definitionbox}

\section{DOS, UNIX, and Linux distributions}
\textbf{DOS} stands for \textbf{Disk Operating System}. MS-DOS was a
command-line operating system: the user typed commands rather than clicking
icons. \textbf{UNIX} is a multi-user, multitasking operating system that
influenced many later systems. \textbf{Linux} is an open-source, UNIX-like
operating system; because its source can be freely used and modified, it is
distributed in many packaged forms called \textbf{distributions}. Well-known
Linux distributions include \textbf{Ubuntu}, \textbf{Fedora}, \textbf{Debian},
and \textbf{Linux Mint}.

This gives another exact-recognition drill. Windows and macOS are operating
systems, but they are \emph{not} Linux distributions. MS-DOS is an operating
system but, again, not a Linux distribution. Only Ubuntu, Fedora, Debian, and
Mint belong on the Linux-distribution list. Similarly, MS Access is a
database application and is \textbf{not} an operating system at all, even
though it is a familiar Microsoft product.

\section{Internal and external hardware}
Hardware is classified by its \emph{position} relative to the system unit
(the cabinet that houses the main machine). \textbf{Internal} hardware sits
inside the system unit: the motherboard, the CPU, RAM, the hard disk, the
power supply, and expansion cards. \textbf{External} hardware, also called a
\textbf{peripheral}, sits outside the system unit and is joined to it through
ports and cables: the monitor, keyboard, mouse, printer, and scanner.

The trap is that a familiar word can appear in the wrong column. A monitor
and a keyboard are external; a motherboard and a RAM module are internal. The
classification follows the position, not the importance of the component.

\begin{definitionbox}{Internal vs external hardware}
Internal hardware is located inside the system unit; external hardware
(a peripheral) is located outside it and connected through ports.
\end{definitionbox}

\section{The Windows Run box and the \texttt{winword} command}
The Windows \textbf{Run} box, opened with the Windows + R shortcut, starts a
program when the user types that program's command. The command
\texttt{winword} opens \textbf{MS Word}. In the same way,
\texttt{notepad} opens Notepad, \texttt{wordpad} opens WordPad, and
\texttt{cmd} opens the Command Prompt. The command must match the program
exactly.

\begin{example}
An item asks which program is opened by typing \texttt{winword} in the Run
box, offering Notepad, MS Word, WordPad, and File Explorer. The answer is
\textbf{MS Word}. Notepad is opened by \texttt{notepad}, WordPad by
\texttt{wordpad}, and File Explorer by \texttt{explorer}. The near-neighbours
are deliberately word-processor or editor names, so the candidate who only
half-remembers the command can be pulled to the wrong option.
\end{example}

\section{Running example: the office desktop}
It helps to hang these facts on one concrete machine. Picture the office
desktop used throughout this course: a multi-core PC running Windows, with
MS Word and MS Excel installed, a printer, and an internet connection.

When the machine is switched on, the Windows OS is loaded from the disk into
primary memory and stays resident there while the machine runs. MS Word and
MS Excel are \emph{application} software that the user opens to do work.
Disk Cleanup and Disk Defragmenter are \emph{utility} software that the user
runs occasionally to keep the machine healthy. The printer does not work
until its \emph{device driver} --- a piece of system software --- is
installed, so that Windows can translate its print commands into the
printer's own instructions. When a large print job or file copy is under way,
a \emph{DMA} transfer can move the data between the device and main memory
while the CPU continues with other work.

\begin{example}
A single scenario can therefore generate several trap options about one
machine. ``Word is system software'' is false --- Word is application
software. ``The printer driver is hardware'' is false --- the driver is
software. ``The OS is stored on the disk and run from there'' is half-true
and therefore defective: it is stored on the disk but \emph{run} from main
memory. Recognising which category each item belongs to settles all three.
\end{example}

\section{Defective options and how to read them}
Some official options are not merely wrong; they are \emph{defective} --- the
wording itself is inconsistent or ambiguous. The knowledge base records this
as a real feature of the papers (Anti-Pattern AP-02). A defective option is
teaching material, not a fact to memorise. The safe method is to quote the
option exactly, mark the defect, and then apply the stable concept.

\begin{example}
A defective option might read, ``The device driver is the hardware that
stores the instructions the operating system uses to control a device.'' The
defect is the phrase ``the hardware that stores the instructions''. A driver
is \emph{software} that holds those instructions; it is not hardware. The
stable concept to apply is TRAP-10: a device driver is software. Once the
defect is spotted, the option can be rejected without any doubt.
\end{example}

The same discipline applies to near-neighbour options that are not defective
but simply close. When two options are both plausible, compare them word by
word against the definitions in this lesson. The distinction is usually a
single term --- \emph{software} versus \emph{hardware}, \emph{device to
memory} versus \emph{CPU to device}, \emph{application} versus
\emph{utility}, or \texttt{winword} versus \texttt{wordpad}. Reading for that
one term is what turns a guess into a correct answer under negative marking.

\section{Terminology ladder for this lesson}
Direct-recall expansion items are the single most common JKSSB format, so the
exact terms from this lesson are collected here with their full forms. Each
row is one fact.

\begin{center}
\begin{tabular}{p{0.13\textwidth}p{0.30\textwidth}p{0.47\textwidth}}
\toprule
\textbf{Term} & \textbf{Full form} & \textbf{One-line meaning} \\
\midrule
OS & Operating System & System software that provides an interface and manages resources \\
DMA & Direct Memory Access & Data moves between an I/O device and main memory with the CPU bypassed \\
RTOS & Real-Time Operating System & An OS that meets fixed time deadlines \\
DOS & Disk Operating System & A command-line operating system (MS-DOS) \\
Driver & Device driver & Software that lets the OS control a specific device \\
Utility & Utility software & Maintenance software such as Disk Cleanup or Defragmenter \\
Peripheral & External hardware & A device outside the system unit, joined by ports \\
\bottomrule
\end{tabular}
\end{center}

\section{Exam retrieval and common confusions}
Six distinctions prevent most errors on this topic.
\begin{enumerate}
  \item A device driver is \textbf{software} (system software); the device it
    controls is hardware (TRAP-10).
  \item The OS is \textbf{system software}. MS Word is application software;
    Disk Cleanup is utility software.
  \item While running, the OS is held in \textbf{primary memory (RAM)}, not
    executed from secondary storage (PYQ-22).
  \item DMA moves data between a \textbf{device and main memory} and
    \textbf{bypasses the CPU}; it is not ``faster CPU copying''.
  \item \textbf{RTOS} = Real-Time Operating System, chosen for fixed
    deadlines.
  \item \textbf{DOS} = Disk Operating System; Linux distributions include
    Ubuntu, Fedora, and Debian. \texttt{winword} opens MS Word.
\end{enumerate}

\subsection*{Exercises}
\begin{enumerate}[label=\arabic*.]
  \item Classify each as system, application, or utility software: Windows,
    MS Word, Disk Defragmenter, a device driver, Disk Cleanup.
  \item Explain why a device driver is described as software even though it
    controls hardware.
  \item Describe a DMA transfer: name the two locations between which data
    moves and state the role of the CPU.
  \item State the full form of RTOS and give one reason an RTOS is used.
  \item Which of these are operating systems, and which are Linux
    distributions: Ubuntu, Windows, Debian, MS Access, Fedora?
  \item What does typing \texttt{winword} in the Windows Run box do, and how
    does it differ from \texttt{wordpad}?
\end{enumerate}

\subsection*{Solutions}
\begingroup\small
\begin{enumerate}[label=\arabic*.]
  \item System software: Windows and the device driver. Application
    software: MS Word. Utility software: Disk Defragmenter and Disk Cleanup.
  \item The driver is a stored set of instructions that the OS loads and
    runs to talk to the device. It has no physical body of its own; the
    hardware is the device. A program that extends OS control of hardware is
    therefore software, and specifically system software.
  \item Data moves directly between the input/output device and main memory.
    A DMA controller performs the transfer; the CPU starts it and is
    interrupted when it finishes, but it does not handle each byte.
  \item RTOS = Real-Time Operating System. It is used where tasks must
    complete within fixed time deadlines, as in embedded and control
    systems.
  \item Operating systems: Ubuntu (also a Linux distribution), Windows,
    Debian (also a Linux distribution), Fedora (also a Linux distribution).
    MS Access is not an operating system; it is a database application.
  \item Typing \texttt{winword} opens MS Word. Typing \texttt{wordpad} opens
    WordPad, a different text editor. The two commands open different
    programs and must not be confused.
\end{enumerate}
\endgroup

\lessonheading{60}{Current OS, Processes, Linux and Utilities}
\section{Deadlock and its four necessary conditions}
A \key{deadlock} is a situation in which a group of processes are each waiting
for a resource that is held by another process in the group, so that none of
them can ever proceed. A simple picture is two processes: process A holds a
printer and waits for a scanner, while process B holds the scanner and waits
for the printer. Each waits for the other, and both stay stuck forever.

For a deadlock to be possible, four conditions must hold \emph{at the same
time}. They are usually listed as follows.

\begin{center}
\begin{tabular}{p{0.28\textwidth}p{0.62\textwidth}}
\toprule
\textbf{Condition} & \textbf{Meaning} \\
\midrule
Mutual exclusion & A resource is held by only one process at a time. \\
Hold and wait & A process holds at least one resource while waiting for more. \\
No preemption & A resource cannot be forcibly taken away from its holder. \\
Circular wait & A cycle of processes exists, each waiting for a resource held by the next. \\
\bottomrule
\end{tabular}
\end{center}

\begin{definitionbox}{Deadlock}
A state in which a set of processes are each blocked waiting for a resource
held by another process in the same set, so no process can make progress.
\end{definitionbox}

\begin{definitionbox}{No preemption (as a condition)}
The rule that a resource already allocated to a process cannot be taken away
by force; it must be released voluntarily. Its opposite, \key{preemption}, is
used to \emph{recover} from a deadlock.
\end{definitionbox}

\section{Deadlock recovery by resource preemption}
Once a deadlock has been detected, the operating system must break it. One
recovery method is \key{resource preemption}: the system takes a resource away
from one of the processes and gives it to another process that needs it,
breaking the circular wait. The chosen process is usually rolled back to a
safe earlier state, or restarted. Preemption is also used to break the
\emph{no-preemption} condition directly, since it removes the guarantee that a
resource stays with its holder.

\begin{example}[The exam contrast]
An item asks for a method of \emph{recovering} from deadlock. The correct
answer is \good{resource preemption} --- the system forcibly reclaims a
resource. A distractor may name ``mutual exclusion'' or ``circular wait'',
but those are \bad{necessary conditions} for a deadlock, not recovery
methods. The other trap is to think preemption is one of the four conditions;
the condition is actually \key{no preemption}. (PYQ-10)
\end{example}

\section{Processes and the scheduler}
A \key{process} is a program in execution. When a program is launched, the
operating system creates a process for it and keeps the information needed to
run it --- its state, its program counter, and the resources it holds. Because
many processes are usually ready at once but only one can run on a given CPU
core at a time, the operating system includes a \key{scheduler} that decides
which ready process gets the CPU next.

Scheduling may be \key{non-preemptive}, where a running process keeps the CPU
until it finishes or blocks, or \key{preemptive}, where the operating system
can interrupt a running process and switch to another. Priorities are one of
the most common ways to choose between processes, and it is the combination of
\emph{priorities} and \emph{preemption} that produces the problems studied in
the next sections: priority inversion and starvation.

\begin{definitionbox}{Process}
A program in execution, together with the state and resources the operating
system maintains for it.
\end{definitionbox}

\begin{definitionbox}{Scheduler}
The operating system component that selects which ready process is given the
CPU next.
\end{definitionbox}

\section{Priority inversion versus priority inheritance}
In a priority-based system, a task with a higher priority is supposed to run
before one with a lower priority. \key{Priority inversion} breaks that rule.
It happens when a \key{high-priority} task is blocked because a
\key{low-priority} task currently holds a resource the high-priority task
needs. While the low-priority task holds that resource, a
\key{medium-priority} task can preempt it and run. The result is that the
high-priority task ends up waiting behind the medium-priority task, even
though it should have run first --- its priority has effectively been
inverted.

The standard fix is \key{priority inheritance}. While the low-priority task
holds the resource, it temporarily \key{inherits} the high priority of the
task that is waiting for it, so it runs at that high priority, releases the
resource quickly, and then returns to its original low priority. This lets the
high-priority task obtain the resource without waiting behind the
medium-priority task.

\begin{center}
\begin{tabular}{p{0.24\textwidth}p{0.32\textwidth}p{0.32\textwidth}}
\toprule
\textbf{Point} & \textbf{Priority inversion} & \textbf{Priority inheritance} \\
\midrule
What it is & A problem in which priority order is reversed. & A solution that temporarily raises the holder's priority. \\
Cause & A low-priority task holds a needed resource. & Applied while a resource is held and awaited. \\
Effect & High-priority task waits behind a medium-priority task. & Low-priority holder finishes fast, so the high-priority task proceeds. \\
\bottomrule
\end{tabular}
\end{center}

\begin{definitionbox}{Priority inversion}
A condition in which a high-priority task is indirectly preempted by a
lower-priority task, typically because a low-priority task holds a shared
resource.
\end{definitionbox}

\begin{example}[What does \emph{not} fix inversion]
An item asks which method prevents priority inversion. \key{Round-robin}
scheduling gives every ready process a turn but knows nothing about resource
holding, so by itself it does \bad{not} prevent inversion. The correct method
is \good{priority inheritance} (or the related priority-ceiling protocol).
(PYQ-11)
\end{example}

\section{Starvation and aging}
\key{Starvation} (indefinite blocking) is different from deadlock. A starved
process is \emph{ready} to run and could run in principle, but it is
\key{repeatedly passed over} by the scheduler, so it may wait for an
unbounded time. A common case is a low-priority process in a strict priority
system: as long as higher-priority processes keep arriving, the low-priority
process never gets the CPU.

The standard remedy for starvation is \key{aging}: the system gradually
\key{raises the priority} of a process the longer it waits, so that eventually
it becomes the highest-priority ready process and runs.

\begin{definitionbox}{Starvation}
A scheduling condition in which a ready process is indefinitely denied the
resources it needs because other processes are repeatedly preferred.
\end{definitionbox}

\begin{center}
\begin{tabular}{p{0.22\textwidth}p{0.34\textwidth}p{0.34\textwidth}}
\toprule
\textbf{Point} & \textbf{Starvation} & \textbf{Deadlock} \\
\midrule
Can the process ever run? & Yes, in principle & No, not without outside action \\
Cause & Scheduling choices (long wait) & A circular resource dependency \\
Cure & Aging (raise priority with waiting time) & Detection + recovery (for example preemption) \\
\bottomrule
\end{tabular}
\end{center}

\section{Preemptive priority scheduling}
In \key{priority scheduling}, the scheduler always picks the ready process with
the \key{highest priority}. The scheme is \key{preemptive} when a process that
is currently running can be \key{interrupted} as soon as a process with a
higher priority becomes ready; the higher-priority process then runs. It is
\key{non-preemptive} when the running process keeps the CPU until it either
finishes or blocks for input/output. Preemptive priority scheduling gives
quick service to important tasks, but with strict priorities it can starve
low-priority tasks unless aging is added.

\section{Real-time operating systems (RTOS)}
An \key{RTOS} is a \key{Real-Time Operating System}. It is designed so that
the system responds within a \key{bounded time}, allowing \key{deadlines} to be
met. Correctness in a real-time system depends not only on the \emph{result}
but also on \emph{when} the result is produced.

\begin{definitionbox}{Real-time operating system (RTOS)}
An operating system designed to process data and events within guaranteed
time limits, so that deadline requirements are satisfied.
\end{definitionbox}

Real-time systems are classified by what happens when a deadline is missed. In
a \key{hard real-time} system a missed deadline is a \bad{failure} and can have
serious or catastrophic consequences. In a \key{soft real-time} system a
missed deadline only \good{degrades} the quality of the result; the result is
still useful. The trap is to think an RTOS is simply ``a very fast operating
system''. Speed helps, but the defining property is \key{meeting deadlines}.

\section{The Linux kernel}
The \key{kernel} is the \key{core of the operating system}. It is the part
that runs with full control of the hardware and provides the basic services
every other program relies on. In Linux, the kernel manages
\key{processes} (through the scheduler), \key{memory}, \key{devices} (through
device drivers), and \key{file systems}.

The kernel runs in a privileged mode called \key{kernel space}, while ordinary
application programs run in \key{user space}. A user program asks the kernel
for a service --- reading a file, for example --- through a \key{system call};
the kernel performs the operation and returns the result. The Linux kernel is
\key{open-source}: its source code can be viewed, modified, and shared.

\begin{definitionbox}{Kernel}
The central component of an operating system that manages the hardware
resources (processes, memory, devices, and file systems) and mediates between
applications and the hardware.
\end{definitionbox}

The \key{shell} is frequently confused with the kernel. The shell is the
\emph{command interpreter} that a user types commands into; it is not the
kernel. The shell passes the user's requests to the kernel, which does the
actual work.

\section{\texttt{/root} versus \texttt{/home}}
Linux stores the personal files of each user in that user's \key{home
directory}. The \key{superuser}, whose account name is \key{root}, has its
home directory at \key{\texttt{/root}}. Ordinary users' home directories are
normally placed under \key{\texttt{/home}}, for example
\texttt{/home/alice}.

It is important not to confuse two different things with similar names:

\begin{center}
\begin{tabular}{p{0.18\textwidth}p{0.70\textwidth}}
\toprule
\textbf{Path} & \textbf{Meaning} \\
\midrule
\texttt{/} & The \key{root of the whole file system} --- the single top directory from which all other directories branch. \\
\texttt{/root} & The \key{home directory of the root (superuser)} account. \\
\texttt{/home} & The parent directory holding the home directories of \key{ordinary users}. \\
\bottomrule
\end{tabular}
\end{center}

The trap is to read ``\texttt{/root}'' as ``the root of the file system''. The
top of the file system is \texttt{/} (a single slash). \texttt{/root} is only
the superuser's own home folder.

\section{File systems}
A \key{file system} is the method an operating system uses to organise, store,
and retrieve files and folders on a storage device. Different systems use
different file systems.

\begin{center}
\begin{tabular}{p{0.16\textwidth}p{0.22\textwidth}p{0.50\textwidth}}
\toprule
\textbf{File system} & \textbf{Full form} & \textbf{Typical use} \\
\midrule
NTFS & New Technology File System & Default for modern Windows installations. \\
FAT32 & File Allocation Table, 32-bit & Older system; widely compatible with small drives and memory cards. \\
ext4 & Fourth Extended file system & Native and common default file system for Linux. \\
\bottomrule
\end{tabular}
\end{center}

The exam cue is the pairing: \key{NTFS = Windows}, \key{ext4 = Linux}, and
\key{FAT32 = broad compatibility}. Full forms are commonly asked directly.

\section{Windows utilities: Store, WordPad/Notepad, Disk Cleanup}
The \key{Windows Store} (now the \key{Microsoft Store}) is the built-in
\key{online marketplace} in Windows. Users \key{search, download, install, and
update} apps from it, so applications come from a trusted catalogue rather
than from arbitrary websites.

Windows includes two simple text editors that are easy to confuse.

\begin{center}
\begin{tabular}{p{0.14\textwidth}p{0.36\textwidth}p{0.36\textwidth}}
\toprule
\textbf{Point} & \textbf{Notepad} & \textbf{WordPad} \\
\midrule
Text handling & \key{Plain text} only; no formatting & \key{Basic formatting}: bold, italic, fonts, colour \\
Typical save & \texttt{.txt} & \texttt{.rtf} (rich text format) \\
Capability & The smallest, fastest text editor & Can insert pictures and objects \\
\bottomrule
\end{tabular}
\end{center}

\key{Notepad} is a plain-text editor: it displays and saves characters without
any formatting. \key{WordPad} supports \emph{basic} rich-text formatting and
can include pictures, but it is still lighter than a full word processor such
as MS Word.

\key{Disk Cleanup} is a built-in Windows utility that \key{frees disk space}.
It deletes files that are not needed, such as temporary files, thumbnail
caches, downloaded program files, and the contents of the \key{Recycle Bin}.
It is the tool to reach for when a drive is running short of free space.

\section{Defragmenter and the SSD caution}
On a traditional \key{hard disk drive (HDD)}, a file can become
\key{fragmented}: its pieces are scattered in different places on the rotating
disk, which makes reading slower. The \key{Disk Defragmenter} rearranges the
pieces so that each file is stored \key{contiguously} (in one run), which
improves access speed.

\textbf{SSD caution.} A \key{solid-state drive (SSD)} has \key{no moving
parts}. Defragmenting it gives no speed benefit and only adds
\key{unnecessary writes}, which reduce the drive's working life. SSDs should
\bad{not} be defragmented; Windows uses the \key{TRIM} command for SSDs
instead.

\begin{definitionbox}{Defragmentation}
The process of rearranging the fragments of files on a hard disk so that each
file occupies contiguous space, improving read performance.
\end{definitionbox}

\section{Imaging Fax and the \texttt{xcopy /s} command}
\key{Windows Fax and Scan} is the built-in \key{imaging and fax} utility in
Windows. With a scanner it can \key{scan} a document and save it as an image;
with a fax modem or fax service it can \key{send and receive faxes}. This is
the tool the exam means by ``Imaging Fax''.

\key{xcopy} is a Windows \key{command-line} command that copies files and
whole directory trees. Its switch \key{\texttt{/s}} copies the named directory
\emph{including its subdirectories}, but it does not copy empty ones. Adding
\key{\texttt{/e}} copies empty subdirectories as well. The exact form to
remember is \key{\texttt{xcopy /s}}: ``copy with subdirectories''.

\begin{center}
\begin{tabular}{p{0.20\textwidth}p{0.68\textwidth}}
\toprule
\textbf{Command / tool} & \textbf{What it does} \\
\midrule
\texttt{xcopy /s} & Copies files and subdirectories (empty ones excluded). \\
\texttt{xcopy /s /e} & Also copies empty subdirectories. \\
Windows Fax and Scan & Scans documents and sends/receives faxes (Imaging Fax). \\
\bottomrule
\end{tabular}
\end{center}

\section{Exam retrieval and common confusions}
Six distinctions prevent most errors on this topic.
\begin{enumerate}
  \item Deadlock's four conditions are mutual exclusion, hold and wait,
    \key{no preemption}, and circular wait. \key{Preemption} is a
    \key{recovery} method, not a condition (PYQ-10).
  \item \key{Priority inversion} is the problem; \key{priority inheritance} is
    the cure. Round-robin by itself does \bad{not} prevent inversion (PYQ-11).
  \item \key{Starvation} is indefinite \emph{scheduling} delay, not deadlock;
    \key{aging} fixes starvation.
  \item An \key{RTOS} is defined by \key{meeting deadlines} (hard vs soft), not
    merely by being fast.
  \item The Linux \key{kernel} is the core; the \key{shell} is the command
    interpreter. \texttt{/root} is the superuser's home; \texttt{/} is the top
    of the file system; \texttt{/home} holds ordinary users' homes.
  \item \key{Notepad} = plain text; \key{WordPad} = basic rich text. Disk
    Cleanup frees space; \key{never defragment an SSD} (use TRIM);
    \texttt{xcopy /s} copies subdirectories.
\end{enumerate}

\begin{example}[Tracing a mixed item]
An item lists statements and asks which are true:
(I) Preemption is one of the four necessary conditions for deadlock.
(II) Priority inheritance reduces priority inversion.
(III) Aging helps prevent starvation.
Statement I is \bad{false}: preemption is a recovery method; the condition is
\emph{no} preemption. Statements II and III are \good{true}. The answer is
``II and III only''.
\end{example}

\subsection*{Exercises}
\begin{enumerate}[label=\arabic*.]
  \item List the four necessary conditions for deadlock, and name one method
    of recovering from a deadlock.
  \item Explain priority inversion and state which technique resolves it.
  \item Distinguish starvation from deadlock, and name the standard remedy for
    starvation.
  \item What is an RTOS? Give the difference between a hard and a soft
    real-time system.
  \item State the meaning of \texttt{/root} and \texttt{/home} in Linux, and
    name the default file system of Windows and of Linux.
  \item Explain why an SSD should not be defragmented, and state what
    \texttt{xcopy /s} does.
\end{enumerate}

\subsection*{Solutions}
\begingroup\small
\begin{enumerate}[label=\arabic*.]
  \item Mutual exclusion, hold and wait, no preemption, and circular wait. A
    deadlock can be recovered from by resource preemption (taking a resource
    away from one process and giving it to another).
  \item Priority inversion is when a high-priority task waits because a
    low-priority task holds a needed resource, while a medium-priority task
    runs. \key{Priority inheritance} resolves it: the low-priority holder
    temporarily inherits the high priority and releases the resource quickly.
  \item A starved process is ready but repeatedly passed over, so it may wait
    indefinitely; a deadlocked process can never proceed because of a circular
    resource dependency. \key{Aging} --- gradually raising the waiting
    process's priority --- prevents starvation.
  \item An RTOS (Real-Time Operating System) guarantees a response within a
    bounded time so deadlines are met. In a \key{hard} real-time system a
    missed deadline is a failure; in a \key{soft} real-time system a missed
    deadline only degrades quality.
  \item \texttt{/root} is the home directory of the superuser (root);
    \texttt{/home} holds the home directories of ordinary users. Windows uses
    \key{NTFS}; Linux uses \key{ext4}.
  \item An SSD has no moving parts, so defragmenting gives no speed benefit and
    only adds wear; SSDs use \key{TRIM} instead. \texttt{xcopy /s} copies a
    directory including its subdirectories, excluding empty ones.
\end{enumerate}
\endgroup

\unitopener{6}{Programming \& Data Structures}
\lessonheading{51}{Programming Foundations, Compiler and Data Structures}
\section{Programming languages and their levels}
A computer can act only on instructions given in its own \textbf{machine
language}, a pattern of binary 0s and 1s. Writing whole programs in that form
is slow, error-prone, and tied to one processor. Programmers therefore write
in languages that are easier for humans, and a translator converts the
resulting program into machine language. Languages are placed on a
\textbf{level} scale: low-level languages (machine and assembly) are close to
the hardware, while high-level languages are closer to human language. The
lower the level, the more closely the code reflects the actual machine
instructions.

\begin{definitionbox}{Level of a programming language}
How close a language is to the underlying machine. Machine and assembly
languages are low-level; C, C++, Java, and Python are high-level.
\end{definitionbox}

\begin{center}
\begin{tabular}{p{0.24\textwidth}p{0.20\textwidth}p{0.46\textwidth}}
\toprule
\textbf{Language} & \textbf{Generation} & \textbf{Main feature} \\
\midrule
Machine language & First & Binary 0s and 1s; only language the CPU runs directly \\
Assembly language & Second & Short mnemonics such as ADD and MOV; machine-dependent \\
High-level language & Third & English-like and mathematical; machine-independent \\
\bottomrule
\end{tabular}
\end{center}

\section{Machine language}
\textbf{Machine language} is the only language a processor executes directly.
It is written entirely in binary, and it is \emph{machine-dependent}: each
processor family has its own instruction set, so a program written for one
machine does not run on another. Machine language needs no translation, so it
is fast to run, but it is very difficult to write, read, and debug by hand. It
is also called the \emph{first-generation language}.

\section{Assembly language}
\textbf{Assembly language} replaces binary instruction codes with short
\emph{mnemonics} such as ADD, SUB, and MOV. It is still
\emph{machine-dependent}, because each processor family has its own assembly
language, and it maps almost one-to-one onto machine instructions. Assembly
code cannot run directly; it must be translated into machine language by an
\textbf{assembler}. Assembly language is called the \emph{second-generation
language} and is a low-level language.

A common exam trap is to call assembly language a high-level language. It is
not: assembly language is a low-level, second-generation language because its
statements correspond closely to machine instructions. High-level languages,
by contrast, are machine-independent and use English-like statements.

\section{High-level languages}
\textbf{High-level languages} use English-like words and mathematical
notation. Examples are C, C++, Java, Python, COBOL, and FORTRAN. They are
\emph{machine-independent}: after translation, the same source program can run
on different machines. They are easier to learn, write, read, and maintain
than low-level languages, but they must be translated by a \textbf{compiler}
or an \textbf{interpreter} before the CPU can run them.

\section{Generations of programming languages}
Languages are also grouped into generations, which run parallel to the level
scale above. The first three generations match the three levels exactly.

\begin{center}
\begin{tabular}{p{0.22\textwidth}p{0.26\textwidth}p{0.38\textwidth}}
\toprule
\textbf{Generation} & \textbf{Example} & \textbf{Main feature} \\
\midrule
First (1GL) & Machine language & Binary 0s and 1s; run directly \\
Second (2GL) & Assembly language & Mnemonics; needs an assembler \\
Third (3GL) & C, C++, Java & High-level; needs a compiler or interpreter \\
Fourth (4GL) & SQL & Non-procedural; states what, not how \\
Fifth (5GL) & Logic / AI languages & Problem-solving and constraints \\
\bottomrule
\end{tabular}
\end{center}

For this exam, the high-yield rows are the first three: machine language is
first generation, assembly language is second generation, and a high-level
language such as C or Java is third generation.

\section{Language translators: assembler, compiler, interpreter}
A \textbf{translator} is a program that converts source code written in one
language into another form, usually machine language. There are three common
translators, each matched to a different input language and method.

\begin{definitionbox}{Assembler}
A translator that converts an \emph{assembly-language} program into machine
code.
\end{definitionbox}

\begin{definitionbox}{Compiler}
A translator that converts a \emph{high-level} program into machine code (object
code) as a whole, translating the entire program before it runs.
\end{definitionbox}

\begin{definitionbox}{Interpreter}
A translator that converts and executes a high-level program
\emph{line by line}, without producing a separate object file.
\end{definitionbox}

\begin{center}
\begin{tabular}{p{0.20\textwidth}p{0.36\textwidth}p{0.36\textwidth}}
\toprule
\textbf{Point} & \textbf{Compiler} & \textbf{Interpreter} \\
\midrule
Input & Whole program at once & One statement at a time \\
Output & Object code file & No separate object file \\
Execution speed & Faster & Slower \\
Errors & Reported after full translation & Reported line by line \\
Debugging & Harder & Easier \\
Typical examples & C, C++ & Python, JavaScript \\
\bottomrule
\end{tabular}
\end{center}

The compiler translates the whole program in one pass and produces object
code, so execution is fast, but a single error can delay the whole build and
debugging is harder. The interpreter takes one statement at a time, translates
it, and executes it immediately; execution is slower, but errors are shown
line by line, which makes debugging easier.

\section{Source code, object code, and machine code}
These three terms describe a program at different stages of translation.

\begin{definitionbox}{Source code}
The program as written by the programmer in a high-level or assembly language;
it is human-readable text.
\end{definitionbox}

\begin{definitionbox}{Object code}
The machine-readable, relocatable code produced by a compiler or assembler
from source code; it is not yet linked into a final runnable program.
\end{definitionbox}

\begin{definitionbox}{Machine code}
The final binary instructions, in the processor's own language, that the CPU
executes directly.
\end{definitionbox}

The path is: the programmer writes \textbf{source code}; the translator
converts it into \textbf{object code}; after linking, the result is
\textbf{machine code} that the processor runs. Source code is human-readable;
object code and machine code are machine-readable.

A common exam trap is to ask which of these is written by the programmer. The
answer is \textbf{source code}. Object code and machine code are produced by
the translator, not written by hand, and they are stored as binary patterns
rather than readable text.

\subsection*{How a program is translated}
The translation of a high-level program follows a fixed sequence:
\begin{enumerate}[label=\arabic*.]
  \item The programmer writes the program as \textbf{source code}.
  \item A \textbf{compiler} translates the whole source file and produces an
    \textbf{object file}.
  \item A \textbf{linker} joins the object file with any library code needed.
  \item The result is a \textbf{machine-code} executable that the CPU runs.
\end{enumerate}
An interpreter shortens this path: it takes each statement of the source
program, translates it, and executes it immediately, so no separate object
file is produced.

\section{Algorithm}
An \textbf{algorithm} is a finite, step-by-step procedure for solving a
problem. It has a clear start and end, and every step is unambiguous so that
anyone following it reaches the same result. An algorithm is
language-independent: it can be written in plain English or in pseudocode,
without committing to a particular programming language. A program is the
algorithm expressed in a programming language.

\begin{definitionbox}{Algorithm}
A finite, ordered set of unambiguous steps that solves a problem, with a
defined start and end.
\end{definitionbox}

\begin{example}
Adding two numbers is a small algorithm: (1) start, (2) read values $a$ and
$b$, (3) compute $sum = a + b$, (4) display $sum$, (5) stop. The same steps
can later be written in any programming language.
\end{example}

\section{Flowchart}
A \textbf{flowchart} is a pictorial, diagrammatic representation of an
algorithm. It uses standard symbols joined by arrows to show the order in
which steps are carried out, which makes the flow of control easy to see at a
glance.

\begin{center}
\begin{tabular}{p{0.34\textwidth}p{0.50\textwidth}}
\toprule
\textbf{Symbol} & \textbf{Meaning} \\
\midrule
Oval (terminator) & Start and end of the process \\
Parallelogram & Input or output \\
Rectangle & Process (a computation or action) \\
Diamond (rhombus) & Decision (yes/no branch) \\
Arrow & Flow line showing direction \\
\bottomrule
\end{tabular}
\end{center}

A flowchart for the add-two-numbers algorithm uses these symbols in order: an
oval for start, a parallelogram to read the two inputs, a rectangle to compute
the sum, a parallelogram to display the result, and an oval for stop. When a
step involves a choice, a diamond is used, and two arrows leave it, one for
each outcome.

\section{Debugging and types of errors}
\textbf{Debugging} is the process of finding and correcting errors, called
\emph{bugs}, in a program. Three kinds of error are commonly distinguished.

\begin{definitionbox}{Syntax error}
A violation of the rules of the language, such as a missing semicolon or a
misspelled keyword. It is reported by the compiler.
\end{definitionbox}

\begin{definitionbox}{Logical error}
An error in the logic of the program: it compiles and runs but produces a
wrong result. It is the hardest kind to find.
\end{definitionbox}

\begin{definitionbox}{Runtime error}
An error that occurs while the program is running, such as dividing by zero
or reading a file that does not exist.
\end{definitionbox}

A compiler catches syntax errors but not logical errors. A program with a
logical error compiles and runs, and only the output reveals that something
is wrong.

\section{Object-oriented programming basics}
\textbf{OOP} stands for \textbf{Object-Oriented Programming}. It organises a
program around objects rather than around a sequence of actions. Two central
terms are the class and the object.

\begin{definitionbox}{Class and object}
A \emph{class} is a blueprint or template that defines the data and methods;
an \emph{object} is a single instance created from that class.
\end{definitionbox}

Four features define the object-oriented approach.

\begin{itemize}
  \item \textbf{Encapsulation}: binding the data and the methods that operate
    on it into one unit, with the data hidden from outside access.
  \item \textbf{Inheritance}: a new class acquires the properties and methods
    of an existing class.
  \item \textbf{Polymorphism}: ``one name, many forms'' --- the same operation
    behaves differently for different objects.
  \item \textbf{Abstraction}: hiding the implementation details and showing
    only the essential features.
\end{itemize}

As an illustration, \emph{Student} can be a class that defines the data
(roll number, name, marks) and the methods (enrol, show result). A particular
student such as one with roll number 12 is then an \emph{object} of that class.
The class describes the shape; each object is one concrete example of it.

\section{Stack and LIFO}
A \textbf{stack} is a linear data structure in which insertion and deletion
take place at only one end, called the \emph{top}. It follows the order
\textbf{LIFO}: Last In, First Out. The last element pushed onto the stack is
the first element popped from it.

\begin{definitionbox}{Stack (LIFO)}
A linear data structure that works on the Last In, First Out rule, with
insertion and deletion performed at the top.
\end{definitionbox}

Insertion into a stack is called \textbf{PUSH} and deletion is called
\textbf{POP}. The other end of the stack is the \emph{bottom}, and no
insertion or deletion happens there.

\begin{example}
A stack of files on a desk behaves in exactly this way. The file placed on top
last is the first one picked up next. That is the LIFO rule in everyday form.
\end{example}

\section{Basic data structures}
A data structure is a way of organising data so that it can be used
efficiently. The common structures tested at this level are the array, stack,
queue, linked list, and tree.

\begin{center}
\begin{tabular}{p{0.20\textwidth}p{0.64\textwidth}}
\toprule
\textbf{Structure} & \textbf{Defining property} \\
\midrule
Array & Fixed-size elements of the same type, stored in contiguous memory and accessed by an index \\
Stack & Last In, First Out (LIFO); push and pop take place at the top \\
Queue & First In, First Out (FIFO); insertion at the rear, deletion at the front \\
Linked list & Nodes that hold data and a pointer to the next node; not contiguous and can grow dynamically \\
Tree & A hierarchical structure with a root node and child nodes \\
\bottomrule
\end{tabular}
\end{center}

To identify a structure in an exam item, look for its defining rule: same type,
fixed size, and index access means an \textbf{array}; Last In, First Out means
a \textbf{stack}; First In, First Out (like people in a queue) means a
\textbf{queue}; data-plus-pointer nodes that grow dynamically mean a
\textbf{linked list}; and a root with branches means a \textbf{tree}.

The two most commonly confused structures are the stack and the queue. They
differ only in the order in which elements leave the structure.

\begin{center}
\begin{tabular}{p{0.20\textwidth}p{0.34\textwidth}p{0.34\textwidth}}
\toprule
\textbf{Point} & \textbf{Stack} & \textbf{Queue} \\
\midrule
Order & LIFO (Last In, First Out) & FIFO (First In, First Out) \\
Insertion & Push, at the top & Enqueue, at the rear \\
Deletion & Pop, at the top & Dequeue, at the front \\
Both ends used? & No, one end only & Yes, front and rear \\
\bottomrule
\end{tabular}
\end{center}

\section{Exam retrieval and common confusions}
The following distinctions prevent most errors on this topic.
\begin{enumerate}
  \item A \textbf{compiler} translates the \textbf{whole program at once}; an
    \textbf{interpreter} works \textbf{line by line}.
  \item An \textbf{assembler} translates \textbf{assembly language}; a
    \textbf{compiler} and an \textbf{interpreter} translate
    \textbf{high-level} languages.
  \item \textbf{Source code} is human-readable; \textbf{object code} and
    \textbf{machine code} are machine-readable.
  \item An \textbf{algorithm} is a step-by-step procedure; a
    \textbf{flowchart} is its pictorial representation.
  \item A \textbf{stack} is \textbf{LIFO}; a \textbf{queue} is
    \textbf{FIFO}. Do not swap them.
  \item A \textbf{logical error} is not caught by the compiler; the program
    runs but the result is wrong.
  \item \textbf{OOP} stands for Object-Oriented Programming; a class is a
    blueprint and an object is an instance of it.
\end{enumerate}

\begin{example}
An item asks which translator produces a separate object file before
execution. Trace the definitions: the compiler translates the whole
high-level program and produces object code, so the answer is the compiler. If
the stem instead says ``translates and executes line by line'', the answer is
the interpreter, and if it says ``translates assembly language'', the answer
is the assembler.
\end{example}

\subsection*{Exercises}
\begin{enumerate}[label=\arabic*.]
  \item Distinguish machine language, assembly language, and a high-level
    language by level and by dependence on the machine.
  \item State the difference between an assembler, a compiler, and an
    interpreter.
  \item Differentiate source code, object code, and machine code.
  \item Define an algorithm and explain how a flowchart represents one.
  \item Distinguish a syntax error, a logical error, and a runtime error.
  \item Explain LIFO and name the two operations of a stack.
  \item Identify the data structure: (a) elements of the same type in
    contiguous memory, (b) Last In First Out, (c) a root with child nodes.
\end{enumerate}

\subsection*{Solutions}
\begingroup\small
\begin{enumerate}[label=\arabic*.]
  \item Machine language is the first-generation, binary, machine-dependent
    language the CPU runs directly. Assembly language is the second
    generation, uses mnemonics, and is also machine-dependent. A high-level
    language is third generation, uses English-like statements, and is
    machine-independent.
  \item An assembler translates assembly language into machine code. A
    compiler translates a high-level program as a whole and produces object
    code. An interpreter translates and executes a high-level program line by
    line.
  \item Source code is the program written by the programmer in a high-level
    or assembly language and is human-readable. Object code is the
    machine-readable code produced by the translator and is not yet linked.
    Machine code is the final binary code the CPU executes directly.
  \item An algorithm is a finite, step-by-step, unambiguous procedure to solve
    a problem. A flowchart represents it pictorially using standard symbols
    (oval, parallelogram, rectangle, diamond, arrow) joined by flow lines.
  \item A syntax error breaks the rules of the language and is reported by
    the compiler. A logical error lets the program run but gives a wrong
    result. A runtime error occurs while the program is running.
  \item LIFO means Last In, First Out: the last element pushed is the first
    element popped. The two operations are push (insertion) and pop
    (deletion), both performed at the top.
  \item (a) Array, (b) Stack, (c) Tree.
\end{enumerate}
\endgroup

\unitopener{7}{Office Applications \& File Formats}
\lessonheading{22}{MS Word Fundamentals}
\section{What MS Word is}
\textbf{MS Word} (Microsoft Word) is a \textbf{word processor}: a program
used to create, edit, format, store, and print text documents. It is one of
the main applications in the \textbf{Microsoft Office} suite. A document
written in Word is saved as a file, and the modern default file format uses
the \textbf{.docx} extension. Word shows the document in a rectangular
window that contains the document area together with several command areas
around it.

\begin{definitionbox}{Word processor}
An application that lets a user type text, change its appearance
(formatting), arrange it on pages, and print it. MS Word is the common
example.
\end{definitionbox}

A word processor separates the \emph{text} from its \emph{appearance}. The
same words can be displayed in many different fonts, sizes, colours, and
layouts, because the formatting is stored alongside the text and can be
changed at any time. This separation is the central idea of the whole
lesson.

\begin{center}
\begin{tabular}{p{0.28\textwidth}p{0.60\textwidth}}
\toprule
\textbf{Part of the window} & \textbf{What it is} \\
\midrule
Document area & The white page where the text is typed and shown. \\
Ribbon & The strip of command tabs across the top of the window. \\
Quick Access Toolbar & The small row of always-visible buttons at the top-left. \\
Status bar & The strip at the bottom showing page, word count, and view buttons. \\
\bottomrule
\end{tabular}
\end{center}

\section{The Ribbon and its tabs}
The \textbf{Ribbon} is the broad band of commands across the top of the Word
window. It replaced the older style of drop-down menus and toolbars. The
Ribbon is organised in three levels:

\begin{enumerate}
  \item \textbf{Tabs} across the top: \textbf{Home, Insert, Design, Layout,
    References, Mailings, Review,} and \textbf{View}.
  \item \textbf{Groups} inside each tab. For example, the Home tab contains
    the \textbf{Font}, \textbf{Paragraph}, and \textbf{Styles} groups.
  \item \textbf{Commands} inside each group, such as the Bold button or the
    alignment buttons.
\end{enumerate}

Some tabs are \textbf{contextual}. They are not shown all the time; they
appear only when the right object is selected. When a table is selected, the
\textbf{Table Tools} tabs appear; when a picture is selected, the
\textbf{Picture Format} tab appears. The \textbf{File} tab is special: it
does not show groups like the other tabs but opens \textbf{Backstage view},
which holds document-level actions such as New, Open, Save, Save As, Print,
and Close.

\begin{definitionbox}{Ribbon}
The command area of the Word window, organised into tabs, groups, and
buttons, which replaced the older menu-and-toolbar layout.
\end{definitionbox}

For exam purposes, the key point is that the \textbf{Home} tab carries the
most-used formatting commands: the \emph{Font} tools that change how
characters look, and the \emph{Paragraph} tools that change alignment,
indentation, spacing, bullets, and numbering.

\section{The Quick Access Toolbar}
The \textbf{Quick Access Toolbar (QAT)} is the small row of buttons in the
top-left corner of the window, just above the Ribbon. It is designed to hold
the few commands a user needs most often, so they stay visible on every tab.
By default it shows \textbf{Save}, \textbf{Undo}, and \textbf{Redo}. A user
can \textbf{customise} it by adding or removing commands, and it can be moved
to sit below the Ribbon if preferred.

\begin{definitionbox}{Quick Access Toolbar}
A small, customisable toolbar at the top-left of the Word window that keeps
frequently used commands (by default Save, Undo, Redo) visible on every tab.
\end{definitionbox}

\section{Creating, opening and saving documents}
The \textbf{File} tab opens \textbf{Backstage view}, the part of Word where
you work with the document as a whole rather than its contents. From here a
user can create a new document, open an existing one, save, print, or close.
The common keyboard shortcuts are worth memorising because direct-recall
items ask for them.

\begin{center}
\begin{tabular}{p{0.24\textwidth}p{0.64\textwidth}}
\toprule
\textbf{Shortcut} & \textbf{Action} \\
\midrule
Ctrl+N & Create a new blank document. \\
Ctrl+O & Open an existing document. \\
Ctrl+S & Save the current document. \\
Ctrl+P & Open the Print screen. \\
\bottomrule
\end{tabular}
\end{center}

\textbf{Save} writes the current document back to its file. \textbf{Save As}
writes a \emph{copy} under a new name, in a new location, or in a different
file format, leaving the original as it was. The default Word format is
\textbf{.docx}; documents made with much older versions use the \textbf{.doc}
extension.

\begin{example}
A student types a letter and presses Ctrl+S, saving it as
\texttt{letter.docx}. To keep the original and also produce a version for
printing, the student uses Save As and saves a second file with a new name.
The first file is unchanged.
\end{example}

\section{Entering and editing text}
Text is typed at the \textbf{insertion point}, which is shown by a blinking
cursor. \textbf{Word wrap} moves text automatically to the next line when
the current line is full; pressing \textbf{Enter} starts a new
\textbf{paragraph}. The \textbf{Backspace} key deletes the character to the
left of the cursor, while the \textbf{Delete} key deletes the character to
the right. \textbf{Ctrl+Z} undoes the last action and \textbf{Ctrl+Y} redoes
it.

Text must be selected before it can be reformatted. A \textbf{word} is
selected by double-clicking it, a \textbf{paragraph} by triple-clicking, and
the \textbf{whole document} with \textbf{Ctrl+A}.

\section{Font formatting}
The \textbf{Font} group on the Home tab changes how individual characters
look. It sets the \textbf{font face} (the typeface, such as Arial or Times
New Roman) and the \textbf{font size}. It provides the style buttons
\textbf{Bold (Ctrl+B)}, \textbf{Italic (Ctrl+I)}, and \textbf{Underline
(Ctrl+U)}, and it sets the \textbf{font colour} and \textbf{highlight}. It
also applies effects such as subscript and superscript. Deeper options are
available in the Font dialog box.

\begin{definitionbox}{Font}
The typeface and its settings (face, size, style, colour) that determine how
characters are displayed. Font formatting changes the appearance of
characters, not their arrangement on the page.
\end{definitionbox}

\section{Paragraph alignment}
\textbf{Alignment} decides where each line of a paragraph sits between the
left and right margins.

\begin{center}
\begin{tabular}{p{0.22\textwidth}p{0.20\textwidth}p{0.46\textwidth}}
\toprule
\textbf{Alignment} & \textbf{Shortcut} & \textbf{Effect} \\
\midrule
Left & Ctrl+L & Text lines up evenly on the left; ragged on the right. This is the default. \\
Centre & Ctrl+E & Each line is placed in the middle of the page. \\
Right & Ctrl+R & Text lines up evenly on the right. \\
Justify & Ctrl+J & Text is flush with both left and right margins. \\
\bottomrule
\end{tabular}
\end{center}

\textbf{Justify} stretches the spaces between words so that both edges of
the paragraph are even, which is why books and newspapers use it.
\textbf{Centre} does not stretch anything; it simply moves each line to the
middle.

\section{Indentation}
\textbf{Indentation} pushes paragraph lines inward, away from the margin.
The \textbf{left indent} and \textbf{right indent} move the whole paragraph
inward from the respective side. A \textbf{first-line indent} pushes only
the first line inward, as in a typical book paragraph. A \textbf{hanging
indent} pushes every line \emph{except} the first inward, which is the
standard layout for a references or bibliography list.

\begin{definitionbox}{Indentation}
The inward offset of paragraph lines from the page margins. First-line
indent affects only the first line; hanging indent affects every line except
the first.
\end{definitionbox}

Indentation must not be confused with \textbf{spacing}. Indentation shifts
lines \emph{horizontally}; spacing changes the \emph{vertical} gaps.

\section{Line and paragraph spacing}
\textbf{Line spacing} sets the vertical gap between lines \emph{inside} a
paragraph; common settings are single, 1.5, and double. \textbf{Paragraph
spacing} sets the gap \emph{before} and \emph{after} a paragraph. Both are
controlled from the Paragraph group on the Home tab, or from the Paragraph
dialog box. Using spacing is better than pressing Enter repeatedly to make
gaps, because spacing stays correct when text is edited.

\section{Bullets and numbering}
\textbf{Bullets} mark list items with a symbol and are used when the order of
the items does not matter. \textbf{Numbering} marks items with numbers or
letters and is used when the order or sequence matters. \textbf{Multilevel
lists} combine levels (for example 1, then 1.1, then 1.1.1) to build nested
outlines. All three are applied from the \textbf{Paragraph} group on the
Home tab.

\begin{example}
A shopping list uses \textbf{bullets} because the order is irrelevant. A set
of instructions to install software uses \textbf{numbering} because the
steps must be followed in order.
\end{example}

\section{Tables}
A \textbf{table} arranges information in \textbf{rows} and \textbf{columns};
each intersection of a row and a column is a \textbf{cell}. A table is
inserted from \textbf{Insert $\rightarrow$ Table}. Cells can be
\textbf{merged} into a single larger cell or \textbf{split} back into
separate cells, and rows and columns can be inserted, deleted, and resized.
When the table is selected, the \textbf{Table Tools} contextual tabs appear.

\begin{definitionbox}{Table}
A grid of rows and columns used to arrange text in cells, inserted in Word
from Insert $\rightarrow$ Table.
\end{definitionbox}

\section{Images and objects}
An image is placed in the document from \textbf{Insert $\rightarrow$
Pictures}. \textbf{Handles} around the image let the user resize and rotate
it. \textbf{Text wrapping} decides how the surrounding text flows around the
image; options include In Line with Text, Square, Tight, and Behind Text.
When the image is selected, the \textbf{Picture Format} contextual tab
appears. Charts, shapes, and \textbf{SmartArt} are inserted from the same
Insert tab.

\begin{definitionbox}{Text wrapping}
The setting that controls how text flows around a picture or object. ``In
Line with Text'' treats the image like a character; the other styles let
text flow around the image.
\end{definitionbox}

\section{Headers and footers}
A \textbf{header} is the area at the \emph{top} of a page and a
\textbf{footer} is the area at the \emph{bottom}. Both repeat on every page
of a section, which makes them ideal for a document title (in the header)
and \textbf{page numbers} (in the footer). They are inserted from
\textbf{Insert $\rightarrow$ Header} and \textbf{Insert $\rightarrow$
Footer}.

\section{Page setup}
Page setup is found on the \textbf{Layout} tab. It controls the overall page
geometry: \textbf{margins} (the blank space between the text and the page
edges), \textbf{orientation} (\textbf{Portrait}, which is tall, or
\textbf{Landscape}, which is wide), and \textbf{page size} (for example A4
or Letter). The same tab sets \textbf{columns} and the page \textbf{breaks}
described next.

\section{Page break vs section break}
A \textbf{Page Break} (\textbf{Ctrl+Enter}) forces the following content to
start on a \emph{new page}, but the page formatting stays the same. A
\textbf{Section Break} starts a new \textbf{section} of the document, which
can carry its own page setup. The most important consequence for the exam is
that only a \textbf{Section Break} allows a part of the document to have
\textbf{different headers and footers}. Both are inserted from
\textbf{Layout $\rightarrow$ Breaks}.

\begin{example}
An item claims, ``A Page Break can be used to give the first page a
different header from the rest of the document.'' This is wrong. A Page
Break only moves content to a new page; to change the header or footer part
way through the document, a \textbf{Section Break} is required. The wrong
option has confused the two breaks (TRAP-13).
\end{example}

\section{Printing a document}
\textbf{Ctrl+P} opens the \textbf{Print} screen in Backstage view. The
\textbf{print preview} shows exactly how the pages will appear on paper. On
this screen the user selects the \textbf{printer}, the number of
\textbf{copies}, and which \textbf{pages} to print, then clicks \textbf{Print}
to send the job to the printer. Checking the preview first catches wrong
margins, stray blank pages, and cut-off text before paper is used.

\section{Exam retrieval and common confusions}
Five distinctions prevent most errors on this topic.
\begin{enumerate}
  \item \textbf{Save} overwrites the current file; \textbf{Save As} stores a
    copy under a new name, location, or format.
  \item \textbf{Ctrl+S} saves; \textbf{Ctrl+P} prints; \textbf{Ctrl+Enter}
    inserts a page break. Do not swap the shortcuts.
  \item \textbf{Bullets} are for unordered lists; \textbf{numbering} is for
    ordered lists.
  \item \textbf{Indentation} shifts lines horizontally; \textbf{spacing}
    changes vertical gaps between lines and paragraphs.
  \item A \textbf{Page Break} does \emph{not} change headers and footers;
    only a \textbf{Section Break} does (TRAP-13).
\end{enumerate}

\begin{example}
An item asks which shortcut opens the Print screen. Trace the File-level
actions: New is Ctrl+N, Open is Ctrl+O, Save is Ctrl+S, and Print is
Ctrl+P. The answer is \textbf{Ctrl+P}. If the stem instead asks for Undo, the
answer is Ctrl+Z, and if it asks for a page break, the answer is Ctrl+Enter.
\end{example}

\subsection*{Exercises}
\begin{enumerate}[label=\arabic*.]
  \item Name the main parts of the MS Word window.
  \item What is the Quick Access Toolbar, and what does it show by default?
  \item Distinguish Save from Save As.
  \item List the four paragraph alignments with their shortcuts.
  \item Explain the difference between a Page Break and a Section Break.
\end{enumerate}

\subsection*{Solutions}
\begingroup\small
\begin{enumerate}[label=\arabic*.]
  \item The document area, the Ribbon (with its tabs and groups), the Quick
    Access Toolbar, and the status bar.
  \item The Quick Access Toolbar is the small customisable row of buttons at
    the top-left of the window. By default it shows Save, Undo, and Redo,
    and it stays visible on every tab.
  \item Save writes the current document back to its existing file. Save As
    writes a copy under a new name, location, or format and leaves the
    original file unchanged.
  \item Left (Ctrl+L), Centre (Ctrl+E), Right (Ctrl+R), and Justify
    (Ctrl+J).
  \item A Page Break (Ctrl+Enter) starts the following content on a new page
    while keeping the same formatting. A Section Break starts a new section
    that can have its own page setup --- including different headers and
    footers. Only a Section Break can change headers and footers part way
    through a document.
\end{enumerate}
\endgroup

\lessonheading{23}{MS Word High-Yield}
\section{Why these Word tools are high-yield}
MS Word is a large program, but JKSSB does not test it evenly. A small set of
review and editing tools is asked again and again, usually as an exact shortcut
or as a near-neighbour distinction. This lesson covers that set: Find and
Replace, spell check, hyperlinks and bookmarks, styles, comments, Track Changes,
Page Break versus Section Break, mail merge, and saving, printing and exporting
to PDF. The aim is not to learn every menu; it is to know each tool's shortcut
and the one thing it is \emph{not}.

\begin{definitionbox}{The high-yield set}
The Word features that competitive-exam items concentrate on: Find/Replace,
spell check and grammar tools, hyperlinks and bookmarks, styles, comments,
Track Changes, Section Break, mail merge, and PDF export.
\end{definitionbox}

\begin{center}
\begin{tabular}{p{0.28\textwidth}p{0.62\textwidth}}
\toprule
\textbf{Term} & \textbf{Meaning in this lesson} \\
\midrule
Find/Replace & The tool to locate text and swap it for other text. \\
Spell check & Word's proofing tool for spelling and grammar. \\
Bookmark & A named location or selection inside a document. \\
Style & A named set of formatting applied to text or a paragraph. \\
Track Changes & A mode that records every edit with its author. \\
Section Break & A break that starts a new section with its own layout. \\
Mail merge & Combining a standard letter with a list of records. \\
PDF & Portable Document Format --- a fixed-layout file. \\
\bottomrule
\end{tabular}
\end{center}

\section{Finding and replacing text}
Word has two related tools. \textbf{Find} locates text in the document and is
opened with \textbf{Ctrl+F}. \textbf{Find and Replace} both locates text and
substitutes new text for it, and is opened with \textbf{Ctrl+H}. Inside the
Find and Replace dialog, \emph{Replace} changes the current match while
\emph{Replace All} changes every match in one step. The same dialog can find
formatting, special characters (such as a paragraph mark) and text that uses a
particular style.

The common confusion is between the two shortcuts. Ctrl+F only finds; it does
not change anything. Ctrl+H is the one that replaces. An item that asks how to
``find and replace'' text expects Ctrl+H.

\begin{example}
A student has typed the word ``recieve'' in ten places and wants ``receive''
everywhere. Press Ctrl+H, type \texttt{recieve} in the Find box and
\texttt{receive} in the Replace box, then choose Replace All. Using Ctrl+F
alone would only highlight the mistakes one by one, without correcting them.
\end{example}

\section{Spelling and grammar tools}
Press \textbf{F7} to start the \textbf{spelling and grammar} check. Word
flags a possible misspelling with a red squiggle and a possible grammar problem
with another underline. Modern Word gathers these checks in the
\textbf{Editor} pane, which groups Spelling, Grammar and Refinements. From the
check you can \emph{Ignore} a flagged word once, \emph{Ignore All} every
occurrence, or \emph{Add} the word to the dictionary so it is accepted from
then on.

Two automatic tools support proofing. \textbf{AutoCorrect} fixes common typing
errors as you type (for example, capitalising the first letter of a sentence),
and the \textbf{thesaurus} (Shift+F7) offers synonyms. The exam usually tests
only the main shortcut, F7. Do not confuse F7 with the PowerPoint slide-show
keys: F5 starts a slide show and F9 has other roles in PowerPoint, but in Word
F7 is proofing.

\begin{definitionbox}{Spell check}
The Word tool, opened with F7, that checks spelling and grammar and offers
corrections, ignoring or adding words to the dictionary.
\end{definitionbox}

\section{Hyperlinks and bookmarks}
A \textbf{hyperlink} is text or an object that jumps to another place when
clicked. Press \textbf{Ctrl+K} to insert one. A hyperlink can point to a web
address, to a file, or to a location inside the same document. When it points
inside the document, it usually points to a \textbf{bookmark}.

A \textbf{bookmark} gives a location or a selection a name, so that text,
a hyperlink or a cross-reference can jump to it. Bookmarks are not printed and
are not visible to the reader by default; they are an internal target. The
clean distinction for the exam is that a hyperlink is the \emph{clickable link},
while a bookmark is the \emph{named target} it points to.

\section{Styles}
A \textbf{style} is a named collection of formatting settings: font, size,
colour, spacing, alignment and indent. Instead of formatting each paragraph by
hand, you apply a style. Applying a \textbf{Heading} style does more than
change the look: it gives the paragraph an outline level. That is why the
\textbf{Navigation Pane} can list the document's headings and why a
\textbf{Table of Contents} can be built and refreshed automatically.

Style-based formatting is consistent and easy to update. If you edit a style,
every paragraph that uses that style changes at once. This is the main
advantage of styles over direct formatting.

\begin{example}
A report uses the Heading 1 style for its chapter titles. Later the author
wants all chapter titles in a larger font. Instead of selecting every title,
the author edits the Heading 1 style once, and every chapter title updates
automatically. The Navigation Pane and the table of contents also stay correct.
\end{example}

\section{Comments and Track Changes}
A \textbf{comment} is a note attached to a point in the document. It lets a
reviewer ask a question or suggest a change \emph{without editing the text
itself}. Comments can be replied to and marked resolved, and they are stored as
metadata shown in the margin or the Reviewing Pane.

\textbf{Track Changes} is different: it records the actual edits. When it is
on, each insertion, deletion and formatting change is marked, together with the
\textbf{author's name} and the time. This is called revision metadata. The
reviewer then \textbf{Accepts} or \textbf{Rejects} each change individually or
all at once. Accepting a change keeps it permanently and, with it, removes the
revision metadata for that change. Rejecting a change removes the edit and
keeps the original text.

\begin{definitionbox}{Track Changes}
A Word mode that records every insertion, deletion and formatting edit along
with the author and time, so a reviewer can accept or reject each change.
\end{definitionbox}

\section{Page Break versus Section Break}
A \textbf{Page Break} pushes the following text onto a new page. It does not
start a new section, so the page keeps the same layout, margins, headers,
footers and page numbering as the rest of the document. A quick way to insert a
page break is \textbf{Ctrl+Enter}.

A \textbf{Section Break} divides the document into independent
\textbf{sections}. Because each section has its own layout settings, a Section
Break is what allows \emph{different headers and footers}, different page
numbering, different margins or a different page orientation within one
document. Section Breaks are inserted from the \textbf{Layout} tab, under
\textbf{Breaks}.

This is one of the most tested Word traps. A Page Break changes only the page; a
Section Break changes the section. Only a Section Break can give a document two
different headers or footers. (TRAP-13)

\begin{example}
An item states, ``Using Page Break allows you to have different headers and
footers in the same document.'' This is wrong. A Page Break keeps the same
section formatting. It is a \textbf{Section Break} that allows different
headers and footers in one document. (PYQ-13)
\end{example}

\section{Mail merge}
\textbf{Mail merge} is the process that combines a \textbf{main document} (a
standard letter or template) with a \textbf{data source} (a list of records
such as names and addresses). The main document contains \textbf{merge fields}
as placeholders, and each record fills those fields to produce one personalised
copy. A common data source is a table of names and addresses held in Excel or
in Outlook contacts.

The typical steps are: create or select the main document; select recipients
(the data source); insert merge fields; preview the results; and finish the
merge by printing the letters or sending them by email. The exam definition to
remember is that mail merge \emph{combines letters with a database of
addresses}. (PYQ-37, PYQ-47)

\begin{definitionbox}{Mail merge}
The Word process that combines a standard document with a data source of
records, producing one personalised copy of the document for each record.
\end{definitionbox}

\section{Saving, printing and exporting to PDF}
Press \textbf{Ctrl+S} to save the document; Word normally saves it in the
\textbf{DOCX} format. Press \textbf{Ctrl+P} to open the print options.
\textbf{Save As} can change the file name, the location and the file format.
Using \textbf{Export} or \textbf{Save As PDF} creates a \textbf{PDF} file.

\textbf{PDF} stands for \textbf{Portable Document Format}. A PDF is a
\emph{fixed-layout} file: it looks the same on every device and is good for
sharing and printing, but it is not meant for easy editing. This pairs with the
separate lesson on file formats, where the extension (for example,
\texttt{.pdf}) and the format (PDF) are kept distinct.

\section{Exam retrieval and common confusions}
Five distinctions prevent most errors on this topic.
\begin{enumerate}
  \item \textbf{Ctrl+F} finds; \textbf{Ctrl+H} finds and replaces. Do not swap
    them.
  \item \textbf{F7} is spelling and grammar in Word; F5 starts a PowerPoint
    slide show. (TRAP-32)
  \item A \textbf{Section Break}, not a Page Break, enables different
    headers/footers in one document. (TRAP-13)
  \item A \textbf{bookmark} is a named target; a \textbf{hyperlink}
    (\textbf{Ctrl+K}) is the clickable link.
  \item \textbf{Mail Merge} combines a letter with an address database; it is
    not Track Changes or AutoCorrect. (TRAP-29)
\end{enumerate}

\begin{example}
An item asks which feature combines letters with a database of addresses.
Trace the meaning: a standard letter plus a list of records is mail merge. If
the item instead asks what records every edit with its author, the answer is
Track Changes, and if it asks what fixes typing errors as you type, the answer
is AutoCorrect. (PYQ-37, PYQ-47)
\end{example}

\subsection*{Exercises}
\begin{enumerate}[label=\arabic*.]
  \item State the shortcuts for Find, Find and Replace, spelling and grammar,
    and inserting a hyperlink.
  \item Explain the difference between a Page Break and a Section Break, and
    say which one allows different headers and footers.
  \item Distinguish a bookmark from a hyperlink.
  \item What does Track Changes record, and what happens when a change is
    accepted?
  \item Define mail merge and name its two parts.
\end{enumerate}

\subsection*{Solutions}
\begingroup\small
\begin{enumerate}[label=\arabic*.]
  \item Find: Ctrl+F. Find and Replace: Ctrl+H. Spelling and grammar: F7.
    Insert a hyperlink: Ctrl+K.
  \item A Page Break moves text to a new page within the same section. A
    Section Break starts a new section with its own layout. Only a Section
    Break allows different headers and footers.
  \item A bookmark is a named location or selection inside the document; it is
    a target and is not printed. A hyperlink is clickable text or an object
    that jumps to a target, which may be a bookmark, a file or a web address.
  \item Track Changes records every insertion, deletion and formatting edit,
    together with the author's name and the time (revision metadata). Accepting
    a change keeps it permanently and removes that change's revision metadata.
  \item Mail merge combines a main document (the standard letter or template)
    with a data source (a list of records such as names and addresses) to
    produce one personalised copy per record.
\end{enumerate}
\endgroup

\lessonheading{24}{Excel Structure}
\section{What Excel is and what a workbook is}
Microsoft Excel is a spreadsheet program. A spreadsheet is a grid of boxes
in which you can store text, numbers, dates, and formulas, and in which the
program can calculate results automatically. Excel is one of the standard
office applications, and its saved document is called a \textbf{workbook}.

A \textbf{workbook} is the Excel \emph{file} --- the single item that you
open, edit, and save with a file name. When you start Excel and see a grid,
you are looking at one \textbf{worksheet} that belongs to a workbook. A
workbook can contain many worksheets, and each worksheet is a separate page
of the same file.

\begin{definitionbox}{Workbook}
The Excel file itself. A single workbook can hold one or more worksheets and
is saved under one file name.
\end{definitionbox}

\begin{definitionbox}{Worksheet}
A single sheet of rows and columns inside a workbook. It is also called a
sheet, and it has a tab at the bottom of the Excel window.
\end{definitionbox}

At the bottom of the Excel window you will see one or more \textbf{sheet
tabs} (for example, Sheet1, Sheet2, Sheet3). A new classic workbook opens
with three worksheets, and you can add or remove sheets. Clicking a tab
brings that worksheet to the front. The key point for the exam is the
relationship: the workbook is the container (the file), and the worksheet is
the content (one page inside it).

\section{Workbook vs worksheet}
These two words are a classic near-neighbour pair, so keep the contrast
sharp. The workbook is the file; the worksheet is a part of the file. One
workbook can contain many worksheets, but a worksheet never contains a
workbook. A worksheet is not a separate saved file --- it lives inside the
workbook and is saved as part of it.

\begin{center}
\begin{tabular}{p{0.20\textwidth}p{0.34\textwidth}p{0.34\textwidth}}
\toprule
\textbf{Point} & \textbf{Workbook} & \textbf{Worksheet} \\
\midrule
What it is & The Excel file & One sheet inside a workbook \\
How many & One file & Many per workbook \\
Saved as & A single file name & Part of that file \\
Tab & --- & A tab at the bottom \\
\bottomrule
\end{tabular}
\end{center}

\begin{example}
An item states, ``A worksheet is the Excel file, and a workbook is a part of
it.'' This is wrong. The relation is the other way round: the workbook is
the file, and the worksheet is one page inside it. The wrong option has
swapped the two terms.
\end{example}

\section{The worksheet grid: rows and columns}
Every worksheet is a large grid made of \textbf{rows} and \textbf{columns}.
Rows run horizontally across the sheet and are identified by
\textbf{numbers}, which appear down the left edge (1, 2, 3, \dots).
Columns run vertically down the sheet and are identified by
\textbf{letters}, which appear across the top (A, B, C, \dots).

The column labels do not stop at Z. After A--Z, the labels continue as
two-letter names: AA, AB, \dots, AZ, then BA, BB, \dots, BZ, and so on. This
scheme continues with three-letter names until the final column. In Excel
2007 and later, a worksheet has:

\begin{center}
\begin{tabular}{ll}
\toprule
\textbf{Property} & \textbf{Value in Excel 2007 and later} \\
\midrule
Total rows & 1,048,576 (numbered 1 to 1,048,576) \\
Total columns & 16,384 (labelled A to XFD) \\
First column & A \\
Last column & XFD \\
\bottomrule
\end{tabular}
\end{center}

These two figures are worth memorising because they are common direct-recall
items. In older Excel versions (Excel 97--2003, the \texttt{.xls} format) the
limits were much smaller: 65,536 rows and 256 columns, with the last column
labelled IV. Modern Excel uses the larger grid ending at column XFD.

\section{The cell and the active cell}
A \textbf{cell} is the small rectangular box formed where a row and a column
meet. It is the basic unit of a worksheet, and it is where you actually type
your data. The whole sheet is just a very large collection of cells.

At any moment, exactly one cell is the \textbf{active cell}. The active cell
is the cell that is currently selected: it has a thick border around it, and
anything you type is entered into it. The active cell's address is shown in
the Name Box, and its contents are shown in the Formula Bar. Even when you
select a group of cells, only one of them --- the active cell --- is the one
the next typed entry goes into.

\begin{definitionbox}{Cell}
The rectangular box formed by the intersection of a row and a column, in
which data is entered.
\end{definitionbox}

\begin{definitionbox}{Active cell}
The cell that is currently selected. It has a thick border, receives the
next typed entry, and has its address shown in the Name Box.
\end{definitionbox}

\section{Cell addresses and references}
Every cell has a unique \textbf{cell address} (also called a cell
reference). The address is written with the \textbf{column letter first},
followed by the \textbf{row number}. For example:

\begin{itemize}
  \item \textbf{B5} means column B, row 5.
  \item \textbf{C7} means column C, row 7.
  \item \textbf{D12} means column D, row 12.
\end{itemize}

The order is fixed: letter then number. Writing 12D, B:5, or B\_5 is not a
valid Excel address. Because column labels can have more than one letter, an
address such as AB100 is also valid: it means column AB, row 100.

\section{Ranges}
A \textbf{range} is a group of two or more cells treated as one block. It is
written by typing the addresses of two diagonally opposite corners and
separating them with a \textbf{colon} (:). For example, \textbf{A1:C10}
means all the cells from column A to column C and from row 1 to row 10. A
range is normally rectangular, and the cells in it are contiguous.

The number of cells in a range is the number of columns multiplied by the
number of rows. For example, the range B2:D5 covers three columns (B, C, D)
and four rows (2, 3, 4, 5), so it contains $3 \times 4 = 12$ cells.

\begin{definitionbox}{Range}
A rectangular group of two or more selected cells, written as two corner
addresses joined by a colon, such as A1:C10.
\end{definitionbox}

\begin{example}
A question lists \texttt{B5} and calls it a range. This is wrong: \texttt{B5}
is a \emph{single} cell, because it has only one address and no colon. A
range must have two corner addresses joined by a colon, such as
\texttt{B5:D9}. The colon is the signal that a range, not a single cell, is
meant.
\end{example}

\section{The formula bar}
The \textbf{Formula Bar} is the long white bar that sits below the Ribbon
and above the worksheet grid. It always shows the \emph{contents} of the
active cell. If the active cell holds plain text or a number, the Formula
Bar shows that text or number. If the active cell holds a formula, the
Formula Bar shows the \emph{formula itself} (for example
\texttt{=A1+A2}), while the cell itself shows the calculated result.

You can also edit the active cell in the Formula Bar by clicking in it. This
is useful when a formula is long. Remember the division of labour: the cell
displays the result, and the Formula Bar displays the underlying formula or
data.

\begin{definitionbox}{Formula Bar}
The bar below the Ribbon that shows, and allows editing of, the contents of
the active cell --- including the formula behind a calculated result.
\end{definitionbox}

\section{The name box}
The \textbf{Name Box} is the small box at the \emph{left end of the Formula
Bar}. It normally shows the \textbf{address of the active cell} (for
example, B5). If a range of cells is selected, it shows the name or the
address range of that selection (for example, A1:C10). You can also type an
address into the Name Box and press Enter to jump directly to that cell.

\begin{definitionbox}{Name Box}
The box at the left of the Formula Bar that shows the address of the active
cell or the name of the selected range.
\end{definitionbox}

Keep the two bars apart. The Name Box (on the left) answers ``which cell
am I in?'' The Formula Bar (the long part) answers ``what is in that
cell?''

\section{Basic data types}
Excel accepts several basic kinds of entry in a cell:

\begin{itemize}
  \item \textbf{Text} (also called a label): letters, words, names, and
    headings. Text is not used in arithmetic.
  \item \textbf{Number}: a numeric value that can be used in calculations.
  \item \textbf{Date and time}: entered and displayed as a date or time, but
    stored internally by Excel as a number.
  \item \textbf{Formula}: an entry that begins with the \textbf{=} sign and
    calculates a result from other cells.
\end{itemize}

Excel also recognises \textbf{logical values} (\texttt{TRUE} and
\texttt{FALSE}) and displays \textbf{error values} such as
\texttt{\#VALUE!} when a calculation cannot be completed. By default, Excel
aligns these data types in a helpful way:

\begin{center}
\begin{tabular}{lll}
\toprule
\textbf{Data type} & \textbf{Default alignment} & \textbf{Example} \\
\midrule
Text & Left & Roll No \\
Number & Right & 4500 \\
Date & Right & 15-08-2026 \\
\bottomrule
\end{tabular}
\end{center}

A quick visual check therefore tells you what Excel thinks the entry is: a
value sitting on the \emph{left} is treated as text, and a value sitting on
the \emph{right} is treated as a number (or a date, which is stored as a
number). This default alignment is a common exam point.

\section{Exam retrieval and common confusions}
Six distinctions prevent most errors on this topic.
\begin{enumerate}
  \item A \textbf{workbook} is the file; a \textbf{worksheet} is one page
    inside it. Do not swap them.
  \item A \textbf{cell} is one box; a \textbf{range} is two or more cells
    joined by a colon.
  \item A cell address is \textbf{column letter first, then row number}
    (B5), never the reverse.
  \item A worksheet in Excel 2007+ has \textbf{1,048,576 rows} and
    \textbf{16,384 columns}; the last column is \textbf{XFD}, not IV (IV was
    the old \texttt{.xls} limit).
  \item The \textbf{Name Box} shows the address; the \textbf{Formula Bar}
    shows the contents.
  \item Only \textbf{one cell is active} at a time, even when a range is
    selected.
\end{enumerate}

\begin{example}
An item asks, ``What is the total number of rows in a modern Excel
worksheet?'' Recall the pair: rows are the larger count, 1,048,576, and
columns are 16,384. If the stem instead asks for the last column, the answer
is XFD; if it asks for the number of cells in a range, multiply columns by
rows.
\end{example}

\subsection*{Exercises}
\begin{enumerate}[label=\arabic*.]
  \item Distinguish a workbook from a worksheet.
  \item What is a cell? How is a cell address written, and give one example.
  \item How many rows and how many columns does a worksheet have in Excel
    2007 and later, and what is the last column label?
  \item What is a range? Write the range that covers columns A to C and rows
    1 to 10, and state how many cells it contains.
  \item State what the Name Box and the Formula Bar each display.
  \item List four basic data types that can be entered in an Excel cell and
    give the default alignment of text and of numbers.
\end{enumerate}

\subsection*{Solutions}
\begingroup\small
\begin{enumerate}[label=\arabic*.]
  \item A workbook is the Excel file itself, saved under one file name; a
    worksheet is a single sheet of rows and columns inside a workbook. A
    workbook can contain many worksheets.
  \item A cell is the box formed where a row and a column intersect. A cell
    address is written as the column letter followed by the row number, for
    example B5 (column B, row 5).
  \item A worksheet has 1,048,576 rows and 16,384 columns; the last column is
    labelled XFD.
  \item A range is a rectangular group of two or more cells written as two
    corner addresses joined by a colon. The range is A1:C10, and it contains
    $3 \times 10 = 30$ cells.
  \item The Name Box shows the address of the active cell (or the name of the
    selected range). The Formula Bar shows the contents of the active cell,
    including the formula behind a calculated result.
  \item Four basic data types are text, number, date/time, and formula. Text
    is aligned to the left by default, and numbers to the right by default.
\end{enumerate}
\endgroup

\lessonheading{25}{Excel Formulas}
\section{What a formula is}
A \textbf{formula} is an instruction that tells Microsoft Excel how to
calculate the value of a cell. Instead of holding a fixed number, the cell
holds a rule, and Excel works out the result of that rule from the other
cells, from constants, or from functions the rule refers to. The formula is
typed into the formula bar; the result is shown in the cell. When any cell
that the formula uses changes, Excel recalculates and shows the new result
automatically.

\begin{definitionbox}{Formula}
An instruction, written in a cell or the formula bar, that calculates a value
from operands such as cell references, constants, and functions. The cell
stores the formula, while the sheet displays the result.
\end{definitionbox}

This lesson covers the four building blocks of a formula: the leading
\texttt{=}, the operators, the cell references, and the way references behave
when a formula is copied. It closes with the standard error codes and the
distinctions that examiners test most often.

\section{The \texttt{=} rule and formula syntax}
Every Excel formula must begin with the equal sign \texttt{=}. This is the
single fact most often tested on this topic. Typing \texttt{=A1+A2} into a
cell is treated as a formula and calculated; typing \texttt{A1+A2} without the
leading \texttt{=} is treated as ordinary text and shown exactly as written.

The same \texttt{=} symbol also serves as a comparison operator that tests
whether two values are equal. The position tells you which role it plays: at
the very start of an entry it begins a formula, while inside a formula it
compares two operands.

\begin{definitionbox}{The \texttt{=} rule}
A formula must start with \texttt{=}. A cell entry that does not start with
\texttt{=} is not a formula.
\end{definitionbox}

A simple formula has three parts: the leading \texttt{=}, the operands, and
the operators. In the formula \texttt{=A1*B1+100}, the operands are the
references \texttt{A1} and \texttt{B1} and the constant \texttt{100}, while
the operators are \texttt{*} (multiply) and \texttt{+} (add).

\section{Formula, function, and constant}
It helps to separate three ideas that often appear together. A
\textbf{constant} is a fixed value written directly into a formula, such as
the number \texttt{100} or a piece of text in quotation marks. A
\textbf{function} is a built-in calculation with a name, such as
\texttt{SUM} or \texttt{AVERAGE}, which is called from inside a formula. A
\textbf{formula} is the whole instruction, and it may combine constants,
references, operators, and functions in one expression.

\begin{definitionbox}{Function}
A named, built-in calculation such as \texttt{SUM(A1:A10)} or
\texttt{AVERAGE(B1:B5)} that is called from within a formula and returns a
single value.
\end{definitionbox}

The distinction matters for exams because a question may ask what a formula
contains rather than what it produces. The formula \texttt{=A1+A2} contains
two relative references and one arithmetic operator. The formula
\texttt{=SUM(A1:A10)} calls a function on a range. Both begin with \texttt{=},
and both return a value to the cell that holds them. The rest of this lesson
concentrates on the references and operators that every formula uses.

\section{Operators}
An operator tells Excel what to do with the values on either side of it.
Excel has four families of operators: arithmetic, comparison, text, and
reference.

\textbf{Arithmetic operators} produce numbers.

\begin{center}
\begin{tabular}{p{0.12\textwidth}p{0.72\textwidth}}
\toprule
\textbf{Operator} & \textbf{Meaning} \\
\midrule
\texttt{+} & Addition \\
\texttt{-} & Subtraction \\
\texttt{*} & Multiplication \\
\texttt{/} & Division \\
\texttt{\textasciicircum} & Exponent (raise to a power) \\
\texttt{\%} & Percent \\
\bottomrule
\end{tabular}
\end{center}

Two symbols are easy to get wrong. Multiplication uses the asterisk
\texttt{*}, not the letter \texttt{x} or a dot; division uses the forward
slash \texttt{/}. The caret \texttt{\textasciicircum} raises a number to a
power, so \texttt{=2\textasciicircum3} is 8. The percent sign \texttt{\%}
divides by 100, so \texttt{=50\%} is 0.5.

\textbf{Comparison operators} compare two values and return \texttt{TRUE} or
\texttt{FALSE}: \texttt{=} (equal to), \texttt{>} (greater than), \texttt{<}
(less than), \texttt{>=} (greater than or equal to), \texttt{<=} (less than or
equal to), and \texttt{<>} (not equal to). The operator \texttt{<>} is the
Excel way of writing ``not equal''; it is not written \texttt{!=}.

A comparison gives a useful result even on its own: if A1 holds 50, the
formula \texttt{=A1>100} shows \texttt{FALSE}. The operators that candidates
confuse most are \texttt{=}, which tests equality when it appears inside a
formula, and \texttt{<>}, which tests inequality. Note that the leading
\texttt{=} that starts every formula and the \texttt{=} comparison operator
are the same character used in two different roles.

\textbf{Text operator.} The ampersand \texttt{\&} joins, or concatenates, two
pieces of text into one. \texttt{="Total: "\&A1} shows the words
\texttt{Total:} followed by the contents of cell A1.

\textbf{Reference operators} build references to groups of cells. The colon
\texttt{:} is the range operator, so \texttt{A1:A10} means every cell from A1
down to A10. The comma \texttt{,} is the union operator, which joins two
ranges into one list. A single space between two references is the
intersection operator, which returns the cells common to both ranges.

The four families can be summarised in one place. Arithmetic operators
(\texttt{+ - * / \textasciicircum \%}) produce numbers. Comparison operators
(\texttt{= > < >= <= <>}) produce \texttt{TRUE} or \texttt{FALSE}. The text
operator \texttt{\&} produces a joined piece of text. The reference operators
\texttt{:}, \texttt{,}, and the space build ranges and lists of cells.
Knowing which family an operator belongs to tells you what kind of result to
expect.

\section{Order of evaluation (precedence)}
When a formula has several operators, Excel does not simply work from left to
right. It follows a fixed order of precedence. Parentheses are evaluated
first. Then the percent sign, then the exponent, then multiplication and
division, then addition and subtraction. After the arithmetic comes the
concatenation operator \texttt{\&}, and finally the comparison operators.

Because multiplication and division are evaluated before addition and
subtraction, \texttt{=2+3*4} is 14, not 20. The multiplication \texttt{3*4} is
done first to give 12, and then 2 is added. In the same way
\texttt{=2+3*4} differs from \texttt{=(2+3)*4}, which is 20. Parentheses are
the reliable way to control the order of evaluation.

\begin{example}
In \texttt{=10+2*5} the multiplication is done first, so the result is 20. If
the intention was to add first, the formula must be written
\texttt{=(10+2)*5}, which gives 60. The exponent binds even tighter than
multiplication, so \texttt{=2*3\textasciicircum2} is 18, because
\texttt{3\textasciicircum2} is evaluated before the multiplication by 2.
\end{example}

\section{Cell references}
A \textbf{cell reference} names a cell by its column letter and row number. In
the default \texttt{A1} reference style, \texttt{A1} is column A, row 1, and
\texttt{B5} is column B, row 5. A reference lets a formula use the current
value of another cell, so the formula stays correct when that value changes.

The crucial property of a reference is not what it points at now but how it
behaves when the formula is \emph{copied}. There are three kinds of
reference: relative, absolute, and mixed.

\section{Relative references}
A reference such as \texttt{A1} is \textbf{relative}. Excel stores it, in
effect, as an offset from the cell that contains the formula. When the formula
is copied to another cell, the reference shifts by the same distance, so it
keeps pointing at the cell that lies in the same relative position.

\begin{example}
If \texttt{=A1} is entered in cell B1 and then copied to B2, the formula
becomes \texttt{=A2}. The reference moved down one row because the copied
formula moved down one row. Copy the same formula from B1 to C1, and it
becomes \texttt{=B1}: the reference moved one column to the right because the
formula did.
\end{example}

Relative references are the default and are what you want for a column of
repeated calculations, such as adding each row's two numbers.

\section{Absolute references}
A reference such as \texttt{\$A\$1} is \textbf{absolute}. The dollar sign
locks the part of the reference that follows it: the first \texttt{\$} locks
the column, and the second \texttt{\$} locks the row. An absolute reference
does not change when the formula is copied; it always points at the same cell.

\begin{example}
If \texttt{=\$A\$1} is entered in B1 and copied to B5, it remains
\texttt{=\$A\$1}. Neither the column nor the row moved. A fixed value such as
a tax rate or a constant kept in one cell should be referenced absolutely, so
that copying the formula does not disturb it.
\end{example}

\section{Mixed references}
A \textbf{mixed} reference locks only one part. In \texttt{\$A1} the column is
locked and the row is relative; copying the formula to another row changes the
row number but never the column. In \texttt{A\$1} the row is locked and the
column is relative; copying the formula to another column changes the column
letter but never the row.

A useful way to remember this is that the \texttt{\$} locks whatever
immediately follows it. A \texttt{\$} before a letter locks the column, and a
\texttt{\$} before a number locks the row.

\section{Copying and filling formulas}
The behaviour of a reference under copying is the heart of the topic. When a
formula is copied, only the unlocked parts of each reference shift; the locked
parts stay put. The table below compares the four reference types when a
formula is copied one column right and one row down, from cell B1 to cell C2.

\begin{center}
\begin{tabular}{p{0.20\textwidth}p{0.30\textwidth}p{0.34\textwidth}}
\toprule
\textbf{Reference} & \textbf{What is locked} & \textbf{After copying B1 to C2} \\
\midrule
\texttt{A1} & nothing & \texttt{B2} \\
\texttt{\$A\$1} & column and row & \texttt{\$A\$1} \\
\texttt{\$A1} & column & \texttt{\$A2} \\
\texttt{A\$1} & row & \texttt{B\$1} \\
\bottomrule
\end{tabular}
\end{center}

Reading the table: the relative reference \texttt{A1} moved on both axes; the
absolute reference \texttt{\$A\$1} moved on neither; \texttt{\$A1} kept its
column but followed the row; and \texttt{A\$1} kept its row but followed the
column.

A second, very common direction is copying a formula straight down a column.
If \texttt{=A1*B1} sits in C1 and is copied down to C2, the formula becomes
\texttt{=A2*B2}: both relative references followed the copy down one row. If
the value in B1 must stay fixed, the formula should instead read
\texttt{=A1*\$B\$1}, so that only A1 follows the copy while \texttt{\$B\$1}
stays anchored.

The same four references behave as follows when copied one column to the
right, and when copied one row down.

\begin{center}
\begin{tabular}{p{0.20\textwidth}p{0.32\textwidth}p{0.32\textwidth}}
\toprule
\textbf{Reference} & \textbf{Copied right (B1$\rightarrow$C1)} & \textbf{Copied down (B1$\rightarrow$B2)} \\
\midrule
\texttt{A1} & \texttt{B1} & \texttt{A2} \\
\texttt{\$A\$1} & \texttt{\$A\$1} & \texttt{\$A\$1} \\
\texttt{\$A1} & \texttt{\$A1} & \texttt{\$A2} \\
\texttt{A\$1} & \texttt{B\$1} & \texttt{A\$1} \\
\bottomrule
\end{tabular}
\end{center}

Reading either column, the rule is the same: a reference moves only along the
axis that is not locked. The relative reference moves on both axes, the
absolute reference on neither, and each mixed reference on exactly one.

\section{Worked examples}

\begin{example}
A price list has prices in column A and a fixed tax rate stored in cell
\texttt{\$B\$1}. The formula in cell C2 is \texttt{=A2*\$B\$1}. When it is
copied down, C3 becomes \texttt{=A3*\$B\$1} and C4 becomes \texttt{=A4*\$B\$1}:
the price reference moves down with each row, while the rate reference stays
locked on \texttt{\$B\$1}. If the formula had been written \texttt{=A2*B1},
copying it down would have moved the rate to B2 and then B3, giving wrong
results.
\end{example}

\begin{example}
In cell B1 the formula is \texttt{=\$A1}. It is copied to C2. The column lock
holds A, and the relative row follows the copy down one row, so C2 contains
\texttt{=\$A2}. If instead the original formula were \texttt{=A\$1}, copying
from B1 to C2 would keep the row 1 and move the column one step from A to B,
giving \texttt{=B\$1}.
\end{example}

\begin{example}
A sheet lists marks in column B. To flag each student as passing, the formula
\texttt{=B2>=40} is entered in C2 and copied down. It returns \texttt{TRUE}
for a mark of 40 or more and \texttt{FALSE} otherwise. Because B2 is relative,
the reference follows the copy down the column, so C3 tests B3 and C4 tests
B4. The comparison operator produces the \texttt{TRUE}/\texttt{FALSE} result,
not a number.
\end{example}

\begin{example}
A score sheet holds marks in B2:B11 and the class total in cell \texttt{\$B\$12}.
To show each student's share as a percentage, the formula in C2 is
\texttt{=B2/\$B\$12*100}. The relative \texttt{B2} follows the copy down the
column, while \texttt{\$B\$12} stays anchored on the total. The operators here
are division and multiplication, evaluated from left to right at the same
precedence, and the result is a plain number. If the total were written as a
relative reference, copying the formula down would corrupt the denominator.
\end{example}

\section{Common formula errors}
When a formula cannot be evaluated, Excel shows an error value in the cell.
The codes tested in exams are:

\begin{center}
\begin{tabular}{p{0.18\textwidth}p{0.68\textwidth}}
\toprule
\textbf{Error} & \textbf{Meaning} \\
\midrule
\texttt{\#DIV/0!} & Division by zero, or by an empty cell. \\
\texttt{\#VALUE!} & The wrong type of operand, such as text where a number is needed. \\
\texttt{\#REF!} & A referenced cell was deleted, so the reference is no longer valid. \\
\texttt{\#NAME?} & An unknown function, or a misspelled name or text entry. \\
\texttt{\#N/A} & The value is not available, as when a lookup finds no match. \\
\texttt{\#NULL!} & The two ranges given do not intersect. \\
\bottomrule
\end{tabular}
\end{center}

Remember the two that examiners use most: \texttt{\#DIV/0!} for division by
zero and \texttt{\#REF!} for a reference that has been broken by deleting a
cell. A related warning is the \textbf{circular reference}, which appears when
a formula refers, directly or through a chain of other formulas, back to its
own cell, so Excel cannot work out a value. Deleting the referenced row or
column is the usual cause of \texttt{\#REF!}, while supplying the wrong type
of operand, such as text where a number is expected, produces
\texttt{\#VALUE!}.

\section{Exam retrieval and common confusions}
Six distinctions prevent most errors on this topic.
\begin{enumerate}
  \item A formula must start with \textbf{\texttt{=}} (TRAP-30, PYQ-38). The
    options \texttt{+}, \texttt{\%}, and \texttt{@} are distractors.
  \item \textbf{\texttt{\$}} before a letter locks the column; \textbf{\texttt{\$}}
    before a number locks the row. Do not swap these.
  \item \textbf{Relative} references change on copy; \textbf{absolute}
    references do not. A \textbf{mixed} reference changes on one axis only.
  \item \texttt{\#REF!} means a broken reference; \texttt{\#DIV/0!} means
    division by zero. Do not swap the error codes.
  \item \texttt{<>} is the Excel operator for ``not equal''; multiplication is
    \texttt{*} and division is \texttt{/}.
  \item A \textbf{function} such as \texttt{SUM} is called from inside a
    formula, and a \textbf{constant} is a fixed value written directly; both
    are operands, not operators.
\end{enumerate}

\begin{example}
An item asks what \texttt{=\$A\$1} becomes when copied from B1 to B5. Because
both the column and the row are locked by the two dollar signs, the reference
does not move: the answer is \texttt{=\$A\$1}. If the question instead gives
\texttt{=\$A1}, the column is locked but the row is relative, so the answer is
\texttt{=\$A5}. The dollar signs, not the row numbers, decide the answer.
\end{example}

\subsection*{Exercises}
\begin{enumerate}[label=\arabic*.]
  \item Which character must every Excel formula begin with, and what happens
    if it is omitted?
  \item List the arithmetic operators for addition, subtraction,
    multiplication, division, and exponent.
  \item Explain the difference between a relative reference \texttt{A1} and an
    absolute reference \texttt{\$A\$1}.
  \item In the mixed reference \texttt{\$A1}, which part is locked and which
    part changes when the formula is copied?
  \item The formula \texttt{=A1} is in cell B1. What does B2 contain after the
    formula is copied to it? What would B2 contain if the formula were
    \texttt{=\$A\$1}?
  \item State the meaning of the errors \texttt{\#DIV/0!} and
    \texttt{\#REF!}.
  \item What is the result of the formula \texttt{=2+3*4}, and why?
  \item Name the Excel operator that joins two pieces of text, and the operator
    that means ``not equal to''.
\end{enumerate}

\subsection*{Solutions}
\begingroup\small
\begin{enumerate}[label=\arabic*.]
  \item A formula must begin with \texttt{=}. Without it, the entry is stored
    and shown as ordinary text rather than calculated.
  \item Addition \texttt{+}, subtraction \texttt{-}, multiplication
    \texttt{*}, division \texttt{/}, and exponent \texttt{\textasciicircum}.
  \item \texttt{A1} is relative: on copying, both its column and its row shift
    by the same offset as the formula. \texttt{\$A\$1} is absolute: the dollar
    signs lock the column and the row, so it does not change on copying.
  \item In \texttt{\$A1} the column A is locked; the row is relative, so the
    row number changes when the formula is copied to a different row.
  \item Copying \texttt{=A1} from B1 to B2 gives \texttt{=A2}, because the
    relative reference follows the copy down one row. Copying \texttt{=\$A\$1}
    to B2 leaves it as \texttt{=\$A\$1}, because both parts are locked.
  \item \texttt{\#DIV/0!} means division by zero or by an empty cell;
    \texttt{\#REF!} means a reference is invalid because a referenced cell was
    deleted.
  \item The result is 14. Multiplication is evaluated before addition, so
    \texttt{3*4} is 12, and \texttt{2+12} is 14.
  \item The ampersand \texttt{\&} joins (concatenates) text, and \texttt{<>}
    means ``not equal to''.
\end{enumerate}
\endgroup

\lessonheading{26}{Excel Functions}
\section{From formulas to functions}
MS Excel is built around the idea that a value can be \emph{calculated}
instead of typed. A cell that begins with the equal sign \texttt{=} holds a
calculation, and Excel works out the result. Such a calculation is called a
\textbf{formula}. A simple formula is \texttt{=A1+A2}, which adds the values
of two cells. A formula can also use operators such as \texttt{+}, \texttt{-},
\texttt{*}, and \texttt{/} to add, subtract, multiply, and divide.

Typing every cell into a formula becomes slow and error-prone when there are
many cells. Excel therefore provides \textbf{functions}: ready-made, named
operations that do a common job for you. \texttt{SUM} adds a range of
numbers, \texttt{AVERAGE} finds their mean, and \texttt{MAX} finds the largest.
A function is always written \emph{inside} a formula. The function name is
followed by parentheses that contain the \emph{arguments} --- the cells or
values the function should work on.

\begin{definitionbox}{Formula}
An expression that begins with the equal sign \texttt{=} and calculates a
value, for example \texttt{=A1+A2} or \texttt{=SUM(A1:A10)}.
\end{definitionbox}

\begin{definitionbox}{Function}
A predefined, named built-in operation, such as \texttt{SUM}, \texttt{COUNT},
or \texttt{IF}, that is used inside a formula to perform a particular
calculation on its arguments.
\end{definitionbox}

The distinction to hold is simple. A \textbf{formula} is the whole
calculation. A \textbf{function} is a built-in tool that performs one part of
that calculation. Every function is written as a formula, but not every
formula contains a function: \texttt{=A1+A2} is a formula with no function in
it, while \texttt{=SUM(A1:A10)} is a formula that uses one function. A
formula must begin with \texttt{=}; if the equal sign is missing, Excel treats
the entry as plain text and shows it as typed.

\section{Function syntax and arguments}
Every function follows the same shape: a name, then a pair of parentheses, and
inside the parentheses the arguments separated by commas. This is written as
\texttt{=FUNCTION(argument1, argument2)}. The \textbf{name} tells Excel which
operation to perform. The \textbf{arguments} tell it which cells or values to
use. For example, in \texttt{=SUM(A1:A10)} the name is \texttt{SUM} and the
single argument is the range \texttt{A1:A10}.

A \textbf{range} is a block of cells written with a colon between the first
and last cell. \texttt{A1:A10} means the cells A1, A2, A3, and so on down to
A10. A function can take more than one argument: \texttt{=SUM(A1:A5, C1:C5)}
adds the two ranges together. The comma separates one argument from the next.
Knowing the name, the parentheses, and the colon is enough to read any of the
basic functions in this lesson.

\begin{center}
\begin{tabular}{p{0.15\textwidth}p{0.28\textwidth}p{0.47\textwidth}}
\toprule
\textbf{Function} & \textbf{Example} & \textbf{What it returns} \\
\midrule
\texttt{SUM} & \texttt{=SUM(A1:A10)} & The total of the numbers in the range. \\
\texttt{AVERAGE} & \texttt{=AVERAGE(A1:A10)} & The arithmetic mean of the numbers. \\
\texttt{COUNT} & \texttt{=COUNT(A1:A10)} & The number of cells holding numbers. \\
\texttt{COUNTA} & \texttt{=COUNTA(A1:A10)} & The number of non-empty cells. \\
\texttt{MAX} & \texttt{=MAX(A1:A10)} & The largest number in the range. \\
\texttt{MIN} & \texttt{=MIN(A1:A10)} & The smallest number in the range. \\
\texttt{IF} & \texttt{=IF(A1>50,"Y","N")} & One of two values, chosen by the test. \\
\bottomrule
\end{tabular}
\end{center}

\section{SUM and AVERAGE}
\texttt{SUM} returns the total of the numbers in its arguments. Its most
common form is \texttt{=SUM(A1:A10)}, which adds all the numbers in the range
A1:A10. Cells that hold text and cells that are blank are ignored; only
numbers are added. SUM accepts several ranges, so \texttt{=SUM(A1:A5,B1:B5)}
adds both blocks in one step.

\texttt{AVERAGE} returns the arithmetic mean of the numbers in its arguments.
\texttt{=AVERAGE(A1:A10)} adds the numbers and divides by how many numbers
there are. The important detail is what counts in the divisor: blank cells and
cells holding text are \emph{not} counted. AVERAGE divides by the number of
\emph{numeric} cells, not by the total number of cells in the range. This is
the source of many wrong answers.

\begin{example}
Suppose A1:A5 holds the values 10, 20, and 30, while A4 and A5 are blank.
\texttt{=SUM(A1:A5)} returns $60$. \texttt{=AVERAGE(A1:A5)} returns
$(10+20+30)/3 = 20$, not $60/5 = 12$. Excel divides by the three cells that
actually hold numbers, because the two blanks are ignored.
\end{example}

\section{COUNT and COUNTA: the near-neighbour pair}
\texttt{COUNT} and \texttt{COUNTA} look almost the same but count different
things, and the difference is a favourite exam trap. \texttt{COUNT(counts)}
only the cells in a range that contain \emph{numbers}. \texttt{COUNTA} (the
``A'' stands for \emph{all}) counts every cell that is \emph{not empty},
whether it holds a number or text. Both functions ignore blank cells.

\begin{definitionbox}{COUNT}
A function that returns the number of cells in a range that contain numeric
values. Text, blanks, and logical values are not counted.
\end{definitionbox}

\begin{definitionbox}{COUNTA}
A function that returns the number of cells in a range that are not empty. It
counts numbers, text, and error values alike; only truly blank cells are
ignored.
\end{definitionbox}

The contrast is easiest to see with a mixed range. Suppose A1:A10 holds six
numbers, two text labels, and two blank cells. \texttt{=COUNT(A1:A10)}
returns 6, because only the numeric cells are counted. \texttt{=COUNTA(A1:A10)}
returns 8, because the six numbers and the two text labels are all non-empty.
Both functions ignore the two blanks. A related function, \texttt{COUNTBLANK},
does the opposite and counts the empty cells, returning 2 for the same range.
For the exam, the pair to memorise is: \textbf{COUNT = numbers only; COUNTA =
all non-empty cells}.

\section{MAX and MIN}
\texttt{MAX} returns the largest number in a range, and \texttt{MIN} returns
the smallest. Their forms are \texttt{=MAX(C1:C10)} and, for the smallest,
\texttt{=MIN(C1:C10)}. Like SUM, AVERAGE, and COUNT, they work on numbers and
ignore text and blank cells. A useful way to remember them is that \texttt{MAX}
is the maximum and \texttt{MIN} is the minimum; do not confuse \texttt{MIN}
with the word ``minimum cell'', which is not a function at all.

\begin{example}
A range holds the marks 42, 67, 55, 91, and 38. \texttt{=MAX(C1:C5)} returns
$91$, the top mark. \texttt{=MIN(C1:C5)} returns $38$, the lowest mark.
\texttt{=AVERAGE(C1:C5)} returns $58.6$, the mean of the five numbers.
\end{example}

\section{The IF function and nested IF}
\texttt{IF} lets a cell show one of two results depending on whether a
condition is true or false. Its syntax is
\texttt{=IF(logical\_test, value\_if\_true, value\_if\_false)}. The first
argument is a \emph{logical test} that can be answered only with true or
false, such as \texttt{A1>50}. The second argument is the value or text to
return when the test is true, and the third is the value to return when it is
false. Thus \texttt{=IF(A1>50,"Pass","Fail")} shows \texttt{Pass} when A1 is
more than 50 and \texttt{Fail} otherwise.

A \textbf{nested IF} places a second IF inside the value-if-false part of the
first, so more than one condition can be tested. A typical form is
\texttt{=IF(A1>50,"High",IF(A1>25,"Medium","Low"))}. Excel checks the
conditions in the order they are written, and the first test that is true
decides the result.

\begin{example}
Take the formula \texttt{=IF(A1>50,"High",IF(A1>25,"Medium","Low"))} and let
A1 be blank. In a numeric comparison Excel reads a blank cell as $0$. Since
$0$ is neither greater than 50 nor greater than 25, both tests are false and
the formula returns the final value, \texttt{Low}. This is the behaviour the
Junior Assistant 2026 paper examined (KB PYQ-14, JA 2026 Q74).
\end{example}

When the input is \emph{non-numeric} text, the behaviour depends on what the
formula does with it. A comparison of text against a number is still answered
by Excel, but if the formula tries to \emph{calculate} with text --- for
example, \texttt{=A1*2} where A1 holds a word --- Excel cannot do the
arithmetic and returns the error value \texttt{\#VALUE!}. The safe exam
reading is: a blank is treated as 0 in a numeric test, and text that is
arithmetic-only produces \texttt{\#VALUE!}.

\section{A worked sheet, function by function}
A small sheet makes the behaviour of every function concrete. Column A holds
five student names --- Anita, Bilal, Chandan, Divya, and Esha --- and column B
holds the marks 40, 55, 70, and 60, with Esha's mark (B5) left blank because
she has not been marked yet. Reading the sheet function by function shows what
each one does with the blank.

\begin{center}
\begin{tabular}{p{0.30\textwidth}p{0.30\textwidth}p{0.28\textwidth}}
\toprule
\textbf{Formula} & \textbf{Result} & \textbf{Why} \\
\midrule
\texttt{=SUM(B1:B5)} & 225 & Adds 40, 55, 70, 60; the blank is ignored. \\
\texttt{=AVERAGE(B1:B5)} & 56.25 & Divides 225 by 4 numeric cells, not by 5. \\
\texttt{=COUNT(B1:B5)} & 4 & Counts the four cells that hold numbers. \\
\texttt{=COUNTA(B1:B5)} & 4 & Counts the four non-empty marks; B5 is blank. \\
\texttt{=MAX(B1:B5)} & 70 & The largest of the four marks. \\
\texttt{=MIN(B1:B5)} & 40 & The smallest of the four marks. \\
\texttt{=IF(B5>=40,"Pass","Fail")} & Fail & B5 is blank, read as 0, so the test fails. \\
\bottomrule
\end{tabular}
\end{center}

The last row links the sheet back to the nested-IF behaviour: a blank cell is
read as zero in a numeric test, so Esha is shown as not yet passing. Note also
that \texttt{COUNTA} on the names column, \texttt{=COUNTA(A1:A5)}, would
return 5, because all five names are present even though only four marks are.
This is exactly how the two count functions are used together to spot missing
data.

\section{Other basic functions worth knowing}
Beyond the core set, a few basic functions appear in Office-suite questions.
\texttt{COUNTIF} counts the cells that meet one condition, and \texttt{SUMIF}
adds the cells that meet one condition; both extend COUNT and SUM. On the
statistics side, \texttt{MEDIAN} returns the middle value of the numbers,
while \texttt{MAX} and \texttt{MIN} return the ends of the range.
\texttt{ROUND} rounds a number to a stated number of digits. \texttt{LEN}
counts the characters in a text value. \texttt{TODAY} returns the current
date, and \texttt{NOW} returns the current date and time.

\begin{center}
\begin{tabular}{p{0.22\textwidth}p{0.62\textwidth}}
\toprule
\textbf{Function} & \textbf{What it returns} \\
\midrule
\texttt{COUNTIF} & The number of cells meeting one condition. \\
\texttt{SUMIF} & The total of the cells meeting one condition. \\
\texttt{MEDIAN} & The middle value of the numbers. \\
\texttt{ROUND} & A number rounded to a set number of digits. \\
\texttt{LEN} & The number of characters in a text value. \\
\texttt{TODAY} & The current date. \\
\bottomrule
\end{tabular}
\end{center}

At JKSSB depth, the exam is satisfied by knowing what each function returns
and that every one of them is written inside a formula that starts with
\texttt{=}.

\section{AutoSum}
\textbf{AutoSum} is a button, not a function, that quickly inserts a
\texttt{SUM} formula. When the active cell is placed below or beside a column
or row of numbers, AutoSum looks at the adjacent numbers and writes a
\texttt{=SUM(...)} formula over the range it selects. Pressing the button
therefore saves the user from typing the range by hand. The keyboard shortcut
for AutoSum is \texttt{Alt + =}. Its drop-down arrow also offers
\texttt{Average}, \texttt{Count Numbers}, \texttt{Max}, and \texttt{Min}, so
the same button can insert any of the core functions, not only SUM.

\begin{definitionbox}{AutoSum}
An Excel button that automatically inserts a \texttt{SUM} formula for the
range next to the active cell; the shortcut is \texttt{Alt + =}.
\end{definitionbox}

\section{Combining functions in one formula}
Functions can be nested inside one another, so the result of one function
becomes the argument of another. For example,
\texttt{=ROUND(AVERAGE(A1:A10),1)} first finds the mean of A1:A10 and then
rounds it to one decimal place. In the same way,
\texttt{=IF(COUNT(B1:B5)=0,0,AVERAGE(B1:B5))} shows 0 when the range contains
no numbers at all, instead of producing an error, and otherwise shows the mean
of the numbers present. Nesting is read from the inside out: the innermost
function is calculated first, and its result is passed outwards to the next
function.

\begin{example}
Consider \texttt{=IF(COUNT(C1:C6)=0,"No data",AVERAGE(C1:C6))}. If C1:C6 holds
only text or blanks, \texttt{COUNT} returns 0, the test is true, and the cell
shows \texttt{No data}. If the range holds numbers, the test is false and the
cell shows their average. This pattern is common on real marks sheets, where a
column may be empty early in the term.
\end{example}

The arguments of a function can also be single cells, several separate cells,
whole ranges, or even constant values such as \texttt{=SUM(A1, 5, B3:B7)}. A
reference such as \texttt{A1} can be relative or absolute (\texttt{\$A\$1});
that choice does not change what the function does, only how the reference
behaves when the formula is copied.

\section{Error values}
When a formula cannot be calculated, Excel displays an \emph{error value}
beginning with the hash sign \texttt{\#}. Each error points to a different
kind of problem, and recognising them is a common direct-recall question.

\begin{center}
\begin{tabular}{p{0.22\textwidth}p{0.60\textwidth}}
\toprule
\textbf{Error value} & \textbf{What it means} \\
\midrule
\texttt{\#DIV/0!} & A number is divided by zero or by a blank cell. \\
\texttt{\#VALUE!} & The wrong type of argument is used --- for example, text where a number is needed. \\
\texttt{\#NAME?} & Excel does not recognise a name in the formula, often a misspelt function name. \\
\texttt{\#REF!} & A cell reference is not valid, usually because the referenced cell was deleted. \\
\texttt{\#N/A} & The value is not available to the formula, as when a lookup finds no match. \\
\texttt{\#NUM!} & A number in the formula is not valid for the operation. \\
\texttt{\#NULL!} & A space (intersection) between two ranges that do not intersect. \\
\bottomrule
\end{tabular}
\end{center}

Two clarifications prevent confusion. First, the string \texttt{\#\#\#\#\#}
that sometimes appears in a cell is \emph{not} an error value; it only means
the column is too narrow to display the number, and it disappears when the
column is widened. Second, \texttt{\#DIV/0!} also appears when the divisor is
a blank cell, because a blank is read as zero in the division. Reading each
error as a clue to what Excel could not do is the fastest way to answer these
items.

\begin{example}
The error value points at the mistake that produced it. A formula
\texttt{=A1/0}, or \texttt{=A1/B1} where B1 is blank, shows
\texttt{\#DIV/0!}. A formula \texttt{=A1+B1} where A1 holds text shows
\texttt{\#VALUE!}. A misspelt function such as \texttt{=SUMM(A1:A5)} shows
\texttt{\#NAME?}. Deleting a cell that a formula refers to turns the reference
into \texttt{\#REF!}. A lookup that finds no match leaves the result
\texttt{\#N/A}, and a mathematically invalid argument such as
\texttt{=SQRT(-1)} gives \texttt{\#NUM!}.
\end{example}

\section{Exam retrieval and common confusions}
Five distinctions prevent most errors on this topic.
\begin{enumerate}
  \item A \textbf{formula} starts with \texttt{=}; a \textbf{function} is a
    named built-in operation used inside a formula. Without the equal sign,
    the entry is text (TRAP-30).
  \item \textbf{COUNT} counts numbers only; \textbf{COUNTA} counts all
    non-empty cells. Do not swap this near-neighbour pair.
  \item \textbf{AVERAGE} divides by the number of numeric cells, not by the
    total number of cells.
  \item A \textbf{blank} cell is read as zero in a numeric test, so a blank
    fails a ``greater than'' condition; text used in arithmetic gives
    \texttt{\#VALUE!} (KB PYQ-14).
  \item \textbf{AutoSum} inserts \texttt{=SUM(...)} and its shortcut is
    \texttt{Alt + =}, not a Ctrl combination.
\end{enumerate}

\begin{example}
An item asks which function counts a range that mixes numbers and labels.
Read the function name carefully: \texttt{COUNT} returns only the numeric
cells, while \texttt{COUNTA} returns every non-empty cell. If the stem asks
for the count of entries of any kind, the answer is \texttt{COUNTA}; if it
asks for the count of numbers, the answer is \texttt{COUNT}. A second item
may ask which error appears when a cell is divided by a blank; the blank
counts as zero, so the error is \texttt{\#DIV/0!}.
\end{example}

\subsection*{Exercises}
\begin{enumerate}[label=\arabic*.]
  \item Explain the difference between a formula and a function in MS Excel.
  \item A range A1:A8 contains 5 numbers, 2 text labels, and 1 blank cell.
    State the values returned by \texttt{=COUNT(A1:A8)} and
    \texttt{=COUNTA(A1:A8)}.
  \item Write the general syntax of the IF function and explain each of its
    three arguments.
  \item A cell holds \texttt{=IF(A1>10,"Yes",IF(A1>0,"Small","No"))} and A1 is
    blank. What is shown, and why?
  \item Match each error value with its cause: \texttt{\#DIV/0!},
    \texttt{\#VALUE!}, \texttt{\#NAME?}, \texttt{\#REF!}.
\end{enumerate}

\subsection*{Solutions}
\begingroup\small
\begin{enumerate}[label=\arabic*.]
  \item A formula is any calculation that begins with \texttt{=}, such as
    \texttt{=A1+A2}. A function is a predefined, named operation such as
    \texttt{SUM} or \texttt{IF} that is used inside a formula. Every function
    is written as a formula, but a formula need not contain a function.
  \item \texttt{=COUNT(A1:A8)} returns 5, because only the numeric cells are
    counted. \texttt{=COUNTA(A1:A8)} returns 7, because the 5 numbers and the
    2 text labels are all non-empty; both ignore the blank.
  \item \texttt{=IF(logical\_test, value\_if\_true, value\_if\_false)}. The
    first argument is a test answered true or false; the second is the value
    returned when it is true; the third is the value returned when it is
    false.
  \item The formula shows \texttt{No}. A blank is read as zero in a numeric
    test, so A1 is neither greater than 10 nor greater than 0, both tests are
    false, and the final value \texttt{No} is returned.
  \item \texttt{\#DIV/0!} --- division by zero or by a blank cell;
    \texttt{\#VALUE!} --- the wrong type of argument, such as text where a
    number is needed; \texttt{\#NAME?} --- an unrecognised name, often a
    misspelt function; \texttt{\#REF!} --- an invalid cell reference after
    the referred cell was deleted.
\end{enumerate}
\endgroup

\lessonheading{27}{Excel Data Tools}
\section{What Excel data tools do}
Once data has been entered into a worksheet, it usually has to be organised,
narrowed, summarised, and presented. Excel groups the features that do this
under what we can call its data tools: \textbf{sort}, \textbf{filter},
\textbf{tables}, \textbf{charts}, \textbf{PivotTables and PivotCharts},
\textbf{printing}, \textbf{protection}, and \textbf{what-if analysis}. These are
the features a JKSSB candidate is expected to recognise by name and by
one-line purpose.

\begin{definitionbox}{Data tools}
The worksheet features that arrange, narrow, summarise, present, print, and
safeguard data without changing the underlying values.
\end{definitionbox}

The single idea that holds the whole lesson together is this: most of these
tools change the \emph{view} of the data. Sorting, filtering, and a PivotTable
reorganise what you see on the screen; the original values in the cells stay
exactly where they were. The exceptions are deliberate edits --- a sort does
physically move rows, and a user can type new data --- but the summary tools
themselves are read-only with respect to the source.

\begin{center}
\begin{tabular}{p{0.22\textwidth}p{0.66\textwidth}}
\toprule
\textbf{Tool} & \textbf{One-line purpose} \\
\midrule
Sort & Arrange rows in a chosen order. \\
Filter & Show only the rows that meet a condition. \\
Excel table & A named, auto-expanding, formatted range of data. \\
Chart & A picture of the data, linked to its source range. \\
PivotTable & A fast summary of a large table of data. \\
PivotChart & A chart built on a PivotTable. \\
Protection & Restrict editing of cells or workbook structure. \\
What-if analysis & Test how the result changes with different inputs. \\
\bottomrule
\end{tabular}
\end{center}

\section{Sorting}
\textbf{Sorting} rearranges rows into an order chosen by the user. \textbf{Ascending}
order goes from A to Z, or from the smallest number to the largest.
\textbf{Descending} order goes from Z to A, or from the largest number to the
smallest. A sort is performed on a \emph{column}, but the whole row moves with
that column so that each record stays intact.

The \textbf{Sort} dialog (\key{Data > Sort}) can order data by \textbf{more than
one level}. For example, a payroll sheet can be sorted by Department first, and
then by Salary within each department. The dialog's \textbf{Add Level} button
adds these secondary and tertiary keys.

A \textbf{custom list} sorts by a defined order rather than alphabetically: days
of the week and months are the standard examples. Excel can also sort by cell
colour, font colour, or cell icon. When the range has titles, the option
\textbf{My data has headers} must be ticked so that the header row is not sorted
along with the data.

\begin{example}
A teacher's marksheet has the columns Roll No., Name, and Marks. Clicking
anywhere in the Marks column and choosing Sort Largest to Smallest reorders the
rows so the highest marks appear first. If only the Marks column is selected and
the rest of the row does not move, the records are scrambled --- the names no
longer match the marks.
\end{example}

\textbf{Trap watch.} Sorting a \textbf{single column} on its own scrambles the
records. Always sort the whole table, or convert the data to an Excel table
first.

\section{Filtering}
\textbf{Filtering} displays only the rows that satisfy a condition and hides the
rest. The hidden rows are \textbf{not deleted}; they are temporarily out of
sight. AutoFilter is switched on and off with \key{Data > Filter}, whose
keyboard shortcut is \key{Ctrl}$+$\key{Shift}$+$\key{L}. Once it is on,
drop-down arrows appear in the header row of each column.

\begin{definitionbox}{Filter}
A tool that shows only the rows matching a condition while hiding --- not
removing --- the other rows.
\end{definitionbox}

Clicking a drop-down arrow allows two broad kinds of narrowing:
\begin{itemize}
  \item a \textbf{value filter}, where one or more ticks are kept from a list of
    the column's values;
  \item a \textbf{condition filter}, where a rule is applied. Text Filters,
    Number Filters, and Date Filters offer common rules such as \textit{Equals},
    \textit{Greater Than}, \textit{Begins With}, and \textit{Between}.
\end{itemize}
\textbf{Top 10} and \textbf{Above Average} are ready-made condition filters.
\textbf{Advanced Filter} is used for more complex criteria and can copy the
matching rows to another location on the sheet.

Clearing the filter with \textbf{Clear Filter} restores every row. The key
contrast to remember is that filtering \emph{hides} rows; it does not remove
them and does not change the data.

\begin{example}
A sales sheet has a Region column. Applying a value filter and ticking only
``North'' hides all other regions. The rows for South, East, and West are still
present in the worksheet; clearing the filter brings them back unchanged.
\end{example}

\section{Excel tables}
An \textbf{Excel table} is a range of data that has been converted with
\key{Insert > Table} (shortcut \key{Ctrl}$+$\key{T}) and given a table style. A
table has a \textbf{name} (by default \textit{Table1}), and it behaves as a
single object rather than as loose cells.

\begin{definitionbox}{Excel table}
A named, formatted range of data that automatically expands to include new
rows and columns and whose formulas use structured references.
\end{definitionbox}

The practical advantages of a table over a plain range are:
\begin{itemize}
  \item it \textbf{auto-expands} when a new row is typed immediately below it;
  \item it shows \textbf{banded rows} that improve readability;
  \item it adds \textbf{filter buttons} to the header row automatically;
  \item it can display an optional \textbf{Total Row}; and
  \item its formulas use \textbf{structured references} such as
    \key{Table1[Sales]} instead of cell addresses such as \key{B2}.
\end{itemize}
Because a structured reference points at the whole column, it automatically
covers new rows that are added to the table. A formula written with cell
addresses would have to be extended by hand.

\begin{center}
\begin{tabular}{lll}
\toprule
\textbf{Point} & \textbf{Plain range} & \textbf{Excel table} \\
\midrule
Name & None & Assigned name (Table1) \\
Grows & Manually & Auto-expands \\
Formulas & Cell addresses & Structured references \\
Filters & Manual AutoFilter & Built-in header buttons \\
\bottomrule
\end{tabular}
\end{center}

A table is \emph{not} a PivotTable. A table is the \textbf{source data}; a
PivotTable is one of the summaries that can be built on top of it.

\section{Charts}
A \textbf{chart} is a graphical picture of worksheet data. Because the chart is
\textbf{linked} to its source range, editing the numbers in the cells updates
the chart automatically. The common chart types and their uses are:

\begin{center}
\begin{tabular}{p{0.26\textwidth}p{0.62\textwidth}}
\toprule
\textbf{Chart type} & \textbf{Best used for} \\
\midrule
Column / bar & Comparing values across categories. \\
Line & Showing a trend, usually over time. \\
Pie / doughnut & Showing the parts of a single whole. \\
Scatter (XY) & Showing the relationship between two numeric variables. \\
\bottomrule
\end{tabular}
\end{center}

The parts of a chart carry their own names, and exam items often test them. The
\textbf{chart area} is the whole chart, including the frame. Inside it, the
\textbf{plot area} is the region where the data is actually drawn. Each set of
values is a \textbf{data series}, and the \textbf{legend} names each series. The
\textbf{category axis} (usually the X-axis) carries the labels, while the
\textbf{value axis} (the Y-axis) carries the numbers. \textbf{Data labels}
display the exact value at each data point, and \textbf{gridlines} help read the
values.

A chart normally floats on the same worksheet as its data, but it can be moved
to a separate \textbf{chart sheet} that contains only the chart.

\begin{example}
A school records the number of students in each class. A column chart compares
the classes at a glance. The same data as a line chart would suggest a trend
across the classes, which may not be meaningful; for parts of one total, such as
the share of boys and girls in the school, a pie chart is the right choice.
\end{example}

\section{PivotTables and PivotCharts}
A \textbf{PivotTable} summarises a large table of data in seconds. Instead of
writing formulas, the user \textbf{drags} fields into four areas:

\begin{center}
\begin{tabular}{p{0.20\textwidth}p{0.66\textwidth}}
\toprule
\textbf{Area} & \textbf{What goes there} \\
\midrule
Rows & Field values listed down the left side of the summary. \\
Columns & Field values spread across the top of the summary. \\
Values & The numbers that are summarised (Sum, Count, Average, and more). \\
Filters & Fields used to filter the whole PivotTable. \\
\bottomrule
\end{tabular}
\end{center}

The \textbf{Values} area chooses the aggregation: by default a numeric field is
summed and a text field is counted, but Count, Average, Max, Min, and other
functions can be chosen. Rearranging the field layout re-computes the summary
instantly --- this quick re-orientation is what gives the PivotTable its name.

\begin{definitionbox}{PivotTable}
An interactive summary that aggregates a large range of data by dragging fields
into Rows, Columns, Values, and Filters areas.
\end{definitionbox}

A \textbf{PivotChart} is a chart built on a PivotTable and linked to it. When
the PivotTable layout changes, the PivotChart changes with it.

The most important contrast is that a PivotTable \textbf{summarises} the source
data; it does not replace it and does not change it. If the source data is
edited, the PivotTable must be \textbf{refreshed} to show the new result.
Deleting the source range breaks the PivotTable, and a PivotTable is never the
same thing as an Excel table.

\section{Printing worksheets}
A worksheet is printed from \key{File > Print}, and its layout is controlled
from the \key{Page Layout} tab. The key settings are:
\begin{itemize}
  \item \textbf{Orientation} --- \textbf{Portrait} for a tall page or
    \textbf{Landscape} for a wide one;
  \item \textbf{Paper size}, \textbf{margins}, and \textbf{headers/footers};
  \item \textbf{Print Area} --- select a range and choose \key{Set Print Area}
    to print only that part of the sheet;
  \item \textbf{Print Titles} --- repeats chosen header rows or columns at the
    top or left of every printed page, so a long table keeps its headings; and
  \item \textbf{Scaling} --- options such as ``Fit Sheet on One Page'' shrink
    the printout to a chosen width or number of pages.
\end{itemize}
\textbf{Gridlines} are the faint cell borders seen on screen. They are
\textbf{not printed by default}; the \key{Print} gridlines option under Sheet
Options must be enabled to show them on paper. \textbf{Page breaks} can be
inserted to control where a new printed page begins.

\begin{example}
A 200-row price list is being printed. Without Print Titles, only the first page
shows the column headings and later pages are unreadable. Setting the heading
row as a print title repeats it on every page. If the list is a little too wide,
landscape orientation or ``Fit All Columns on One Page'' keeps it tidy.
\end{example}

\section{Protecting worksheets and workbooks}
\textbf{Protection} restricts changes that can be made to a workbook.
\textbf{Protect Sheet} (\key{Review > Protect Sheet}) prevents cells on that
sheet from being edited, while allowing selected actions such as selecting
locked cells or formatting. \textbf{Protect Workbook} protects the
\textbf{structure} of the workbook --- preventing sheets from being added,
deleted, renamed, or moved --- and can also protect its windows.

\begin{definitionbox}{Protection}
A restriction placed on a worksheet or workbook that prevents specified changes,
optionally guarded by a password.
\end{definitionbox}

Cell locking is the detail most candidates miss. Every cell is \textbf{Locked}
by default in the Format Cells dialog, but this lock has \textbf{no effect until
the sheet is protected}. To let a user edit some cells and not others, the
chosen cells must first be \textbf{unlocked}, and then the sheet must be
protected.

Protection is not the same as encryption. Protection restricts editing inside
Excel; it is not the password that encrypts the whole file (that is a separate
file-level password).

\section{What-if analysis}
\textbf{What-if analysis} answers the question ``what would the result be if an
input changed?''. Excel provides three tools for this at the depth a JKSSB
candidate needs:

\begin{center}
\begin{tabular}{p{0.26\textwidth}p{0.62\textwidth}}
\toprule
\textbf{Tool} & \textbf{What it does} \\
\midrule
Goal Seek & Finds the input value needed to reach a target result. \\
Scenario Manager & Stores and compares several sets of input values. \\
Data Table & Shows how an output changes as one or two inputs change. \\
\bottomrule
\end{tabular}
\end{center}

All three are found under \key{Data > What-If Analysis}. Goal Seek works
backwards from a target: for example, it finds the selling price that makes the
profit exactly a desired figure. Scenario Manager keeps named alternatives such
as ``Best case'' and ``Worst case''. A Data Table is a grid that shows a range
of results for one or two changing inputs.

\section{Near-neighbour contrast: COUNTIFS vs SUMPRODUCT}
To \textbf{count rows that satisfy several conditions at the same time}, Excel
offers two tools that students often confuse.

\textbf{COUNTIFS} is the direct multi-criterion counting function. Its arguments
are written as \textbf{paired} ranges and criteria:
\[
\texttt{COUNTIFS(range1, crit1, range2, crit2, \ldots)}
\]
It applies \textbf{AND} logic: a row is counted only if it meets \emph{every}
range-and-criteria pair. For example,
\texttt{COUNTIFS(A2:A100,"North",B2:B100,">500")} counts the rows that are in
the North region \emph{and} have a value greater than 500.

\textbf{SUMPRODUCT} is a general array function. It can also count
multi-criterion rows by multiplying boolean results:
\[
\texttt{SUMPRODUCT((A2:A100="North")*(B2:B100>500))}
\]
Each TRUE/FALSE comparison becomes 1 or 0 when multiplied, and SUMPRODUCT adds
the products, giving the same count. Because the multiplication already coerces
the boolean values to numbers, SUMPRODUCT does not need the \key{--} (double
unary) that is sometimes seen when boolean arrays are fed to other functions.

For \textbf{OR} logic --- a row that meets either of two conditions --- the two
matches are \emph{added}, for example
\texttt{COUNTIF(range,"North")+COUNTIF(range,"South")}.

\begin{example}
A marksheet has a Stream column and a Marks column. To count students who are in
the Science stream \emph{and} scored more than 80, write
\[ \texttt{COUNTIFS(Stream,"Science",Marks,">80")}. \]
The same count can be obtained with
\[ \texttt{SUMPRODUCT((Stream="Science")*(Marks>80))}. \]
Confusing the paired arguments of COUNTIFS, or thinking COUNTIFS cannot handle
two conditions at once, is the classic error (TRAP-14, PYQ-15).
\end{example}

\section{Exam retrieval and common confusions}
Six distinctions prevent most errors on this topic.
\begin{enumerate}
  \item \textbf{Sort} rearranges rows into an order; \textbf{filter} shows only
    matching rows and \textbf{hides} the rest --- it does not delete them.
  \item A \textbf{PivotTable} summarises the source data but is \textbf{not} the
    source table; it must be \textbf{refreshed} after the data changes.
  \item A PivotTable is \textbf{not} an Excel table. The table is the source;
    the PivotTable is a summary built on it.
  \item \textbf{COUNTIFS} counts with AND across \textbf{paired}
    range-and-criteria arguments; \textbf{SUMPRODUCT} does the same with
    boolean arrays (TRAP-14).
  \item \textbf{Locked} cells are protected only \textbf{after} the sheet is
    protected; the default lock has no effect on its own.
  \item \textbf{Gridlines} are \textbf{not printed} unless Print gridlines is
    enabled.
\end{enumerate}

\begin{example}
An item states, ``A PivotTable changes the source data and updates on its own.''
Both halves are wrong. A PivotTable never changes the source; and after the
source changes it does not update by itself --- the user must refresh it. The
stable concept is: PivotTable = summary, source = untouched, refresh = manual.
\end{example}

\subsection*{Exercises}
\begin{enumerate}[label=\arabic*.]
  \item Distinguish between sorting and filtering, and state what happens to the
    rows that a filter does not show.
  \item Name two features that an Excel table provides but a plain range does
    not.
  \item Match each chart to its best use: column, line, pie, scatter.
  \item List the four areas of a PivotTable and say what goes into each.
  \item Explain why a PivotTable must be refreshed after the source data
    changes.
  \item What is the difference between Protect Sheet and Protect Workbook?
    When does a locked cell actually become protected?
  \item Write the COUNTIFS formula that counts rows in the North region with a
    value greater than 500, and give the equivalent SUMPRODUCT formula.
\end{enumerate}

\subsection*{Solutions}
\begingroup\small
\begin{enumerate}[label=\arabic*.]
  \item Sorting rearranges rows into a chosen order; filtering shows only the
    rows matching a condition. A filter \emph{hides} the other rows temporarily;
    it does not delete them, and clearing the filter restores them.
  \item An Excel table has a name and auto-expands when new rows are typed; it
    also provides structured references, banded rows, and built-in filter
    buttons. A plain range has none of these by default.
  \item A column chart compares values across categories; a line chart shows a
    trend over time; a pie chart shows parts of one whole; a scatter chart shows
    the relationship between two numeric variables.
  \item Rows --- field values down the left; Columns --- field values across the
    top; Values --- the numbers that are summarised; Filters --- fields used to
    filter the whole PivotTable.
  \item A PivotTable stores only a summary view; it does not track the source
    automatically. After the source data changes, the user must Refresh the
    PivotTable so the summary is recalculated.
  \item Protect Sheet stops cells on a sheet from being edited; Protect Workbook
    protects the workbook structure (adding, deleting, renaming, or moving
    sheets). A cell is Locked by default, but the lock takes effect only after
    the sheet is protected.
  \item \texttt{COUNTIFS(A2:A100,"North",} \texttt{B2:B100,">500")}; the
    equivalent is \texttt{SUMPRODUCT((A2:A100="North")}
    \texttt{*(B2:B100>500))}.
\end{enumerate}
\endgroup

\lessonheading{28}{Access Database Basics}
\section{What a database is}
A \textbf{database} is an organised collection of related data, stored so
that it can be searched, updated, and retrieved easily. The data is not kept
as a loose pile of separate files; it is arranged in a planned structure.
A simple example is a school's information: the list of students, their
classes, and their marks all belong together as related information, and a
database keeps them that way.

A useful comparison is a spreadsheet. A spreadsheet (such as Excel) is built
for calculation on a simple grid. A database is built for storing and linking
large amounts of related data, and it can hold many tables that are related
to one another. When a record system grows large, a database keeps the data
consistent and avoids unnecessary duplication.

\begin{definitionbox}{Database}
An organised collection of related data, stored so that it can be searched,
updated, and retrieved easily.
\end{definitionbox}

\section{Database vs DBMS}
A database and a DBMS are not the same thing, and the difference is a common
exam point. The \textbf{database} is the collection of data itself. A
\textbf{DBMS} (Database Management System) is the software that manages that
data. The database is data; the DBMS is a program.

\begin{center}
\begin{tabular}{p{0.18\textwidth}p{0.35\textwidth}p{0.37\textwidth}}
\toprule
\textbf{Point} & \textbf{Database} & \textbf{DBMS} \\
\midrule
What it is & The organised collection of data & The software that manages the data \\
Nature & Data & Program \\
Example & A school's records & MS Access, Oracle, MySQL \\
\bottomrule
\end{tabular}
\end{center}

\begin{definitionbox}{DBMS (Database Management System)}
The software that stores a database and lets users insert, update, delete,
and retrieve data while controlling access and keeping the data consistent.
\end{definitionbox}

A \textbf{relational database management system} (RDBMS) is a DBMS that keeps
data in related tables. Microsoft Access is an example of an RDBMS.

\section{What a DBMS does}
A DBMS performs several jobs. It \emph{stores} the data and controls how it is
kept on disk. It lets users \emph{insert}, \emph{update}, \emph{delete}, and
\emph{retrieve} records. It controls \emph{access}, so that only permitted
users can read or change the data. It also keeps the data
\emph{consistent} --- different parts of the database are not allowed to
contradict each other --- and protects it from corruption.

Common DBMS software includes MS Access, Oracle, MySQL, and SQL Server. For
this course, the important point is that Access is a DBMS: it is the program
that manages a database, not the data itself.

\section{Microsoft Access in the Office suite}
\textbf{Microsoft Access} is the database program of the Microsoft Office
suite. In the suite, Word handles text documents, Excel handles spreadsheets,
PowerPoint handles presentations, and Access handles databases. Access is a
relational DBMS: it stores its data in tables, and it also holds the queries,
forms, and reports that work on that data. All of these objects are kept
inside a single database file for each database.

\begin{example}
A student is asked, ``Which MS Office program is used to store and manage
databases?'' The answer is \textbf{Access}. Excel is a spreadsheet program,
Word is a word processor, and PowerPoint is for presentations; none of them
is a DBMS.
\end{example}

\section{The objects in an Access database}
An Access database contains several kinds of object, each with one job.

\begin{center}
\begin{tabular}{p{0.16\textwidth}p{0.70\textwidth}}
\toprule
\textbf{Object} & \textbf{Job} \\
\midrule
Table & Stores the actual data in rows and columns \\
Query & Asks a question and retrieves selected data \\
Form & Provides a screen to enter and view data \\
Report & Presents data in a printable layout \\
\bottomrule
\end{tabular}
\end{center}

The data itself always lives in \textbf{tables}. The other objects work on
that data: a query retrieves it, a form is used to enter and view it, and a
report presents it for printing.

\section{Tables, records, and fields}
A \textbf{table} is the basic object where data is stored. Inside a table,
data is arranged in rows and columns. Each \textbf{column} is one
\textbf{field} --- one kind of information, such as Name or Marks. Each
\textbf{row} is one \textbf{record} --- one complete entry, such as the full
details of one student. A table therefore holds many records, and every
record has the same set of fields.

\begin{center}
\begin{tabular}{p{0.16\textwidth}p{0.34\textwidth}p{0.30\textwidth}}
\toprule
\textbf{Term} & \textbf{Meaning} & \textbf{Position in a table} \\
\midrule
Field & One kind of information & A column \\
Record & One complete entry & A row \\
Table & A collection of related records & The whole grid \\
\bottomrule
\end{tabular}
\end{center}

\begin{definitionbox}{Record}
One complete entry in a table, held in a single row and made up of one value
for each field.
\end{definitionbox}

\begin{definitionbox}{Field}
One kind of information in a table, held in a single column; each field has a
name and a data type.
\end{definitionbox}

\begin{example}
In a Students table the columns are RollNo, Name, Marks, and Passed. One row
--- for example, ``1, Ayaan, 78, Yes'' --- is one record. The Name column is
the Name field. The whole set of rows together is the table.
\end{example}

\section{Data types in Access}
Every field has a \textbf{data type}, which is the kind of value it can hold.
Common Access data types are:

\begin{center}
\begin{tabular}{p{0.22\textwidth}p{0.60\textwidth}}
\toprule
\textbf{Data type} & \textbf{Used for} \\
\midrule
Short Text & Names, codes, and other text \\
Number & Values used in calculations \\
Date/Time & Dates and times \\
Yes/No & Only True or False (a Boolean value) \\
AutoNumber & A unique number added automatically for each new record \\
\bottomrule
\end{tabular}
\end{center}

Because a field accepts only values that match its data type, the data stays
clean. A field set to Number will not accept ordinary text, and a Yes/No
field accepts only a true-or-false value.

\section{Field properties and validation}
A \textbf{field property} is a setting that describes how a field behaves.
Some important properties are \textbf{Field Size}, which limits how much a
field can hold; \textbf{Default Value}, which fills in a value
automatically; \textbf{Required}, which, when set to Yes, stops the field
from being left blank; and \textbf{Input Mask}, which fixes the pattern of
entry, such as the format of a phone number.

\textbf{Validation} means checking that entered data is sensible and allowed.
A \textbf{Validation Rule} states the condition a value must meet. For
example, the rule \texttt{Between 0 And 100} accepts only values in that
range. A \textbf{Validation Text} is the message shown when the rule is
violated.

\begin{definitionbox}{Validation}
The checking of entered data against a rule so that only sensible, allowed
values are stored in a field.
\end{definitionbox}

\begin{example}
On a Marks field, the Validation Rule \texttt{Between 0 And 100} is set. A
user types 150. Access rejects the entry because it fails the rule, and shows
the Validation Text as a message. A valid value such as 78 is accepted
normally.
\end{example}

\section{Primary key and index}
A \textbf{primary key} is a field that uniquely identifies each record. Its
value is different for every record and is never left blank. In the Students
table, RollNo can serve as the primary key because no two students share a
roll number. An AutoNumber field is often used as a ready-made primary key,
because Access fills it with a fresh unique number for every new record.

An \textbf{index} is a special structure that makes searching and sorting a
table faster. It works like the index of a book: instead of reading every
page, it points straight to the right place. A primary key is automatically
indexed. An index takes extra space and can slightly slow down the work of
adding data, so it is created only where it is genuinely useful.

\begin{definitionbox}{Primary key}
A field, or set of fields, whose value uniquely identifies every record in a
table; it cannot contain duplicates and cannot be left blank.
\end{definitionbox}

\section{Exam retrieval and common confusions}
Five distinctions prevent most errors on this topic.
\begin{enumerate}
  \item A \textbf{database} is the data; a \textbf{DBMS} is the software that
    manages it. Do not swap them.
  \item A \textbf{field} is a column; a \textbf{record} is a row; a
    \textbf{table} is the whole collection of related records.
  \item \textbf{Validation} checks the value entered; an \textbf{index}
    speeds up searching. They do different jobs.
  \item A \textbf{primary key} is unique and cannot be blank; it cannot
    contain duplicate values.
  \item Microsoft \textbf{Access} is a DBMS and part of the MS Office suite;
    it is not a spreadsheet.
\end{enumerate}

\begin{example}
An item asks which object stores the actual data in an Access database.
Trace the objects: a query retrieves data, a form is for entry, and a report
is for printing, but only the \textbf{table} stores the data. If the stem
instead asks which field uniquely identifies a record, the answer is the
\textbf{primary key}; if it asks which field setting stops a blank entry, the
answer is \textbf{Required}.
\end{example}

\subsection*{Exercises}
\begin{enumerate}[label=\arabic*.]
  \item Define a database and explain how it differs from a DBMS.
  \item Distinguish a field, a record, and a table, stating the position of
    each in a table.
  \item Name any four Access data types and give one use of each.
  \item Explain the difference between a validation rule and an index.
  \item What is a primary key, and why can it not contain duplicate values?
\end{enumerate}

\subsection*{Solutions}
\begingroup\small
\begin{enumerate}[label=\arabic*.]
  \item A database is an organised collection of related data. A DBMS is the
    software that manages that data --- it stores it, allows insert, update,
    delete, and retrieve operations, and controls access. The database is the
    data; the DBMS is the program.
  \item A field is one kind of information and is a column in a table. A
    record is one complete entry and is a row in a table. A table is the
    whole collection of related records, arranged in rows and columns.
  \item Any four of: Short Text (names and codes), Number (values used in
    calculations), Date/Time (dates and times), Yes/No (True or False),
    AutoNumber (a unique number for each new record).
  \item A validation rule checks that entered data meets a condition, such as
    \texttt{Between 0 And 100}, and rejects values that fail it. An index is a
    structure that speeds up searching and sorting; it does not check values.
  \item A primary key is a field whose value uniquely identifies each record.
    It cannot contain duplicates because two records with the same key value
    could not be told apart, and uniqueness is the whole purpose of the key.
\end{enumerate}
\endgroup

\lessonheading{29}{Access Keys and Relationships}
\section{Why keys and relationships matter}
A database keeps data in \textbf{tables}. A table is a grid: each
\textbf{record} is one row, and each \textbf{field} is one column that holds a
single kind of value. When a table holds many records, values such as a name,
a city, or a date can repeat, so a row cannot be recognised by those values
alone. A database therefore gives each row a reliable \textbf{identifier},
called a key, and it links rows that belong together using
\textbf{relationships}. Keys provide identity; relationships join tables that
share data; referential integrity keeps those joins correct.

\begin{definitionbox}{Table, record, field}
A \textbf{table} stores data about one subject as rows and columns. A
\textbf{record} is one complete row. A \textbf{field} is one column, holding
one attribute for every record.
\end{definitionbox}

\begin{center}
\begin{tabular}{p{0.28\textwidth}p{0.60\textwidth}}
\toprule
\textbf{Term} & \textbf{Meaning in this lesson} \\
\midrule
Primary key (PK) & A field or set of fields that uniquely identifies each record. \\
Foreign key (FK) & A field that refers to the primary key of another table. \\
Composite key & A primary key made of two or more fields together. \\
Candidate key & Any field that could serve as the primary key. \\
Alternate key & A candidate key that was not chosen as the primary key. \\
Junction table & A table that links two tables in a many-to-many relationship. \\
Referential integrity & The rule that a foreign-key value must match an existing primary-key value. \\
\bottomrule
\end{tabular}
\end{center}

The rest of the lesson uses these terms in one consistent sense. A
\textbf{primary key} identifies a record inside its own table. A
\textbf{foreign key} points from one table to another table's primary key. A
\textbf{composite key} is simply a primary key built from more than one field.
A \textbf{junction table} is the third table that implements a many-to-many
relationship.

\section{The primary key}
A \textbf{primary key (PK)} is a field, or a set of fields, that
\textbf{uniquely identifies} each record in a table. It has three firm rules.
First, its value must be \textbf{unique}: no two records may share the same
primary-key value. Second, it \textbf{cannot be left blank} (it cannot contain
a null value), because a blank value cannot identify anything. Third, each
table has \textbf{at most one} primary key, although that one key may be built
from several fields.

\begin{definitionbox}{Primary key}
A field (or set of fields) whose value is unique for every record and is never
blank, so it can identify each record in the table.
\end{definitionbox}

\begin{example}
In an Employee table, the field \texttt{EmployeeID} is a good primary key: every
employee receives a different number, and no employee is stored without one.
The employee's name would be a poor choice, because two employees can share a
name.
\end{example}

Access can generate a primary key automatically with the \textbf{AutoNumber}
data type, which inserts a unique number for each new record. A primary key is
what makes a table the \textbf{``one'' side} of a one-to-many relationship.

\section{Candidate keys, alternate keys and composite keys}
Any field that could serve as the primary key is called a \textbf{candidate
key}. When a table has more than one candidate key, only one is selected as
the primary key; the remaining candidates are called \textbf{alternate keys}.
For example, a Student table might have both \texttt{RollNumber} and
\texttt{Email} as candidate keys; if \texttt{RollNumber} is chosen, then
\texttt{Email} becomes an alternate key.

A \textbf{composite key} (also called a \textbf{compound key}) is a primary key
made of \textbf{two or more fields used together}. The combination is unique
even when each field on its own is not. Composite keys appear most often in
junction tables, where two foreign keys together form the primary key.

\begin{definitionbox}{Composite key}
A primary key that consists of two or more fields, whose combined value is
unique for every record.
\end{definitionbox}

\begin{example}
An OrderDetails table records each product line within an order. Neither
\texttt{OrderID} nor \texttt{ProductID} is unique by itself, but the pair
(\texttt{OrderID}~+~\texttt{ProductID}) is unique for every line item, so the
pair can be the primary key. This is a composite key.
\end{example}

A good primary key is also \textbf{stable} and \textbf{meaningless}: it should
not need to change later. This is why Access often uses an AutoNumber field
rather than a value such as a name or a phone number, which a person may edit.

\begin{center}
\begin{tabular}{p{0.22\textwidth}p{0.62\textwidth}}
\toprule
\textbf{Key} & \textbf{What it does} \\
\midrule
Primary key & Uniquely identifies a record; unique and not blank; at most one per table. \\
Candidate key & Any field eligible to become the primary key. \\
Alternate key & A candidate key that was not chosen as the primary key. \\
Composite key & A primary key made of two or more fields together. \\
Foreign key & A field that refers to another table's primary key; may repeat. \\
\bottomrule
\end{tabular}
\end{center}

\section{The foreign key}
A \textbf{foreign key (FK)} is a field in one table that refers to the
\textbf{primary key} of another table. The foreign key is the link that
creates a relationship between the two tables. Unlike a primary key, a foreign
key \textbf{may repeat} and \textbf{may be blank}; it is not unique. The table
that holds the foreign key is the \textbf{child} (or related) table, and the
table whose primary key is referenced is the \textbf{parent} table.

\begin{definitionbox}{Foreign key}
A field in one table whose values must match the primary-key values of another
table, used to link the two tables.
\end{definitionbox}

\begin{example}
In an Employee table, the field \texttt{DepartmentID} is a foreign key that
refers to the \texttt{DepartmentID} primary key in the Department table. Many
employees can share the same \texttt{DepartmentID}, so the foreign key repeats;
that is expected and correct.
\end{example}

\begin{center}
\begin{tabular}{p{0.24\textwidth}p{0.32\textwidth}p{0.32\textwidth}}
\toprule
\textbf{Point} & \textbf{Primary key} & \textbf{Foreign key} \\
\midrule
Purpose & Uniquely identifies a record & Links to another table \\
Uniqueness & Must be unique & May repeat \\
Blank allowed & No & Yes \\
Count per table & At most one & Can be more than one \\
\bottomrule
\end{tabular}
\end{center}

\section{Types of relationships}
A \textbf{relationship} works between two tables that share a common field ---
usually a primary key in one table and a foreign key in the other. Access
supports three relationship types: \textbf{one-to-one},
\textbf{one-to-many}, and \textbf{many-to-many}. Relationships are created and
edited in the \textbf{Relationships} window.

\begin{center}
\begin{tabular}{p{0.20\textwidth}p{0.38\textwidth}p{0.32\textwidth}}
\toprule
\textbf{Type} & \textbf{Meaning} & \textbf{How it is built} \\
\midrule
One-to-one & Each record in one table matches at most one in the other & A shared field that is unique in both tables \\
One-to-many & One record matches many records in the other table & Primary key on the ``one'' side, foreign key on the ``many'' side (most common) \\
Many-to-many & Many records match many records & Two one-to-many relationships through a junction table \\
\bottomrule
\end{tabular}
\end{center}

The rule to hold on to is simple: the \textbf{``one'' side} always contributes
the \textbf{primary key}, and the \textbf{``many'' side} always carries the
\textbf{foreign key}. A one-to-one relationship is just a special case in
which both sides are unique.

\subsection{One-to-one}
In a \textbf{one-to-one} relationship, each record in the first table matches
at most one record in the second table, and vice versa. It is used to split a
table with many fields into smaller parts, or to keep sensitive data in a
separate table that can be secured on its own. It is the \textbf{least common}
of the three types.

\begin{example}
An Employee table and an EmployeeConfidential table are linked one-to-one:
each employee has exactly one confidential record, and each confidential
record belongs to exactly one employee.
\end{example}

\subsection{One-to-many}
In a \textbf{one-to-many} relationship, one record in the first table matches
\textbf{many} records in the second table. The ``one'' side holds the
\textbf{primary key}, and the ``many'' side holds the \textbf{foreign key}.
This is the \textbf{most common} relationship type used in Access.

\begin{example}
One Department has many Employees, but each Employee belongs to only one
Department. The Department table is the ``one'' side and the Employee table is
the ``many'' side, linked by \texttt{DepartmentID}.
\end{example}

\subsection{Many-to-many and the junction table}
In a \textbf{many-to-many} relationship, many records in the first table match
many records in the second. Access \textbf{cannot} store a many-to-many
relationship directly. Instead it is split into \textbf{two one-to-many}
relationships by adding a third table called a \textbf{junction table} (also
known as an \textbf{associative} or \textbf{linking} table). The junction
table contains a foreign key for each of the two other tables.

\begin{definitionbox}{Junction table}
A table that connects two tables in a many-to-many relationship, holding a
foreign key for each of them.
\end{definitionbox}

A junction table's own primary key is usually a \textbf{composite key} formed
from the two foreign keys. It may also carry extra fields that describe the
link itself, such as a quantity, a grade, or a date.

\begin{example}
Students and Courses have a many-to-many relationship: a student takes many
courses, and a course has many students. Access implements this with an
Enrollments junction table that holds \texttt{StudentID} and \texttt{CourseID}
as its composite primary key, together with a \texttt{Grade} field.
\end{example}

\section{Creating a relationship in Access}
In the Access interface, relationships are built and edited in a small number
of steps. On the \textbf{Database Tools} tab, the \textbf{Relationships} button
opens the \textbf{Relationships} window. The tables to be linked are added to
this window, and then the \textbf{primary-key field} of one table is
\textbf{dragged onto} the matching foreign-key field of the other table.

\begin{enumerate}
  \item Open the \textbf{Relationships} window from the \textbf{Database Tools}
    tab.
  \item Add the two tables that will be linked.
  \item Drag the primary key of the ``one'' table onto the foreign key of the
    ``many'' table.
  \item In the \textbf{Edit Relationships} dialog box, choose
    \textbf{Enforce Referential Integrity}.
  \item Optionally tick \textbf{Cascade Update} and \textbf{Cascade Delete},
    then confirm.
\end{enumerate}

\begin{example}
To link Departments and Employees, drag \texttt{DepartmentID} from the
Department (parent) table onto \texttt{DepartmentID} in the Employee (child)
table, then enforce referential integrity. Access now refuses to store an
employee whose \texttt{DepartmentID} does not exist in the Department table.
\end{example}

\section{Referential integrity}
\textbf{Referential integrity} is the rule that a foreign-key value must
\textbf{match an existing primary-key value} in the related table (or be
blank). It prevents \textbf{orphan records}, which are child rows that point to
a parent row that does not exist. In Access, referential integrity is switched
on in the \textbf{Edit Relationships} dialog box (opened from the
\textbf{Relationships} window).

\begin{definitionbox}{Referential integrity}
A rule that keeps the foreign key in a child table consistent with the primary
key in the parent table: a child value must match an existing parent value or
be null.
\end{definitionbox}

When referential integrity is enforced, Access offers two related options.
\textbf{Cascade Update} means that if the parent primary-key value changes,
Access automatically updates the matching foreign-key values in the child
table. \textbf{Cascade Delete} means that if a parent record is deleted, Access
automatically deletes the matching child records as well.

\begin{example}
If a Department record is deleted and Cascade Delete is on, all Employees that
belong to that Department are deleted too. Without referential integrity,
Access would refuse to delete the Department while those Employees still point
to it.
\end{example}

\section{Access file formats: \texttt{.mdb} vs \texttt{.accdb}}
The file \textbf{extension} shows which Access format a database uses.
\texttt{.mdb} (\textit{Microsoft Database}) is the \textbf{older} Access file
format, used by Access 2003 and earlier versions. \texttt{.accdb}
(\textit{Access Database}) is the \textbf{newer, default} format introduced
with Access 2007 and used since. Access 2007 and later can still open an old
\texttt{.mdb} file, but new databases are saved as \texttt{.accdb} by default.

\begin{center}
\begin{tabular}{p{0.18\textwidth}p{0.30\textwidth}p{0.40\textwidth}}
\toprule
\textbf{Extension} & \textbf{Access version} & \textbf{Notes} \\
\midrule
\texttt{.mdb} & Access 2003 and earlier & The older format; supports user-level security. \\
\texttt{.accdb} & Access 2007 and later & The newer default; supports attachment, multi-valued and calculated fields. \\
\bottomrule
\end{tabular}
\end{center}

The newer \texttt{.accdb} format added \textbf{attachment} fields,
\textbf{multi-valued} fields, and \textbf{calculated} fields, and it integrates
better with SharePoint and Outlook. Keep the pairing straight: \texttt{.mdb}
is the \emph{old} format and \texttt{.accdb} is the \emph{new} format.
Reversing ``old'' and ``new'' is a common exam trap.

\section{Exam retrieval and common confusions}
Six distinctions clear up most errors on this topic.
\begin{enumerate}
  \item A primary key is \textbf{unique} and \textbf{not blank}; a foreign key
    \textbf{may repeat} and \textbf{may be blank}.
  \item A table has \textbf{at most one} primary key, but it can hold
    \textbf{several} foreign keys.
  \item A foreign key refers to the \textbf{primary key of another table}, not
    to its own table's primary key.
  \item \textbf{One-to-many} is the most common relationship; a
    \textbf{many-to-many} relationship is never stored directly but is built
    with a \textbf{junction table}.
  \item A junction table commonly has a \textbf{composite primary key} made
    from the two foreign keys.
  \item \texttt{.mdb} is the \textbf{older} Access format; \texttt{.accdb} is
    the \textbf{newer} format.
\end{enumerate}

\begin{example}
An item asks which key may contain duplicate values. Walk through the rules: a
primary key must be unique and not blank, so it cannot repeat; a foreign key
may repeat and may be blank; a composite key is still a primary key and so is
unique as a combination. The answer is the foreign key.
\end{example}

\begin{example}
An item states that ``a many-to-many relationship is created by adding a
foreign key to one of the two tables.'' This is wrong. A many-to-many
relationship cannot be stored in either table directly; a \textbf{junction
table} is added, and two one-to-many relationships are built through it. The
wrong option has skipped the junction table.
\end{example}

\subsection*{Quick self-check}
Before the exercise set, confirm you can state each of these in one line.
\begin{itemize}
  \item A primary key is unique, never blank, and there is at most one per
    table.
  \item A foreign key links to another table's primary key and may repeat.
  \item A composite key is a primary key of two or more fields.
  \item One-to-many is the most common relationship; the ``one'' side supplies
    the primary key and the ``many'' side the foreign key.
  \item A many-to-many relationship is always built with a junction table.
  \item Referential integrity keeps foreign-key values matching primary-key
    values and offers Cascade Update and Cascade Delete.
  \item \texttt{.mdb} is the older format; \texttt{.accdb} is the newer one.
\end{itemize}

\subsection*{Exercises}
\begin{enumerate}[label=\arabic*.]
  \item Define a primary key and state its two main rules about uniqueness and
    blank values.
  \item Explain the difference between a primary key and a foreign key.
  \item What is a composite key? Give one example where it is used.
  \item Name the three relationship types in Access and say which is the most
    common.
  \item Why does a many-to-many relationship need a junction table?
  \item State the rule of referential integrity and name its two cascade
    options.
  \item Distinguish the Access file formats \texttt{.mdb} and \texttt{.accdb}.
\end{enumerate}

\subsection*{Solutions}
\begingroup\small
\begin{enumerate}[label=\arabic*.]
  \item A primary key is a field (or set of fields) that uniquely identifies
    each record. Its value must be unique --- no two records may share it ---
    and it cannot be left blank (no null value). A table has at most one
    primary key.
  \item A primary key uniquely identifies records inside its own table and
    cannot repeat or be blank. A foreign key is a field that refers to another
    table's primary key; it may repeat and may be blank, and it creates the
    relationship between the two tables.
  \item A composite key is a primary key made of two or more fields used
    together, where the combination is unique. It is used in a junction table
    --- for example, an Enrollments table whose primary key is
    (\texttt{StudentID}~+~\texttt{CourseID}).
  \item The three types are one-to-one, one-to-many, and many-to-many. The
    \textbf{one-to-many} relationship is the most common.
  \item Access cannot store a many-to-many relationship directly. It is split
    into two one-to-many relationships by a junction table that holds a
    foreign key for each of the two tables.
  \item Referential integrity means a foreign-key value must match an existing
    primary-key value in the related table (or be blank), which prevents
    orphan records. Its two options are \textbf{Cascade Update} and
    \textbf{Cascade Delete}.
  \item \texttt{.mdb} is the older Access format used by Access 2003 and
    earlier. \texttt{.accdb} is the newer, default format used from Access
    2007 onward, adding attachment, multi-valued and calculated fields.
\end{enumerate}
\endgroup

\lessonheading{30}{Access Queries, Forms and Reports}
\section{The Access objects and the object sequence}
An Access database is built from a small set of \textbf{objects}. The
\textbf{table} stores the actual data. The \textbf{query} asks a question
and retrieves selected data. The \textbf{form} provides a screen for
entering and viewing data. The \textbf{report} presents data in a formatted,
printable layout. These four objects are used together in a natural order:
data is stored in a table, selected by a query, entered through a form, and
printed by a report. That order --- \textbf{Table} $\rightarrow$
\textbf{Query} $\rightarrow$ \textbf{Form} $\rightarrow$ \textbf{Report}
--- is the object sequence this lesson is built around.

\begin{definitionbox}{Database object}
A named component of an Access database that performs one job. The main
objects are the table, query, form and report.
\end{definitionbox}

\begin{center}
\begin{tabular}{p{0.16\textwidth}p{0.68\textwidth}}
\toprule
\textbf{Object} & \textbf{Job} \\
\midrule
Table & Stores the actual data in rows and columns. \\
Query & Asks a question and retrieves selected data. \\
Form & Provides an on-screen interface to enter and view data. \\
Report & Presents data in a formatted, printable layout. \\
\bottomrule
\end{tabular}
\end{center}

The rest of this lesson uses these terms in one consistent sense. A
\textbf{table} is the store. A \textbf{query} is the question. A
\textbf{form} is the on-screen entry screen. A \textbf{report} is the printed
output.

\section{The table is the source of data}
A table stores data in \textbf{fields} (columns) and \textbf{records}
(rows). Every other object draws its data from a table or from a query. A
form or report built on a table shows that table's data; built on a query,
it shows only the rows the query selected. No data is stored in a query, a
form or a report --- those objects only \emph{present} what the table holds.
This is why the table is called the foundation of the database.

\section{What a query is}
A \textbf{query} is a question asked of the database. It retrieves data from
one or more tables, usually by a condition. A query does not normally store
new data of its own; it shows a \emph{result set} drawn from the tables each
time it is opened. Queries can also be used to \emph{change} data; those are
called \textbf{action queries}, and they are discussed below.

\begin{definitionbox}{Query}
An Access object that selects, filters, or acts on records according to the
rules the user sets, drawing its data from one or more tables.
\end{definitionbox}

\section{The Select query: retrieving on a condition}
A \textbf{Select query} retrieves the records that satisfy a
\textbf{condition} (also called a \emph{criterion}). It is the most common
query type in Access. The condition decides \emph{which rows appear}, and
the chosen fields decide \emph{which columns appear}. For example, a Select
query on a Students table could show only those students whose
\texttt{Marks} are greater than 60.

\begin{example}
An item asks: ``What is a query that retrieves data on the basis of a
condition called?'' The query is only \emph{retrieving} data --- it is not
changing it --- so it is a \textbf{Select query}. This is the concept tested
by PYQ-39 (Website Operator 2026, query retrieving data on a condition).
\end{example}

\section{Building a query: Design view and criteria}
Queries are usually built in \textbf{Query Design} view, or typed directly in
\textbf{SQL} view. The design grid is called \textbf{QBE} ---
\textbf{Query By Example}. Each column of the grid holds one \textbf{field},
and the \textbf{Criteria} row holds the condition for that field. A criterion
of \texttt{>60} on the \texttt{Marks} field, for instance, keeps only the
rows in which marks are above 60.

\begin{definitionbox}{QBE (Query By Example)}
The grid-based design view in Access in which a query is built by choosing
fields and typing criteria, without writing the SQL by hand.
\end{definitionbox}

\begin{example}
For a Students table, place \texttt{Marks} in the field row and \texttt{>60}
in the Criteria row. When the query runs, Access returns only the students
who scored more than 60.
\end{example}

\section{Action queries: Update, Append and Delete}
An \textbf{action query} changes the stored data. The three common types are:
\begin{itemize}
  \item An \textbf{Update query} changes the existing values in records.
  \item An \textbf{Append query} adds records from one table to another.
  \item A \textbf{Delete query} removes the records that match a condition.
\end{itemize}
Because these queries alter the table itself, they are treated with care in
Access. The distinction between retrieving data and changing data is the key
one: a Select query only reads, while Update, Append and Delete queries
write.

\begin{center}
\begin{tabular}{p{0.18\textwidth}p{0.68\textwidth}}
\toprule
\textbf{Query} & \textbf{What it does} \\
\midrule
Select & Retrieves records that satisfy a condition; does not change data. \\
Update & Changes existing values in records. \\
Append & Adds records from one table to another. \\
Delete & Removes records that match a condition. \\
\bottomrule
\end{tabular}
\end{center}

\section{Select versus action queries}
The nearest confusion in this topic is between a Select query and the action
queries. A \textbf{Select query} \textbf{retrieves} data on a condition and
leaves the tables unchanged. An \textbf{action query} (Update, Append,
Delete) \textbf{modifies} the stored data. If a question asks which query
\emph{changes} records, the answer is an action query; if it asks which query
\emph{shows} records that meet a condition, the answer is a Select query.
This is the distinction recorded as TRAP-31.

\begin{example}
An item lists Select, Update, Append and Delete queries and asks which one
does \textbf{not} modify the data. Select is the only one that leaves the
tables unchanged, so Select is the answer.
\end{example}

\section{Lookup fields}
A \textbf{lookup} field lets the user \emph{pick a value from a list} instead
of typing it. The list usually comes from a table or a query --- for example,
a list of department names from a Departments table. A lookup is often used
for a \textbf{foreign key}, so that the value entered matches an existing
value in the related table. The benefit is fewer typing errors and more
consistent data. A lookup field only \emph{chooses} a value from a list; it
does not create a new table or a new query.

\begin{definitionbox}{Lookup field}
A field that displays a list of values, drawn from a table or query, so the
user can select one instead of typing it.
\end{definitionbox}

\section{Basic SQL awareness}
\textbf{SQL} stands for \textbf{Structured Query Language}. It is the
standard language for working with relational databases, and Access runs SQL
behind the scenes whenever a query is run. A retrieval statement has the
general shape \texttt{SELECT fields FROM table WHERE condition;}. The
\texttt{SELECT} keyword names the fields to show, \texttt{FROM} names the
table, and \texttt{WHERE} gives the condition. To retrieve only the students
with marks above 60, one would write
\texttt{SELECT Name FROM Students WHERE Marks > 60;}.

\begin{definitionbox}{SQL}
Structured Query Language --- the standard language used to define and
retrieve data in a relational database.
\end{definitionbox}

\section{Forms: entering and viewing data}
A \textbf{form} is an object that provides an on-screen interface for
entering, viewing and editing data. A form is usually based on a table or a
query. It shows one record at a time (or a few), with clearly labelled
fields, and makes data entry easier and safer than typing directly into a
table grid. Where the table is the raw store, the form is the friendly
screen the user actually works with.

\section{Reports: printable output}
A \textbf{report} presents data in a formatted, printable layout. It can
group records and show totals and summaries. A report is \textbf{read-only}:
it is meant for printing and sharing, not for editing data. Like a form, a
report can be based on a table or a query. The report is the output object of
the database.

\begin{example}
A school can keep all student marks in a table, use a Select query to pick
the students who passed, enter new students through a form, and print the
pass list as a report. The same data moves, in order, through the
Table~$\rightarrow$~Query~$\rightarrow$~Form~$\rightarrow$~Report sequence.
\end{example}

\section{Forms versus reports}
The nearest confusion between the two presentation objects is worth stating
plainly. A \textbf{form} is used \emph{on screen} to enter and view data, and
it allows editing. A \textbf{report} is \emph{read-only} and is used to
\emph{print} data and present summaries. A form answers the question ``how do
I put data in and look at it?''; a report answers ``how do I print it?''

\begin{center}
\begin{tabular}{p{0.18\textwidth}p{0.36\textwidth}p{0.34\textwidth}}
\toprule
\textbf{Point} & \textbf{Form} & \textbf{Report} \\
\midrule
Purpose & Enter and view data on screen & Print and present data \\
Data change & Allows editing & Read-only \\
Typical use & Data-entry screen & Printed summary \\
\bottomrule
\end{tabular}
\end{center}

\section{Exam retrieval and common confusions}
Six distinctions clear up most errors on this topic.
\begin{enumerate}
  \item A \textbf{Select} query \textbf{retrieves} data on a condition;
    \textbf{Update}, \textbf{Append} and \textbf{Delete} queries
    \textbf{modify} data (TRAP-31).
  \item Data is stored in a \textbf{table}, \textbf{not} in a query, form or
    report.
  \item A \textbf{form} is for on-screen entry and viewing; a \textbf{report}
    is a printable, read-only output.
  \item The object sequence is \textbf{Table} $\rightarrow$ \textbf{Query}
    $\rightarrow$ \textbf{Form} $\rightarrow$ \textbf{Report}.
  \item \textbf{SQL} stands for \textbf{Structured Query Language};
    \texttt{SELECT ... FROM ... WHERE} retrieves data on a condition.
  \item A \textbf{lookup} field picks a value from a list supplied by a table
    or query.
\end{enumerate}

\begin{example}
An item states, ``A report is used to enter new records into the
database.'' This is wrong. Entering records is the job of a \textbf{form}; a
report is \textbf{read-only} and is used for printing. The wrong option has
swapped the form and the report.
\end{example}

\subsection*{Exercises}
\begin{enumerate}[label=\arabic*.]
  \item Name the four main objects of an Access database and state the job of
    each.
  \item Give the correct object sequence for storing, selecting, entering and
    printing data.
  \item What is a Select query, and how does it differ from an Update query?
  \item Name the three common action queries and say what each one does.
  \item What is a lookup field, and why is it useful?
  \item Expand SQL and write the general form of a retrieval statement.
  \item Distinguish a form from a report.
\end{enumerate}

\subsection*{Solutions}
\begingroup\small
\begin{enumerate}[label=\arabic*.]
  \item \textbf{Table} stores the data; \textbf{Query} asks a question and
    retrieves selected data; \textbf{Form} provides an on-screen interface to
    enter and view data; \textbf{Report} presents data in a formatted,
    printable layout.
  \item \textbf{Table} $\rightarrow$ \textbf{Query} $\rightarrow$
    \textbf{Form} $\rightarrow$ \textbf{Report}. The table stores the data,
    the query selects it, the form enters it, and the report prints it.
  \item A Select query \textbf{retrieves} records that satisfy a condition
    and leaves the tables unchanged. An Update query \textbf{changes} the
    existing values in records, so it modifies the stored data.
  \item \textbf{Update} changes existing values in records;
    \textbf{Append} adds records from one table to another;
    \textbf{Delete} removes records that match a condition.
  \item A lookup field displays a list of values, drawn from a table or
    query, so the user selects one instead of typing it. It reduces typing
    errors and keeps values consistent, and it is often used for a foreign
    key.
  \item SQL is \textbf{Structured Query Language}. A retrieval statement has
    the form \texttt{SELECT fields FROM table WHERE condition;}, where
    \texttt{SELECT} names the fields, \texttt{FROM} names the table, and
    \texttt{WHERE} gives the condition.
  \item A \textbf{form} is an on-screen interface used to enter and view
    data, and it allows editing. A \textbf{report} is a read-only, formatted
    output used for printing and summarising data.
\end{enumerate}
\endgroup

\lessonheading{31}{PowerPoint Basics}
\section{What PowerPoint is}
MS PowerPoint is the presentation program of the MS Office suite. A
\textbf{presentation} is a file made of one or more pages, and each page is
called a \textbf{slide}. The slides are shown one after another as a
\textbf{slide show} on a screen or projector. The official syllabus writes
this unit as ``Powerpoint Presentation'', and the usual short form is
\textbf{PPT}. A single slide can hold text, pictures, charts, tables, audio,
and video.

\begin{definitionbox}{Presentation, slide, slide show}
A presentation is the whole file; a slide is one page of it; a slide show is
the full-screen sequential display of those slides to an audience.
\end{definitionbox}

\begin{center}
\begin{tabular}{p{0.24\textwidth}p{0.66\textwidth}}
\toprule
\textbf{Term} & \textbf{Meaning in this lesson} \\
\midrule
Presentation & The whole PowerPoint file, made of slides. \\
Slide & One page of the presentation. \\
Placeholder & The dotted-line box on a slide that holds content. \\
Layout & The arrangement of placeholders on a slide. \\
Theme & A coordinated set of colours, fonts, and effects. \\
Template & A ready-made starter file with a theme, layouts, and sample content. \\
Slide Master & The top-level slide that controls all linked layouts and slides. \\
Notes & The speaker's reminders below a slide, unseen by the audience. \\
Handout & A printed sheet carrying several slides for the audience. \\
\bottomrule
\end{tabular}
\end{center}

The rest of this lesson uses these terms in one consistent sense. A
\textbf{slide} is the page, a \textbf{placeholder} is the box that holds
content on that page, and a \textbf{layout} decides where those boxes sit. A
\textbf{theme} restyles the file, a \textbf{template} starts a new file from a
ready design, and the \textbf{Slide Master} controls the design of every
linked slide at once.

\section{Full forms and file extensions}
Three full forms carry direct-recall marks. \textbf{PPT} stands for PowerPoint
Presentation. \textbf{PPTX} stands for PowerPoint XML Presentation and is the
default format since Office 2007. \textbf{PPSX} stands for PowerPoint Slide
Show. Two more extensions complete the family: \textbf{POTX} is the PowerPoint
Template, and \textbf{PPTM} is the macro-enabled presentation.

\begin{definitionbox}{PowerPoint file extensions}
The editable working file uses \texttt{.pptx}; the template uses
\texttt{.potx}; the show file that opens directly as a slide show uses
\texttt{.ppsx}; the macro-enabled file uses \texttt{.pptm}. Presentations can
also be saved or exported as \texttt{.pdf}.
\end{definitionbox}

The working distinction is simple. Double-clicking a \texttt{.pptx} file opens
it for editing, while double-clicking a \texttt{.ppsx} file launches the
slide show directly. A \texttt{.potx} file is never the finished talk; it is
the reusable pattern from which new talks are started. Confusing these three
extensions is one of the most common wrong answers in this unit.

\section{Slides, placeholders, and layouts}
A slide is a single page of a presentation. Content is not typed freely onto
the page background; it goes into \textbf{placeholders}, the dotted-line boxes
marked Title, Text, Picture, Chart, or Table. Clicking inside a placeholder
allows typing, and the small icons inside a content placeholder insert a
picture, chart, table, or SmartArt graphic directly.

\begin{definitionbox}{Placeholder}
A dotted-line container on a slide that holds one kind of content (title,
text, picture, chart, table, or media).
\end{definitionbox}

A \textbf{layout} is the arrangement of those placeholders. The common layouts
are Title Slide, Title and Content, Section Header, Two Content, Comparison,
Title Only, and Blank. A new slide is added from the Home tab through New
Slide, which also offers the layout gallery; the layout of an existing slide
is changed through Home tab to Layout. The Blank layout carries no
placeholders at all, so every object on it must be inserted by hand.

\begin{example}
A student builds a five-slide class talk on the office desktop. The first
slide uses the Title Slide layout (one title plus one subtitle placeholder);
each following point uses Title and Content; the closing slide reuses Section
Header. Changing the layout of slide 3 from Title and Content to Two Content
rearranges only slide 3; nothing else in the file moves.
\end{example}

\section{Themes, templates, and the Slide Master}
A \textbf{theme} is a coordinated set of colours, fonts, and effects applied
from the Design tab. One click restyles every slide in the file, because the
theme is a design rule, not content. A \textbf{template} is a ready-made
pattern file (extension \texttt{.potx}) that bundles a theme with a set of
layouts and sometimes sample text. Starting a new presentation from a template
reuses that whole design instead of building it again.

\begin{definitionbox}{Theme and template}
A theme is a coordinated set of colours, fonts, and effects. A template is a
reusable starter file (\texttt{.potx}) carrying a theme plus layouts and
sample content.
\end{definitionbox}

The \textbf{Slide Master} is the top-level slide that controls the design of
all layouts and slides linked to it. Editing the font, the logo position, or
the footer on the Slide Master changes every linked slide at once, which is
far faster than editing forty slides one by one. It is opened from the View
tab through Slide Master and left through Close Master View. PowerPoint has
three masters in total: the Slide Master for on-screen slides, the
\textbf{Handout Master} for printed handout layouts, and the \textbf{Notes
Master} for printed notes pages.

\begin{example}
An item states, ``The Handout Master controls the design of all on-screen
slides.'' This is wrong. The Slide Master controls the on-screen slides; the
Handout Master controls only the printed handout layouts. The wrong option has
swapped two of the three masters.
\end{example}

\section{Views: Normal, Outline, Sorter, Reading, Show, Notes Page}
PowerPoint offers several views, each with one job. \textbf{Normal view} is
the default editing view, with three areas: the large slide pane in the
centre, the slides or outline pane with thumbnails on the left, and the notes
pane below. All typing, formatting, and inserting happen here. \textbf{Outline
View} shows only the text outline of every slide, which suits restructuring
long talks.

\textbf{Slide Sorter view} shows thumbnails of all slides together, so slides
can be reordered by drag and drop and checked, deleted, or duplicated at a
glance. \textbf{Reading View} shows the show almost full-screen for reading on
the author's own screen, with simple navigation but without editing tools.
\textbf{Slide Show view} is the real full-screen presentation to an audience.
\textbf{Notes Page view} shows one slide above its speaker-notes text, one
page per slide, exactly as it will print.

\begin{definitionbox}{Normal view}
The default editing view of PowerPoint, combining the slide pane, the
slides or outline pane, and the notes pane.
\end{definitionbox}

\begin{example}
A student has twelve slides in the wrong order. The fastest fix is Slide
Sorter view: drag the thumbnails into the right sequence. Opening each slide
in Normal view and cutting and pasting would work but wastes time; choosing
Reading View would not help at all, because Reading View cannot reorder
anything.
\end{example}

\section{Notes, handouts, and printing}
\textbf{Notes} (speaker notes) are reminders typed in the notes pane below a
slide. The audience never sees them during the show; they exist so the speaker
does not forget a figure, a definition, or the next point. During the show the
speaker reads them only in Presenter View on the speaker's own screen. Notes
print through File to Print to Notes Pages.

\begin{definitionbox}{Speaker notes and handouts}
Speaker notes help the speaker and stay hidden from the audience. Handouts
help the audience and are printed sheets carrying several slides per page.
\end{definitionbox}

\textbf{Handouts} are those printed sheets: 1, 2, 3, 4, 6, or 9 slides per
page. The most examined choice is 6 slides per page, and the 3-per-page layout
adds ruled lines beside each slide so the audience can write. Their design is
controlled by the Handout Master. The full print path is File to Print, with
choices for Full Page Slides, Handouts, Notes Pages, and Outline View.

\section{Running and navigating the slide show}
Two shortcuts start the show. \textbf{F5} begins from the first slide, and
\textbf{Shift+F5} begins from the current slide. This pair is a confirmed past
item (WO 2026 Q80 tested F5 for ``from the beginning''). During the show, N,
Space, Right arrow, or Enter moves forward; P or Left arrow moves back; Esc
ends the show. Pressing B blanks the screen to black and pressing W blanks it
to white, which is useful while the speaker talks without a slide. The Slide
Show tab holds timing, rehearsal, and monitor settings.

A \textbf{transition} is the movement between one slide and the next (Fade,
Push, Wipe, Morph), set on the Transitions tab with one transition per slide.
An \textbf{animation} is the movement of an object on a slide, set on the
Animations tab: Entrance, Emphasis, Exit, and Motion Paths, managed in the
Animation Pane. SmartArt conversion behaviour (which content cannot convert
without re-entry) has also been examined (JA 2026 Q77), so candidates should
practise the Insert to SmartArt path on text and picture lists rather than
assume every object converts.

\section{Exam retrieval and common confusions}
Six distinctions prevent most errors on this topic.
\begin{enumerate}
  \item The \textbf{slide} is the page; the \textbf{placeholder} is the
    dotted box holding content; the \textbf{layout} arranges those boxes.
  \item \textbf{.pptx} is the editable file, \textbf{.ppsx} is the show file,
    and \textbf{.potx} is the template. Do not swap them.
  \item A \textbf{theme} restyles (colours, fonts, effects); a
    \textbf{template} is the reusable starter file.
  \item The \textbf{Slide Master} controls on-screen slides; the Handout and
    Notes Masters control print layouts.
  \item \textbf{Normal} is the default editing view; \textbf{Slide Sorter}
    reorders; Reading reviews; \textbf{notes} serve the speaker and
    \textbf{handouts} serve the audience.
  \item \textbf{F5} starts from the beginning and \textbf{Shift+F5} from the
    current slide; a \textbf{transition} moves between slides while an
    \textbf{animation} moves objects on a slide.
\end{enumerate}

\begin{example}
An item asks which view is the default editing view and offers Normal, Slide
Sorter, Reading, and Slide Show. Trace the definitions: Normal is the editing
view with the slide, thumbnail, and notes panes; Slide Sorter manages order;
Reading reviews; Slide Show presents. The answer is Normal view. If the stem
instead asks which view reorders slides by drag and drop, the answer changes
to Slide Sorter.
\end{example}

\subsection*{Exercises}
\begin{enumerate}[label=\arabic*.]
  \item Distinguish a slide, a placeholder, and a layout, giving one example
    of each.
  \item Expand PPT, PPTX, and PPSX, and state the file extension of a
    PowerPoint template.
  \item Distinguish a theme from a template.
  \item What does the Slide Master control, and how do the Handout and Notes
    Masters differ from it?
  \item List four PowerPoint views and state the job of each.
  \item Distinguish speaker notes from handouts, and state which viewers each
    serves.
  \item Which keys start a show from the beginning and from the current slide,
    and which key ends the show?
  \item Distinguish a transition from an animation.
\end{enumerate}

\subsection*{Solutions}
\begingroup\small
\begin{enumerate}[label=\arabic*.]
  \item A slide is one page of the presentation (for example, slide 2 of a
    class talk). A placeholder is the dotted-line box holding content on that
    page (for example, the Title placeholder). A layout is the arrangement of
    those boxes (for example, Title and Content).
  \item PPT is PowerPoint Presentation; PPTX is PowerPoint XML Presentation;
    PPSX is PowerPoint Slide Show. A PowerPoint template uses
    \texttt{.potx}.
  \item A theme is a coordinated set of colours, fonts, and effects that
    restyles a file; a template is a reusable starter file
    (\texttt{.potx}) that bundles a theme with layouts and sample content.
  \item The Slide Master controls the design of all linked on-screen slides
    and layouts. The Handout Master controls printed handout layouts and the
    Notes Master controls printed notes pages; neither restyles the on-screen
    slides.
  \item Normal is the default editing view; Slide Sorter shows thumbnails for
    reordering; Reading View shows the show for on-screen review; Slide Show
    presents full-screen; Notes Page shows one slide above its notes text.
    Any four with correct jobs earn full credit.
  \item Speaker notes are the speaker's hidden reminders below a slide; they
    serve the speaker. Handouts are printed multi-slide sheets; they serve
    the audience.
  \item F5 starts from the beginning; Shift+F5 starts from the current slide;
    Esc ends the show.
  \item A transition is the movement between one slide and the next; an
    animation is the movement of an object on a slide (Entrance, Emphasis,
    Exit, Motion Paths).
\end{enumerate}
\endgroup

\lessonheading{32}{PowerPoint Animation and Transition}
\section{Animation and transition: the core distinction}
PowerPoint can make a presentation move, but it uses two different tools for
two different kinds of movement. An \textbf{animation} is an effect applied to
one object or to text \emph{inside} a single slide. A \textbf{transition} is an
effect applied to the movement \emph{from one slide to the next}. The simplest
way to keep them apart is to ask where the movement happens: movement inside a
slide is animation; movement between slides is a transition.

Both kinds of effect are added from the PowerPoint ribbon. Animation is
managed on the \textbf{Animations} tab, and transitions are managed on the
\textbf{Transitions} tab. Neither tool changes the content of the slides; they
change only how that content appears and how the deck moves from slide to
slide.

\begin{definitionbox}{Animation}
An effect applied to a single object or to text on a slide --- for example
making a picture appear, or making a title grow. Animation acts within a
slide.
\end{definitionbox}

\begin{definitionbox}{Transition}
An effect applied to the change from one slide to the next, such as Fade,
Push, or Morph. A transition acts between slides.
\end{definitionbox}

\begin{center}
\begin{tabular}{p{0.22\textwidth}p{0.32\textwidth}p{0.34\textwidth}}
\toprule
\textbf{Point} & \textbf{Animation} & \textbf{Transition} \\
\midrule
Where it applies & An object or text on a slide & The change between two slides \\
Ribbon tab & Animations & Transitions \\
Nearest confusion & Transitions & Entrance/Exit effects \\
\bottomrule
\end{tabular}
\end{center}

\section{The four animation effects}
When you select an object and open the Animations tab, PowerPoint groups the
effects into four categories. Each category answers a different question about
the object.

\begin{center}
\begin{tabular}{p{0.20\textwidth}p{0.64\textwidth}}
\toprule
\textbf{Effect} & \textbf{What it does} \\
\midrule
Entrance & Makes an object \textbf{appear} on the slide. \\
Exit & Makes an object \textbf{leave} the slide. \\
Emphasis & Draws attention to an object that is \textbf{already visible},
staying on the slide. \\
Motion Path & Moves an object \textbf{along a defined path}, such as a line or
curve. \\
\bottomrule
\end{tabular}
\end{center}

\textbf{Entrance} effects are used to reveal content step by step: the title or
a bullet point appears when the presenter is ready. \textbf{Exit} effects are
the opposite --- they take an object off the slide once it is no longer needed.
An entrance and an exit are not the same as an emphasis: an emphasis effect
does not add or remove the object, it only makes the object stand out while it
remains on the slide. A \textbf{motion path} is different again: it changes the
object's position, moving it from one place to another along a drawn path.

\begin{example}
A presenter wants a logo to travel from the left edge to the corner of the
slide. That is a \textbf{Motion Path}, because the object's position changes
during the show. If instead the logo simply spins once and stops in the same
place, that is an \textbf{Emphasis} effect. If the logo is not on the slide at
all until clicked, that is an \textbf{Entrance} effect.
\end{example}

Animations can be set to start when the slide appears or when the presenter
clicks, and each effect has a duration. The order and timing of the effects on
one slide are controlled from the \textbf{Animation Pane}.

\section{Transitions and Morph}
A transition plays as the presentation moves from the current slide to the
next. Common transitions include \textbf{Fade}, \textbf{Push}, \textbf{Wipe},
and \textbf{Dissolve}. Each transition can be given a duration and an optional
sound. From the Transitions tab, the \textbf{Apply To All} button copies the
chosen transition to every slide in the deck.

Transitions are often grouped by the strength of their effect. \textbf{Subtle}
transitions, such as Fade and Wipe, are gentle and professional.
\textbf{Exciting} transitions, such as Push and Split, use stronger movement.
\textbf{Dynamic Content} transitions respond to the slide content itself.

\begin{definitionbox}{Morph}
A modern transition that smoothly moves and animates objects as one slide
changes into the next. Morph works best when two slides share similar objects,
so the objects appear to glide from their old position to their new one.
\end{definitionbox}

\begin{example}
Slide 1 shows a circle on the left side; slide 2 shows the same circle on the
right side. With \textbf{Morph} applied to slide 2, the circle appears to
slide smoothly across the screen instead of jumping. Because Morph is a
transition, it is applied from the \textbf{Transitions} tab, not the Animations
tab.
\end{example}

\section{Starting and controlling the slide show}
The slide show has a small set of shortcut keys that JKSSB papers ask about
directly. The most important is \textbf{F5}, which starts the slide show from
the beginning --- that is, from the first slide. \textbf{Shift+F5} starts the
slide show from the current slide --- the one selected. \textbf{Alt+F5} starts
the presentation in \textbf{Presenter View}. \textbf{Esc} ends a running slide
show.

\begin{center}
\begin{tabular}{p{0.22\textwidth}p{0.60\textwidth}}
\toprule
\textbf{Key} & \textbf{Action} \\
\midrule
F5 & Start the slide show from the \textbf{beginning}. \\
Shift+F5 & Start the slide show from the \textbf{current} slide. \\
Alt+F5 & Start the presentation in \textbf{Presenter View}. \\
Esc & \textbf{End} the slide show. \\
\bottomrule
\end{tabular}
\end{center}

Once the show is running, more controls are available. The keys \textbf{N},
\textbf{Enter}, \textbf{Space}, or the right arrow advance to the next slide or
animation; \textbf{P} or the left arrow go back to the previous slide.
Pressing \textbf{B} turns the screen black and \textbf{W} turns it white, which
is useful for pausing discussion. Typing a slide number and pressing
\textbf{Enter} jumps straight to that slide.

\begin{example}
An item asks, ``Which key starts the slide show from the beginning?'' Trace
the key set: F5 is the full show from the first slide, Shift+F5 is from the
current slide, Alt+F5 opens Presenter View, and F7 is the spelling and grammar
check. The answer is \textbf{F5}. A distractor that offers F7 is testing
whether you confuse the show shortcut with spell check.
\end{example}

\section{Speaker notes, views and handouts}
The \textbf{Notes pane} sits below the slide in Normal view and holds
\textbf{speaker notes} --- the reminders the presenter reads while presenting.
Notes are not shown on the projected slide during a normal slide show; in
\textbf{Presenter View} the notes are visible only to the presenter, along with
the current slide, the next slide, and a timer.

PowerPoint provides several views for different jobs. \textbf{Normal view} is
used to edit one slide, with an outline on the side and the notes pane below.
\textbf{Outline view} works with the text of all slides as an outline.
\textbf{Slide Sorter view} shows small \textbf{thumbnails} of every slide and
is the view for rearranging slides. \textbf{Notes Page view} shows a slide
together with its speaker notes. \textbf{Reading View} reviews the show in a
window without going full screen, and \textbf{Slide Show view} presents the
deck full screen.

A \textbf{handout} is the printed version of the presentation that is given to
the \textbf{audience}. Handouts usually place several slides on one page ---
two, three, four, six, or nine per page --- and some handout layouts include
lined space for the audience to write. The key contrast is simple: the notes
pane is for the presenter, and a handout is for the audience.

\begin{definitionbox}{Handout}
A printed or exported layout of the presentation for the audience, normally
showing several slides on a single page, often with room for notes.
\end{definitionbox}

\section{SmartArt conversion and its limits}
\textbf{SmartArt} is a graphic that arranges text into a designed layout such
as a list, process, or hierarchy. A SmartArt graphic can be converted into
\textbf{shapes} or into \textbf{text}. Converting to \textbf{text} turns the
graphic into a plain bulleted list, so the graphical layout is lost. Converting
to \textbf{shapes} keeps individual shapes but removes the SmartArt
arrangement, and once the conversion is done the graphic cannot be turned back
into SmartArt --- the content that depended on the layout must be
\textbf{re-entered} or re-applied by hand.

\begin{example}
A student has a SmartArt process diagram and converts it to text to reuse the
wording in a document. The words survive, but the arrows, boxes, and arranged
layout do not. To get the diagram back, the SmartArt must be created again,
because a text conversion has dropped the design. The JA 2026 paper tested this
kind of limit --- which content cannot convert cleanly without re-entry --- as
its SmartArt item (PYQ-17).
\end{example}

\section{Exam retrieval and common confusions}
Five distinctions prevent most errors on this topic.
\begin{enumerate}
  \item Animation applies to an \textbf{object inside a slide}; a transition
    applies between \textbf{slides}. Movement inside a slide is animation;
    movement between slides is a transition.
  \item \textbf{Entrance}, \textbf{Exit}, \textbf{Emphasis}, and
    \textbf{Motion Path} are animation effects, not transitions.
  \item \textbf{F5} starts from the beginning; \textbf{Shift+F5} starts from
    the current slide; \textbf{Alt+F5} opens Presenter View; \textbf{Esc} ends
    the show. Do not confuse F5 with F7 (spell check) (TRAP-32).
  \item The \textbf{Notes pane} is for the presenter; a \textbf{handout} is for
    the audience; \textbf{Slide Sorter} is for reordering slides.
  \item \textbf{SmartArt} converted to text loses its graphical layout and may
    need the content re-entered.
\end{enumerate}

\begin{example}
An item lists four effects and asks which one is a transition. Read each
option against the definition: Entrance, Exit, and Emphasis all act on an
object within a slide, while Fade acts on the move between slides. Fade is the
transition. The same logic applies when an option offers Morph: Morph is a
transition, so it belongs with Fade, not with the animation effects.
\end{example}

\subsection*{Exercises}
\begin{enumerate}[label=\arabic*.]
  \item Distinguish an animation from a transition, giving one example of each.
  \item Name the four animation effect categories and state what each does.
  \item State the action of F5, Shift+F5, Alt+F5, and Esc during a slide show.
  \item Explain what the Notes pane is for and how a handout differs from it.
  \item What happens to a SmartArt graphic when it is converted to text?
\end{enumerate}

\subsection*{Solutions}
\begingroup\small
\begin{enumerate}[label=\arabic*.]
  \item An animation is an effect on one object or text inside a slide (for
    example, a picture entering). A transition is an effect on the change
    between two slides (for example, a Fade).
  \item Entrance makes an object appear; Exit makes an object leave; Emphasis
    draws attention to an object that stays on the slide; Motion Path moves an
    object along a defined path.
  \item F5 starts the show from the beginning; Shift+F5 starts from the current
    slide; Alt+F5 starts Presenter View; Esc ends the show.
  \item The Notes pane holds speaker notes for the presenter and is not shown
    to the audience during a normal show. A handout is the printed layout given
    to the audience, usually with several slides per page and room to write.
  \item It becomes a plain bulleted list; the graphical layout is lost, and the
    design must be re-applied or the SmartArt re-created.
\end{enumerate}
\endgroup

\lessonheading{33}{PDF and File Formats}
\section{What a file format is}
Every file on a computer is stored as a long series of bits. On its own, a
string of bits says nothing about whether it is a letter, a photograph, a
spreadsheet, or a slide show. The \textbf{file format} is the rule that tells
the software how those bits are arranged --- where the text begins, where the
image data sits, and how the page should be drawn. When a program opens a
file, it reads the format to know how to interpret the contents.

The \textbf{extension} is the short suffix at the end of a file name, such as
\texttt{.pdf}, \texttt{.docx}, or \texttt{.png}. The extension is a label
that helps the operating system choose which program should open the file.
It is convenient but only a name: the real identity of the file is its format.

\begin{definitionbox}{File format}
The defined structure or rule by which data is stored in a file, so that
software can correctly read, display, and save its contents.
\end{definitionbox}

\begin{definitionbox}{Extension}
The suffix at the end of a file name, usually written with a dot, that
indicates the file's intended type (for example \texttt{.pdf}, \texttt{.docx},
\texttt{.png}).
\end{definitionbox}

\begin{center}
\begin{tabular}{p{0.28\textwidth}p{0.58\textwidth}}
\toprule
\textbf{Term} & \textbf{Meaning in this lesson} \\
\midrule
Format & The actual structure of the data inside the file. \\
Extension & The dot-suffix in the file name, such as \texttt{.docx}. \\
Fixed layout & A format that stores exact positions, so the page looks the same everywhere. \\
Plain text & A file of ordinary characters with no formatting. \\
\bottomrule
\end{tabular}
\end{center}

\section{Extension versus format (the core trap)}
The single most repeated confusion in this topic is treating the extension and
the format as the same thing. They are not. In the file name
\texttt{report.docx}, the base name is \texttt{report}, the extension is
\texttt{docx} (written \texttt{.docx}), and the format is DOCX, the Office
Open XML word-processing format. Likewise, \texttt{.png} is an extension,
while PNG (Portable Network Graphics) is the format.

The extension can also be changed easily, and doing so does not convert the
file. Renaming a photograph from \texttt{pic.jpg} to \texttt{pic.png} does
not turn a JPEG image into a PNG image; it only changes the label. This is
why the exam stresses the distinction: the extension is a name, the format is
the content.

\begin{example}
An item states, ``\texttt{.png} is an image file format.'' The statement
mixes the two ideas. \texttt{.png} is the \textbf{extension}; PNG is the
\textbf{format}. A correct statement would be, ``PNG is an image file format,
and \texttt{.png} is its usual extension.''
\end{example}

\section{Common document formats}
The Office applications each save their work in a standard format. These are
the Office Open XML formats, and each has its own extension.

\begin{center}
\begin{tabular}{p{0.16\textwidth}p{0.30\textwidth}p{0.40\textwidth}}
\toprule
\textbf{Extension} & \textbf{Format} & \textbf{Typical use} \\
\midrule
\texttt{.docx} & Word document & Text, formatting, and images (MS Word) \\
\texttt{.xlsx} & Excel workbook & Worksheets, formulas, and charts (MS Excel) \\
\texttt{.pptx} & PowerPoint presentation & Slides and slide shows (PowerPoint) \\
\texttt{.csv} & Comma-Separated Values & Plain-text rows with values separated by commas \\
\texttt{.txt} & Plain text & Unformatted characters only \\
\bottomrule
\end{tabular}
\end{center}

A \textbf{word-processing document} (DOCX) stores formatted text. A
\textbf{workbook} (XLSX) stores worksheets with formulas and numbers. A
\textbf{presentation} (PPTX) stores slides. A \textbf{CSV} file is plain text
in which each line is a record and commas separate the fields; it is often
used to move data between programs. A \textbf{TXT} file holds plain
characters and no formatting at all.

\section{Common image extensions}
Images also come in many formats, each suited to a different purpose.
\textbf{PNG} (Portable Network Graphics) uses lossless compression and
supports transparency, so it is good for logos and screenshots.
\textbf{JPG/JPEG} uses lossy compression, so it is good for photographs where
a small loss of quality is acceptable. \textbf{GIF} (Graphics Interchange
Format) handles simple images and short animations. \textbf{BMP} (Bitmap)
stores a raw image and is usually larger. \textbf{TIFF} (Tagged Image File
Format) is used for high-quality scanning and printing. \textbf{SVG}
(Scalable Vector Graphics) stores a drawing as lines and shapes, so it scales
without blurring.

\begin{center}
\begin{tabular}{p{0.16\textwidth}p{0.34\textwidth}p{0.36\textwidth}}
\toprule
\textbf{Extension} & \textbf{Format} & \textbf{Character} \\
\midrule
\texttt{.png} & Portable Network Graphics & Lossless, transparency \\
\texttt{.jpg} & JPEG image & Lossy, photos \\
\texttt{.gif} & Graphics Interchange Format & Simple images, animation \\
\texttt{.bmp} & Bitmap & Raw, larger \\
\texttt{.tiff} & Tagged Image File Format & High quality, scanning/print \\
\texttt{.svg} & Scalable Vector Graphics & Vector, scales without blur \\
\bottomrule
\end{tabular}
\end{center}

\section{PDF: the Portable Document Format}
\textbf{PDF} stands for \textbf{Portable Document Format}. It is a
\textbf{fixed-layout} format: every item on the page is placed at a fixed
position. Because the file also carries the fonts and images it needs, a PDF
looks the \textbf{same on every device and printer}. This is why PDF is used
for e-books, forms, official letters, admit cards, and result sheets, where
the exact appearance matters.

Contrast this with a word-processing document, which is a \emph{re-flowable}
format: text re-flows to fit the screen, so the layout can shift when the
font or device changes. A PDF, by design, does not re-flow; it is a picture of
a finished page expressed as instructions.

\begin{definitionbox}{PDF}
Portable Document Format --- a fixed-layout file format that preserves the
appearance of a document so it renders identically across devices and
printers.
\end{definitionbox}

\section{Text and scanned PDFs}
Not every PDF contains real text. It is important to separate two kinds.

A \textbf{text PDF} (sometimes called an editable-source PDF) is generated
from a source document such as a Word file. Its pages contain genuine text and
vector graphics. The text can be selected, copied, and searched, and the file
is usually smaller.

A \textbf{scanned PDF} (image-only PDF) is produced by scanning printed pages
or photographing them. Each page is stored as an \emph{image}, much like a
photograph. The text is readable to the human eye but not to the computer: it
cannot be selected, copied, or searched. Scanned files are also usually
larger.

\begin{example}
An item says, ``In a scanned PDF, you can search for a word using
\texttt{Ctrl+F}.'' This is wrong. A scanned PDF is only a picture of the
page, so its text is not searchable until OCR is applied. In a text PDF, the
same search works immediately.
\end{example}

\section{OCR: Optical Character Recognition}
\textbf{OCR} stands for \textbf{Optical Character Recognition}. It is the
technique that reads the \emph{picture of the text} and converts it into
\textbf{machine-readable text}. Once OCR has been run, the text in a scanned
document can be selected, searched, copied, and edited just like the text in a
text PDF.

OCR is used on scanned pages and on photographs of documents. It is the
bridge that turns an image-only PDF into a searchable one. The recognition is
not always perfect, so the result should be checked.

\begin{definitionbox}{OCR (Optical Character Recognition)}
The process of converting an image of text into machine-readable, editable,
and searchable text.
\end{definitionbox}

\section{Encryption and password protection}
A PDF can be protected so that its contents are confidential.
\textbf{Encryption} scrambles the file so that it cannot be read without the
correct key or password. Two kinds of password are common. A \textbf{password
to open} (user password) is required before the file will open at all. A
\textbf{permissions password} (owner password) restricts what a reader may
do --- for example, printing, copying, or editing.

Encryption provides \textbf{confidentiality}: only the authorised reader can
see the contents. It does not prove who created the file, and it is not the
same as a digital signature.

\section{Digital signature}
A \textbf{digital signature} is attached to a file to prove \textbf{who
signed it}. It also proves \textbf{integrity}: if the file is altered after
it has been signed, the signature no longer matches and the change is
detected. A digital signature uses a \textbf{digital certificate} issued by a
trusted authority.

A digital signature provides \textbf{authenticity} (the signer is verified)
and \textbf{integrity} (the file has not changed). It does \emph{not} hide
the contents. This is the distinction the exam most often tests: a digital
signature is about \emph{who} and \emph{whether it changed}, while encryption
is about \emph{who can read it}.

\begin{definitionbox}{Digital signature}
An electronic signature that verifies the identity of the signer and detects
any change made to the file after signing, providing authenticity and
integrity.
\end{definitionbox}

\begin{center}
\begin{tabular}{lll}
\toprule
\textbf{Point} & \textbf{Encryption} & \textbf{Digital signature} \\
\midrule
Main purpose & Confidentiality & Authenticity and integrity \\
Result & Contents unreadable & Signer verified, tampering detected \\
Needs & Key / password & Digital certificate \\
\bottomrule
\end{tabular}
\end{center}

\begin{example}
An item states, ``A digital signature encrypts the PDF so that only the
recipient can read it.'' The statement is wrong. The encryption that keeps
others from reading the file is \textbf{encryption}. A digital signature does
not hide the content; it confirms the signer and shows whether the file has
been changed.
\end{example}

\section{Redaction and PDF/A}
\textbf{Redaction} means \textbf{permanently removing} sensitive content
before a file is shared. Names, account numbers, addresses, or photographs
may be redacted. A proper redaction deletes the underlying text or image; it
is not enough to draw a black rectangle over it, because hidden text can
often still be selected and copied. A careful reader should always check that
the removed text is truly gone.

\textbf{PDF/A} is a special ISO-standard version of PDF designed for
\textbf{long-term archiving}. It embeds the fonts and colour information
inside the file so that the document will render exactly the same years later,
even on future software. The ``A'' in PDF/A stands for \emph{archival}. It is
used by libraries, courts, and government records, and is meant for
preservation rather than everyday editing.

\section{Exam retrieval and common confusions}
Five distinctions prevent most errors on this topic.
\begin{enumerate}
  \item The \textbf{extension} is the file-name suffix; the \textbf{format}
    is the structure of the data. \texttt{.png} is an extension; PNG is the
    format (TRAP-34).
  \item A \textbf{scanned PDF} is an image; its text cannot be selected or
    searched until \textbf{OCR} is applied. A text PDF is searchable at once.
  \item \textbf{Encryption} gives confidentiality; a \textbf{digital
    signature} gives authenticity and integrity. They are not the same
    (TRM-11, TRAP-15).
  \item \textbf{Redaction} removes sensitive content permanently; a black
    box drawn over text is not sufficient.
  \item \textbf{PDF/A} is for long-term \textbf{archiving}; ordinary PDF is
    the everyday fixed-layout format.
\end{enumerate}

\begin{example}
An item asks which format keeps a fixed layout and looks the same on every
device. The answer is PDF. If the stem instead asks which format stores
spreadsheet data, the answer is XLSX; if it asks which is a plain-text data
format, the answer is CSV; and if it asks which format is meant for long-term
archiving, the answer is PDF/A.
\end{example}

\subsection*{Exercises}
\begin{enumerate}[label=\arabic*.]
  \item Distinguish a file extension from a file format, with one example.
  \item State the full form of PDF and explain why PDF is called a
    fixed-layout format.
  \item Explain the difference between a text PDF and a scanned PDF, and
    where OCR fits in.
  \item Distinguish encryption from a digital signature in terms of the
    security goal each provides.
  \item What is redaction, and why is drawing a black box over text not a
    safe redaction? What is PDF/A used for?
\end{enumerate}

\subsection*{Solutions}
\begingroup\small
\begin{enumerate}[label=\arabic*.]
  \item The extension is the dot-suffix in the file name, such as
    \texttt{.docx}; it is only a label. The format is the actual structure of
    the data inside the file, such as the DOCX word-processing format. The
    same extension can be renamed without changing the format.
  \item PDF stands for Portable Document Format. It is called a fixed-layout
    format because every item is placed at a fixed position on the page, and
    the fonts and images are carried in the file, so the document looks the
    same on every device and printer.
  \item A text PDF contains real, machine-readable text, so it can be
    selected, copied, and searched. A scanned PDF stores each page as an
    image, so its text cannot be selected or searched. OCR (Optical Character
    Recognition) converts the image of the text into machine-readable text,
    after which the scanned PDF becomes searchable.
  \item Encryption provides confidentiality: it scrambles the file so only
    someone with the correct key or password can read it. A digital signature
    provides authenticity and integrity: it proves who signed the file and
    detects any change made after signing. It does not hide the contents.
  \item Redaction is the permanent removal of sensitive content before a file
    is shared. Drawing a black box over text is not safe because the hidden
    text may still be selectable and copyable underneath. PDF/A is a
    standardised PDF version for long-term archiving that embeds fonts and
    colour information so the document renders the same in the future.
\end{enumerate}
\endgroup

\lessonheading{54}{Office Shortcut and Product Trap Lab}
\section{Why this lesson is a trap drill}
The computer section of JKSSB papers does not reward vague recognition. It
tests the \textbf{exact key} and the \textbf{exact product category}. A wrong
option is usually not random; it is the nearest neighbour to the right answer
--- the key that differs by one letter, the same function key in a different
program, or a familiar software name placed in the wrong category. The safe
method is to read every option as a claim to be tested and to keep the exact
English term: \texttt{F5}, \texttt{F7}, \texttt{Ctrl+H}, \texttt{\$A\$1}.
This lesson gathers the Word, PowerPoint, Excel, and product-category traps
that JKSSB reuses.

\section{MS Word: exact shortcut keys}
Word has a small set of shortcuts that JKSSB asks again and again. Learn each
key with its full action.

\begin{center}
\begin{tabular}{p{0.26\textwidth}p{0.62\textwidth}}
\toprule
\textbf{Shortcut} & \textbf{Action in MS Word} \\
\midrule
F7 & Spelling \& Grammar (spell check) \\
Shift+F7 & Thesaurus \\
Ctrl+H & Find and Replace \\
Ctrl+F & Find only (opens the Navigation pane) \\
Ctrl+K & Insert Hyperlink \\
Ctrl+X & Cut the selection \\
Ctrl+C & Copy the selection \\
Ctrl+V & Paste the Clipboard contents \\
\bottomrule
\end{tabular}
\end{center}

\begin{definitionbox}{Find and Replace}
A Word tool, opened with \textbf{Ctrl+H}, that searches for a word or phrase
and replaces it with another across the document.
\end{definitionbox}

The one distinction to fix first is \textbf{Ctrl+H} versus \textbf{Ctrl+F}.
\texttt{Ctrl+H} opens Find and Replace, so it can change text;
\texttt{Ctrl+F} only finds text. The second is \textbf{F7} versus
\textbf{Shift+F7}: \texttt{F7} checks spelling and grammar, while
\texttt{Shift+F7} opens the Thesaurus. A third is \textbf{Ctrl+K}, which
inserts a hyperlink, not a comment or a caption.

\section{Word views and review tools}
Word offers several views of the same document. \textbf{Print Layout} shows
the document exactly as it will appear on the printed page, with margins,
headers, and footers in place. The other common views are Web Layout, Outline,
and Draft. For exam purposes, the key phrase ``as it will appear when printed''
points to Print Layout.

\textbf{Track Changes} is Word's review tool. When it is switched on, every
insertion, deletion, and formatting change is recorded rather than applied
silently. Each marked change can later be \textbf{accepted} or
\textbf{rejected} by a reviewer. Track Changes also stores the author's name
and the time of the edit as document metadata. The important trap is that
Track Changes does \emph{not} delete the original text permanently: the
original remains available until a reviewer accepts or rejects the change
(PYQ-12).

\begin{example}
A candidate reads the option ``Track Changes permanently deletes the original
text, so it can never be restored.'' This is false. Track Changes preserves
the original text and marks the edit; the reviewer decides whether to accept
or reject it. The trap works because ``records changes'' sounds close to
``deletes the old version.''
\end{example}

\section{Word: Section Break versus Page Break}
A \textbf{Page Break} starts a new page but keeps the document in the same
section. A \textbf{Section Break} starts a new section. The distinction
matters because only a Section Break lets the pages that follow use a
\emph{different} header or footer. This is a repeated trap: the wrong option
claims that a Page Break changes the header or footer, which it does not
(PYQ-13; TRAP-13).

\begin{definitionbox}{Section Break}
A break that divides a document into sections, so each section can have its
own page setup, including a different header and footer.
\end{definitionbox}

\section{Word feature trap: Mail Merge}
\textbf{Mail Merge} combines a \textbf{main document} (the common letter or
label) with a \textbf{data source} such as an address database, and produces
many personalised copies. It is a Word feature. The trap is to confuse it
with Track Changes or AutoCorrect. The reliable clue in the question is the
pair ``combines letters with an address database'', which points to Mail
Merge and not to any other Word tool (PYQ-37; TRAP-29).

\section{PowerPoint: the show keys}
PowerPoint's presentation shortcuts are a favourite near-neighbour set.

\begin{center}
\begin{tabular}{p{0.26\textwidth}p{0.60\textwidth}}
\toprule
\textbf{Shortcut} & \textbf{Action in PowerPoint} \\
\midrule
F5 & Start the slide show from the beginning \\
Shift+F5 & Start the slide show from the current slide \\
Alt+F5 & Start the presentation in Presenter View \\
Esc & End the running slide show \\
\bottomrule
\end{tabular}
\end{center}

The classic trap is to call \texttt{F5} the spell-check key. In PowerPoint,
\texttt{F5} starts the show from the first slide; \texttt{F7} checks
spelling. \texttt{Shift+F5} begins from the slide you are on, and
\texttt{Alt+F5} opens Presenter View, which shows the speaker notes on the
presenter's screen while the audience sees only the slides. \texttt{Esc} ends
the show but does not jump to the first slide (PYQ-40; TRAP-32).

\section{Handouts, Slide Master, and orientation}
\textbf{Handouts} are a print layout that places several slides on one page
for the audience, usually with lined space for notes. The \textbf{Handout
Master} controls the layout of handouts and notes pages.

The \textbf{Slide Master} is the design template behind the whole
presentation. It controls the fonts, colours, background, and placeholder
positions of every slide that uses it. Editing the Slide Master changes all
dependent slides at once, which is exactly why it is used for a consistent
look.

\textbf{Orientation} is the direction of the page. \textbf{Portrait} is taller
than it is wide; \textbf{Landscape} is wider than it is tall. Slide
orientation is set in PowerPoint, and notes and handouts have their own
orientation setting.

\section{PowerPoint: converting to SmartArt}
\textbf{SmartArt} turns text into a visual diagram such as a process, cycle,
or hierarchy. A \textbf{bulleted list}, a \textbf{block of text}, and
\textbf{indented text} can all be converted directly to SmartArt. A
\textbf{table of data} cannot be converted directly, because a table is a
structured object rather than a plain text string; the Convert to SmartArt
command is disabled for it. To use table content, the text must first be
copied into a text box (PYQ-17).

\section{Excel: the formula palette}
Every Excel formula begins with the \textbf{=} sign (PYQ-38; TRAP-30). The
\textbf{formula palette} --- the Insert Function dialog --- is the helper that
lets you build a formula. It lists the function and shows each
\textbf{argument} in its own box, so you can fill in the values one at a
time. It is opened from the \textbf{fx} button next to the formula bar. The
palette guides formula entry; it does not resize cells or change data.

\begin{definitionbox}{Formula palette}
The Excel dialog box that lets a user select a function and enter its
arguments while building a formula.
\end{definitionbox}

\section{Excel: cell references and the \$ sign}
A cell reference tells Excel which cell to use. The reference type decides
what happens when the formula is copied to another cell.

\begin{center}
\begin{tabular}{p{0.20\textwidth}p{0.28\textwidth}p{0.40\textwidth}}
\toprule
\textbf{Reference} & \textbf{Type} & \textbf{Behaviour when copied} \\
\midrule
A1 & Relative & Both column and row shift \\
\$A\$1 & Absolute & Neither column nor row changes \\
\$A1 & Mixed & Column fixed, row shifts \\
A\$1 & Mixed & Row fixed, column shifts \\
\bottomrule
\end{tabular}
\end{center}

The \textbf{\$} sign locks the part immediately after it. In \texttt{\$A\$1}
both the column and the row are locked, so the reference is
\textbf{absolute}: it always points to the same cell no matter where the
formula is copied. \texttt{\$A1} locks the column only, and \texttt{A\$1}
locks the row only; both are \textbf{mixed} references. While editing a
reference, pressing \textbf{F4} cycles through relative, absolute, and mixed
forms.

\begin{example}
A formula in \texttt{D4} reads \texttt{=B4*C4}. Copied down to \texttt{D5}, it
becomes \texttt{=B5*C5}, because the references are relative. If instead the
formula is \texttt{=\$B\$4*\$C\$4}, it stays \texttt{=\$B\$4*\$C\$4} in
every cell, because the \texttt{\$} signs make both references absolute.
\end{example}

\section{Product-category traps: Word, Excel, Access, MS-DOS}
JKSSB often places four familiar names together and asks which one belongs to
a given category. Match each name to its exact category.

\begin{center}
\begin{tabular}{p{0.26\textwidth}p{0.58\textwidth}}
\toprule
\textbf{Program} & \textbf{Correct category} \\
\midrule
MS Word & Word processor (text documents) \\
MS Excel & Spreadsheet (cells and formulas) \\
MS Access & Database management system (tables, queries, reports) \\
MS-DOS & Disk Operating System (a command-line operating system) \\
\bottomrule
\end{tabular}
\end{center}

The trap is that all four are familiar computing names. \textbf{MS Access} is
a database, \emph{not} a spreadsheet; \textbf{MS Excel} is a spreadsheet,
\emph{not} a database; and none of Word, Excel, or Access is an operating
system. \textbf{MS-DOS} is the only operating system in the group.

\section{Exam retrieval and common confusions}
\begin{enumerate}
  \item Word \textbf{F7} is spell check; \textbf{Ctrl+H} is Find and Replace;
    \textbf{Ctrl+K} is Insert Hyperlink.
  \item \textbf{Section Break}, not Page Break, allows a different
    header/footer on later pages (TRAP-13).
  \item PowerPoint \textbf{F5} starts from the beginning, \textbf{Shift+F5}
    from the current slide, \textbf{Alt+F5} is Presenter View, and
    \textbf{Esc} ends the show (TRAP-32).
  \item A \textbf{table} cannot be converted directly to SmartArt.
  \item Excel \textbf{\$A\$1} is an absolute reference; every formula starts
    with \textbf{=}.
  \item \textbf{MS Access} is a database; \textbf{MS-DOS} is an operating
    system.
\end{enumerate}

\begin{example}
A question asks which key starts a PowerPoint slide show from the beginning.
Trace the set: \texttt{F5} begins the show, \texttt{Shift+F5} starts from the
current slide, \texttt{Alt+F5} opens Presenter View, and \texttt{Esc} ends the
show. The answer is \texttt{F5}. If the stem instead says ``from the current
slide'', the answer is \texttt{Shift+F5}.
\end{example}

\subsection*{Exercises}
\begin{enumerate}[label=\arabic*.]
  \item State the action of F7, Ctrl+H, and Ctrl+K in MS Word.
  \item Distinguish Print Layout from the other Word views.
  \item Explain what Track Changes records and why it does not permanently
    delete the original text.
  \item Distinguish a Page Break from a Section Break, and say which one
    allows a different header.
  \item Give the PowerPoint keys for starting a show from the beginning, from
    the current slide, for Presenter View, and for ending the show.
  \item Explain the difference between the relative reference A1 and the
    absolute reference \$A\$1 in Excel.
  \item Classify MS Word, MS Excel, MS Access, and MS-DOS by category.
\end{enumerate}

\subsection*{Solutions}
\begingroup\small
\begin{enumerate}[label=\arabic*.]
  \item \texttt{F7} opens Spelling \& Grammar; \texttt{Ctrl+H} opens Find and
    Replace; \texttt{Ctrl+K} inserts a hyperlink.
  \item Print Layout shows the document as it will appear on the printed
    page, with margins and headers in place. Web Layout, Outline, and Draft
    are the other common views and are not print-accurate.
  \item Track Changes records insertions, deletions, and formatting edits,
    along with author and time metadata. It marks the edits instead of
    applying them silently, so the original text is preserved until a
    reviewer accepts or rejects each change.
  \item A Page Break starts a new page within the same section; a Section
    Break starts a new section. Only a Section Break allows a different
    header or footer on the following pages.
  \item \texttt{F5} from the beginning; \texttt{Shift+F5} from the current
    slide; \texttt{Alt+F5} for Presenter View; \texttt{Esc} to end the show.
  \item A1 is relative, so both its column and row shift when the formula is
    copied. \$A\$1 is absolute, so neither the column nor the row changes; the
    reference always points to the same cell.
  \item MS Word is a word processor; MS Excel is a spreadsheet; MS Access is
    a database management system; MS-DOS is an operating system.
\end{enumerate}
\endgroup

\lessonheading{55}{File Formats and Named Technology Associations}
\section{Extension and format are two different things}
Every file stored on a computer follows some fixed plan for arranging its
data. That plan is the \textbf{format} of the file. The format decides which
program can open the file and what the file is able to store: a picture, a
document, a table, or a database. An \textbf{image format} stores a picture;
a \textbf{document format} stores a printable page; a \textbf{plain-text
format} stores only characters.

A file name usually ends with an \textbf{extension} --- the letters written
after the last dot. In the name \texttt{report.png}, the extension is
\texttt{png}. The extension is only a label; it tells the reader which format
is expected inside. The extension and the format are therefore related but
not the same thing. The exam tests this distinction directly, so it is stated
here in its own section.

\begin{definitionbox}{File format}
A fixed way of arranging the data in a file so that a particular program can
read it, write it, and display it correctly.
\end{definitionbox}

\begin{definitionbox}{File extension}
The letters after the last dot in a file name, such as \texttt{png} in
\texttt{report.png}. The extension signals the file's format.
\end{definitionbox}

\begin{center}
\begin{tabular}{p{0.26\textwidth}p{0.62\textwidth}}
\toprule
\textbf{Term} & \textbf{Meaning in this lesson} \\
\midrule
Extension & The letters after the last dot in a file name. \\
Format & The actual arrangement of data inside the file. \\
PNG & Portable Network Graphics, an image format. \\
PDF & Portable Document Format, a fixed-layout document format. \\
CSV & Comma-Separated Values, a table stored as plain text. \\
TXT & A plain text file with no formatting. \\
\texttt{.mdb} / \texttt{.accdb} & Older and newer Microsoft Access database formats. \\
\bottomrule
\end{tabular}
\end{center}

\section{Image formats}
The best-known image format for this course is \textbf{PNG}, which stands for
\textbf{Portable Network Graphics}. PNG is a raster image format: it stores a
picture as a grid of pixels. It uses \emph{lossless} compression, meaning the
picture is saved without permanently throwing away detail, and it supports
transparency.

Other common image formats are \textbf{JPEG} (written \texttt{.jpg} or
\texttt{.jpeg}), a \emph{lossy} format often used for photographs, and
\textbf{GIF} (\texttt{.gif}), which also supports simple animation. Typical
image extensions are \texttt{.png}, \texttt{.jpg}, \texttt{.jpeg},
\texttt{.gif}, \texttt{.bmp}, and \texttt{.tiff}. The key exam point is
simple: PNG, JPEG, and GIF are all \textbf{image formats}, and none of them
is a document, spreadsheet, or database format.

\begin{definitionbox}{Raster image}
An image stored as a grid of pixels. PNG and JPEG are raster formats.
\end{definitionbox}

\section{PDF: the Portable Document Format}
\textbf{PDF} stands for \textbf{Portable Document Format}, a format created
by Adobe. Its defining property is a \textbf{fixed layout}: a PDF keeps the
exact page layout, fonts, and formatting on every device, so the document
looks the same whether it is opened on a desktop, a phone, or a printer. This
is why PDF is the standard format for sharing a finished document that must
not shift or reflow.

A PDF opens in a \textbf{PDF reader}, the best-known being \textbf{Adobe
Acrobat Reader}. A \textbf{web browser} can also display a PDF, but reading a
PDF is not the same as editing an editable word-processing document. A PDF
may contain real, selectable text, or it may be a \textbf{scanned image} of
each page. A scanned PDF is only a picture; to make its text searchable and
selectable, \textbf{OCR} (Optical Character Recognition) must be applied.
\textbf{PDF/A} is an archival PDF standard meant for long-term preservation.

\begin{example}
A notice is scanned on a photocopier and saved as a PDF. The PDF now contains
pictures of the pages, not typed text. Selecting a sentence with the mouse
does nothing, and searching for a word finds nothing. Running OCR on the file
creates a text layer, after which the same words can be searched and copied.
The format is still PDF; only the internal content has changed from an image
to searchable text.
\end{example}

\section{The Office document family}
Microsoft Office files use a matching family of extensions that are easy to
remember as a group.

\begin{center}
\begin{tabular}{p{0.16\textwidth}p{0.34\textwidth}p{0.36\textwidth}}
\toprule
\textbf{Extension} & \textbf{Format} & \textbf{Program} \\
\midrule
\texttt{.docx} & Word document & Microsoft Word \\
\texttt{.xlsx} & Excel workbook & Microsoft Excel \\
\texttt{.pptx} & PowerPoint presentation & Microsoft PowerPoint \\
\bottomrule
\end{tabular}
\end{center}

These are the modern \textbf{Office Open XML} formats. The older, pre-2007
extensions were \texttt{.doc} (Word), \texttt{.xls} (Excel), and \texttt{.ppt}
(PowerPoint). For an exam, the pairing to hold is: \textbf{.docx = Word},
\textbf{.xlsx = Excel}, \textbf{.pptx = PowerPoint}; each extension always
belongs to the same program, in the modern and the older naming alike.

\section{Plain-data formats: CSV and TXT}
\textbf{CSV} stands for \textbf{Comma-Separated Values}. A CSV file stores a
table as plain text, with the values in each row separated by commas. Because
it is only text, a CSV file can be opened by a spreadsheet program such as
Excel, by a database, or by a simple text editor. The commas are the field
separators; the file itself carries no formatting such as bold text or
colours.

\textbf{TXT} denotes a \textbf{plain text} file. It stores only characters,
with no formatting and no table structure. Both CSV and TXT are simple,
vendor-neutral formats: unlike a \texttt{.docx} or \texttt{.xlsx} file, they
do not belong to one program. The distinction to keep is that \textbf{CSV
holds a table written as text separated by commas}, while \textbf{TXT holds
plain characters with no inherent structure}.

\section{Microsoft Access databases: \texttt{.mdb} and \texttt{.accdb}}
Microsoft Access is the database program of the Office family, and it uses
two extensions that candidates often confuse. \texttt{.mdb} is the
\textbf{older} Access database format, used by Access 2003 and earlier.
\texttt{.accdb} is the \textbf{newer} Access database format, used from
Access 2007 onward. Both files can hold the standard database objects:
tables, queries, forms, reports, macros, and modules.

The important point is that \textbf{both} \texttt{.mdb} and \texttt{.accdb}
are \textbf{Microsoft Access database} files. Neither is a Word document, an
Excel workbook, or a PowerPoint presentation.

\section{Named technology associations}
This section is the high-yield part of the lesson: matching a named
technology to the organisation that produced it or the purpose it serves.

\textbf{Firefox} is a \textbf{web browser} developed by \textbf{Mozilla}.
Mozilla is the organisation, and Firefox is its browser product; Mozilla also
produces \textbf{Thunderbird}, an email client. Firefox must not be confused
with Microsoft's browser, \textbf{Edge} (or the older Internet Explorer), nor
with a PDF reader.

\textbf{C-DAC} stands for the \textbf{Centre for Development of Advanced
Computing}. It is an Indian research and development organisation under the
\textbf{Ministry of Electronics and Information Technology (MeitY)}.
\textbf{PARAM} is a family of \textbf{supercomputers} developed by C-DAC;
\textbf{PARAM 8000} was India's first indigenous supercomputer. The pairing
to remember is \textbf{PARAM = Indian supercomputer, C-DAC = the organisation
that developed it}.

\section{Software roles: browser, PDF reader, and remote support}
Software is often tested by the role it performs rather than by its maker.

\begin{center}
\begin{tabular}{p{0.30\textwidth}p{0.56\textwidth}}
\toprule
\textbf{Role} & \textbf{Example} \\
\midrule
Web browser & Firefox, Google Chrome, Microsoft Edge \\
PDF reader & Adobe Acrobat Reader \\
Remote-support tool & TeamViewer, AnyDesk, Chrome Remote Desktop \\
\bottomrule
\end{tabular}
\end{center}

A \textbf{web browser} is used to visit web pages; a \textbf{PDF reader} is
used to open PDF documents; a \textbf{remote-support tool} lets one computer
control or view another over a network. Firefox is a browser and Adobe
Acrobat Reader is a PDF reader; neither is a remote-support tool. Match the
software to its \emph{job}, not to a similar-sounding name.

\begin{example}
A help-desk officer needs to see a distant user's screen to fix a problem.
The right tool is a remote-support program such as TeamViewer or AnyDesk,
\emph{not} a browser and \emph{not} a PDF reader. If the user then sends a
manual as a fixed-layout file, it will arrive as a PDF and open in a PDF
reader. If the user sends a spreadsheet of data, it will be an XLSX (Excel)
or a CSV. Each task has its own format and its own tool.
\end{example}

\section{Exam retrieval and common confusions}
Six distinctions prevent most errors on this topic.
\begin{enumerate}
  \item An \textbf{extension} is the letters after the dot; a \textbf{format}
    is how the data is stored. ``Extension = format'' is wrong (TRM-09,
    TRAP-34).
  \item \textbf{PNG, JPEG, and GIF are image formats}. CSV is \textbf{not} an
    image format; it is a plain-text table format.
  \item \textbf{PDF} is a fixed-layout document format. A scanned PDF is a
    picture and needs \textbf{OCR} to become searchable.
  \item \textbf{.docx = Word}, \textbf{.xlsx = Excel}, \textbf{.pptx =
    PowerPoint}; the extension always matches the program.
  \item \textbf{.mdb} and \textbf{.accdb} are \textbf{both Access database}
    files --- older and newer respectively.
  \item \textbf{Firefox is a browser by Mozilla}; \textbf{PARAM is a
    supercomputer by C-DAC}. Do not swap the makers.
\end{enumerate}

\subsection*{Exercises}
\begin{enumerate}[label=\arabic*.]
  \item Distinguish a file \emph{extension} from a file \emph{format}, using
    \texttt{report.png} as the example.
  \item Name any three image formats and state one property of PNG.
  \item Explain what is meant by a fixed-layout format, and name the format
    created by Adobe that behaves this way.
  \item Match each extension to its program: \texttt{.docx},
    \texttt{.xlsx}, \texttt{.pptx}, \texttt{.accdb}.
  \item What does CSV stand for, and how is its data arranged?
  \item (a) Who develops Firefox? (b) What is PARAM, and which organisation
    develops it?
\end{enumerate}

\subsection*{Solutions}
\begingroup\small
\begin{enumerate}[label=\arabic*.]
  \item The extension is the suffix after the last dot --- here
    \texttt{png}. The format is the way the data inside the file is actually
    stored --- here PNG, the Portable Network Graphics image format. The
    extension signals the format but is not the format itself.
  \item PNG (Portable Network Graphics), JPEG (\texttt{.jpg/.jpeg}), and GIF
    (\texttt{.gif}), among others. PNG uses lossless compression and supports
    transparency.
  \item A fixed-layout format keeps the exact page layout, fonts, and
    formatting on every device, so the document looks the same wherever it is
    opened. PDF (Portable Document Format) is the format created by Adobe
    that behaves this way.
  \item \texttt{.docx} $\rightarrow$ Microsoft Word; \texttt{.xlsx}
    $\rightarrow$ Microsoft Excel; \texttt{.pptx} $\rightarrow$ Microsoft
    PowerPoint; \texttt{.accdb} $\rightarrow$ Microsoft Access (newer
    database format).
  \item CSV stands for Comma-Separated Values. Its data is stored as a table
    in plain text, with the values in each row separated by commas.
  \item (a) Firefox is developed by Mozilla. (b) PARAM is a family of
    supercomputers, developed by C-DAC (the Centre for Development of
    Advanced Computing); PARAM 8000 was India's first indigenous
    supercomputer.
\end{enumerate}
\endgroup

\lessonheading{61}{Current Office, Email and Statement Evaluation}
\section{Word: structured content and content controls}
Most Word documents are free-form, but many office documents --- forms,
templates, and records --- need \emph{structure}: a fixed place for a name, a
date, a choice from a list, or a check box. Word provides \textbf{content
controls} for exactly this task. A content control is an individual element
that you insert at a chosen point in a document or template; the user can type
or choose only what the control allows.

\begin{definitionbox}{Content control}
A single, structured element inserted into a Word document or template to
capture one piece of information in a controlled way.
\end{definitionbox}

Common content controls include plain text, rich text, picture, drop-down
list, combo box, date picker, and check box. They are inserted from the
\textbf{Developer} tab, which is hidden by default and is enabled through
Word Options $\rightarrow$ Customize Ribbon. A content control can restrict
editing at that spot and can enforce limited \textbf{data validation}: a
drop-down list allows only the listed items, and a date picker allows only a
date. This is the fact behind the JA 2026 word item (PYQ-12): a content
control \emph{can} enforce data validation, so the statement that it cannot is
the incorrect one.

Content controls should not be confused with \emph{placeholders}, which belong
to PowerPoint slide layouts, or with form \emph{fields} from older Word
versions. In this course, ``content control'' means the Word Developer-tab
element that structures a document.

\section{First-line and hanging indents}
Indentation decides how far a paragraph's text sits from the left (or right)
margin. Word distinguishes two special paragraph indents.

\begin{center}
\begin{tabular}{p{0.22\textwidth}p{0.66\textwidth}}
\toprule
\textbf{Indent} & \textbf{Meaning} \\
\midrule
First-line indent & Only the \textbf{first line} of the paragraph is indented;
the remaining lines start at the left margin. \\
Hanging indent & Every line \textbf{except the first} is indented; the first
line hangs out to the left. Used for references, bibliographies, and bulleted
lists. \\
\bottomrule
\end{tabular}
\end{center}

\begin{definitionbox}{Hanging indent}
A paragraph indent in which all lines except the first are pushed inward, so
the first line ``hangs'' at the left margin.
\end{definitionbox}

These are paragraph-level formatting settings. Two paragraphs can carry
different indents even when they use the \emph{same style}, because a
paragraph style supplies a default that individual paragraphs can override.
The usual wrong statement claims that a hanging indent shifts the first line;
in fact the first line is the one that stays out, and all the others move
inward.

\section{The Navigation Pane}
The \textbf{Navigation Pane} is a Word panel opened from the \textbf{View}
tab. It is a browsing aid rather than a formatting tool. It presents three
views or tabs.

\begin{enumerate}
  \item \textbf{Headings}: an outline of the document built from paragraphs
    that carry a \emph{Heading style} (Heading 1, Heading 2, and so on).
  \item \textbf{Pages}: thumbnail pictures of each page.
  \item \textbf{Results}: the matches for text typed into the search box.
\end{enumerate}

Because the Headings view is driven by heading styles, applying a Heading
style does two things at once: it gives the paragraph an outline level and it
makes the paragraph appear in the Navigation Pane (and in an automatic Table
of Contents).

\begin{example}
A student writes a long report. Applying \emph{Heading 1} to each chapter
title makes every chapter appear in the Navigation Pane, so the reader can
jump straight to any chapter. A reviewer's \emph{comment} on a sentence does
\emph{not} appear in the Navigation Pane; comments are read in the
markup/comments pane. That distinction is the trap in JA 2026 (PYQ-12).
\end{example}

\section{Track Changes and revision metadata}
\textbf{Track Changes} is Word's review tool. When it is switched on, every
edit is recorded as a \emph{revision} instead of being applied silently. Three
kinds of change are tracked: insertions, deletions, and formatting changes.
Each recorded change stores \textbf{metadata} --- the \textbf{author's name}
and the \textbf{date and time} of the edit.

\begin{definitionbox}{Revision metadata with Track Changes}
The information Word stores alongside a tracked change: who made it and when.
\end{definitionbox}

A tracked change can be \textbf{accepted} or \textbf{rejected}. Accepting a
change keeps the new text as final and \textbf{removes the revision metadata}
attached to that change; rejecting a change restores the original text. Before
sharing a document, Word's \emph{Inspect Document} feature can be used to
strip hidden personal information such as author names.

\section{Section Break versus Page Break}
Word documents are divided into \textbf{sections}, and page-level settings
belong to a section. The two breaks behave quite differently.

\begin{center}
\begin{tabular}{p{0.18\textwidth}p{0.36\textwidth}p{0.36\textwidth}}
\toprule
\textbf{Break} & \textbf{What it does} & \textbf{What it changes} \\
\midrule
Page Break & Moves the following text to the next page. & Nothing else; the
document stays in the same section with the same headers, footers, and
margins. \\
Section Break & Starts a \textbf{new section}. & Each section may have its own
headers and footers, margins, page orientation, columns, and page numbering. \\
\bottomrule
\end{tabular}
\end{center}

\begin{definitionbox}{Section Break}
A break that ends the current section and begins a new one, so that
page-level formatting can differ from that point onward.
\end{definitionbox}

\begin{example}
A report needs no header on the title page but a running header on every later
page. Inserting a single Page Break will not do it, because both pages remain
in one section. The correct tool is a \textbf{Section Break} after the title
page, after which the new section can carry its own header. This is the exact
distinction tested in JA 2026 (PYQ-13): Section Break, not Page Break, allows
different headers and footers in one document.
\end{example}

\section{Excel: nested IF, blank and non-numeric input}
An \textbf{IF} function chooses between two values:
\texttt{=IF(condition, value\_if\_true, value\_if\_false)}. A \textbf{nested
IF} places a second IF inside one of those two positions, which lets one
formula test several conditions in turn:
\[
\texttt{=IF(A1>50,"High",IF(A1>25,"Medium","Low"))}.
\]
Excel evaluates the conditions from left to right and returns the value of the
\textbf{first} condition that is TRUE; once a condition is satisfied, the
remaining tests are skipped.

\begin{example}
With A1 = 70 the first condition (70 > 50) is TRUE, so the result is
\emph{High}. With A1 = 30 the first condition is FALSE but the second
(30 > 25) is TRUE, so the result is \emph{Medium}. With A1 = 10 both are
FALSE, so the result is \emph{Low}.
\end{example}

The behaviour on \textbf{blank} and \textbf{non-numeric} input is the trap
(PYQ-14). In a numeric comparison Excel treats a \textbf{blank cell as zero}.
Therefore, if A1 is empty, \texttt{A1>50} is FALSE and
\texttt{A1>25} is also FALSE, so the formula falls through to the final value
--- \emph{Low}. A blank cell does not make the formula stop or produce an
error; it is simply read as 0.

Non-numeric \textbf{text} is different again. Text is not a number, so:
\begin{itemize}
  \item arithmetic on it fails --- for example \texttt{=A1+5} with text in A1
    returns the error \texttt{\#VALUE!};
  \item \texttt{ISNUMBER} returns FALSE for a text cell and for a blank cell.
\end{itemize}
For this reason, robust formulas guard the input, for instance with
\texttt{ISBLANK}, \texttt{ISNUMBER}, or \texttt{IFERROR}, before using the
value in a calculation.

\section{COUNTIFS and SUMPRODUCT}
Counting rows that meet \emph{several} conditions at once is a frequent task.
Excel offers two routes.

\begin{definitionbox}{COUNTIFS}
A function that counts the rows satisfying \emph{all} of the given
range-and-criteria pairs (AND logic):
\texttt{COUNTIFS(range1, criteria1, range2, criteria2, \ldots)}.
\end{definitionbox}

The \textbf{COUNTIFS} function takes \emph{paired} arguments: a range followed
by the criterion that applies to it, repeated for each condition. All the
ranges must be the same size. The singular \textbf{COUNTIF} handles one
condition; the plural \textbf{COUNTIFS} handles several. Text criteria can use
the wildcards \texttt{*} (any sequence of characters) and \texttt{?} (any one
character).

\begin{definitionbox}{SUMPRODUCT}
A function that multiplies corresponding elements of the given arrays and
returns the sum of those products.
\end{definitionbox}

On its own, \texttt{SUMPRODUCT(A1:A3,B1:B3)} returns
\texttt{A1*B1 + A2*B2 + A3*B3}. It can also be used to \emph{count}, by
building boolean arrays and multiplying them:

\begin{center}
\texttt{=SUMPRODUCT((A1:A10>50)*(B1:B10="Yes"))}
\end{center}

\noindent counts the rows in which the first condition AND the second
condition are both true, because the product is 1 only when both parts are 1.
The double-unary form \texttt{=SUMPRODUCT(--(A1:A10>50))} is a common way to
coerce TRUE/FALSE into 1/0.

\begin{example}
To count how many students scored above 50 \emph{and} belong to the Science
stream, either
\texttt{=COUNTIFS(Marks,">50",Stream,"Science")} or
\texttt{=SUMPRODUCT((Marks>50)*(Stream="Science"))} will work. The stable
distinction tested in JA 2026 (PYQ-15, TRAP-14) is that COUNTIFS takes several
criteria directly, while SUMPRODUCT does arithmetic on arrays and can
therefore both \emph{add} and \emph{count} --- it does not ``only add.''
\end{example}

\section{Email: filters, signatures, POP3 and privacy}
\subsection*{Filters and rules}
An email \textbf{filter}, usually called a \textbf{rule} in desktop mail
clients, acts automatically on messages. A rule has two parts: a
\textbf{condition} (sender, recipient, subject, keywords, presence of an
attachment) and an \textbf{action} (move to a folder, flag, forward, delete,
mark as read). A typical rule sends all mail from one mailing list into a
dedicated folder, keeping the inbox tidy.

\subsection*{Digital signature versus encryption}
These two security features are constantly confused (PYQ-16, TRAP-15).

\begin{center}
\begin{tabular}{p{0.22\textwidth}p{0.34\textwidth}p{0.34\textwidth}}
\toprule
\textbf{Point} & \textbf{Digital signature} & \textbf{Encryption} \\
\midrule
Purpose & Authenticity and integrity & Confidentiality \\
What it establishes & Who sent the message, and that it was not altered in
transit & That only the intended reader can read the contents \\
Does it hide the content? & No & Yes \\
\bottomrule
\end{tabular}
\end{center}

A \textbf{digital signature} does \emph{not} encrypt the message body. It
proves the sender's identity and detects tampering. It is \textbf{encryption}
that scrambles the content so that only the holder of the correct key can read
it.

\subsection*{POP3 and IMAP}
\textbf{POP3} stands for \textbf{Post Office Protocol version 3}. It is a mail
\emph{retrieval} protocol: it \textbf{downloads} messages from the mail server
to the local device. By default it typically \textbf{removes} the messages
from the server, so once downloaded the mail lives only on that machine and is
not kept in sync across devices. \textbf{IMAP} works differently: it leaves
the mail on the server and keeps all devices synchronised. The stable
distinction (PYQ-16, TRAP-17) is that POP3 downloads to the local machine,
whereas IMAP syncs. \textbf{SMTP}, in contrast, is the protocol for
\emph{sending} mail.

\subsection*{BCC and recipient privacy}
An email has three recipient fields.

\begin{center}
\begin{tabular}{p{0.14\textwidth}p{0.22\textwidth}p{0.52\textwidth}}
\toprule
\textbf{Field} & \textbf{Full form} & \textbf{Visibility} \\
\midrule
To & --- & Primary recipients; visible to everyone. \\
Cc & Carbon Copy & Copied recipients; visible to everyone. \\
Bcc & Blind Carbon Copy & Addresses hidden from the \emph{other} recipients. \\
\bottomrule
\end{tabular}
\end{center}

The key point (PYQ-16, TRAP-16) is what Bcc hides. Bcc hides the blind
recipients from \emph{each other} and from the To/Cc recipients; it does
\emph{not} hide them from the \textbf{sender}, who always sees the complete
recipient list. A common wrong statement says that Bcc hides the sender and
recipients from one another --- the sender is never hidden.

\section{PowerPoint SmartArt conversion}
\textbf{SmartArt} turns a plain list into a diagram. PowerPoint can convert
content in both directions, with one firm limit.

\begin{itemize}
  \item \textbf{Text to SmartArt}: select a bulleted list and choose
    \textbf{Home} $\rightarrow$ Convert to SmartArt Graphic; choose a layout.
    The text is preserved as the diagram's content.
  \item \textbf{SmartArt to text}: with the graphic selected, use SmartArt
    Tools Design $\rightarrow$ Convert $\rightarrow$ Convert to Text. You get
    a bulleted list back; the shapes and any pictures in the graphic are
    \emph{lost}.
  \item \textbf{SmartArt to shapes}: Convert to Shapes breaks the graphic
    into ordinary shapes. Once that is done, the shapes \textbf{cannot} be
    converted back into SmartArt.
\end{itemize}

\begin{example}
A presenter has six bullet points and wants a cycle diagram. She selects the
list and uses \emph{Convert to SmartArt}. Later she needs the plain list
again, so she uses \emph{Convert to Text} and the six points return as
bullets. What cannot be undone is the step to individual shapes: after
\emph{Convert to Shapes}, rebuilding the SmartArt would mean starting over
(the ``re-entry'' point behind PYQ-17).
\end{example}

The direct source of Convert to SmartArt is \emph{text}. A picture or a chart
is not transformed into SmartArt in one step; you would instead insert a
picture-type SmartArt layout and place the images by hand.

\section{Exam retrieval and common confusions}
The following distinctions prevent most errors in this unit.
\begin{enumerate}
  \item Different headers and footers require a \textbf{Section Break}, not a
    Page Break (TRAP-13).
  \item A \textbf{hanging indent} shifts every line except the first; a
    first-line indent shifts only the first line.
  \item A \textbf{blank cell} is treated as 0 in a comparison; it does not
    cause an error (PYQ-14).
  \item \textbf{COUNTIFS} counts with AND logic over paired arguments;
    \textbf{SUMPRODUCT} multiplies arrays and can both add and count
    (TRAP-14).
  \item A \textbf{digital signature} gives authenticity and integrity, not
    confidentiality; \textbf{encryption} gives confidentiality (TRAP-15).
  \item \textbf{POP3 downloads} to the local machine; \textbf{IMAP} syncs
    (TRAP-17).
  \item \textbf{Bcc} hides recipients from each other, never from the sender
    (TRAP-16).
  \item \textbf{Navigation Pane} shows headings, pages, and search results,
    not comments.
\end{enumerate}

\begin{example}
An item states, ``A digital signature encrypts the message so that only the
recipient can read it.'' This is wrong. The signature establishes the sender
and integrity; encryption is the feature that hides the content. The wrong
option has merged two distinct security services.
\end{example}

\subsection*{Exercises}
\begin{enumerate}[label=\arabic*.]
  \item Distinguish a first-line indent from a hanging indent, with one use
    for each.
  \item State why a Section Break, and not a Page Break, is needed for
    different headers and footers.
  \item Explain what a nested IF returns when several conditions are true at
    once.
  \item What does \texttt{=IF(A1>50,"High","Low")} display when A1 is blank,
    and why?
  \item Give the general form of COUNTIFS, and state one thing SUMPRODUCT can
    do that COUNTIFS cannot.
  \item Distinguish a digital signature from encryption in one sentence.
  \item What does POP3 do with mail, and how does that differ from IMAP?
  \item In a message, who can see the addresses placed in the Bcc field?
  \item Describe two directions of SmartArt conversion and one that cannot be
    reversed.
\end{enumerate}

\subsection*{Solutions}
\begingroup\small
\begin{enumerate}[label=\arabic*.]
  \item A first-line indent indents only the first line (normal paragraph
    openings); a hanging indent indents every line except the first
    (references and bulleted lists).
  \item Page-level settings such as headers and footers belong to a section.
    A Page Break keeps the same section, so its header/footer cannot differ;
    a Section Break starts a new section that can carry its own header and
    footer.
  \item A nested IF returns the value of the \emph{first} condition that is
    TRUE; once a condition is satisfied the remaining tests are skipped.
  \item It displays \emph{Low}, because a blank cell is treated as 0, and 0 is
    not greater than 50, so the FALSE branch is taken.
  \item \texttt{COUNTIFS(range1, criteria1, range2, criteria2, \ldots)}.
    SUMPRODUCT can multiply corresponding array elements and return their sum
    (arithmetic), so it can add as well as count; COUNTIFS only counts.
  \item A digital signature proves the sender's identity and that the message
    was not altered, whereas encryption hides the content so only the
    intended reader can read it.
  \item POP3 downloads mail from the server to the local device and by
    default removes it from the server; IMAP leaves mail on the server and
    keeps devices synchronised.
  \item The sender can see all Bcc addresses, and each Bcc recipient is hidden
    from the other recipients (To, Cc, and the other Bcc recipients).
  \item Text converts to SmartArt (Convert to SmartArt), and SmartArt
    converts back to text (Convert to Text). Converting SmartArt to shapes is
    \emph{not} reversible --- the shapes cannot be turned back into SmartArt.
\end{enumerate}
\endgroup

\unitopener{8}{Internet, Web \& Email}
\lessonheading{34}{Internet, WWW and Browser}
\section{The Internet: the network itself}
The Internet is a global network of connected computers and networks. It is
best described as a \textbf{network of networks}: smaller networks joined
together so that any connected device can reach any other. It carries many
services --- the World Wide Web, e-mail, file transfer, and more. The
Internet is the physical infrastructure underneath all of these: the cables,
routers, servers, and the devices connected to them.

\begin{definitionbox}{Internet}
The global network of interconnected computers and networks over which
services such as the Web and e-mail run.
\end{definitionbox}

\begin{center}
\begin{tabular}{p{0.24\textwidth}p{0.66\textwidth}}
\toprule
\textbf{Term} & \textbf{Meaning in this lesson} \\
\midrule
Internet & The global network of networks; the infrastructure. \\
WWW (World Wide Web) & A service on the Internet: linked web pages and sites. \\
Website & A collection of related webpages under one address. \\
Webpage & A single document on the Web, shown in a browser. \\
Homepage & The main or starting page of a website. \\
Browser & Software that opens and displays websites and webpages. \\
Search engine & A website that finds other pages matching typed keywords. \\
Client & The device or program that requests a service or resource. \\
Server & The computer or program that supplies services to clients. \\
ISP (Internet Service Provider) & The company or organisation that provides the Internet connection. \\
Hyperlink & Clickable text or image that jumps to another page or section. \\
URL (Uniform Resource Locator) & The address of a webpage or resource on the Web. \\
\bottomrule
\end{tabular}
\end{center}

The rest of this lesson uses these terms in one consistent sense. The
\textbf{Internet} is the network. The \textbf{WWW} is the linked-pages
service running on that network. A \textbf{webpage} is one document, a
\textbf{website} is the collection, and the \textbf{homepage} is the
collection's front door. A \textbf{browser} opens pages; a \textbf{search
engine} finds pages. A \textbf{client} requests and a \textbf{server}
supplies. An \textbf{ISP} provides the connection itself. A
\textbf{hyperlink} points to another page, and a \textbf{URL} is a page's
address.

\section{Internet vs WWW}
The single most tested confusion in this lesson is Internet versus WWW. The
Internet is the whole global network; the WWW --- World Wide Web --- is one
service that runs on the Internet, namely the system of linked web pages and
sites. E-mail is another service on the same Internet, and it is not part of
the Web. Access follows the same layering: a device reaches the Internet
through an ISP connection, and it reaches the Web through a browser running
over that Internet connection.

\begin{definitionbox}{World Wide Web (WWW)}
The system of linked web pages and websites accessed over the Internet
through a browser; one service among several that the Internet carries.
\end{definitionbox}

\begin{example}
An item states, ``The Internet and the World Wide Web are the same thing.''
This is wrong. The Internet is the global network of networks, while the WWW
is one service running on it. The wrong option has collapsed the
infrastructure and one of its services into a single term.
\end{example}

\section{Website, webpage, and homepage}
A \textbf{webpage} is a single document on the Web, shown inside a browser.
A \textbf{website} is a collection of related webpages grouped under one
address (domain). The \textbf{homepage} is the main or starting page of a
website --- the page that loads when the site's base address is typed. In
order of size: the homepage is one page, any single document is a webpage,
and the whole collection is the website.

\begin{definitionbox}{Website, webpage, homepage}
A webpage is one document on the Web; a website is the collection of related
webpages under one address; the homepage is the website's main or starting
page.
\end{definitionbox}

Each webpage of a website has its own URL, and the homepage URL is the
site's base address. Typing that address in the browser's address bar loads
the homepage first, and hyperlinks on it lead to the site's remaining pages.

\section{Browser}
A \textbf{browser} is software used to open and view websites and webpages.
It sends a request for a page and displays the page that comes back. Common
browsers are Google Chrome, Mozilla Firefox, Microsoft Edge, and Apple
Safari. A browser is software installed on the user's device; it is not the
Internet itself, and it is not a search engine.

\begin{definitionbox}{Browser}
Software that requests webpages and displays them; examples are Google
Chrome, Mozilla Firefox, Microsoft Edge, and Apple Safari.
\end{definitionbox}

\section{Search engine}
A \textbf{search engine} is a website (a service) that finds other webpages
matching keywords typed by the user. Examples are Google, Bing, and Yahoo.
It is reached \emph{through} a browser: the browser is the tool open on the
device, and the search-engine site is loaded inside it. Searching therefore
needs both layers --- a browser open on the device, and a search-engine site
loaded in that browser.

\begin{definitionbox}{Search engine}
A website that returns the pages matching a user's keywords; examples are
Google, Bing, and Yahoo.
\end{definitionbox}

The browser-versus-search-engine table is worth memorising as a whole. A
browser is software (an application) whose role is to open and display pages;
a search engine is a website whose role is to find pages for keywords. Chrome,
Firefox, Edge, and Safari belong in the first column; Google, Bing, and
Yahoo belong in the second. The classic trap swaps them: it calls Chrome a
search engine or Google a browser.

\section{Client and server}
A \textbf{client} is the device or program that requests a service or
resource. The office PC's browser asking for a webpage is acting as a
client: it asks, and it does not supply the page itself. A \textbf{server}
is the computer (or program) that supplies services or resources to clients.
A web server stores websites and sends pages to the browsers that ask for
them. One server can serve many clients at the same time, and servers stay
switched on and connected so that clients can reach them at any time.

\begin{definitionbox}{Client and server}
The client requests a service or resource; the server supplies it. On the
Web, the browser is the client and the web server holds the pages.
\end{definitionbox}

The two always appear as a pair, in three steps. First, the client (the
browser) sends a request for a page. Second, the server receives the request
and sends the page back. Third, the browser displays the received page. This
request-and-response pair is the client-server model of the Web: typing a
site address and pressing Enter sends a request, the site's server returns
the homepage, and the browser shows it.

\section{ISP: the connection provider}
\textbf{ISP} stands for Internet Service Provider. It is the company or
organisation that provides the Internet connection to homes and offices.
Without an ISP connection, the device cannot reach the Internet at all, so
the ISP link comes before everything else: no ISP means no browsing, no
search, and no e-mail. This is the tested fact of WO 2026 Q74 (PYQ-34): ISP
expands to Internet Service Provider, and no other expansion is accepted.

\begin{definitionbox}{ISP (Internet Service Provider)}
The company or organisation that provides the Internet connection to homes
and offices.
\end{definitionbox}

\section{Hyperlink and URL}
A \textbf{hyperlink} (link) is clickable text or an image that jumps to
another page or section. Clicking it loads the linked webpage, which may sit
on the same site or on a different one. Links are what make the Web a
``web'': pages joined to pages. Note the division of labour: a hyperlink
points the way, and the browser follows it when clicked.

A \textbf{URL} --- Uniform Resource Locator --- is the address of a webpage
or resource on the Web. Typing the site's URL in the browser's address bar
loads its homepage. The URL names the page; the hyperlink jumps to it.

\begin{definitionbox}{Hyperlink and URL}
A hyperlink is clickable text or an image that leads to another page; a URL
(Uniform Resource Locator) is the address of a webpage or resource.
\end{definitionbox}

\section{Full-form ladder of this lesson}
Direct-recall full-form items are among the most common JKSSB computer
formats, so learn each expansion exactly and give the full form on first
mention. \textbf{ISP} --- Internet Service Provider --- gives the Internet
connection. \textbf{WWW} --- World Wide Web --- is the linked system of
pages and sites. \textbf{URL} --- Uniform Resource Locator --- is the address
of a page or resource. \textbf{HTTP} --- HyperText Transfer Protocol --- is
the set of rules browsers and servers follow when requesting and delivering
pages. \textbf{HTML} --- HyperText Markup Language --- is the language in
which webpages are written.

\section{Exam retrieval and common confusions}
Five distinctions prevent most errors on this topic.
\begin{enumerate}
  \item The Internet is \textbf{not} the Web; the WWW is \textbf{one
    service} running on the Internet.
  \item A browser is \textbf{not} a search engine. Chrome is a
    \textbf{browser}; Google is a \textbf{search engine}.
  \item A website is the \textbf{collection}; a webpage is \textbf{one
    document}; the homepage is the \textbf{first page}.
  \item The client \textbf{requests} and the server \textbf{supplies}. The
    ISP --- Internet Service Provider --- gives the connection (PYQ-34).
  \item A hyperlink \textbf{points}; the browser \textbf{opens}. A URL is the
    page's address.
\end{enumerate}

\begin{example}
An item asks which software is used to view websites. Trace the roles: the
browser opens and displays pages, while the search engine finds pages and
the ISP provides the connection. The answer is the browser (for example
Chrome). If the stem instead asks which site finds pages for keywords, the
answer is the search engine (for example Google), and if it asks who
provides the connection, the answer is the ISP.
\end{example}

\subsection*{Exercises}
\begin{enumerate}[label=\arabic*.]
  \item Distinguish the Internet from the World Wide Web.
  \item Distinguish a website from a webpage and from a homepage.
  \item Distinguish a browser from a search engine, giving two examples of
    each.
  \item Explain the roles of client, server, and ISP when a page is opened.
  \item Define hyperlink and URL, and explain how the two work together.
\end{enumerate}

\subsection*{Solutions}
\begingroup\small
\begin{enumerate}[label=\arabic*.]
  \item The Internet is the global network of networks --- the
    infrastructure. The WWW (World Wide Web) is one service running on that
    infrastructure: the system of linked web pages and sites.
  \item A webpage is one document on the Web; a website is the collection of
    related webpages under one address; the homepage is the website's main or
    starting page.
  \item A browser is software that opens and displays pages (Chrome, Firefox,
    Edge, Safari). A search engine is a website that finds pages for keywords
    (Google, Bing, Yahoo); it is reached through a browser.
  \item The client (browser) requests the page, the server (web server)
    supplies it, and the ISP (Internet Service Provider) provides the
    Internet connection over which the request and response travel.
  \item A hyperlink is clickable text or an image that jumps to another page;
    a URL (Uniform Resource Locator) is a page's address. Clicking the
    hyperlink tells the browser to load the page named by that URL.
\end{enumerate}
\endgroup

\lessonheading{35}{URL, DNS, HTTP and HTTPS}
\section{What a URL is}
A URL (Uniform Resource Locator) is the full address of one resource on the
Internet: one page, file, image, or service. Typing a URL tells the browser
both how to connect and exactly what to fetch.

\begin{definitionbox}{URL}
The complete address of a single Internet resource, naming the protocol
(scheme), the host, and the location of the resource on that host,
including the optional port, path, and query.
\end{definitionbox}

\begin{center}
\begin{tabular}{p{0.24\textwidth}p{0.66\textwidth}}
\toprule
\textbf{Term} & \textbf{Meaning in this lesson} \\
\midrule
URL & The full address of one resource, such as a page or file. \\
Scheme & The protocol word before the colon, such as https or http. \\
Domain name & The readable host name, such as www.example.in. \\
TLD & Top-Level Domain: the last label, such as .in or .com. \\
Subdomain & A prefix section of the name, such as www or mail. \\
Port & The door number on the server (80 for HTTP, 443 for HTTPS). \\
Path & The location of the page inside the site. \\
Query & Parameters after a ?, such as ?topic=dns. \\
IP address & The numeric address of a device (IPv4 is 32-bit, IPv6 is 128-bit). \\
DNS & Domain Name System: translates domain names into IP addresses. \\
HTTP & Hypertext Transfer Protocol: carries web pages. \\
HTTPS & Hypertext Transfer Protocol Secure: HTTP protected by TLS. \\
TLS & Transport Layer Security: encrypts the exchange (successor of SSL). \\
\bottomrule
\end{tabular}
\end{center}

The rest of this lesson uses these terms in one consistent sense. A
\textbf{URL} is the whole address; the \textbf{domain name} is only the host
part inside it. \textbf{DNS} translates that host name into an \textbf{IP
address}. \textbf{HTTP} moves the page, and \textbf{HTTPS} is HTTP carried
inside a \textbf{TLS}-encrypted channel.

\section{Anatomy of a URL}
Consider the example used throughout this lesson:
\begin{center}
\texttt{https://www.example.in:443/study/computer?topic=dns\&lang=en}
\end{center}
Read it left to right. The \textbf{scheme} is \texttt{https}: the protocol
the browser must speak. The \texttt{://} is only a separator between the
scheme and the host. The host is \texttt{www.example.in}, in which
\texttt{www} is the \textbf{subdomain}, \texttt{example} is the registered
domain, and \texttt{.in} is the \textbf{TLD} (Top-Level Domain). Common TLDs
include \texttt{.com}, \texttt{.org}, \texttt{.edu}, \texttt{.gov}, and
country labels such as \texttt{.in}.

After the host comes the optional \textbf{port}, \texttt{:443}: the door
number on the server. The \textbf{path}, \texttt{/study/computer}, locates
the page inside the site. The \textbf{query}, \texttt{?topic=dns\&lang=en},
passes extra parameters to the server. Port and query are optional: when the
port is missing, the browser fills in the default for the scheme.

\begin{definitionbox}{TLD (Top-Level Domain)}
The last label of a domain name, such as .com, .org, .edu, .gov, or a
country label such as .in.
\end{definitionbox}

\begin{example}
Split \texttt{mail.example.com}. The subdomain is \texttt{mail}, the
registered domain is \texttt{example}, and the TLD is \texttt{.com}. One
registered domain can carry many such subdomains (\texttt{www},
\texttt{mail}, \texttt{blog}), each pointing to a different service, while
the domain itself stays the same.
\end{example}

\section{Ports, paths, and queries}
The \textbf{port} selects which service on the server receives the
connection. Port \textbf{80} is the default for HTTP and port \textbf{443}
is the default for HTTPS, which is why both numbers are usually hidden in
the address bar. The \textbf{path} names the page or file within the site,
and the \textbf{query} carries name--value pairs after a \texttt{?} mark.
Neither the port nor the query changes which host is contacted; the host is
fixed by the domain name alone.

\section{IP addresses: the real destination}
An IP (Internet Protocol) address is the numeric address every device uses
to send and receive data on the network. People type names; machines route
by numbers.

\begin{definitionbox}{IP address}
The numeric address of a device on a TCP/IP network, used for routing data
to that device.
\end{definitionbox}

\textbf{IPv4} (Internet Protocol version 4) is \textbf{32-bit} and is
written as four decimal groups, for example \texttt{192.0.2.44}.
\textbf{IPv6} (Internet Protocol version 6) is \textbf{128-bit} and is
written in hexadecimal groups. It is never 256-bit; that figure is a
repeated wrong option. IPv6 exists because the IPv4 space is too small for
the number of connected devices.

\section{DNS: the phone book of the Internet}
DNS (Domain Name System) is the distributed naming system that translates
domain names into IP addresses. It does not carry page content and does not
assign addresses to devices; it only answers the question ``which IP address
belongs to this name?''

\begin{definitionbox}{DNS}
The hierarchical, distributed system that maps each domain name to the IP
address of its server, with answers cached to speed up repeat visits.
\end{definitionbox}

Resolution runs as a short chain. The browser first checks its own cache
for the name. If the name is unknown, the \textbf{resolver} (usually run by
the ISP --- Internet Service Provider) asks the DNS hierarchy: the root
servers, then the TLD servers for labels such as \texttt{.in}, then the
domain's own authoritative server. That authoritative server returns the IP
address, the resolver caches the answer for a fixed time, and the browser
receives the number it can connect to.

\begin{example}
On the office desktop, typing \texttt{www.example.in} first triggers a DNS
lookup. The resolver walks the hierarchy, learns the address (say
\texttt{192.0.2.44}), caches it, and only then does the browser open a
connection to that number. Typing the IP address directly would skip DNS,
but nobody memorises numbers for every site.
\end{example}

\section{HTTP: carrying web pages}
HTTP (Hypertext Transfer Protocol) is the request--response protocol of the
web. The browser sends a request naming the path (and query); the server
returns the page. HTTP runs by default over port 80. Plain HTTP provides no
encryption: the request and the page travel in readable form, so anyone on
the path between browser and server can read them.

\begin{definitionbox}{HTTP}
The protocol by which a web client requests a resource and the server
returns it. Plain HTTP does not encrypt the exchange.
\end{definitionbox}

\section{HTTPS and TLS: the secure form}
HTTPS (Hypertext Transfer Protocol Secure) is HTTP protected by TLS. TLS
(Transport Layer Security) is the successor of the older SSL (Secure
Sockets Layer); the phrase ``SSL/TLS'' survives only as a name for the same
family of protection. Before any page moves, the browser and server run a
TLS handshake: the server proves its identity with a certificate, and the
two sides agree on encryption keys. The HTTP request and response then
travel inside that encrypted channel.

\begin{definitionbox}{HTTPS}
HTTP carried inside a TLS-encrypted channel, by default over port 443 and
shown with a padlock in the address bar.
\end{definitionbox}

TLS provides three things at once: \textbf{confidentiality} through
encryption (only the two ends can read the data), \textbf{integrity}
(tampering is detected), and server \textbf{authentication} (the certificate
ties the site to its claimed name). HTTPS therefore protects the trip; it
is not a different page language.

\section{The full flow: from typed URL to secure page}
Put the pieces together for \texttt{https://www.example.in/study/computer}.
The browser parses the URL: scheme \texttt{https}, host
\texttt{www.example.in}, path \texttt{/study/computer}. DNS resolves the
host to an IP address. The browser opens a connection to that IP on port
443, the HTTPS default. The TLS handshake secures the channel, and the HTTP
request for the path travels inside it. The server returns the page, and the
browser renders it with the padlock shown.

\section{Exam retrieval and common confusions}
Six distinctions prevent most errors on this topic.
\begin{enumerate}
  \item A \textbf{URL} addresses one resource; the \textbf{domain name} is
    only the host part inside it.
  \item \textbf{DNS} translates names into IP addresses. It does not assign
    addresses and does not carry page content.
  \item \textbf{IPv4} is 32-bit; \textbf{IPv6} is 128-bit, never 256-bit.
  \item \textbf{HTTP} uses port 80; \textbf{HTTPS} uses port 443. Do not
    reverse them.
  \item \textbf{HTTPS} is HTTP over \textbf{TLS}. HTTP alone gives no
    encryption.
  \item The \textbf{scheme} is the protocol word (\texttt{https}), not the
    domain name that follows it.
\end{enumerate}

\begin{example}
An item asks which part of \texttt{https://mail.example.com/inbox} is the
subdomain. Trace the host \texttt{mail.example.com} from the right: the TLD
is \texttt{.com}, the domain is \texttt{example}, and the remaining prefix
\texttt{mail} is the subdomain. If the stem instead asks for the scheme,
the answer is \texttt{https}; if it asks for the default port of that
scheme, the answer is 443.
\end{example}

\subsection*{Exercises}
\begin{enumerate}[label=\arabic*.]
  \item Expand URL, DNS, HTTP, HTTPS, TLS, IP, TLD, and ISP.
  \item Label every part of \texttt{https://blog.example.org:443/notes/net?ch=1}.
  \item Distinguish a URL from a domain name.
  \item Describe the steps of DNS resolution when a name is not cached.
  \item Contrast HTTP with HTTPS on security, default port, and address-bar
    form.
\end{enumerate}

\subsection*{Solutions}
\begingroup\small
\begin{enumerate}[label=\arabic*.]
  \item URL --- Uniform Resource Locator; DNS --- Domain Name System; HTTP
    --- Hypertext Transfer Protocol; HTTPS --- Hypertext Transfer Protocol
    Secure; TLS --- Transport Layer Security; IP --- Internet Protocol; TLD
    --- Top-Level Domain; ISP --- Internet Service Provider.
  \item Scheme \texttt{https}; subdomain \texttt{blog}; domain
    \texttt{example}; TLD \texttt{.org}; port \texttt{443}; path
    \texttt{/notes/net}; query \texttt{?ch=1}.
  \item A URL is the full address of one resource (scheme, host, port,
    path, query). A domain name is only the readable host part inside it.
  \item The browser checks its cache; then the resolver asks the DNS
    hierarchy (root, then TLD, then the authoritative server), which returns
    the IP address; the resolver caches it and hands it to the browser.
  \item HTTP carries pages without encryption over port 80
    (\texttt{http://}); HTTPS is HTTP protected by TLS over port 443
    (\texttt{https://} with a padlock).
\end{enumerate}
\endgroup

\lessonheading{36}{Web Browsing, Search, Upload and Download}
\section{The browser, tabs, bookmarks, and history}
A \textbf{web browser} is the software used to open and view web pages on the
internet. Common examples are Google Chrome, Mozilla Firefox, and Microsoft
Edge. The browser sends a request for a page and then displays the returned
text, images, and links. The address typed to reach a page is the \textbf{URL}
--- Uniform Resource Locator --- and it is typed in the \textbf{address bar}
at the top of the browser window.

\begin{definitionbox}{Web browser}
Software that requests web pages from servers and displays them to the user.
\end{definitionbox}

\begin{center}
\begin{tabular}{p{0.24\textwidth}p{0.66\textwidth}}
\toprule
\textbf{Term} & \textbf{Meaning in this lesson} \\
\midrule
Web browser & Software (Chrome, Firefox, Edge) that opens and displays web pages. \\
Tab & One open page inside the browser window; many tabs can stay open together. \\
Bookmark & A saved shortcut to a page for quick return without retyping the address. \\
Browsing history & The date-wise list of pages visited earlier, kept by the browser. \\
URL & Uniform Resource Locator: the address of a page, typed in the address bar. \\
Search engine & A website (Google, Bing) that finds pages matching typed words. \\
\bottomrule
\end{tabular}
\end{center}

Each open page lives in its own \textbf{tab}, so several pages can stay open
in one window. A new tab is opened with \textbf{Ctrl+T}. A \textbf{bookmark}
is a saved shortcut to a page: instead of typing the address again, the saved
bookmark is clicked. It is created with \textbf{Ctrl+D}. \textbf{Browsing
history} is the list of pages visited earlier, stored by date, and it is
opened with \textbf{Ctrl+H}. The distinction is simple: history answers
``where was I'', while a bookmark answers ``where do I want to return''.

\begin{example}
While preparing a report on the office desktop, the clerk keeps the office
portal open in one tab and the search results in a second tab, and bookmarks
the portal page so it can be opened directly the next morning. The history
list that evening shows both pages in order.
\end{example}

\section{How a search engine works}
A \textbf{search engine} is a website, such as Google or Bing, that finds
pages matching the words typed by the user. It works in three ordered stages:
crawling, indexing, and ranking.

\begin{definitionbox}{Crawling}
The collection step: automatic programs called crawlers or spiders visit
pages and follow their links to discover new and changed pages.
\end{definitionbox}

\begin{definitionbox}{Indexing}
The storage step: the crawled pages are kept in a large index, like the
index at the back of a book, so matches can be found quickly.
\end{definitionbox}

\begin{definitionbox}{Ranking}
The ordering step: when words are typed, matching pages are picked from the
index and ordered by relevance on the search results page.
\end{definitionbox}

Crawling comes first because without it the engine would never learn that a
new page exists. Indexing comes second: the collected pages are stored in the
index before any search is typed. Ranking comes last: only after the user
types words does the engine select matches from that ready index and order
them. The anchor to remember is the order itself: \textbf{crawl, then index,
then rank}.

\section{Sponsored results and organic results}
The matches chosen by relevance are called \textbf{organic results}. In
addition, the results page shows \textbf{sponsored results}: paid placements
that appear because an advertiser paid for them, not because they ranked
highest on relevance. A sponsored result always carries a visible
\textbf{Ad} or \textbf{Sponsored} label. The practical consequence is direct:
being listed first on the page does not always mean being the most relevant
match, since the top position may be a paid placement.

\begin{example}
A search for a printer shows two blocks: the first carries a small ``Ad''
label and links to a seller who paid for the position, while the entries
below it are the organic results ordered by relevance. The topmost entry is
therefore not necessarily the best match.
\end{example}

\section{Cookies, cache, and history}
These three stores are the most confused triple in this lesson. Each stores a
different thing, and each is stored by a different party.

\begin{definitionbox}{Cookie}
A small file a website stores on the computer to remember that visitor:
login state, cart items, and preferences such as language.
\end{definitionbox}

\begin{definitionbox}{Browser cache}
Copies of page parts (images, styles, scripts) kept by the browser on the
computer so repeat visits load faster.
\end{definitionbox}

\begin{definitionbox}{Browsing history}
The list of visited page addresses, kept by the browser by date.
\end{definitionbox}

\begin{center}
\begin{tabular}{p{0.18\textwidth}p{0.36\textwidth}p{0.36\textwidth}}
\toprule
\textbf{Item} & \textbf{What it stores} & \textbf{Who stores it} \\
\midrule
Cookie & Visitor data: login, cart, choices & The website \\
Cache & Copies of page files: images, styles & The browser \\
History & List of visited page addresses & The browser \\
\bottomrule
\end{tabular}
\end{center}

Cookies are set by websites and read back on the next visit, which is why a
user stays logged in. The cache is managed by the browser: on the next visit
the saved copies load from the disk instead of being fetched again, so the
page opens faster. Clearing them has different effects. Deleting
\textbf{cookies} signs the user out of sites and clears saved preferences.
Clearing the \textbf{cache} does not sign the user out; it only makes the
next visit slightly slower. Clearing \textbf{history} removes the record of
visited addresses but changes neither logins nor saved page copies.

\section{Downloading and uploading}
\textbf{Downloading} means copying a file from the internet to the user's
computer: saving a form PDF from the office portal, or saving a photo from a
web page. \textbf{Uploading} means sending a file from the user's computer to
the internet: attaching a document to an email, or submitting a scanned form
on a portal. The direction is always judged from the user's machine: download
moves data \textbf{down} to the machine, upload moves data \textbf{up} from
the machine to a server.

\begin{definitionbox}{Downloading}
Copying a file from the internet (a server) to the user's computer.
\end{definitionbox}

\begin{definitionbox}{Uploading}
Sending a file from the user's computer to the internet (a server).
\end{definitionbox}

\begin{center}
\begin{tabular}{p{0.22\textwidth}p{0.34\textwidth}p{0.34\textwidth}}
\toprule
\textbf{Point} & \textbf{Download} & \textbf{Upload} \\
\midrule
Direction & Server to computer & Computer to server \\
Typical use & Opening pages, saving files & Mail attachments, form submission \\
Typical speed & Faster & Slower \\
\bottomrule
\end{tabular}
\end{center}

Home and office connections usually provide a faster download speed than
upload speed, because receiving pages and files is the more common need. A
downloaded file, once saved, opens without needing the internet; an upload
needs the connection at the moment of sending.

\section{File conversion}
\textbf{File conversion} means changing a file from one format to another
while keeping the same content. The common office need is converting a Word
document (\textbf{.docx}) to \textbf{PDF} before submitting it on a portal.
Conversion changes the \textbf{format}, not the transfer direction: a
converted file must still be uploaded or downloaded separately. The
\textbf{extension} (such as .pdf, .docx, .jpg) is the name tag that says
which format the file is stored in, while the \textbf{format} (such as PDF,
PNG) is the storage method itself. The two words must not be conflated: the
extension names the format; it is not the format.

\begin{example}
The clerk finishes the report in Word, converts the .docx file to PDF, and
then uploads the PDF on the portal. The conversion changed the format; the
upload changed the location. If the portal instead needed a photograph, a
.JPG image would be converted to the needed picture format the same way.
\end{example}

\section{Exam retrieval and common confusions}
Five distinctions prevent most errors on this topic.
\begin{enumerate}
  \item The search order is \textbf{crawl}, then \textbf{index}, then
    \textbf{rank}. The index is built before the search is typed.
  \item \textbf{Sponsored} means a paid placement with an Ad label;
    \textbf{organic} means ranked by relevance. The top of the page is not
    always the best match.
  \item \textbf{Cookie} (website; visitor data) vs \textbf{cache} (browser;
    page copies) vs \textbf{history} (browser; visited list). Do not swap
    the three.
  \item \textbf{Download} means internet to computer; \textbf{upload} means
    computer to internet. Judge the direction from the user's machine.
  \item \textbf{Conversion} changes the format; \textbf{transfer} changes
    the location. The extension (.png) names the format (PNG); it is not
    the format itself.
\end{enumerate}

\subsection*{Exercises}
\begin{enumerate}[label=\arabic*.]
  \item List the three stages of a search engine in order and say what each stage does.
  \item Distinguish a sponsored result from an organic result.
  \item Distinguish a cookie from the browser cache and from browsing history.
  \item Explain the difference between downloading and uploading, giving one office example of each.
  \item What does file conversion change, and how is the extension related to the format?
\end{enumerate}

\subsection*{Solutions}
\begingroup\small
\begin{enumerate}[label=\arabic*.]
  \item Crawling (crawlers visit pages and follow links to collect them),
    then indexing (pages stored in a large index), then ranking (matches
    picked from the index and ordered by relevance after words are typed).
  \item A sponsored result is a paid placement carrying an Ad or Sponsored
    label; an organic result is a normal match chosen and ordered by
    relevance.
  \item A cookie is visitor data (login, cart, choices) stored by the
    website; the cache is copies of page files (images, styles) stored by
    the browser for faster repeat visits; history is the browser's
    date-wise list of visited addresses.
  \item Downloading copies a file from the internet to the computer, such as
    saving a form PDF from the portal; uploading sends a file from the
    computer to the internet, such as attaching a document to an email.
  \item Conversion changes the format while keeping the content (for example
    .docx to PDF). The extension (.pdf, .docx) is the name tag for the
    format; the format is the storage method itself.
\end{enumerate}
\endgroup

\lessonheading{37}{Email Basics}
\section{What email is}
Electronic mail (email) is a message sent from one electronic mailbox to
another over a network such as the Internet. Each user owns a mailbox
identified by a unique email address. The sender composes a message, directs
it with the recipient fields, and sends it; the message then waits in each
recipient's mailbox until it is opened. Because delivery does not require both
parties to be present at the same time, email is called a store-and-forward
service.

\begin{definitionbox}{Email address}
A unique identifier of a mailbox, written as a user name, then the @ symbol,
then the domain name, for example \texttt{assistant@gmail.com}.
\end{definitionbox}

\begin{center}
\begin{tabular}{p{0.24\textwidth}p{0.66\textwidth}}
\toprule
\textbf{Term} & \textbf{Meaning in this lesson} \\
\midrule
User name & The part of the address left of the @; the mailbox on the server. \\
Domain name & The part of the address right of the @; the mail server holding the mailbox. \\
To / Cc / Bcc & Recipient fields: primary recipient, visible copy, hidden copy. \\
Subject / body & The topic line and the main written content. \\
Attachment & A file sent along with the message. \\
Inbox / Sent / Drafts & Folders for received, sent, and unfinished mail. \\
Spam / Trash & Folders for filtered junk and for deleted messages. \\
Webmail / email client & Mail used in a browser vs mail used in an installed app. \\
\bottomrule
\end{tabular}
\end{center}

The rest of this lesson uses these terms in one consistent sense. The
\textbf{user name} is whatever stands left of the @ and the \textbf{domain
name} is whatever stands right of it. \textbf{To} is the main recipient,
\textbf{Cc} (Carbon Copy) is a visible additional recipient, and \textbf{Bcc}
(Blind Carbon Copy) is a hidden additional recipient. \textbf{Subject} is the
topic line, \textbf{body} is the message text, and an \textbf{attachment} is a
separate file carried with the message.

\section{The email address}
Every email address contains exactly one @ symbol, read as ``at''. The @
separates the user name on the left from the domain name on the right. In
\texttt{assistant@gmail.com}, \texttt{assistant} is the user name and
\texttt{gmail.com} is the domain name. The domain name identifies which mail
server holds the mailbox; the user name identifies which mailbox on that
server. This is the point tested when an item asks what the @ separates: the
answer is the user name and the domain name (PYQ-35).

\begin{definitionbox}{Domain name}
The part of an email address right of the @ that identifies the mail server or
organisation holding the mailbox, such as \texttt{gmail.com}.
\end{definitionbox}

Addresses are written without spaces, and the @ appears exactly once.
Anything before it belongs to the user name; anything after it belongs to the
domain. Options that say the @ separates the domain and the host, or two
parts of the domain, are therefore wrong.

\section{Recipient fields: To, Cc and Bcc}
The three recipient fields decide who receives the message and who can see
whom. \textbf{To} names the primary recipient, the person the message is
mainly for. \textbf{Cc} (Carbon Copy) names additional recipients who receive
a visible copy; every recipient can see the To and Cc lists. \textbf{Bcc}
(Blind Carbon Copy) names additional recipients whose addresses are hidden
from the other recipients.

\begin{definitionbox}{Bcc}
A recipient field whose addresses are hidden from the other recipients of the
same message; each Bcc recipient is concealed from the To, Cc, and fellow Bcc
recipients.
\end{definitionbox}

The word ``blind'' refers to this concealment. A Bcc recipient still receives
the message normally. The other recipients simply cannot see that address.
The sender, who composed the message, always knows the full list. To and Cc
addresses, by contrast, are visible to everyone who receives the message.
This is the distinction behind the common trap: Bcc hides recipients from each
other, not from the sender, and it does not hide the To or Cc lists from
anyone (TRAP-16).

\begin{example}
The officer sends a notice To the assistant and Cc the clerk, with the head
office on Bcc. The assistant and the clerk see each other's addresses but
cannot see the head office address. The head office receives the same notice.
The sender (the officer) knows all three addresses.
\end{example}

For the examination, read every Bcc option word by word. An option that says
Bcc hides the sender is wrong; an option that says Bcc hides the To list from
the recipients is wrong; the correct statement is that Bcc recipients are
hidden from each other while To and Cc remain visible to all.

\section{Subject, body and attachments}
A message has two content parts and an optional file part. The
\textbf{subject} is a single short line stating the topic, such as ``Leave
application for 12 October''. It is displayed in the message list before the
message is opened, so a clear subject helps the recipient identify the
message. The \textbf{body} is the main written content: the greeting, the
message itself, and the closing.

An \textbf{attachment} is a file sent along with the message, usually shown
with a paperclip icon. Common attachments are documents, spreadsheets,
presentations, PDFs, and images. Attaching a file sends a copy to the
recipient; the sender keeps the original. Two practical points matter for the
examination. First, mail servers impose a size limit, so a very large file is
rejected or must be shared as a link instead. Second, an attachment can carry
a virus, so an unexpected attachment should not be opened blindly.

\section{Folders: Inbox, Sent, Drafts, Spam, Trash}
Received mail arrives in the \textbf{Inbox}. Every account has exactly one
Inbox; it is the default view when the mailbox is opened. A copy of every
message that is sent is kept in the \textbf{Sent} folder (sometimes labelled
Sent Items or Sent Mail). A message that has been started but not yet sent can
be saved in \textbf{Drafts}; it waits there until it is completed and sent.

\begin{definitionbox}{Drafts}
The folder holding composed messages that have been saved but not yet sent.
\end{definitionbox}

Unwanted or bulk mail is separated into the \textbf{Spam} folder (also called
Junk). The mail service filters suspected junk out of the Inbox automatically;
a wanted message found there should be moved back or marked ``not spam''.
Messages the user deletes go to the \textbf{Trash} folder (also called Deleted
Items). They remain recoverable until the folder is emptied; emptying Trash
or Spam permanently removes those messages.

The five folders are most easily confused in pairs. Inbox receives while Sent
keeps sent copies; Drafts holds work not yet sent while Sent holds work
already sent; Spam holds filtered junk while Trash holds what the user deleted.
No other folder does another's job: drafts never send themselves, sent copies
never return to drafts, and spam is not the delete bin.

\section{Reply, Reply All and Forward}
Three actions dispose of a received message. \textbf{Reply} composes an answer
addressed only to the sender. \textbf{Reply All} composes an answer addressed
to the sender plus all the To and Cc recipients of the original message.
\textbf{Forward} sends the received message onward to new recipients chosen
by the forwarder, usually with a short added note on top.

\begin{definitionbox}{Reply All}
A response addressed to the sender and to every visible (To and Cc)
recipient of the original message; Bcc recipients of the original are never
re-added.
\end{definitionbox}

Two consequences follow. First, Reply All can only reach the addresses visible
on the original message, so hidden Bcc addresses are never included. Second,
a forwarded message carries the original content (and often its attachments)
inside the new message; the forwarder becomes the new sender. For the
examination, match the verb to the action: ``answers only the sender'' is
Reply; ``answers everyone visible'' is Reply All; ``sends it onward to new
recipients'' is Forward.

\section{Webmail versus an email client}
Email can be used in two ways. \textbf{Webmail} means using mail through a
web browser, with the mailbox kept on the provider's server; Gmail, Yahoo
Mail, and Outlook.com are webmail services. An \textbf{email client} is an
installed application used for mail, such as Microsoft Outlook or Mozilla
Thunderbird; it can keep or download copies on the local machine.

Closely tied to the client are the incoming-mail protocols. \textbf{POP3}
(Post Office Protocol version 3) downloads mail to the local machine, usually
removing it from the server. \textbf{IMAP} (Internet Message Access Protocol)
keeps mail synced on the server, so a phone, an office computer, and a laptop
all show the same folders. Outgoing mail in both arrangements is normally
sent with \textbf{SMTP} (Simple Mail Transfer Protocol). The one-line
contrast to remember is: POP3 downloads, IMAP syncs, SMTP sends.

\begin{example}
The assistant reads office mail in two places: through Gmail in the browser
(webmail) and through the Outlook app on the office computer (email client)
configured with IMAP, so a message read on one device shows as read on the
other. Had the client used POP3, the downloaded message would sit on that one
computer instead of staying synced on the server.
\end{example}

\section{Exam retrieval and common confusions}
Six distinctions prevent most errors on this topic.
\begin{enumerate}
  \item The @ separates the \textbf{user name} (left) from the \textbf{domain
    name} (right). It does not separate domain and host (PYQ-35, TRAP-28).
  \item \textbf{Bcc} hides recipients from \textbf{each other}; the sender
    always knows, and To/Cc stay visible to all (TRAP-16).
  \item \textbf{Cc} is Carbon Copy (visible); \textbf{Bcc} is Blind Carbon Copy
    (hidden). Do not swap ``blind'' and ``visible''.
  \item \textbf{Reply} answers the sender; \textbf{Reply All} answers all
    visible recipients; \textbf{Forward} sends the message onward.
  \item \textbf{Drafts} holds unsent work, \textbf{Sent} holds sent copies,
    \textbf{Inbox} holds received mail, \textbf{Spam} holds filtered junk,
    and \textbf{Trash} holds deleted mail.
  \item \textbf{Webmail} runs in a browser; an \textbf{email client} is
    installed software. \textbf{POP3} downloads mail to the machine while
    \textbf{IMAP} keeps it synced on the server.
\end{enumerate}

\subsection*{Exercises}
\begin{enumerate}[label=\arabic*.]
  \item In the address \texttt{assistant@gmail.com}, identify the user name
    and the domain name. What does the @ separate?
  \item Distinguish To, Cc, and Bcc by stating who can see whom.
  \item Distinguish the subject, the body, and an attachment.
  \item State which folder holds each of the following: an unfinished message,
    a copy of a sent message, a newly arrived message, filtered junk, a
    deleted message.
  \item Distinguish Reply, Reply All, and Forward.
  \item Distinguish webmail from an email client, and POP3 from IMAP.
\end{enumerate}

\subsection*{Solutions}
\begingroup\small
\begin{enumerate}[label=\arabic*.]
  \item The user name is \texttt{assistant} and the domain name is
    \texttt{gmail.com}. The @ separates the user name (left) from the domain
    name (right).
  \item To is the primary recipient and Cc is a visible copy; both lists are
    seen by all recipients. Bcc is a hidden copy: each Bcc address is concealed
    from the other recipients, while the sender knows the full list.
  \item The subject is the short topic line shown in the message list; the
    body is the main written content; an attachment is a separate file carried
    with the message (shown with a paperclip icon).
  \item Unfinished: Drafts. Sent copy: Sent. Newly arrived: Inbox. Filtered
    junk: Spam (junk). Deleted: Trash (deleted items), recoverable until
    emptied.
  \item Reply answers only the sender; Reply All answers the sender plus all
    visible To and Cc recipients (never Bcc); Forward sends the received
    message onward to new recipients.
  \item Webmail is mail used through a browser with the mailbox on the
    server (Gmail, Yahoo Mail, Outlook.com); an email client is installed
    software (Outlook, Thunderbird). POP3 downloads mail to the local
    machine, usually removing it from the server; IMAP keeps mail synced on
    the server so all devices show the same folders.
\end{enumerate}
\endgroup

\lessonheading{38}{SMTP, POP and IMAP}
\section{Why email needs protocols}
Sending and reading mail involve different programs and different servers.
The sender writes a message in a mail program, that program hands it to a
mail server, the message travels between servers, waits in the recipient's
mailbox, and is finally collected by the recipient's program. For all these
independent programs to work together, they follow fixed rules called
\textbf{protocols}. One protocol sends mail out, and two other protocols
bring it in, each in its own way. Secure transmission and protection against
mail threats sit on top of these three protocols.

\begin{definitionbox}{Email protocol}
A fixed set of rules that a mail program and a mail server follow when
sending, forwarding, or collecting mail, so any standard program can work
with any standard server.
\end{definitionbox}

\begin{center}
\begin{tabular}{p{0.24\textwidth}p{0.66\textwidth}}
\toprule
\textbf{Term} & \textbf{Meaning in this lesson} \\
\midrule
SMTP & Simple Mail Transfer Protocol; the sending protocol. \\
POP3 & Post Office Protocol version 3; a receiving protocol that downloads mail to one machine. \\
IMAP & Internet Message Access Protocol; a receiving protocol that syncs mail across devices. \\
SSL/TLS & Secure Sockets Layer / Transport Layer Security; the secure channel for mail transmission. \\
Spam & Bulk unwanted mail, usually junk or advertising. \\
Phishing & A fake message impersonating a trusted sender to steal secrets. \\
Spoofing & Forging the sender address so mail looks as if it came from someone trusted. \\
\bottomrule
\end{tabular}
\end{center}

The rest of this lesson uses these terms in one consistent sense.
\textbf{SMTP} always means the sending step. \textbf{POP3} always means
downloading to one local machine. \textbf{IMAP} always means keeping every
device in step with the server copy. \textbf{SSL/TLS} always means the
secure transmission channel. \textbf{Spam}, \textbf{phishing}, and
\textbf{spoofing} are three distinct threats defined above.

\section{SMTP: the sending protocol}
\textbf{SMTP} (Simple Mail Transfer Protocol) is the protocol used to send
mail. When the sender presses send, the mail program hands the message to
the sender's mail server using SMTP, and the servers forward it toward the
recipient's mail server using SMTP. SMTP moves the message out of the
sender and along the path between servers until it reaches the recipient's
mailbox.

\begin{definitionbox}{SMTP}
The protocol that sends mail from the sender's program to the mail server
and forwards it between mail servers toward the recipient's mailbox.
\end{definitionbox}

Two limits matter for the exam. First, SMTP only \emph{sends}; it does not
deliver mail into the recipient's hands for reading. Second, the receiving
step that follows is the job of POP3 or IMAP, described next. Any option
that says SMTP downloads mail to the inbox, or that POP3 sends mail,
reverses the direction and is wrong.

\section{POP3: download to one machine}
\textbf{POP3} (Post Office Protocol version 3) is a receiving protocol.
It pulls mail from the mail server down to one local machine. After the
download, the message is typically removed from the server, so the only
copy lives on that one computer. Mail downloaded this way can be read
offline, but a second device will not see the same mailbox, because the
server copy is gone.

\begin{definitionbox}{POP3}
A receiving protocol that downloads mail from the server to one local
machine, usually removing it from the server.
\end{definitionbox}

POP3 suits a single computer where mail is collected and read in one place.
Its weakness is exactly its method: because the mail leaves the server,
bookmarks such as read marks, folders, and deletions do not travel to other
devices.

\begin{example}
A clerk collects office mail only on the office desktop. With POP3, each
morning's mail downloads to that desktop and is removed from the server.
The clerk can read it offline at the desk, but the same messages do not
appear on the clerk's phone, because the server copy is gone.
\end{example}

\section{IMAP: sync across devices}
\textbf{IMAP} (Internet Message Access Protocol) is also a receiving
protocol, but it works differently from POP3. The mail stays on the
server, and the mail program shows a synced view of it: folders, read and
unread marks, deletions, and moves are all mirrored between the device and
the server. When the reader acts on one device, every other device shows
the same change, because each of them reflects the same server copy.

\begin{definitionbox}{IMAP}
A receiving protocol that keeps mail on the server and syncs each device's
view --- folders, read marks, deletions --- with that server copy.
\end{definitionbox}

IMAP suits a reader with several devices, such as a phone and a desktop,
who must see the same inbox everywhere. Its cost is that the mailbox lives
on the server, so reading new mail needs a connection, and the server
mailbox must have room.

\section{POP3 vs IMAP and the full mail flow}
Both POP3 and IMAP receive mail; neither sends it. The difference is the
method. POP3 downloads to one machine and usually clears the server; IMAP
keeps the server copy and syncs every device with it.

\begin{center}
\begin{tabular}{p{0.20\textwidth}p{0.34\textwidth}p{0.36\textwidth}}
\toprule
\textbf{Point} & \textbf{POP3} & \textbf{IMAP} \\
\midrule
Job & Receives mail & Receives mail \\
Method & Downloads to one machine & Syncs a view of the server copy \\
Server copy & Usually removed & Kept on the server \\
Many devices & Poor fit & Designed for it \\
Offline reading & Natural (copy is local) & Needs a connection for new mail \\
\bottomrule
\end{tabular}
\end{center}

The full journey of one message combines the sending protocol with one of
the receiving protocols. The sender's program sends the message by SMTP to
the sender's mail server; that server forwards it by SMTP to the
recipient's mail server, where it waits in the mailbox; the recipient's
program then collects it with POP3 (download) or IMAP (sync). In one line:
sender, then SMTP, then the recipient server, then POP3 or IMAP, then the
reader. Sending is always SMTP; the receiving side is POP3 or IMAP.

\section{To, Cc, BCC, filters, and signatures}
The \textbf{To} field names the main recipients of a message, and
\textbf{Cc} (carbon copy) names others kept in the loop. Everyone who
receives the message can see both lists. \textbf{BCC} (blind carbon copy)
names recipients who are hidden from each other: nobody on the To or Cc
list sees who is on BCC, and BCC recipients do not see each other. The
sender always knows the full list; BCC hides recipients from each other,
not from the sender.

\textbf{Filters}, also called \textbf{rules}, sort incoming mail
automatically. A filter examines the headers and content of each message,
such as the sender, the subject, or words in the body, and then files,
flags, or deletes matching mail --- for example, moving all newsletters to
one folder. A mail \textbf{signature} is different: it is a fixed block of
sender details, such as name, post, and phone number, added at the end of
outgoing messages. It identifies the sender but does not encrypt or
otherwise protect the message.

\section{Digital signature vs encryption, and secure transmission}
A \textbf{digital signature} proves who sent a message and that it was not
changed on the way: it gives authenticity and integrity. \textbf{Encryption}
does something different: it scrambles the message so that only the
intended recipient can read it, giving confidentiality. A signed message
can still be read by anyone --- a signature is not secrecy. An encrypted
message can still be forwarded --- encryption hides the content but does
not by itself prove who sent it. Confusing the two is a tested trap
(TRAP-15): a signature never encrypts the body.

\begin{definitionbox}{Digital signature vs encryption}
A digital signature provides authenticity (who sent it) and integrity (it
was not altered); encryption provides confidentiality (only the intended
recipient can read it).
\end{definitionbox}

When mail travels between programs and servers, the channel itself can be
protected by \textbf{SSL} (Secure Sockets Layer) and its successor
\textbf{TLS} (Transport Layer Security). Together \textbf{SSL/TLS} provide
the secure channel over which mail is transmitted, so that others on the
path cannot read it. For the exam, the answer is direct: the secure email
transmission protocol is SSL/TLS (PYQ-36). The official Junior Assistant
item behind PYQ-16 separately confirms the POP3, BCC, filter, and
signature-versus-encryption facts used in this lesson.

\begin{example}
An item states, ``A digitally signed email can be read only by its
recipient.'' This is wrong. The signature proves the sender and detects
tampering, but anyone can still read the body. Only encryption makes the
body readable solely by the intended recipient. The wrong option has
assigned encryption's job to the signature.
\end{example}

\section{Spam, phishing, spoofing, and malicious content}
\textbf{Spam} is bulk unwanted mail, usually advertising or junk sent to
many addresses at once. Its defining mark is that it is bulk and unwanted;
it clutters the inbox and may also carry fraud or malware. Spam filters move
suspected junk to a separate folder using the sender's reputation and
content rules.

\textbf{Phishing} is a fake message that pretends to be a trusted sender
--- a bank, an office, or a colleague --- in order to steal passwords or
money. It typically creates urgency (``verify now'', ``account blocked'',
``claim within 24 hours'') and asks the reader to click a link or open an
attachment and hand over secrets. The defence is to check the real sender
address and never enter passwords through a link reached from mail.

\textbf{Spoofing} forges the sender address so that a message looks as if
it came from someone trusted. Phishing messages often use spoofing, but
the two terms are distinct: phishing is the fraud attempt, and spoofing is
the forged-sender trick it may use. A \textbf{malicious link} leads to a
fake login page or a malware download; hovering over it shows an address
that does not match the claimed sender. A \textbf{malicious attachment}
carries malware behind an innocent-looking file name. The safe rule is to
verify an unexpected message through another channel before clicking any
link or opening any attachment.

\section{Exam retrieval and common confusions}
Six distinctions prevent most errors on this topic.
\begin{enumerate}
  \item SMTP \textbf{sends}; POP3 and IMAP \textbf{receive}. Reversing the
    direction is a common wrong answer.
  \item POP3 \textbf{downloads} mail to one local machine (usually clearing
    the server); IMAP \textbf{syncs} every device with the server copy.
    Do not swap them (TRAP-17).
  \item A digital signature proves the sender and detects tampering; it
    does \textbf{not} encrypt the body. Encryption alone provides secrecy
    (TRAP-15).
  \item BCC hides recipients \textbf{from each other}; the sender always
    knows the full list, and To/Cc are visible to all (PYQ-16).
  \item Secure transmission of email is \textbf{SSL/TLS} (PYQ-36). A
    signature is not a transmission protocol.
  \item Spam is \textbf{bulk unwanted} mail; phishing is \textbf{fake
    trusted mail that steals secrets}; spoofing is the \textbf{forged
    sender address} either may use.
\end{enumerate}

\begin{example}
An item asks which protocol suits reading the same office mailbox from a
phone and a desktop. Trace the methods: POP3 downloads to one machine and
clears the server, so the second device misses the mail; IMAP keeps the
server copy and syncs every device with it. The answer is IMAP. If the stem
instead asks which protocol sends the message out, the answer is SMTP, and
if it asks which channel makes transmission secure, the answer is SSL/TLS.
\end{example}

\subsection*{Exercises}
\begin{enumerate}[label=\arabic*.]
  \item State the full form of SMTP, POP3, IMAP, SSL, and TLS, and give the
    job of each of the first three.
  \item Explain the difference between POP3 and IMAP when the same mailbox
    is read from two devices.
  \item Trace one message from the sender to the reader, naming the
    protocol used at each step.
  \item Distinguish a digital signature from encryption, and state which
    protocol secures email transmission.
  \item Define spam, phishing, and spoofing, and state the safe rule for
    unexpected links and attachments.
\end{enumerate}

\subsection*{Solutions}
\begingroup\small
\begin{enumerate}[label=\arabic*.]
  \item SMTP is Simple Mail Transfer Protocol and sends mail; POP3 is Post
    Office Protocol version 3 and downloads received mail to one machine;
    IMAP is Internet Message Access Protocol and syncs received mail across
    devices; SSL is Secure Sockets Layer and TLS is Transport Layer
    Security, together the secure transmission channel.
  \item With POP3 the mail downloads to the first machine and usually
    leaves the server, so the second device does not see it; with IMAP the
    mail stays on the server and both devices show the same synced view of
    folders, read marks, and deletions.
  \item The sender's program sends the message by SMTP to the sender's
    server; the servers forward it by SMTP to the recipient's server, where
    it waits; the recipient collects it with POP3 (download) or IMAP
    (sync).
  \item A digital signature provides authenticity and integrity (who sent
    it, unaltered); encryption provides confidentiality (only the recipient
    can read it). Secure transmission uses SSL/TLS.
  \item Spam is bulk unwanted mail; phishing is a fake trusted message
    that steals secrets through urgency plus a link or attachment;
    spoofing forges the sender address. Verify an unexpected message by
    another channel before clicking any link or opening any attachment.
\end{enumerate}
\endgroup

\lessonheading{39}{E-Banking}
\section{What e-banking covers}
Electronic banking, or e-banking, means using banking services without
visiting the branch. The account itself stays in the bank; only the channel
used to reach it changes. \textbf{Net banking} is banking through the bank's
website in a browser. \textbf{Mobile banking} is banking through the bank's
app on the phone. \textbf{UPI} (Unified Payments Interface) apps move money
directly between bank accounts from the phone. An ATM, the net-banking site,
the mobile app, and a UPI app are therefore four doors into the same
account.

\begin{definitionbox}{E-banking}
The use of banking services through electronic channels --- net banking,
mobile banking, UPI, and ATMs --- instead of visiting the branch in person.
\end{definitionbox}

\begin{center}
\begin{tabular}{p{0.24\textwidth}p{0.66\textwidth}}
\toprule
\textbf{Term} & \textbf{Meaning in this lesson} \\
\midrule
Net banking & Banking through the bank's website in a browser. \\
Mobile banking & Banking through the bank's official app. \\
UPI & Unified Payments Interface; instant phone-based payments between accounts. \\
VPA & Virtual Payment Address (for example name@bank); identifies the account without the account number. \\
ATM PIN & Personal Identification Number used with the physical card. \\
UPI PIN & Personal Identification Number authorising UPI payments. \\
CVV & Card Verification Value printed on the card. \\
OTP & One-Time Password sent by SMS, valid for one transaction only. \\
MFA & Multi-Factor Authentication; checking two or more factors together. \\
\bottomrule
\end{tabular}
\end{center}

The rest of this lesson uses these terms in one consistent sense. A
\textbf{PIN} is a secret typed by the holder. \textbf{CVV} is printed on the
card. \textbf{OTP} arrives by SMS and expires. \textbf{MFA} is the
combination of two or more different proofs of identity.

\section{UPI in detail}
UPI stands for Unified Payments Interface. Each user is identified by a
\textbf{VPA} (Virtual Payment Address), such as name@bank, so money can be
sent without sharing the account number or IFSC. Transfers are
\textbf{instant} and work \textbf{24 $\times$ 7}, including Sundays and
holidays. Every payment is authorised on the payer's own phone by entering
the \textbf{UPI PIN}, which is set separately for each bank account inside
the UPI app. The central rule of UPI safety is short: entering the UPI PIN
always sends money, and receiving money never needs any PIN.

Other electronic transfer modes appear in nearby options and should be kept
in one line each. \textbf{IMPS} (Immediate Payment Service) is an instant
bank-to-bank transfer available 24 $\times$ 7. \textbf{NEFT} (National
Electronic Funds Transfer) settles in batches and has no minimum amount.
\textbf{RTGS} (Real Time Gross Settlement) settles large-value transfers
instantly and carries a minimum amount. UPI and IMPS are the instant,
phone-friendly modes; NEFT is batched; RTGS is for large values.

\begin{definitionbox}{UPI}
Unified Payments Interface: an instant, round-the-clock system for moving
money between bank accounts from a phone, identified by VPA and authorised by
a UPI PIN.
\end{definitionbox}

\section{The four secret codes}
Four codes look alike and are deliberately mixed in distractors. Each lives
in a different place and does a different job.

\textbf{ATM PIN.} The Personal Identification Number used with the physical
card at an ATM or shop terminal. It is memorised, typed on the machine, and
authorises card transactions on that channel. It never authorises a UPI
payment.

\textbf{UPI PIN.} The Personal Identification Number typed in the phone to
authorise a UPI payment. It is set per bank account in the UPI app and can
differ from the ATM PIN. Approving any screen with the UPI PIN sends money.

\textbf{CVV.} The Card Verification Value: the 3-digit number printed on the
back of the card (4 digits on the front for some cards). It is entered with
the card number and expiry date for online card payments and proves the
card is in hand. It is fixed and reusable for the life of the card.

\textbf{OTP.} The One-Time Password: a short numeric code sent by SMS that
is valid for one transaction only and for a few minutes. It is generated
fresh each time and expires after use or after the time window.

\begin{center}
\begin{tabular}{p{0.16\textwidth}p{0.24\textwidth}p{0.22\textwidth}p{0.28\textwidth}}
\toprule
\textbf{Code} & \textbf{Where it lives} & \textbf{How long it lasts} & \textbf{Typical use} \\
\midrule
ATM PIN & In memory & Until changed & Card at ATM or shop terminal \\
UPI PIN & In memory & Until changed & Authorising a UPI payment \\
CVV & Printed on card & Life of the card & Online card payment with card number \\
OTP & Sent by SMS & One transaction, minutes & Second approval for a login or payment \\
\bottomrule
\end{tabular}
\end{center}

\begin{example}
A question states, ``The 3-digit number on the back of the card was sent by
SMS and expired after one use.'' This mixes two codes. The number printed on
the back of the card is the \textbf{CVV}, which is fixed and reusable; the
code that arrives by SMS and expires after one use is the \textbf{OTP}. The
stem has attached OTP behaviour to the CVV name.
\end{example}

\section{Authentication and multi-factor authentication}
\textbf{Authentication} means proving ``I am the account holder'' before the
bank acts. It uses three kinds of factor. \textbf{Something you know} is a
password or PIN. \textbf{Something you have} is the phone, the card, or the
OTP that arrives on the registered mobile number. \textbf{Something you are}
is a fingerprint or face scan.

\begin{definitionbox}{Multi-Factor Authentication (MFA)}
Authentication that checks two or more different factors together, for
example a password (know) plus an SMS OTP (have), or a PIN plus a
fingerprint.
\end{definitionbox}

A login password plus an SMS OTP is therefore \textbf{two-factor}
authentication, not one: the password proves knowledge and the OTP proves
possession of the registered phone. Asking for the OTP a second time does
not add a new factor; a fingerprint or card would. The exam trap is to call a
password-plus-OTP login ``single-factor'' or to call two passwords
``multi-factor''; factors must be of \emph{different} kinds.

\section{QR codes and collect requests}
A \textbf{QR} (Quick Response) code is the square barcode the merchant
displays. Scanning it with the UPI app fills in the merchant's payment
address; the payer then checks the \textbf{merchant name and amount} on
their own screen and approves with the UPI PIN. A QR never asks for a PIN or
OTP by itself. A common fraud is a fake QR sticker pasted over the real one,
so the displayed name on the payer's screen is the check that matters: if
the name does not match the shop, stop.

A \textbf{collect request} is a message asking the receiver to pay: it
appears as ``pay Rs.\ X to Y, approve with UPI PIN''. Approving it and
entering the UPI PIN sends the receiver's money to the requester. No refund,
cashback, prize, or buyer ever needs the seller to enter a UPI PIN; entering
the PIN always means paying. Unexpected collect requests, especially ones
with remarks like ``refund'' or ``cashback'', must be declined or ignored.
To receive money, the receiver only shares a VPA or QR and enters nothing.

\begin{example}
A buyer sends a Rs.\ 500 collect request to a seller with the remark
``advance, approve to receive''. The seller's screen shows ``pay Rs.\ 500'',
and approving with the UPI PIN would send Rs.\ 500 to the buyer. The remark
is false: receiving needs no PIN. The safe action is to decline and ask the
buyer to send the payment instead.
\end{example}

\section{Safe-banking rules}
The rules below are tested almost verbatim, so they are stated exactly.

\begin{enumerate}
  \item The bank \textbf{never} asks for OTP, CVV, ATM PIN, or UPI PIN by
    call, SMS, email, or link. Any such request is fraud.
  \item \textbf{Never share} an OTP for any reason --- to ``cancel'', to
    ``verify'', or to ``receive'' a payment. Sharing it completes the
    transaction the fraudster started.
  \item Do not open \textbf{KYC-update or account-blocked} links from SMS.
    Open the bank's site or app by typing the address or using the official
    app.
  \item \textbf{Never install} a screen-sharing or remote-access app at a
    caller's request; it exposes OTPs, PINs, and balances on the screen.
  \item Use \textbf{typed addresses and official apps}, and check the
    \textbf{padlock (https)} on the banking site before logging in.
  \item Avoid \textbf{public Wi-Fi} for banking; prefer mobile data or a
    trusted network.
  \item \textbf{Log out} after net banking, lock the UPI app with the device
    lock, and keep \textbf{SMS and email alerts on} so an unknown debit is
    noticed at once.
\end{enumerate}

\section{Exam retrieval and common confusions}
Six distinctions prevent most errors on this topic.
\begin{enumerate}
  \item ATM PIN is for the \textbf{card} channel; UPI PIN is for the
    \textbf{UPI} channel. They are set and used separately.
  \item CVV is \textbf{printed on the card} and reusable; OTP is
    \textbf{sent by SMS} for one transaction and expires.
  \item Entering the UPI PIN \textbf{always pays}; receiving money needs
    \textbf{no PIN}.
  \item A password plus an SMS OTP is \textbf{two-factor} authentication
    (know + have); two passwords are still one factor.
  \item The bank \textbf{never asks} for OTP, CVV, or any PIN. Any caller
    who asks is not the bank.
  \item A QR is verified by the \textbf{name and amount on your own
    screen}; a pasted-over sticker is defeated by that check.
\end{enumerate}

\begin{example}
An item asks what is needed to receive a UPI refund. Trace the rule:
receiving needs a VPA or QR and no PIN; any screen asking for the UPI PIN is
a payment screen. The answer is that no PIN or OTP is entered to receive;
the ``approve to receive'' request is a collect-request fraud.
\end{example}

\subsection*{Exercises}
\begin{enumerate}[label=\arabic*.]
  \item Expand UPI, OTP, CVV, PIN, QR, and MFA.
  \item Distinguish the ATM PIN from the UPI PIN: where is each used?
  \item Distinguish CVV from OTP: where does each live, and how long does
    each last?
  \item Explain why a login password plus an SMS OTP counts as two-factor
    authentication.
  \item A seller receives a collect request marked ``refund, approve to
    receive''. What should the seller do, and why?
\end{enumerate}

\subsection*{Solutions}
\begingroup\small
\begin{enumerate}[label=\arabic*.]
  \item UPI --- Unified Payments Interface; OTP --- One-Time Password; CVV
    --- Card Verification Value; PIN --- Personal Identification Number; QR
    --- Quick Response code; MFA --- Multi-Factor Authentication.
  \item The ATM PIN is typed on a machine with the physical card at an ATM
    or shop terminal. The UPI PIN is typed in the phone to authorise a UPI
    payment, set per bank account in the UPI app. Neither works on the
    other's channel.
  \item CVV is printed on the card and is reusable for the life of the card;
    it is entered with the card number for online card payments. OTP is sent
    by SMS, is valid for one transaction only and a few minutes, and then
    expires.
  \item The password is something the user knows and the OTP on the
    registered phone is something the user has. Two different factor kinds
    are checked, so the login is two-factor (a form of MFA).
  \item Decline or ignore it. A collect request approved with the UPI PIN
    sends the seller's money; receiving a refund needs no PIN. The remark is
    a fraud pattern.
\end{enumerate}
\endgroup

\lessonheading{53}{Networking Legacy Services and Web Standards}
\section{How to read this lesson}
This is a trap-drill lesson. Almost every concept here has a near-neighbour:
another term that sounds similar or sits next to it in a list but does a
different job. The exam builds its wrong options out of exactly those
neighbours, and it often asks the question with a negative stem (NOT, EXCEPT,
or an incorrect statement). The defence is always the same: state what the
term is, what it is not, and which neighbour it is confused with.

\begin{center}
\begin{tabular}{p{0.24\textwidth}p{0.66\textwidth}}
\toprule
\textbf{Term} & \textbf{Meaning in this lesson} \\
\midrule
Circuit switching & A communication method that reserves a dedicated path for the whole session before data moves. \\
Packet switching & A communication method that breaks data into packets, each travelling independently with no reserved path. \\
Telnet & TELecommunication NETwork; a remote-login service for running commands on a distant computer. \\
FTP & File Transfer Protocol; the service for uploading and downloading files. \\
SMTP & Simple Mail Transfer Protocol; sends mail from the sender toward the receiver's server. \\
POP3 & Post Office Protocol version 3; downloads mail to the local machine. \\
IMAP & Internet Message Access Protocol; keeps mail on the server and syncs it across devices. \\
W3C & World Wide Web Consortium; the body that develops open web standards. \\
DHTML & Dynamic HyperText Markup Language; HTML combined with styling and scripting for dynamic pages. \\
URL & Uniform Resource Locator; the full address of a web resource. \\
DNS & Domain Name System; translates domain names into IP addresses. \\
IP address & The numeric address of a device on the network. \\
OSI model & The seven-layer networking reference model. \\
Physical port & A hardware socket on the machine. \\
Logical port & A number identifying a service or process in networking software. \\
\bottomrule
\end{tabular}
\end{center}

\section{Packet switching vs circuit switching}
\textbf{Circuit switching} establishes a dedicated communication path between
the two ends before any data flows, and that path stays reserved for the whole
session. The classic example is the old telephone call: the line is yours
until you hang up, whether or not anyone is speaking.

\begin{definitionbox}{Circuit switching}
A switching method in which a dedicated path is set up and reserved for the
entire communication session before data transfer begins.
\end{definitionbox}

\textbf{Packet switching} works the opposite way. Data is broken into small
units called packets, and each packet travels independently through the
network. No path is reserved in advance; packets from many senders share the
same links. The Internet runs on packet switching.

\begin{definitionbox}{Packet switching}
A switching method in which data is divided into packets that travel
independently over shared links, with no dedicated path reserved in advance.
\end{definitionbox}

\begin{example}
An item states, ``The Internet reserves a dedicated circuit for each browsing
session.'' This is wrong. The Internet uses packet switching: pages arrive as
independent packets over shared links, and no path is reserved for the session.
The wrong option has described circuit switching instead.
\end{example}

\section{Telnet and FTP}
\textbf{Telnet} (TELecommunication NETwork) is a remote-login service. It lets
a user connect to another computer over the network and run commands there as
if sitting at its keyboard. Telnet sends its traffic, including passwords, in
plain text. It is not a file-transfer service and not a mail service.

\textbf{FTP} (File Transfer Protocol) is the file-transfer service. It moves
files between computers: uploading from the local machine to a server and
downloading back. It is not a remote-login shell and not a mail protocol.

The exam tests the triple together: Telnet logs in remotely, FTP moves files,
and SMTP sends mail. Any option that swaps these roles --- ``Telnet transfers
files'' or ``FTP sends mail'' --- is a planted near-neighbour distractor.

\begin{definitionbox}{Telnet}
A network service that provides remote login to another computer so the user
can run commands on it.
\end{definitionbox}

\begin{definitionbox}{FTP}
A network protocol used to transfer files between computers on a network.
\end{definitionbox}

\section{SMTP, POP3, and IMAP}
Mail handling is split across three protocols, and the exam tests the split.
\textbf{SMTP} (Simple Mail Transfer Protocol) \emph{sends} mail: it carries a
message from the sender's machine toward the receiver's mail server.
\textbf{POP3} (Post Office Protocol version 3) \emph{receives by download}:
it pulls mail down to the local machine, typically removing it from the
server. \textbf{IMAP} (Internet Message Access Protocol) \emph{receives by
sync}: it keeps mail on the server so the same mailbox can be read from
several devices.

\begin{definitionbox}{SMTP / POP3 / IMAP roles}
SMTP sends mail; POP3 downloads mail to the local machine; IMAP keeps mail on
the server and synchronises it across devices.
\end{definitionbox}

Three attached traps come from real papers. First, an option saying
``POP3 keeps mail synced on the server'' is wrong; that synced behaviour is
IMAP (TRAP-17). Second, ``a digital signature encrypts the body so only the
recipient can read it'' is wrong: a digital signature proves authenticity and
integrity, while encryption provides confidentiality (TRAP-15). Third,
``BCC hides the sender and the recipients from each other'' is wrong: BCC
hides recipients from each other, the sender always knows everyone, and To
and Cc names remain visible to all recipients (TRAP-16). Finally, in any mail
address the \textbf{@} sign separates the user name from the domain name
(TRAP-28).

\begin{example}
An item asks which protocol downloads mail to the local machine. The answer
is POP3. If the stem instead asks which keeps the mailbox synced on the
server across a phone and a desktop, the answer is IMAP, and if it asks
which sends the message in the first place, the answer is SMTP.
\end{example}

\section{W3C and DHTML}
The \textbf{W3C} (World Wide Web Consortium) develops the open standards of
the World Wide Web, such as HTML and CSS. It is a standards body --- not a
browser, not a search engine, and not an Internet service provider. The
near-neighbour is the \textbf{ISP} (Internet Service Provider), which sells
Internet access; W3C writes the standards that the access carries.

\begin{definitionbox}{W3C}
The World Wide Web Consortium: the international body that develops open
standards for the World Wide Web.
\end{definitionbox}

\textbf{HTML} (HyperText Markup Language) marks up the static structure of a
web page. \textbf{DHTML} (Dynamic HyperText Markup Language) is not a
separate replacement language; it is the combination of HTML with styling and
scripting that makes pages respond and change. The trap option calls DHTML a
brand-new language that replaces HTML. The correct statement is that DHTML
builds on HTML to add dynamic behaviour.

\section{URL, DNS, and IP addresses}
A \textbf{URL} (Uniform Resource Locator) is the full address of a resource
on the web. A \textbf{domain name} is the human-readable part inside it. The
\textbf{DNS} (Domain Name System) translates that human-readable name into a
numeric address. The \textbf{IP address} is that numeric address: it
identifies a device on the network.

The chain in one line is: type a URL, DNS resolves its domain name, and the
browser reaches the IP address. The exam swaps the links: an option saying
``DNS is the numeric address'' confuses DNS with the IP address, and one
saying ``the URL is the numeric address'' confuses the URL with the IP
address.

\section{IPv4 vs IPv6}
\textbf{IPv4} uses 32-bit addresses, written as four decimal groups. The
address space proved too small, so \textbf{IPv6} was introduced with 128-bit
addresses and a fixed header of 40 bytes (PYQ-18).

\begin{definitionbox}{IPv4 / IPv6 address size}
IPv4 addresses are 32 bits long; IPv6 addresses are 128 bits long.
\end{definitionbox}

The single most repeated distractor on this topic claims IPv6 is 256-bit
(TRAP-18, from JA 2026 Q78). It is always wrong: IPv6 is 128-bit, never
256-bit. A second attached trap from the same paper concerns companion
protocols: \textbf{ICMP} is a control and diagnostic protocol that carries no
application data and uses no ports (TRAP-19); \textbf{BGP} routes between
autonomous systems (inter-domain) and does not route inside one local network
(TRAP-20); and \textbf{HTTP/3} runs over QUIC over UDP, not over TCP
(TRAP-21, from JA 2026 Q79).

\section{The seven OSI layers}
The \textbf{OSI model} (Open Systems Interconnection) is the seven-layer
reference model for networking. From top to bottom the layers are
Application, Presentation, Session, Transport, Network, Data Link, and
Physical. A memory hook is ``All People Seem To Need Data Processing''.

\begin{definitionbox}{OSI model}
A seven-layer reference model for network communication: Application,
Presentation, Session, Transport, Network, Data Link, Physical.
\end{definitionbox}

The near-neighbour trap names four layers. Four is the layer count of the
TCP/IP model, not the OSI model. When the stem says OSI, the number is
always seven; an option offering four has imported the wrong model.

\section{Physical vs logical ports}
A \textbf{physical port} is a hardware socket on the machine: the place where
a cable or device plugs in. A \textbf{logical port} is a number in networking
software that identifies a service or process. The trap swaps them: a logical
port number is not a hardware socket, and a physical socket is not a service
number.

\begin{definitionbox}{Physical vs logical port}
A physical port is a hardware socket for cables or devices; a logical port is
a number identifying a network service or process.
\end{definitionbox}

\section{Exam retrieval and common confusions}
Nine distinctions prevent most errors on this topic.
\begin{enumerate}
  \item Circuit switching reserves a dedicated path; packet switching sends
    independent packets with no reserved path.
  \item Telnet gives remote login; FTP transfers files; SMTP sends mail. Do
    not swap the three roles.
  \item POP3 downloads mail to the local machine; IMAP keeps mail synced on
    the server (TRAP-17).
  \item A digital signature proves authenticity and integrity; encryption
    provides confidentiality (TRAP-15).
  \item BCC hides recipients from each other; the sender always knows all
    recipients (TRAP-16).
  \item IPv6 is 128-bit with a 40-byte header, never 256-bit (TRAP-18).
  \item ICMP carries no application data; BGP is inter-domain; HTTP/3 uses
    QUIC over UDP (TRAP-19, TRAP-20, TRAP-21).
  \item The OSI model has seven layers; four layers is the TCP/IP model.
  \item A physical port is a hardware socket; a logical port is a service
    number.
\end{enumerate}

\begin{example}
An item asks which statement is incorrect. Work each option against its
near-neighbour: does the switching option reserve or share the path, does
the protocol option match the service to its role, does the address option
use 128-bit for IPv6, does the layer option say seven for OSI? The option
that breaks its pair is the incorrect statement the stem asks for.
\end{example}

\subsection*{Exercises}
\begin{enumerate}[label=\arabic*.]
  \item Distinguish circuit switching from packet switching in one sentence
    each.
  \item Match Telnet, FTP, and SMTP to their roles.
  \item State the roles of SMTP, POP3, and IMAP, and say which one keeps mail
    synced on the server.
  \item Expand W3C, DHTML, URL, DNS, FTP, and SMTP.
  \item State the bit lengths of IPv4 and IPv6, and name the OSI layers from
    top to bottom.
\end{enumerate}

\subsection*{Solutions}
\begingroup\small
\begin{enumerate}[label=\arabic*.]
  \item Circuit switching reserves a dedicated path for the whole session
    before data moves; packet switching breaks data into independently
    travelling packets with no reserved path.
  \item Telnet gives remote login to another computer; FTP transfers files
    between computers; SMTP sends mail toward the receiver's server.
  \item SMTP sends mail, POP3 downloads it to the local machine, and IMAP
    keeps it on the server and syncs it; the synced-on-server protocol is
    IMAP.
  \item W3C --- World Wide Web Consortium; DHTML --- Dynamic HyperText Markup
    Language; URL --- Uniform Resource Locator; DNS --- Domain Name System;
    FTP --- File Transfer Protocol; SMTP --- Simple Mail Transfer Protocol.
  \item IPv4 is 32-bit and IPv6 is 128-bit; the OSI layers from the top are
    Application, Presentation, Session, Transport, Network, Data Link,
    Physical.
\end{enumerate}
\endgroup

\lessonheading{63}{Current Web, HTML, JavaScript, VBScript and ASP.NET}
\section{Static pages, markup, scripts, and the server}
A web page can be static or dynamic. A \textbf{static} page shows the same content to every visitor. A \textbf{dynamic} page builds its content, or changes it, at view time. Three layers combine to produce the modern page. \textbf{Markup} describes structure: HTML lays out headings, paragraphs, and links. \textbf{Scripting} adds behaviour: small programs in JavaScript or VBScript run inside the page and react to the visitor. \textbf{Server technology} such as ASP.NET builds the page on the server before the browser ever sees it.

\begin{definitionbox}{Markup language}
A language that tags text to say what each part is (heading, paragraph, link), rather than giving computational instructions.
\end{definitionbox}

\begin{center}
\begin{tabular}{p{0.24\textwidth}p{0.66\textwidth}}
\toprule
\textbf{Term} & \textbf{Meaning in this lesson} \\
\midrule
HTML & HyperText Markup Language; structures the page; stored as ASCII text. \\
Tag & A named element in angle brackets, e.g.\ \texttt{<p>}. \\
Attribute & Detail inside the opening tag as \texttt{name="value"}. \\
DHTML & Dynamic HTML: HTML plus styles plus scripts for interactive pages. \\
XML & eXtensible Markup Language: carries data with user-defined tags. \\
JavaScript & Netscape browser scripting language; \texttt{+} joins strings. \\
VBScript & Microsoft scripting language; variables declared with \texttt{Dim}. \\
ASP.NET & Microsoft server framework; \texttt{Session} stores one user's data. \\
Meta-search engine & Queries other search engines; Dogpile is the classic example. \\
\bottomrule
\end{tabular}
\end{center}

The rest of this lesson uses these terms in one consistent sense. A \textbf{tag} names an element; an \textbf{attribute} refines it. \textbf{DHTML} is not a new language but a combination. \textbf{XML} carries data while HTML displays it. \textbf{Java} and \textbf{JavaScript} are different languages despite the similar names.

\section{HTML: a markup language in an ASCII-text document}
\textbf{HTML} (HyperText Markup Language) is a markup language, not a programming language. It does not compute; it labels. Each piece of text is wrapped in tags that say what it is: a heading, a paragraph, a link. The browser reads these tags and renders the page, and the tags themselves are not displayed.

An HTML document is an \textbf{ASCII-text document}. It is plain text and can be written and read in any text editor. No special software is needed to create one, and nothing about it is compiled to binary. This is a favourite direct-recall fact: HTML files are plain ASCII text.

\begin{definitionbox}{HTML}
HyperText Markup Language: the markup language that structures web pages, stored as plain ASCII text and rendered by the browser.
\end{definitionbox}

\begin{example}
A student opens a page source and sees \texttt{<html>}, \texttt{<body>}, and \texttt{<p>Hello</p>}. These are tags, readable as text in the source, but the visitor only sees the rendered words. The document on disk is ordinary text, which is why the lesson calls HTML an ASCII-text document.
\end{example}

\section{Tags, attributes, and validity}
A \textbf{tag} names an element and is written in angle brackets. Most elements have three parts: an \textbf{opening tag} such as \texttt{<p>}, the content, and a \textbf{closing tag} with a slash such as \texttt{</p>}. An \textbf{attribute} adds detail and always sits inside the opening tag in the form \texttt{name="value"}. In \texttt{<body bgcolor="yellow">}, \texttt{body} is the tag and \texttt{bgcolor="yellow"} is the attribute. In \texttt{<font size="4">}, \texttt{font} is the tag and \texttt{size} is the attribute.

A tag or attribute is \textbf{valid} when the language defines it. An invented name is simply not recognised by the browser. Exam items test this by mixing real names (\texttt{h1}, \texttt{p}, \texttt{marquee}, \texttt{font}, \texttt{body}, \texttt{html} as tags; \texttt{size}, \texttt{bgcolor} as attributes) with invented ones. The method is always the same: match the name against the defined set, and never confuse a tag with an attribute.

\section{The \texttt{h1} heading and the \texttt{marquee} tag}
Headings run from \texttt{h1} through \texttt{h6}. \textbf{\texttt{h1}} is the largest, top-level heading, and the numbers grow as the headings shrink, so \texttt{h6} is the smallest. It is written as \texttt{<h1>Title</h1>}. The single most common error is reversing this order, so fix the anchor now: \textbf{\texttt{h1} is biggest}.

The \textbf{\texttt{marquee}} tag makes text scroll across the page, written as \texttt{<marquee>Welcome</marquee>}. It is a presentation tag: it moves text but adds no structure. It belongs to older, browser-specific HTML rather than strict modern structure, which is exactly why exam setters like it as a recall item.

\begin{definitionbox}{Heading levels}
Six levels \texttt{h1} to \texttt{h6}; \texttt{h1} renders largest and \texttt{h6} smallest.
\end{definitionbox}

\section{DHTML: the dynamic combination}
\textbf{DHTML} (Dynamic HyperText Markup Language) is not a separate language. It is the combination of HTML with style sheets and scripts working together, so that a page can change after it has loaded --- for example, text that reacts when the mouse passes over it. Plain HTML is static; a DHTML page is interactive without needing a reload.

\begin{definitionbox}{DHTML}
Dynamic HyperText Markup Language: HTML combined with style sheets and scripts so a page can change after loading.
\end{definitionbox}

The exam contrast is always the same. If the stem says the page is fixed once sent, the answer is HTML. If it says the page reacts in the browser without reloading, the answer is DHTML.

\section{XML: data, not display}
\textbf{XML} (eXtensible Markup Language) carries and stores data; it does not lay out a page the way HTML does. Its tags are \textbf{user-defined}: the author invents tags such as \texttt{<city>} to label the data being carried. HTML uses a fixed tag set for display; XML uses an open tag set for data. That is the whole classification: HTML displays, XML transports.

\begin{definitionbox}{XML}
eXtensible Markup Language: a markup language for storing and transporting data with author-defined tags.
\end{definitionbox}

\begin{example}
An item asks which language is used to transport structured data between systems with custom tags. HTML is wrong because its tags are fixed and presentational. The answer is XML, whose tags the author defines for the data at hand.
\end{example}

\section{Java and JavaScript are different languages}
\textbf{Java} is a full programming language: code is compiled to bytecode and runs on the Java platform. \textbf{JavaScript} is a lightweight scripting language: code is interpreted in the browser to make pages interactive. They share part of a name and nothing else. Every exam cycle reuses this trap, so treat any option that calls JavaScript a version or subset of Java as wrong on sight.

\textbf{JavaScript} was developed at \textbf{Netscape} for the Navigator browser in the mid-1990s. It first shipped under other names (Mocha, then LiveScript) before the JavaScript name was settled. For exam purposes the anchor is one line: JavaScript means Netscape, just as VBScript means Microsoft.

\section{JavaScript string concatenation with \texttt{+}}
In JavaScript the \textbf{\texttt{+}} operator joins (concatenates) strings. Joining \texttt{"Good"} and \texttt{"Morning"} gives \texttt{"GoodMorning"}. Mixing text and a number with \texttt{+} also joins as text. For example, joining \texttt{"Total: "} with \texttt{5} gives \texttt{"Total: 5"}. The same symbol adds numbers, but wherever a string is involved it concatenates. Items test this by asking what an expression evaluates to; the method is mechanical: join the pieces as text.

\begin{example}
\texttt{"Web" + "Page"} evaluates to \texttt{"WebPage"} --- joined, not added. An option showing a numeric sum for two quoted strings has confused the two roles of \texttt{+}.
\end{example}

\section{VBScript and the \texttt{Dim} statement}
\textbf{VBScript} (Visual Basic Scripting Edition) is Microsoft's scripting language in the Basic family, used for page scripting and small Windows tasks. Variables are declared with the \textbf{\texttt{Dim}} statement, short for Dimension: \texttt{Dim city} declares a variable named \texttt{city}. The exam anchor is direct: see \texttt{Dim} and think VBScript variable declaration. No other language in this lesson declares variables that way.

\section{ASP.NET and the \texttt{Session} object}
\textbf{ASP.NET} is Microsoft's framework for building dynamic web pages on the server. The server runs the code, builds the page, and sends plain HTML to the browser. The \textbf{\texttt{Session}} object stores data for \textbf{one user across many pages}, such as a login name kept while the visitor moves from page to page. Its nearest neighbour is the \textbf{Application} object, which is shared by all users. The distinction is per-user versus shared: \texttt{Session} keeps one visitor's data; it is never shared across visitors.

\begin{definitionbox}{Session object}
The ASP.NET object that holds data for a single user across multiple page requests.
\end{definitionbox}

\section{Meta-search engines: Dogpile}
A \textbf{search engine} such as Google or Bing keeps its own index of pages. A \textbf{meta-search engine} keeps no index of its own; it sends the query to other search engines and combines their results. \textbf{Dogpile} is the classic exam example of a meta-search engine. The memory hook is in the name: Dogpile piles up other engines' results and shows them together.

\section{Exam retrieval and common confusions}
Six distinctions prevent most errors on this topic.
\begin{enumerate}
  \item HTML is an \textbf{ASCII-text} document. An option calling it a compiled binary or a database file is wrong.
  \item \textbf{\texttt{h1}} is the largest heading; \textbf{\texttt{marquee}} scrolls text. Do not reverse the heading order.
  \item A \textbf{tag} names the element; an \textbf{attribute} sits inside the opening tag as \texttt{name="value"}. \texttt{font} is a tag; \texttt{size} is its attribute.
  \item \textbf{Java} is a full compiled language; \textbf{JavaScript} is a Netscape browser scripting language. Similar names, different languages.
  \item In JavaScript \textbf{\texttt{+}} concatenates strings; in VBScript \textbf{\texttt{Dim}} declares a variable.
  \item The ASP.NET \textbf{\texttt{Session}} object is per-user, not shared; \textbf{Dogpile} is a meta-search engine, not a search engine with its own index.
\end{enumerate}

\begin{example}
An item asks what \texttt{<h1>} produces and offers ``smallest heading'' as an option. Trace the heading ladder: \texttt{h1} is the top level and renders largest, with \texttt{h6} smallest. The option has reversed the ladder. If the stem instead asks which tag scrolls text, scan for \texttt{marquee}; if it asks which object keeps one user's login across pages, scan for \texttt{Session}.
\end{example}

\subsection*{Exercises}
\begin{enumerate}[label=\arabic*.]
  \item Expand HTML, DHTML, XML, and ASP.NET, and state in one line what each does.
  \item Explain why an HTML document is called an ASCII-text document.
  \item Distinguish a tag from an attribute, giving one example of each.
  \item Which heading tag is largest, and which tag scrolls text across the page?
  \item Distinguish Java from JavaScript, and state where JavaScript originated.
  \item What does the \texttt{+} operator do to two strings in JavaScript? What does \texttt{Dim} do in VBScript?
  \item What does the ASP.NET \texttt{Session} object store, and how does it differ from Application state?
  \item What is a meta-search engine, and why is Dogpile an example?
\end{enumerate}

\subsection*{Solutions}
\begingroup\small
\begin{enumerate}[label=\arabic*.]
  \item HTML --- HyperText Markup Language: structures web pages. DHTML --- Dynamic HyperText Markup Language: HTML plus styles plus scripts for interactive pages. XML --- eXtensible Markup Language: carries data with user-defined tags. ASP.NET: Microsoft's server framework for building dynamic pages.
  \item It is plain text: it can be written and read in any text editor, with no compilation step. The browser renders the tags at view time.
  \item A tag names an element in angle brackets (e.g.\ \texttt{<p>}); an attribute sits inside the opening tag as \texttt{name="value"} (e.g.\ \texttt{size="4"} in \texttt{<font size="4">}).
  \item \texttt{h1} is the largest heading (\texttt{h6} smallest); \texttt{marquee} scrolls text.
  \item Java is a full compiled programming language running on the Java platform; JavaScript is a lightweight browser scripting language developed at Netscape.
  \item \texttt{+} joins (concatenates) the strings as text. \texttt{Dim} declares a variable in VBScript.
  \item \texttt{Session} stores data for one user across many pages (e.g.\ a login name). Application state is shared by all users; \texttt{Session} is per-user.
  \item A meta-search engine keeps no index and combines results from other search engines. Dogpile does exactly this, so it is the classic example.
\end{enumerate}
\endgroup

\unitopener{9}{Networking}
\lessonheading{40}{Networking Basics}
\section{What a network is: nodes, hosts, clients, servers}
A computer network is formed when two or more computers and related devices
are connected so that they can share data and resources. The shared
resources include files, printers, an internet connection, and application
software. On the office desktop, opening a shared folder, printing to a
shared printer, and browsing the web all use the same underlying network.

\begin{definitionbox}{Computer network}
A set of two or more connected computing devices (nodes) that can exchange
data and share resources such as files, printers, and an internet
connection.
\end{definitionbox}

Every participant in a network is called a \textbf{node}. A node can be a
computer, a printer, a switch, or any device with a network address. A
\textbf{host} is a node that can send or receive data over the network, such
as a desktop computer, a laptop, or a network printer. In practice every host
is a node, but the word host is used when the device is an end point that
produces or consumes data.

\begin{center}
\begin{tabular}{p{0.20\textwidth}p{0.70\textwidth}}
\toprule
\textbf{Term} & \textbf{Meaning in this lesson} \\
\midrule
Network & Two or more connected devices that share data and resources. \\
Node & Any participant device in the network. \\
Host & A node (computer, printer) that sends or receives data. \\
Client & The device or program that requests a service. \\
Server & The computer that provides the service and answers requests. \\
\bottomrule
\end{tabular}
\end{center}

The client--server pair is the most examined relationship in this block. The
\textbf{client} is the device (or program) that requests a service --- a web
page, a file, or a print job. The \textbf{server} is the computer that
provides that service: it stores the shared files or pages and answers the
requests of many clients. When the office desktop opens a file stored on the
department's shared machine, the desktop is the client and the shared machine
is the server.

\begin{definitionbox}{Client and server}
The client requests a service; the server stores the shared resource and
provides it in response to client requests.
\end{definitionbox}

\section{Networks by area: PAN, LAN, MAN, WAN}
Networks are classified by the geographical area they cover. Four names must
be learned together with their full forms, because the exam tests both the
expansion and the coverage in the same item.

\begin{center}
\begin{tabular}{p{0.12\textwidth}p{0.30\textwidth}p{0.48\textwidth}}
\toprule
\textbf{Type} & \textbf{Full form} & \textbf{Coverage and example} \\
\midrule
PAN & Personal Area Network & Around one person: phone, laptop, and earphones linked by Bluetooth. \\
LAN & Local Area Network & One building or campus: an office, school, or laboratory network. \\
MAN & Metropolitan Area Network & One city: a bank's branches linked across a town. \\
WAN & Wide Area Network & A country or the world: the Internet is the largest WAN. \\
\bottomrule
\end{tabular}
\end{center}

A \textbf{PAN} (Personal Area Network) covers the smallest area --- the
devices around a single person, such as a phone connected to a laptop and
wireless earphones. A \textbf{LAN} (Local Area Network) covers one building
or campus, such as all the computers of one office floor connected through a
switch. A \textbf{MAN} (Metropolitan Area Network) spans a city, for example
the branches of one organisation linked across a town. A \textbf{WAN} (Wide
Area Network) spans a country or the whole world; the Internet is the largest
example of a WAN.

\begin{definitionbox}{PAN, LAN, MAN, WAN}
PAN covers one person; LAN covers one building or campus; MAN covers one
city; WAN covers a country or the world.
\end{definitionbox}

\begin{example}
A question gives four coverages --- one person, one building, one city, and
the whole world --- and asks for the matching names. Read the area first:
person means PAN, building means LAN, city means MAN, and worldwide means
WAN. The area decides the answer before the full form is even expanded.
\end{example}

\section{WLAN, intranet, Internet, and ISP}
Three further names complete the vocabulary. A \textbf{WLAN} (Wireless Local
Area Network) is a LAN built with wireless signals instead of cables: the
office Wi-Fi through which laptops connect without a cable is a WLAN. An
\textbf{intranet} is a private network inside one organisation, usable only
by its own staff --- for example the company's internal portal. The
\textbf{Internet} is the worldwide public network of connected networks,
open to anyone with access. The two words differ by one prefix and are a
standing trap: intranet is private to one organisation, while the Internet
is public and worldwide.

An \textbf{ISP} (Internet Service Provider) is the company that provides a
customer with access to the Internet. The home or office connects through its
ISP --- over a phone line, cable, fibre, or mobile link --- and the ISP
carries the traffic onward to the Internet.

\begin{definitionbox}{WLAN, intranet, Internet, ISP}
A WLAN is a wireless LAN; an intranet is a private organisational network;
the Internet is the worldwide public network; an ISP (Internet Service
Provider) gives customers access to the Internet.
\end{definitionbox}

\section{Topology: bus, star, ring, mesh, tree}
The word \textbf{topology} means the arrangement by which nodes are connected
in a network. Five topologies are examined: bus, star, ring, mesh, and tree.
Each one decides how data travels and what happens when one cable or device
fails, and the exam repeatedly asks exactly that: which topology has a
central device, and which one fails fully when the main cable breaks.

In a \textbf{bus} topology, all nodes share one single main cable called the
bus. It is simple and cheap for a small network. Data sent by one node
travels along the whole bus and is accepted by the addressed node. Its
weakness follows from its shape: if the main cable breaks, the whole network
fails.

In a \textbf{star} topology, every node connects to one central device (a
hub or a switch), and all data passes through that centre. If one spoke cable
fails, only that node is affected and the rest keep working. But if the
central device itself fails, the whole network stops. The ordinary office LAN
--- one switch with many desktops plugged into it --- is a star, and it is
the most common topology in exam examples.

In a \textbf{ring} topology, each node connects to exactly two neighbours,
forming a closed circle. Data travels around the ring, usually in one
direction, passed from node to node until it reaches its destination. Each
intermediate node repeats the data it is not addressed to. A break at one
point disturbs the whole ring, unless a second (dual) ring provides a backup
path.

In a \textbf{mesh} topology, every node has a direct link to many (in a full
mesh, all) other nodes. It is the most reliable topology: if one path fails,
another path carries the data. It is also the costliest, because it needs the
most cable and the most ports.

In a \textbf{tree} topology, several star networks are joined through a
central backbone, like the branches of a tree. It suits a large campus: each
department is a star, and the stars join at the centre. Its weakness is the
backbone: if the backbone fails, whole branches lose connection.

\begin{center}
\begin{tabular}{p{0.12\textwidth}p{0.30\textwidth}p{0.48\textwidth}}
\toprule
\textbf{Topology} & \textbf{Shape} & \textbf{What one failure does} \\
\midrule
Bus & One shared cable & Main cable breaks: the whole network fails. \\
Star & Centre with spokes & One spoke fails: one node affected; centre fails: all fail. \\
Ring & Closed circle & One break disturbs the ring. \\
Mesh & Many direct links & One path fails: traffic uses another path. \\
Tree & Stars joined on a backbone & Backbone fails: whole branches fail. \\
\bottomrule
\end{tabular}
\end{center}

\begin{definitionbox}{Topology}
The physical or logical arrangement of nodes in a network: bus (one cable),
star (central device), ring (closed circle), mesh (many direct links), and
tree (stars joined on a backbone).
\end{definitionbox}

\begin{example}
An item states, ``In which topology does the failure of the central device
bring down the whole network?'' Trace the table: only the star routes every
message through one centre, so the answer is star. If the stem instead says
the main cable broke and everything stopped, the shape is the single shared
cable, so the answer is bus.
\end{example}

\section{Transmission media: guided and unguided}
\textbf{Transmission media} are the paths that carry data signals from one
device to another. They divide into two families. \textbf{Guided (wired)
media} confine the signal inside a physical cable. \textbf{Unguided
(wireless) media} carry the signal through air or space without a cable.

There are three guided media. \textbf{Twisted-pair cable} consists of pairs
of copper wires twisted together; it is the ordinary telephone and office-LAN
cable --- cheap and suitable for short ranges. \textbf{Coaxial cable} has a
central copper wire with insulation and a metal shield around it; it is used
for cable TV and older LANs and has better shielding than twisted pair.
\textbf{Fibre-optic cable} uses thin glass or plastic strands carrying light
signals; it is the fastest medium with the longest range and is immune to
electrical noise, but it is also the costliest. Ranked by speed and range,
fibre-optic comes first, then coaxial, then twisted pair.

There are four unguided media at this level. \textbf{Radio waves} travel in
all directions and are used for FM radio and for short-range WLAN and
Bluetooth links. \textbf{Microwaves} travel in a straight line and need line
of sight; they are used for point-to-point links between towers.
\textbf{Satellite links} are microwave signals relayed through a satellite;
they cover very long distances but introduce delay. \textbf{Infrared} works
over a very short range, is blocked by walls, and is used for TV remotes.

\begin{definitionbox}{Guided and unguided media}
Guided media carry signals inside a cable (twisted pair, coaxial,
fibre-optic); unguided media carry signals through air or space (radio
waves, microwaves, satellite links, infrared).
\end{definitionbox}

\section{Networking devices and their roles}
Ten devices are examined, and device--role matching is the highest-yield item
type of this lesson. Each device has one role to memorise.

A \textbf{NIC} (Network Interface Card) is the card --- or the built-in
circuit on the motherboard --- that lets a computer physically connect to a
network. Without it the machine has no network port to plug into.

A \textbf{modem} (modulator-demodulator) converts digital computer signals
into analogue signals for a telephone line and back again. It is the device
that connects a home or office to the ISP over a line that cannot carry
digital signals directly.

A \textbf{hub} connects several nodes in a star and broadcasts every
incoming signal to all of its ports; the addressed node accepts the data and
the rest ignore it. A \textbf{switch} also connects nodes in a star, but it
forwards each message only to the addressed port. The switch is therefore
smarter and faster than the hub, because it does not flood every port with
every message.

A \textbf{repeater} regenerates a weak signal so that it can travel further
along a long cable. It works at the signal level only: it does not choose a
path and does not filter traffic. A \textbf{bridge} connects two LAN
segments and filters traffic, passing across only the data meant for the
other side.

A \textbf{router} connects different networks --- for example a home LAN to
the ISP's network --- and chooses the best path for each packet of data. A
\textbf{gateway} connects networks that use different protocols or data
formats, translating between them so the two sides can understand each
other. An \textbf{access point} is the wireless base station that lets Wi-Fi
devices join a wired network. A \textbf{firewall} blocks unauthorised access
while letting legitimate traffic pass; it is a security device that can be
hardware, software, or both.

\begin{center}
\begin{tabular}{p{0.22\textwidth}p{0.68\textwidth}}
\toprule
\textbf{Device} & \textbf{Role (one line)} \\
\midrule
NIC (Network Interface Card) & Connects a computer to the network. \\
Modem & Converts digital to analogue and back; links the site to the ISP. \\
Hub & Broadcasts every incoming signal to all ports. \\
Switch & Forwards each message only to the addressed port. \\
Repeater & Regenerates a weak signal for a longer run. \\
Bridge & Joins two LAN segments and filters traffic between them. \\
Router & Joins different networks and chooses the best path. \\
Gateway & Translates between networks using different protocols. \\
Access point & Lets wireless (Wi-Fi) devices join the wired network. \\
Firewall & Blocks unauthorised access; a security device. \\
\bottomrule
\end{tabular}
\end{center}

\begin{definitionbox}{The four most confused pairs}
Hub sends to all, switch sends to the addressed port; repeater regenerates
the signal, router chooses the path; bridge joins two similar LAN segments,
gateway translates between different protocols; modem connects to the line,
firewall guards the entry.
\end{definitionbox}

\section{Exam retrieval and common confusions}
Six distinctions prevent most errors on this topic.
\begin{enumerate}
  \item The client \textbf{requests} and the server \textbf{provides}. Do not
    swap the roles.
  \item Read the area first: one person is PAN, one building is LAN, one
    city is MAN, worldwide is WAN.
  \item \textbf{Intranet} is private to one organisation; the
    \textbf{Internet} is public and worldwide.
  \item A hub broadcasts to \textbf{all} ports; a switch forwards only to
    the \textbf{addressed} port.
  \item A repeater only \textbf{regenerates} the signal; the router
    \textbf{chooses the path}. A bridge joins two LAN segments; a gateway
    translates between \textbf{different} protocols.
  \item Fibre-optic is the \textbf{fastest} guided medium; twisted pair is
    the slowest of the three cables.
\end{enumerate}

\begin{example}
An item asks which device ``connects different networks and selects the best
path for each packet.'' Walk the device table: the hub and switch stay inside
one network, the repeater only strengthens the signal, and the gateway
translates protocols --- only the router joins different networks and
chooses paths, so the answer is router.
\end{example}

\subsection*{Exercises}
\begin{enumerate}[label=\arabic*.]
  \item Distinguish a node, a host, a client, and a server.
  \item Expand PAN, LAN, MAN, WAN, WLAN, NIC, and ISP, and state the coverage
    of each network type.
  \item Distinguish an intranet from the Internet.
  \item Describe the five topologies and state, for each, what a single
    failure does.
  \item Distinguish guided from unguided transmission media, giving the three
    cables and the four wireless media of this lesson.
  \item Match each of the ten devices of this lesson to its one-line role.
\end{enumerate}

\subsection*{Solutions}
\begingroup\small
\begin{enumerate}[label=\arabic*.]
  \item A node is any participant device in the network. A host is a node
    (computer, printer) that sends or receives data. A client is the device
    or program that requests a service; a server is the computer that
    provides the service and answers client requests.
  \item PAN --- Personal Area Network (around one person); LAN --- Local Area
    Network (one building or campus); MAN --- Metropolitan Area Network (one
    city); WAN --- Wide Area Network (a country or the world); WLAN ---
    Wireless Local Area Network (a LAN built with wireless signals); NIC ---
    Network Interface Card (lets a computer connect to a network); ISP ---
    Internet Service Provider (gives customers Internet access).
  \item An intranet is a private network inside one organisation, usable only
    by its staff. The Internet is the worldwide public network of connected
    networks, open to anyone with access.
  \item Bus: one shared cable; a break in the main cable fails the whole
    network. Star: every node joins one centre; a spoke failure affects one
    node, a centre failure stops all. Ring: closed circle of neighbours; one
    break disturbs the ring. Mesh: many direct links; one failed path leaves
    others. Tree: stars joined on a backbone; a backbone failure drops whole
    branches.
  \item Guided media confine the signal in a cable: twisted-pair (cheap,
    short range), coaxial (shielded, e.g.\ cable TV), and fibre-optic
    (light signals; fastest, longest range, costliest). Unguided media use
    air or space: radio waves (all directions, short range), microwaves
    (straight line, towers), satellite links (long distance, with delay),
    and infrared (very short range, blocked by walls).
  \item NIC connects the computer; modem converts digital--analogue for the
    line to the ISP; hub broadcasts to all ports; switch forwards to the
    addressed port; repeater regenerates a weak signal; bridge joins two LAN
    segments with filtering; router joins different networks and chooses the
    path; gateway translates between different protocols; access point admits
    Wi-Fi devices; firewall blocks unauthorised access.
\end{enumerate}
\endgroup

\lessonheading{41}{Networking Protocols}
\section{Why protocols and addresses matter}
Computers on a network can understand each other only because they follow
agreed rules. A \textbf{protocol} is such a set of rules: it fixes what
messages look like and in which order the machines exchange them. Different
jobs need different protocols. Opening a web page, moving a file, sending an
email, translating a name into an address, and handing out addresses to new
devices are five different jobs, so they use five different protocols. The
exam tests each protocol by its \textbf{exact English name} together with the
single role it performs.

\begin{definitionbox}{Protocol}
A fixed set of rules governing communication between devices on a network,
covering message format and exchange order. TCP, IP, DNS, DHCP, HTTP, HTTPS,
FTP, Telnet, and SMTP are the protocols of this lesson.
\end{definitionbox}

\begin{center}
\begin{tabular}{p{0.24\textwidth}p{0.66\textwidth}}
\toprule
\textbf{Term} & \textbf{Meaning in this lesson} \\
\midrule
TCP & Transmission Control Protocol: reliable, ordered, connection-oriented delivery. \\
UDP & User Datagram Protocol: fast, connectionless, best-effort delivery. \\
IP & Internet Protocol: addressed, best-effort carriage of packets across networks. \\
ICMP & Internet Control Message Protocol: control and diagnostic messages only. \\
DNS & Domain Name System: translates names into IP addresses. \\
DHCP & Dynamic Host Configuration Protocol: automatically assigns addresses. \\
HTTP / HTTPS & HyperText Transfer Protocol (plain) and its Secure version for the web. \\
FTP & File Transfer Protocol: moves files between computers. \\
Telnet & Plain-text remote login protocol. \\
SMTP & Simple Mail Transfer Protocol: sends outgoing mail. \\
OSI model & The seven-layer reference model from Physical to Application. \\
Port & A number naming the service on a machine (HTTP 80, HTTPS 443). \\
\bottomrule
\end{tabular}
\end{center}

The rest of this lesson uses these terms in one consistent sense. \textbf{TCP}
and \textbf{UDP} are the two transport protocols. \textbf{IP} names the
addressing protocol with its two versions, IPv4 and IPv6. \textbf{ICMP} is the
helper that reports errors, never the carrier of user data. \textbf{DNS} and
\textbf{DHCP} are the two address services: one resolves names, the other
assigns addresses. \textbf{HTTP}, \textbf{HTTPS}, \textbf{FTP},
\textbf{Telnet}, and \textbf{SMTP} are application protocols, each with one
job and one well-known port.

\section{TCP and UDP}
The transport layer offers two opposite bargains. \textbf{TCP} (Transmission
Control Protocol) is connection-oriented: it sets up the connection, numbers
every segment, retransmits whatever is lost, and delivers the bytes complete
and in order. What TCP guarantees is \emph{ordering and correctness}; it does
not promise when the data will arrive. \textbf{UDP} (User Datagram Protocol)
is connectionless: it sends each datagram independently, with no setup, no
ordering, and no retransmission. What UDP offers in return is speed and low
overhead.

\begin{definitionbox}{TCP}
A connection-oriented transport protocol that delivers data reliably and in
order through sequencing, acknowledgement, and retransmission.
\end{definitionbox}

\begin{definitionbox}{UDP}
A connectionless transport protocol that sends independent datagrams quickly
with no delivery, ordering, or error-recovery guarantee.
\end{definitionbox}

Choose between them by what the task values most. A file, a web page, or an
email must arrive complete and correct, so it travels over TCP. Live voice, a
video stream, or a quick DNS lookup values speed over perfection, so it
travels over UDP, where a lost fragment is simply skipped.

\begin{example}
An item states, ``TCP guarantees fast delivery of every packet.'' This is
wrong. TCP guarantees that the bytes arrive complete and in order (PYQ-18);
speed is exactly what TCP sacrifices. The fast, best-effort protocol is UDP.
\end{example}

\section{IPv4 and IPv6}
The \textbf{Internet Protocol (IP)} carries addressed packets from the sender
to the receiver across networks. Each device holds an IP address, and routers
read that address to forward the packet hop by hop. IP by itself is
unreliable: delivery is best-effort, and it is TCP above it that adds
reliability.

\textbf{IPv4} uses 32-bit addresses written as four decimal octets, for
example 192.168.1.10. Its header is variable in length (20 to 60 bytes).
\textbf{IPv6} uses 128-bit addresses written as eight hexadecimal groups, for
example 2001:0db8::1, and its header is fixed at 40 bytes (PYQ-18, TRM-10).
The longer address space is the whole point of IPv6: it supplies vastly more
addresses than IPv4.

\begin{definitionbox}{IPv4}
The 32-bit version of the Internet Protocol, with dotted-decimal addresses
and a variable-length header.
\end{definitionbox}

\begin{definitionbox}{IPv6}
The 128-bit version of the Internet Protocol, with hexadecimal addresses and
a fixed 40-byte header.
\end{definitionbox}

\section{ICMP: the control channel}
\textbf{ICMP} (Internet Control Message Protocol) carries control and
diagnostic messages about the network itself: error reports such as
``destination unreachable'' and reachability probes. Its two familiar tools
are \textbf{ping}, which asks whether a host is reachable, and
\textbf{traceroute}, which shows the path packets take. ICMP transfers no
application data and uses no port numbers (TRAP-19). That last point is the
exam's favourite trap: any option giving ICMP a port, a file transfer, or a
mail delivery is wrong.

\section{DNS and DHCP}
\textbf{DNS} (Domain Name System) translates a human name such as
\texttt{jkssb.nic.in} into the IP address the network needs. Without it every
user would have to memorise numeric addresses. DNS queries travel mainly over
UDP on port 53, falling back to TCP for large zone transfers.

\textbf{DHCP} (Dynamic Host Configuration Protocol) automatically assigns an
IP address, subnet mask, default gateway, and DNS server address to each
device that joins the network. Without it every address would have to be
typed in by hand. The home router and the office network each run a DHCP
server; the joining laptop or phone acts as the DHCP client, using UDP ports
67 and 68.

\begin{definitionbox}{DNS}
The distributed naming service that resolves a domain name into its IP
address. It answers ``what number goes with this name?''
\end{definitionbox}

\begin{definitionbox}{DHCP}
The service that automatically assigns IP configuration (address, mask,
gateway, DNS server) to devices joining a network. It answers ``what address
may this newcomer use?''
\end{definitionbox}

\begin{example}
An item asks which service gives a newly connected office desktop its IP
address. Trace the pair: DNS resolves names the user already has, while DHCP
hands out addresses to newcomers. The answer is DHCP. If the stem instead
asks what turns \texttt{jkssb.nic.in} into an address, the answer is DNS.
\end{example}

\section{Application protocols: web, files, remote login, mail}
\textbf{HTTP} (HyperText Transfer Protocol, port 80) fetches web pages in
plain text. \textbf{HTTPS} (HyperText Transfer Protocol Secure, port 443) is
the same protocol with an SSL/TLS encryption layer, shown by
\texttt{https://} and the padlock. The S stands for Secure (WO 2026 Q76
territory).

\textbf{FTP} (File Transfer Protocol, port 21) moves files between computers;
it is neither email nor web browsing. \textbf{Telnet} (port 23) opens a
plain-text remote login session and is today replaced by the encrypted SSH.
\textbf{SMTP} (Simple Mail Transfer Protocol, port 25) sends outgoing mail to
the server. Receiving is a separate job handled by \textbf{POP3} (Post Office
Protocol version 3, port 110), which downloads mail to the local machine
(PYQ-16, TRAP-17), and \textbf{IMAP} (Internet Message Access Protocol, port
143), which keeps mail synced on the server.

\section{The seven OSI layers}
The \textbf{OSI model} (Open Systems Interconnection) is a seven-layer
reference frame for describing network communication:

\begin{center}
\begin{tabular}{p{0.10\textwidth}p{0.22\textwidth}p{0.58\textwidth}}
\toprule
\textbf{No.} & \textbf{Layer} & \textbf{Job} \\
\midrule
7 & Application & Services the user sees: HTTP, FTP, SMTP, DNS, DHCP. \\
6 & Presentation & Translation, encryption, compression. \\
5 & Session & Opens, manages, and closes sessions between machines. \\
4 & Transport & End-to-end delivery: TCP and UDP. \\
3 & Network & Addressing and routing: IP (IPv4, IPv6) and ICMP. \\
2 & Data Link & Frames on a single link: switch, MAC address. \\
1 & Physical & Bits on the wire, cable, or radio signal. \\
\bottomrule
\end{tabular}
\end{center}

Data moves down the layers on sending and up the layers on receiving, each
layer adding (or reading) its own header. For this exam, three placements
carry almost all the marks: the application names (HTTP, FTP, SMTP, DNS,
DHCP) at layer 7, TCP and UDP at layer 4, and IP and ICMP at layer 3.

\begin{example}
Opening a web page uses HTTP at the Application layer, TCP at the Transport
layer, and IP at the Network layer: one protocol per layer, working
together. An item asking ``which layer does IP belong to'' is answered by
the same table: the Network layer (layer 3).
\end{example}

\section{Port awareness}
An \textbf{IP address} names the machine; a \textbf{port number} names the
service on that machine. Ports 0--1023 are the well-known ports. The exam set
is: FTP 21, Telnet 23, SMTP 25, DNS 53, HTTP 80, POP3 110, IMAP 143, HTTPS
443. DHCP uses UDP ports 67 and 68, and ICMP uses no ports at all. The
recommended drill is to cover the service names and recall the numbers in the
chain 21--23--25--53--80--443.

\section{Exam retrieval and common confusions}
Seven distinctions prevent most errors on this topic.
\begin{enumerate}
  \item IPv6 is \textbf{128-bit}, never 256-bit; its header is fixed at 40
    bytes (TRAP-18, TRM-10).
  \item ICMP is \textbf{control/diagnostic} only: no application data, no
    ports (TRAP-19).
  \item TCP promises \textbf{order}, not timing; UDP promises neither
    (PYQ-18).
  \item \textbf{SMTP sends} mail; \textbf{POP3/IMAP receive} it (TRAP-17).
  \item \textbf{DNS resolves} names to addresses; \textbf{DHCP assigns}
    addresses to joining devices.
  \item HTTP is port \textbf{80} and plain; HTTPS is port \textbf{443} with
    SSL/TLS.
  \item TCP and UDP live at the \textbf{Transport} layer; IP and ICMP at the
    \textbf{Network} layer; HTTP, FTP, SMTP, DNS, and DHCP at the
    \textbf{Application} layer.
\end{enumerate}

\subsection*{Exercises}
\begin{enumerate}[label=\arabic*.]
  \item Distinguish TCP from UDP on reliability, connection setup, and one
    typical use each.
  \item Contrast IPv4 with IPv6 on address length, notation, and header size.
  \item Explain why ICMP is described as a control protocol and name its two
    common tools.
  \item Distinguish DNS from DHCP with one sentence on each role.
  \item List the seven OSI layers in order with the protocols of this lesson
    placed at their layers.
\end{enumerate}

\subsection*{Solutions}
\begingroup\small
\begin{enumerate}[label=\arabic*.]
  \item TCP is connection-oriented, reliable, and ordered: it numbers
    segments and retransmits losses; typical uses are web pages, email, and
    file transfer. UDP is connectionless and best-effort: no setup, ordering,
    or retransmission; typical uses are streaming, voice calls, and DNS
    queries.
  \item IPv4 is 32-bit, dotted-decimal (four octets), with a variable 20--60
    byte header. IPv6 is 128-bit, hexadecimal (eight groups), with a fixed
    40-byte header.
  \item ICMP carries control and diagnostic messages about the network
    (errors, reachability) rather than user data, and it uses no ports. Its
    tools are ping (reachability) and traceroute (path).
  \item DNS translates a domain name into its IP address. DHCP automatically
    assigns an IP address and related configuration to a device joining the
    network.
  \item 7 Application (HTTP, HTTPS, FTP, Telnet, SMTP, POP3, IMAP, DNS,
    DHCP); 6 Presentation; 5 Session; 4 Transport (TCP, UDP); 3 Network (IP,
    ICMP); 2 Data Link; 1 Physical.
\end{enumerate}
\endgroup

\lessonheading{62}{Current Internet, Networking and Cloud IaaS}
\section{TCP: ordered and reliable, but with no timing promise}
A running Internet conversation must survive lost, duplicated, and reordered
packets. \textbf{TCP} (Transmission Control Protocol) is the transport
protocol that hides this mess from the application. It is
\emph{connection-oriented}: the two ends set up a connection before any data
moves. Once the connection exists, TCP numbers every byte, acknowledges what
arrives, and retransmits whatever is lost, so the receiving application sees
the bytes in exactly the order they were sent. That is the double guarantee
to memorise: \textbf{ordered} delivery and \textbf{reliable} delivery.

What TCP does not promise is timing. It never claims the data will arrive
within a fixed time, because retransmissions take as long as they take. This
``no timing guarantee'' point is tested directly (PYQ-18): an option that
claims TCP promises delivery by a deadline is wrong.

\begin{definitionbox}{TCP}
A connection-oriented transport protocol that delivers bytes in order and
retransmits lost segments, with no guarantee about delivery time.
\end{definitionbox}

\begin{center}
\begin{tabular}{p{0.24\textwidth}p{0.66\textwidth}}
\toprule
\textbf{Term} & \textbf{Meaning in this lesson} \\
\midrule
TCP & Transmission Control Protocol; ordered and reliable transport. \\
UDP & User Datagram Protocol; connectionless, fast, unreliable transport. \\
ICMP & Internet Control Message Protocol; control and diagnostic messages only. \\
BGP & Border Gateway Protocol; inter-domain routing between autonomous systems. \\
IPv6 & 128-bit next-generation Internet addressing; no broadcast. \\
QUIC & Quick UDP Internet Connections; the UDP-based transport under HTTP/3. \\
DNS & Domain Name System; translates domain names into IP addresses. \\
PIR & Public Interest Registry; operates the .org registry. \\
ICANN & Internet Corporation for Assigned Names and Numbers; global name/number coordination. \\
ARPANET & Advanced Research Projects Agency Network (1969); ancestor of the Internet. \\
IaaS & Infrastructure as a Service; provider manages hardware, customer manages OS and above. \\
\bottomrule
\end{tabular}
\end{center}

The rest of this lesson uses these terms in one consistent sense. \textbf{TCP}
means ordered and reliable transport. \textbf{UDP} means the fast,
connectionless alternative with neither guarantee. \textbf{ICMP} means the
control-message protocol that never carries application data.
\textbf{BGP} means routing between autonomous systems, never inside one LAN.
\textbf{IPv6} means 128-bit addresses with no broadcast. \textbf{QUIC} means
the UDP-based transport that HTTP/3 uses instead of TCP. \textbf{DNS} means
the name-to-address directory. \textbf{PIR} runs the .org registry under
\textbf{ICANN} coordination. \textbf{ARPANET} is the 1969 packet-switched
ancestor of the Internet. \textbf{IaaS} splits responsibility so that the
provider manages networking hardware while the customer manages the operating
system, applications, and data.

\section{UDP and the TCP/UDP contrast}
\textbf{UDP} (User Datagram Protocol) is TCP's near-neighbour and the most
common distractor for it. UDP is \emph{connectionless}: each datagram is sent
independently with no setup. It provides no ordering guarantee, no
reliability guarantee, and no retransmission. Its header is small, so it is
faster and cheaper per packet than TCP.

\begin{definitionbox}{UDP}
A connectionless transport protocol that sends independent datagrams with no
ordering, reliability, or retransmission guarantee.
\end{definitionbox}

\begin{center}
\begin{tabular}{p{0.20\textwidth}p{0.35\textwidth}p{0.35\textwidth}}
\toprule
\textbf{Point} & \textbf{TCP} & \textbf{UDP} \\
\midrule
Connection & Connection-oriented & Connectionless \\
Ordering & Guaranteed & Not guaranteed \\
Reliability & Retransmits lost data & No retransmission \\
Speed & Slower (acknowledgements, state) & Faster (light header) \\
Typical use & Web pages, mail, file transfer & Live voice, video, games, DNS lookups \\
\bottomrule
\end{tabular}
\end{center}

The exam tests this pair by swapping one property: an option that says UDP
guarantees order, or that TCP is connectionless, is always wrong. The anchor
is the direction of the trade: TCP buys correctness at the cost of speed,
while UDP buys speed at the cost of correctness.

\begin{example}
An item states, ``UDP guarantees ordered delivery of datagrams.'' This is
wrong. Ordering is a TCP guarantee; UDP provides no ordering at all. The
wrong option has moved a TCP property onto UDP.
\end{example}

\section{ICMP: the network's own messenger}
While TCP and UDP carry user data, \textbf{ICMP} (Internet Control Message
Protocol) carries messages \emph{about} the network. Routers and hosts use it
to report errors --- destination unreachable, time exceeded, parameter
problems --- and diagnostic tools such as \textbf{ping} (echo request and
reply) and \textbf{traceroute} are built on ICMP messages.

\begin{definitionbox}{ICMP}
The control and diagnostic protocol of the Internet layer; reports errors and
supports ping and traceroute, and never carries application data.
\end{definitionbox}

Three facts together close every ICMP trap (PYQ-18, PYQ-19, TRAP-19). First,
ICMP transfers no application data: it never moves a file, a mail message,
or a web page. Second, it uses no ports: ports belong to TCP and UDP, and an
option that assigns ICMP a port number is wrong. Third, it is neither a
routing protocol like BGP nor a mail protocol like SMTP; it only reports on
conditions the other protocols encounter.

\section{BGP: routing between autonomous systems}
The Internet is not one network but a federation of large independently
managed networks called \textbf{autonomous systems} (an ISP, a university
backbone, a government network). \textbf{BGP} (Border Gateway Protocol) is
the protocol that routes \emph{between} them. It exchanges reachability
information across autonomous-system borders and selects the paths that
traffic takes over the Internet backbone.

\begin{definitionbox}{BGP}
The inter-domain routing protocol that exchanges reachability between
autonomous systems on the Internet backbone.
\end{definitionbox}

The tested contrast is BGP against the \emph{interior} protocols
\textbf{OSPF} (Open Shortest Path First) and \textbf{RIP} (Routing
Information Protocol), which route \emph{inside} a single domain (PYQ-19,
TRAP-20). Any option that places BGP inside one small local network has
swapped the two scopes. The one-word shortcut is ``inter-domain'': whenever
the stem says inter-domain or between autonomous systems, the answer is BGP.

\section{IPv6: address space, header, and no broadcast}
\textbf{IPv4} uses 32-bit addresses (about 4.3 billion), which ran out under
the growth of the Internet. \textbf{IPv6} replaces them with \textbf{128-bit}
addresses, written as eight groups of hexadecimal digits, giving about
$3.4 \times 10^{38}$ addresses. This vast space is the standard reason given
for IPv6 removing the need for address-sharing workarounds such as
\textbf{NAT} (Network Address Translation), which lets many private addresses
share one public IPv4 address (PYQ-19).

\begin{definitionbox}{IPv6 address}
A 128-bit Internet-layer address, written as eight hexadecimal groups, with a
fixed 40-byte header.
\end{definitionbox}

Two further IPv6 facts are tested. The fixed IPv6 header is \textbf{40 bytes}
long (PYQ-18). And IPv6 has \textbf{no broadcast}: IPv4's three delivery
styles (unicast to one receiver, multicast to a group, broadcast to everyone)
become unicast, multicast, and anycast (nearest of several receivers) in
IPv6, with a multicast to the all-nodes group replacing the broadcast. An
option mentioning an ``IPv6 broadcast address'' is therefore self-defeating.
The 256-bit figure that appears as a distractor is always wrong (TRAP-18).

\begin{example}
An item offers ``IPv6 uses 256-bit addresses'' as one option. This is wrong.
IPv6 uses 128-bit addresses; 256 is double the correct size and matches no
Internet-layer standard. The option confuses IPv6 with an unrelated larger
number.
\end{example}

\section{HTTP/3 and QUIC}
Web traffic travelled over TCP for decades: HTTP/1.1 and HTTP/2 both run
over TCP. \textbf{HTTP/3} breaks that pattern. It runs over \textbf{QUIC}
(Quick UDP Internet Connections), which itself runs over \textbf{UDP}, not
over TCP. QUIC rebuilds the needed reliability, congestion control, and
security on top of UDP, gaining faster connection setup and independent
streams without head-of-line blocking across streams.

\begin{definitionbox}{QUIC}
A UDP-based transport protocol carrying HTTP/3, providing its own reliability
and security instead of using TCP.
\end{definitionbox}

The trap writes itself: any option claiming HTTP/3 relies on TCP has the
protocol stack backwards (PYQ-19, TRAP-21). The correct chain is HTTP/3 over
QUIC over UDP.

\section{DNS and the .org / PIR / ICANN chain}
\textbf{DNS} (Domain Name System) is the Internet's directory: it translates
human-readable domain names such as \texttt{example.org} into numeric IP
addresses that routers can use. Resolution runs hierarchically from root
servers down through top-level-domain servers to the domain's own
authoritative servers. DNS is a lookup service, not a routing protocol and
not a mail-transfer protocol, and confusing it with either is a standard
wrong answer.

The governance half of the topic names three actors. \textbf{.org} is a
generic top-level domain originally meant for non-profit organisations.
Its registry is operated by \textbf{PIR} (Public Interest Registry).
\textbf{ICANN} (Internet Corporation for Assigned Names and Numbers)
coordinates domain names and IP addresses at the global level. The exam
distinction is role-based: PIR \emph{runs} the .org registry, while ICANN
\emph{coordinates} the naming system as a whole.

\section{ARPANET and packet switching}
\textbf{ARPANET} (Advanced Research Projects Agency Network), live in
\textbf{1969}, is the ancestor of the Internet. Its contribution was proving
that \textbf{packet switching} works across long distances: splitting data
into addressed packets that travel independently and are reassembled at the
destination, with the network routing around failed links. Later ARPANET work
produced the TCP/IP protocol family.

\begin{definitionbox}{Packet switching}
Splitting data into addressed packets that may take different routes and are
reassembled at the destination.
\end{definitionbox}

The tested contrast is packet switching against \textbf{circuit switching},
the older telephone model that holds one dedicated path for the whole call.
The Internet uses packet switching; the traditional telephone call used
circuit switching. Options that assign a dedicated reserved circuit to the
Internet have the two models swapped.

\section{The router: joining networks by reading IP addresses}
A \textbf{router} connects different networks and forwards packets between
them. It does this by reading each packet's destination \textbf{IP address}
and choosing the next hop from its \textbf{routing table}. In
assertion--reason form: the assertion (a router joins networks) and the
reason (it reads the destination IP and consults its table) are both true,
and the reason correctly explains the assertion, because address-reading is
\emph{how} the router forwards between networks.

The near-neighbours are the \textbf{switch}, which forwards \emph{inside} one
network using MAC addresses, and the \textbf{hub}, which merely repeats
everything it receives. An option that says a router forwards within one
network using MAC addresses describes a switch, not a router.

\section{IaaS: the provider manages the hardware}
Cloud service models split responsibility between provider and customer.
In \textbf{IaaS} (Infrastructure as a Service), the \textbf{provider} manages
the physical infrastructure: servers, storage hardware, \textbf{networking
hardware}, and virtualisation. The \textbf{customer} manages everything above
that layer: the operating system, runtime, applications, and data. The exam
anchor is the provider's side: in IaaS the provider manages the networking
hardware (PYQ-20).

\begin{definitionbox}{IaaS responsibility split}
The provider manages infrastructure (networking and storage hardware,
virtualisation); the customer manages the operating system, runtime,
applications, and data.
\end{definitionbox}

Moving down the ladder, the provider manages more. In \textbf{PaaS}
(Platform as a Service) the provider additionally manages the operating
system and runtime, leaving the customer its applications and data. In
\textbf{SaaS} (Software as a Service) the provider manages everything up to
the finished application, and the customer only manages its own data and
settings. The trap claims the IaaS provider manages the customer's OS or
data; it never does (TRAP-22).

\begin{example}
An item states, ``In IaaS the provider manages the customer's operating
system and data.'' This is wrong. Those belong to the customer; the provider
stops at the infrastructure layer (networking and storage hardware plus
virtualisation). The option describes SaaS, not IaaS.
\end{example}

\section{Exam retrieval and common confusions}
Seven distinctions prevent most errors on this topic.
\begin{enumerate}
  \item TCP guarantees \textbf{order and reliability} but no timing; UDP
    guarantees \textbf{neither} (PYQ-18).
  \item ICMP carries \textbf{control messages only}: no application data, no
    ports (TRAP-19).
  \item BGP is \textbf{inter-domain} (between autonomous systems); OSPF and
    RIP are intra-domain (TRAP-20).
  \item IPv6 is \textbf{128-bit} with a \textbf{40-byte} header; 256-bit is
    always wrong (PYQ-18, TRAP-18).
  \item IPv6 has \textbf{no broadcast}; the replacement is
    \textbf{multicast} to the all-nodes group (PYQ-19).
  \item HTTP/3 uses \textbf{QUIC over UDP}, never TCP (TRAP-21).
  \item In IaaS the \textbf{provider} manages networking hardware; the
    \textbf{customer} manages the OS, applications, and data (TRAP-22).
\end{enumerate}

\begin{example}
An item asks which protocol routes between autonomous systems on the
Internet backbone. Trace the scope word: ``between autonomous systems'' is
the definition of inter-domain routing, and the inter-domain protocol is BGP.
OSPF and RIP are interior protocols, and ICMP is not a routing protocol at
all, so all three are eliminated by scope.
\end{example}

\subsection*{Exercises}
\begin{enumerate}[label=\arabic*.]
  \item State the two guarantees TCP provides and the one it does not.
  \item Distinguish TCP from UDP on connection, ordering, reliability, and
    speed.
  \item What does ICMP carry, and what two things does it never do?
  \item Distinguish BGP from OSPF/RIP by scope.
  \item State the IPv6 address size and header size, and explain what replaces
    the broadcast.
  \item Explain the HTTP/3 protocol stack in one chain of three names.
  \item Distinguish the roles of PIR and ICANN for the .org domain.
  \item State the IaaS responsibility split in two sentences.
\end{enumerate}

\subsection*{Solutions}
\begingroup\small
\begin{enumerate}[label=\arabic*.]
  \item TCP guarantees ordered delivery and reliable delivery (lost segments
    are retransmitted). It provides no timing guarantee.
  \item TCP is connection-oriented, ordered, and reliable but slower; UDP is
    connectionless, unordered, and unreliable (no retransmission) but faster.
  \item ICMP carries control and diagnostic messages (errors, ping,
    traceroute). It never transfers application data and never uses ports.
  \item BGP routes between autonomous systems (inter-domain, Internet
    backbone); OSPF and RIP route inside a single domain.
  \item IPv6 uses 128-bit addresses with a fixed 40-byte header. It has no
    broadcast; a multicast to the all-nodes group replaces it.
  \item HTTP/3 runs over QUIC, which runs over UDP: HTTP/3 over QUIC over
    UDP.
  \item PIR operates the .org registry; ICANN coordinates domain names and
    IP addresses globally.
  \item The provider manages the infrastructure (networking and storage
    hardware, virtualisation); the customer manages the operating system,
    runtime, applications, and data.
\end{enumerate}
\endgroup

\unitopener{10}{Cybersecurity \& IT Law}
\lessonheading{42}{Malware}
\section{What ``malware'' means}
Malware is short for \textbf{malicious software}: any program written to
harm a computer, steal from it, or misuse it. The official syllabus names
\textbf{computer virus and latest anti-virus} explicitly, and the exam's
computer block tests the near neighbours as well: worm, Trojan, ransomware,
spyware, adware, rootkit, keylogger, bot and botnet, backdoor, logic bomb,
and cryptojacking.

\begin{definitionbox}{Malware}
Any software designed to damage a computer, steal data from it, or misuse
its resources without the owner's consent.
\end{definitionbox}

\begin{center}
\begin{tabular}{p{0.24\textwidth}p{0.66\textwidth}}
\toprule
\textbf{Term} & \textbf{Meaning in this lesson} \\
\midrule
Virus & Malware that attaches to a host file and needs human action to run. \\
Worm & Standalone, self-spreading malware that travels over networks. \\
Trojan (Trojan horse) & Malware disguised as useful software; does not replicate itself. \\
Ransomware & Malware that locks files and demands payment for release. \\
Spyware & Malware that secretly watches the user and steals data. \\
Adware & Malware that floods the machine with unwanted advertisements. \\
Keylogger & A spying tool that records every key pressed. \\
Rootkit & Malware that hides itself and other malware deep in the system. \\
Backdoor & A secret entry path into the system bypassing normal checks. \\
Bot / botnet & A remotely controlled machine / a network of such machines. \\
Logic bomb & Malicious code that waits for a trigger, then strikes. \\
Cryptojacking & Secret use of the victim's computer to mine cryptocurrency. \\
Anti-virus & Software that detects, blocks, and removes known malware. \\
\bottomrule
\end{tabular}
\end{center}

The rest of this lesson uses these terms in one consistent sense. A
\textbf{virus} always needs a host file and a human launch. A \textbf{worm}
needs neither. A \textbf{Trojan} pretends to be legitimate and does not
replicate. \textbf{Ransomware} encrypts and demands money.
\textbf{Spyware} hides and steals; \textbf{adware} shows ads. A
\textbf{keylogger} records keystrokes. A \textbf{rootkit} conceals; a
\textbf{backdoor} re-admits. A \textbf{botnet} is a controlled network (robot
+ network). A \textbf{logic bomb} waits for its trigger.
\textbf{Cryptojacking} mines on the victim's machine.

\section{Virus, worm, and Trojan}
A \textbf{virus} attaches itself to a host program or file and stays inert
until a person opens, runs, or shares that file. Common carriers are email
attachments, pen drives, and downloaded files. Once launched it can damage,
delete, or corrupt files. The defining limit is dependence: no host file
and no human action, no spread.

\begin{definitionbox}{Virus}
Malware that attaches to a host program or file and requires human action
to execute and spread.
\end{definitionbox}

A \textbf{worm} is standalone and self-replicating. It copies itself from
machine to machine over network connections, exploiting security gaps,
without any host file and without any user click. Its main harm is
consumption: networks clog, systems slow, and services jam under the load
of endless self-copies.

\begin{definitionbox}{Worm}
Standalone malware that replicates and spreads across networks by itself,
without a host file or human action.
\end{definitionbox}

A \textbf{Trojan}, or Trojan horse, takes its name from deception: it
arrives looking like something useful --- a free game, a handy tool, a
harmless document. The user installs or opens it willingly, and the hidden
malicious part runs alongside. Unlike a virus or a worm, a Trojan does not
replicate itself; each new machine needs its own act of trickery. Once
inside, it can steal data, install spyware or ransomware, or open a backdoor
for the attacker's later return.

\begin{definitionbox}{Trojan (Trojan horse)}
Malware disguised as legitimate software that tricks the user into
installing it; it does not replicate itself.
\end{definitionbox}

The three are best remembered as one table: a virus needs a host file and a
user action; a worm needs neither; a Trojan needs the user's trust and does
not self-replicate. File damage points to a virus, network-wide slowdown
with no clicks points to a worm, and a ``useful'' download that betrays the
machine points to a Trojan.

\begin{example}
An item describes a program that ``spreads from computer to computer over
the network without any user opening a file.'' Trace the table: needs no
host and no click means worm. If instead the stem says the malware ``was
installed as part of a free utility the user downloaded,'' the deception
plus no replication means Trojan.
\end{example}

\section{Ransomware}
\textbf{Ransomware} locks the victim out of their own data, usually by
encrypting files (or the whole screen), and then displays a demand: pay the
ransom to receive the unlock key. The ``pay to unlock'' message, often with
a countdown or a payment address, is its signature.

Ransomware usually arrives the Trojan way --- a harmful attachment or link
the user opens --- and then encrypts documents, photos, and databases in
minutes. The expected defences are regular \textbf{offline backups}, updated
operating system and \textbf{anti-virus} software, and
caution with attachments. The exam's standing line is that paying the ransom
does not guarantee recovery; the backup is what defeats the demand.

\begin{definitionbox}{Ransomware}
Malware that encrypts the victim's files (or locks the screen) and demands
payment for the key to restore them.
\end{definitionbox}

\section{Spyware, adware, and the keylogger}
\textbf{Spyware} is the silent watcher. It hides on the machine, observes
browsing habits, logins, and personal details, and sends them to a third
party. It deliberately causes no visible damage so the victim suspects
nothing. It commonly arrives bundled with free downloads or deceptive
installers, and its purpose is theft by observation: passwords, bank
details, and identity data.

\textbf{Adware} is the noisy opposite. It floods the screen with unwanted
advertisements --- pop-ups, banners, redirected searches --- to earn money
from ad traffic. Heavy adware also slows the machine and may carry spyware
along with it, so the exam treats persistent pop-ups as more than a
nuisance.

A \textbf{keylogger} is a focused spying tool: it records every key pressed
and forwards the log to the attacker, harvesting usernames, passwords, and
card numbers. Keyloggers exist as software and as small hardware devices
fitted between the keyboard and the computer.

\begin{definitionbox}{Spyware vs adware}
Spyware hides and steals information quietly; adware displays unwanted
advertisements openly. A keylogger is a spying tool that records keystrokes.
\end{definitionbox}

\begin{example}
An item reports ``constant pop-up advertisements and redirected searches.''
The visible ads mean adware. If instead the item reports ``passwords stolen
although nothing on screen looked wrong,'' the silent theft means spyware,
and if it specifies ``every keystroke was recorded,'' the tool is a
keylogger.
\end{example}

\section{Rootkit and backdoor}
A \textbf{rootkit} is malware whose job is concealment. It lodges deep in
the system, often with administrator-level control, and hides files,
processes, and network activity from the user and from ordinary detection
tools --- including hiding other malware.

A \textbf{backdoor} is the secret entry a rootkit or Trojan protects: a
path into the system that bypasses normal login and security checks, left
open so the attacker can return at will. The pairing to remember is that
Trojans and rootkits commonly install backdoors; the rootkit hides the
intruder, and the backdoor keeps the door open.

\begin{definitionbox}{Rootkit and backdoor}
A rootkit conceals an intruder and its malware deep in the system; a
backdoor is a hidden entry path that bypasses normal authentication.
\end{definitionbox}

\section{Bot and botnet, logic bomb, cryptojacking}
A \textbf{bot} is an ordinary computer that has been compromised and is now
obeyed remotely; a \textbf{botnet} is a network of such machines (the word
joins robot + network) under one master. Botnets send spam, spread further
malware, and launch large-scale attacks that no single machine could mount.
The owner typically notices only slowness.

A \textbf{logic bomb} is malicious code already planted that waits for a
trigger --- a date, an event, or a specific action --- and strikes only when
the condition is met. It does not spread; it waits.

\textbf{Cryptojacking} is theft of computing power: the attacker secretly
runs cryptocurrency mining on the victim's machine, so the computer runs hot
and slow while the attacker collects the earnings. Like a bot, its outward
sign is an otherwise unexplained slowdown.

\section{How malware reaches a machine, and how to stop it}
Malware reaches users through a short list of routes: harmful email
attachments and links, infected pen drives, bundled free-software
installers, downloads from untrusted sites, and unpatched system software
with known gaps.

The defences form a matching list. \textbf{Anti-virus} software detects,
blocks, quarantines, and removes known malware; both it and the operating
system must be kept updated, since each update closes gaps and refreshes the
list of known threats. A \textbf{firewall} filters incoming and outgoing
network traffic and blocks suspicious connections, but it is a barrier, not
a cleaner: it does not remove malware already inside. Safe habits complete
the set: do not open unknown attachments or links, scan pen drives before
use, download only from trusted sources, and keep regular \textbf{offline
backups} so ransomware loses its leverage.

\section{Exam retrieval and common confusions}
Six distinctions prevent most errors on this topic.
\begin{enumerate}
  \item A virus needs a host file \textbf{and} a human action; a worm needs
    \textbf{neither}; a Trojan disguises itself and \textbf{does not
    replicate}.
  \item Spyware \textbf{hides and steals}; adware \textbf{shows ads}. A
    keylogger records \textbf{keystrokes}.
  \item Ransomware \textbf{encrypts and demands payment}; its signature is
    the ``pay to unlock'' message, and the defence is the \textbf{offline
    backup}.
  \item A rootkit \textbf{hides} the intruder; a backdoor is the
    \textbf{secret re-entry}. Trojans commonly install backdoors.
  \item A botnet is a \textbf{network} of remotely controlled machines
    (robot + network), used for spam and large attacks.
  \item A logic bomb \textbf{waits for a trigger}; cryptojacking
    \textbf{mines cryptocurrency} on the victim's machine.
  \item Anti-virus \textbf{removes} malware; a firewall \textbf{filters}
    traffic and does not clean an infected machine.
\end{enumerate}

\begin{example}
An item asks which malware ``encrypts the user's files and demands money
for the decryption key.'' The lock plus the demand means ransomware ---
not a worm (self-spreading), not adware (ads), not spyware (silent theft).
If the stem instead says the program ``waits until a fixed date and then
deletes files,'' the waiting trigger means a logic bomb.
\end{example}

\subsection*{Exercises}
\begin{enumerate}[label=\arabic*.]
  \item Distinguish a virus from a worm from a Trojan on three points: host
    file, self-spreading, and user action.
  \item Contrast spyware with adware. Where does a keylogger fit?
  \item Explain what ransomware does and name the defence that defeats its
    demand.
  \item Distinguish a rootkit from a backdoor, and explain how they work
    together.
  \item Define bot/botnet, logic bomb, and cryptojacking in one sentence
    each, naming each type's outward sign.
\end{enumerate}

\subsection*{Solutions}
\begingroup\small
\begin{enumerate}[label=\arabic*.]
  \item A virus needs a host file and a human launch and does not spread by
    itself; a worm is standalone and self-spreading over networks with no
    host and no click; a Trojan pretends to be legitimate software, is
    installed willingly, and does not replicate itself.
  \item Spyware hides and secretly steals user data; adware openly floods
    the machine with unwanted ads. A keylogger is a focused spying tool that
    records every keystroke to steal credentials.
  \item Ransomware encrypts the victim's files (or locks the screen) and
    demands payment for the unlock key. Regular offline backups defeat the
    demand, since the data can be restored without paying.
  \item A rootkit hides the intruder and its malware deep in the system; a
    backdoor is a secret entry bypassing normal checks. Trojans and rootkits
    install backdoors so the attacker can return unseen.
  \item A bot is a remotely controlled compromised machine and a botnet is a
    network of them (sign: slowness during spam or attacks); a logic bomb
    waits for a trigger such as a date and then strikes (sign: a timed
    surprise); cryptojacking secretly mines cryptocurrency on the victim's
    machine (sign: hot, slow computer).
\end{enumerate}
\endgroup

\lessonheading{43}{Hacking, Phishing and Spoofing}
\section{Why attacks succeed}
Most attacks on ordinary computer users do not defeat the machine by
force. They persuade the person in front of the screen to hand over the
key: a password, a One-Time Password, a click, or a file download. The
attacker needs a weakness to use and a moment in which the victim acts
without verifying. The weakness may be technical, such as a short reused
password or outdated software, or human, such as trust in an
authoritative-sounding caller, hurry, or fear of an account being
blocked. This lesson stays at a safe awareness level: it names each
attack, shows how to recognise it, and teaches the correct defensive
response. It does not teach how to carry out any attack.

\begin{definitionbox}{Attack chain}
A \textbf{vulnerability} is a weakness in software, settings, or human
habit that can be misused. An \textbf{exploit} is a method that takes
advantage of such a weakness. An \textbf{attack} is the complete harmful
act against a person or system. The \textbf{payload} is the harmful
element the attack delivers, such as stolen credentials, stolen money,
or a malicious file.
\end{definitionbox}

\begin{center}
\begin{tabular}{p{0.24\textwidth}p{0.66\textwidth}}
\toprule
\textbf{Term} & \textbf{Meaning in this lesson} \\
\midrule
Hacking & Gaining unauthorised access to a computer, account, or network. \\
Phishing & A fake message, mail, or call imitating a trusted sender to steal credentials or money. \\
Spoofing & Faking an identity, such as a sender address, caller number, or website address. \\
Social engineering & Tricking a person, rather than breaking software, into giving access. \\
Password attack & An attempt to discover a password by guessing, a word list, exhaustive trial, or observation. \\
Malicious link & A web address that leads to a fake page or an unwanted download. \\
Malicious attachment & A file that harms the computer or steals data when opened. \\
OTP & One-Time Password: a single-use code; never shared with anyone. \\
\bottomrule
\end{tabular}
\end{center}

The rest of this lesson uses these terms in one consistent sense.
\textbf{Hacking} is the broad act of unauthorised access.
\textbf{Phishing} is fraud by imitation that asks the victim to act.
\textbf{Spoofing} is the forged identity that makes the fraud look
genuine. \textbf{Social engineering} is the human manipulation behind
both. A \textbf{vulnerability} is the weakness, the \textbf{exploit} is
the method, and the \textbf{payload} is the harm delivered.

\section{Hacking}
\textbf{Hacking} means gaining \textbf{unauthorised access} to a
computer system, an online account, or a network. A \textbf{hacker} is
the person who performs this act. Access without permission is unlawful
even when nothing is copied or damaged; the absence of permission is
what defines the act, not the amount of harm that follows.

An \textbf{ethical hacker} is a security tester who probes systems
\textbf{with written permission} from the owner in order to find and fix
weaknesses before criminals exploit them. The permission is the entire
difference: the same technical step performed with written authorisation
is a security test, and performed without it is an attack.

\begin{definitionbox}{Hacking}
Gaining access to a computer, account, or network without the owner's
permission.
\end{definitionbox}

\begin{example}
A student opens a colleague's mailbox by guessing the password and
reads the mail. Even though nothing is deleted, this is hacking,
because the access was unauthorised. A bank that hires a tester to try
its login page, with a signed agreement, is instead conducting an
ethical security test.
\end{example}

\section{Password attacks}
Passwords remain the most common lock on an account, so attackers try
to discover them first. Four methods appear repeatedly in examinations.

\textbf{Guessing} tries obvious choices: the user's name, birth year,
phone number, or common strings such as \texttt{password123}.
\textbf{Dictionary attack} tries every entry in a list of common
passwords automatically, which succeeds quickly against short familiar
words. \textbf{Brute-force attack} tries every possible combination of
characters until one works; it always succeeds eventually, but a longer
password makes the search impractically slow. \textbf{Shoulder surfing}
is direct observation: watching the keyboard or screen while the victim
types a password or PIN.

The defence follows from the methods. Use a \textbf{long} password that
mixes letters, digits, and symbols, and never reuse one password across
accounts. Never disclose a password or an \textbf{OTP (One-Time
Password)} to anyone, including a caller who claims to represent the
bank or the office. Lock the screen when leaving the desk, shield the
keypad while typing, and switch on \textbf{two-step verification} (also
called two-factor authentication) wherever it is offered, so that the
password alone is not enough to enter the account.

\begin{definitionbox}{Brute-force vs dictionary attack}
A \textbf{dictionary attack} tries each word in a ready list of common
passwords. A \textbf{brute-force attack} tries every possible character
combination in turn.
\end{definitionbox}

\section{Social engineering}
\textbf{Social engineering} is the practice of tricking a person into
granting access, instead of breaking the software technically. Two
named forms recur. \textbf{Pretexting} invents a false identity and
story --- a bank officer verifying records, an IT support engineer
fixing the account --- to extract a password, OTP, or remote-access
permission. \textbf{Baiting} offers a lure --- a prize, a free recharge,
free software --- that triggers the victim to click a link or download
a file.

The standard signal is manufactured urgency: words such as
\textbf{urgent}, \textbf{blocked}, \textbf{suspended}, and \textbf{verify
now}, combined with a demand to act immediately without checking. The
safe response is fixed: slow down, verify the claim through an
independently known channel (the published helpline or the real
website), and never disclose a password or OTP under pressure.

\begin{definitionbox}{Social engineering}
Manipulating a person, rather than a machine, into revealing
credentials or granting access.
\end{definitionbox}

\section{Phishing}
\textbf{Phishing} is a fake message, email, text, or call that
\textbf{imitates a trusted sender} --- a bank, an office, a delivery
service --- in order to steal passwords, OTPs, card details, or money.
The message always contains a call to action: click the enclosed link,
open the attached file, or call back the given number. The link leads
to a look-alike login page that records whatever is typed; the
attachment installs the payload when opened.

Named channels help classification. Fraudulent text messages are called
\textbf{smishing} (SMS phishing) and fraudulent voice calls are called
\textbf{vishing} (voice phishing), but both are phishing by the same
definition: imitation plus a request for action. The recognition signs
are stable: an urgent or threatening tone, spelling and grammar errors,
a sender address that does not match the claimed organisation, and a
link whose real address differs from the genuine website.

\begin{definitionbox}{Phishing}
A fraudulent message that imitates a trusted sender and asks the victim
to click, open, reply, or pay, in order to steal credentials or money.
\end{definitionbox}

\begin{example}
A clerk receives a mail headed ``Salary slip ready --- download within
24 hours or the account will be blocked.'' The sender shows the office
name, but the reply address is an unknown mailbox and the enclosed
link points to a misspelled portal. The mail is phishing: imitation (the
office name), pressure (blocking within 24 hours), and a call to action
(the link). The safe response is to ignore the link, open the real
portal by typing its known address, and report the message.
\end{example}

\section{Spoofing}
\textbf{Spoofing} means \textbf{faking an identity} so that a message,
call, or website appears to come from someone it does not. Three forms
cover the examination scope. \textbf{Email spoofing} forges the
``From'' address so the mail displays a trusted name that never sent
it. \textbf{Caller ID spoofing} forges the number shown on the phone so
the call appears to come from a bank or office. \textbf{Website (URL)
spoofing} registers a look-alike address --- for instance
\texttt{bank-verlfy.com} in place of the genuine
\texttt{bank-verify.com} --- and hosts a fake login page there.

Spoofing is best understood as a disguise rather than a complete fraud.
A phishing message very often \textbf{uses} spoofing --- a forged
sender address --- to look trustworthy, but spoofing can also appear
without any message at all, as in a forged website waiting for
visitors. The check is the same in every case: confirm the identity
through a known channel instead of trusting what is displayed.

\begin{definitionbox}{Spoofing}
Presenting a forged sender address, caller number, or website address
as though it were genuine.
\end{definitionbox}

\section{Phishing vs spoofing vs hacking}
These three terms form the mandatory near-neighbour set of the lesson,
and examination options deliberately mix them. \textbf{Hacking} is the
broadest: any unauthorised access. \textbf{Phishing} is fraud by
imitation that asks the victim to act: click, open, reply, or pay.
\textbf{Spoofing} is the forged identity itself: the faked address,
number, or site. A single incident can contain all three --- a forged
sender address (spoofing) carries a fake bank mail asking for the OTP
(phishing), and the stolen OTP is then used to enter the account
(hacking) --- but each word answers a different question: was access
taken without permission, was a fake message asking for action, or was
an identity forged?

\section{Malicious links and attachments}
A \textbf{malicious link} is a web address that leads to a fake login
page or starts an unwanted download. Before clicking, hover over the
link with the mouse to reveal the real address, or long-press it on a
phone to preview the destination. Safer still, ignore the link entirely
and type the website's known address by hand. Shortened addresses,
misspelled brand names, and extra words inserted into a familiar
address, arriving inside an unexpected or urgent message, are warning
signs.

A \textbf{malicious attachment} is a file that harms the computer or
steals data once opened. Never open an unexpected attachment, even when
the displayed sender is a known name, because that name may itself be
spoofed; confirm with the sender through a separate channel first. Be
especially cautious with documents that ask the reader to ``enable
macros'' or ``enable content'' on opening, since that step is what
activates many payloads. In short: the link is the address, the
attachment is the file, and both can carry the payload.

\section{Safe habits and response after targeting}
Four habits block the large majority of these attacks. First, verify
every urgent request through a \textbf{known channel} --- the published
helpline or the manually typed website --- before acting. Second, keep
the operating system, browser, and \textbf{antivirus} software updated,
and scan external drives before opening their files. Third, keep a
regular \textbf{backup} of important files so that one incident cannot
destroy the only copy. Fourth, use the mail client's \textbf{report
phishing} option on a suspicious message rather than replying to it or
forwarding it to colleagues.

If a suspicious page has already been opened, do not type any password,
OTP, or card details into it. Close the page, delete the message, and
do not call back any number printed inside it. Change the affected
password by visiting the real website directly, inform the bank or
office through its published helpline, and preserve a record --- a
screenshot and the sender's address --- for the complaint. The order is
fixed: stop, verify on the real channel, change credentials, report.

\section{Exam retrieval and common confusions}
Six distinctions prevent most errors on this topic.
\begin{enumerate}
  \item Phishing \textbf{imitates} a trusted sender to steal data or
    money; spoofing \textbf{forges} the identity it wears; hacking is
    the broader \textbf{unauthorised access}.
  \item OTP stands for \textbf{One-Time Password}. No genuine bank,
    office, or support desk ever asks for it by call or message.
  \item A \textbf{dictionary} attack tries a ready word list; a
    \textbf{brute-force} attack tries every combination.
  \item A \textbf{vulnerability} is the weakness, an \textbf{exploit} is
    the method that uses it, and the \textbf{payload} is the harm
    delivered.
  \item Ethical hacking requires \textbf{written permission}; without
    it, access is unauthorised.
  \item The safe response to any urgent message is fixed: do not click,
    verify through a known channel, and report.
\end{enumerate}

\begin{example}
An item describes a call in which the displayed number matches the
bank's helpline and the caller demands the OTP to ``unblock'' the
account. Trace the definitions: the displayed number is forged
(spoofing), the call imitates the bank and demands the OTP (phishing by
vishing, using social engineering), and any resulting entry into the
account is unauthorised access (hacking). The correct answer names
phishing, and the correct action is to refuse, disconnect, and verify
through the published helpline.
\end{example}

\subsection*{Exercises}
\begin{enumerate}[label=\arabic*.]
  \item Define hacking, and state what separates an ethical security
    test from an attack.
  \item Distinguish a dictionary attack from a brute-force attack.
  \item Define phishing and list four recognition signs of a phishing
    message.
  \item Distinguish phishing from spoofing in one sentence each, and
    explain how they combine.
  \item State the safe response when an unexpected message demands an
    OTP urgently.
\end{enumerate}

\subsection*{Solutions}
\begingroup\small
\begin{enumerate}[label=\arabic*.]
  \item Hacking is gaining unauthorised access to a computer, account,
    or network. An ethical security test is performed with the owner's
    written permission in order to find and fix weaknesses; without
    that permission the same act is an attack.
  \item A dictionary attack tries each entry of a ready list of common
    passwords; a brute-force attack tries every possible character
    combination in turn, so longer passwords resist it.
  \item Phishing is a fraudulent message imitating a trusted sender in
    order to steal credentials or money. Signs: urgent or threatening
    tone, spelling errors, a sender address that does not match the
    claimed organisation, and a link whose real address differs from
    the genuine site.
  \item Phishing is fraud by imitation that asks the victim to act;
    spoofing is the forged identity itself. A phishing message often
    wears a spoofed sender address to look trustworthy.
  \item Do not click, reply, or disclose the OTP; verify through an
    independently known channel (published helpline or manually typed
    website), change credentials on the real site if exposed, and
    report the message.
\end{enumerate}
\endgroup

\lessonheading{44}{Cybersecurity Prevention}
\section{Prevention and the CIA triad}
Cleaning an infected machine is slow and uncertain, so the first goal of
security is prevention: small repeated habits plus a few tools that each do
one distinct job. The standard way to judge whether protection is complete is
the \textbf{CIA triad}: \textbf{confidentiality}, \textbf{integrity}, and
\textbf{availability}.

\begin{definitionbox}{Confidentiality}
Only authorised persons can read the data. Passwords, encryption, and MFA
protect it.
\end{definitionbox}

\begin{definitionbox}{Integrity}
The data is accurate and has not been changed by unauthorised means.
Clean backups, digital signatures, and timely patching protect it.
\end{definitionbox}

\begin{definitionbox}{Availability}
The data and systems are reachable when needed. Backups, antivirus software,
and firewalls protect it.
\end{definitionbox}

\begin{center}
\begin{tabular}{p{0.24\textwidth}p{0.66\textwidth}}
\toprule
\textbf{Term} & \textbf{Meaning in this lesson} \\
\midrule
Antivirus & Software that detects, blocks, and removes malicious programs. \\
Signature & A stored pattern of a known virus, used for recognition. \\
Real-time scanning & Checking files at the moment they are opened, downloaded, or copied. \\
Quarantine & Holding a suspicious file in an isolated area so it cannot run. \\
Firewall & A filter that allows or blocks network traffic by rule. \\
Patch & A small vendor fix for a known security weakness. \\
MFA --- Multi-Factor Authentication & Proving identity with two or more independent factors. \\
Backup & A separate clean copy of data kept for recovery. \\
HTTPS & The secure, encrypted form of web browsing. \\
WPA2/WPA3 & Encryption standards that protect a Wi-Fi link. \\
Autorun & The feature that launches programs automatically when a drive is connected. \\
\bottomrule
\end{tabular}
\end{center}

One attack can break more than one leg of the triad. Stealing a password over
open Wi-Fi breaks confidentiality; altering marks in a result file breaks
integrity; ransomware that locks every document breaks availability. The rest
of this lesson maps each preventive control to the goal it serves.

\section{Antivirus: signatures, scanning, and quarantine}
\textbf{Antivirus} software detects, blocks, and removes malicious programs
(malware). Its oldest method is the \textbf{signature}: a stored pattern of a
known virus, like a fingerprint kept on record. The antivirus compares each
file against its signature database and flags anything that matches. Because
new malware appears every day, the signature database must be updated
regularly; an outdated antivirus misses new threats.

\begin{definitionbox}{Signature}
A stored pattern identifying a known malicious program, against which files
are compared.
\end{definitionbox}

Signatures catch known malware, but they cannot recognise a brand-new
(zero-day) threat. Behaviour-based detection watches what a program
\emph{does} --- for example, trying to encrypt every document --- and stops
it even when no signature matches.

Two scanning modes matter. \textbf{Real-time scanning} checks files at the
moment they are opened, downloaded, or copied, before harm is done. A
\textbf{full (on-demand) scan} checks the whole system when the user asks for
it. When either mode finds something suspicious, the antivirus moves the file
to \textbf{quarantine}: an isolated area where the file cannot run or spread.
Quarantine is not deletion. The file is held safely until the user deletes it
or restores it (for example, when a clean file was flagged by mistake).

\begin{example}
A pen drive is plugged into the office desktop and a file on it matches a
known virus signature. Real-time scanning blocks the file the moment it is
opened and moves it to quarantine, so it cannot run. The user then deletes the
quarantined file. If instead the file had been a clean office template flagged
by mistake, quarantine would allow restoring it --- deletion would not.
\end{example}

\section{Firewall: the network gatekeeper}
A \textbf{firewall} monitors incoming and outgoing network traffic and blocks
whatever its rules forbid. Its job is to stop unauthorised network access: a
stranger on the internet reaching into the office computer, or a malicious
program on the computer sending data out. A firewall can be software running
on the computer or hardware placed at the network boundary; offices commonly
use both.

\begin{definitionbox}{Firewall}
A hardware or software filter that allows or blocks network traffic
according to a set of rules.
\end{definitionbox}

The firewall--antivirus boundary is a favourite exam trap. A firewall blocks
bad \emph{traffic}; it does not remove viruses already saved on the disk.
Antivirus removes bad \emph{files}; it does not, by itself, patrol the network
boundary. The two complement each other and neither replaces the other.

\section{Patching and updates}
A \textbf{patch} is a small fix released by the vendor for a known security
weakness, and a \textbf{software update} delivers such fixes (plus features).
Attackers study published weaknesses, so any delay in installing patches
leaves known doors open. The operating system, the web browser, and the
antivirus itself must all be updated; each has flaws attackers can exploit.
Automatic updates are the safest setting for most users.

\begin{example}
The office desktop offers an operating-system update. Installing it closes a
published weakness the same week it became known. Postponing it for months
means every attacker who read the same notice knows exactly which door is
still open on that machine.
\end{example}

\section{Strong passwords and MFA}
A \textbf{strong password} is long and mixes uppercase letters, lowercase
letters, numbers, and symbols. It must be unique per site, never shared, and
never left at its factory default: routers and office devices ship with
well-known default passwords that must be changed immediately. A
\textbf{password manager} stores one distinct strong password per site so the
user only memorises a single master secret.

\textbf{MFA} stands for Multi-Factor Authentication. It demands two or more
independent proofs of identity: something you know (the password), plus
something you have (a one-time code on the phone) or something you are (a
fingerprint). Even when the password leaks, the attacker still lacks the
second factor, so the account holds. Note the distinction the exam tests: a
longer password is still one factor; MFA is the extra check, not a stronger
password.

\begin{definitionbox}{MFA (Multi-Factor Authentication)}
An authentication method that requires two or more independent factors ---
for example, a password plus a phone code or fingerprint.
\end{definitionbox}

\section{Backup: the availability control}
A \textbf{backup} is a separate clean copy of important data, kept so work
can continue when the original is lost, damaged, or locked by ransomware.
Backup is the great protector of availability. Good backups are regular (so
little work is lost), kept offline or off-site (so one fire, theft, or
ransomware attack cannot destroy both copies), and tested by trial restores. A
backup that cannot be restored is not a backup.

The backup--antivirus boundary is another tested pair. Backup restores
\emph{lost} data; antivirus removes \emph{malicious} data. A machine can have
fresh antivirus and still lose a year's work to a failed disk if there is no
backup.

\section{Safe browsing}
Most infections arrive through the browser, so browsing habits are a control
in their own right. Look for \textbf{HTTPS} and the padlock symbol before
entering passwords or bank details; HTTPS encrypts the session between the
browser and the site. Treat unexpected links and attachments as hostile:
\textbf{phishing} messages imitate trusted senders (a bank, an office, a
colleague) to steal credentials. Download software only from official or
trusted sources, log out on shared computers, and never save passwords on a
machine others use.

\section{Secure Wi-Fi}
The wireless link needs its own protection. Home and office networks must use
\textbf{WPA2 or WPA3} encryption with a strong password, and the router's
default admin password must be changed. Open (password-free) public hotspots
expose traffic, so sensitive logins --- banking, email, office accounts ---
should wait for mobile data, a trusted network, or a \textbf{VPN} (Virtual
Private Network), which encrypts traffic across untrusted links. Two extra
habits help: turn off automatic connection to unknown hotspots, and keep
software updated so known wireless flaws are patched. Note the layering the
exam tests: HTTPS protects the website session, while WPA2/WPA3 protects the
wireless link itself; both layers matter.

\begin{definitionbox}{VPN (Virtual Private Network)}
An encrypted tunnel that carries traffic safely across an untrusted network.
\end{definitionbox}

\section{Removable-media safety}
USB drives and memory cards carry malware between machines. Scan every drive
with antivirus \emph{before} opening its files, and never plug an unknown
found drive into the office computer. Disable \textbf{autorun} so a
newly connected drive cannot launch programs by itself. Handle drives gently
at the end as well: use the safe-eject (safely remove hardware) option before
pulling the device out, so files are not corrupted.

\section{Which control does what}
This table is the highest-yield page in the lesson; most questions are
matching items built from its rows.

\begin{center}
\begin{tabular}{p{0.32\textwidth}p{0.58\textwidth}}
\toprule
\textbf{Control} & \textbf{Main job} \\
\midrule
Antivirus + signatures & Detect and remove known malware files. \\
Real-time scanning + quarantine & Check on access; isolate suspects before harm. \\
Firewall & Block unauthorised network traffic. \\
Patching / updates & Close known software weaknesses. \\
Strong password + MFA & Prove identity; block account theft. \\
Backup & Restore data after loss or ransomware. \\
HTTPS + safe browsing & Protect sessions; dodge phishing and bad downloads. \\
WPA2/WPA3 + VPN & Protect data travelling on the network link. \\
Scan USB + disable autorun & Stop malware entering on removable drives. \\
\bottomrule
\end{tabular}
\end{center}

A useful habit: read the stem's verb. ``Remove'' points to antivirus,
``block traffic'' to the firewall, ``prove identity'' to passwords and MFA,
``restore'' to backup, and ``protect the link'' to WPA2/WPA3 or a VPN.

\section{Exam retrieval and common confusions}
Six distinctions prevent most errors on this topic.
\begin{enumerate}
  \item The firewall blocks \textbf{traffic}; the antivirus removes
    \textbf{files}. Do not swap them.
  \item Quarantine \textbf{isolates}; deletion \textbf{removes}. A
    quarantined file can be restored; a deleted one cannot.
  \item MFA adds a \textbf{second proof of identity}; a longer password is
    still one factor.
  \item Backup restores \textbf{availability}; encryption protects
    \textbf{confidentiality}. A backup is not an antivirus and an antivirus
    is not a backup.
  \item HTTPS protects the \textbf{website session}; WPA2/WPA3 protects the
    \textbf{Wi-Fi link}.
  \item Confidentiality keeps secrets secret, integrity keeps data correct,
    and availability keeps systems reachable. Map each attack to the leg it
    breaks before choosing the control.
\end{enumerate}

\begin{example}
An item asks which control lets the office resume work after ransomware locks
every document. Trace the verb: ``resume work after loss'' means restoring
availability, so the answer is the offline backup. Antivirus (removes malware)
and the firewall (blocks traffic) cannot unlock or recreate the lost files.
\end{example}

\subsection*{Exercises}
\begin{enumerate}[label=\arabic*.]
  \item State the three goals of the CIA triad in one line each.
  \item Explain what a virus signature is and why it must be updated.
  \item Distinguish quarantine from deletion.
  \item Distinguish a firewall from antivirus software.
  \item Explain why a long password alone is not MFA, and what MFA adds.
\end{enumerate}

\subsection*{Solutions}
\begingroup\small
\begin{enumerate}[label=\arabic*.]
  \item Confidentiality: only authorised persons can read the data.
    Integrity: the data is accurate and unaltered by unauthorised means.
    Availability: the data and systems are reachable when needed.
  \item A signature is a stored pattern of a known virus used for recognition.
    New malware appears daily, so the database must be updated or new threats
    will not match anything on record.
  \item Quarantine isolates a suspicious file in a safe area where it cannot
    run, and the file can later be restored; deletion removes the file, and
    it cannot be recovered from quarantine-style holding.
  \item A firewall monitors network traffic and blocks what its rules forbid;
    it does not clean infected files. Antivirus detects and removes malicious
    files on the machine; it does not by itself patrol the network boundary.
  \item A long password is still a single factor (something you know). MFA
    (Multi-Factor Authentication) adds at least one independent factor, such
    as a phone code or fingerprint, so a leaked password alone is not enough.
\end{enumerate}
\endgroup

\lessonheading{56}{Cybersecurity, IT Act and Security Controls}
\section{Why security, and the CIA triad}
A working computer holds records worth protecting: student marks, account
details, personal files, and official correspondence on the office desktop.
\textbf{Cybersecurity} is the practice of protecting computers, networks,
and data against theft, damage, unauthorised access, and misuse. Three
goals organise the whole topic, remembered as the \textbf{CIA triad}:
\textbf{confidentiality} (only authorised persons can read the data),
\textbf{integrity} (the data is not altered by unauthorised means), and
\textbf{availability} (the system stays usable when needed). Every threat
in this lesson attacks one of these goals, and every control defends one.

\begin{definitionbox}{Cybersecurity}
The protection of computer systems, networks, and data against theft,
damage, disruption, and unauthorised access, so that confidentiality,
integrity, and availability are preserved.
\end{definitionbox}

\begin{center}
\begin{tabular}{p{0.24\textwidth}p{0.66\textwidth}}
\toprule
\textbf{Term} & \textbf{Meaning in this lesson} \\
\midrule
Malware & Umbrella term for malicious software written to harm or exploit. \\
Virus & Malware that attaches to a host file and needs human action to spread. \\
Worm & Standalone self-copying malware that spreads over a network unaided. \\
Trojan horse & Malware disguised as useful software; opens a backdoor; does not replicate. \\
Ransomware & Malware that locks or encrypts files and demands payment. \\
Encryption & Coding data so only authorised persons can read it. \\
SSL / TLS & Secure Sockets Layer / Transport Layer Security; secure transmission protocols. \\
Firewall & Barrier filtering network traffic against rules. \\
Quarantine & Isolation of a suspicious file so it cannot run. \\
MFA & Multi-Factor Authentication; two or more identity checks at login. \\
\bottomrule
\end{tabular}
\end{center}

The rest of this lesson uses these terms in one consistent sense.
\textbf{Malware} is always the general term, and virus, worm, Trojan, and
ransomware are its specific types. \textbf{Encryption} always means hiding
content (confidentiality), while \textbf{digital signature} always means
proving the sender and intactness (authenticity and integrity). The exam
tests exactly these pairings, so the wording here is kept fixed.

\section{The Information Technology Act, 2000}
The \textbf{Information Technology Act, 2000} is India's law for the
electronic world. Its scope has three parts: it gives \textbf{legal
validity to electronic records and electronic signatures}, so that online
transactions and electronic governance stand on a legal footing; it enables
\textbf{electronic governance} (official records, notices, and transactions
in electronic form); and it defines \textbf{cyber contraventions and
offences}, such as unauthorised access to a computer, damage to computer
systems and data, and fraud or cheating committed through a computer.

\begin{definitionbox}{Information Technology Act, 2000}
India's legislation of the year 2000 that recognises electronic records
and electronic signatures in law, supports electronic governance and
electronic commerce, and provides for cyber contraventions and offences.
\end{definitionbox}

Two exam points follow. First, the year is \textbf{2000}; later changes to
the law were amendments to the 2000 Act, not a fresh Act, so an option
calling the Act itself ``the IT Act of 2008'' is wrong. Second, the scope
words are \textbf{electronic records, electronic governance, and
cybercrime}: an option that limits the Act to hardware prices, typing
tests, or office furniture has left the scope entirely. No section numbers
are needed for this lesson; the exam fact is the Act's year and its three
branches of coverage.

\section{Malware: the umbrella and its members}
\textbf{Malware}, short for malicious software, is any program written to
harm a computer, steal data, or disrupt work. It is the umbrella term: the
question stem may say ``malware'' when it means the general class, and name
virus, worm, Trojan, or ransomware when it means one member. Two further
nuisance programs belong nearby: \textbf{spyware}, which secretly watches
the user's activity and reports it, and \textbf{adware}, which displays
unwanted advertisements.

\begin{definitionbox}{Malware}
Malicious software: any program designed to damage, disrupt, or gain
unauthorised access to a computer system. Virus, worm, Trojan horse, and
ransomware are types of malware.
\end{definitionbox}

A \textbf{virus} attaches itself to a host file or program and spreads when
that file is run or shared. It needs human action --- opening an infected
attachment, running an infected file, or carrying it on a pen drive --- and
once it runs it can corrupt, delete, or alter files. A \textbf{worm}, by
contrast, is standalone: it copies itself and spreads across a network
without any host file and without waiting for a user to click. Worms
consume bandwidth and system resources, so their signature effect is a
slowed or crashed network. A \textbf{Trojan horse} pretends to be useful
software --- a free utility, a game, a handy tool --- but hides a harmful
purpose inside. It does not replicate by itself; instead it opens a
\textbf{backdoor} through which an attacker can steal passwords, spy on the
user, or control the machine. \textbf{Ransomware} encrypts the victim's
files or locks the screen and demands payment (a ransom) for release.
Paying does not guarantee recovery; the reliable defence is a clean
\textbf{backup} kept separate from the machine.

\begin{example}
An item describes a program that ``looks like a useful download, does not
copy itself, but lets an attacker read passwords.'' Trace the three
contrasts: it needs no host file to be described (so not a virus), it does
not self-spread (so not a worm), and it demands no payment (so not
ransomware). The false-useful disguise plus the backdoor is the Trojan
horse.
\end{example}

\section{Antivirus, quarantine, and patching}
\textbf{Antivirus} software detects, blocks, and removes malware. It works
by recognising known malware signatures and by flagging suspicious
behaviour, and it must be updated regularly because new malware appears
every day. When the antivirus meets a file it cannot yet judge, it places
it in \textbf{quarantine}: an isolated holding area where the file cannot
run or spread while the user or administrator decides its fate. Quarantine
is isolation, not deletion and not a backup --- the file is still present
but harmless while isolated.

\begin{definitionbox}{Quarantine}
The isolation of a suspicious file by security software so that it cannot
execute or spread, pending inspection, removal, or restoration.
\end{definitionbox}

\textbf{Patching} is the companion hygiene step: installing the
vendor-supplied updates that fix known security flaws in the operating
system and applications. An unpatched machine remains open to attacks that
the vendor has already fixed, so ``install updates promptly'' is the exam's
standard correct advice. Antivirus cleans what is already inside,
quarantine holds what is suspect, and patching closes the holes through
which the next attack would enter.

\section{Phishing and anti-phishing}
\textbf{Phishing} is not malware at all but deception: a fake message ---
mail, SMS, or website --- that impersonates a trusted sender (a bank, an
office, a delivery firm) to steal passwords, OTPs, or money. Its warning
signs are urgency and threat (``your account will be blocked''), unexpected
attachments, links whose real destination does not match the claimed
sender, and any request for a password or OTP. The defence, called
\textbf{anti-phishing} practice, is procedural: verify the sender through a
separate channel, hover over or long-press links to inspect them before
clicking, never share an OTP or password by message, and report the message
to the service provider or office.

\begin{definitionbox}{Phishing}
A fraudulent message impersonating a trusted sender, designed to trick the
recipient into revealing passwords, OTPs, or financial details.
\end{definitionbox}

\begin{example}
An item shows a mail reading ``Dear customer, your account is blocked;
click here and enter your OTP to restore it.'' The correct response is to
treat it as phishing: do not click, do not enter the OTP, and verify with
the bank through its known number or site. A genuine bank or office never
asks for a password or OTP by mail or message, so any option advising
compliance is the trap.
\end{example}

\section{Encryption, digital signatures, and SSL/TLS}
\textbf{Encryption} converts readable data (plaintext) into coded form
(ciphertext) so that only authorised persons holding the key can read it.
Its goal is confidentiality. A \textbf{digital signature}, by contrast,
proves who sent a message and that it was not altered in transit; its goals
are authenticity and integrity. The pair must never be swapped: a signature
does not hide the message (anyone can still read it), and encryption by
itself does not prove who sent it. This is the same distinction tested in
the email lesson (TRAP-15), and it returns here because ransomware's
``encryption'' step uses the same idea against the victim.

\textbf{SSL} (Secure Sockets Layer) and its successor \textbf{TLS}
(Transport Layer Security) provide a secure, encrypted channel over a
network. They protect data \textbf{in transit}: mail submission, web
logins, and online forms, signalled by the padlock in the browser address
bar. The lesson's anchor fact, already seen in WO 2026 Q76, is that the
secure email transmission protocol is \textbf{SSL/TLS}. Note the limit:
SSL/TLS protect the journey, not the endpoints --- they do not clean an
infected file and do not replace antivirus.

\begin{definitionbox}{Encryption}
The conversion of plaintext into ciphertext with a key, so that only
authorised holders of the key can recover the original content.
\end{definitionbox}

\section{Firewall, MFA, and the matching map}
A \textbf{firewall} is a barrier, implemented in hardware or software (or
both), that filters incoming and outgoing network traffic against a set of
rules. It blocks unauthorised access while letting permitted traffic pass,
like a gatekeeper at the office network door checking every packet against
the list. Its partner inside the machine is the antivirus: the firewall
filters traffic at the boundary, while the antivirus cleans malware already
inside. One does not replace the other.

\textbf{MFA} (Multi-Factor Authentication) protects logins by demanding two
or more independent checks --- for example, a password plus an OTP, or a
password plus a fingerprint. It closes the door against stolen passwords:
even a leaked password is useless without the second factor. Strong,
distinct passwords remain the first factor; MFA adds factors rather than
replacing them.

\begin{center}
\begin{tabular}{p{0.28\textwidth}p{0.62\textwidth}}
\toprule
\textbf{Control} & \textbf{Job (the stem verb to watch)} \\
\midrule
Firewall & Filters network traffic; blocks unauthorised access. \\
Antivirus & Detects and removes malware on the machine. \\
Quarantine & Isolates a suspicious file so it cannot run. \\
Encryption / SSL-TLS & Hides content; protects data in transit. \\
Digital signature & Proves the sender and that content is intact. \\
Patching & Fixes known security flaws by installing updates. \\
MFA & Adds extra login checks beyond the password. \\
Backup & Recovers data after ransomware, deletion, or failure. \\
\bottomrule
\end{tabular}
\end{center}

The matching technique is mechanical: read the verb in the stem.
``Filters traffic'' points to the firewall; ``cleans malware'' to the
antivirus; ``isolates'' to quarantine; ``hides content'' to encryption;
``proves sender'' to digital signature; ``fixes flaws'' to patching;
``extra login check'' to MFA; ``recovers after loss'' to backup.

\section{Exam retrieval and common confusions}
Seven distinctions prevent most errors on this topic.
\begin{enumerate}
  \item The IT Act is of \textbf{2000} and covers electronic records,
    e-governance, and cybercrime; later changes were amendments.
  \item Malware is the \textbf{umbrella}; virus, worm, Trojan, and
    ransomware are specific types under it.
  \item Virus needs a \textbf{host file plus human action}; worm is
    standalone and \textbf{self-spreading}; Trojan \textbf{does not
    replicate}; ransomware \textbf{locks files and demands payment}.
  \item A digital signature \textbf{proves} (authenticity/integrity); it
    does not \textbf{hide} --- encryption provides confidentiality
    (TRAP-15).
  \item \textbf{SSL/TLS} protect data \textbf{in transit}; they do not
    clean infected files.
  \item A \textbf{firewall} filters \textbf{traffic} at the boundary;
    antivirus cleans \textbf{malware} inside; quarantine
    \textbf{isolates} without deleting.
  \item \textbf{MFA} adds login factors; \textbf{patching} fixes software
    flaws; \textbf{backup} is the recovery answer to ransomware.
\end{enumerate}

\begin{example}
An item asks which control ``proves that a message came from the claimed
sender and was not altered.'' The verb ``proves'' selects the digital
signature, not encryption (which hides) and not the firewall (which
filters). If the stem instead says ``hides the content from
eavesdroppers,'' the same map selects encryption; if it says ``blocks
unauthorised packets,'' it selects the firewall.
\end{example}

\subsection*{Exercises}
\begin{enumerate}[label=\arabic*.]
  \item State the year and the three branches of scope of the Information
    Technology Act covered in this lesson.
  \item Define malware and place virus, worm, Trojan horse, and ransomware
    under it in one sentence each.
  \item Distinguish a virus from a worm on the two points of host file and
    human action.
  \item Why is a Trojan horse dangerous even though it does not replicate?
  \item Explain why backup, rather than payment, is the correct response to
    ransomware.
  \item Contrast encryption with a digital signature: which goal does each
    serve?
  \item Using the matching map, assign one control to each verb: filters,
    isolates, hides, proves, fixes.
\end{enumerate}

\subsection*{Solutions}
\begingroup\small
\begin{enumerate}[label=\arabic*.]
  \item The Information Technology Act is of \textbf{2000}. It gives legal
    validity to electronic records and electronic signatures, enables
    electronic governance, and defines cyber contraventions and offences.
  \item Malware is malicious software written to harm or exploit. A virus
    attaches to a host file and needs human action; a worm self-copies
    across networks unaided; a Trojan disguises itself as useful software
    and opens a backdoor without replicating; ransomware locks or encrypts
    files and demands payment.
  \item A virus needs a host file and a human act (opening, running,
    sharing) to spread; a worm is standalone, copies itself, and spreads
    over the network without a host file or user action.
  \item The Trojan's disguise persuades the user to install it, and the
    hidden backdoor then lets the attacker steal passwords, spy, or
    control the machine --- damage that needs no replication.
  \item Payment does not guarantee recovery and rewards the attacker,
    while a clean separate backup restores the data without negotiation;
    hence backup is the reliable defence.
  \item Encryption serves \textbf{confidentiality} (hides content);
    a digital signature serves \textbf{authenticity and integrity}
    (proves sender and intactness).
  \item Filters --- firewall; isolates --- quarantine; hides ---
    encryption (SSL/TLS in transit); proves --- digital signature;
    fixes --- patching (installing updates).
\end{enumerate}
\endgroup

\lessonheading{64}{Current Cybersecurity, Cryptography and Cyber Law}
\section{Scope and study method}
This lesson covers the current-affairs edge of the official Computer Awareness
scope: computer virus and latest anti-virus, terms and abbreviations used in
IT, and the role of IT in governance. It treats cybersecurity threats,
cryptography basics, firewall types, and cyber law at awareness level. The
exam depth is recall and classification, with at most a one-step scenario.
The study method for every term is the near-neighbour contrast: what it is,
what it is not, and the nearest confusion listed alongside it.

\begin{center}
\begin{tabular}{p{0.24\textwidth}p{0.66\textwidth}}
\toprule
\textbf{Term} & \textbf{Meaning in this lesson} \\
\midrule
Malware & Hostile software, including virus, worm, Trojan horse, ransomware, spyware. \\
Parasitic / file-infector virus & A virus that attaches itself to an executable file. \\
Phishing & Fraudulent messages that trick the user into revealing credentials. \\
Ransomware & Malicious code that locks or encrypts files and demands payment. \\
DDoS & Distributed Denial of Service: flood of traffic from many machines. \\
Hacking & Unauthorised access to a computer or account. \\
DES --- Data Encryption Standard & Symmetric cipher working on 64-bit blocks. \\
Encryption & Conversion of plaintext to ciphertext for confidentiality. \\
Digital signature & Authenticity and integrity proof; not confidentiality. \\
Firewall & Barrier between a trusted and an untrusted network. \\
Personal data & Data identifying a person, protected by consent and purpose limits. \\
Copyright & Owner's exclusive rights over an original work. \\
\bottomrule
\end{tabular}
\end{center}

\section{Malware and the parasitic virus}
\textbf{Malware} is the general name for hostile software. A \textbf{virus}
is malware that needs a host program and some human action, such as opening
an attachment, to spread. A \textbf{worm} spreads on its own across a
network without a host program. A \textbf{Trojan horse} looks like useful
software but hides a harmful payload, and it does not self-replicate.

\begin{definitionbox}{Parasitic (file-infector) virus}
A virus that attaches itself to an executable file. When the infected program
runs, the virus code runs first and then hands control back to the host
program, spreading further as infected files are shared or copied.
\end{definitionbox}

The parasitic virus is the file-infecting member of the family: it lives
inside executable files, activates when the host program runs, and travels
with shared files. An \textbf{anti-virus} detects, blocks, and removes such
infections, and its database must be kept updated against the latest strains.

\begin{example}
A student copies a free game from a friend's pen drive and runs it on the
office desktop. The game file carries a file infector: the virus runs first,
attaches to other executable files, and spreads when those files are shared
further. Updating the anti-virus and scanning the drive before opening the
files is the correct defence.
\end{example}

\section{Phishing, ransomware, DDoS, and hacking}
These four threats differ in mechanism, and the exam swaps their definitions.
\textbf{Phishing} attacks the user: a fraudulent mail or message, often styled
as a bank or KYC notice, tricks the user into revealing a password or OTP on
a fake login page. \textbf{Ransomware} attacks the data: malicious code
locks or encrypts the victim's files and demands payment for the decryption
key. \textbf{DDoS} (Distributed Denial of Service) attacks availability: a
flood of traffic from many machines overwhelms a server so genuine users
cannot reach it. A flood from a single machine is only DoS; the distributed
prefix matters. \textbf{Hacking} is the broad term for unauthorised access to
a computer or account; phishing, ransomware, and DDoS are specific methods or
goals within that broad idea.

\begin{definitionbox}{Ransomware}
Malicious code that locks or encrypts the victim's files and demands payment
for their release.
\end{definitionbox}

\begin{example}
An item describes a clerk whose files are suddenly unreadable, with a message
demanding payment for the key. Tracing the mechanism --- files locked,
payment demanded --- gives ransomware. If the stem instead described a fake
bank mail asking for an OTP, the same trace gives phishing; a server slowed to
a halt under flood traffic gives DDoS.
\end{example}

\section{Cryptography and DES}
\textbf{Cryptography} protects information by transforming it so only the
intended recipient can read it. \textbf{Plaintext} is the readable message,
\textbf{ciphertext} is the transformed message, and a \textbf{key} controls
the transformation. \textbf{Encryption} converts plaintext to ciphertext and
\textbf{decryption} converts it back. In \textbf{symmetric} cryptography the
same key encrypts and decrypts; in \textbf{asymmetric} cryptography a public
key and a private key form a pair.

\begin{definitionbox}{DES (Data Encryption Standard)}
A symmetric cipher that processes data in 64-bit blocks. Its key is 64 bits
long, of which 56 bits are effective, the remaining 8 bits being parity.
\end{definitionbox}

The two numbers for DES must be kept apart. The \textbf{block size} is 64
bits: that is the chunk of data processed at a time. The \textbf{effective
key} is 56 bits. An option offering 128 bits confuses DES with a later
cipher's block or key size.

\section{Encryption vs digital signature}
This is the TRAP-15 distinction and it recurs in papers. \textbf{Encryption}
provides \textbf{confidentiality}: only the intended recipient can read the
message. A \textbf{digital signature} provides \textbf{authenticity and
integrity}: it proves who sent the message and that it was not altered in
transit. A signature does not hide the message content, and encryption by
itself does not prove the sender's identity.

\begin{definitionbox}{Digital signature}
A cryptographic proof attached to a message that establishes the sender's
authenticity and the message's integrity, without by itself providing
confidentiality.
\end{definitionbox}

\begin{example}
An item states, ``A digital signature encrypts the message body so only the
recipient can read it.'' This is wrong. Hiding the body is encryption
(confidentiality). The signature answers two different questions: who sent it
and whether it was changed. The wrong option has merged the two security
verbs.
\end{example}

\section{Firewalls: three types}
A \textbf{firewall} is a barrier between a trusted internal network and an
untrusted external network. It permits or denies traffic according to a rule
set. It filters packets and connections; it does not disinfect
file-infecting viruses inside downloads that its rules allow through --- that
is the anti-virus's job.

A \textbf{packet-filtering} firewall examines each packet's headers ---
source and destination address, port, and protocol --- and decides packet by
packet. It is the fastest and simplest type but cannot judge a session as a
whole. A \textbf{stateful-inspection} firewall examines the same headers plus
the \textbf{connection state}: it remembers open sessions and drops packets
that do not belong to a known session. An \textbf{application (proxy)}
firewall sits between the client and the server as a go-between and inspects
application-layer content and commands, such as web or mail requests.

\begin{definitionbox}{Firewall}
A network barrier that permits or denies traffic between a trusted internal
network and an untrusted external network according to a rule set.
\end{definitionbox}

The ordering shorthand is: header only, then header plus state, then full
content via proxy. Filtering checks the envelope, stateful inspection checks
the conversation, and the proxy reads the letter inside.

\section{Internet governance, privacy, and data protection}
The Internet has no single owner; names and numbers are coordinated
internationally, for example by ICANN (Internet Corporation for Assigned
Names and Numbers). Governance concerns at awareness level are four: the
\textbf{security} of networks, the \textbf{privacy} of personal data, fair
\textbf{access} to the network, and which country's \textbf{law} applies when
an act crosses borders. The positive side is the \textbf{role of IT in
governance} (e-governance): online services, transparency, and delivery of
benefits to citizens.

\textbf{Personal data} is any data that identifies a person, such as a name
combined with a phone number, biometrics, or account details. Protection
principles at awareness level are: collect data with \textbf{consent}, use it
only for the stated \textbf{purpose}, keep it \textbf{secure}, and retain it
no longer than needed. Privacy settings, strong passwords, and OTP discipline
are the user-side defences. Misuse or leakage of personal data can attract
legal consequences under cyber law.

\section{Cyber law and copyright: awareness-level map}
Cyber law at this level means two things: unauthorised access, fraud, and
data misuse attract punishment, and creators hold enforceable rights over
their works. The \textbf{Copyright Act} protects original literary, artistic,
musical, and software works from copying without permission. The
awareness-level map matches each concept to its area: the \textbf{meaning of
copyright} (what exclusive rights the owner holds), \textbf{first ownership}
(who owns the work first), \textbf{infringement} (what counts as violation),
\textbf{fair dealing} (limited permitted uses such as private study), and
\textbf{offences} (punishment for wilful infringement). Software piracy ---
copying licensed software without permission --- is the standard infringement
example. Section numbers are recall only where the question itself supplies
them; otherwise the concept match is what the exam tests.

\section{Exam retrieval and common confusions}
Seven distinctions prevent most errors on this topic.
\begin{enumerate}
  \item DES block size is \textbf{64 bits}; the effective key is 56 bits. Do
    not offer 128 for either (DES fact).
  \item Digital signature $\ne$ encryption: authenticity/integrity vs
    confidentiality (TRAP-15/TRM-11).
  \item A file infector attaches to \textbf{executable files}; it is not a
    worm (self-spreading) or a Trojan (non-replicating disguise).
  \item Packet filter checks headers only; stateful adds \textbf{connection
    state}; proxy inspects \textbf{application content}.
  \item Phishing steals \textbf{credentials}; ransomware locks \textbf{files};
    DDoS blocks \textbf{access}; hacking is unauthorised access generally.
  \item DDoS needs \textbf{many} machines (distributed); a single-machine
    flood is only DoS.
  \item Copyright protects the owner against \textbf{infringement}, with
    limited \textbf{fair-dealing} exceptions such as study --- piracy is the
    standard violation example.
\end{enumerate}

\begin{example}
An item asks which firewall remembers open sessions and drops stray packets.
Trace the triple: the packet filter decides per packet, the proxy inspects
application content, and only stateful inspection tracks connection state.
The answer is the stateful-inspection firewall. If the stem instead asked
which device disinfects an infected executable, the answer leaves firewalls
entirely and gives the anti-virus.
\end{example}

\subsection*{Exercises}
\begin{enumerate}[label=\arabic*.]
  \item Define a parasitic (file-infector) virus and distinguish it from a worm and a Trojan horse.
  \item State the DES block size and effective key size, and expand the abbreviation DES.
  \item Distinguish packet-filtering, stateful-inspection, and application/proxy firewalls.
  \item Distinguish phishing, ransomware, DDoS, and hacking in one line each.
  \item Explain why a digital signature is not encryption, using TRAP-15.
  \item List the four Internet-governance concerns and the four data-protection principles from this lesson.
  \item Match each copyright concept (meaning, first ownership, infringement, fair dealing, offences) to what it covers.
\end{enumerate}

\subsection*{Solutions}
\begingroup\small
\begin{enumerate}[label=\arabic*.]
  \item A parasitic virus attaches to an executable file, runs when the host
    program runs, and spreads with shared files. A worm spreads on its own
    across a network without a host; a Trojan horse disguises itself as useful
    software and does not self-replicate.
  \item DES (Data Encryption Standard) is a symmetric cipher with a 64-bit
    block size and a 56-bit effective key (64 bits with parity).
  \item Packet-filtering checks per-packet headers only; stateful inspection
    adds connection-state tracking; the application/proxy firewall inspects
    application-layer content as a go-between.
  \item Phishing tricks the user into revealing credentials; ransomware locks
    or encrypts files for ransom; DDoS floods a service from many machines so
    genuine users cannot reach it; hacking is unauthorised access generally.
  \item Encryption provides confidentiality (who can read it); a digital
    signature provides authenticity and integrity (who sent it, whether it was
    changed). Merging them is TRAP-15.
  \item Governance concerns: security, privacy, access, jurisdiction
    (applicable law). Data-protection principles: consent, stated purpose,
    security, limited retention.
  \item Meaning: the exclusive rights the owner holds; first ownership: who
    owns the work first; infringement: what counts as violation (for example
    software piracy); fair dealing: limited permitted uses such as private
    study; offences: punishment for wilful infringement.
\end{enumerate}
\endgroup

\unitopener{11}{Governance \& Integration}
\lessonheading{47}{Humanware and IT Governance}
\section{What humanware means}
A computer system is not only metal and programs. It also needs people. The
human element of a computer system is called \textbf{humanware}, and it is
also known as \textbf{liveware} or \textbf{peopleware}. Humanware means the
people who design, build, operate, manage and use computers. It is
\emph{neither hardware nor software}; it is the people component of the
system.

\begin{definitionbox}{Humanware}
The people involved with a computer system --- its users, operators,
administrators, programmers and support staff. Also called liveware or
peopleware.
\end{definitionbox}

\begin{center}
\begin{tabular}{p{0.24\textwidth}p{0.66\textwidth}}
\toprule
\textbf{Term} & \textbf{Meaning in this lesson} \\
\midrule
Hardware & The physical parts of the computer. \\
Software & The programs and instructions the computer runs. \\
Humanware & The people who use and manage the hardware and software. \\
User & The person who uses the finished system to do a task. \\
Administrator & The person who manages the system, network or database. \\
\bottomrule
\end{tabular}
\end{center}

The rest of this lesson uses these terms in one consistent sense.
\textbf{Hardware} is the machinery. \textbf{Software} is the set of
instructions. \textbf{Humanware} is the people.

\section{The people in a computer system}
Humanware is not a single role. Many different people take part, and each has
a distinct job. The main categories are the following.

\begin{itemize}
  \item \textbf{User (or end user)} --- the person who uses the finished
    system to perform a task. A clerk entering data, a student reading a
    portal and a customer paying a bill are all users.
  \item \textbf{Operator} --- the person who runs the computer system, feeds
    jobs into it and monitors routine operation.
  \item \textbf{Administrator} --- the person who manages the system. A
    \textbf{system administrator} manages servers, user accounts and access
    rights; a \textbf{network administrator} manages the network; a
    \textbf{database administrator} manages the database.
  \item \textbf{Programmer (developer)} --- the person who writes, tests and
    maintains the software.
  \item \textbf{IT support (help desk) and technician} --- the people who
    answer user problems, troubleshoot faults, and install, repair and
    maintain the hardware.
\end{itemize}

The distinction to hold is simple. Administrators \emph{manage} the system,
operators \emph{run} it, and users \emph{use} its output.

\section{Hardware, software and humanware}
A working computer system needs all three components. Hardware is the
physical machine. Software is the set of programs that tell the hardware what
to do. Humanware is the people who design, run and use both. Hardware without
software does nothing, and both are useless without the people who operate
them. This is why humanware is taught as a full component, not as an
afterthought.

\begin{example}
A government office installs new computers (hardware) and an accounting
program (software). The system still does not work by itself. A programmer
must maintain the program, a system administrator must create user accounts,
an operator must run the daily jobs, and the office staff (users) must enter
the data. Those people are the humanware.
\end{example}

\section{The role of IT in governance}
\textbf{Governance} is the process of decision-making and its implementation
in an organisation or a state. \textbf{IT} (Information Technology) and
\textbf{ICT} (Information and Communication Technology) let government carry
out that work over computers and networks. The role of IT in governance is to
make government faster, cheaper, more transparent and more accessible to the
people it serves.

The main benefits of using IT in governance are the following.

\begin{itemize}
  \item \textbf{Transparency} --- citizens can see the rules, the status of
    their applications and the decisions taken.
  \item \textbf{Accountability} --- actions are recorded and can be traced.
  \item \textbf{Efficiency} --- less paperwork and faster processing.
  \item \textbf{Convenience} --- services are available from home and around
    the clock.
  \item \textbf{Lower cost} and \textbf{reduced corruption}, because manual
    handling and physical counters are replaced by digital processes.
\end{itemize}

The overall aim is \textbf{citizen-centric} service: the citizen is placed at
the centre of every process.

\section{IT governance versus e-governance}
These two terms are a common near-neighbour pair, and the difference matters.

\begin{itemize}
  \item \textbf{IT governance} is the process by which an organisation
    directs and controls its own IT resources so that they support its goals.
    It is \emph{governance of IT}.
  \item \textbf{E-governance} is the use of IT by government to govern and
    deliver public services. It is \emph{IT in governance}.
\end{itemize}

The trap is to swap them. IT governance is about managing IT well inside an
organisation; e-governance is about using IT to govern the state and serve
citizens.

\section{E-governance basics}
\textbf{E-governance} is the application of ICT by government to deliver
services, share information and transact with citizens, businesses and other
government agencies. It uses the internet, databases and digital records in
place of paper and physical counters. It makes government services faster,
wider-reaching and easier to audit.

\begin{definitionbox}{E-governance}
The use of information and communication technology to improve the way
government serves and interacts with the public.
\end{definitionbox}

E-governance is usually described by the parties who interact. The four
standard types are the following.

\begin{center}
\begin{tabular}{p{0.14\textwidth}p{0.32\textwidth}p{0.44\textwidth}}
\toprule
\textbf{Type} & \textbf{Full form} & \textbf{Interaction} \\
\midrule
G2C & Government to Citizen & Services to the general public. \\
G2G & Government to Government & Between government departments. \\
G2B & Government to Business & With companies and traders. \\
G2E & Government to Employee & With government staff. \\
\bottomrule
\end{tabular}
\end{center}

This lesson keeps the treatment at awareness level. Named platforms and
detailed G2x matching belong to Lesson 57.

\section{Digital public services}
A \textbf{digital public service} is a government service that a citizen can
obtain online. Common examples are certificates, licences, land records, bill
payment, applications and grievance redressal. Citizens reach these services
through service portals and through \textbf{Common Service Centres (CSCs)},
which provide a physical access point where help is needed. Typical features
are online application, status tracking, digital payment and downloadable
documents. A service counts as digital when the citizen can complete the task
without visiting an office in person.

\section{Electronic records}
An \textbf{electronic record} is information created, stored and retrieved in
electronic (digital) form. Examples include e-documents, e-files, scanned
papers, database entries and email. Electronic records are easy to store,
search, retrieve, share and back up, and they replace large volumes of paper.
An electronic record can have legal validity when it is authenticated, for
example with a digital signature.

\begin{definitionbox}{Electronic record}
Information that is created, stored and retrieved in electronic form, such as
an e-document, an e-file or a database entry.
\end{definitionbox}

\section{Authentication}
\textbf{Authentication} is the process of verifying that a user, device or
system is who or what it claims to be. It answers the question: ``Who are
you?'' Common methods are a password, a PIN, an \textbf{OTP} (One-Time
Password), a biometric such as a fingerprint or iris, and a digital
certificate. Authentication must succeed before the system grants access.

\begin{definitionbox}{Authentication}
The process of confirming a claimed identity before access to a system is
allowed.
\end{definitionbox}

Authentication is often confused with \textbf{authorization}.
\textbf{Authorization} is the process of granting or denying permission ---
it decides what a verified identity is allowed to do. Authentication proves
identity; authorization grants or refuses access rights. Authentication
always comes first.

Authentication relies on \emph{factors}. There are three standard types:
something you \textbf{know} (a password or PIN), something you \textbf{have}
(an OTP on your phone, a token or a smart card), and something you \textbf{are}
(a fingerprint, iris or face). When two different factors are combined, it is
called \textbf{two-factor authentication (2FA)}; when two or more are used, it
is \textbf{multi-factor authentication (MFA)}. Two passwords are not two
factors, because they are the same factor used twice.

\begin{example}
A candidate logs in to a government portal with a password and is then asked
for an OTP sent to the registered mobile number. The password is something the
candidate knows; the OTP is something the candidate has. Together they form
two-factor authentication. The portal then decides which forms the candidate
may open --- that decision is authorization, not authentication.
\end{example}

\section{Exam retrieval and common confusions}
Five distinctions prevent most errors on this topic.
\begin{enumerate}
  \item Humanware means \textbf{people}. It is neither hardware nor software;
    do not confuse it with firmware (which is software stored in hardware).
  \item \textbf{IT governance} is governance \emph{of} IT;
    \textbf{e-governance} is IT \emph{in} governance. Do not swap them.
  \item \textbf{G2C} is Government to \textbf{Citizen}, not Government to
    Company. The business interaction is G2B.
  \item \textbf{Authentication} proves identity; \textbf{authorization}
    grants permission. Authentication comes first.
  \item An \textbf{electronic record} is digital information. A printed paper
    file is a physical record, not an electronic one.
\end{enumerate}

\begin{example}
An item asks which component of a computer system consists of the people who
operate and manage it. Hardware is the machinery and software is the set of
instructions, so both are wrong. The correct answer is humanware. If the stem
instead asks what improves transparency and accountability in government, the
answer is the role of IT in governance, and if it asks what verifies a user's
identity, the answer is authentication.
\end{example}

\subsection*{Exercises}
\begin{enumerate}[label=\arabic*.]
  \item Define humanware and name four categories of people who form it.
  \item Distinguish hardware, software and humanware, and explain why all
    three are needed.
  \item Give three benefits of using IT in governance.
  \item Distinguish IT governance from e-governance.
  \item What are the four types of e-governance interaction, and what does
    each connect?
  \item Define an electronic record and give two examples.
  \item Distinguish authentication from authorization, and name the three
    standard authentication factors.
\end{enumerate}

\subsection*{Solutions}
\begingroup\small
\begin{enumerate}[label=\arabic*.]
  \item Humanware (also called liveware or peopleware) is the human element of
    a computer system. Four categories are users (end users), operators,
    administrators, and programmers (with IT support and technicians).
  \item Hardware is the physical machine, software is the set of programs,
    and humanware is the people. Hardware without software does nothing, and
    both need people to design, run and use them, so all three are required.
  \item Any three of: transparency, accountability, efficiency, convenience,
    lower cost, reduced corruption, and citizen-centric service.
  \item IT governance is the direction and control of an organisation's own
    IT resources --- governance of IT. E-governance is the use of IT by
    government to deliver public services --- IT in governance.
  \item G2C (Government to Citizen) serves the public; G2G (Government to
    Government) links departments; G2B (Government to Business) links
    companies; G2E (Government to Employee) links government staff.
  \item An electronic record is information created, stored and retrieved in
    electronic form. Examples include an e-document, an e-file, a scanned
    paper or a database entry.
  \item Authentication verifies identity (``who are you?''); authorization
    grants or denies permission (``what may you do?''). The three factors are
    something you know, something you have, and something you are.
\end{enumerate}
\endgroup

\lessonheading{57}{Digital Platforms and E-Governance}
\section{What e-Governance means}
A running government delivers many services: certificates, licences,
payments, benefits, and information. \textbf{e-Governance} is the use of
Information and Communication Technology (ICT) by government to deliver
these services. Instead of standing at a physical counter, a citizen reaches
the service through an online portal. The aim is to make services
accessible, efficient, transparent, and reliable.

e-Governance is about governance, not merely about technology. The
technology is the tool; the purpose is better delivery of public services.

\begin{definitionbox}{e-Governance}
The use of Information and Communication Technology (ICT) by government to
deliver services to citizens, businesses, and other parts of government in
an accessible, efficient, transparent, and reliable way.
\end{definitionbox}

\begin{center}
\begin{tabular}{p{0.22\textwidth}p{0.68\textwidth}}
\toprule
\textbf{Term} & \textbf{Meaning in this lesson} \\
\midrule
e-Governance & ICT-based delivery of government services. \\
G2x & The interaction categories: G2C, G2G, G2B, G2E. \\
Digital India & The flagship Government of India programme launched on 1 July 2015. \\
NeGP & National e-Governance Plan, approved in May 2006. \\
e-Kranti & Pillar 5 of Digital India and the upgraded NeGP (NeGP 2.0). \\
NeGD & National e-Governance Division, created in 2009 under MeitY. \\
e-NAM & National Agriculture Market, the pan-India agricultural trading portal. \\
\bottomrule
\end{tabular}
\end{center}

\section{The G2x classification}
e-Governance interactions are named by who takes part. The ``G'' always
stands for Government, and the letter after it says who the government is
interacting with.

\begin{center}
\begin{tabular}{p{0.12\textwidth}p{0.30\textwidth}p{0.48\textwidth}}
\toprule
\textbf{Type} & \textbf{Full form} & \textbf{Who is served / examples} \\
\midrule
G2C & Government to Citizen & The general public. UMANG, DigiLocker, certificates, utility bills. \\
G2G & Government to Government & Government departments and levels. e-Office, PFMS. \\
G2B & Government to Business & Businesses and companies. GeM, GSTN. \\
G2E & Government to Employee & Government employees. An employee service book or salary portal. \\
\bottomrule
\end{tabular}
\end{center}

The classification is the highest-yield part of this topic. An item can
name a platform and ask which G2x category it belongs to, or name a category
and ask which platform fits. The safe method is to ask ``who is the
government serving?'' A citizen means G2C; another government body means
G2G; a company means G2B; an employee means G2E.

\begin{example}
The stem asks, ``GeM (Government e-Marketplace) is an example of which
category?'' GeM is used by government to buy goods and services from
businesses, so the government is serving a business. The answer is G2B, not
G2C. A student who sees only that ``government uses it'' may wrongly choose
G2C; the deciding question is who is on the other side.
\end{example}

\section{Digital India}
\textbf{Digital India} is a flagship programme of the Government of India.
It was launched on \textbf{1 July 2015} by the Prime Minister, with the
vision to transform India into a \textbf{digitally empowered society and
knowledge economy}.

The programme rests on \textbf{three vision areas}:
\begin{enumerate}[label=(\roman*)]
  \item Digital Infrastructure as a Core Utility to Every Citizen;
  \item Governance and Services on Demand;
  \item Digital Empowerment of Citizens.
\end{enumerate}

Under these vision areas the programme is built on \textbf{nine pillars}:
Broadband Highways; Universal Access to Mobile Connectivity; Public Internet
Access Programme; e-Governance: Reforming Government through Technology;
e-Kranti: Electronic Delivery of Services; Information for All; Electronics
Manufacturing; IT for Jobs; and Early Harvest Programmes. Pillar 5,
e-Kranti, is the service-delivery pillar and is treated separately below.

\begin{definitionbox}{Digital India}
A flagship Government of India programme, launched on 1 July 2015, with the
vision of transforming India into a digitally empowered society and
knowledge economy, built on three vision areas and nine pillars.
\end{definitionbox}

\section{NeGP, e-Kranti, and NeGD: three names, three things}
These three names are close and are a standing source of confusion. They are
kept apart here by asking what each one \emph{is}.

The \textbf{National e-Governance Plan (NeGP)} is the older umbrella
\emph{plan}. It was approved in \textbf{May 2006}, with the vision to make
all government services accessible to the common man in his locality,
through common service delivery outlets, at affordable cost. It is built
around \textbf{Mission Mode Projects (MMPs)} --- 27 MMPs with supporting
components.

\textbf{e-Kranti} means ``electronic delivery of services''. It is
\textbf{Pillar 5} of the Digital India programme and is also called
\textbf{NeGP 2.0}, an upgraded approach to the National e-Governance Plan.
Its vision is ``Transforming e-Governance for Transforming Governance'',
and its mission is to deliver all government services electronically through
integrated and interoperable systems.

The \textbf{National e-Governance Division (NeGD)} is not a plan at all. It
is an \emph{organisation}. It was created in \textbf{2009} by the Ministry
of Electronics and Information Technology (MeitY) as an independent business
division under the Digital India Corporation. Its role is to support MeitY
in the programme management and implementation of e-Governance projects, and
to give technical and advisory support to ministries and departments.

\begin{center}
\begin{tabular}{p{0.12\textwidth}p{0.40\textwidth}p{0.10\textwidth}p{0.22\textwidth}}
\toprule
\textbf{Name} & \textbf{What it is} & \textbf{Year} & \textbf{One word} \\
\midrule
NeGP & Umbrella e-Governance plan with 27 MMPs & 2006 & The plan \\
e-Kranti & Pillar 5 of Digital India; NeGP 2.0 & 2015 & The approach \\
NeGD & Division under MeitY supporting projects & 2009 & The body \\
\bottomrule
\end{tabular}
\end{center}

\begin{example}
An item states, ``NeGD is the National e-Governance Plan approved in 2006.''
This is wrong on two counts. The plan approved in 2006 is \textbf{NeGP}
(National e-Governance Plan). \textbf{NeGD} is the National e-Governance
\emph{Division}, an organisation created in 2009. The distractor swaps a
plan for a body.
\end{example}

\section{e-NAM and agricultural marketing}
\textbf{e-NAM} stands for \textbf{National Agriculture Market} (also written
eNAM). It is a pan-India \textbf{electronic trading portal} for agricultural
commodities, launched on \textbf{14 April 2016}. It is implemented by the
\textbf{Small Farmers Agribusiness Consortium (SFAC)} under the Ministry of
Agriculture and Farmers Welfare. The portal is available at
\textbf{enam.gov.in}.

The key idea is that e-NAM \textbf{networks the existing Agriculture Produce
Market Committee (APMC) mandis} to create a \textbf{unified national market}
for agricultural commodities. Because trade happens online, farmers get
\textbf{online competitive bidding and transparent price discovery}, and
direct sale to buyers reduces the role of intermediaries.

\begin{definitionbox}{e-NAM}
The National Agriculture Market, a pan-India electronic trading portal,
launched on 14 April 2016, that networks the existing APMC mandis to create
a unified national market for agricultural commodities.
\end{definitionbox}

A near-neighbour to keep separate is \textbf{AGMARKNET}. AGMARKNET is a
portal that provides agricultural marketing \emph{information} such as
market prices. e-NAM, by contrast, is the \emph{trading} platform where the
sale itself takes place. An item that asks for the agricultural marketing
\emph{trading} platform wants e-NAM.

\section{Citizen portal, electronic records, and authentication}
The \textbf{National Portal of India} is \textbf{india.gov.in}. It is a
\textbf{single-window} gateway from which a citizen can reach the
information and services of the Government of India. A related platform,
\textbf{MyGov}, is the citizen-engagement platform for participative
governance.

Electronic governance depends on giving electronic documents the same
standing as paper. The \textbf{Information Technology Act, 2000} gives legal
recognition to \textbf{electronic records}. An electronic record can be
stored, retrieved, and used in electronic form. To be trustworthy, it must
be authenticated.

\textbf{Authentication} proves that a message or document is genuine and
comes from the claimed sender. A \textbf{digital signature} provides
authenticity, integrity, and non-repudiation. \textbf{Encryption}, by
contrast, provides confidentiality: it keeps the content secret from
unauthorised readers. \textbf{Aadhaar} is used as a digital identity for
authentication of a person.

\begin{definitionbox}{Digital signature}
A mechanism used to authenticate an electronic record. It provides
authenticity, integrity, and non-repudiation, and is distinct from
encryption, which provides confidentiality.
\end{definitionbox}

\begin{example}
An item states, ``A digital signature encrypts the document so that only the
recipient can read it.'' This confuses two ideas. Encrypting for secrecy is
\textbf{encryption}, which provides confidentiality. A \textbf{digital
signature} proves who sent the document and that it was not altered; it
provides authenticity and integrity. The wrong option has used the
signature's name for encryption's job.
\end{example}

\section{Exam retrieval and common confusions}
Six distinctions prevent most errors on this topic.
\begin{enumerate}
  \item e-NAM = agricultural marketing and trading; it is \textbf{not}
    general e-commerce.
  \item GeM = G2B procurement, \textbf{not} a G2C service. e-Office and PFMS
    are G2G.
  \item NeGP (plan, 2006), e-Kranti (pillar / NeGP 2.0, 2015), and NeGD
    (division, 2009) are three different things; do not swap them.
  \item Digital signature = authenticity and integrity; encryption =
    confidentiality. They are \textbf{not} the same (TRAP-15).
  \item The citizen portal of the Government of India is india.gov.in, not
    the Digital India or NeGD site.
  \item Digital India is built on three vision areas and nine pillars;
    e-Kranti is Pillar 5, the service-delivery pillar.
\end{enumerate}

\begin{example}
An item asks which body supports MeitY in the programme management of
e-Governance projects. Trace the three names: NeGP is a plan, e-Kranti is a
pillar, and only NeGD is an organisation. The answer is NeGD.
\end{example}

\subsection*{Exercises}
\begin{enumerate}[label=\arabic*.]
  \item Expand e-NAM and state its purpose in one line.
  \item Distinguish G2C, G2G, G2B, and G2E, giving one example of each.
  \item State the launch date and vision of the Digital India programme.
  \item Explain the difference between NeGP, e-Kranti, and NeGD.
  \item Define authentication and explain how a digital signature differs
    from encryption.
\end{enumerate}

\subsection*{Solutions}
\begingroup\small
\begin{enumerate}[label=\arabic*.]
  \item e-NAM is the National Agriculture Market: a pan-India electronic
    trading portal that networks the existing APMC mandis to create a
    unified national market for agricultural commodities.
  \item G2C (Government to Citizen) serves the public, e.g.\ UMANG; G2G
    (Government to Government) serves government bodies, e.g.\ e-Office;
    G2B (Government to Business) serves businesses, e.g.\ GeM; G2E
    (Government to Employee) serves government employees, e.g.\ an employee
    service book.
  \item Digital India was launched on 1 July 2015, with the vision to
    transform India into a digitally empowered society and knowledge
    economy.
  \item NeGP is the National e-Governance Plan, the umbrella plan approved
    in May 2006 with 27 MMPs. e-Kranti is Pillar 5 of Digital India and the
    upgraded NeGP (NeGP 2.0). NeGD is the National e-Governance Division, an
    organisation created in 2009 under MeitY to support e-Governance
    projects.
  \item Authentication proves that a message or document is genuine and
    comes from the claimed sender. A digital signature provides authenticity,
    integrity, and non-repudiation; encryption provides confidentiality by
    keeping the content secret. They are different functions.
\end{enumerate}
\endgroup


\end{document}
